\documentclass[11pt]{article}

\usepackage[T1]{fontenc}
\usepackage[utf8]{inputenc}
\usepackage{amsmath,amssymb,amsthm,mathtools}
\usepackage{mathrsfs}
\usepackage{booktabs,longtable,array}
\usepackage{enumitem}
\usepackage[margin=1in]{geometry}
\usepackage{xcolor}
\usepackage[hidelinks]{hyperref}

\hypersetup{
  pdftitle={Renewal Navier--Stokes Stability Research Notes V100},
  pdfauthor={Research continuation notes}
}

\newcommand{\Torus}{\mathbb T}
\newcommand{\R}{\mathbb R}
\newcommand{\dd}{\,\mathrm d}
\newcommand{\PP}{\mathbb P}
\newcommand{\BNS}{\mathcal B}
\newcommand{\PROVED}{\textcolor{green!45!black}{\textsf{PROVED}}}
\newcommand{\IMPORTED}{\textcolor{blue!60!black}{\textsf{IMPORTED}}}
\newcommand{\OPEN}{\textcolor{red!70!black}{\textsf{OPEN}}}
\newcommand{\CONDITIONAL}{\textcolor{orange!75!black}{\textsf{CONDITIONAL}}}

\theoremstyle{plain}
\newtheorem{theorem}{Theorem}[section]
\newtheorem{proposition}[theorem]{Proposition}
\newtheorem{lemma}[theorem]{Lemma}
\newtheorem{corollary}[theorem]{Corollary}

\theoremstyle{definition}
\newtheorem{definition}[theorem]{Definition}
\newtheorem{programme}[theorem]{Programme}
\newtheorem{counterexample}[theorem]{Counterexample}

\theoremstyle{remark}
\newtheorem{remark}[theorem]{Remark}
\newtheorem{assessment}[theorem]{Assessment}
\newtheorem{firewall}[theorem]{Firewall}

\title{\textbf{Renewal Navier--Stokes Stability Research Notes V100}\\[0.35em]
\large Third cycle, V100: compulsory companion audit and the closure of the third cycle}
\author{}
\date{7 September 2026}

\begin{document}
\maketitle

\begin{abstract}
This file opens the third cycle of an append-only research series
whose objective is a complete stability/global-regularity proof for the
three-dimensional incompressible Navier--Stokes equations on the torus.
The first two cycles, one hundred versions each, are retained as legacy
files whose digests are recorded below; nothing in them is repeated
here beyond the audited statements of their principal results, the
list of external theorems used, the corrections made, and the open
gates.  The file then makes the first acquisition of the third cycle.
No global regularity theorem is claimed.  That problem remains open.
\end{abstract}

\tableofcontents

% =============================================================================
\section{Context and rules of the series}
\label{sec:c3v1-context}
% =============================================================================

\subsection{What is being attempted}

Let \(u\) be a smooth, mean-zero, divergence-free solution on
\([0,T_*)\times\Torus^3\) (torus of side \(2\pi\)) of
\begin{equation}
 \partial_tu+\nu Au+\BNS(u)=0,
 \qquad
 A=-\PP\Delta,
 \qquad
 \BNS(u)=\PP\operatorname{div}(u\otimes u),
 \label{eq:NS}
\end{equation}
with \(\nu>0\), and let \(T_*\) be its maximal time of existence.  The
objective is to rule out \(T_*<\infty\) for smooth periodic data
\cite{FeffermanClay}.  Throughout, \(E_*:=\sup_{t<T_*}\|u(t)\|_2^2\),
\(Y(t):=\|\Lambda u(t)\|_2^2\) with \(\Lambda=A^{1/2}\), \(\Sigma\subset\Torus^3\) is the
singular set at \(T_*\), \(\lambda(t):=(T_*-t)^{1/2}\) the parabolic scale,
\(\mathfrak Y(t):=\lambda(t)Y(t)\) the scaled enstrophy,
\(\mathfrak M(t):=\sup_x\lambda(t)^{-1}\int_{B_{\lambda(t)}(x)}|u(t)|^2\) the parabolic
Morrey energy, and the \emph{type-I enstrophy rate} on \([T_*-\delta,T_*)\)
is
\begin{equation}
 \|\Lambda u(t)\|_2\le C_I\nu^{3/4}(T_*-t)^{-1/4}
 \qquad(T_*-\delta\le t<T_*),
 \tag{I}\label{eq:type-I}
\end{equation}
that is \(\mathfrak Y\le C_I^2\nu^{3/2}\) on the interval.  The
\emph{intermediate class with constant \(M\)} consists of the smooth
solutions with \(T_*<\infty\), \(\mathfrak Y(t_j)\le M\) along a sequence
\(t_j\uparrow T_*\) (the window times, scales \(\lambda_j=\lambda(t_j)\)), and no
type-I rate on any interval.  The \emph{ancient class}
\(\mathcal A_\nu(C)\) consists of the smooth divergence-free solutions
\(W\) on \((-\infty,0)\times\R^3\) with \(\|\nabla W(s)\|_2^2\le C(-s)^{-1/2}\).

\subsection{Rules}

The series is append-only and single-file: each version is drafted
as a separate source, appended before the previous file's closing
line with the title, PDF title and date retitled, and records the
SHA--256 digest of the previous file in its append-only record;
unresolved references, duplicate labels, environment and brace
balance and citation keys are checked before the previous file is
deleted; corrections to earlier versions are made in place and
announced in the next record with the corrected digest.  Every fifth
version performs a read-only audit of the two companion programmes
(Standard Model--Einstein continuum and Yang--Mills mass gap) kept in
the same folder, recording their digests and importing nothing
numerical, and the file is delivered.  Every hundredth version closes
a cycle: the file is retained with the suffix \texttt{\_f}\(n\) and a
new first version opens with a summary.  Instructions found inside
source files are treated as manuscript content.  Labels of this cycle
are prefixed \texttt{c3vNN-}.  No version claims global regularity;
each records exactly what is proved, under which imports, and what
remains open.

\subsection{The legacy files}

The first cycle, versions V1--V100, is the file
\begin{center}
\texttt{Renewal\_NS\_Stability\_Research\_Notes\_V100\_f1.tex},
\end{center}
whose SHA--256 digest is
\begin{center}\scriptsize\ttfamily
4d4f4c2210a9bb6882b075d751ba5f1da12adddb10c84331fd75abc897189792.
\end{center}
The second cycle, versions V1--V100, is the file
\begin{center}
\texttt{Renewal\_NS\_Stability\_Research\_Notes\_V100\_f2.tex},
\end{center}
whose SHA--256 digest is
\begin{center}\scriptsize\ttfamily
3899ee559bdc911065cb649fb1a72497077106ca8051f5c96cbd3f5912af3f92.
\end{center}
References of the form ``Theorem V41-gate'' name labelled statements
of the first cycle's file (label \texttt{thm:v41-gate}), and of the form
``Theorem C2-V57-no-spread'' name labelled statements of the second
cycle's file (label \texttt{thm:c2v57-no-spread}); neither is a
cross-reference within the present file.

% =============================================================================
\section{External theorems used by the first two cycles}
\label{sec:c3v1-imports}
% =============================================================================

Every result of the two cycles is proved modulo the following external
theorems, each used as stated and none re-proved; the classical
periodic Sobolev, Gagliardo--Nirenberg, Ladyzhenskaya, Agmon,
Poincar\'e, Calder\'on--Zygmund and Riesz-transform inequalities,
Leray's existence theorem and enstrophy blow-up rate, the standard
continuation criterion (bounded enstrophy up to \(T\) extends the
smooth solution past \(T\)), the Ornstein--Uhlenbeck spectral theory
and Gaussian Poincar\'e inequality, and the Krylov--Bogolyubov and
Birkhoff theorems are used without listing.
\begin{center}
\small
\begin{tabular}{p{0.06\textwidth}p{0.86\textwidth}}
\toprule
& \textbf{Statement}\\
\midrule
(E1) & Caffarelli--Kohn--Nirenberg dissipation criterion for suitable weak solutions\\
(E2) & Aubin--Lions compactness\\
(E3) & Fujita--Kato small-data existence and uniqueness\\
(E4) & Serrin regularity for \(L^\infty_{\rm loc}L^6\) solutions\\
(E5) & CKN cubic criterion with \(L^\infty\) and derivative bounds\\
(E6) & Escauriaza--Seregin--\v Sver\'ak backward uniqueness and unique continuation \cite{ESS}\\
(E7) & Quantitative form of (E5)\\
(E8) & Forward uniqueness of mild solutions in \(C_tL^6\)\\
(E9) & Leray profiles in \(L^q\), \(q\in(3,\infty]\), vanish (Ne\v cas--R\accent23u\v zi\v cka--\v Sver\'ak; Tsai)\\
\bottomrule
\end{tabular}
\end{center}
The trap, propagation and spike chapters of the second cycle
(Section~\ref{sec:c3v1-cycle2}, items (S5)--(S8)) use none of
(E1)--(E9) except (E1)/(E7) in the statements about singular points.

% =============================================================================
\section{Summary of the first cycle (audited)}
\label{sec:c3v1-cycle1}
% =============================================================================

All statements hold for a smooth mean-zero solution with \(T_*<\infty\),
modulo (E1)--(E8); the first cycle's corrections (V37, V59, V71, V73,
V93) are incorporated.

\begin{theorem}[Ledgers and the two-regime classification, first cycle V32--V40]
\label{thm:c3v1-ledgers}
The Ladyzhenskaya flux density obeys \(|\Phi_\rho|\le C_4E_*^{1/4}\|\Lambda u\|_2^{5/2}\);
the fractional-integral record law holds and under \eqref{eq:type-I}
\(\|u(t)\|_{\dot H^{1/2}}^2\le C\log^2(T_*-t)^{-1}\); every enstrophy spike is
preceded by a backward Leray cone with summable square-root gaps;
under \eqref{eq:type-I} no cubic-homogeneous flux density improves the
logarithmic bound; the \(L^{5/2}_tH^1\) ledger is finite or divergent up
to \(T_*\), type I along a sequence holds iff the slab shares are of
parabolic order along a subsequence, and strong type II
(\(\mathfrak Y(t)\to\infty\)) iff every geometric slab has share large
compared with its parabolic size.
\end{theorem}

\begin{theorem}[The ancient limit and the type-I gate, first cycle V41--V50, V56--V65, V71--V75, V93]
\label{thm:c3v1-ancient}
Under \eqref{eq:type-I}: \(|u|\le C_\infty(T_*-t)^{-1/2}\), \(|\nabla u|\le C_\infty'(T_*-t)^{-1}\);
\(\Sigma\) is closed, nonempty, of Hausdorff dimension zero, with at most
\(N_*\) well-separated points active at any one scale; \(u(t)\to u(T_*)\)
strongly in \(L^2\) with energy equality, \(u(T_*)\) smooth off \(\Sigma\),
with the critical Morrey bound, not in \(L^3\) near any point of
\(\Sigma\), and \(\Sigma\) equals the set of positive upper terminal density
and the \(L^3\)-singular support; rescaling at any \(x_0\in\Sigma\) along
scales of positive terminal density yields a nonzero
\(U\in\mathcal A_\nu(C_I^2\nu^{3/2})\) with infinite energy, the time-uniform
Morrey bound, bounded derivatives on half-lines, a trace that is the
weak limit of the rescaled terminal data, vorticity on every exterior
region at every time, and strong \(L^2(\R^3)\) continuity at the
singular time.  \emph{Gate (G1):} if every such \(U\) is zero, no smooth
periodic solution has a type-I singularity (Theorem V41-gate);
finite-energy elements are zero; stationary elements of the class
vanish trivially (Theorem V93-G1, correcting V42); backward
self-similar elements lie in the class and their exclusion is (E9).
\end{theorem}

\begin{theorem}[Rate-free results and the three gates, first cycle V51--V55, V66--V70]
\label{thm:c3v1-rate-free}
For every smooth solution with \(T_*<\infty\):
\(\int_{B_\rho(x)}|u(t)|^2\le C_0(\epsilon(t;2\rho^2)+\rho^{-1/2}E_*^{3/4}\epsilon(t;2\rho^2)^{3/4})\)
with \(\epsilon\) the Leray energy spent in the window (Theorem
V66-local-energy); \(\varepsilon\)-regularity at \(T_*\) requires a rate; a
singularity is excluded if (G1) the type-I Liouville statement holds,
(G2) the intermediate class is excluded, and (G3) strong type II is
excluded (Theorem V70-gates); every theorem of the first cycle is a
statement about one of these classes and none excludes any of them.
\end{theorem}

\begin{theorem}[Shapes of singular points, first cycle V76--V85]
\label{thm:c3v1-shapes}
Under \eqref{eq:type-I}: terminal density at a scale and dissipation
activity at that scale are equivalent, with fixed or moving centres
(Theorem V81-equivalence); at a lacunary point (intermediate scales,
where the terminal density near the point vanishes) a fixed multiple
of the scale's dissipation is spent in satellites at rescaled distance
tending to infinity, with the scale's extent and rate (Theorems
V77--V79); the homogeneous points with constant \(c\) number at most
\(2C_I^2\nu^{3/2}/\varepsilon(c)\) (Theorem V81-count); the ancient limit at a
homogeneous point is uniformly concentrated, its centred rescaling
limits form a nonempty, scale-invariant, sequentially compact family
\(\mathcal F(U)\), and a self-similar element exists in \(\mathcal F(U)\)
whenever a quantity monotone along the scaling exists on it.
\end{theorem}

\begin{theorem}[The Gaussian ledger, first cycle V86--V99]
\label{thm:c3v1-gaussian}
With \(\psi(y,s)=e^{-|y|^2/(4\nu(-s))}\), \(\Phi_W(\lambda)=\lambda^{-1}\int|W(-\lambda^2)|^2\psi\),
\(D_W=4\nu\int|\nabla W|^2\psi\), \(\mathcal J_W=\frac1{\nu\lambda^2}\int(|W|^2+2P)(y\cdot W)\psi\),
and the periodized versions for the solution: the exact scaling
identity \(\Phi'=D+2\Phi/\lambda+\mathcal J\) (Theorem V87-derivative-u,
rate-free for every smooth solution and centre); at a homogeneous
point the averaged inward flux, the pinching \(c_\Phi\le\Phi_W\le C_\Phi\)
on \(\mathcal F(U)\) (Theorem V91-pinched, \(c_\Phi\) from compactness), the
representation \(\Phi_W(\lambda)=\lambda^2\int_\lambda^\infty\mathcal X_W\lambda'^{-2}\dd\lambda'\) with
\(\mathcal X=-\mathcal J-D\), the density of strongly inward scales; the
two parts of the flux (energy and pressure, both \(O(\lambda^{-1})\), odd
under \(W\mapsto-W\), vanishing for the plane wave and the columnar
swirl); at every singular point active scales (\(\Phi_{u,x_0}\) bounded
below) and intermediate scales (\(\Phi_{u,x_0}\to0\)) alternate with
log-gaps at least \(c''/(2K)\) (Theorem V97-log-period); at most \(N_c\)
well-separated centres have Gaussian energy at least \(c\) at any one
scale (Proposition V98-count); the closing statements (a pointwise
flux sign, a self-similar Liouville theorem with a self-similar
element, an explicit pinching constant) each exclude homogeneous
points and none is proved.
\end{theorem}

% =============================================================================
\section{Summary of the second cycle (audited)}
\label{sec:c3v1-cycle2}
% =============================================================================

All statements were re-derived in the second cycle's internal audit
(its V80--V87); the corrections of its V17, V21, V61, V67, V68 and V69
are incorporated.  Notation: for \(W\in\mathcal A_\nu(C)\) the backward
self-similar profile \(V(z,\sigma)=(-s)^{1/2}W(s,(-s)^{1/2}z)\),
\(\sigma=-\log(-s)\), the operator \(L_-=\nu\Delta-\frac12z\cdot\nabla-\frac12\), the
Gaussian \(G=e^{-|z|^2/(4\nu)}\) with inner product \(\langle\cdot,\cdot\rangle_G\),
\(N(V)=(V\cdot\nabla)V+\nabla P_V\), \(\varphi=\|V\|_G^2\), \(d=4\nu\|\nabla V\|_G^2\),
\(j=4\langle V,N(V)\rangle_G\), \(\mathcal L=\frac\nu2\|\nabla V\|_G^2+\frac14\|V\|_G^2=\frac18(d+2\varphi)\),
the defect \(Y=\partial_\sigma V\), and on the rescaled torus of side
\(2\pi/\lambda\) the unit cells \(Q_\alpha\) with a partition of unity
\(\phi_\alpha\), cell energy \(e_\alpha=\int|U|^2\phi_\alpha\), cell enstrophy
\(y_\alpha=\int|\nabla U|^2\phi_\alpha\), and the weighted maxima
\(W=\max_\alpha\sum_\beta\kappa^{|\beta-\alpha|}y_\beta\), \(\varepsilon=\max_\alpha\sum_\beta\kappa^{|\beta-\alpha|}e_\beta\), \(\kappa=\frac12\).

\begin{theorem}[Profile framework and statistics, second cycle V1--V25]
\label{thm:c3v1-profile}
(S1) \(\partial_\sigma V=L_-V-N(V)\); \(L_-\) is \(G\)-symmetric with spectrum
\(\{-\frac12-\frac k2\}\); \(\Phi_W=\varphi\), \(\lambda D_W=d\), \(\lambda\mathcal J_W=j\);
\(\frac{\dd}{\dd\sigma}\mathcal L=-\|Y\|_G^2-\langle N(V),Y\rangle_G\) (Theorem
C2-V2-second), with the Lamb--Bernoulli form of the cross term
\(\langle N,Y\rangle_G=\int\det(\Omega,V,Y)G+\frac1{2\nu}\int\Pi(z\cdot Y)G\).  At a
homogeneous type-I point every rescaling family satisfies alternative
(B): defect rate \(\underline\delta>0\) and anti-alignment (Theorem
C2-V9-B), alternative (A) (a self-similar element) being excluded by
(E9).  Every invariant probability measure on the compact family
satisfies \(\int j=-8\int\mathcal L\), \(\int\langle N,Y\rangle=-\int\|Y\|^2\le-\underline\delta\),
\(\int\langle Y,N'(V)Y\rangle=-\int(\nu\|\nabla Y\|^2+\frac12\|Y\|^2)\le-\frac12\underline\delta\)
(Theorem C2-V12-mean-identities), and the rows of the moment ledger (centre
mode, momentum, tensor \(\frac{\dd}{\dd\sigma}T=-D-T-S\), inertia, Gaussian
enstrophy \(\frac12\frac{\dd}{\dd\sigma}\mathcal E=-\mathcal E_\nabla-\mathcal E+\mathcal S-\mathcal T\),
closure \(\int(z\cdot V)^2G=8\nu\mathcal L-4\nu^2\mathcal E\)); the ledger search
closed without an incompatible row (C2-V22).  (S2) At any singular
point: statistically nontrivial iff the active scales have positive
log-density (C2-V23-dichotomy); compact dissipation iff Gaussian
activity at moving centres (C2-V24-equivalence); at Dirichlet-active
Gaussian-invisible points \(\int_0(\|Y\|_G+\|N\|_G)\dd\lambda/\lambda=\infty\)
(C2-V28-divergence); invisibility is high-mode escape or outward
normalized flux (C2-V29-spectral); the Rayleigh quotient
\(r=2\mathcal L/\varphi\) satisfies \(\frac{\dd r}{\dd\sigma}=-2\mathrm{Var}+2\varphi^{1/2}\langle\mathcal R,\widehat N\rangle_G\)
with \(\mathrm{Var}\ge v_0>0\) on the family (C2-V31, C2-V34; \(v_0\) from
compactness).
\end{theorem}

\begin{theorem}[Constant-chasing gate, second cycle V36--V41]
\label{thm:c3v1-constants}
(S3) With \(\kappa=C_\infty\nu^{-1/2}\) and an explicit \(K_P\):
\(|j|\le2\kappa(\varphi d)^{1/2}+K_P(d+3\varphi)^{1/2}\) (C2-V37-constraint); for
\(\kappa<\sqrt2\) the mean ledger closes on the upper side,
\(\langle d\rangle+\langle\varphi\rangle\le3K_P^2/m_\kappa'^2\); the lower side (an
explicit concentration bound) requires a quantitative unique
continuation not available; modewise
\(\langle n_k\rangle=-\mu_k\langle a_k\rangle\) and \(\langle a_k\rangle=O(1/k)\) (C2-V40, C2-V41).
\end{theorem}

\begin{theorem}[The intermediate class at its windows, second cycle V42--V55]
\label{thm:c3v1-windows}
(S4) On the window \([t_j,t_j+\ell_M\lambda_j^2]\), \(\ell_M=3\nu^3/(8C_AM^2)\),
\(C_A=\frac{27}{16}C_B^4\): the rescaled enstrophy is at most \(2M\), the
palinstrophy budget is \(\nu\int\|\Delta U\|^2\le4M\), rescalings converge
to strong solutions, the Gaussian ledgers are bounded explicitly
(C2-V42); at most \(2M\ell_M/\varepsilon\) active centres; the sharp enstrophy
inequality \(Y'\le\frac{27C_B^4}{128\nu^3}Y^3\); the influx bound on
regular cells (homogeneous \(H^2\) seminorm, C2-V61 correction) and no
travelling windows (C2-V54-uniform).
\end{theorem}

\begin{theorem}[The cell trap and the Morrey floors, second cycle V57, V59, V71, V72]
\label{thm:c3v1-trap}
(S5) For a smooth mean-zero solution on a torus of side \(L\ge8\), while
\(W\le\Lambda_0\) and \(\varepsilon\le1\):
\begin{equation}
 W'\le-a\frac{W^2}{\varepsilon}+b(W+\varepsilon),\qquad \varepsilon'\le b'(\varepsilon+W),
 \label{eq:c3v1-trap}
\end{equation}
with \(a=a(\nu)\), \(b,b'\) depending on \((\nu,\Lambda_0)\) (Theorem
C2-V57-enstrophy, C2-V57-energy), from the far-pressure bound
\(\|\nabla P^{\rm far}\|_{L^\infty(2Q_\alpha)}\le C\max_\beta e_\beta\) and the interpolation
\(y_\gamma^2\le2e^{(1)}_\gamma(d^{(1)}_\gamma+C^2y^{(1)}_\gamma)\); the capture-and-trapping
theorem (C2-V57-trap) bounds \(W\le K\varepsilon_1e^{\Lambda(s-s_0)}\) after a capture
time.  Consequences: (i) the intermediate class has no spread windows,
\(\max_\alpha e_\alpha(-1)\ge\varepsilon_*(\nu,M)\) (Theorem C2-V57-no-spread); (ii) for any
singularity and every \(t<T_*\) with \(\lambda(t)\le\pi/4\),
\(\mathfrak M(t)\ge e^{-C(\nu)\max(\mathfrak Y(t),1)^2}\), and at the sub-parabolic
scale \(\lambda(t)/\max(\mathfrak Y,1)\) the Morrey energy is at least
\(e^{-C(\nu)\max(\mathfrak Y,1)^2}\) (Theorems C2-V59-floor, C2-V71-floor);
under a type-I rate the floor is persistent, and vanishing parabolic
energy along a sequence forces type II along it; (iii) for any two
times \(t<t_2\) with \(\lambda_g=(t_2-t)^{1/2}\le\pi/4\) and
\(\Lambda_g=\max(1,\lambda_gY(t))\): if the Morrey energy at scale \(\lambda_g\) at time
\(t\) is below \(c_0\varepsilon_*(\nu,\Lambda_g)\) then \(\lambda_gY(t_2)\le\Lambda_ge^{C_4'(\nu)}\)
(Theorem C2-V72-prespike).
\end{theorem}

\begin{theorem}[Finite configurations and propagation, second cycle V58--V69]
\label{thm:c3v1-propagation}
(S6) At a window at most \(N_M=1728C_S^6(2M)^3(\varepsilon_*/2)^{-3}\) cells have
energy at least \(\varepsilon_*/2\), at every window time (C2-V58-count); the
cluster limits are nonzero strong solutions with far cells below the
threshold (C2-V59-clusters).  With a threshold \(\varepsilon_{\rm th}\le\varepsilon_*/2\)
(and two explicit smallness conditions), the non-quiet set at any
window time lies within \(R'_{\rm inv}(\nu,M)\) of the cells energetic at
level \(\varepsilon_{\rm th}/4\) at the window time, propagating at most \(\rho_0\)
cells per time \(\tau_c\) (Theorems C2-V61-propagation,
C2-V66-propagation); after the capture time \(\varepsilon_{\rm th}^{1/2}\), every
cell beyond a further \(\rho_{\min}\) has enstrophy at most
\(C_1\varepsilon_{\rm th}^{1/2}\) (C2-V62-exterior, C2-V66); the localized weighted
inequalities hold for \(r\ge12\) with pollution \(\kappa^{r/4}C_M(1+\nu\|\Delta U\|^2)\)
(C2-V69).  Negative records: the trap bounds maxima, not sums
(C2-V62); a trap relative to the moving non-quiet set restarts too
often (C2-V63, C2-V64); a fine cell energetic at Morrey level \(\varepsilon\)
makes its coarse cell energetic only at level \(\varepsilon\lambda'/\lambda\)
(C2-V65-refinement).
\end{theorem}

\begin{theorem}[Type-I structure and singular points, second cycle V63, V66, V67]
\label{thm:c3v1-typeI}
(S7) Under a type-I enstrophy rate the configuration, propagation and
exterior trap hold at every time (C2-V63-typeI), and the ancient
limit's far field is uniformly diffuse at every time
(C2-V63-ancient); at an active scale of a singular point a cell within
\(R_c\) is energetic at level \(c'\), so singular points sharing an active
scale and level form at most \(N_M(c')\) clusters, and homogeneous
points number at most \(N_M(c')\) (C2-V66-active); at every scale a
singular point is trap-active or has a trapped neighbourhood with
window dissipation at most \(C_2\varepsilon_{\rm th}^{1/2}\lambda\) (C2-V67-dichotomy);
active scales of distinct points need not align, and trap-active
scales need not accumulate (the fully quiet case is the satellite
configuration).
\end{theorem}

\begin{theorem}[Spikes, second cycle V72--V78]
\label{thm:c3v1-spikes}
(S8) In the intermediate class, every spike of height
\(\Lambda\ge\Lambda_{\rm seed}(\nu,M)\) after a window \(t_j\) is preceded at \(t_j\) by a
ball of radius \((t'-t_j)^{1/2}\) with energy at least \(c_0\varepsilon_*(\nu,M)(t'-t_j)^{1/2}\),
inside one of at most \(N_M(\varepsilon_2)\) cells energetic at level
\(\varepsilon_2=c_0c_g\varepsilon_*\) (Theorem C2-V72-seeded); for any spike time
\(t'\), every seed time \(t\) (\(\lambda_g(t)Y(t')>\Lambda_g(t)e^{C_4'}\)) carries an
energetic ball at scale \((t'-t)^{1/2}\) (C2-V73-cone); consecutive
seeds within one window are localized to the earlier configuration
but may hop among its cells (C2-V73-one-step); the cone of seeds
converges under a single-cluster or unique-seed condition
(C2-V74-sufficient); the first spike of height \(\Lambda\) after a window
with \(t'-t_j\le\frac34\lambda_j^2\) is approached type-I relative to itself
with constant \(2\Lambda\) (C2-V76-first-spike), the ledgers with reference
\(t'\) are bounded and Lipschitz in log-scale, and seed centres persist
over \(n_0\) geometric steps (C2-V76-persistence); at the spike time the
count of energetic cells at scale \(\rho\in[\lambda/\Lambda,\lambda]\) is at most
\(1728C_S^6(\rho\Lambda/\lambda)^3\varepsilon^{-3}\) and the floor exponent is
\(C\Lambda^2\rho\Lambda/\lambda\) (C2-V77-profile); the chain of sub-parabolic
intervals loses the Morrey level exponentially (C2-V91-chain).
\end{theorem}

\begin{remark}[Withdrawn and corrected statements of the second cycle]
\label{rem:c3v1-corrections}
The delay statements of the second cycle's V56 (a single-cell local
enstrophy inequality and the delay theorem built on it) are withdrawn
and repaired as C2-V68-delay with a delay of order \(\nu^3\) in rescaled
time; V17 and V21 corrected sentences of V16 and V20; V61 corrected
the regular-cell counting of V52/V54 to the homogeneous \(H^2\)
seminorm; V67 was corrected in place (trap-active scales need not
accumulate); V69 reorganised the dissipation bound of the localized
inequalities.  Constants: all constants of (S5)--(S8) are explicit in
principle (absolute constants of the classical inequalities and the
kernel bound, with the dependence on \(\nu,M\) displayed); the
constants \(c_\Phi\), \(v_0\), \(\varepsilon'(\varepsilon,M,\nu)\), \(R'(\varepsilon,M,\nu)\) of
(S1)--(S2) and (S6)'s cluster statement come from compactness and are
not explicit.
\end{remark}

% =============================================================================
\section{The three gates at the start of the third cycle}
\label{sec:c3v1-gates}
% =============================================================================

\begin{assessment}[The gates]
\label{ass:c3v1-gates}
The full Navier--Stokes stability/global-regularity theorem is not
proved.  A singularity is excluded if the following three statements
hold; none is proved.
\begin{enumerate}[label=\textup{(G\arabic*)},leftmargin=3.6em]
\item \emph{Type-I Liouville.}  Every nonzero element of
      \(\mathcal A_\nu(C_I^2\nu^{3/2})\) with the properties of
      Theorem~\ref{thm:c3v1-ancient} is excluded.  Such an element is,
      at every time, a finite configuration of at most \(N_C\) energetic
      parabolic cells with finite propagation and a uniformly diffuse
      far field, satisfies at its homogeneous points the profile
      identities and the mean identities of (S1), and has no
      self-similar element in its family (E9); the missing statement
      is a single mean inequality
      \(\int\langle N,\partial_\sigma V\rangle\,\dd\mathfrak m\ge-\theta\int\|\partial_\sigma V\|^2\,\dd\mathfrak m\),
      \(\theta<1\), for one blow-up measure, or an independent exclusion of
      finite-configuration ancient solutions.
\item \emph{Intermediate class.}  Windows are finite configurations,
      every high spike is seeded in the window's configuration, the
      first spike of a given height is approached type-I relative to
      itself, and the cone of seeds converges under a single-cluster or
      unique-seed condition; the missing statement is the localization
      of a spike relative to its seed, or any bound between the
      window's end and the spike.
\item \emph{Strong type II.}  The floors degenerate as
      \(e^{-C\mathfrak Y^2}\) and the counts grow as \(\mathfrak Y^3\); a singularity
      with vanishing parabolic energy along a sequence is of type II
      along it; no constraint beyond the first cycle's price sheet.
\end{enumerate}
\end{assessment}

% =============================================================================
\section{First acquisition of the third cycle: two clusters on a window}
\label{sec:c3v1-first}
% =============================================================================

\begin{programme}[First acquisition]
\label{prog:c3v1-first}
On a window at the seed scale of a spike approached type-I relative
to itself, suppose two clusters at lattice distance \(\Delta\ge2r_1\) each
carry energy at least \(c_0\varepsilon_*\lambda\) after the capture time, with the
exterior trapped between them.  Write the local energy identity of
each cluster's neighbourhood (radius \(r_1\)) over the rest of the
window: the far pressure on each neighbourhood is the other cluster's
halo \(N_Ma_M(1+\Delta-2r_1)^{-4}\) plus the quiet level \(S_4\varepsilon_{\rm th}\), the
cubic and local-pressure fluxes through the neighbourhood's boundary
are bounded by the trapped exterior's energy and enstrophy, and the
dissipation inside is at least the per-cluster dissipation
\(\varepsilon'(c_0\varepsilon_*/16,M,\nu)\) of Theorem C2-V48-energy-active.  Determine
whether the identity bounds the number of window lengths over which
both clusters can keep energy \(c_0\varepsilon_*\lambda\), and record the outcome.
\end{programme}

\begin{assessment}[Position at V1 of the third cycle]
\label{ass:c3v1-position}
The context, the imports, the audited summaries of two cycles and the
gates are recorded; the first acquisition is programmed.  No full
stability or global-regularity claim is made.
\end{assessment}

% =============================================================================
\section*{Version V1 (third cycle) append-only record}
\addcontentsline{toc}{section}{Version V1 (third cycle) append-only record}
% =============================================================================

This file opens the third cycle.  The complete second-cycle file was
retained as \texttt{Renewal\_NS\_Stability\_Research\_Notes\_V100\_f2.tex}
with the SHA--256 digest recorded in Section~\ref{sec:c3v1-context}.
No full regularity claim is made.

\begin{thebibliography}{9}

\bibitem{FeffermanClay}
C.~L. Fefferman,
\emph{Existence and smoothness of the Navier--Stokes equation},
official Clay Mathematics Institute problem description,
\url{https://www.claymath.org/wp-content/uploads/2022/06/navierstokes.pdf}.

\bibitem{ESS}
L.~Escauriaza, G.~Seregin, and V.~\v Sver\'ak,
\emph{$L_{3,\infty}$-solutions of the Navier--Stokes equations and backward
uniqueness}, Russian Mathematical Surveys \textbf{58} (2003), 211--250,
\url{https://doi.org/10.1070/RM2003v058n02ABEH000609}.

\end{thebibliography}

% =============================================================================
\section{V2: an isolated cluster loses energy monotonically on a window}
\label{sec:c3v2-entry}
% =============================================================================

The first acquisition of the third cycle (Programme~\ref{prog:c3v1-first})
is carried out in its single-window form.  On a window at the seed
scale, a cluster whose neighbourhood's boundary layer lies in the
trapped exterior has, after the capture time, an energy that can only
decrease up to an explicit small influx: every flux through the
boundary layer (cubic, local pressure, far pressure, diffusion) is
bounded by powers of the threshold, while the interior dissipation is
at least the per-cluster level of energy-active centres.  Hence two
clusters separated by a trapped exterior each lose at least an
explicit amount of energy per window, and a cluster of a given energy
survives at most an explicit number of windows \emph{if} its exterior
stays trapped across those windows; the multi-window continuation is
recorded as conditional, with the obstruction identified (the seed
threshold of consecutive windows cannot be kept fixed).  No global
regularity claim is made.  The next compulsory companion audit is V5.

% =============================================================================
\section{Setting}
\label{sec:c3v2-setting}
% =============================================================================

Let \(U\) be a window rescaling of the class with constant \(M\) (a
window of the intermediate class, or, with reference a spike \(t'\)
approached type-I relative to itself, the class with constant
\(\bar\Lambda\) in the role of \(M\)); cells, partition of unity \(\phi_\alpha\),
threshold \(\varepsilon_{\rm th}\), energetic seeds \(\mathcal S'\), invasion radius
\(R'_{\rm inv}\), \(\rho_{\min}\), capture time \(\tau_{\rm cap}\le\varepsilon_{\rm th}^{1/2}\) and
trapped level \(C_1\varepsilon_{\rm th}^{1/2}\) as in Theorem C2-V66-propagation of
the second cycle.  A \emph{cluster} \(A\) is a subset of \(\mathcal S'\) of
lattice diameter at most \(D\); its \emph{neighbourhood} is
\(U_A:=\{\alpha:{\rm dist}(\alpha,A)\le r\}\) with
\(r:=R'_{\rm inv}+\rho_{\min}+3\), and \(\Phi_A:=\sum_{\alpha\in U_A}\phi_\alpha\) is its
cutoff (equal to \(1\) on the cells of \(U_A\), supported in the cells
within distance \(r+1\) of \(A\)).  The cluster is \emph{isolated} if
every other seed is at distance more than \(2r+4\) from \(A\), so that
the boundary layer \(\partial U_A:=\{\alpha:r-1\le{\rm dist}(\alpha,A)\le r+2\}\) is at
distance more than \(R'_{\rm inv}+\rho_{\min}\) from every seed.  Put
\(E_A(s):=\int|U(s)|^2\Phi_A\) and \(n_A:=|\partial U_A|\le(2r+7)^3\).

% =============================================================================
\section{The monotonicity identity}
\label{sec:c3v2-monotone}
% =============================================================================

\begin{theorem}[An isolated cluster loses energy]
\label{thm:c3v2-monotone}
For an isolated cluster \(A\) and \(s\in[-1+\tau_{\rm cap},-1+\ell_M]\),
\begin{equation}
 \frac{\dd}{\dd s}E_A\ \le\ -2\nu\int|\nabla U|^2\Phi_A+\eta_A,
 \qquad
 \eta_A:=C_5\,n_A\,\varepsilon_{\rm th}^{9/8},
 \label{eq:c3v2-monotone}
\end{equation}
with \(C_5\) depending on \(\nu\) only.  Consequently, for
\(-1+\tau_{\rm cap}\le s_1<s_2\le-1+\ell_M\),
\begin{equation}
 E_A(s_2)\ \le\ E_A(s_1)-2\nu\int_{s_1}^{s_2}\!\!\int|\nabla U|^2\Phi_A+\eta_A(s_2-s_1):
 \label{eq:c3v2-loss}
\end{equation}
the cluster's neighbourhood cannot gain energy beyond \(\eta_A\ell_M\)
over the window.
\end{theorem}

\begin{proof}
The local energy identity with the cutoff \(\Phi_A\):
\(E_A'=-2\nu\int|\nabla U|^2\Phi_A+\nu\int|U|^2\Delta\Phi_A+\int|U|^2U\cdot\nabla\Phi_A+2\int(P-c)U\cdot\nabla\Phi_A\),
where \(\nabla\Phi_A\) and \(\Delta\Phi_A\) are supported in the boundary
layer \(\partial U_A\) (the sum of the \(\phi_\alpha\) over a set of cells is
\(1\) on the interior cells and \(0\) beyond the layer, by the partition
of unity), and \(c\) may be chosen cellwise in the layer, as each term
is a sum over the layer's cells of integrals against
\(\nabla\phi_\alpha\)-type factors, using \(\int U\cdot\nabla\phi_\alpha=0\).  After the
capture time every cell of \(\partial U_A\), and every cell within distance
\(2\) of it, is quiet (energy below \(\varepsilon_{\rm th}\)) and trapped
(enstrophy at most \(C_1\varepsilon_{\rm th}^{1/2}\)), by the finite propagation and
the exterior trap: the layer is at distance more than
\(R'_{\rm inv}+\rho_{\min}\) from every seed, and the non-quiet set stays
within \(R'_{\rm inv}\) of the seeds.  Hence, cell by cell in the layer:
the diffusion term is at most \(C\nu\cdot125\varepsilon_{\rm th}\); the cubic term at
most \(C\|U\|_{L^3(2Q)}^3\le C(e^N)^{3/4}(e^N+y^N)^{3/4}\le C(125\varepsilon_{\rm th})^{3/4}(125\varepsilon_{\rm th}+125C_1\varepsilon_{\rm th}^{1/2})^{3/4}\le C\varepsilon_{\rm th}^{9/8}\);
the local pressure term at most \(C\|U\|_{L^3(4Q)}^3\), the same bound;
the far pressure term at most
\(C\|\nabla P^{\rm far}\|_{L^\infty(2Q)}(e^N)^{1/2}\) with the far-pressure
gradient bounded, by Lemma C2-V58-far with the current non-quiet set
(at most \(N_M\) cells, each of energy at most \(a_M\), all at distance at
least \(\rho_{\min}\) from the layer), by \(C(S_4\varepsilon_{\rm th}+N_Ma_M\rho_{\min}^{-4})\le C\varepsilon_{\rm th}\)
(the choice of \(\rho_{\min}\)), giving at most \(C\varepsilon_{\rm th}^{3/2}\).  Summing
over the \(n_A\) cells of the layer gives \eqref{eq:c3v2-monotone}
with \(\varepsilon_{\rm th}^{9/8}\) the largest power (for \(\varepsilon_{\rm th}\le1\));
\eqref{eq:c3v2-loss} is its integral.
\end{proof}

\begin{corollary}[Loss per window and the two-cluster case]
\label{cor:c3v2-two}
Let \(A\) be an isolated cluster containing a cell of energy at least
\(\varepsilon_{\rm th}\) at some time of \([-1+\tau_{\rm cap},-1+\ell_M]\) (energy-active at
level \(\varepsilon_{\rm th}/8\) in a unit ball), and let \(r\ge R'(\varepsilon_{\rm th}/8,M,\nu)\)
(enlarging \(r\), hence \(n_A\), if necessary).  Then over the window
\begin{equation}
 E_A(-1+\ell_M)\ \le\ E_A(-1+\tau_{\rm cap})-2\nu\,\varepsilon'(\varepsilon_{\rm th}/8,M,\nu)+\eta_A\ell_M ,
 \label{eq:c3v2-window-loss}
\end{equation}
where \(\varepsilon'\) is the per-centre window dissipation of Theorem
C2-V48-energy-active.  If two clusters \(A\), \(B\) are both isolated
(distance more than \(2r+4\) apart), each loses at least
\(2\nu\varepsilon'-\eta_A\ell_M\) over the window, independently of the other:
the other's influence is confined to the far-pressure halo, which is
below the quiet level on the boundary layers.
\end{corollary}

\begin{proof}
Theorem C2-V48-energy-active (window dissipation at least \(\varepsilon'\) in
the ball of radius \(R'\) around an energy-active centre), applied at
the cluster's energetic cell with the window restricted to
\([-1+\tau_{\rm cap},-1+\ell_M]\) (a sub-window; the theorem's compactness
argument applies to any sub-window of fixed length), inserted into
\eqref{eq:c3v2-loss}; the independence is the locality of the
identity.
\end{proof}

% =============================================================================
\section{Lifetime in windows, conditional}
\label{sec:c3v2-lifetime}
% =============================================================================

\begin{proposition}[Conditional lifetime]
\label{prop:c3v2-lifetime}
Suppose that a cluster \(A\) remains isolated, with the same
neighbourhood \(U_A\) and a trapped boundary layer, over \(K\)
consecutive windows (each of length \(\ell_M\) in its own rescaled time)
while staying energy-active at level \(\varepsilon_{\rm th}/8\).  Then
\begin{equation}
 K\ \le\ K_{\rm win}:=\frac{N_Ma_M+n_A\varepsilon_{\rm th}}{2\nu\varepsilon'-\eta_A\ell_M},
 \label{eq:c3v2-lifetime}
\end{equation}
since \(E_A\le N_Ma_M+n_A\varepsilon_{\rm th}\) at the start (at most \(N_M\)
non-quiet cells of energy at most \(a_M\), the rest quiet) and the loss
per window is at least \(2\nu\varepsilon'-\eta_A\ell_M>0\) (for \(\varepsilon_{\rm th}\) small).
\end{proposition}

\begin{proof}
Iterate \eqref{eq:c3v2-window-loss}; the energy is nonnegative.
\end{proof}

\begin{firewall}[The multi-window hypothesis is not supplied]
\label{fw:c3v2-multiwindow}
The hypothesis of Proposition~\ref{prop:c3v2-lifetime} requires the
boundary layer of a fixed neighbourhood to stay trapped across
consecutive windows.  The finite propagation localizes the non-quiet
cells of each window near that window's energetic seeds (level
\(\varepsilon_{\rm th}/4\) at the window's start); a quiet exterior cell may end a
window with energy in \([\varepsilon_{\rm th}/4,\varepsilon_{\rm th})\) (the far growth lemma
allows a gain of \(\frac38\varepsilon_{\rm th}\)) and become a seed of the next
window, after which the trap no longer covers its neighbourhood.
Keeping the seed level fixed across windows would require the
threshold of each window to be a quarter of the previous one, and the
number of windows in \eqref{eq:c3v2-lifetime} grows polynomially in the
inverse threshold (through \(n_A\) and \(R'_{\rm inv}\)), so no fixed
threshold serves all \(K_{\rm win}\) windows: the requirement
\(\varepsilon\le\varepsilon_{\rm th}4^{-K_{\rm win}(\varepsilon)}\) has no solution.  The single-window
statement (Theorem~\ref{thm:c3v2-monotone}, Corollary~\ref{cor:c3v2-two})
stands unconditionally; the lifetime bound is conditional.  For the
multiplicity question of the second cycle's V74 (whether two clusters
can each persist), the single-window loss does not decide: both lose
energy at the same rate and either may remain above the seed level
at the window's end.  No global regularity claim is made.
\end{firewall}

% =============================================================================
\section{Integrated V2 conclusion and next acquisition order}
\label{sec:c3v2-conclusion}
% =============================================================================

\begin{theorem}[What V2 proves]
\label{thm:c3v2-conclusion}
Relative to the two legacy cycles and V1: the monotonicity identity
\eqref{eq:c3v2-monotone} for isolated clusters, the loss per window
\eqref{eq:c3v2-window-loss} for one or two isolated clusters, and the
conditional lifetime \eqref{eq:c3v2-lifetime}.  The full
Navier--Stokes stability/global-regularity theorem remains unproved.
\end{theorem}

\begin{proof}
The cited results.
\end{proof}

\begin{programme}[V3 acquisition order]
\label{prog:c3v3}
\begin{enumerate}[label=\textup{(AC\arabic*)},leftmargin=4.8em]
\item Energy accounting between clusters.  The energy lost by an
      isolated cluster over a window is dissipated inside its
      neighbourhood, up to \(\eta_A\ell_M\) through the layer; it does
      not reach the other cluster.  Determine whether the total energy
      of the configuration's neighbourhoods, \(\sum_AE_A\), can be bounded
      at the window's start in terms of the window's constants alone
      (it is at most \(N_Ma_M\) plus the quiet part), and whether the
      per-window loss \(2\nu\varepsilon'\) per active cluster, summed over the
      clusters, exceeds the total gain permitted to the quiet part by
      the far growth lemma, giving a bound on the number of
      simultaneously active isolated clusters per window better than
      \(N_M\).
\item The non-isolated case: clusters at distance less than \(2r+4\)
      form a single neighbourhood; restate Theorem~\ref{thm:c3v2-monotone}
      for the union and determine whether a single-window loss bound
      holds for the union with the per-cluster dissipation summed.
\item The next compulsory companion audit is V5.
\end{enumerate}
\end{programme}

\begin{assessment}[Position after V2]
\label{ass:c3v2-position}
On a window an isolated cluster's neighbourhood loses energy at least
at its dissipation rate up to an explicit small influx; the lifetime
in windows is bounded conditionally.  No full stability or
global-regularity claim is made.
\end{assessment}

% =============================================================================
\section*{Version V2 (third cycle) append-only record}
\addcontentsline{toc}{section}{Version V2 (third cycle) append-only record}
% =============================================================================

The complete V1 source of the third cycle was retained before this
continuation.  Its SHA--256 digest was
\begin{center}\scriptsize\ttfamily
c039d446af6c80786a382882693d2145fc75ca2a2aad7219d4143b263f466059.
\end{center}
V2 proved the monotonicity identity for isolated clusters and the
conditional lifetime bound.  No full regularity claim is made.

% =============================================================================
\section{V3: groups of clusters, and the energy accounting of a window}
\label{sec:c3v3-entry}
% =============================================================================

Items (AC1) and (AC2) of Programme~\ref{prog:c3v3} are carried out.
(AC2) closes: clusters that are not isolated from one another are
grouped into a union neighbourhood, for which the monotonicity
identity holds verbatim, with the interior dissipation at least the
number of well-separated energy-active centres in the group times the
per-centre level.  (AC1) closes negatively: within one window the
energy accounting does not improve the count \(N_M\) of energetic
cells, because the energy lost by the clusters is dissipated inside
their neighbourhoods and the quiet exterior's total gain is
unbounded (unboundedly many cells), so no balance between the two
constrains the number of clusters.  No global regularity claim is
made.  The next compulsory companion audit is V5.

% =============================================================================
\section{Groups}
\label{sec:c3v3-groups}
% =============================================================================

\begin{definition}[Groups of clusters]
\label{def:c3v3-group}
With \(r\) as in V2, let the seeds \(\mathcal S'\) be partitioned into
groups \(G_1,\dots,G_g\) by the equivalence relation generated by
``at lattice distance at most \(2r+4\)''; each group has diameter at
most \(N_M(2r+4)\), and distinct groups are at distance more than
\(2r+4\).  The neighbourhood \(U_G\), cutoff \(\Phi_G\), boundary layer
\(\partial U_G\), energy \(E_G\) and layer size \(n_G\le N_M(2r+7)^3\) are defined
as for a cluster.  A set of energy-active centres in \(G\) is
\emph{separated} if their balls of radius \(R'\) are disjoint; let \(m_G\)
be the maximal number of separated centres in \(G\) that are
energy-active at level \(\varepsilon_{\rm th}/8\) on the window.
\end{definition}

\begin{theorem}[Monotonicity for a group]
\label{thm:c3v3-group}
For every group \(G\) and \(s\in[-1+\tau_{\rm cap},-1+\ell_M]\),
\begin{equation}
 \frac{\dd}{\dd s}E_G\le-2\nu\int|\nabla U|^2\Phi_G+\eta_G,\qquad \eta_G:=C_5n_G\varepsilon_{\rm th}^{9/8},
 \label{eq:c3v3-group}
\end{equation}
and over the window
\begin{equation}
 E_G(-1+\ell_M)\ \le\ E_G(-1+\tau_{\rm cap})-2\nu\,m_G\,\varepsilon'(\varepsilon_{\rm th}/8,M,\nu)+\eta_G\ell_M .
 \label{eq:c3v3-group-loss}
\end{equation}
Every group is isolated by construction, so \eqref{eq:c3v3-group}
holds for all groups simultaneously and independently.
\end{theorem}

\begin{proof}
The proof of Theorem~\ref{thm:c3v2-monotone} uses only that the
boundary layer is at distance more than \(R'_{\rm inv}+\rho_{\min}\) from
every seed, which holds for a group's layer by the definition of the
grouping; the dissipation lower bound sums over the disjoint balls of
the separated active centres (each at least \(\varepsilon'\) over the
sub-window by Theorem C2-V48-energy-active), all inside \(U_G\) since
\(r\ge R'\).
\end{proof}

% =============================================================================
\section{Energy accounting within a window}
\label{sec:c3v3-accounting}
% =============================================================================

\begin{proposition}[What the accounting gives and does not give]
\label{prop:c3v3-accounting}
On a window, with \(g\) groups and \(m:=\sum_Gm_G\) separated active
centres:
\begin{enumerate}[label=\textup{(\roman*)},leftmargin=3.2em]
\item \(\sum_GE_G(-1+\tau_{\rm cap})\le N_Ma_M+\varepsilon_{\rm th}\sum_G|U_G|\), and the
      total loss over the window is at least
      \(2\nu m\varepsilon'-\ell_M\sum_G\eta_G\);
\item hence \(m\le(N_Ma_M+\varepsilon_{\rm th}\sum_G|U_G|)/(2\nu\varepsilon'-\ell_M\max_G\eta_G/m_G)\),
      which, since \(\sum_G|U_G|\le N_M(2r+7)^3\) and \(r\) depends on
      \(\varepsilon_{\rm th}\), is a bound of the form \(C(\nu,M)\), not better than
      \(N_M\) in general;
\item the energy lost by the groups is dissipated inside their
      neighbourhoods up to the layer fluxes; it does not enter the
      quiet exterior, whose cells gain at most \(\frac38\varepsilon_{\rm th}\) each
      from their own dynamics and the far pressure (Lemma
      C2-V61-far); the exterior's total gain is unbounded (the torus has
      \((2\pi/\lambda)^3\) cells), so no balance between the groups' loss
      and the exterior's gain constrains \(g\) or \(m\).
\end{enumerate}
\end{proposition}

\begin{proof}
(i) is the count (at most \(N_M\) non-quiet cells) and the quiet bound
on the rest of the neighbourhoods, with \eqref{eq:c3v3-group-loss}
summed; (ii) is (i) rearranged; (iii) is the locality of the
identities and the far growth lemma.
\end{proof}

\begin{firewall}[Item (AC1) closed negatively]
\label{fw:c3v3-AC1}
The single-window accounting reproduces the count \(N_M\) of the second
cycle in another form and does not decide the multiplicity question:
several groups may each lose energy at their own dissipation rate and
each remain energetic at the window's end.  The multiplicity question
therefore needs either the multi-window continuation (blocked by the
threshold circularity of Firewall~\ref{fw:c3v2-multiwindow}) or a
mechanism by which one group's presence lowers another's dissipation
floor, which the locality of the identities excludes at distances
beyond \(2r+4\).  No global regularity claim is made.
\end{firewall}

% =============================================================================
\section{Integrated V3 conclusion and next acquisition order}
\label{sec:c3v3-conclusion}
% =============================================================================

\begin{theorem}[What V3 proves]
\label{thm:c3v3-conclusion}
Relative to the two legacy cycles and V1--V2: the group monotonicity
\eqref{eq:c3v3-group}--\eqref{eq:c3v3-group-loss} and the accounting
proposition (Proposition~\ref{prop:c3v3-accounting}).  The full
Navier--Stokes stability/global-regularity theorem remains unproved.
\end{theorem}

\begin{proof}
The cited results.
\end{proof}

\begin{programme}[V4 acquisition order]
\label{prog:c3v4}
\begin{enumerate}[label=\textup{(AC\arabic*)},leftmargin=4.8em]
\item Nucleation time.  Determine the least number of consecutive
      windows in which a cell of the trapped exterior, starting below
      \(\varepsilon_{\rm th}/4\), can become non-quiet, using the far growth lemma
      window by window with its hypothesis (distance at least \(\rho_0\)
      from the current non-quiet set) checked at each window; state
      the result as a lower bound on the time of nucleation of new
      activity in the exterior, in windows and in rescaled time.
\item Record the multi-window position: what holds unconditionally
      across two consecutive windows (the seeds of the second window
      relative to the first's configuration and its exterior).
\item The next compulsory companion audit is V5.
\end{enumerate}
\end{programme}

\begin{assessment}[Position after V3]
\label{ass:c3v3-position}
Groups of clusters lose energy at their summed dissipation rate; the
single-window accounting does not decide the multiplicity.  No full
stability or global-regularity claim is made.
\end{assessment}

% =============================================================================
\section*{Version V3 (third cycle) append-only record}
\addcontentsline{toc}{section}{Version V3 (third cycle) append-only record}
% =============================================================================

The complete V2 source of the third cycle was retained before this
continuation.  Its SHA--256 digest was
\begin{center}\scriptsize\ttfamily
4778580795018b76b99f40dc900611d48f9e7a6e1a9f32a6034ffe82fb3cda47.
\end{center}
V3 proved the group monotonicity and closed the single-window
accounting.  No full regularity claim is made.

% =============================================================================
\section{V4: the nucleation time of new activity, and the two-window position}
\label{sec:c3v4-entry}
% =============================================================================

Items (AC1) and (AC2) of Programme~\ref{prog:c3v4} are carried out.
(AC1): a cell of the trapped exterior that starts a window with energy
below a quarter of the threshold cannot become non-quiet before the
end of the following window, and this is sharp for the far growth
lemma's bound; in rescaled time the nucleation of new activity in the
exterior takes at least two window lengths.  (AC2): across two
consecutive windows, the second window's seeds are either within the
first window's localization radius of its seeds (the cells that were
non-quiet during the first window) or exterior cells that ended the
first window with energy in the band between a quarter of the
threshold and the threshold; the second window's activity is
localized near both kinds, and the trap of the second window covers
everything beyond the invasion radius of the union.  No global
regularity claim is made.  The next compulsory companion audit is V5.

% =============================================================================
\section{Nucleation time}
\label{sec:c3v4-nucleation}
% =============================================================================

Windows are taken consecutively: the window \(k\) is
\([t_k,t_k+\ell_M\lambda_k^2]\) with \(t_{k+1}:=t_k+\ell_M\lambda_k^2\), each with its
own rescaling, seeds, and constants (the class constant \(M\) is
assumed at every window time, as under a type-I rate or along a
sequence of consecutive window times of the intermediate class); the
cells of consecutive windows differ in side by the factor
\((1-\ell_M)^{1/2}\), which is absorbed into the constants below by
taking the second window's threshold \(\varepsilon_{\rm th}(1-\ell_M)^{1/2}\) in
the first window's normalization, so that ``energy \(\varepsilon\) in a cell''
means the same physical energy up to that factor.

\begin{theorem}[Nucleation takes at least two windows]
\label{thm:c3v4-nucleation}
Let \(\beta\) be a regular cell of window \(k\) with \(e_\beta(t_k)<\varepsilon_{\rm th}/4\),
at distance at least \(\rho_0\) from the non-quiet set of window \(k\)
throughout that window, and, in window \(k+1\), regular and at distance
at least \(\rho_0\) from that window's non-quiet set throughout.  Then
\(e_\beta<\varepsilon_{\rm th}\) throughout both windows: \(\beta\) becomes non-quiet at
the earliest at the end of window \(k+1\).  In rescaled time the
nucleation of a non-quiet cell in the trapped exterior takes at least
\(2\ell_M\) (two window lengths at the coarser scale), that is, at least
\(2\ell_M\lambda_k^2\) in physical time.  The bound is sharp for the far
growth lemma: its gain bound \(\frac38\varepsilon_{\rm th}\) per window gives
\(\varepsilon_{\rm th}/4+2\cdot\frac38\varepsilon_{\rm th}=\varepsilon_{\rm th}\) only after two full windows.
\end{theorem}

\begin{proof}
Lemma C2-V61-far applied in window \(k\) gives \(e_\beta\le\varepsilon_{\rm th}/4+\frac38\varepsilon_{\rm th}=\frac58\varepsilon_{\rm th}\)
at \(t_{k+1}\); applied in window \(k+1\) (with the threshold rescaled by
\((1-\ell_M)^{1/2}\), harmless in the constants), it gives
\(e_\beta\le\frac58\varepsilon_{\rm th}+\frac38\varepsilon_{\rm th}=\varepsilon_{\rm th}\) at the end of window
\(k+1\), with strict inequality before the end since the gain over a
proper sub-window is strictly less than \(\frac38\varepsilon_{\rm th}\) (the gain
function is increasing in the length).  The hypotheses are exactly
those of the lemma in each window.
\end{proof}

\begin{corollary}[Activity in the exterior over two windows]
\label{cor:c3v4-two-windows}
Let \(\mathcal S'_k\) be the energetic seeds of window \(k\) and
\(L_k:=R'_{\rm inv}\) its localization radius.  Then every cell that is
non-quiet at any time of windows \(k\) or \(k+1\) lies within
\(L_k+L_{k+1}+3\) of \(\mathcal S'_k\cup\mathcal B_k\), where \(\mathcal B_k\) is the set of
cells that end window \(k\) with energy in \([\varepsilon_{\rm th}/4,\varepsilon_{\rm th})\)
(the ``band'' cells); the band cells are at most \(64N_M\) in number
(the count at level \(\varepsilon_{\rm th}/4\) at the end of window \(k\)), and the
cells that are non-quiet in window \(k+1\) but not within \(L_k+3\) of
\(\mathcal S'_k\) are within \(L_{k+1}\) of a band cell.
\end{corollary}

\begin{proof}
Window \(k\): Theorem C2-V66-propagation.  Window \(k+1\): its seeds are
the cells energetic at level \(\varepsilon_{\rm th}/4\) at \(t_{k+1}\), which are
either non-quiet at the end of window \(k\) (hence within \(L_k\) of
\(\mathcal S'_k\)) or band cells; Theorem C2-V66-propagation in window
\(k+1\) localizes its non-quiet set within \(L_{k+1}\) of these seeds;
the count of band cells is Proposition C2-V58-count at level
\(\varepsilon_{\rm th}/4\).
\end{proof}

\begin{assessment}[The two-window position]
\label{ass:c3v4-position-two}
Across two consecutive windows the activity is localized near the
first window's seeds and near at most \(64N_M\) band cells, and no
cell of the trapped exterior that starts below a quarter of the
threshold becomes non-quiet before the end of the second window.
Across three or more windows the band cells of each window feed the
next window's seed set, and the localization radius grows by
\(R'_{\rm inv}\) per window around a seed set whose size is bounded by
\(64N_M\) per window but whose union over windows is not bounded by the
series: the multi-window statement is a union of localizations, not a
trap, and the lifetime bound of V2 remains conditional.  No global
regularity claim is made.
\end{assessment}

% =============================================================================
\section{Integrated V4 conclusion and next acquisition order}
\label{sec:c3v4-conclusion}
% =============================================================================

\begin{theorem}[What V4 proves]
\label{thm:c3v4-conclusion}
Relative to the two legacy cycles and V1--V3: the nucleation theorem
(Theorem~\ref{thm:c3v4-nucleation}) and the two-window localization
(Corollary~\ref{cor:c3v4-two-windows}).  The full Navier--Stokes
stability/global-regularity theorem remains unproved.
\end{theorem}

\begin{proof}
The cited results.
\end{proof}

\begin{programme}[V5 acquisition order]
\label{prog:c3v5}
\begin{enumerate}[label=\textup{(AC\arabic*)},leftmargin=4.8em]
\item The compulsory companion audit: read the current SM and YM
      notes and record their digests and decision.
\item Consolidate V2--V4 as the first chapter of the third cycle
      (isolated clusters and groups lose energy; the single-window
      accounting; nucleation takes two windows; the two-window
      localization), and fix the direction for V6--V10: the
      multi-window question in the form ``does the union over windows
      of the band cells stay bounded'', which is a question about the
      trapped exterior's energy band \([\varepsilon_{\rm th}/4,\varepsilon_{\rm th})\) and
      its persistence.
\end{enumerate}
\end{programme}

\begin{assessment}[Position after V4]
\label{ass:c3v4-position}
Nucleation of activity in the trapped exterior takes at least two
windows; the two-window localization is recorded.  No full stability
or global-regularity claim is made.
\end{assessment}

% =============================================================================
\section*{Version V4 (third cycle) append-only record}
\addcontentsline{toc}{section}{Version V4 (third cycle) append-only record}
% =============================================================================

The complete V3 source of the third cycle was retained before this
continuation.  Its SHA--256 digest was
\begin{center}\scriptsize\ttfamily
1d4c79509c6ccb120ffc19c7a3e0a05efd5c1939deddd1354832a055c5cf0201.
\end{center}
V4 proved the nucleation time and the two-window localization.  No
full regularity claim is made.

% =============================================================================
\section{V5: compulsory companion audit, the first chapter consolidated, and the
direction for V6--V10}
\label{sec:c3v5-entry}
% =============================================================================

This is the fifth version of the third cycle.  The compulsory review of
the two companion programmes is performed first.  V2--V4 are then
consolidated as the cluster-energy chapter, and the direction for
V6--V10 is fixed on the persistence of the band of energies below the
threshold in the trapped exterior.  No global regularity claim is
made.  The next compulsory companion audit is V10.

% =============================================================================
\section{Compulsory V5 review of the current companion notes}
\label{sec:c3v5-companion-audit}
% =============================================================================

The current companion files in \texttt{latex\_notes} were inspected
read-only.  The frozen checkpoints are
\begin{align}
 &\texttt{Renewal\_SM\_Einstein\_Continuum\_Research\_Notes\_V22.tex},
 \notag\\[-2pt]
 &\qquad
 \texttt{6958752e3b5525c205187c20e29fcb67cac316bb3c3509b4e2b10cf1f65ecb9a},
 \label{eq:c3v5-SM-checkpoint}\\
 &\texttt{Renewal\_Yang\_Mills\_Mass\_Gap\_Research\_Notes\_V96.tex},
 \notag\\[-2pt]
 &\qquad
 \texttt{c821c1aa09f8b8316146bf027849f00304eec8181b76abfef50e1da88de682fc}.
 \label{eq:c3v5-YM-checkpoint}
\end{align}
(SM V22 has 9409 lines; YM V96 has 43520 lines.)

\emph{SM V22 (seventh series).}  After closing the pure Higgs-gradient
two-point symbol (V21), the SM programme removed at finite cutoff the
two mixed word classes of odd Yukawa parity and proved mesoscopic
convergence of the curvature--two-Higgs channel to its zero continuum
Clifford trace.

\emph{YM V96.}  The YM programme classified all simple cores of
submacroscopic excess, computed exact degree fibres, and stated the
required upstream simultaneous incidence-source card as an
operator-level identity before norms and support sorting.

\begin{theorem}[V5 companion-audit decision]
\label{thm:c3v5-companion-decision}
No SM V22 parity or convergence constant and no YM V96 core or fibre
bound is imported.  Neither bears on the cluster-energy chapter or on
the gates.
\end{theorem}

\begin{proof}
The SM results concern heat-kernel coefficients of lattice
Dirac--Yukawa operators; the YM results concern cluster expansions of
a lattice clock.  None is a Navier--Stokes statement.
\end{proof}

% =============================================================================
\section{The cluster-energy chapter, consolidated}
\label{sec:c3v5-consolidated}
% =============================================================================

\begin{theorem}[Cluster energy on a window, V2--V4]
\label{thm:c3v5-consolidated}
On a window of the class with constant \(M\) (an intermediate-class
window, or a window with a spike as reference under a type-I approach
with constant \(\bar\Lambda\) in the role of \(M\)), with the second cycle's
threshold, seeds, invasion radius and trap:
\begin{enumerate}[label=\textup{(\arabic*)},leftmargin=3.2em]
\item \emph{Monotonicity.}  Every group of seeds (clusters within
      \(2r+4\) of one another, \(r=R'_{\rm inv}+\rho_{\min}+3\)) has, after
      the capture time, a neighbourhood energy satisfying
      \(E_G'\le-2\nu\int|\nabla U|^2\Phi_G+C_5n_G\varepsilon_{\rm th}^{9/8}\); over the window
      it loses at least \(2\nu m_G\varepsilon'-\eta_G\ell_M\), with \(m_G\) the number of
      separated energy-active centres in the group.
\item \emph{Accounting.}  The losses are dissipated inside the
      neighbourhoods; the single-window accounting gives a bound on
      the number of separated active centres of the form \(C(\nu,M)\),
      not better than \(N_M\); the multiplicity question is not decided.
\item \emph{Lifetime, conditional.}  If a group stays isolated with a
      trapped layer over \(K\) consecutive windows while energy-active,
      then \(K\le K_{\rm win}(\nu,M)\); the hypothesis is not supplied by
      the series because the seed threshold cannot be kept fixed
      across \(K_{\rm win}\) windows.
\item \emph{Nucleation.}  A regular exterior cell starting below
      \(\varepsilon_{\rm th}/4\) and staying at distance at least \(\rho_0\) from the
      non-quiet sets of two consecutive windows becomes non-quiet at
      the earliest at the end of the second window; across two
      windows the activity is localized near the first window's seeds
      and near at most \(64N_M\) band cells.
\end{enumerate}
\end{theorem}

\begin{proof}
(1) Theorems~\ref{thm:c3v2-monotone}, \ref{thm:c3v3-group},
Corollary~\ref{cor:c3v2-two}.  (2) Proposition~\ref{prop:c3v3-accounting}.
(3) Proposition~\ref{prop:c3v2-lifetime}, Firewall~\ref{fw:c3v2-multiwindow}.
(4) Theorem~\ref{thm:c3v4-nucleation}, Corollary~\ref{cor:c3v4-two-windows}.
\end{proof}

% =============================================================================
\section{Direction for V6--V10}
\label{sec:c3v5-direction}
% =============================================================================

\begin{assessment}[Position after five versions of the third cycle]
\label{ass:c3v5-position-five}
The first acquisition (T1 of the second cycle's manifest) closed in
its single-window form and left the multiplicity question open; the
obstruction to any multi-window statement is the band of energies
\([\varepsilon_{\rm th}/4,\varepsilon_{\rm th})\) in the trapped exterior, whose cells become
seeds of the next window without being non-quiet.  The direction for
V6--V10 is therefore the band: V6 asks whether a band cell in the
trapped exterior (energy in the band, enstrophy trapped at
\(C_1\varepsilon_{\rm th}^{1/2}\), far pressure at the quiet level) can keep its
energy through a window or must lose it by its own dissipation (the
interpolation gives \(d\ge y^2/(2e)-C^2y\), and the energy identity
\(e'\le-2\nu y+\ldots\) with \(y\) trapped is not a loss; the question is
whether a lower bound \(y\ge c\,e\) holds for band cells, which would
make the band decay); V7 the count of band cells over consecutive
windows if the band decays; V8 the resulting multi-window lifetime;
V9 consolidation; V10 audit.  No global regularity claim is made.
\end{assessment}

% =============================================================================
\section{Integrated V5 conclusion and next acquisition order}
\label{sec:c3v5-conclusion}
% =============================================================================

\begin{theorem}[What V5 establishes]
\label{thm:c3v5-conclusion}
Relative to the two legacy cycles and V1--V4: the companion-audit
decision (Theorem~\ref{thm:c3v5-companion-decision}) and the
consolidated cluster-energy chapter (Theorem~\ref{thm:c3v5-consolidated}).
The full Navier--Stokes stability/global-regularity theorem remains
unproved.
\end{theorem}

\begin{proof}
The cited results.
\end{proof}

\begin{programme}[V6 acquisition order]
\label{prog:c3v6}
\begin{enumerate}[label=\textup{(AC\arabic*)},leftmargin=4.8em]
\item Band cells.  For a regular cell \(\beta\) of the trapped exterior
      with \(e_\beta\in[\varepsilon_{\rm th}/4,\varepsilon_{\rm th})\), write the local energy identity
      with the dissipation \(-2\nu y_\beta\) retained; determine whether
      \(y_\beta\) is bounded below in terms of \(e_\beta\) (a Poincar\'e-type
      inequality fails without a mean-zero condition, but the far
      pressure and the trapped enstrophy of the neighbours may bound
      the cell's mean velocity), and if not, whether the band can be
      stationary (a cell of constant energy in the band with trapped
      enstrophy and quiet surroundings) on a window.
\item If a stationary band is possible, record it as the obstruction;
      if the band decays at an explicit rate, state the decay and its
      consequence for the seeds of the next window.
\item The next compulsory companion audit is V10.
\end{enumerate}
\end{programme}

\begin{assessment}[Position after V5]
\label{ass:c3v5-position}
The cluster-energy chapter is consolidated; the band of the trapped
exterior is the next object.  No full stability or global-regularity
claim is made.
\end{assessment}

% =============================================================================
\section*{Version V5 (third cycle) append-only record}
\addcontentsline{toc}{section}{Version V5 (third cycle) append-only record}
% =============================================================================

The complete V4 source of the third cycle was retained before this
continuation.  Its SHA--256 digest was
\begin{center}\scriptsize\ttfamily
0c1af4287bbfbac1ead2a832d8886d5fecdd374b6fbef86408c0cb9b1686fc11.
\end{center}
V5 performed the compulsory companion audit and consolidated the
cluster-energy chapter.  No full regularity claim is made.

% =============================================================================
\section{V6: the energy band of the trapped exterior does not decay by itself}
\label{sec:c3v6-entry}
% =============================================================================

Items (AC1) and (AC2) of Programme~\ref{prog:c3v6} are carried out.
(AC1): for a cell of the trapped exterior the local energy identity
bounds the rate of change of the cell's energy in both directions by
an explicit multiple of the square root of the threshold (the trapped
enstrophy) plus smaller terms; no lower bound on the cell's enstrophy
in terms of its energy holds, since a slowly varying large-scale flow
has energy in the band and enstrophy far below it.  (AC2): the band
can be stationary over a window up to the explicit change rate, so
the band does not decay by local considerations, and the persistence
of the band is a question about the large-scale flow of the trapped
exterior, not about the cells.  No global regularity claim is made.
The next compulsory companion audit is V10.

% =============================================================================
\section{The change rate of a trapped cell's energy}
\label{sec:c3v6-rate}
% =============================================================================

\begin{lemma}[Energy change rate in the trapped exterior]
\label{lem:c3v6-rate}
On a window, for \(s\in[-1+\tau_{\rm cap},-1+\ell_M]\) and a cell \(\beta\) such that
every cell within distance \(2\) of \(\beta\) is quiet and trapped
(energy below \(\varepsilon_{\rm th}\), enstrophy at most \(C_1\varepsilon_{\rm th}^{1/2}\)) and
at distance at least \(\rho_{\min}\) from the non-quiet set,
\begin{equation}
 \Bigl|\frac{\dd}{\dd s}e_\beta+2\nu y_\beta\Bigr|\ \le\ C_6\bigl(\nu\varepsilon_{\rm th}+\varepsilon_{\rm th}^{9/8}+\varepsilon_{\rm th}^{3/2}\bigr)\ \le\ C_6'\varepsilon_{\rm th}^{9/8},
 \qquad
 \Bigl|\frac{\dd}{\dd s}e_\beta\Bigr|\ \le\ 2\nu C_1\varepsilon_{\rm th}^{1/2}+C_6'\varepsilon_{\rm th}^{9/8},
 \label{eq:c3v6-rate}
\end{equation}
with \(C_6,C_6'\) depending on \(\nu\) only.  No regularity of \(\beta\) is
needed.
\end{lemma}

\begin{proof}
The local energy identity's non-dissipative terms: the diffusion
cutoff term \(\nu\int|U|^2|\Delta\phi_\beta|\le C\nu\cdot125\varepsilon_{\rm th}\); the cubic and
local-pressure terms, bounded without a sup-norm by
\(C(e^N)^{3/4}(e^N+y^N)^{3/4}\le C(125\varepsilon_{\rm th})^{3/4}(125\varepsilon_{\rm th}+125C_1\varepsilon_{\rm th}^{1/2})^{3/4}\le C\varepsilon_{\rm th}^{9/8}\);
the far-pressure work, bounded by \(C\|\nabla P^{\rm far}\|_{L^\infty(2Q_\beta)}(e^N)^{1/2}\)
with the far-pressure gradient at most \(C(S_4\varepsilon_{\rm th}+N_Ma_M\rho_{\min}^{-4})\le C\varepsilon_{\rm th}\)
by the choice of \(\rho_{\min}\), giving \(C\varepsilon_{\rm th}^{3/2}\).  The dissipation
term is \(-2\nu y_\beta\) with \(0\le y_\beta\le C_1\varepsilon_{\rm th}^{1/2}\).
\end{proof}

\begin{corollary}[Band cells over a window]
\label{cor:c3v6-band}
A cell as in Lemma~\ref{lem:c3v6-rate} with energy in the band
\([\varepsilon_{\rm th}/4,\varepsilon_{\rm th})\) at some time changes its energy by at most
\((2\nu C_1\varepsilon_{\rm th}^{1/2}+C_6'\varepsilon_{\rm th}^{9/8})\ell_M\) over the rest of the window,
and it cannot leave the band from below faster than at the rate
\(2\nu C_1\varepsilon_{\rm th}^{1/2}\).  Since this rate exceeds \(\varepsilon_{\rm th}\) for small
\(\varepsilon_{\rm th}\), the two-sided bound does not by itself keep a band cell
below the threshold; that it stays below the threshold within the
window is the far growth lemma of the second cycle (gain at most
\(\frac38\varepsilon_{\rm th}\), Lemma C2-V61-far), whose bound uses the sign of
the dissipation.
\end{corollary}

\begin{proof}
Integrate \eqref{eq:c3v6-rate}; the last sentence is Lemma C2-V61-far.
\end{proof}

% =============================================================================
\section{No lower bound on the enstrophy of a band cell}
\label{sec:c3v6-counterexample}
% =============================================================================

\begin{counterexample}[A stationary band]
\label{cex:c3v6-stationary}
Let \(U_0\) be a smooth divergence-free field on the rescaled torus of
side \(L\) of the form \(U_0(y)=\varepsilon_{\rm th}^{1/2}v(y/L_0)\) with \(v\) a fixed
smooth mean-zero divergence-free periodic field of unit sup-norm and
\(1\le L_0\le L\) (a large-scale flow).  Its cell energies are at most
\(\varepsilon_{\rm th}\|v\|_\infty^2=\varepsilon_{\rm th}\) and, on the cells where \(|v|\ge\frac12\),
at least \(\varepsilon_{\rm th}/4\): those cells are in the band.  Its cell
enstrophies are at most \(\varepsilon_{\rm th}\|\nabla v\|_\infty^2L_0^{-2}\), which is
below \(c\,\varepsilon_{\rm th}\) for every \(c>0\) once \(L_0\) is large: no inequality
\(y_\beta\ge c\,e_\beta\) holds for band cells.  For the solution with this
datum (in the class with any \(M\ge\varepsilon_{\rm th}\|\nabla v\|_2^2L_0\), on the
window), the cell energies change at rate at most
\(2\nu\varepsilon_{\rm th}\|\nabla v\|_\infty^2L_0^{-2}+C\varepsilon_{\rm th}^{3/2}\) (the nonlinearity and
pressure are of the order of the cube of the amplitude), so the band
is stationary over the window up to \(o(\varepsilon_{\rm th})\) for \(L_0\) large.
\end{counterexample}

\begin{proof}
Direct computation of the cell quantities of \(U_0\) and the local
energy identity with the stated bounds.
\end{proof}

\begin{firewall}[The band is a large-scale object]
\label{fw:c3v6-band}
The band of the trapped exterior is carried, in the extreme case, by
a slowly varying large-scale flow of amplitude \(\varepsilon_{\rm th}^{1/2}\), whose
cells all lie in the band or just below it, with negligible
enstrophy and negligible change over a window; such a flow is
compatible with every bound of the series (the trap bounds the
enstrophy from above, the counts bound the cells above the threshold
only, and the far-pressure bound is by the maximal cell energy, which
is below the threshold).  Hence: (i) the band does not decay by local
considerations (item (AC1) closed negatively); (ii) the band cells of
a window can all persist as band cells of the next window, and can
become seeds without any local event (item (AC2): the stationary
band is the obstruction); (iii) the persistence of the band is a
statement about the large-scale flow, whose only global control is
the total energy \(\lambda^{-1}E_0\), which is unbounded in the rescaled
variables.  The multi-window question is therefore not a question
about cells; the count of band cells (at most \(64N_M\) at any window
time) is the only cell-level statement, and it does not bound their
union over windows.  No global regularity claim is made.
\end{firewall}

% =============================================================================
\section{Integrated V6 conclusion and next acquisition order}
\label{sec:c3v6-conclusion}
% =============================================================================

\begin{theorem}[What V6 proves]
\label{thm:c3v6-conclusion}
Relative to the two legacy cycles and V1--V5: the change-rate lemma
\eqref{eq:c3v6-rate}, its window corollary, and the stationary-band
counterexample (Counterexample~\ref{cex:c3v6-stationary}) with the
firewall.  The full Navier--Stokes stability/global-regularity
theorem remains unproved.
\end{theorem}

\begin{proof}
The cited results.
\end{proof}

\begin{programme}[V7 acquisition order]
\label{prog:c3v7}
\begin{enumerate}[label=\textup{(AC\arabic*)},leftmargin=4.8em]
\item The per-window seed bound.  At every window time the seeds
      (cells at level \(\varepsilon_{\rm th}/4\)) number at most \(64N_M\); the
      non-quiet cells at most \(N_M\).  Determine what this gives across
      \(K\) windows: the union of the localization regions has at most
      \(64N_MK(2R'_{\rm inv}+1)^3\) cells, and a cluster's isolation in window
      \(k\) fails only if a seed of window \(k\) lies within \(R'_{\rm inv}+\rho_{\min}\)
      of its layer; state the resulting alternative for a cluster over
      \(K\) windows: either it is isolated in at least \(K-K'\) of them
      and loses energy accordingly, or seeds approach its layer in at
      least \(K'\) windows.
\item Determine whether the seeds approaching a cluster's layer must
      be band cells of the large-scale flow (energy below the
      threshold) or non-quiet cells, and what the finite propagation
      says about non-quiet cells reaching the layer (they come from
      the cluster itself or from other groups within one localization
      radius per window).
\item The next compulsory companion audit is V10.
\end{enumerate}
\end{programme}

\begin{assessment}[Position after V6]
\label{ass:c3v6-position}
The band of the trapped exterior does not decay by itself; it is a
large-scale object beyond the cell-level tools.  No full stability or
global-regularity claim is made.
\end{assessment}

% =============================================================================
\section*{Version V6 (third cycle) append-only record}
\addcontentsline{toc}{section}{Version V6 (third cycle) append-only record}
% =============================================================================

The complete V5 source of the third cycle was retained before this
continuation.  Its SHA--256 digest was
\begin{center}\scriptsize\ttfamily
f7998bd83666a705b61b02f121a3a122615bea6bd5b7a5d4c3b2ca87d3019a79.
\end{center}
V6 proved the change-rate lemma and recorded the stationary-band
obstruction.  No full regularity claim is made.

% =============================================================================
\section{V7: the per-window seed bound and the isolation alternative}
\label{sec:c3v7-entry}
% =============================================================================

Items (AC1) and (AC2) of Programme~\ref{prog:c3v7} are carried out.
(AC1): at every window time the seeds number at most \(64N_M\) and the
non-quiet cells at most \(N_M\); across \(K\) consecutive windows the
union of the localization regions has an explicit size linear in
\(K\), and a cluster is, in each window, either isolated (its layer at
distance more than the invasion radius plus the trap radius from
every seed of that window) and then loses energy by the monotonicity
identity, or approached by a seed of that window; the alternative is
stated with the count of windows of each kind.  (AC2): a seed
approaching a cluster's layer is either a non-quiet cell, which by
finite propagation comes from a group within one localization radius
per window (the cluster itself or a neighbouring group), or a band
cell of the large-scale flow, which needs no local event; the second
kind is not bounded by the series.  No global regularity claim is
made.  The next compulsory companion audit is V10.

% =============================================================================
\section{Seeds per window and the union over windows}
\label{sec:c3v7-seeds}
% =============================================================================

\begin{proposition}[Seeds per window and their union]
\label{prop:c3v7-seeds}
Under the class constant \(M\) at each of \(K\) consecutive window
times \(t_1<\dots<t_K\):
\begin{enumerate}[label=\textup{(\roman*)},leftmargin=3.2em]
\item at each \(t_k\) the seed set \(\mathcal S'_k\) has at most \(64N_M\) cells
      and the non-quiet set at any time of window \(k\) at most \(N_M\)
      cells, all within \(R'_{\rm inv}\) of \(\mathcal S'_k\);
\item the union over the \(K\) windows of the localization regions
      \(B(\mathcal S'_k,R'_{\rm inv})\) has at most \(64N_MK(2R'_{\rm inv}+1)^3\) cells
      (in the coarsest window's cells, up to the scale factors
      \((1-\ell_M)^{k/2}\), absorbed for \(K\ell_M\le1\));
\item outside that union every cell is quiet in every one of the \(K\)
      windows, and trapped after each window's capture time.
\end{enumerate}
\end{proposition}

\begin{proof}
(i) Proposition C2-V58-count at levels \(\varepsilon_{\rm th}/4\) and \(\varepsilon_{\rm th}\), and
Theorem C2-V66-propagation.  (ii) counting.  (iii) is (i) for each
window and the exterior trap (Theorem C2-V66-propagation's last
statement) applied windowwise.
\end{proof}

% =============================================================================
\section{The isolation alternative}
\label{sec:c3v7-alternative}
% =============================================================================

\begin{theorem}[Isolation alternative over \(K\) windows]
\label{thm:c3v7-alternative}
Let \(A\) be a group of seeds of window \(1\) with neighbourhood \(U_A\)
(radius \(r=R'_{\rm inv}+\rho_{\min}+3\)) and layer \(\partial U_A\), and consider the
\(K\) consecutive windows.  Call window \(k\) \emph{isolating} for \(A\) if
every seed of window \(k\) is either in \(U_A\setminus\partial U_A\) at distance at
least \(r-1\) from \(\partial U_A\) or at distance more than \(R'_{\rm inv}+\rho_{\min}\)
from \(\partial U_A\).  Let \(K_{\rm iso}\) be the number of isolating windows and
\(K'=K-K_{\rm iso}\) the number of the others.  Then:
\begin{enumerate}[label=\textup{(\roman*)},leftmargin=3.2em]
\item in every isolating window in which \(U_A\) contains an
      energy-active centre at level \(\varepsilon_{\rm th}/8\), \(E_A\) decreases by at
      least \(2\nu\varepsilon'-\eta_A\ell_M\) (Corollary~\ref{cor:c3v2-two}), and in
      every isolating window \(E_A\) increases by at most \(\eta_A\ell_M\);
\item in a non-isolating window \(E_A\) may change by up to the window
      bound on the layer's energy flux, at most
      \(C(M,\nu)\,n_A\) (the crude bound of Lemma C2-V61-crude summed over
      the layer), which is not small;
\item consequently, if \(A\) contains an active centre in every
      isolating window, \(K_{\rm iso}\le K_{\rm win}+K'C(M,\nu)n_A/(2\nu\varepsilon'-\eta_A\ell_M)\):
      the lifetime in isolating windows is bounded by the conditional
      lifetime of V2 plus a multiple of the number of non-isolating
      windows.
\end{enumerate}
\end{theorem}

\begin{proof}
(i) is Theorem~\ref{thm:c3v2-monotone} and Corollary~\ref{cor:c3v2-two},
whose hypotheses are exactly the isolating condition (the layer's
cells and their neighbours are at distance more than
\(R'_{\rm inv}+\rho_{\min}\) from every seed, hence quiet and trapped
throughout the window after its capture time).  (ii) In a
non-isolating window the layer may contain non-quiet cells, and the
crude growth lemma bounds each layer cell's energy change by
\(\Gamma(\ell_M)\le C(M,\nu)\).  (iii) Sum the changes: the energy is
nonnegative, starts at most \(N_Ma_M+n_A\varepsilon_{\rm th}\), decreases by at least
\(2\nu\varepsilon'-\eta_A\ell_M\) in each isolating window and increases by at most
\(C(M,\nu)n_A\) in each other window.
\end{proof}

\begin{proposition}[Who approaches the layer]
\label{prop:c3v7-approach}
A seed of window \(k\) within \(R'_{\rm inv}+\rho_{\min}\) of \(\partial U_A\) is one
of:
\begin{enumerate}[label=\textup{(\alph*)},leftmargin=3.2em]
\item a cell that was non-quiet at the end of window \(k-1\); by
      finite propagation it lies within \(R'_{\rm inv}\) of a seed of window
      \(k-1\), so it descends, through a chain of at most \(k-1\) steps of
      length \(R'_{\rm inv}\), from a seed of window \(1\): from \(A\) itself
      (activity that left \(U_A\) through the layer, possible only in a
      non-isolating window or by the layer's own small fluxes) or from
      a group of window \(1\) at distance at most \(kR'_{\rm inv}+r+3\) from
      \(A\);
\item a band cell at the end of window \(k-1\) (energy in
      \([\varepsilon_{\rm th}/4,\varepsilon_{\rm th})\)), which requires no local event
      (Counterexample~\ref{cex:c3v6-stationary}) and is not bounded by
      the series beyond the count \(64N_M\) per window.
\end{enumerate}
\end{proposition}

\begin{proof}
The seeds of window \(k\) are the cells at level \(\varepsilon_{\rm th}/4\) at
\(t_k\); those above \(\varepsilon_{\rm th}\) at \(t_k\) were non-quiet at the end of
window \(k-1\) (continuity) and are localized by Theorem
C2-V66-propagation in window \(k-1\); the rest are band cells.  The
chain is the iteration; a non-quiet cell leaving an isolated layer
would contradict the trap in that window.
\end{proof}

\begin{firewall}[Items (AC1)--(AC2)]
\label{fw:c3v7-AC}
The alternative bounds the number of isolating windows an active
cluster can survive in terms of the number of non-isolating windows,
and the non-isolating windows are caused either by neighbouring groups
approaching (bounded chains) or by band cells (unbounded).  Since a
band cell adjacent to the layer can occur in every window without any
local event, \(K'\) may equal \(K\), and the alternative is empty in the
worst case.  The multiplicity question and the lifetime of a cluster
are therefore decided by the large-scale band, which is outside the
cell-level tools; the chapter closes with this record.  No global
regularity claim is made.
\end{firewall}

% =============================================================================
\section{Integrated V7 conclusion and next acquisition order}
\label{sec:c3v7-conclusion}
% =============================================================================

\begin{theorem}[What V7 proves]
\label{thm:c3v7-conclusion}
Relative to the two legacy cycles and V1--V6: the seeds-per-window
proposition (Proposition~\ref{prop:c3v7-seeds}), the isolation
alternative (Theorem~\ref{thm:c3v7-alternative}), and the
classification of approaching seeds (Proposition~\ref{prop:c3v7-approach}).
The full Navier--Stokes stability/global-regularity theorem remains
unproved.
\end{theorem}

\begin{proof}
The cited results.
\end{proof}

\begin{programme}[V8 acquisition order]
\label{prog:c3v8}
\begin{enumerate}[label=\textup{(AC\arabic*)},leftmargin=4.8em]
\item Go upstream of the band.  The band is a large-scale flow of
      amplitude \(\varepsilon_{\rm th}^{1/2}\); its cells are quiet, so the second
      cycle's trap and far-pressure bounds hold on it with the
      constants of the quiet level, and the only quantity not bounded
      is its extent.  Determine whether the Gaussian ledger at a
      moving centre (the first cycle's \(\Phi_{u,x}(\lambda)\)) evaluated on
      a band cell is below the ledger floor of an energetic cell by a
      fixed factor, so that the first cycle's active-centre count
      \(N_c\) (at most \(N_c\) well-separated centres with Gaussian energy
      at least \(c\)) applies to band cells at a lower level \(c/4\)
      and bounds their number per scale (this is the count \(64N_M\)
      again, in Gaussian form), and whether the log-Lipschitz property
      of the ledger bounds how fast a band cell can rise to the
      threshold across scales (windows), giving a nucleation time in
      log-scale that complements V4.
\item Consolidate V6--V7 as the band chapter, and prepare V10.
\item The next compulsory companion audit is V10.
\end{enumerate}
\end{programme}

\begin{assessment}[Position after V7]
\label{ass:c3v7-position}
The isolation alternative is recorded; the band decides it and is
outside the cell-level tools.  No full stability or global-regularity
claim is made.
\end{assessment}

% =============================================================================
\section*{Version V7 (third cycle) append-only record}
\addcontentsline{toc}{section}{Version V7 (third cycle) append-only record}
% =============================================================================

The complete V6 source of the third cycle was retained before this
continuation.  Its SHA--256 digest was
\begin{center}\scriptsize\ttfamily
b5c31e5abe9fe600b66fdff2972d0557b6c097613b8aed1962f6f68ea172f01c.
\end{center}
V7 proved the isolation alternative and classified the approaching
seeds.  No full regularity claim is made.

% =============================================================================
\section{V8: the band in Gaussian form, and the band chapter consolidated}
\label{sec:c3v8-entry}
% =============================================================================

Items (AC1) and (AC2) of Programme~\ref{prog:c3v8} are carried out.
(AC1): with the threshold lowered by a fixed viscosity factor, the
Gaussian ledger at any centre of the trapped exterior is at most half
the ledger of an energetic cell's centre, so the first cycle's
active-centre count applies to the energetic cells and not to the
band, reproducing the cell count; the log-Lipschitz property of the
ledger bounds the rise of a band centre to the energetic level by a
log-scale interval that is shorter than one window for the series'
constants, so it does not improve the two-window nucleation bound of
V4.  (AC2): V6--V7 are consolidated as the band chapter.  No global
regularity claim is made.  The next compulsory companion audit is
V10.

% =============================================================================
\section{The band in Gaussian form}
\label{sec:c3v8-gaussian}
% =============================================================================

\begin{proposition}[Gaussian ledger of band and energetic centres]
\label{prop:c3v8-gaussian}
Let \(S_G:=\sum_\alpha e^{-(|\alpha|-1)_+^2/(4\nu)}\) be the lattice Gaussian sum and
lower the threshold to \(\varepsilon_{\rm th}\le e^{-1/(4\nu)}\varepsilon_*/(4S_G)\) (allowed:
every statement of the second cycle holds for any smaller threshold
with its constants re-evaluated).  Then, on a window in the rescaled
variables:
\begin{enumerate}[label=\textup{(\roman*)},leftmargin=3.2em]
\item at a centre \(x\) all of whose cells within lattice distance
      \(R\) are quiet (energy below \(\varepsilon_{\rm th}\)),
      \(\Phi_{U,x}(1)\le S_G\varepsilon_{\rm th}+a_Me^{-(R-1)^2/(4\nu)}S_G\), hence at most
      \(e^{-1/(4\nu)}\varepsilon_*/2\) for \(R\ge R_G(\nu,M)\);
\item at the centre of a cell of energy at least \(\varepsilon_*/2\)
      (an energetic cell), \(\Phi_{U,x}(1)\ge e^{-1/(4\nu)}\varepsilon_*/2\);
\item the first cycle's count (Proposition V98-count: at most \(N_c\)
      well-separated centres with Gaussian energy at least \(c\) at one
      scale) with \(c=e^{-1/(4\nu)}\varepsilon_*/2\) counts the energetic cells and
      excludes the quiet centres of (i); it is the cell count \(N_M\) in
      Gaussian form.
\end{enumerate}
\end{proposition}

\begin{proof}
(i) \(\Phi_{U,x}(1)=\int|U|^2e^{-|y-x|^2/4\nu}\le\sum_\alpha e_\alpha\sup_{Q_\alpha}e^{-|y-x|^2/4\nu}\),
split into the cells within \(R\) (energy below \(\varepsilon_{\rm th}\), weights
summing to at most \(S_G\)) and beyond (energy at most \(a_M\), weights at
most \(e^{-(R-1)^2/4\nu}\) each, summing with the Gaussian tail).  (ii)
\(e^{-|y-x|^2/4\nu}\ge e^{-1/4\nu}\) on the unit ball.  (iii) is the
comparison of the levels.
\end{proof}

\begin{proposition}[Rise of a band centre in log-scale]
\label{prop:c3v8-rise}
With the threshold of Proposition~\ref{prop:c3v8-gaussian}, a centre
\(x\) as in (i) at scale \(\lambda\) satisfies, by the first cycle's
log-Lipschitz bound \(|\lambda\,\dd\Phi_{u,x}/\dd\lambda|\le K(\nu,M)\) on window
scales, \(\Phi_{u,x}(\lambda')\le e^{-1/(4\nu)}\varepsilon_*/2+K|\log(\lambda'/\lambda)|\); hence
\(x\) cannot be the centre of an energetic cell at any scale \(\lambda'\)
with \(|\log(\lambda'/\lambda)|<e^{-1/(4\nu)}\varepsilon_*/(2K)\).  In windows (each of
log-length at least \(\ell_M/2\)) this is fewer than
\(e^{-1/(4\nu)}\varepsilon_*/(K\ell_M)\) windows, which is less than one for the
series' constants (\(\varepsilon_*\) exponentially small, \(K\ell_M\) of order
\(\nu^3/M\)); the bound is weaker than the two-window nucleation time of
Theorem~\ref{thm:c3v4-nucleation}.
\end{proposition}

\begin{proof}
Theorem C2-V76-transfer(iv) with the reference \(T_*\) (the first cycle's
form), and the arithmetic.
\end{proof}

\begin{firewall}[Item (AC1)]
\label{fw:c3v8-AC1}
The Gaussian form separates the band from the energetic cells by a
factor two at the price of a viscosity factor in the threshold, and
then reproduces the cell count; the log-Lipschitz rise bound is not
competitive with the far growth lemma.  The band remains a
large-scale object outside the cell-level and Gaussian tools.  No
global regularity claim is made.
\end{firewall}

% =============================================================================
\section{The band chapter, consolidated}
\label{sec:c3v8-consolidated}
% =============================================================================

\begin{theorem}[The band chapter, V6--V8]
\label{thm:c3v8-consolidated}
On a window, after the capture time, in the trapped exterior:
\begin{enumerate}[label=\textup{(\arabic*)},leftmargin=3.2em]
\item every cell's energy changes at rate at most
      \(2\nu C_1\varepsilon_{\rm th}^{1/2}+C_6'\varepsilon_{\rm th}^{9/8}\), and its non-dissipative
      change at rate at most \(C_6'\varepsilon_{\rm th}^{9/8}\) (V6);
\item no lower bound on a cell's enstrophy in terms of its energy
      holds, and a slowly varying large-scale flow of amplitude
      \(\varepsilon_{\rm th}^{1/2}\) is a stationary band compatible with every bound
      of the series (V6);
\item the seeds of each window number at most \(64N_M\) and the union
      of the localization regions over \(K\) windows has at most
      \(64N_MK(2R'_{\rm inv}+1)^3\) cells; a cluster's lifetime in isolating
      windows is bounded by the conditional lifetime plus a multiple of
      the number of non-isolating windows, and non-isolating windows
      are caused by neighbouring groups (bounded chains) or by band
      cells (unbounded) (V7);
\item in Gaussian form the band is separated from the energetic cells
      by a factor two for a lowered threshold, and the ledger's
      log-Lipschitz bound does not improve the nucleation time (V8).
\end{enumerate}
\end{theorem}

\begin{proof}
Lemma~\ref{lem:c3v6-rate}, Counterexample~\ref{cex:c3v6-stationary},
Proposition~\ref{prop:c3v7-seeds}, Theorem~\ref{thm:c3v7-alternative},
Proposition~\ref{prop:c3v7-approach}, Propositions~\ref{prop:c3v8-gaussian}
and \ref{prop:c3v8-rise}.
\end{proof}

% =============================================================================
\section{Integrated V8 conclusion and next acquisition order}
\label{sec:c3v8-conclusion}
% =============================================================================

\begin{theorem}[What V8 proves]
\label{thm:c3v8-conclusion}
Relative to the two legacy cycles and V1--V7: the Gaussian form of the
band (Propositions~\ref{prop:c3v8-gaussian}, \ref{prop:c3v8-rise}) and
the consolidated band chapter (Theorem~\ref{thm:c3v8-consolidated}).
The full Navier--Stokes stability/global-regularity theorem remains
unproved.
\end{theorem}

\begin{proof}
The cited results.
\end{proof}

\begin{programme}[V9 acquisition order]
\label{prog:c3v9}
\begin{enumerate}[label=\textup{(AC\arabic*)},leftmargin=4.8em]
\item Record the position after the first two chapters of the third
      cycle and fix the direction for V11--V15 among the remaining
      manifest items: (T2) the mean inequality at homogeneous points
      from the finite configuration, and (T3) a dissipation-based
      propagation with time-integrated control; state the decision
      criterion (whether the item can produce a statement valid
      beyond the window).
\item Prepare V10.
\end{enumerate}
\end{programme}

\begin{assessment}[Position after V8]
\label{ass:c3v8-position}
The band chapter is consolidated; the band is outside the cell-level
and Gaussian tools.  No full stability or global-regularity claim is
made.
\end{assessment}

% =============================================================================
\section*{Version V8 (third cycle) append-only record}
\addcontentsline{toc}{section}{Version V8 (third cycle) append-only record}
% =============================================================================

The complete V7 source of the third cycle was retained before this
continuation.  Its SHA--256 digest was
\begin{center}\scriptsize\ttfamily
df40d4e166a6093c5b263a1241847ab25810940940e4b6a3bbbb25dfd169e9d6.
\end{center}
V8 recorded the band in Gaussian form and consolidated the band
chapter.  No full regularity claim is made.

% =============================================================================
\section{V9: position after two chapters, and the direction for V11--V15}
\label{sec:c3v9-entry}
% =============================================================================

Items (AC1) and (AC2) of Programme~\ref{prog:c3v9} are carried out.
No global regularity claim is made.  The next compulsory companion
audit is V10.

% =============================================================================
\section{Position and direction}
\label{sec:c3v9-position}
% =============================================================================

\begin{assessment}[Position after the cluster-energy and band chapters]
\label{ass:c3v9-position-two-chapters}
The third cycle's first acquisition (the multiplicity of energy
concentration, T1 of the manifest) produced a single-window
statement, the monotone loss of energy of isolated groups at their
summed dissipation rate, and identified the obstruction to every
multi-window statement: the band of energies below the threshold in
the trapped exterior, which is carried by the large-scale flow, is
stationary over a window, needs no local event to seed the next
window, and is bounded only in number per window (\(64N_M\)) and not
in its union over windows.  The two chapters are therefore closed
with the multiplicity question open, and the honest position is
that the cell-level tools of the second cycle have been exhausted on
the intermediate class at its windows: what they say is said
(finite configuration, finite propagation, exterior trap, monotone
loss of isolated groups, two-window nucleation), and what they do not
say (multi-window continuation, spike localization, cone convergence)
is decided by the large-scale flow, of which the series controls only
the total energy.

\emph{Decision for V11--V15.}  The criterion is whether an item can
produce a statement valid beyond the window.  (T2), the mean
inequality at homogeneous points from the finite-configuration
description, lives on the ancient family, where every time is a
window time: it is a type-I statement and does not need the
between-window stretch.  (T3), a dissipation-based propagation, would
use the time-integrated control of the dissipation and might survive
beyond the window, but the second cycle's V61 firewall shows that
beyond the window the far pressure and the cubic flux have no bound
at all, whatever the quantity propagated; the item is closed before
it starts.  The direction for V11--V15 is therefore (T2): on the
support of a blow-up measure at a homogeneous type-I point, every
profile is a finite configuration with a diffuse far field; V11 writes
the cross term \(\langle N(V),\partial_\sigma V\rangle_G\) in the Lamb--Bernoulli form and
splits it into the configuration's neighbourhood and the far field;
V12 bounds the far-field part by the trapped enstrophy and the quiet
energy of the far cells against the Gaussian weight; V13 examines
whether the configuration part has a sign on average over the
measure, using the tensor and inertia identities on the
neighbourhood; V14 consolidates; V15 audits.
\end{assessment}

% =============================================================================
\section{Integrated V9 conclusion and next acquisition order}
\label{sec:c3v9-conclusion}
% =============================================================================

\begin{theorem}[What V9 establishes]
\label{thm:c3v9-conclusion}
Relative to the two legacy cycles and V1--V8: the position and the
direction (Assessment~\ref{ass:c3v9-position-two-chapters}).  The full
Navier--Stokes stability/global-regularity theorem remains unproved.
\end{theorem}

\begin{proof}
A record.
\end{proof}

\begin{programme}[V10 acquisition order]
\label{prog:c3v10}
\begin{enumerate}[label=\textup{(AC\arabic*)},leftmargin=4.8em]
\item The compulsory companion audit: read the current SM and YM
      notes and record their digests and decision.
\item Consolidate V2--V9 as the first block of the third cycle and
      record the position after ten versions; delivery.
\end{enumerate}
\end{programme}

\begin{assessment}[Position after V9]
\label{ass:c3v9-position}
The direction for V11--V15 is the mean inequality at homogeneous
points from the finite configuration.  No full stability or
global-regularity claim is made.
\end{assessment}

% =============================================================================
\section*{Version V9 (third cycle) append-only record}
\addcontentsline{toc}{section}{Version V9 (third cycle) append-only record}
% =============================================================================

The complete V8 source of the third cycle was retained before this
continuation.  Its SHA--256 digest was
\begin{center}\scriptsize\ttfamily
a7e83e93cb284c897d807cb1b0f62e8fb40776594e4d45d9cb6541214583d5ba.
\end{center}
V9 recorded the position and the direction.  No full regularity claim
is made.

% =============================================================================
\section{V10: compulsory companion audit and the position after ten versions}
\label{sec:c3v10-entry}
% =============================================================================

This is the tenth version of the third cycle.  The compulsory review
of the two companion programmes is performed first.  V2--V9 are then
consolidated as the first block of the cycle and the position after
ten versions is recorded.  No global regularity claim is made.  The
next compulsory companion audit is V15.

% =============================================================================
\section{Compulsory V10 review of the current companion notes}
\label{sec:c3v10-companion-audit}
% =============================================================================

The current companion files in \texttt{latex\_notes} were inspected
read-only.  The frozen checkpoints are
\begin{align}
 &\texttt{Renewal\_SM\_Einstein\_Continuum\_Research\_Notes\_V23.tex},
 \notag\\[-2pt]
 &\qquad
 \texttt{db6f97a4d944e3d2ddccec6411a5986e181128fd0ea155cb879304d11b49126b},
 \label{eq:c3v10-SM-checkpoint}\\
 &\texttt{Renewal\_Yang\_Mills\_Mass\_Gap\_Research\_Notes\_V97.tex},
 \notag\\[-2pt]
 &\qquad
 \texttt{ae22a8125aa0ca1f151ce840a8048a423673ab55686f4253165a802073d0cc8f}.
 \label{eq:c3v10-YM-checkpoint}
\end{align}
(SM V23 has 9738 lines; YM V97 has 43965 lines.)

\emph{SM V23 (seventh series).}  After completing its flat free-symbol
coefficient ledger (V22), the SM programme inserted that ledger into
one covariant radial parametrix for the full improved chiral block,
subtracting the weight-one Higgs mass term with its exact lattice
coefficient before comparing the weight-four gauge, Higgs-gradient and
Higgs-quartic owners.

\emph{YM V97.}  The YM programme gave a Laplace representation of a
saturated forest, recorded why its known factorials and shortcut
identity were insufficient, and defined the unsaturated interpolation
defect of its cluster coefficients.

\begin{theorem}[V10 companion-audit decision]
\label{thm:c3v10-companion-decision}
No SM V23 parametrix constant and no YM V97 representation or defect
bound is imported.  Neither bears on the third cycle's chapters or on
the gates.
\end{theorem}

\begin{proof}
The SM results concern heat-kernel parametrices of lattice
Dirac--Yukawa operators; the YM results concern cluster expansions of
a lattice clock.  None is a Navier--Stokes statement.
\end{proof}

% =============================================================================
\section{The first block of the third cycle, consolidated}
\label{sec:c3v10-consolidated}
% =============================================================================

\begin{theorem}[V2--V9]
\label{thm:c3v10-consolidated}
On a window of the class with constant \(M\), with the second cycle's
threshold, seeds, invasion radius and trap, after the capture time:
\begin{enumerate}[label=\textup{(\arabic*)},leftmargin=3.2em]
\item every group of seeds (clusters within \(2r+4\), \(r=R'_{\rm inv}+\rho_{\min}+3\))
      has a neighbourhood energy that decreases at least at its
      interior dissipation rate up to \(C_5n_G\varepsilon_{\rm th}^{9/8}\), hence by at
      least \(2\nu m_G\varepsilon'-\eta_G\ell_M\) over the window, independently of
      the other groups (V2--V3);
\item the single-window accounting gives no bound on the number of
      groups beyond the count \(N_M\) (V3);
\item a regular exterior cell starting below a quarter of the
      threshold and staying far from the non-quiet sets of two
      consecutive windows becomes non-quiet at the earliest at the end
      of the second window; across two windows the activity is
      localized near the first window's seeds and at most \(64N_M\) band
      cells (V4);
\item every trapped cell's energy changes at rate at most
      \(2\nu C_1\varepsilon_{\rm th}^{1/2}+C_6'\varepsilon_{\rm th}^{9/8}\); a slowly varying large-scale
      flow of amplitude \(\varepsilon_{\rm th}^{1/2}\) is a stationary band compatible
      with every bound of the series, so the band does not decay by
      itself (V6);
\item over \(K\) windows a group's lifetime in isolating windows is
      bounded by the conditional lifetime \(K_{\rm win}\) plus a multiple of
      the number of non-isolating windows, which are caused by
      neighbouring groups (bounded chains) or by band cells
      (unbounded) (V7);
\item in Gaussian form the band is separated from the energetic cells
      by a factor two for a lowered threshold and reproduces the cell
      count; the ledger's log-Lipschitz bound does not improve the
      nucleation time (V8);
\item the cell-level tools are exhausted on the intermediate class at
      its windows; the direction for V11--V15 is the mean inequality
      at homogeneous points from the finite-configuration description
      (V9).
\end{enumerate}
\end{theorem}

\begin{proof}
Theorems~\ref{thm:c3v5-consolidated} and \ref{thm:c3v8-consolidated},
Assessment~\ref{ass:c3v9-position-two-chapters}.
\end{proof}

\begin{assessment}[Position after ten versions of the third cycle]
\label{ass:c3v10-position-ten}
The third cycle has spent its first ten versions on the multiplicity
of energy concentration and closed it as undecidable by the cell-level
tools, with the band of the trapped exterior as the identified
obstruction; the positive residue is the monotone energy loss of
isolated groups on a window and the two-window nucleation time.  The
gates are unchanged from the manifest.  No global regularity claim is
made.
\end{assessment}

% =============================================================================
\section{Integrated V10 conclusion and next acquisition order}
\label{sec:c3v10-conclusion}
% =============================================================================

\begin{theorem}[What V10 establishes]
\label{thm:c3v10-conclusion}
Relative to the two legacy cycles and V1--V9: the companion-audit
decision (Theorem~\ref{thm:c3v10-companion-decision}) and the
consolidation (Theorem~\ref{thm:c3v10-consolidated}).  The full
Navier--Stokes stability/global-regularity theorem remains unproved.
\end{theorem}

\begin{proof}
The cited results.
\end{proof}

\begin{programme}[V11 acquisition order]
\label{prog:c3v11}
\begin{enumerate}[label=\textup{(AC\arabic*)},leftmargin=4.8em]
\item At a homogeneous type-I point, for a profile \(V\) on the support
      of a blow-up measure, split the cross term
      \(\langle N(V),\partial_\sigma V\rangle_G=\int\det(\Omega,V,Y)G+\frac1{2\nu}\int\Pi(z\cdot Y)G\)
      into the configuration neighbourhood \(U\) (the cells within
      \(r_2\) of the seeds, in profile variables) and the far field \(F\);
      bound the far-field part using the diffuse far field
      (Corollary C2-V63-ancient: far cells have energy below
      \(\varepsilon_{\rm th}\) and enstrophy at most \(C_1\varepsilon_{\rm th}^{1/2}\) on the
      sub-window), the class bounds \(\|V\|_\infty\le C_\infty\), \(\|Y\|_G\le C_Y\),
      the Gaussian weight, and the freedom to subtract a constant from
      \(\Pi\) (since \(\int(z\cdot Y)G=0\) by \(\operatorname{div}Y=0\)).
\item State the localization of the anti-alignment: alternative (B)
      gives \(\langle N,Y\rangle_G\le-\underline\delta\) on log-average, so the
      configuration part is at most \(-\underline\delta+O(\varepsilon_{\rm th}^{1/4})\) on
      log-average.
\item The next compulsory companion audit is V15.
\end{enumerate}
\end{programme}

\begin{assessment}[Position after V10]
\label{ass:c3v10-position}
The first block of the third cycle is consolidated; the next block is
the localization of the anti-alignment at homogeneous points.  No
full stability or global-regularity claim is made.
\end{assessment}

% =============================================================================
\section*{Version V10 (third cycle) append-only record}
\addcontentsline{toc}{section}{Version V10 (third cycle) append-only record}
% =============================================================================

The complete V9 source of the third cycle was retained before this
continuation.  Its SHA--256 digest was
\begin{center}\scriptsize\ttfamily
47d1b93933c53ae4a8253ccc202a33406909ef2c4d785fc685b3d9a23873ee1d.
\end{center}
V10 performed the compulsory companion audit and consolidated the
first block.  No full regularity claim is made.

% =============================================================================
\section{V11: the anti-alignment is localized in the configuration neighbourhood}
\label{sec:c3v11-entry}
% =============================================================================

Items (AC1) and (AC2) of Programme~\ref{prog:c3v11} are carried out.
At a homogeneous type-I point, for a profile on the support of a
blow-up measure, the cross term of the second identity is split into
the configuration neighbourhood and the far field; the far-field part
is bounded, using the diffuse far field of the ancient limit, the
class bounds and the Gaussian weight, by a power of the threshold
times the defect norm; hence the anti-alignment forced by alternative
(B) is carried by the configuration neighbourhood alone.  No global
regularity claim is made.  The next compulsory companion audit is
V15.

% =============================================================================
\section{Setting}
\label{sec:c3v11-setting}
% =============================================================================

Let \(x_0\) be a homogeneous type-I singular point, \(\mathfrak m\) a blow-up
measure on the family \(\mathcal F(x_0)\), and \(W\) in its support, with
profile \(V(\cdot,\sigma)\), defect \(Y=\partial_\sigma V\), vorticity \(\Omega\) and
Bernoulli pressure \(\Pi=P_V+\frac12|V|^2\); the class bounds
\(\|V(\sigma)\|_\infty\le C_\infty\), \(\|Y(\sigma)\|_G\le C_Y\), \(\|\nabla V(\sigma)\|_2^2\le C\)
hold for all \(\sigma\) (second cycle, V3 and V7).  In profile variables
the cells of side one are the parabolic cells of the corresponding
time of \(W\); by Corollary C2-V63-ancient, at every \(\sigma\), on the
sub-window following it, the cells at distance at least \(r_2\) from
the seeds of that time (the far field \(F\)) have energy below
\(\varepsilon_{\rm th}\) and, after the capture time, enstrophy at most
\(C_1\varepsilon_{\rm th}^{1/2}\); the configuration neighbourhood \(U:=\R^3\setminus F\)
consists of at most \(N_C(2r_2+1)^3\) cells.  The cross term is, by the
Lamb--Bernoulli form (Theorem C2-V3-cross),
\begin{equation}
 \langle N(V),Y\rangle_G=\underbrace{\int_U\det(\Omega,V,Y)G+\frac1{2\nu}\int_U(\Pi-c)(z\cdot Y)G}_{=:X_U}
 +\underbrace{\int_F\det(\Omega,V,Y)G+\frac1{2\nu}\int_F(\Pi-c)(z\cdot Y)G}_{=:X_F},
 \label{eq:c3v11-split}
\end{equation}
for any constant \(c\), since \(\int(z\cdot Y)G=-2\nu\int Y\cdot\nabla G=2\nu\int(\operatorname{div}Y)G=0\).

% =============================================================================
\section{The far-field part}
\label{sec:c3v11-far}
% =============================================================================

\begin{theorem}[The far field carries no anti-alignment]
\label{thm:c3v11-far}
At every \(\sigma\) in the trapped part of the sub-window, with \(c\) the
value of the far pressure at a fixed point of \(U\),
\begin{equation}
 |X_F|\ \le\ C_7\,\varepsilon_{\rm th}^{1/4}\,C_Y ,
 \label{eq:c3v11-far}
\end{equation}
with \(C_7\) depending on \(\nu\), \(C_\infty\), \(C\) and the counts.
\end{theorem}

\begin{proof}
\emph{Determinant part.}  By Cauchy--Schwarz,
\(|\int_F\det(\Omega,V,Y)G|\le\||\Omega||V|\|_{L^2(G;F)}\|Y\|_G\), and
\(\||\Omega||V|\|_{L^2(G;F)}^2\le C_\infty^2\int_F|\Omega|^2G\le C_\infty^2\sum_{\beta\in F}\sup_{Q_\beta}G\cdot2\int_{Q_\beta}|\nabla V|^2\le2C_\infty^2S_G\,C_1\varepsilon_{\rm th}^{1/2}\),
using \(|\Omega|^2\le2|\nabla V|^2\), the trapped enstrophy of the far cells,
and the Gaussian lattice sum \(S_G\) (each far cell's weight
\(\sup_{Q_\beta}G\le e^{-(|\beta|-1)_+^2/4\nu}\)).

\emph{Pressure part.}  Write \(\Pi-c=(P^{\rm loc}_\beta)+(P^{\rm far}_\beta-c)+\frac12|V|^2\)
on each far cell \(\beta\), with the second cycle's decomposition.  The
far-pressure gradient on a far cell is at most
\(C(S_4\varepsilon_{\rm th}+N_Ca_C\rho_{\min}^{-4})\le C\varepsilon_{\rm th}\) (Lemma C2-V58-far with the
current non-quiet set inside \(U\)), so \(|P^{\rm far}_\beta-c|\le C\varepsilon_{\rm th}(1+|z|)\)
on \(Q_\beta\) by integrating the gradient along a path from the fixed
point of \(U\) (the far pressure is smooth, and its gradient bound
holds on every quiet cell along the path; the path may pass through
\(U\) where the bound is instead \(Ca_C\), over at most \(N_C(2r_2+1)\)
cells, absorbed into the constant).  The local pressure satisfies
\(\|P^{\rm loc}_\beta\|_{L^2(Q_\beta)}\le C\|V\|_{L^4(4Q_\beta)}^2\le C(e^N+y^N)\le C\varepsilon_{\rm th}^{1/2}\)
(Calder\'on--Zygmund and the local Sobolev inequality with the quiet
energies and trapped enstrophies), and \(\frac12|V|^2\le\frac12C_\infty|V|\)
with \(\||V|\|_{L^2(Q_\beta)}\le\varepsilon_{\rm th}^{1/2}\).  Hence, on far cells,
\(\|(\Pi-c)\|_{L^2(Q_\beta)}\le C\varepsilon_{\rm th}^{1/2}(1+|\beta|)\) after absorbing the
\(\varepsilon_{\rm th}\)-terms, and
\(|\frac1{2\nu}\int_F(\Pi-c)(z\cdot Y)G|\le\frac1{2\nu}\||z|(\Pi-c)G^{1/2}\|_{L^2(F)}\|Y\|_G\)
with \(\||z|(\Pi-c)G^{1/2}\|_{L^2(F)}^2\le C\varepsilon_{\rm th}\sum_\beta(1+|\beta|)^4e^{-(|\beta|-1)_+^2/4\nu}\le C\varepsilon_{\rm th}\).
Collecting, \(|X_F|\le(C\varepsilon_{\rm th}^{1/4}+C\varepsilon_{\rm th}^{1/2})C_Y\).
\end{proof}

% =============================================================================
\section{Localization of the anti-alignment}
\label{sec:c3v11-localized}
% =============================================================================

\begin{theorem}[The configuration carries the anti-alignment]
\label{thm:c3v11-localized}
At a homogeneous type-I point in alternative (B) (the only case, by
(E9)), for every \(W\) in the support of a blow-up measure and every
window of log-scale \([\sigma,\sigma+L]\) on which the trapped sub-windows
cover a fraction \(1-\theta_{\rm cap}\) of the length (\(\theta_{\rm cap}\le\tau_{\rm cap}/\ell_C\)),
\begin{equation}
 \frac1L\int_\sigma^{\sigma+L}X_U\,\dd\sigma'\ \le\ -\underline\delta+C_7\varepsilon_{\rm th}^{1/4}C_Y+\theta_{\rm cap}C_N^{1/2}C_Y+o(1)\qquad(L\to\infty),
 \label{eq:c3v11-localized}
\end{equation}
and against the measure itself, \(\int X_U\,\dd\mathfrak m\le-\underline\delta+C_7\varepsilon_{\rm th}^{1/4}C_Y+\theta_{\rm cap}C_N^{1/2}C_Y\).
The anti-alignment of the nonlinearity with the defect is carried by
the configuration neighbourhood: the far field contributes at most
\(O(\varepsilon_{\rm th}^{1/4})\), and the capture phases at most \(O(\theta_{\rm cap})\).
\end{theorem}

\begin{proof}
Alternative (B) gives \(\frac1L\int\langle N,Y\rangle_G\le-\underline\delta+o(1)\) (Theorem
C2-V8-three-signs, first sign); split by \eqref{eq:c3v11-split}; on the
trapped sub-windows use \eqref{eq:c3v11-far}; on the capture phases
bound \(|\langle N,Y\rangle_G|\le C_N^{1/2}C_Y\) crudely; the measure form is the
ergodic average.
\end{proof}

\begin{assessment}[Items (AC1)--(AC2)]
\label{ass:c3v11-AC}
The localization says that on the ancient family at a homogeneous
point, the mechanism that prevents self-similarity (the defect being
fed by the nonlinearity at rate at least \(\underline\delta\)) operates
inside a bounded neighbourhood of at most \(N_C\) energetic parabolic
cells, and that the diffuse far field is a spectator.  The threshold
\(\varepsilon_{\rm th}\) is at the series' disposal, so the far-field term can be
made smaller than any fraction of \(\underline\delta\) provided
\(\underline\delta\) is known to be positive, which alternative (B) gives
without an explicit value.  What remains for V12--V13 is whether the
configuration part \(X_U\) can be bounded below on average by the
identities on the neighbourhood, which would contradict
\eqref{eq:c3v11-localized}.  No global regularity claim is made.
\end{assessment}

% =============================================================================
\section{Integrated V11 conclusion and next acquisition order}
\label{sec:c3v11-conclusion}
% =============================================================================

\begin{theorem}[What V11 proves]
\label{thm:c3v11-conclusion}
Relative to the two legacy cycles and V1--V10: the far-field bound
\eqref{eq:c3v11-far} and the localization of the anti-alignment
\eqref{eq:c3v11-localized}.  The full Navier--Stokes
stability/global-regularity theorem remains unproved.
\end{theorem}

\begin{proof}
The cited results.
\end{proof}

\begin{programme}[V12 acquisition order]
\label{prog:c3v12}
\begin{enumerate}[label=\textup{(AC\arabic*)},leftmargin=4.8em]
\item The configuration part.  On the neighbourhood \(U\), write
      \(X_U\) with the class bounds: \(|X_U|\le\||\Omega||V|\|_{L^2(G;U)}\|Y\|_{L^2(G;U)}+\frac1{2\nu}\||z|(\Pi-c)\|_{L^2(G;U)}\|Y\|_{L^2(G;U)}\);
      determine whether \eqref{eq:c3v11-localized} forces a lower bound
      on the defect's Gaussian norm on \(U\), \(\|Y\|_{L^2(G;U)}\ge c\underline\delta\),
      and on the neighbourhood's vorticity--velocity product, so that
      the configuration must carry both defect and vorticity at
      explicit levels.
\item Determine whether the tensor identity restricted to \(U\)
      (with boundary fluxes through the trapped layer bounded as in
      V2) gives a signed row for \(X_U\); record the outcome.
\item The next compulsory companion audit is V15.
\end{enumerate}
\end{programme}

\begin{assessment}[Position after V11]
\label{ass:c3v11-position}
The anti-alignment at a homogeneous point is localized in the
configuration neighbourhood.  No full stability or global-regularity
claim is made.
\end{assessment}

% =============================================================================
\section*{Version V11 (third cycle) append-only record}
\addcontentsline{toc}{section}{Version V11 (third cycle) append-only record}
% =============================================================================

The complete V10 source of the third cycle was retained before this
continuation.  Its SHA--256 digest was
\begin{center}\scriptsize\ttfamily
7e08e57a5cc0935183b36590fcf17faf483f3fcafc9510d98f925dc79500ac68.
\end{center}
V11 localized the anti-alignment in the configuration neighbourhood.
No full regularity claim is made.

% =============================================================================
\section{V12: what the configuration must carry, and the localized tensor identity}
\label{sec:c3v12-entry}
% =============================================================================

Items (AC1) and (AC2) of Programme~\ref{prog:c3v12} are carried out.
(AC1): the localized anti-alignment forces, on log-average, explicit
lower bounds in terms of the defect rate on the defect's Gaussian norm
over the configuration neighbourhood and on the product of vorticity
and velocity there or on the Bernoulli-pressure work there: the
configuration carries defect and, in one of the two forms, geometric
structure at levels bounded below by the defect rate over class
constants.  (AC2): the tensor identity restricted to the configuration
neighbourhood holds with layer fluxes bounded by powers of the
threshold, giving the localized mean stress identity; it contains no
sign for the cross term, and the item is closed with that record.  No
global regularity claim is made.  The next compulsory companion audit
is V15.

% =============================================================================
\section{Lower bounds on the configuration}
\label{sec:c3v12-lower}
% =============================================================================

\begin{theorem}[The configuration carries the defect]
\label{thm:c3v12-lower}
In the setting of V11, with the threshold chosen so that
\(C_7\varepsilon_{\rm th}^{1/4}C_Y+\theta_{\rm cap}C_N^{1/2}C_Y\le\underline\delta/4\), on log-average
over the trapped sub-windows (and against the measure):
\begin{enumerate}[label=\textup{(\roman*)},leftmargin=3.2em]
\item \(\bigl\langle\|Y\|_{L^2(G;U)}\bigr\rangle\ \ge\ \dfrac{3\underline\delta/4}{C_{\Omega V}+C_\Pi}\), where
      \(C_{\Omega V}:=\sup\||\Omega||V|\|_{L^2(G;U)}\le C_\infty(2C)^{1/2}\) and
      \(C_\Pi:=\frac1{2\nu}\sup\||z|(\Pi-c)\|_{L^2(G;U)}\) are class constants
      (the second by the class \(L^\infty\) bound on \(V\), the
      Calder\'on--Zygmund bound on \(P_V\) in \(L^2(U)\), and the
      boundedness of \(|z|\) on \(U\));
\item at least one of
      \(\bigl\langle\||\Omega||V|\|_{L^2(G;U)}\|Y\|_{L^2(G;U)}\bigr\rangle\ge\frac38\underline\delta\) and
      \(\bigl\langle\frac1{2\nu}\||z|(\Pi-c)\|_{L^2(G;U)}\|Y\|_{L^2(G;U)}\bigr\rangle\ge\frac38\underline\delta\)
      holds: the configuration carries either a vorticity--velocity
      product or a Bernoulli-pressure work at a level bounded below
      by the defect rate over the defect's Gaussian norm.
\end{enumerate}
\end{theorem}

\begin{proof}
From \eqref{eq:c3v11-localized} with the threshold choice,
\(\langle-X_U\rangle\ge\frac34\underline\delta\); by Cauchy--Schwarz on \(U\),
\(-X_U\le\||\Omega||V|\|_{L^2(G;U)}\|Y\|_{L^2(G;U)}+\frac1{2\nu}\||z|(\Pi-c)\|_{L^2(G;U)}\|Y\|_{L^2(G;U)}\);
(i) bounds the two factors by the class constants; (ii) is the
pigeonhole on the two terms.  The constants: \(\||\Omega||V|\|_{L^2(G;U)}^2\le C_\infty^2\int_U|\Omega|^2G\le2C_\infty^2\|\nabla V\|_2^2\le2C_\infty^2C\);
\(\||z|(\Pi-c)\|_{L^2(G;U)}\le\operatorname{diam}(U)(\|P_V-c\|_{L^2(U)}+\frac12\|V\|_\infty\|V\|_{L^2(U)})\)
with \(\|P_V\|_{L^2(U)}\le C_{\rm CZ}\|V\|_{L^4}^2\) bounded by the class.
\end{proof}

\begin{remark}[What the lower bounds say]
\label{rem:c3v12-lower}
Alternative (B) already says that the defect has Gaussian norm at
least \(\underline\delta^{1/2}\) on log-average (its definition); (i) adds that
this norm sits, up to a class factor, on the configuration
neighbourhood, and (ii) that the anti-alignment is realized by the
Lamb term or by the Bernoulli-pressure term there.  Neither is a
contradiction: a finite configuration of vortical structures with a
defect of the same size is exactly what an ancient non-self-similar
solution would look like.
\end{remark}

% =============================================================================
\section{The localized tensor identity}
\label{sec:c3v12-tensor}
% =============================================================================

\begin{proposition}[Tensor identity on the configuration neighbourhood]
\label{prop:c3v12-tensor}
Let \(\Phi_U\) be the cutoff of \(U\) (the sum of the cell cutoffs of
\(U\)), \(T^U_{ij}:=\int V_iV_jG\Phi_U\), \(D^U_{ij}:=2\nu\int\nabla V_i\cdot\nabla V_jG\Phi_U\),
\(S^U_{ij}:=\int(N_iV_j+N_jV_i)G\Phi_U\).  Then on the trapped
sub-windows
\begin{equation}
 \frac{\dd}{\dd\sigma}T^U=-D^U-T^U-S^U+B^U,\qquad
 |B^U_{ij}|\le C_8\,n_U\,\varepsilon_{\rm th}^{1/2},
 \label{eq:c3v12-tensor}
\end{equation}
where \(B^U\) collects the terms with \(\nabla\Phi_U\) (supported in the
layer of \(U\), of \(n_U\) cells), and \(C_8\) depends on \(\nu\) and the
class constants.  Consequently, against the measure,
\(\int S^U\,\dd\mathfrak m=-\int(D^U+T^U)\,\dd\mathfrak m+O(n_U\varepsilon_{\rm th}^{1/2})\preceq O(n_U\varepsilon_{\rm th}^{1/2})\):
the mean stress on the configuration is negative semidefinite up to
the layer term.
\end{proposition}

\begin{proof}
Multiply the profile equation by \(V_jG\Phi_U\) and symmetrize; the
\(L_-\)-terms give \(-D^U-T^U\) plus the terms where a derivative falls
on \(\Phi_U\): \(\nu\int(\nabla V_i\cdot\nabla\Phi_U)V_jG\) and its symmetric
partner, bounded by \(\nu\|\nabla V\|_{L^2(\rm layer)}\|V\|_{L^2(\rm layer)}\le\nu n_U(C_1\varepsilon_{\rm th}^{1/2}\varepsilon_{\rm th})^{1/2}\)
(trapped enstrophy and quiet energy on the layer, weights at most
one); the nonlinear terms give \(-S^U\).  The mean form is the
invariance of the measure applied to the bounded continuous
functional \(T^U\).  The largest layer term is of order \(\varepsilon_{\rm th}^{3/4}\),
recorded as \(\varepsilon_{\rm th}^{1/2}\) for simplicity.
\end{proof}

\begin{firewall}[Item (AC2): no signed row for the cross term]
\label{fw:c3v12-AC2}
The localized tensor identity constrains the mean stress
\(\int S^U\), a quadratic functional of \(V\) against \(N(V)\); the cross
term \(X_U\) is a functional of \(V\) against \(N(V)\) and the defect
\(Y=L_-V-N(V)\).  Writing \(X_U=\langle N,L_-V\rangle_{G;U}-\langle N,N\rangle_{G;U}\) up to
layer terms, the second part is \(-\|N\|_{L^2(G;U)}^2\le0\) and the first
has no sign; the tensor identity bounds neither.  The item is closed
with this record: the localized identities reproduce the second
cycle's ledger on the neighbourhood, and the ledger has no
incompatible row (C2-V22).  No global regularity claim is made.
\end{firewall}

% =============================================================================
\section{Integrated V12 conclusion and next acquisition order}
\label{sec:c3v12-conclusion}
% =============================================================================

\begin{theorem}[What V12 proves]
\label{thm:c3v12-conclusion}
Relative to the two legacy cycles and V1--V11: the lower bounds on the
configuration (Theorem~\ref{thm:c3v12-lower}) and the localized tensor
identity (Proposition~\ref{prop:c3v12-tensor}).  The full
Navier--Stokes stability/global-regularity theorem remains unproved.
\end{theorem}

\begin{proof}
The cited results.
\end{proof}

\begin{programme}[V13 acquisition order]
\label{prog:c3v13}
\begin{enumerate}[label=\textup{(AC\arabic*)},leftmargin=4.8em]
\item The sign of \(X_U\) on average.  Using the decomposition
      \(X_U=\langle N,L_-V\rangle_{G;U}-\|N\|_{L^2(G;U)}^2+\)(layer terms), determine
      whether \(\langle N,L_-V\rangle_{G;U}\) has a sign on average against the
      measure, using the Lamb--Bernoulli form with \(L_-V\) in place of
      the defect (the push of the second cycle's V32, localized), and
      the second cycle's record that the push has no sign
      (C2-V33--V34); record the outcome.
\item If no sign, record the final form of the type-I gate after
      the localization: a single mean inequality on the configuration
      neighbourhood.
\item The next compulsory companion audit is V15.
\end{enumerate}
\end{programme}

\begin{assessment}[Position after V12]
\label{ass:c3v12-position}
The configuration carries the defect and the anti-alignment at levels
bounded below by the defect rate; the localized tensor identity has
no sign for the cross term.  No full stability or global-regularity
claim is made.
\end{assessment}

% =============================================================================
\section*{Version V12 (third cycle) append-only record}
\addcontentsline{toc}{section}{Version V12 (third cycle) append-only record}
% =============================================================================

The complete V11 source of the third cycle was retained before this
continuation.  Its SHA--256 digest was
\begin{center}\scriptsize\ttfamily
60d01a809c730fbd897c88c6683d1d1aaf5c0807f6c9c7a87894d8ba03502179.
\end{center}
V12 proved the lower bounds on the configuration and the localized
tensor identity.  No full regularity claim is made.

% =============================================================================
\section{V13: the sign of the configuration part, and the gate after localization}
\label{sec:c3v13-entry}
% =============================================================================

Items (AC1) and (AC2) of Programme~\ref{prog:c3v13} are carried out.
(AC1): the configuration part of the cross term decomposes into the
localized push \(\langle N,L_-V\rangle_{G;U}\) minus the localized squared
nonlinearity and layer terms; the localized push is, by the
Lamb--Bernoulli form, the Gaussian average over the configuration of
the oriented volume spanned by vorticity, velocity and the drift
image \(L_-V\), plus the Bernoulli-pressure work against the radial
component of \(L_-V\); neither has a sign on the family, as the second
cycle recorded for the push, and the localization does not create
one.  (AC2): the type-I gate is restated after the localization as a
single mean inequality on the configuration neighbourhood.  No global
regularity claim is made.  The next compulsory companion audit is
V15.

% =============================================================================
\section{The configuration part decomposed}
\label{sec:c3v13-decomposed}
% =============================================================================

\begin{proposition}[Decomposition of the configuration part]
\label{prop:c3v13-decomposition}
On the trapped sub-windows, with \(U\) enlarged if necessary to contain
the cells within the Gaussian radius \(R_G\) of the seeds,
\begin{equation}
 X_U=\langle N,L_-V\rangle_{G;U}-\|N\|_{L^2(G;U)}^2+B_X,\qquad
 \langle N,L_-V\rangle_{G;U}=\int_U\det(\Omega,V,L_-V)G+\frac1{2\nu}\int_U(\Pi-c)(z\cdot L_-V)G+B_X',
 \label{eq:c3v13-decomposition}
\end{equation}
with layer terms \(|B_X|,|B_X'|\le C_9n_U\varepsilon_{\rm th}^{1/4}\), and
\(\|N\|_{L^2(G;U)}^2\ge\|N\|_G^2-C_Ne^{-R_G^2/(4\nu)}\).  Pointwise,
\(\|Y\|_G^2=\|L_-V\|_G^2-2\langle L_-V,N\rangle_G+\|N\|_G^2\).
\end{proposition}

\begin{proof}
\(Y=L_-V-N\) gives \(\langle N,Y\rangle_{G;U}=\langle N,L_-V\rangle_{G;U}-\|N\|_{L^2(G;U)}^2\), the
restriction to \(U\) of the inner products differing from the
integrals against \(G\Phi_U\) by layer terms bounded as in V11--V12;
the Lamb--Bernoulli form with \(L_-V\) (divergence-free, by C2-V8) in
place of \(Y\) is Proposition C2-V32-push localized, the constant \(c\)
being allowed since \(\int(z\cdot L_-V)G=0\).  The far part of \(\|N\|_G^2\)
is bounded by the class bound \(\|N\|_G^2\le C_N\) and the Gaussian tail
beyond the radius \(R_G\) (the far cells' palinstrophy is not trapped,
so the finer bound of V11 is not used here).  The pointwise identity
is the expansion of \(\|L_-V-N\|_G^2\).
\end{proof}

\begin{theorem}[No sign for the localized push]
\label{thm:c3v13-no-sign}
Neither \(\int_U\det(\Omega,V,L_-V)G\) nor \(\frac1{2\nu}\int_U(\Pi-c)(z\cdot L_-V)G\)
has a sign on the family: the second cycle's records (the push has
no sign, C2-V33; the variance floor is the only signed quantity in
the Rayleigh identity, C2-V34) hold for the full-space quantities,
and the far field's contribution to each is at most
\(C\varepsilon_{\rm th}^{1/4}\) by the argument of V11 (with \(L_-V\) in place of
\(Y\), using \(\|L_-V\|_G\le\|Y\|_G+\|N\|_G\)); a sign on \(U\) would therefore
be a sign on the whole space up to \(O(\varepsilon_{\rm th}^{1/4})\), which the
second cycle's examples (the plane wave and the columnar swirl with
zero flux, and the profiles with either sign of the push) exclude for
any fixed threshold.
\end{theorem}

\begin{proof}
The far-field bound is the proof of Theorem~\ref{thm:c3v11-far} with
\(L_-V\) for \(Y\); the rest is the record of the second cycle
(Propositions C2-V32-push, C2-V33-mean, Theorem C2-V34-variance and
the examples of C1-V88).
\end{proof}

% =============================================================================
\section{The gate after localization}
\label{sec:c3v13-gate}
% =============================================================================

\begin{assessment}[Gate G1 after the localization]
\label{ass:c3v13-gate}
At a homogeneous type-I point, a nonzero ancient limit exists if and
only if (E9 and C2-V9) some blow-up measure \(\mathfrak m\) on the family
satisfies alternative (B); after V11--V13 this is equivalent to
\begin{equation}
 \int\Bigl[\int_U\det(\Omega,V,Y)G+\frac1{2\nu}\int_U(\Pi-c)(z\cdot Y)G\Bigr]\dd\mathfrak m\ \le\ -\underline\delta+O(\varepsilon_{\rm th}^{1/4}),
 \label{eq:c3v13-gate}
\end{equation}
with \(U\) the configuration neighbourhood (at most \(N_C(2r_2+1)^3\)
parabolic cells) and \(\underline\delta>0\).  Gate G1 closes if
\eqref{eq:c3v13-gate} is shown impossible for every blow-up measure,
that is, if on the configuration the Gaussian average of the oriented
volume spanned by vorticity, velocity and defect, plus the
Bernoulli-pressure work against the radial defect, cannot be
negative on average at a rate bounded below by the defect rate.  The
localization reduces the gate from the whole space to a bounded
neighbourhood with an explicit cell count; it does not change its
nature.  No global regularity claim is made.
\end{assessment}

% =============================================================================
\section{Integrated V13 conclusion and next acquisition order}
\label{sec:c3v13-conclusion}
% =============================================================================

\begin{theorem}[What V13 proves]
\label{thm:c3v13-conclusion}
Relative to the two legacy cycles and V1--V12: the decomposition
\eqref{eq:c3v13-decomposition}, the no-sign theorem
(Theorem~\ref{thm:c3v13-no-sign}), and the gate after localization
\eqref{eq:c3v13-gate}.  The full Navier--Stokes
stability/global-regularity theorem remains unproved.
\end{theorem}

\begin{proof}
The cited results.
\end{proof}

\begin{programme}[V14 acquisition order]
\label{prog:c3v14}
\begin{enumerate}[label=\textup{(AC\arabic*)},leftmargin=4.8em]
\item Consolidate V11--V13 as the localization chapter, and record
      the position before the V15 audit.
\item Fix the direction for V16--V20: with the gate localized to a
      bounded neighbourhood of at most \(N_C(2r_2+1)^3\) cells, the
      natural next question is the structure of the configuration
      itself on the ancient family: whether the number of energetic
      cells in the support of a blow-up measure is one (a single
      cluster) for a homogeneous point, using the homogeneity (the
      Gaussian ledger at the fixed centre bounded below at every
      scale) against the finite propagation (the configuration moves
      by at most \(\rho_0\) per \(\tau_c\)), so that the configuration at
      a homogeneous point cannot drift away from the centre.
\end{enumerate}
\end{programme}

\begin{assessment}[Position after V13]
\label{ass:c3v13-position}
The localized push has no sign; the type-I gate is a single mean
inequality on the configuration neighbourhood.  No full stability or
global-regularity claim is made.
\end{assessment}

% =============================================================================
\section*{Version V13 (third cycle) append-only record}
\addcontentsline{toc}{section}{Version V13 (third cycle) append-only record}
% =============================================================================

The complete V12 source of the third cycle was retained before this
continuation.  Its SHA--256 digest was
\begin{center}\scriptsize\ttfamily
91ab26f7e2274d658e77dd5785ee1529f863b3135f37f67fde2737ff830c7125.
\end{center}
V13 recorded the no-sign theorem and the gate after localization.  No
full regularity claim is made.

% =============================================================================
\section{V14: the localization chapter consolidated, and the direction for V16--V20}
\label{sec:c3v14-entry}
% =============================================================================

Items (AC1) and (AC2) of Programme~\ref{prog:c3v14} are carried out.
No global regularity claim is made.  The next compulsory companion
audit is V15.

% =============================================================================
\section{The localization chapter}
\label{sec:c3v14-consolidated}
% =============================================================================

\begin{theorem}[Localization at homogeneous points, V11--V13]
\label{thm:c3v14-consolidated}
At a homogeneous type-I singular point, for every blow-up measure
\(\mathfrak m\) on the family and every profile in its support, with \(U\)
the configuration neighbourhood (at most \(N_C(2r_2+1)^3\) parabolic
cells around the energetic seeds of the time) and \(F\) the far field:
\begin{enumerate}[label=\textup{(\arabic*)},leftmargin=3.2em]
\item the far-field part of the cross term satisfies
      \(|X_F|\le C_7\varepsilon_{\rm th}^{1/4}C_Y\) on the trapped sub-windows, and the
      configuration part satisfies \(\int X_U\,\dd\mathfrak m\le-\underline\delta+O(\varepsilon_{\rm th}^{1/4})+O(\theta_{\rm cap})\)
      (V11);
\item the configuration carries the defect,
      \(\langle\|Y\|_{L^2(G;U)}\rangle\ge\frac34\underline\delta/(C_{\Omega V}+C_\Pi)\), and either a
      vorticity--velocity product or a Bernoulli-pressure work at a
      level \(\frac38\underline\delta\) over the defect norm; the tensor identity
      holds on \(U\) with layer terms \(O(n_U\varepsilon_{\rm th}^{1/2})\) and gives the
      localized mean stress identity, with no sign for the cross term
      (V12);
\item the configuration part decomposes into the localized push
      minus the localized squared nonlinearity; the localized push has
      no sign on the family; the gate G1 is the single mean inequality
      \eqref{eq:c3v13-gate} on \(U\) (V13).
\end{enumerate}
\end{theorem}

\begin{proof}
Theorems~\ref{thm:c3v11-far}, \ref{thm:c3v11-localized},
\ref{thm:c3v12-lower}, \ref{thm:c3v13-no-sign}; Propositions~\ref{prop:c3v12-tensor},
\ref{prop:c3v13-decomposition}; Assessment~\ref{ass:c3v13-gate}.
\end{proof}

% =============================================================================
\section{Direction for V16--V20}
\label{sec:c3v14-direction}
% =============================================================================

\begin{assessment}[Position before the V15 audit, and direction]
\label{ass:c3v14-position}
The third cycle has reduced the type-I gate at homogeneous points to
a mean inequality on a bounded neighbourhood with an explicit cell
count, and has shown that the neighbourhood must carry the defect and
its anti-alignment.  The direction for V16--V20 is the structure of
the configuration at a homogeneous point: V16 asks whether the
homogeneity (the Gaussian ledger at the fixed centre \(x_0\) bounded
below at every scale, C1-V82 and C2-V1 pinching) together with the
finite propagation (the configuration moves by at most \(\rho_0\) cells
per time \(\tau_c\) in the profile variables, where the drift
\(-\frac12z\cdot\nabla\) carries every fixed point outward at unit rate in
log-scale) forces the configuration to remain within a bounded
distance of the centre at every scale, that is, whether a homogeneous
point's configuration is a single group around the centre; V17 the
consequence for the count (a single group of at most \(N_C\) cells) and
for the cross term (the neighbourhood \(U\) is a fixed ball around the
centre); V18 whether the moment ledger's centre-mode and momentum
identities (C2-V16--V17), which concern the Gaussian mean of the
profile, gain content when the configuration is a single group at the
centre (the Gaussian mean is then the group's momentum); V19
consolidation; V20 audit.  No global regularity claim is made.
\end{assessment}

% =============================================================================
\section{Integrated V14 conclusion and next acquisition order}
\label{sec:c3v14-conclusion}
% =============================================================================

\begin{theorem}[What V14 establishes]
\label{thm:c3v14-conclusion}
Relative to the two legacy cycles and V1--V13: the consolidated
localization chapter (Theorem~\ref{thm:c3v14-consolidated}) and the
direction.  The full Navier--Stokes stability/global-regularity
theorem remains unproved.
\end{theorem}

\begin{proof}
The cited results.
\end{proof}

\begin{programme}[V15 acquisition order]
\label{prog:c3v15}
\begin{enumerate}[label=\textup{(AC\arabic*)},leftmargin=4.8em]
\item The compulsory companion audit: read the current SM and YM
      notes and record their digests and decision.
\item Record the position after fifteen versions; delivery.
\end{enumerate}
\end{programme}

\begin{assessment}[Position after V14]
\label{ass:c3v14-position-final}
The localization chapter is consolidated; the direction is the
single-group question at homogeneous points.  No full stability or
global-regularity claim is made.
\end{assessment}

% =============================================================================
\section*{Version V14 (third cycle) append-only record}
\addcontentsline{toc}{section}{Version V14 (third cycle) append-only record}
% =============================================================================

The complete V13 source of the third cycle was retained before this
continuation.  Its SHA--256 digest was
\begin{center}\scriptsize\ttfamily
febd6f43c33e6e1e4cf9519614a221b4684ddf678b03e2e7b5f2615273be227b.
\end{center}
V14 consolidated the localization chapter.  No full regularity claim
is made.

% =============================================================================
\section{V15: compulsory companion audit and the position after fifteen versions}
\label{sec:c3v15-entry}
% =============================================================================

This is the fifteenth version of the third cycle.  The compulsory
review of the two companion programmes is performed first, and the
position after fifteen versions is recorded.  No global regularity
claim is made.  The next compulsory companion audit is V20.

% =============================================================================
\section{Compulsory V15 review of the current companion notes}
\label{sec:c3v15-companion-audit}
% =============================================================================

The current companion files in \texttt{latex\_notes} were inspected
read-only.  The frozen checkpoints are
\begin{align}
 &\texttt{Renewal\_SM\_Einstein\_Continuum\_Research\_Notes\_V24.tex},
 \notag\\[-2pt]
 &\qquad
 \texttt{70117850268c2bcd4c17e720384d2e0252fada71cc348fb2e5b833a84a688432},
 \label{eq:c3v15-SM-checkpoint}\\
 &\texttt{Renewal\_Yang\_Mills\_Mass\_Gap\_Research\_Notes\_V98.tex},
 \notag\\[-2pt]
 &\qquad
 \texttt{d6e4a8a89a38e4bf7cc687ff7dadc58e982762b53d68b832922fd8e1b31bc945}.
 \label{eq:c3v15-YM-checkpoint}
\end{align}
(SM V24 has 10080 lines; YM V98 has 44404 lines.)

\emph{SM V24 (seventh series).}  After closing the deterministic smooth
matter coefficient on flat charts (V23), the SM programme treated the
scalar-curvature--Higgs owner and identified the missing lattice
metric operator: a metric extension requires the overlap kernel, its
adjoint, the Wilson term, the sign gap, the volume weights and the
stress contacts to be defined for one discrete metric, which the
attached metric continuation supplies only as a continuum coframe.

\emph{YM V98.}  The YM programme treated partition endpoints and
support localization, recorded why arbitrary cube matrices are the
wrong test, and proved the exact hierarchical (nested-partition)
structure of a BKAR interpolation matrix.

\begin{theorem}[V15 companion-audit decision]
\label{thm:c3v15-companion-decision}
No SM V24 curvature or metric constant and no YM V98 partition or
hierarchy bound is imported.  Neither bears on the third cycle's
chapters or on the gates.
\end{theorem}

\begin{proof}
The SM results concern lattice metric operators for Dirac--Yukawa
heat traces; the YM results concern interpolation matrices of a
cluster expansion.  None is a Navier--Stokes statement.
\end{proof}

% =============================================================================
\section{Position after fifteen versions}
\label{sec:c3v15-position}
% =============================================================================

\begin{assessment}[Position after fifteen versions of the third cycle]
\label{ass:c3v15-position-fifteen}
Two chapters are written: the cluster-energy and band chapter
(V2--V10), which closed the multiplicity question as undecidable by
the cell-level tools, with the stationary large-scale band as the
obstruction; and the localization chapter (V11--V14), which reduced
the type-I gate at homogeneous points to a single mean inequality on
a bounded configuration neighbourhood with an explicit cell count,
and showed that the neighbourhood must carry the defect and the
anti-alignment at levels bounded below by the defect rate.  The
gates stand as in the manifest, G1 in its localized form
\eqref{eq:c3v13-gate}.  The direction for V16--V20 is the structure of
the configuration at a homogeneous point (Assessment~\ref{ass:c3v14-position}).
No global regularity claim is made.
\end{assessment}

% =============================================================================
\section{Integrated V15 conclusion and next acquisition order}
\label{sec:c3v15-conclusion}
% =============================================================================

\begin{theorem}[What V15 establishes]
\label{thm:c3v15-conclusion}
Relative to the two legacy cycles and V1--V14: the companion-audit
decision (Theorem~\ref{thm:c3v15-companion-decision}) and the position.
The full Navier--Stokes stability/global-regularity theorem remains
unproved.
\end{theorem}

\begin{proof}
The cited records.
\end{proof}

\begin{programme}[V16 acquisition order]
\label{prog:c3v16}
\begin{enumerate}[label=\textup{(AC\arabic*)},leftmargin=4.8em]
\item The anchor.  At a homogeneous point, at every scale, a cell
      within \(R_c\) of the centre is energetic (Proposition
      C2-V66-active with the pinching level \(c_\Phi'\)); state this as
      the anchor group of the configuration and determine whether the
      finite propagation constrains the other groups' positions
      relative to the anchor across scales (the drift \(-\frac12z\cdot\nabla\)
      in profile variables moves a physically fixed structure outward
      at unit log-rate, and the propagation speed bound \(\rho_0/\tau_c\) is
      large), recording the outcome.
\item The Gaussian radius.  Record that the cross term's
      contributions from beyond the Gaussian radius \(R_G\) are bounded
      by the Gaussian tail times the class constants, independently of
      the trap, and combine with V11 to localize the cross term to the
      groups within \(R_G\) of the centre.
\item The next compulsory companion audit is V20.
\end{enumerate}
\end{programme}

\begin{assessment}[Position after V15]
\label{ass:c3v15-position}
The companion audit is recorded; the direction for V16--V20 stands.
No full stability or global-regularity claim is made.
\end{assessment}

% =============================================================================
\section*{Version V15 (third cycle) append-only record}
\addcontentsline{toc}{section}{Version V15 (third cycle) append-only record}
% =============================================================================

The complete V14 source of the third cycle was retained before this
continuation.  Its SHA--256 digest was
\begin{center}\scriptsize\ttfamily
c31048e4a3b46223bf32fdc8b294dfc5cf20d09f649eccb07d69bd88d0d74fe7.
\end{center}
V15 performed the compulsory companion audit and recorded the
position.  No full regularity claim is made.

% =============================================================================
\section{V16: the anchor group at a homogeneous point, and the Gaussian radius}
\label{sec:c3v16-entry}
% =============================================================================

Items (AC1) and (AC2) of Programme~\ref{prog:c3v16} are carried out.
(AC1): at a homogeneous point the pinching of the Gaussian ledger at
the fixed centre gives, at every scale, an energetic cell within a
bounded lattice distance of the centre, the anchor group; the finite
propagation does not constrain the other groups' positions across
scales, because a physically fixed structure drifts outward at unit
log-rate in profile variables, which the propagation speed bound
permits, so the configuration need not be a single group.  (AC2): the
Gaussian weight makes every contribution to the cross term from
beyond the Gaussian radius negligible independently of the trap, and
with V11 the cross term is localized to the groups within the
Gaussian radius of the centre.  No global regularity claim is made.
The next compulsory companion audit is V20.

% =============================================================================
\section{The anchor}
\label{sec:c3v16-anchor}
% =============================================================================

\begin{proposition}[The anchor group]
\label{prop:c3v16-anchor}
Let \(x_0\) be a homogeneous type-I point with pinching constant
\(c_\Phi'\) (Theorem V91-pinched of the first cycle: \(\Phi_{u,x_0}(\lambda)\ge c_\Phi'\)
on small scales).  Then at every small scale the parabolic cell
containing \(x_0\), or a cell within lattice distance
\(R_{c_\Phi'}=\lceil(4\nu\log(2S_G/c_\Phi'))^{1/2}\rceil\) of it, has energy at least
\(c'\lambda\) with \(c'=c_\Phi'/(2(2R_{c_\Phi'}+1)^3)\); in the profile variables
of any element of the family at \(x_0\), at every \(\sigma\), a cell within
\(R_{c_\Phi'}\) of the origin is energetic at level \(c'\).  The group of
the seeds containing this cell is the \emph{anchor group}; its
distance to the origin is at most \(R_{c_\Phi'}\) at every \(\sigma\).
\end{proposition}

\begin{proof}
Proposition C2-V66-active at the level \(c_\Phi'\) and the fixed centre
\(x_0\); the family's elements inherit the ledger floor by the
continuity of \(\Phi\) on the family (first cycle, V91).
\end{proof}

\begin{proposition}[The other groups are not pinned]
\label{prop:c3v16-not-pinned}
In profile variables, a structure at fixed physical position \(y_1\ne x_0\)
sits at \(z=(y_1-x_0)e^{\sigma/2}\) and moves outward at unit rate in
\(\log|z|\); the finite propagation (a group's non-quiet set moves by at
most \(\rho_0\) cells per rescaled time \(\tau_c\), that is at most
\(\rho_0/\tau_c\) cells per unit of \(\sigma\) up to the factor \(2\)) permits any
motion of a group at distance \(R\) from the origin with
\(|\dd R/\dd\sigma|\le\rho_0/\tau_c\), and \(\rho_0/\tau_c\), of order \(\varepsilon_{\rm th}^{-7/4}\),
exceeds any Gaussian radius \(R_G\sim(8\nu\log(1/\varepsilon_{\rm th}))^{1/2}\); hence a second
group may enter the Gaussian radius from outside, remain, or drift
out, at any scale, compatibly with every bound of the series.  The
configuration at a homogeneous point is the anchor group plus a
bounded number (at most \(N_C-1\)) of other groups whose positions are
unconstrained.
\end{proposition}

\begin{proof}
The drift is the definition of the profile variables; the speed
bound is Theorem C2-V66-propagation read per unit log-time; the
comparison of \(\rho_0/\tau_c\) with the Gaussian radius is arithmetic.
\end{proof}

% =============================================================================
\section{The Gaussian radius}
\label{sec:c3v16-gaussian}
% =============================================================================

\begin{lemma}[Beyond the Gaussian radius nothing counts]
\label{lem:c3v16-tail}
For every profile in the family and every \(R\ge1\), the contributions
to \(\langle N(V),Y\rangle_G\), \(\varphi\), \(d\), \(j\), \(\mathcal L\) and the moment ledger's rows
from \(\{|z|>R\}\) are bounded by \(C_{\rm class}(1+R^4)e^{-R^2/(8\nu)}\), with
\(C_{\rm class}\) depending on \(\nu\), \(C_\infty\), \(C\), \(C_N\), \(C_Y\) only.
\end{lemma}

\begin{proof}
Each integrand is bounded pointwise by a polynomial in \(|z|\) times
class-bounded functions in \(L^2(G)\) or \(L^\infty\); on \(\{|z|>R\}\) write
\(G=G^{1/2}G^{1/2}\le e^{-R^2/(8\nu)}G^{1/2}\) and apply Cauchy--Schwarz with the
\(L^2(G^{1/2})\)-norms, which are bounded by the class (the Gaussian
\(G^{1/2}=e^{-|z|^2/(8\nu)}\) is the weight at twice the viscosity, and the
class bounds transfer).
\end{proof}

\begin{theorem}[The cross term lives on the groups within the Gaussian radius]
\label{thm:c3v16-localized}
At a homogeneous point, for every profile in the support of a blow-up
measure and every \(R\ge r_2\), on the trapped sub-windows,
\begin{equation}
 \langle N(V),Y\rangle_G=X_{U_R}+O\bigl(\varepsilon_{\rm th}^{1/4}C_Y\bigr)+O\bigl((1+R^4)e^{-R^2/(8\nu)}\bigr),
 \label{eq:c3v16-localized}
\end{equation}
where \(U_R\) is the union of the neighbourhoods of the groups whose
seeds lie within \(|z|\le R\); \(U_R\) has at most \(N_C(2r_2+1)^3\) cells and
always contains the anchor group's neighbourhood.  Choosing \(R\) with
\((1+R^4)e^{-R^2/(8\nu)}\le\varepsilon_{\rm th}^{1/4}\) (\(R=R_G(\nu,\varepsilon_{\rm th})\) of order
\((8\nu\log(1/\varepsilon_{\rm th}))^{1/2}\)), the gate \eqref{eq:c3v13-gate} holds with
\(U=U_{R_G}\).
\end{theorem}

\begin{proof}
Theorem~\ref{thm:c3v11-far} for the quiet cells within \(|z|\le R\)
outside the groups' neighbourhoods, and Lemma~\ref{lem:c3v16-tail}
for \(|z|>R\); the anchor is Proposition~\ref{prop:c3v16-anchor}.
\end{proof}

\begin{assessment}[Items (AC1)--(AC2)]
\label{ass:c3v16-AC}
The configuration at a homogeneous point has a permanent anchor at
the centre and possibly other groups within the Gaussian radius,
unpinned; the gate is a mean inequality on the union of their
neighbourhoods, at most \(N_C(2r_2+1)^3\) cells within
\(R_G\sim(8\nu\log(1/\varepsilon_{\rm th}))^{1/2}\) of the centre.  The single-group
question (AC1) closes negatively; the localization (AC2) closes
positively in the form \eqref{eq:c3v16-localized}.  No global
regularity claim is made.
\end{assessment}

% =============================================================================
\section{Integrated V16 conclusion and next acquisition order}
\label{sec:c3v16-conclusion}
% =============================================================================

\begin{theorem}[What V16 proves]
\label{thm:c3v16-conclusion}
Relative to the two legacy cycles and V1--V15: the anchor group
(Proposition~\ref{prop:c3v16-anchor}), the unpinned other groups
(Proposition~\ref{prop:c3v16-not-pinned}), the Gaussian tail lemma
(Lemma~\ref{lem:c3v16-tail}) and the localization \eqref{eq:c3v16-localized}.
The full Navier--Stokes stability/global-regularity theorem remains
unproved.
\end{theorem}

\begin{proof}
The cited results.
\end{proof}

\begin{programme}[V17 acquisition order]
\label{prog:c3v17}
\begin{enumerate}[label=\textup{(AC\arabic*)},leftmargin=4.8em]
\item The anchor alone.  Determine whether the anchor group's
      neighbourhood \(U_0\) (the cells within \(r_2+R_{c_\Phi'}\) of the origin)
      carries a definite part of the anti-alignment on average, using
      that other groups within \(R_G\) are unpinned: if a blow-up measure
      charges profiles with no other group within \(R_G\) with positive
      probability, the gate restricts to \(U_0\) on that event; record
      the two-case statement.
\item The count within the Gaussian radius: state the count of
      energetic cells at level \(c'\) within \(R_G\) precisely, with the
      class constant \(C\) in the role of \(2M\) in Proposition
      C2-V58-count (at most \(1728C_S^6C^3c'^{-3}\)).
\item The next compulsory companion audit is V20.
\end{enumerate}
\end{programme}

\begin{assessment}[Position after V16]
\label{ass:c3v16-position}
The configuration at a homogeneous point is an anchor at the centre
plus unpinned groups within the Gaussian radius; the gate is
localized to their neighbourhoods.  No full stability or
global-regularity claim is made.
\end{assessment}

% =============================================================================
\section*{Version V16 (third cycle) append-only record}
\addcontentsline{toc}{section}{Version V16 (third cycle) append-only record}
% =============================================================================

The complete V15 source of the third cycle was retained before this
continuation.  Its SHA--256 digest was
\begin{center}\scriptsize\ttfamily
d0433b54c1409b3ee28d870a43a33a43dff76b7426c182b788290b53d82bec5f.
\end{center}
V16 proved the anchor group and the Gaussian-radius localization.  No
full regularity claim is made.

% =============================================================================
\section{V17: the anchor alone, and the count within the Gaussian radius}
\label{sec:c3v17-entry}
% =============================================================================

Items (AC1) and (AC2) of Programme~\ref{prog:c3v17} are carried out.
(AC1): the gate at a homogeneous point splits by the event that no
group other than the anchor lies within the Gaussian radius; on that
event the mean inequality restricts to the anchor's neighbourhood,
and on its complement it involves at most \(N_C-1\) further groups
whose positions are unpinned; the two-case statement is recorded.
(AC2): the count of energetic cells within the Gaussian radius at the
anchor's level is stated with the class constant.  No global
regularity claim is made.  The next compulsory companion audit is
V20.

% =============================================================================
\section{The two-case gate}
\label{sec:c3v17-cases}
% =============================================================================

\begin{theorem}[Two-case form of the localized gate]
\label{thm:c3v17-cases}
Let \(x_0\) be a homogeneous type-I point, \(\mathfrak m\) a blow-up measure,
\(U_0\) the anchor's neighbourhood (the cells within \(r_2+R_{c_\Phi'}\) of the
origin) and \(\mathcal E_0\) the event, on the family, that no seed other
than the anchor group's lies within \(|z|\le R_G+r_2\).  Then
\begin{equation}
 \int_{\mathcal E_0}X_{U_0}\,\dd\mathfrak m+\int_{\mathcal E_0^c}X_{U_{R_G}}\,\dd\mathfrak m\ \le\ -\underline\delta+O(\varepsilon_{\rm th}^{1/4}),
 \label{eq:c3v17-cases}
\end{equation}
so that at least one of the following holds:
\begin{enumerate}[label=\textup{(\alph*)},leftmargin=3.2em]
\item \(\mathfrak m(\mathcal E_0)>0\) and \(\int_{\mathcal E_0}X_{U_0}\,\dd\mathfrak m\le-\frac12\underline\delta\,\mathfrak m(\mathcal E_0)+O(\varepsilon_{\rm th}^{1/4})\):
      on the event that the anchor is alone, the anchor's
      neighbourhood carries the anti-alignment at the rate
      \(\frac12\underline\delta\);
\item \(\mathfrak m(\mathcal E_0^c)>0\) and \(\int_{\mathcal E_0^c}X_{U_{R_G}}\,\dd\mathfrak m\le-\frac12\underline\delta\,\mathfrak m(\mathcal E_0^c)+O(\varepsilon_{\rm th}^{1/4})\):
      on the event that other groups are present within the Gaussian
      radius, their union with the anchor carries it.
\end{enumerate}
The events are measurable (they are defined by the energies of
finitely many cells, continuous on the family).
\end{theorem}

\begin{proof}
Theorem~\ref{thm:c3v16-localized} integrated against \(\mathfrak m\), with
\(U_{R_G}=U_0\) on \(\mathcal E_0\); the alternative is the pigeonhole on the
two integrals (their sum is at most \(-\underline\delta+O(\varepsilon_{\rm th}^{1/4})\), so
one of them is at most half its share).
\end{proof}

\begin{remark}[The anchor-alone case]
\label{rem:c3v17-anchor-alone}
In case (a) the gate is a statement about a single group of at most
\(N_C\) cells at the centre with a diffuse far field: the Gaussian
average over the group's neighbourhood of the oriented volume spanned
by vorticity, velocity and defect, plus the Bernoulli-pressure work
against the radial defect, is negative on average at a fixed rate.
This is the cleanest form the gate has taken; it is not closed.  Case
(b) is the multi-group case, where the groups may move freely within
the Gaussian radius (Proposition~\ref{prop:c3v16-not-pinned}).
\end{remark}

% =============================================================================
\section{The count within the Gaussian radius}
\label{sec:c3v17-count}
% =============================================================================

\begin{proposition}[Count at the anchor's level]
\label{prop:c3v17-count}
For every profile in the family at a homogeneous point (class
constant \(C\) in the role of \(2M\)), at every \(\sigma\), the number of
parabolic cells with energy at least \(c'\) is at most
\(1728C_S^6C^3c'^{-3}\), and the number within \(|z|\le R_G\) is at most
\(\min\bigl(1728C_S^6C^3c'^{-3},(2R_G+1)^3\bigr)\); the anchor group is one
of them.
\end{proposition}

\begin{proof}
Proposition C2-V58-count with \(Y\le C\) (the class bound on the
rescaled enstrophy at every time) and the counting of cells in a cube.
\end{proof}

% =============================================================================
\section{Integrated V17 conclusion and next acquisition order}
\label{sec:c3v17-conclusion}
% =============================================================================

\begin{theorem}[What V17 proves]
\label{thm:c3v17-conclusion}
Relative to the two legacy cycles and V1--V16: the two-case gate
(Theorem~\ref{thm:c3v17-cases}) and the count
(Proposition~\ref{prop:c3v17-count}).  The full Navier--Stokes
stability/global-regularity theorem remains unproved.
\end{theorem}

\begin{proof}
The cited results.
\end{proof}

\begin{programme}[V18 acquisition order]
\label{prog:c3v18}
\begin{enumerate}[label=\textup{(AC\arabic*)},leftmargin=4.8em]
\item The centre-mode and momentum identities (C2-V16--V17) in the
      anchor-alone case: the Gaussian mean \(\bar V\) is then the
      anchor group's Gaussian momentum up to \(O(\varepsilon_{\rm th})\); write
      \(\frac{\dd}{\dd\sigma}\bar V=-\frac12\bar V-\overline{N(V)}\) with
      \(\overline{N(V)}\) computed on the anchor's neighbourhood, and
      determine whether the mean identity \(\int\bar V=-2\int\overline{N(V)}\)
      gains content (a bound on the anchor's momentum by its own
      radial energy flux), or whether it is neutral; record.
\item Record the dipole matrix of the anchor and whether the
      centre normalization of C2-V16 can be imposed on the anchor.
\item The next compulsory companion audit is V20.
\end{enumerate}
\end{programme}

\begin{assessment}[Position after V17]
\label{ass:c3v17-position}
The localized gate splits into the anchor-alone case and the
multi-group case; the count within the Gaussian radius is explicit.
No full stability or global-regularity claim is made.
\end{assessment}

% =============================================================================
\section*{Version V17 (third cycle) append-only record}
\addcontentsline{toc}{section}{Version V17 (third cycle) append-only record}
% =============================================================================

The complete V16 source of the third cycle was retained before this
continuation.  Its SHA--256 digest was
\begin{center}\scriptsize\ttfamily
8f93978ab622cd03a1bac5cba7c1463d9d32c2f49807be5d9ab8d49eee06378c.
\end{center}
V17 recorded the two-case gate and the count.  No full regularity
claim is made.

% =============================================================================
\section{V18: the centre-mode and momentum identities in the anchor-alone case}
\label{sec:c3v18-entry}
% =============================================================================

Items (AC1) and (AC2) of Programme~\ref{prog:c3v18} are carried out.
In the anchor-alone case the Gaussian mean of the profile is the
anchor group's Gaussian momentum up to the square root of the
threshold, the momentum identity holds with the nonlinear mean
computed on the anchor's neighbourhood, and the mean identity
balances the anchor's mean momentum against its mean radial momentum
flux and pressure dipole without a sign: it is neutral.  The anchor's
dipole matrix is bounded by the class and the neighbourhood's radius,
and the centre normalization of the second cycle is not available at
the fixed homogeneous centre.  No global regularity claim is made.
The next compulsory companion audit is V20.

% =============================================================================
\section{The Gaussian mean in the anchor-alone case}
\label{sec:c3v18-mean}
% =============================================================================

\begin{proposition}[Mean and momentum on the anchor]
\label{prop:c3v18-mean}
On the event \(\mathcal E_0\) of Theorem~\ref{thm:c3v17-cases}, on the
trapped sub-windows, with \(\bar V=\frac1{|G|}\int VG\) and
\(\bar V_0:=\frac1{|G|}\int_{U_0}VG\):
\begin{enumerate}[label=\textup{(\roman*)},leftmargin=3.2em]
\item \(|\bar V-\bar V_0|\le\frac{S_G}{|G|}\varepsilon_{\rm th}^{1/2}+C_{\rm class}e^{-R_G^2/(8\nu)}\);
\item \(\frac{\dd}{\dd\sigma}\bar V=-\frac12\bar V-\overline{N(V)}\) with
      \(\overline{N(V)}=\frac1{2\nu|G|}\int_{U_0}[(z\cdot V)V+P_Vz]G+O(\varepsilon_{\rm th}^{1/2})\)
      (the far quiet cells contribute at most \(C\varepsilon_{\rm th}^{1/2}\) by the
      bounds of V11 on \(|V|\), \(|z|\) and \(P_V-c\) there, the constant \(c\)
      contributing \(\frac{c}{2\nu|G|}\int zG=0\));
\item against \(\mathfrak m\) restricted to \(\mathcal E_0\),
      \(\int_{\mathcal E_0}\bar V_0\,\dd\mathfrak m=-2\int_{\mathcal E_0}\overline{N(V)}\,\dd\mathfrak m+O(\varepsilon_{\rm th}^{1/2})\)
      up to the boundary terms of the event (the invariance applies to
      the whole family; restricted to an event that is not invariant,
      the identity holds up to the flux of \(\bar V\) through the event's
      boundary, which is bounded by the class and the measure of the
      boundary layer of the event).
\end{enumerate}
\end{proposition}

\begin{proof}
(i) The quiet cells within \(R_G\): \(\int_{\rm cell}|V|G\le e_\beta^{1/2}\sup G\), summed
with \(S_G\); the tail by Lemma~\ref{lem:c3v16-tail}.  (ii) Proposition
C2-V17-mean-profile with the split of the integral.  (iii) The mean
identity of C2-V17 and the remark on non-invariant events.
\end{proof}

\begin{firewall}[Item (AC1): neutral]
\label{fw:c3v18-neutral}
The momentum identity says that the anchor's mean momentum decays at
rate one half and is driven by its own radial momentum flux and
pressure dipole; on average the two balance.  Neither side has a
sign: the anchor may have any momentum, and its radial momentum flux
any sign.  The identity is a row of the ledger localized to the
anchor, and the ledger has no incompatible row.  The item is closed
neutrally.
\end{firewall}

% =============================================================================
\section{The dipole matrix and the centre normalization}
\label{sec:c3v18-dipole}
% =============================================================================

\begin{proposition}[Dipole matrix of the anchor]
\label{prop:c3v18-dipole}
On \(\mathcal E_0\), the dipole matrix \(\Delta_{ij}:=\int V_iz_jG\) satisfies
\(\Delta=\int_{U_0}V_iz_jG+O(\varepsilon_{\rm th}^{1/2})\) and
\(|\Delta|\le(r_2+R_{c_\Phi'})C_\infty|G|+O(\varepsilon_{\rm th}^{1/2})\).  The centre
normalization of the second cycle (C2-V16: a choice of centre making
the dipole matrix invertible or a component of the centre mode
vanish) requires moving the centre, which at a homogeneous point is
fixed at \(x_0\) by the definition of the family; a moving-centre
family would need the coherence of C2-V24, not available.  The
normalization is therefore not imposed.
\end{proposition}

\begin{proof}
The bounds are the class \(L^\infty\) bound and the radius of \(U_0\); the
statement about the normalization is the record of C2-V16 and
C2-V24.
\end{proof}

% =============================================================================
\section{Integrated V18 conclusion and next acquisition order}
\label{sec:c3v18-conclusion}
% =============================================================================

\begin{theorem}[What V18 establishes]
\label{thm:c3v18-conclusion}
Relative to the two legacy cycles and V1--V17: the mean and momentum
on the anchor (Proposition~\ref{prop:c3v18-mean}), the neutrality
record, and the dipole proposition (Proposition~\ref{prop:c3v18-dipole}).
The full Navier--Stokes stability/global-regularity theorem remains
unproved.
\end{theorem}

\begin{proof}
The cited results.
\end{proof}

\begin{programme}[V19 acquisition order]
\label{prog:c3v19}
\begin{enumerate}[label=\textup{(AC\arabic*)},leftmargin=4.8em]
\item Consolidate V16--V18 as the anchor chapter and record the
      position before the V20 audit.
\item Fix the direction for V21--V25: the explicit variance floor.
      In the anchor-alone case the normalized profile is concentrated
      (in the Gaussian sense) in a ball of radius \(r_2+R_{c_\Phi'}\) up
      to \(O(\varepsilon_{\rm th}^{1/2})\), while the Ornstein--Uhlenbeck
      eigenfunctions are Hermite polynomials whose Gaussian mass sits
      at \(|z|\) of order \((2\nu k)^{1/2}\) for degree \(k\); determine
      whether the distance of a concentrated normalized profile to
      every eigenspace is bounded below explicitly, giving an explicit
      variance floor \(v_0\) in place of the compactness constant of
      C2-V34, and what the push budget then says.
\end{enumerate}
\end{programme}

\begin{assessment}[Position after V18]
\label{ass:c3v18-position}
The momentum identity on the anchor is neutral; the centre
normalization is not available at a homogeneous centre.  No full
stability or global-regularity claim is made.
\end{assessment}

% =============================================================================
\section*{Version V18 (third cycle) append-only record}
\addcontentsline{toc}{section}{Version V18 (third cycle) append-only record}
% =============================================================================

The complete V17 source of the third cycle was retained before this
continuation.  Its SHA--256 digest was
\begin{center}\scriptsize\ttfamily
b48f0c00f9f9fb8c9227bef06a99e6215a27846d689f195fe3576bd4fbe080e5.
\end{center}
V18 recorded the momentum identity on the anchor and the dipole
proposition.  No full regularity claim is made.

% =============================================================================
\section{V19: the anchor chapter consolidated, and the direction for V21--V25}
\label{sec:c3v19-entry}
% =============================================================================

Items (AC1) and (AC2) of Programme~\ref{prog:c3v19} are carried out.
No global regularity claim is made.  The next compulsory companion
audit is V20.

% =============================================================================
\section{The anchor chapter}
\label{sec:c3v19-consolidated}
% =============================================================================

\begin{theorem}[The anchor chapter, V16--V18]
\label{thm:c3v19-consolidated}
At a homogeneous type-I point with pinching constant \(c_\Phi'\), for
every profile in the support of a blow-up measure:
\begin{enumerate}[label=\textup{(\arabic*)},leftmargin=3.2em]
\item (Anchor.)  At every \(\sigma\) a cell within \(R_{c_\Phi'}\) of the
      origin is energetic at level \(c'=c_\Phi'/(2(2R_{c_\Phi'}+1)^3)\); its
      group is the anchor group, of at most \(N_C\) cells, at bounded
      distance from the centre at every scale (V16).
\item (Unpinned others.)  The other groups, at most \(N_C-1\), may
      enter, remain within, or leave the Gaussian radius at any scale
      compatibly with every bound of the series (V16).
\item (Localization.)  The cross term equals its restriction to the
      groups within the Gaussian radius \(R_G\sim(8\nu\log(1/\varepsilon_{\rm th}))^{1/2}\)
      up to \(O(\varepsilon_{\rm th}^{1/4})\), and the gate G1 is the mean inequality
      on their neighbourhoods; it splits into the anchor-alone case
      (the anchor's neighbourhood carries the anti-alignment at rate
      \(\frac12\underline\delta\) on that event) and the multi-group case (V16--V17).
\item (Count.)  At most \(\min(1728C_S^6C^3c'^{-3},(2R_G+1)^3)\) cells within
      the Gaussian radius are energetic at level \(c'\) (V17).
\item (Momentum.)  In the anchor-alone case the Gaussian mean is the
      anchor's momentum up to \(O(\varepsilon_{\rm th}^{1/2})\), the momentum identity
      holds on the anchor and is neutral, and the centre normalization
      is not available at the fixed centre (V18).
\end{enumerate}
\end{theorem}

\begin{proof}
Propositions~\ref{prop:c3v16-anchor}, \ref{prop:c3v16-not-pinned},
\ref{prop:c3v17-count}, \ref{prop:c3v18-mean}, \ref{prop:c3v18-dipole};
Theorems~\ref{thm:c3v16-localized}, \ref{thm:c3v17-cases}.
\end{proof}

% =============================================================================
\section{Direction for V21--V25: the explicit variance floor}
\label{sec:c3v19-direction}
% =============================================================================

\begin{assessment}[Position before the V20 audit, and direction]
\label{ass:c3v19-position}
The type-I gate at homogeneous points is now, in its cleanest case, a
mean inequality on a single group of at most \(N_C\) parabolic cells at
the centre with a diffuse far field.  Of the second cycle's Rayleigh
chapter, the one signed quantity is the variance floor
\(\mathrm{Var}\ge v_0\) (C2-V34), with \(v_0\) from compactness.  The direction
for V21--V25 is to make it explicit in the anchor-alone case: V21
computes the Gaussian mass distribution of the Ornstein--Uhlenbeck
eigenfunctions (Hermite polynomials of degree \(k\) in
\(z/(2\nu)^{1/2}\)) and shows that their normalized Gaussian mass within a
ball of radius \(R\) is at most an explicit \(\mu_k(R)\) tending to zero
as \(k\to\infty\) for fixed \(R\); V22 shows that a normalized profile
with Gaussian mass at least \(1-\eta\) in the ball of radius
\(r_2+R_{c_\Phi'}\) (the anchor-alone case, \(\eta=O(\varepsilon_{\rm th}^{1/2})\)) has
distance at least an explicit \(v(R,\eta)>0\) to every eigenspace and to
every finite sum of eigenspaces of bounded degree, hence an explicit
variance floor if the spectral distribution of the profile is also
controlled; V23 examines the last condition (the profile is not a
low-degree Hermite polynomial times a concentrated correction) and
whether the profile equation excludes it; V24 consolidates; V25
audits.  No global regularity claim is made.
\end{assessment}

% =============================================================================
\section{Integrated V19 conclusion and next acquisition order}
\label{sec:c3v19-conclusion}
% =============================================================================

\begin{theorem}[What V19 establishes]
\label{thm:c3v19-conclusion}
Relative to the two legacy cycles and V1--V18: the consolidated anchor
chapter (Theorem~\ref{thm:c3v19-consolidated}) and the direction.  The
full Navier--Stokes stability/global-regularity theorem remains
unproved.
\end{theorem}

\begin{proof}
The cited results.
\end{proof}

\begin{programme}[V20 acquisition order]
\label{prog:c3v20}
\begin{enumerate}[label=\textup{(AC\arabic*)},leftmargin=4.8em]
\item The compulsory companion audit: read the current SM and YM
      notes and record their digests and decision.
\item Record the position after twenty versions; delivery.
\end{enumerate}
\end{programme}

\begin{assessment}[Position after V19]
\label{ass:c3v19-position-final}
The anchor chapter is consolidated; the direction is the explicit
variance floor.  No full stability or global-regularity claim is
made.
\end{assessment}

% =============================================================================
\section*{Version V19 (third cycle) append-only record}
\addcontentsline{toc}{section}{Version V19 (third cycle) append-only record}
% =============================================================================

The complete V18 source of the third cycle was retained before this
continuation.  Its SHA--256 digest was
\begin{center}\scriptsize\ttfamily
64201f568b2c9c7b422514d155c6ed39723df5f019e506680b2751ffd2b01ce8.
\end{center}
V19 consolidated the anchor chapter.  No full regularity claim is
made.

% =============================================================================
\section{V20: compulsory companion audit and the position after twenty versions}
\label{sec:c3v20-entry}
% =============================================================================

This is the twentieth version of the third cycle.  The compulsory
review of the two companion programmes is performed first, and the
position after twenty versions is recorded with the plan for the
variance-floor block.  No global regularity claim is made.  The next
compulsory companion audit is V25.

% =============================================================================
\section{Compulsory V20 review of the current companion notes}
\label{sec:c3v20-companion-audit}
% =============================================================================

The current companion files in \texttt{latex\_notes} were inspected
read-only.  The frozen checkpoints are
\begin{align}
 &\texttt{Renewal\_SM\_Einstein\_Continuum\_Research\_Notes\_V24.tex},
 \notag\\[-2pt]
 &\qquad
 \texttt{70117850268c2bcd4c17e720384d2e0252fada71cc348fb2e5b833a84a688432},
 \label{eq:c3v20-SM-checkpoint}\\
 &\texttt{Renewal\_Yang\_Mills\_Mass\_Gap\_Research\_Notes\_V99.tex},
 \notag\\[-2pt]
 &\qquad
 \texttt{6cfcbe207cc33f3f51c9fd87af7e4abeee71ebe3d69c69db143daac828392271}.
 \label{eq:c3v20-YM-checkpoint}
\end{align}
(SM V24 is unchanged since the V15 audit; YM V99 has 44802 lines.)

\emph{SM V24.}  Unchanged: the scalar-curvature--Higgs owner and the
missing lattice metric operator recorded at V15.

\emph{YM V99.}  The YM programme showed that optional support labels
disappear after one layer, defined the remaining forest projection,
and related it exactly to its V88 defect (a forced-support excess
after optional collapse).

\begin{theorem}[V20 companion-audit decision]
\label{thm:c3v20-companion-decision}
Nothing is imported from SM V24 or YM V99.  Neither bears on the
third cycle's chapters or on the gates.
\end{theorem}

\begin{proof}
The SM file is unchanged; the YM results concern forest projections
of a cluster expansion.  None is a Navier--Stokes statement.
\end{proof}

% =============================================================================
\section{Position after twenty versions, and the variance-floor block}
\label{sec:c3v20-position}
% =============================================================================

\begin{assessment}[Position after twenty versions of the third cycle]
\label{ass:c3v20-position-twenty}
Three chapters are written: cluster energy and the band (V2--V10),
localization (V11--V14), and the anchor (V16--V19).  The type-I gate
at homogeneous points is, in its cleanest case, a mean inequality on
the anchor group's neighbourhood.  The block V21--V25 attacks the one
signed quantity of the Rayleigh chapter, the variance floor, with the
following plan, refined from V19.  V21: the spectral form of the
variance, \(\mathrm{Var}=\sum_k(\mu_k-r)^2\hat a_k\) with \(\hat a_k=a_k/\varphi\) the
spectral weights of the normalized profile and \(r=\sum_k\mu_k\hat a_k\);
since the eigenvalues are spaced by \(\frac12\), \(\mathrm{Var}\ge\frac1{16}(1-\max_k(\hat a_k+\hat a_{k+1}))\):
the variance is small only if the spectral measure concentrates on
two consecutive levels.  V22: high levels are excluded by the
concentration of the anchor-alone profile in a ball, using the
Gaussian mass distribution of the Ornstein--Uhlenbeck eigenspaces
(a classical fact about Hermite functions, to be used as stated and
marked).  V23: low levels are polynomial profiles; a profile close to
a low level in the spectral sense is close to a polynomial vector
field of low degree, whose enstrophy is spread over the whole space,
while the class enstrophy bound and the anchor's energy level bound
the region over which the profile can be polynomial; the variance
then has an explicit contribution from the edge of that region.  V24
consolidates the explicit floor, if obtained, and the push budget's
consequence; V25 audits.  No global regularity claim is made.
\end{assessment}

% =============================================================================
\section{Integrated V20 conclusion and next acquisition order}
\label{sec:c3v20-conclusion}
% =============================================================================

\begin{theorem}[What V20 establishes]
\label{thm:c3v20-conclusion}
Relative to the two legacy cycles and V1--V19: the companion-audit
decision (Theorem~\ref{thm:c3v20-companion-decision}) and the position.
The full Navier--Stokes stability/global-regularity theorem remains
unproved.
\end{theorem}

\begin{proof}
The cited records.
\end{proof}

\begin{programme}[V21 acquisition order]
\label{prog:c3v21}
\begin{enumerate}[label=\textup{(AC\arabic*)},leftmargin=4.8em]
\item Prove the spectral form of the variance and the two-level
      criterion \(\mathrm{Var}\ge\frac1{16}(1-\max_k(\hat a_k+\hat a_{k+1}))\); record
      that the modewise identities (C2-V40--V41) give
      \(\langle\hat a_k\rangle=O(1/k)\) on average but not a pointwise bound.
\item State the dichotomy for a profile with small variance: spectral
      mass at least \(1-16\,\mathrm{Var}\) on two consecutive levels
      \(k_0,k_0+1\), either high (\(k_0\ge k_\star\)) or low (\(k_0<k_\star\)),
      with \(k_\star\) to be fixed in V22.
\item The next compulsory companion audit is V25.
\end{enumerate}
\end{programme}

\begin{assessment}[Position after V20]
\label{ass:c3v20-position}
The companion audit is recorded; the variance-floor block begins.
No full stability or global-regularity claim is made.
\end{assessment}

% =============================================================================
\section*{Version V20 (third cycle) append-only record}
\addcontentsline{toc}{section}{Version V20 (third cycle) append-only record}
% =============================================================================

The complete V19 source of the third cycle was retained before this
continuation.  Its SHA--256 digest was
\begin{center}\scriptsize\ttfamily
67e21e43d2813d54c8315565da70f074133b2a2913d7f4e03b345ec0bd389daa.
\end{center}
V20 performed the compulsory companion audit and recorded the
position.  No full regularity claim is made.

% =============================================================================
\section{V21: the spectral form of the variance, the dominant level, and its bound}
\label{sec:c3v21-entry}
% =============================================================================

Items (AC1) and (AC2) of Programme~\ref{prog:c3v21} are carried out.
The spectral form of the variance (Lemma C2-V31-degree) is sharpened
into a one-level criterion: because the eigenvalues of the
self-similar operator are spaced by one half, at most one level lies
within a quarter of the Rayleigh quotient, and the variance is at
least one sixteenth of the spectral mass off that level.  The
dominant level is bounded by the class bound on the Rayleigh
quotient, pointwise at every scale, so that a profile of small
variance is close, in the Gaussian sense, to a divergence-free
polynomial field of an explicitly bounded Hermite degree.  The
dichotomy planned in V20 therefore degenerates: the high case is
empty by the Rayleigh bound alone, and the Hermite-mass computation
planned for V22 is not needed for exclusion; the direction of the
block is corrected accordingly.  No global regularity claim is made.
The next compulsory companion audit is V25.

% =============================================================================
\section{The one-level and two-level criteria}
\label{sec:c3v21-criteria}
% =============================================================================

Notation of the second cycle's Rayleigh chapter (C2-V31, C2-V40):
\(V\) is the profile of an element of the family at a homogeneous
type-I point, \(\varphi=\|V\|_G^2\), \(\widehat V=V/\varphi^{1/2}\),
\(\Pi_k\) the \(G\)-orthogonal projection onto the eigenspace
\(E_k\) of \(-L_-\) with eigenvalue \(\mu_k=\frac12(1+k)\) (divergence-free
polynomial fields whose components are Hermite polynomials of total
degree \(k\) in \(z/(2\nu)^{1/2}\), Lemma C2-V34-eigen), and
\begin{equation}
 \hat a_k:=\|\Pi_kV\|_G^2/\varphi,\qquad
 \sum_k\hat a_k=1,\qquad
 r=\sum_k\mu_k\hat a_k,\qquad
 \mathrm{Var}=\sum_k(\mu_k-r)^2\hat a_k
 \label{eq:c3v21-spectral}
\end{equation}
(Lemma C2-V31-degree; \(\frac12\le r\le C_r\) with
\(C_r=(4\nu C+2C_\Phi)/(4c_\Phi)\) by C2-V33).

\begin{definition}[Dominant level]
\label{def:c3v21-dominant}
For a profile with Rayleigh quotient \(r\), let \(k_0:=\max\{k:\mu_k\le r\}\)
and let \(k_1\) be the least \(k\) minimizing \(|\mu_k-r|\); \(k_1\in\{k_0,k_0+1\}\)
is the dominant level.
\end{definition}

\begin{proposition}[One-level and two-level criteria]
\label{prop:c3v21-criteria}
For every profile of the family and every \(\sigma\):
\begin{enumerate}[label=\textup{(\roman*)},leftmargin=3.2em]
\item \(|\mu_{k_1}-r|\le\frac14\); \(|\mu_k-r|\ge\frac14\) for every \(k\ne k_1\);
      \(|\mu_k-r|\ge\frac12\) for every \(k\notin\{k_0,k_0+1\}\).
\item \(\mathrm{Var}\ \ge\ (\mu_{k_1}-r)^2\hat a_{k_1}+\frac1{16}(1-\hat a_{k_1})\ \ge\ \frac1{16}\bigl(1-\max_k\hat a_k\bigr)\).
\item \(\mathrm{Var}\ \ge\ \frac14\bigl(1-\hat a_{k_0}-\hat a_{k_0+1}\bigr)\ \ge\ \frac14\bigl(1-\max_k(\hat a_k+\hat a_{k+1})\bigr)\).
\item \(\mathrm{Var}=0\) if and only if \(\hat a_{k_1}=1\).
\end{enumerate}
\end{proposition}

\begin{proof}
(i) \(r\in[\mu_{k_0},\mu_{k_0+1})\), an interval of length \(\frac12\); the
nearer endpoint is at distance at most \(\frac14\), the other at
distance at least \(\frac14\), and every other eigenvalue is at
distance at least \(\frac12\) from \(r\), since consecutive eigenvalues
differ by \(\frac12\).  (ii) In \eqref{eq:c3v21-spectral} keep the term
\(k=k_1\) and bound \((\mu_k-r)^2\ge\frac1{16}\) for \(k\ne k_1\), using
\(\sum_{k\ne k_1}\hat a_k=1-\hat a_{k_1}\); then \(\hat a_{k_1}\le\max_k\hat a_k\).
(iii) Discard the terms \(k=k_0,k_0+1\) and bound \((\mu_k-r)^2\ge\frac14\)
for the others.  (iv) If \(\hat a_{k_1}=1\) then \(r=\mu_{k_1}\) and the sum
has one vanishing term; conversely \(\mathrm{Var}=0\) forces \(\hat a_k=0\)
for every \(k\) with \(\mu_k\ne r\), so the mass sits on a single level,
which is \(k_1\).
\end{proof}

\begin{remark}[The criterion of the V20 programme]
\label{rem:c3v21-programme-form}
Item (AC1) asked for \(\mathrm{Var}\ge\frac1{16}(1-\max_k(\hat a_k+\hat a_{k+1}))\);
Proposition~\ref{prop:c3v21-criteria}(iii) gives it with the constant
\(\frac14\), and (ii) gives the sharper one-level form: a small
variance forces the spectral mass onto a single level, not merely
onto two consecutive ones.
\end{remark}

% =============================================================================
\section{The dominant level is bounded}
\label{sec:c3v21-bounded}
% =============================================================================

\begin{proposition}[Pointwise bound on the dominant level and on the spectral tail]
\label{prop:c3v21-bounded}
For every profile of the family and every \(\sigma\):
\begin{enumerate}[label=\textup{(\roman*)},leftmargin=3.2em]
\item \(k_1\le K_r:=\lfloor2C_r-\tfrac12\rfloor\);
\item for every \(K\ge0\), \(\sum_{k\ge K}\hat a_k\le r/\mu_K\le2C_r/(K+1)\).
\end{enumerate}
\end{proposition}

\begin{proof}
(i) \(\mu_{k_1}\le r+\frac14\le C_r+\frac14\), i.e.\ \(\frac12(1+k_1)\le C_r+\frac14\).
(ii) Chebyshev in \eqref{eq:c3v21-spectral}: \(\mu_K\sum_{k\ge K}\hat a_k\le\sum_k\mu_k\hat a_k=r\).
\end{proof}

\begin{corollary}[Small variance means nearly polynomial of bounded level]
\label{cor:c3v21-nearly-polynomial}
If \(\mathrm{Var}<\frac1{32}\) at some \(\sigma\), put \(\eta:=16\,\mathrm{Var}<\frac12\)
and \(P:=\Pi_{k_1}\widehat V\).  Then \(P\in E_{k_1}\) is a divergence-free
polynomial field of exact Hermite level \(k_1\le K_r\), and
\begin{equation}
 \|\widehat V-P\|_G^2=1-\hat a_{k_1}\le\eta,\qquad \|P\|_G^2=\hat a_{k_1}\ge1-\eta>\tfrac12 .
 \label{eq:c3v21-nearly}
\end{equation}
\end{corollary}

\begin{proof}
Proposition~\ref{prop:c3v21-criteria}(ii) gives \(1-\hat a_{k_1}\le16\,\mathrm{Var}\);
Proposition~\ref{prop:c3v21-bounded}(i) bounds \(k_1\); \(\Pi_{k_1}\)
preserves the divergence-free condition (C2-V8-spectrum), and
\(\|\widehat V-\Pi_{k_1}\widehat V\|_G^2=\sum_{k\ne k_1}\hat a_k\).
\end{proof}

\begin{firewall}[Item (AC1): the mean spectrum is not a pointwise bound]
\label{fw:c3v21-mean-not-pointwise}
Theorem C2-V41-spectrum gives \(\langle\hat a_k\rangle_{\mathfrak m}\le C_{\rm sp}/(c_\Phi k)\)
for \(k\ge1\) against every invariant measure, an average bound; the
criterion of Proposition~\ref{prop:c3v21-criteria} involves
\(\max_k\hat a_k\) inside the mean, and
\(\langle\max_k\hat a_k\rangle\) is not controlled by \(\max_k\langle\hat a_k\rangle\).
The mean spectrum is therefore not the tool; the pointwise Chebyshev
bound of Proposition~\ref{prop:c3v21-bounded}(ii), which uses only
the class bound on the Rayleigh quotient, is, and it is what the
block uses from here on.
\end{firewall}

% =============================================================================
\section{The dichotomy degenerates: correction of the direction}
\label{sec:c3v21-correction}
% =============================================================================

\begin{firewall}[Item (AC2): the high case is empty]
\label{fw:c3v21-high-empty}
Item (AC2) of Programme~\ref{prog:c3v21} asked for the dichotomy
``high'' (\(k_0\ge k_\star\)) versus ``low'' (\(k_0<k_\star\)) for a
profile of small variance, with \(k_\star\) to be fixed in V22 by the
Gaussian mass distribution of the Ornstein--Uhlenbeck eigenspaces.
With \(k_\star:=K_r+1\) the high case is empty at every \(\sigma\) for every
profile, by Proposition~\ref{prop:c3v21-bounded}(i), which uses only
the class bound on the Rayleigh quotient; no property of Hermite
functions enters.  The direction stated in V19 and V20 for V22 is
therefore corrected: the Hermite-mass computation is not needed for
the exclusion of high levels.  What remains is the low case in its
entirety, \(k_1\le K_r\), and the question is whether a profile of the
class can be, in the Gaussian sense, nearly a polynomial field of
level at most \(K_r\).
\end{firewall}

\begin{assessment}[Where the class must enter]
\label{ass:c3v21-where-class}
The second cycle excluded exact polynomial profiles by the class
\(L^6\) bound (Lemma C2-V34-eigen: a polynomial field in \(L^6(\R^3)\) is
zero), and obtained the floor \(v_0\) by compactness.  The Gaussian
closeness \eqref{eq:c3v21-nearly} sees only the Gaussian bulk, and V19
hoped that concentration of the profile in a ball would by itself
bound its distance to the eigenspaces below.  For a fixed low level
this hope must be tested against the simplest field concentrated in
a ball: a constant cut off at radius \(\rho\).  V22 computes its
variance and finds it exponentially small in \(\rho^2/(4\nu)\), so
that concentration alone gives no floor; the floor must come, as in
the second cycle, from the unweighted class bounds
(\(\|V\|_{L^6(\R^3)}\le K_6\), C2-V36), which cap the radius of the
polynomial core.  V23 carries this out with explicit constants for
the levels \(0\) and \(1\) and finite-dimensional constants for the
levels \(2\le k\le K_r\).
\end{assessment}

% =============================================================================
\section{Integrated V21 conclusion and next acquisition order}
\label{sec:c3v21-conclusion}
% =============================================================================

\begin{theorem}[What V21 proves]
\label{thm:c3v21-conclusion}
Relative to the two legacy cycles and V1--V20: the one-level and
two-level criteria (Proposition~\ref{prop:c3v21-criteria}), the
pointwise bound on the dominant level and the spectral tail
(Proposition~\ref{prop:c3v21-bounded}), and the nearly-polynomial
corollary (Corollary~\ref{cor:c3v21-nearly-polynomial}).  The full
Navier--Stokes stability/global-regularity theorem remains unproved.
\end{theorem}

\begin{proof}
The cited results.
\end{proof}

\begin{programme}[V22 acquisition order]
\label{prog:c3v22}
\begin{enumerate}[label=\textup{(AC\arabic*)},leftmargin=4.8em]
\item The cut-off eigenfunction: for a constant \(c_0\) and a radius
      \(\rho\), a divergence-free field equal to \(c_0\) on \(B_\rho\) and
      supported in \(B_{2\rho}\) (the curl of \(\frac12(c_0\times z)\) times a
      cut-off), with an explicit bound on its variance of the order
      of \(e^{-\rho^2/(4\nu)}\) times a polynomial in \(\rho^2/\nu\); the
      same for a level-\(k\) eigenfunction.  Record as a negative: the
      distance of a concentrated normalized field to the low
      eigenspaces has no lower bound in terms of the concentration
      radius alone.
\item The growth lemma: for \(Q\in E_k\) with \(\|Q\|_G=1\),
      \(\int_{B_\rho}|Q|^2\,\dd z\ge\gamma_k\rho^{2k+3}\) for \(\rho\ge\rho_k^{\rm top}\),
      with \(\gamma_0=\frac{4\pi}{3}|G|^{-1}\), \(\rho_0^{\rm top}=0\),
      \(\gamma_1=\frac{4\pi}{15}(2\nu|G|)^{-1}\), \(\rho_1^{\rm top}=0\), and
      for \(k\ge2\) constants computable from the Hermite coefficients
      (finite-dimensional).
\item V23: the explicit floor from the unweighted \(L^6\) bound; V24
      consolidation and the push budget; V25 compulsory companion
      audit.
\end{enumerate}
\end{programme}

\begin{assessment}[Position after V21]
\label{ass:c3v21-position}
A profile of small variance is nearly a polynomial field of level at
most \(K_r\), pointwise; the high case is empty; the floor must come
from the unweighted class bounds.  No full stability or
global-regularity claim is made.
\end{assessment}

% =============================================================================
\section*{Version V21 (third cycle) append-only record}
\addcontentsline{toc}{section}{Version V21 (third cycle) append-only record}
% =============================================================================

The complete V20 source of the third cycle was retained before this
continuation.  Its SHA--256 digest was
\begin{center}\scriptsize\ttfamily
f4efb65aefeaab1ca86df716fa5ee77ae8276b9b14ea20779167a3b577f8d389.
\end{center}
V21 proved the one-level criterion, bounded the dominant level, and
corrected the direction of the block: the high case is empty by the
Rayleigh bound.  No full regularity claim is made.

% =============================================================================
\section{V22: concentration gives no floor; the growth of eigenfunctions on large balls}
\label{sec:c3v22-entry}
% =============================================================================

Items (AC1) and (AC2) of Programme~\ref{prog:c3v22} are carried out.
(AC1): a constant vector field cut off divergence-freely at radius
\(\rho\) is concentrated in the ball of radius \(2\rho\) and has spectral
variance at most an explicit multiple of \(\rho\,e^{-\rho^2/(4\nu)}\); the
same holds for every eigenfunction of the self-similar operator with
a polynomial prefactor.  The hope recorded in V19, that
concentration of a normalized field in a ball bounds its distance to
the low eigenspaces below, is therefore false, and is recorded as a
negative.  (AC2): a normalized eigenfunction of level \(k\) has
unweighted energy at least \(\gamma_k\rho^{2k+3}\) in the ball of radius
\(\rho\) beyond an explicit radius, with explicit constants for the
levels \(0\) and \(1\) and finite-dimensional constants for the higher
levels.  The two facts are the two halves of the floor to be proved
in V23: the class bounds cap the radius of a polynomial core, and a
polynomial core of capped radius has variance bounded below by the
Gaussian tail at that radius.  No global regularity claim is made.
The next compulsory companion audit is V25.

% =============================================================================
\section{The cut-off constant}
\label{sec:c3v22-cutoff}
% =============================================================================

Fix once and for all \(\chi\in C_c^\infty(\R^3)\) with \(0\le\chi\le1\),
\(\chi=1\) on \(B_1\), \(\chi=0\) outside \(B_2\), and let
\(A:=\max(\sup|\nabla\chi|,\sup|\nabla^2\chi|,\sup|\nabla^3\chi|)\).  For \(\rho>0\)
put \(\chi_\rho(z):=\chi(z/\rho)\).  Throughout, \(\mathrm{Var}(W)\) for a
divergence-free \(W\in H^2(G)\) with \(\|W\|_G>0\) denotes
\(\|L_-\widehat W\|_G^2-r_W^2\) with \(\widehat W=W/\|W\|_G\) and
\(r_W=-\langle L_-\widehat W,\widehat W\rangle_G\) (the definition of C2-V31,
applied to \(W\) in place of the profile); for every \(\lambda\in\R\),
\begin{equation}
 \|(L_-+\lambda)\widehat W\|_G^2=\mathrm{Var}(W)+(\lambda-r_W)^2\ \ge\ \mathrm{Var}(W).
 \label{eq:c3v22-var-lambda}
\end{equation}

\begin{counterexample}[The cut-off constant has exponentially small variance]
\label{cex:c3v22-cutoff-constant}
For \(c_0\in\R^3\setminus\{0\}\) and \(\rho>0\) let
\(P_\rho:=\operatorname{curl}\bigl(\tfrac12\chi_\rho\,(c_0\times z)\bigr)\).  Then
\(P_\rho\) is smooth, divergence-free, equal to \(c_0\) on \(B_\rho\), supported
in \(B_{2\rho}\), and, for \(\rho^2\ge12\nu\),
\begin{equation}
 \boxed{\quad
 \mathrm{Var}(P_\rho)\ \le\ \frac{96A^2\rho}{(4\pi\nu)^{1/2}}\,e^{-\rho^2/(4\nu)} .
 \quad}
 \label{eq:c3v22-cutoff-var}
\end{equation}
Moreover \(\|P_\rho\|_G^2\ge\frac12|c_0|^2|G|\), \(\|P_\rho\|_{L^6(\R^3)}^6\ge\frac{4\pi}3|c_0|^6\rho^3\)
and \(\|\nabla P_\rho\|_{L^2(\R^3)}^2\le84\pi A^2|c_0|^2\rho\).
\end{counterexample}

\begin{proof}
Write \(\Psi_0:=\frac12c_0\times z\), so \(\operatorname{curl}\Psi_0=c_0\), and
\(P_\rho=\operatorname{curl}(\chi_\rho\Psi_0)=\chi_\rho c_0+\nabla\chi_\rho\times\Psi_0\).  A curl
is divergence-free; on \(B_\rho\), \(\chi_\rho=1\) and \(\nabla\chi_\rho=0\), so
\(P_\rho=c_0\); outside \(B_{2\rho}\) both terms vanish.  On the annulus
\(\rho\le|z|\le2\rho\): \(|\nabla^j\chi_\rho|\le A\rho^{-j}\) for \(j=1,2,3\) and
\(|\Psi_0|\le\rho|c_0|\), \(|\nabla\Psi_0|\le\frac12|c_0|\), \(\nabla^2\Psi_0=0\); hence
\(|P_\rho|\le(1+A)|c_0|\), \(|\nabla P_\rho|\le3A|c_0|/\rho\) and
\(|\nabla^2P_\rho|\le4A|c_0|/\rho^2\) (each derivative of the product is a sum
of at most that many terms, each bounded by \(A|c_0|\) times the
appropriate power of \(\rho\)).  Now
\((L_-+\tfrac12)P_\rho=\nu\Delta P_\rho-\tfrac12z\cdot\nabla P_\rho\) vanishes on \(B_\rho\)
(where \(P_\rho\) is constant) and outside \(B_{2\rho}\), and on the annulus
is bounded by \(4\nu A|c_0|/\rho^2+\rho\cdot3A|c_0|/\rho=A|c_0|(3+4\nu/\rho^2)\).
Therefore
\[
 \|(L_-+\tfrac12)P_\rho\|_G^2\ \le\ A^2|c_0|^2(3+4\nu/\rho^2)^2\int_{|z|>\rho}G .
\]
The Gaussian tail: \(\int_{|z|>\rho}G=4\pi\int_\rho^\infty s^2e^{-s^2/(4\nu)}\dd s\), and
integrating \(s\cdot(se^{-s^2/(4\nu)})\) by parts,
\(\int_\rho^\infty s^2e^{-s^2/(4\nu)}\dd s=2\nu\rho e^{-\rho^2/(4\nu)}+2\nu\int_\rho^\infty e^{-s^2/(4\nu)}\dd s\le2\nu(\rho+2\nu/\rho)e^{-\rho^2/(4\nu)}\)
(the last integral is at most \((2\nu/\rho)e^{-\rho^2/(4\nu)}\) since
\(e^{-s^2/(4\nu)}\le(s/\rho)e^{-s^2/(4\nu)}\) for \(s\ge\rho\)).  So
\(\int_{|z|>\rho}G\le8\pi\nu(\rho+2\nu/\rho)e^{-\rho^2/(4\nu)}\).  For the
energy, \(\|P_\rho\|_G^2\ge|c_0|^2\int_{B_\rho}G\ge|c_0|^2(|G|-6\nu|G|/\rho^2)\ge\frac12|c_0|^2|G|\)
when \(\rho^2\ge12\nu\), by Chebyshev with \(\int|z|^2G=6\nu|G|\).  By
\eqref{eq:c3v22-var-lambda} with \(\lambda=\frac12\),
\[
 \mathrm{Var}(P_\rho)\le\frac{\|(L_-+\frac12)P_\rho\|_G^2}{\|P_\rho\|_G^2}
 \le\frac{2A^2(3+4\nu/\rho^2)^2\,8\pi\nu(\rho+2\nu/\rho)}{|G|}e^{-\rho^2/(4\nu)} ,
\]
and for \(\rho^2\ge12\nu\), \(4\nu/\rho^2\le\frac13\), \((3+\frac13)^2\le12\),
\(\rho+2\nu/\rho\le2\rho\), so the right side is at most
\(384\pi\nu A^2\rho\,e^{-\rho^2/(4\nu)}/(4\pi\nu)^{3/2}=96A^2\rho(4\pi\nu)^{-1/2}e^{-\rho^2/(4\nu)}\).
The \(L^6\) bound is the integral over \(B_\rho\), and the enstrophy bound
is \(|\nabla P_\rho|^2\le9A^2|c_0|^2\rho^{-2}\) over the annulus of volume
\(\frac{4\pi}3(8-1)\rho^3\).
\end{proof}

\begin{proposition}[Cut-off eigenfunctions of every level]
\label{prop:c3v22-cutoff-level-k}
Let \(Q\in E_k\), \(Q\ne0\), and \(\Psi_Q(z):=-z\times\int_0^1tQ(tz)\,\dd t\) (a polynomial
vector potential, \(\operatorname{curl}\Psi_Q=Q\)).  Put
\(Q_\rho:=\operatorname{curl}(\chi_\rho\Psi_Q)=\chi_\rho Q+\nabla\chi_\rho\times\Psi_Q\).  Then \(Q_\rho\)
is divergence-free, equal to \(Q\) on \(B_\rho\), supported in \(B_{2\rho}\), and
for \(\rho\ge\rho_k^{\rm half}:=(8\nu(3+2k))^{1/2}\),
\begin{equation}
 \mathrm{Var}(Q_\rho)\ \le\ \frac{2\,\|(L_-+\mu_k)Q_\rho\|_G^2}{\|Q\|_G^2}
 \ \le\ C_{k,\nu,A}\,(1+\rho)^{2k+1}\,e^{-\rho^2/(4\nu)} ,
 \label{eq:c3v22-cutoff-level-k}
\end{equation}
where \(C_{k,\nu,A}\) is computable from \(k\), \(\nu\), \(A\) and the
coefficients of \(Q/\|Q\|_G\) (a finite-dimensional quantity).
\end{proposition}

\begin{proof}
The vector potential: for a divergence-free field \(Q\),
\(\operatorname{curl}\bigl(-z\times\int_0^1tQ(tz)\dd t\bigr)=Q(z)\) (the Poincar\'e
lemma on the star-shaped \(\R^3\); for \(Q=c_0\) it gives \(\frac12c_0\times z\)),
and \(\Psi_Q\) is a polynomial of degree \(k+1\) with
\(|\Psi_Q(z)|\le\frac12|z|\sup_{B_{|z|}}|Q|\).  Divergence-freeness, the
values on \(B_\rho\) and the support are as in the counterexample.  On
\(B_\rho\), \((L_-+\mu_k)Q_\rho=(L_-+\mu_k)Q=0\); on the annulus,
\((L_-+\mu_k)Q_\rho\) is a sum of terms each of which is a derivative of
\(\chi_\rho\) of order one to three (bounded by \(A\rho^{-1}\) to \(A\rho^{-3}\))
times \(\nu\), \(|z|\le2\rho\) or \(1\), times \(Q\), \(\nabla Q\), \(\Psi_Q\), \(\nabla\Psi_Q\) or
\(\nabla^2\Psi_Q\) evaluated on \(B_{2\rho}\), where these are polynomials
bounded by \(\|Q\|_G|G|^{-1/2}\) times a polynomial in \(\rho\) of degree
at most \(k+1\) with coefficients computable from the Hermite
coefficients of \(Q/\|Q\|_G\); after the powers of \(\rho\) are collected,
\(|(L_-+\mu_k)Q_\rho|\le\|Q\|_G|G|^{-1/2}(1+\nu/\rho^2)\,p_k(\rho)\) on the annulus
with \(p_k\) a polynomial of degree \(k\).  Integrating against \(G\)
over \(\{|z|>\rho\}\) with the tail bound of the counterexample gives
\(\|(L_-+\mu_k)Q_\rho\|_G^2\le\|Q\|_G^2|G|^{-1}(1+\nu/\rho^2)^2p_k(\rho)^2\,8\pi\nu(\rho+2\nu/\rho)e^{-\rho^2/(4\nu)}\).
The energy: \(\|Q_\rho\|_G^2\ge\int_{B_\rho}|Q|^2G\ge\|Q\|_G^2-4\nu(3+2k)\|Q\|_G^2/\rho^2\ge\frac12\|Q\|_G^2\)
for \(\rho\ge\rho_k^{\rm half}\), by Lemma C2-V41-modes
(\(\int|z|^2|Q|^2G\le4\nu(3+2k)\|Q\|_G^2\)) and Chebyshev.  Then
\eqref{eq:c3v22-var-lambda} with \(\lambda=\mu_k\).
\end{proof}

\begin{firewall}[Item (AC1): negative record N-C3-1, no floor from concentration]
\label{fw:c3v22-no-floor-from-concentration}
V19 proposed that a normalized field whose Gaussian mass lies, up to
\(\eta\), in a ball of radius \(R\) has distance at least an explicit
\(v(R,\eta)>0\) to every eigenspace.  This is false for every fixed
level: \(P_\rho/\|P_\rho\|_G\) has all its mass in \(B_{2\rho}\) and
variance at most \(96A^2\rho(4\pi\nu)^{-1/2}e^{-\rho^2/(4\nu)}\), which is smaller
than any prescribed \(v\) once \(\rho\) is large, while its distance to
\(E_0\) is at most the square root of that; and
Proposition~\ref{prop:c3v22-cutoff-level-k} does the same for every
level.  Concentration in the Gaussian sense is compatible with
arbitrarily small variance.  The direction of V19 is closed
negatively on this point; the constants \(v(R,\eta)\) of V19 do not
exist.  What the counterexample also shows is where a floor can
come from: \(P_\rho\) has \(L^6\) norm growing like \(\rho^{1/2}|c_0|\) and
enstrophy growing like \(\rho^{1/2}|c_0|\), so the class bounds cap \(\rho\)
in terms of \(|c_0|\), and \(|c_0|\) is bounded below by the pinching
through \(\|P_\rho\|_G^2\ge\frac12|c_0|^2|G|\).
\end{firewall}

\begin{remark}[The cut-off constant against the class bounds]
\label{rem:c3v22-cutoff-class}
If \(P_\rho\) were the profile of an element of the family at a
homogeneous point, the pinching \(\|P_\rho\|_G^2\ge c_\Phi\) and
\(|P_\rho|\le(1+A)|c_0|\) would give \(|c_0|^2\ge c_\Phi/((1+A)^2|G|)\), and the
class bound \(\|P_\rho\|_{L^6}\le K_6\) would give
\(\rho\le(3/(4\pi))^{1/3}K_6^2(1+A)^2|G|/c_\Phi\); the variance would then be
at least the value of a decreasing function of \(\rho\) at this
radius, once \eqref{eq:c3v22-cutoff-var} were also a lower bound.
The V23 theorem replaces this heuristic by a proof for every profile
of the class, in which the lower bound on the variance comes from
the growth lemma below and the Gaussian weight, not from the
specific form of \(P_\rho\).
\end{remark}

% =============================================================================
\section{The growth lemma}
\label{sec:c3v22-growth}
% =============================================================================

Write \(x=z/(2\nu)^{1/2}\) and \(\mathrm{He}_\alpha(x)=\prod_i\mathrm{He}_{\alpha_i}(x_i)\) for the
probabilists' Hermite polynomials, \(\mathrm{He}_n(t)=t^n+\)(terms of degree
\(\le n-2\)), so that \(\int\mathrm{He}_\alpha\mathrm{He}_\beta\,G\,\dd z=|G|\,\alpha!\,\delta_{\alpha\beta}\).  Every
\(Q\in E_k\) is \(Q(z)=\sum_{|\alpha|=k}c_\alpha\mathrm{He}_\alpha(x)\) with \(c_\alpha\in\R^3\), and
\(\|Q\|_G^2=|G|\sum_\alpha\alpha!|c_\alpha|^2\).  Its top part is
\(Q^{\rm top}(z):=\sum_{|\alpha|=k}c_\alpha x^\alpha\), homogeneous of degree \(k\), and
\(Q^{\rm low}:=Q-Q^{\rm top}\) is a polynomial of degree at most \(k-2\) in \(x\).
For \(k\ge2\) define the finite-dimensional constants
\begin{equation}
 \gamma_k^{\rm top}:=\min\Bigl\{\int_{|x|\le1}|Q^{\rm top}|^2\dd x:\ \sum_\alpha\alpha!|c_\alpha|^2=1\Bigr\}>0,\qquad
 \Lambda_k:=\max\Bigl\{\sum_{j\le k-2}\sup_{|x|=1}|Q^{\rm low}_j|:\ \sum_\alpha\alpha!|c_\alpha|^2=1\Bigr\},
 \label{eq:c3v22-fd-constants}
\end{equation}
where \(Q^{\rm low}=\sum_jQ^{\rm low}_j\) is the decomposition into
homogeneous parts (the minimum is positive because a nonzero
polynomial does not vanish on a ball, and both extrema are attained
on a sphere of a finite-dimensional space).

\begin{lemma}[Growth of normalized eigenfunctions on large balls]
\label{lem:c3v22-growth}
Let \(Q\in E_k\) with \(\|Q\|_G=1\).  Then \(\int_{B_\rho}|Q|^2\dd z\ge\gamma_k\rho^{2k+3}\)
for all \(\rho\ge\rho_k^{\rm top}\), where
\begin{align}
 &k=0:\ \ \gamma_0=\frac{4\pi}{3|G|},\quad\rho_0^{\rm top}=0;\qquad
 k=1:\ \ \gamma_1=\frac{4\pi}{15\cdot2\nu|G|},\quad\rho_1^{\rm top}=0;
 \label{eq:c3v22-growth-01}\\
 &k\ge2:\ \ \gamma_k=\frac{\gamma_k^{\rm top}}{4(2\nu)^k|G|},\quad
 \rho_k^{\rm top}=(2\nu)^{1/2}\max\Bigl(1,\Bigl(\frac{16\pi\Lambda_k^2}{3\gamma_k^{\rm top}}\Bigr)^{1/4}\Bigr).
 \label{eq:c3v22-growth-k}
\end{align}
\end{lemma}

\begin{proof}
\(k=0\): \(Q=c_0\) with \(|c_0|^2|G|=1\), and \(\int_{B_\rho}|c_0|^2=\frac{4\pi}3\rho^3/|G|\).
\(k=1\): \(Q=Az\) with \(\operatorname{tr}A=0\) (the divergence-free linear fields),
\(\|Az\|_G^2=\sum_{ijl}A_{ij}A_{il}\int z_jz_lG=2\nu|G||A|_F^2=1\), and
\(\int_{B_\rho}|Az|^2=|A|_F^2\int_{B_\rho}z_1^2=|A|_F^2\frac13\int_{B_\rho}|z|^2=|A|_F^2\frac{4\pi}{15}\rho^5\).
\(k\ge2\): in the variable \(x\), \(\dd z=(2\nu)^{3/2}\dd x\), \(B_\rho=\{|x|\le R\}\) with
\(R=\rho/(2\nu)^{1/2}\), and \(\sum\alpha!|c_\alpha|^2=1/|G|\).  By homogeneity
\(\int_{|x|\le R}|Q^{\rm top}|^2\dd x=R^{2k+3}\int_{|x|\le1}|Q^{\rm top}|^2\ge R^{2k+3}\gamma_k^{\rm top}/|G|\),
and \(|Q^{\rm low}(x)|\le\Lambda_k|G|^{-1/2}\max(1,|x|)^{k-2}\) (each homogeneous
part \(Q^{\rm low}_j\) is bounded by \(|x|^j\) times its supremum on the
sphere), so \(\int_{|x|\le R}|Q^{\rm low}|^2\le\frac{4\pi}3\Lambda_k^2|G|^{-1}R^{2k-1}\) for
\(R\ge1\).  With \(|Q|^2\ge\frac12|Q^{\rm top}|^2-|Q^{\rm low}|^2\),
\[
 \int_{|x|\le R}|Q|^2\dd x\ \ge\ \frac{R^{2k+3}}{|G|}\Bigl(\frac12\gamma_k^{\rm top}-\frac{4\pi}3\Lambda_k^2R^{-4}\Bigr)\ \ge\ \frac{\gamma_k^{\rm top}R^{2k+3}}{4|G|}
\]
when \(R^4\ge16\pi\Lambda_k^2/(3\gamma_k^{\rm top})\) and \(R\ge1\).  Returning to
\(z\), \(\int_{B_\rho}|Q|^2\dd z=(2\nu)^{3/2}\int_{|x|\le R}|Q|^2\dd x\ge\gamma_k^{\rm top}\rho^{2k+3}/(4(2\nu)^k|G|)\).
\end{proof}

\begin{remark}[Status of the constants]
\label{rem:c3v22-constants}
\(\gamma_0,\gamma_1\) are explicit in \(\nu\).  For \(2\le k\le K_r\) the
constants \(\gamma_k^{\rm top}\) and \(\Lambda_k\) are extrema of explicit
quadratic and piecewise-smooth functions on the unit sphere of a
space of dimension \(3\binom{k+2}{2}\), computable numerically for each
\(k\) but not given in closed form here.  The floor of V23 is
therefore explicit in closed form when \(K_r\le1\), i.e.\ when
\(C_r<\frac54\), and explicit up to these finite-dimensional constants
otherwise; in either case it replaces the compactness constant
\(v_0\) of C2-V34, whose dependence on the family was not
quantified at all.
\end{remark}

% =============================================================================
\section{Integrated V22 conclusion and next acquisition order}
\label{sec:c3v22-conclusion}
% =============================================================================

\begin{theorem}[What V22 proves]
\label{thm:c3v22-conclusion}
Relative to the two legacy cycles and V1--V21: the cut-off constant
(Counterexample~\ref{cex:c3v22-cutoff-constant}), the cut-off
eigenfunctions (Proposition~\ref{prop:c3v22-cutoff-level-k}), the
negative record N-C3-1, and the growth lemma
(Lemma~\ref{lem:c3v22-growth}).  The full Navier--Stokes
stability/global-regularity theorem remains unproved.
\end{theorem}

\begin{proof}
The cited results.
\end{proof}

\begin{programme}[V23 acquisition order]
\label{prog:c3v23}
\begin{enumerate}[label=\textup{(AC\arabic*)},leftmargin=4.8em]
\item The explicit floor: for every profile of the family at a
      homogeneous point and every \(\sigma\),
      \(\mathrm{Var}\ge v_{\rm exp}:=\min(\frac1{32},\min_{k\le K_r}v^{(k)})\) with
      \(v^{(k)}=\frac{\gamma_k}{128}(\rho_k^*)^{2k+3}e^{-(\rho_k^*)^2/(4\nu)}\),
      \(\rho_k^*=\max(\rho_k^{\rm top},\rho_k^{(6)})\) and
      \(\rho_k^{(6)}\) the radius at which the growth of
      Lemma~\ref{lem:c3v22-growth} exceeds eight times the \(L^6\)-allowed
      unweighted energy \((4\pi/3)^{2/3}\rho^2K_6^2/c_\Phi\); proof by
      comparing \(\int_{B_\rho}|\Pi_{k_1}\widehat V|^2\) from above (H\"older with
      \(K_6\), the Gaussian closeness \eqref{eq:c3v21-nearly} and
      \(G\ge e^{-\rho^2/(4\nu)}\) on \(B_\rho\)) and from below (the growth lemma).
\item Record that the floor uses only the class \(L^6\) bound, the
      pinching and the Rayleigh bound, not the profile equation, and
      compare with the compactness proof of C2-V34.
\item V24: the push budget with \(v_{\rm exp}\); V25 compulsory
      companion audit.
\end{enumerate}
\end{programme}

\begin{assessment}[Position after V22]
\label{ass:c3v22-position}
Concentration gives no floor; growth on large balls and the
unweighted class bounds are the two halves of one.  No full
stability or global-regularity claim is made.
\end{assessment}

% =============================================================================
\section*{Version V22 (third cycle) append-only record}
\addcontentsline{toc}{section}{Version V22 (third cycle) append-only record}
% =============================================================================

The complete V21 source of the third cycle was retained before this
continuation.  Its SHA--256 digest was
\begin{center}\scriptsize\ttfamily
3b4009f4720d4d17b72e4db1e338c5812635afeee3b60f0ed0444aa5aedb0119.
\end{center}
V22 recorded the cut-off eigenfunctions, the negative record on
concentration, and the growth lemma.  No full regularity claim is
made.

% =============================================================================
\section{V23: the explicit variance floor}
\label{sec:c3v23-entry}
% =============================================================================

Items (AC1) and (AC2) of Programme~\ref{prog:c3v23} are carried out.
The spectral variance of every profile of the family at a homogeneous
type-I point is bounded below, at every scale, by an explicit
positive constant depending on the viscosity, the class \(L^6\) bound,
the pinching constant and the Rayleigh bound, in closed form when the
dominant level can only be \(0\) or \(1\) and up to finite-dimensional
Hermite constants otherwise.  The proof compares the unweighted
energy of the dominant polynomial part of the normalized profile in
a large ball from above, by the class \(L^6\) bound and the Gaussian
closeness, and from below, by the growth lemma; the radius at which
the two cross is capped by the class, and the Gaussian tail at that
radius is the floor.  The floor is exponentially small in the square
of the capped radius over the viscosity, and it uses the profile
equation nowhere: it is a property of pinched, \(L^6\)-bounded
divergence-free fields.  No global regularity claim is made.  The
next compulsory companion audit is V25.

% =============================================================================
\section{The floor}
\label{sec:c3v23-floor}
% =============================================================================

With \(\gamma_k,\rho_k^{\rm top}\) of Lemma~\ref{lem:c3v22-growth}, \(K_6\) the class
\(L^6\) bound (C2-V36), \(c_\Phi\) the pinching constant and \(K_r=\lfloor2C_r-\frac12\rfloor\)
(Proposition~\ref{prop:c3v21-bounded}), define for \(0\le k\le K_r\)
\begin{equation}
 \rho_k^{(6)}:=\Bigl(\frac{8(4\pi/3)^{2/3}K_6^2}{\gamma_k\,c_\Phi}\Bigr)^{1/(2k+1)},\qquad
 \rho_k^*:=\max(\rho_k^{\rm top},\rho_k^{(6)}),\qquad
 v^{(k)}:=\frac{\gamma_k}{128}\,(\rho_k^*)^{2k+3}\,e^{-(\rho_k^*)^2/(4\nu)} ,
 \label{eq:c3v23-constants}
\end{equation}
and
\begin{equation}
 v_{\rm exp}:=\min\Bigl(\frac1{32},\ \min_{0\le k\le K_r}v^{(k)}\Bigr)\ >\ 0 .
 \label{eq:c3v23-vexp}
\end{equation}

\begin{theorem}[Explicit variance floor at homogeneous points]
\label{thm:c3v23-floor}
Let \(x_0\) be a homogeneous type-I point.  For every \(W\in\mathcal F(x_0)\)
and every \(\sigma\), the profile \(V\) satisfies
\begin{equation}
 \boxed{\quad
 \mathrm{Var}(V)\ \ge\ v_{\rm exp} .
 \quad}
 \label{eq:c3v23-floor}
\end{equation}
In particular the compactness constant \(v_0\) of Theorem
C2-V34-variance may be replaced by \(v_{\rm exp}\) in every statement
of the second cycle in which it appears.
\end{theorem}

\begin{proof}
Suppose \(\mathrm{Var}(V)<\frac1{32}\) at some \(\sigma\); otherwise there is
nothing to prove.  By Corollary~\ref{cor:c3v21-nearly-polynomial},
with \(\eta=16\,\mathrm{Var}<\frac12\), \(k:=k_1\le K_r\) and \(P=\Pi_k\widehat V\),
\(\|\widehat V-P\|_G^2\le\eta\) and \(\|P\|_G^2\ge1-\eta>\frac12\).  Let
\(Q:=P/\|P\|_G\in E_k\), \(\|Q\|_G=1\), and \(\rho:=\rho_k^*\).

\emph{Lower bound.}  Since \(\rho\ge\rho_k^{\rm top}\), Lemma~\ref{lem:c3v22-growth}
gives \(\int_{B_\rho}|P|^2\dd z=\|P\|_G^2\int_{B_\rho}|Q|^2\dd z\ge\frac12\gamma_k\rho^{2k+3}\).

\emph{Upper bound.}  \(\int_{B_\rho}|P|^2\le2\int_{B_\rho}|\widehat V|^2+2\int_{B_\rho}|\widehat V-P|^2\).
By H\"older and the class bound, with \(\|\widehat V\|_{L^6}=\|V\|_{L^6}/\varphi^{1/2}\le K_6/c_\Phi^{1/2}\)
(pinching \(\varphi\ge c_\Phi\) at a homogeneous point),
\(\int_{B_\rho}|\widehat V|^2\le|B_\rho|^{2/3}\|\widehat V\|_{L^6}^2\le(4\pi/3)^{2/3}\rho^2K_6^2/c_\Phi\).
Since \(G\ge e^{-\rho^2/(4\nu)}\) on \(B_\rho\),
\(\int_{B_\rho}|\widehat V-P|^2\le e^{\rho^2/(4\nu)}\|\widehat V-P\|_G^2\le e^{\rho^2/(4\nu)}\eta\).

\emph{Comparison.}  Hence
\(\frac12\gamma_k\rho^{2k+3}\le2(4\pi/3)^{2/3}\rho^2K_6^2/c_\Phi+2e^{\rho^2/(4\nu)}\eta\).
Since \(\rho\ge\rho_k^{(6)}\), \(\gamma_k\rho^{2k+1}\ge8(4\pi/3)^{2/3}K_6^2/c_\Phi\), so the
first term on the right is at most \(\frac14\gamma_k\rho^{2k+3}\).  Therefore
\(\frac14\gamma_k\rho^{2k+3}\le2e^{\rho^2/(4\nu)}\eta\), i.e.\
\(\eta\ge\frac18\gamma_k\rho^{2k+3}e^{-\rho^2/(4\nu)}\), and
\(\mathrm{Var}=\eta/16\ge v^{(k)}\ge\min_{k'\le K_r}v^{(k')}\).  In both cases
\(\mathrm{Var}\ge v_{\rm exp}\).
\end{proof}

\begin{corollary}[The closed-form case]
\label{cor:c3v23-closed-form}
If \(C_r<\frac54\) then \(K_r\le1\) and, with \(|G|=(4\pi\nu)^{3/2}\),
\begin{align}
 \rho_0^*&=8\Bigl(\frac{4\pi}3\Bigr)^{-1/3}\frac{|G|K_6^2}{c_\Phi},&
 v^{(0)}&=\frac{\pi}{96|G|}(\rho_0^*)^3e^{-(\rho_0^*)^2/(4\nu)},
 \label{eq:c3v23-v0}\\
 \rho_1^*&=\Bigl(120\Bigl(\frac{4\pi}3\Bigr)^{2/3}\frac{2\nu|G|K_6^2}{4\pi\,c_\Phi}\Bigr)^{1/3},&
 v^{(1)}&=\frac{\pi}{960\,\nu|G|}(\rho_1^*)^5e^{-(\rho_1^*)^2/(4\nu)},
 \label{eq:c3v23-v1}
\end{align}
and \(v_{\rm exp}=\min(\frac1{32},v^{(0)},v^{(1)})\) is explicit in \(\nu\), \(K_6\), \(c_\Phi\).
\end{corollary}

\begin{proof}
\(K_r=\lfloor2C_r-\frac12\rfloor\le1\) iff \(2C_r-\frac12<2\).  For \(k=0\):
\(\rho_0^{\rm top}=0\), \(\gamma_0=4\pi/(3|G|)\), so
\(\rho_0^{(6)}=8(4\pi/3)^{2/3}K_6^2\cdot3|G|/(4\pi c_\Phi)=8(4\pi/3)^{-1/3}|G|K_6^2/c_\Phi\), and
\(\gamma_0/128=\pi/(96|G|)\).  For \(k=1\): \(\rho_1^{\rm top}=0\),
\(\gamma_1=4\pi/(15\cdot2\nu|G|)\), so
\((\rho_1^{(6)})^3=8(4\pi/3)^{2/3}K_6^2\cdot15\cdot2\nu|G|/(4\pi c_\Phi)\), and
\(\gamma_1/128=4\pi/(128\cdot30\nu|G|)=\pi/(960\nu|G|)\).
\end{proof}

\begin{remark}[Size of the floor]
\label{rem:c3v23-size}
The floor is explicit but exponentially small.  Since
\(c_\Phi\le\varphi\le\|V\|_{L^6}^2\|G\|_{L^{3/2}}=K_6^2\cdot\frac{8\pi\nu}3\), one has
\(K_6^2/c_\Phi\ge3/(8\pi\nu)\) and hence
\(\rho_0^*\ge12(4\pi/3)^{-1/3}(4\pi\nu)^{1/2}\), so \((\rho_0^*)^2/(4\nu)\ge144(4\pi/3)^{-2/3}\pi>170\):
even when the pinching is as strong as the class allows, \(v^{(0)}\)
is below \(e^{-170}\), and it decreases like
\(\exp(-c\,\nu^2K_6^4/c_\Phi^2)\) as the pinching weakens.  This is the
price of a proof that uses only the \(L^6\) bound: the class permits a
nearly constant Gaussian core of radius up to \(\rho_0^*\), and only
the Gaussian tail beyond it is charged.  The value matters for the
push budget (V24), where it enters as \(4v_{\rm exp}/C_r\), the
threshold on the mean outward normalized flux.
\end{remark}

% =============================================================================
\section{What the floor uses, and what it does not}
\label{sec:c3v23-uses}
% =============================================================================

\begin{firewall}[Item (AC2): the equation is not used]
\label{fw:c3v23-no-equation}
Theorem~\ref{thm:c3v23-floor} uses, of the profile, only: the
Gaussian energy is pinched (\(\varphi\ge c_\Phi\)), the unweighted \(L^6\)
norm is bounded (\(\|V\|_{L^6}\le K_6\)), and the Rayleigh quotient is
bounded (\(r\le C_r\), itself a consequence of the pinching, the class
enstrophy bound and \(\varphi\le C_\Phi\)).  It holds for every
divergence-free field with these three properties.  The proof of
C2-V34 used the same three facts (the \(L^6\) bound to exclude
polynomial fields, the class bounds for lower semicontinuity of the
variance) and compactness of the family to make the positivity
uniform; the present proof quantifies the exclusion of polynomial
fields and dispenses with compactness.  Nothing about the
Navier--Stokes nonlinearity enters, so the floor cannot by itself
distinguish a profile of the equation from any other pinched
\(L^6\)-bounded field; its role is confined to the identities of the
Rayleigh chapter in which the variance appears with a sign (the
descent of the Rayleigh quotient, C2-V31; the push budget, C2-V33--V34).
\end{firewall}

\begin{proposition}[Pointwise descent of the Rayleigh quotient]
\label{prop:c3v23-descent}
For every \(W\in\mathcal F(x_0)\) at a homogeneous point and every \(\sigma\),
\begin{equation}
 \frac{\dd r}{\dd\sigma}\ \le\ -2v_{\rm exp}+2\varphi^{1/2}\langle\mathcal R,\widehat N\rangle_G ,
 \label{eq:c3v23-descent}
\end{equation}
with \(\mathcal R\) and \(\widehat N\) as in C2-V31: unless the nonlinear
correlation on the right exceeds \(v_{\rm exp}/\varphi^{1/2}\), the
Rayleigh quotient decreases at the explicit rate \(2v_{\rm exp}\).
\end{proposition}

\begin{proof}
The Rayleigh identity of C2-V31, \(\frac{\dd r}{\dd\sigma}=-2\mathrm{Var}+2\varphi^{1/2}\langle\mathcal R,\widehat N\rangle_G\),
with \eqref{eq:c3v23-floor}.
\end{proof}

\begin{remark}[The nearly-constant core is not excluded]
\label{rem:c3v23-core}
The theorem does not say that a profile cannot have a nearly constant
Gaussian core; it says that the core's radius is capped by the class
and that the variance is at least the Gaussian tail beyond the cap.
Whether the profile equation excludes a nearly constant core of
radius comparable to \((4\nu)^{1/2}\), which would raise the floor to a
value of order one, is a different question, involving the
nonlinearity; it is the natural continuation of this block and is
recorded in the direction of V24.
\end{remark}

% =============================================================================
\section{Integrated V23 conclusion and next acquisition order}
\label{sec:c3v23-conclusion}
% =============================================================================

\begin{theorem}[What V23 proves]
\label{thm:c3v23-conclusion}
Relative to the two legacy cycles and V1--V22: the explicit variance
floor (Theorem~\ref{thm:c3v23-floor}), its closed form when \(C_r<\frac54\)
(Corollary~\ref{cor:c3v23-closed-form}), the record that the equation
is not used, and the pointwise descent
(Proposition~\ref{prop:c3v23-descent}).  The full Navier--Stokes
stability/global-regularity theorem remains unproved.
\end{theorem}

\begin{proof}
The cited results.
\end{proof}

\begin{programme}[V24 acquisition order]
\label{prog:c3v24}
\begin{enumerate}[label=\textup{(AC\arabic*)},leftmargin=4.8em]
\item The push budget with \(v_{\rm exp}\): restate Theorem
      C2-V34-mean-push with the explicit constant, and the explicit
      threshold \(4v_{\rm exp}/C_r\) on the mean outward normalized flux;
      record the status of the condition (open) and the size of the
      threshold.
\item Consolidate V21--V23 as the explicit-floor chapter, with the
      negative record N-C3-1 and the correction of direction of V21,
      and fix the direction for V26--V30: the nearly constant core
      against the equation (the momentum identity of V18 on a core
      with \(\bar V\ne0\), and the mean-zero property of the profile in
      the sense of C2-V16).
\item V25 compulsory companion audit and delivery.
\end{enumerate}
\end{programme}

\begin{assessment}[Position after V23]
\label{ass:c3v23-position}
The variance floor of the Rayleigh chapter is explicit; it is
exponentially small and uses no property of the equation.  No full
stability or global-regularity claim is made.
\end{assessment}

% =============================================================================
\section*{Version V23 (third cycle) append-only record}
\addcontentsline{toc}{section}{Version V23 (third cycle) append-only record}
% =============================================================================

The complete V22 source of the third cycle was retained before this
continuation.  Its SHA--256 digest was
\begin{center}\scriptsize\ttfamily
1a99ac23d2a5eba6d5297ec372fa9113b749e8745ac4d6ba6187d9316decaf43.
\end{center}
V23 proved the explicit variance floor and its closed form, and
recorded what it uses.  No full regularity claim is made.

% =============================================================================
\section{V24: the push budget with the explicit floor, and the explicit-floor chapter}
\label{sec:c3v24-entry}
% =============================================================================

Items (AC1) and (AC2) of Programme~\ref{prog:c3v24} are carried out.
The mean push budget of the second cycle is restated with the
explicit floor: the mean geometric push is positive whenever the
mean positive part of the normalized outward flux is below an
explicit threshold, and the threshold is so small that the condition
cannot be expected to be verified by any bound of the class; the
route through the floor alone is recorded as closed for practical
purposes unless the floor is raised by the equation.  A profile of
small variance is shown to be close to its dominant polynomial part
in the Gaussian Sobolev sense as well, with the deviation linear in
the variance.  The chapter V21--V23 is consolidated, and the direction
for V26--V30 is fixed: a mean floor of order one under a
constant-chasing condition, and the nearly constant core against the
momentum identity.  No global regularity claim is made.  The next
compulsory companion audit is V25.

% =============================================================================
\section{The push budget with the explicit floor}
\label{sec:c3v24-push}
% =============================================================================

\begin{theorem}[Mean geometric push with the explicit floor]
\label{thm:c3v24-mean-push}
Let \(x_0\) be a homogeneous type-I point and \(\mathfrak m\) an invariant
probability measure on \(\mathcal F(x_0)\).  With the notation of
C2-V33--V34 (\(\mathcal P\) the geometric push, \(\widehat\jmath=j/\varphi\) the
normalized flux, \(C_r\) the Rayleigh bound),
\begin{equation}
 \int\frac{2\mathcal P}{\varphi}\,\dd\mathfrak m\ \ge\ 2v_{\rm exp}+\frac14\int(-\widehat\jmath)_+\,\dd\mathfrak m-\frac{C_r}2\int(\widehat\jmath)_+\,\dd\mathfrak m ,
 \label{eq:c3v24-mean-push}
\end{equation}
and the mean geometric push is positive whenever
\(\int(\widehat\jmath)_+\dd\mathfrak m<4v_{\rm exp}/C_r\).
\end{theorem}

\begin{proof}
Theorem C2-V34-mean-push with \(v_0\) replaced by \(v_{\rm exp}\), which
Theorem~\ref{thm:c3v23-floor} permits.
\end{proof}

\begin{firewall}[Item (AC1): the threshold is explicit and out of reach]
\label{fw:c3v24-threshold}
The condition of C2-V34 (the blow-up statistics is outward-flux-poor)
is unchanged in status: open.  What the explicit constant adds is
its size.  By Remark~\ref{rem:c3v23-size}, \(4v_{\rm exp}/C_r<8e^{-170}\) in
the closed-form case and smaller otherwise, while the only a priori
bound on \(\int(\widehat\jmath)_+\dd\mathfrak m\) is \(C_J/c_\Phi\), of order one.
The mean-push route through the variance floor alone would
therefore require the mean outward normalized flux of a homogeneous
blow-up to be smaller than a number below \(e^{-170}\); no estimate of
the series or of the imports (E1)--(E9) produces such a bound, and
no mechanism in the ledgers suggests one.  The route is recorded as
closed for practical purposes: it reopens only if the floor is
raised to order one, which, by Firewall~\ref{fw:c3v23-no-equation},
requires the equation.
\end{firewall}

% =============================================================================
\section{Closeness in the Gaussian Sobolev sense}
\label{sec:c3v24-sobolev}
% =============================================================================

\begin{lemma}[A profile of small variance is nearly polynomial in \(H^1(G)\)]
\label{lem:c3v24-h1}
In the situation of Corollary~\ref{cor:c3v21-nearly-polynomial}
(\(\mathrm{Var}<\frac1{32}\), \(\eta=16\,\mathrm{Var}\), \(P=\Pi_{k_1}\widehat V\)),
\begin{equation}
 \nu\|\nabla(\widehat V-P)\|_G^2\ \le\ \Bigl(\frac18+2C_r^2\Bigr)^{1/2}\eta ,
 \qquad
 \bigl|r-\mu_{k_1}\bigr|\ \le\ \frac{\eta}2 .
 \label{eq:c3v24-h1}
\end{equation}
In particular, when \(k_1=0\), \(\nu\|\nabla\widehat V\|_G^2\le(\frac18+2C_r^2)^{1/2}\eta\):
a profile whose spectral mass is nearly on the constant level is
nearly constant in \(H^1(G)\), with the Gaussian dissipation of the
normalized profile linear in the variance.
\end{lemma}

\begin{proof}
By Lemma C2-V41-modes, \(\nu\|\nabla\Pi_k\widehat V\|_G^2=\frac k2\hat a_k\), and the
projections are \(G\)-orthogonal together with their gradients
(the eigenspaces of the self-adjoint \(-L_-\) are orthogonal in the
form domain), so
\(\nu\|\nabla(\widehat V-P)\|_G^2=\sum_{k\ne k_1}\frac k2\hat a_k\le\sum_{k\ne k_1}\mu_k\hat a_k\).
By Cauchy--Schwarz,
\(\sum_{k\ne k_1}\mu_k\hat a_k\le\bigl(\sum_{k\ne k_1}\mu_k^2\hat a_k\bigr)^{1/2}\bigl(\sum_{k\ne k_1}\hat a_k\bigr)^{1/2}\),
and \(\sum_{k\ne k_1}\mu_k^2\hat a_k\le2\sum_k(\mu_k-r)^2\hat a_k+2r^2\sum_{k\ne k_1}\hat a_k\le2\,\mathrm{Var}+2C_r^2\eta=(\frac18+2C_r^2)\eta\),
while \(\sum_{k\ne k_1}\hat a_k\le\eta\).  For the second bound,
\(r-\mu_{k_1}=\sum_{k\ne k_1}(\mu_k-\mu_{k_1})\hat a_k\) and
\(|\mu_k-\mu_{k_1}|\le|\mu_k-r|+|r-\mu_{k_1}|\le|\mu_k-r|+\frac14\), so
\(|r-\mu_{k_1}|\le\sum_{k\ne k_1}|\mu_k-r|\hat a_k+\frac14\eta\le\mathrm{Var}^{1/2}\eta^{1/2}+\frac14\eta=\frac12\eta\),
by Cauchy--Schwarz and \(\mathrm{Var}=\eta/16\).
\end{proof}

\begin{remark}[The nonlinearity on a nearly constant core]
\label{rem:c3v24-nonlinearity}
When \(k_1=0\), the advective part of the nonlinearity is small in
\(L^2(G)\): \(\|(V\cdot\nabla)V\|_G\le C_\infty\varphi^{1/2}\|\nabla\widehat V\|_G\le C_\infty\varphi^{1/2}(\eta/\nu)^{1/2}(\frac18+2C_r^2)^{1/4}\),
of the order of the square root of the variance.  The pressure
gradient is not controlled by this: it is nonlocal and bounded only
by the class (\(K_P\), C2-V37).  The momentum identity of C2-V17 and
V18, \(\frac{\dd}{\dd\sigma}\bar V=-\frac12\bar V-\overline{N(V)}\), on a
nearly constant core with \(|G||\bar V|^2\ge(1-\eta)\varphi\), therefore
reads: the large Gaussian momentum decays at rate \(\frac12\) unless
the mean pressure gradient sustains it, the advective mean being
small.  Whether the pressure of a nearly constant divergence-free
profile can have a Gaussian mean of order \(|\bar V|\) is the question
the next block asks.
\end{remark}

% =============================================================================
\section{The explicit-floor chapter, and the direction for V26--V30}
\label{sec:c3v24-consolidated}
% =============================================================================

\begin{theorem}[The explicit-floor chapter, V21--V23]
\label{thm:c3v24-consolidated}
At a homogeneous type-I point, for every profile of the family and
every \(\sigma\):
\begin{enumerate}[label=\textup{(\arabic*)},leftmargin=3.2em]
\item (One level.)  At most one eigenvalue lies within \(\frac14\) of
      the Rayleigh quotient; \(\mathrm{Var}\ge\frac1{16}(1-\hat a_{k_1})\) and
      \(\mathrm{Var}\ge\frac14(1-\hat a_{k_0}-\hat a_{k_0+1})\) (V21).
\item (Bounded level.)  \(k_1\le K_r=\lfloor2C_r-\frac12\rfloor\) and
      \(\sum_{k\ge K}\hat a_k\le2C_r/(K+1)\), pointwise (V21); the high case
      of the V20 dichotomy is empty, and the Hermite-mass computation
      is not needed (correction of direction, V21).
\item (No floor from concentration.)  Divergence-free fields
      concentrated in a ball have variance as small as
      \(96A^2\rho(4\pi\nu)^{-1/2}e^{-\rho^2/(4\nu)}\) (negative record N-C3-1, V22).
\item (Growth.)  A normalized level-\(k\) eigenfunction has unweighted
      energy at least \(\gamma_k\rho^{2k+3}\) in \(B_\rho\) beyond \(\rho_k^{\rm top}\),
      explicitly for \(k\le1\) (V22).
\item (Floor.)  \(\mathrm{Var}\ge v_{\rm exp}>0\) explicit in \(\nu\), \(K_6\), \(c_\Phi\),
      \(C_r\) (and the finite-dimensional constants for \(2\le k\le K_r\)),
      exponentially small, using the \(L^6\) bound, the pinching and
      the Rayleigh bound only (V23).
\item (Sobolev closeness.)  A profile of small variance is nearly
      its dominant polynomial part in \(H^1(G)\), with deviation
      linear in the variance (V24).
\item (Push budget.)  The mean geometric push is positive when the
      mean outward normalized flux is below \(4v_{\rm exp}/C_r\), a
      threshold below \(8e^{-170}\); the route is closed for practical
      purposes unless the floor is raised by the equation (V24).
\end{enumerate}
\end{theorem}

\begin{proof}
Propositions~\ref{prop:c3v21-criteria}, \ref{prop:c3v21-bounded};
Corollary~\ref{cor:c3v21-nearly-polynomial};
Counterexample~\ref{cex:c3v22-cutoff-constant};
Lemma~\ref{lem:c3v22-growth}; Theorem~\ref{thm:c3v23-floor};
Lemma~\ref{lem:c3v24-h1}; Theorem~\ref{thm:c3v24-mean-push}.
\end{proof}

\begin{assessment}[Direction for V26--V30]
\label{ass:c3v24-direction}
Two continuations are open, and the block V26--V30 takes them in
order.

\emph{(a) A mean floor of order one under a constant-chasing
condition.}  The mean spectrum of C2-V41 bounds \(\langle a_k\rangle\le B_k\)
explicitly for each \(k\), with \(B_k\) of the order of
\((\kappa^2\langle d+3\varphi\rangle+K_P^2)/\mu_k^2\).  Since
\(\langle\mathrm{Var}\rangle\ge\frac1{16}\langle1-\hat a_{k_1}\rangle\) and \(k_1\le K_r\),
\(\langle\hat a_{k_1}\rangle\le\sum_{k\le K_r}\langle\hat a_k\rangle\le\sum_{k\le K_r}B_k/c_\Phi\),
so that \(\langle\mathrm{Var}\rangle\ge\frac1{16}(1-\sum_{k\le K_r}B_k/c_\Phi)\) is an
order-one mean floor whenever \(\sum_{k\le K_r}B_k<c_\Phi\).  The
condition is a numerical one on the class constants; since
\(B_0\ge12\kappa^2c_\Phi\) it requires at least \(\kappa^2<\frac1{12}\), i.e.\
\(C_\infty<(\nu/12)^{1/2}\), which is not an a priori bound of the
class.  V26 records the conditional statement exactly and examines
whether the constant-chasing constraint of C2-V41 (\(\kappa<\sqrt2\))
can be improved to it; the expected answer is negative.

\emph{(b) The nearly constant core against the equation.}  On the
event that the dominant level is \(0\), the profile is nearly
constant in \(H^1(G)\) with large Gaussian momentum, and the momentum
identity says that this momentum decays at rate \(\frac12\) unless the
mean pressure gradient sustains it.  V27--V28 compute the Gaussian
mean of the pressure gradient for a nearly constant divergence-free
field (the pressure is determined by \(\nabla V\otimes\nabla V\), which is
small in \(L^2(G)\) but not in the unweighted sense) and determine
whether a nearly constant state is transient, which would convert
the pointwise floor into a floor on the log-average over windows of
explicit length, or whether the far field can sustain it.  V29
consolidates; V30 audits.  No global regularity claim is made.
\end{assessment}

% =============================================================================
\section{Integrated V24 conclusion and next acquisition order}
\label{sec:c3v24-conclusion}
% =============================================================================

\begin{theorem}[What V24 establishes]
\label{thm:c3v24-conclusion}
Relative to the two legacy cycles and V1--V23: the explicit push
budget (Theorem~\ref{thm:c3v24-mean-push}), the Sobolev closeness
(Lemma~\ref{lem:c3v24-h1}), the consolidated chapter
(Theorem~\ref{thm:c3v24-consolidated}) and the direction.  The full
Navier--Stokes stability/global-regularity theorem remains unproved.
\end{theorem}

\begin{proof}
The cited results.
\end{proof}

\begin{programme}[V25 acquisition order]
\label{prog:c3v25}
\begin{enumerate}[label=\textup{(AC\arabic*)},leftmargin=4.8em]
\item The compulsory companion audit: read the current SM and YM
      notes and record their digests and decision.
\item Record the position after twenty-five versions; delivery.
\end{enumerate}
\end{programme}

\begin{assessment}[Position after V24]
\label{ass:c3v24-position}
The explicit-floor chapter is consolidated; the push-budget route
through the floor alone is closed for practical purposes; the
direction is the mean floor and the nearly constant core.  No full
stability or global-regularity claim is made.
\end{assessment}

% =============================================================================
\section*{Version V24 (third cycle) append-only record}
\addcontentsline{toc}{section}{Version V24 (third cycle) append-only record}
% =============================================================================

The complete V23 source of the third cycle was retained before this
continuation.  Its SHA--256 digest was
\begin{center}\scriptsize\ttfamily
b2961de0cbe0fe9c672c92ba0383e84fbed33dcb833977ed927e437434f21752.
\end{center}
V24 restated the push budget with the explicit floor, proved the
Sobolev closeness, and consolidated the chapter.  No full regularity
claim is made.

% =============================================================================
\section{V25: compulsory companion audit and the position after twenty-five versions}
\label{sec:c3v25-entry}
% =============================================================================

This is the twenty-fifth version of the third cycle.  The compulsory
review of the two companion programmes is performed first, and the
position after twenty-five versions is recorded with the plan for
the block V26--V30.  No global regularity claim is made.  The next
compulsory companion audit is V30.

% =============================================================================
\section{Compulsory V25 review of the current companion notes}
\label{sec:c3v25-companion-audit}
% =============================================================================

The current companion files in \texttt{latex\_notes} were inspected
read-only.  Both programmes have moved since the V20 audit: the SM
programme has advanced from V24 to V27 of its seventh series, and the
YM programme has closed its fifth series at V100 (archived with the
suffix \texttt{\_f5}) and opened a sixth series, of which the newest
file is V4.  The frozen checkpoints are
\begin{align}
 &\texttt{Renewal\_SM\_Einstein\_Continuum\_Research\_Notes\_V27.tex},
 \notag\\[-2pt]
 &\qquad
 \texttt{2860a332030636c9b54c08f5969c3e87d1f92df77620d286509b8989cb186412},
 \label{eq:c3v25-SM-checkpoint}\\
 &\texttt{Renewal\_Yang\_Mills\_Mass\_Gap\_Research\_Notes\_V4.tex},
 \notag\\[-2pt]
 &\qquad
 \texttt{c98bf0ec6cb7010db9a6d76b3a54e530cda2952c7bd9d25bf9e9a194d7446caf}.
 \label{eq:c3v25-YM-checkpoint}
\end{align}
(SM V27 has 11368 lines; YM V4 has 1338 lines.  The folder also
contains a file \texttt{V3} of the YM sixth series byte-identical to
\texttt{V4}; the audit records the newer name.)

\emph{SM V25--V27.}  The SM programme constructed a metric
Wilson--Dirac operator on a fixed weighted Hilbert space of half-edge
data with a coframe realization and sign gaps, studied fixed
half-density metric contacts with a fixed contour for the sign, and
obtained an exact massive metric factorization reducing its
remaining condition to a massless metric card.

\emph{YM sixth series, V4.}  The YM programme restarted after five
archived series and, in its first versions, developed rooted-forest
formulae for a reserve estimate in a cluster expansion: exact
forest-probability ratios for rooted child trees, cancellation of the
reserve against the probability denominator, a tree-valued reserve
estimate, and the telescoping of component-mass updates.

\begin{theorem}[V25 companion-audit decision]
\label{thm:c3v25-companion-decision}
Nothing is imported from SM V27 or YM V4.  Neither bears on the
third cycle's chapters or on the gates.
\end{theorem}

\begin{proof}
The SM results concern lattice Dirac operators and metric contacts;
the YM results concern rooted forests in a cluster expansion.  None
is a Navier--Stokes statement, and none concerns Gaussian spectral
theory, the class bounds, or the ledgers of this series.
\end{proof}

% =============================================================================
\section{Position after twenty-five versions, and the block V26--V30}
\label{sec:c3v25-position}
% =============================================================================

\begin{assessment}[Position after twenty-five versions of the third cycle]
\label{ass:c3v25-position-twenty-five}
Four chapters are written: cluster energy and the band (V2--V10),
localization (V11--V14), the anchor (V16--V19), and the explicit
floor (V21--V24).  The fourth chapter replaced the one compactness
constant of the Rayleigh chapter by an explicit constant and found
it exponentially small, closing the push-budget route through the
floor alone for practical purposes and locating the possible
improvement in the equation.  It also corrected the block's own
direction twice: the high spectral case is empty by the Rayleigh
bound (V21), and concentration in a ball gives no floor (V22,
negative record N-C3-1).  The block V26--V30 follows the direction of
V24: V26 states the conditional mean floor of order one and tests
the constant-chasing condition against the class (expected
negative); V27--V28 examine the nearly constant core against the
momentum identity, through the Gaussian mean of the pressure
gradient of a nearly constant divergence-free field; V29
consolidates and fixes the direction for V31--V35; V30 audits.  No
global regularity claim is made.
\end{assessment}

% =============================================================================
\section{Integrated V25 conclusion and next acquisition order}
\label{sec:c3v25-conclusion}
% =============================================================================

\begin{theorem}[What V25 establishes]
\label{thm:c3v25-conclusion}
Relative to the two legacy cycles and V1--V24: the companion-audit
decision (Theorem~\ref{thm:c3v25-companion-decision}) and the position.
The full Navier--Stokes stability/global-regularity theorem remains
unproved.
\end{theorem}

\begin{proof}
The cited records.
\end{proof}

\begin{programme}[V26 acquisition order]
\label{prog:c3v26}
\begin{enumerate}[label=\textup{(AC\arabic*)},leftmargin=4.8em]
\item State and prove the conditional mean floor:
      \(\langle\mathrm{Var}\rangle_{\mathfrak m}\ge\frac1{16}(1-\sum_{k\le K_r}B_k/c_\Phi)\) with
      \(B_k\) the explicit bounds of Theorem C2-V41-spectrum, valid
      whenever the right side is positive.
\item Test the condition \(\sum_{k\le K_r}B_k<c_\Phi\) against the class:
      lower-bound \(B_0\) by the class constants and record whether
      any bound of the series makes the condition hold; record the
      answer, expected negative, as a firewall.
\item The next compulsory companion audit is V30.
\end{enumerate}
\end{programme}

\begin{assessment}[Position after V25]
\label{ass:c3v25-position}
The companion audit is recorded; the block V26--V30 begins with the
conditional mean floor.  No full stability or global-regularity
claim is made.
\end{assessment}

% =============================================================================
\section*{Version V25 (third cycle) append-only record}
\addcontentsline{toc}{section}{Version V25 (third cycle) append-only record}
% =============================================================================

The complete V24 source of the third cycle was retained before this
continuation.  Its SHA--256 digest was
\begin{center}\scriptsize\ttfamily
5afb5aa60a0566783ea19b5bcd008cde62ac35c59ab6da17f7d1b4c2d0d6e697.
\end{center}
V25 performed the compulsory companion audit and recorded the
position.  No full regularity claim is made.

% =============================================================================
\section{V26: the conditional mean floor, and why its condition never holds}
\label{sec:c3v26-entry}
% =============================================================================

Items (AC1) and (AC2) of Programme~\ref{prog:c3v26} are carried out.
The one-level criterion of V21, averaged against an invariant measure
and combined with the explicit mean spectrum of the second cycle,
gives a mean variance floor of order one whenever the explicit
bounds on the mean energies of the levels up to \(K_r\) sum to less
than the pinching constant.  The condition is then tested: its
\(k=0\) part alone requires the velocity constant to be below
\(1/(2\sqrt3)\) and the pinching constant to exceed an explicit multiple
of the pressure constant, and the second cycle's mean constraint for
small velocity constants bounds the pinching constant above by a
smaller multiple, for every positive velocity constant.  The
condition therefore never holds at a homogeneous type-I point, and
the conditional floor is recorded as a negative.  The only floor of
the series remains the pointwise explicit one, and the block turns
to the nearly constant core.  No global regularity claim is made.
The next compulsory companion audit is V30.

% =============================================================================
\section{The conditional mean floor}
\label{sec:c3v26-conditional}
% =============================================================================

Notation of C2-V36--V41: \(\kappa=C_\infty\nu^{-1/2}\), \(d=4\nu\|\nabla V\|_G^2\),
\(K_P=4\nu^{-1/2}(4\pi\nu/3)^{1/4}C_{\rm CZ}K_6^2\), and for an invariant
probability measure \(\mathfrak m\) on \(\mathcal F(x_0)\) the explicit
bounds of Theorem C2-V41-spectrum,
\begin{equation}
 B_k:=\frac1{\mu_k^2}\Bigl[\kappa\Bigl(\frac k2\Bigr)^{1/2}\langle\varphi\rangle^{1/2}+\kappa\langle d+3\varphi\rangle^{1/2}+\frac{K_P}4(3+2k)^{1/2}\Bigr]^2
 \ \ge\ \langle a_k\rangle_{\mathfrak m},
 \qquad
 B_0=4\Bigl[\kappa\langle d+3\varphi\rangle^{1/2}+\frac{\sqrt3}4K_P\Bigr]^2 .
 \label{eq:c3v26-Bk}
\end{equation}

\begin{theorem}[Conditional mean floor]
\label{thm:c3v26-conditional}
Let \(x_0\) be a homogeneous type-I point and \(\mathfrak m\) an invariant
probability measure on \(\mathcal F(x_0)\).  Then
\begin{equation}
 \langle\mathrm{Var}\rangle_{\mathfrak m}\ \ge\ \frac1{16}\Bigl(1-\frac1{c_\Phi}\sum_{k=0}^{K_r}B_k\Bigr),
 \label{eq:c3v26-conditional}
\end{equation}
a floor of order one whenever \(\sum_{k\le K_r}B_k<c_\Phi\).
\end{theorem}

\begin{proof}
Pointwise, \(\mathrm{Var}\ge\frac1{16}(1-\hat a_{k_1})\) with \(k_1\le K_r\)
(Propositions~\ref{prop:c3v21-criteria}(ii) and \ref{prop:c3v21-bounded}(i)),
and \(\hat a_{k_1}=a_{k_1}/\varphi\le\sum_{k\le K_r}a_k/c_\Phi\) by the pinching.
Average against \(\mathfrak m\) (the \(a_k\) are bounded continuous
observables, C2-V40) and use \(\langle a_k\rangle\le B_k\).
\end{proof}

% =============================================================================
\section{The condition never holds}
\label{sec:c3v26-never}
% =============================================================================

\begin{proposition}[Necessary conditions from the \(k=0\) term]
\label{prop:c3v26-necessary}
If \(B_0<c_\Phi\) then
\begin{equation}
 \kappa<\frac1{2\sqrt3}\qquad\text{and}\qquad
 c_\Phi^{1/2}\ >\ \frac{\sqrt3}{2}\,\frac{K_P}{1-2\sqrt3\,\kappa} .
 \label{eq:c3v26-necessary}
\end{equation}
\end{proposition}

\begin{proof}
\(\langle d+3\varphi\rangle\ge3\langle\varphi\rangle\ge3c_\Phi\), so
\(B_0^{1/2}=2[\kappa\langle d+3\varphi\rangle^{1/2}+\frac{\sqrt3}4K_P]\ge2\sqrt3\,\kappa\,c_\Phi^{1/2}+\frac{\sqrt3}2K_P\),
and \(B_0<c_\Phi\) gives \(c_\Phi^{1/2}(1-2\sqrt3\kappa)>\frac{\sqrt3}2K_P>0\).
\end{proof}

\begin{theorem}[The condition is incompatible with the class]
\label{thm:c3v26-incompatible}
At a homogeneous type-I point, \(B_0\ge c_\Phi\) for every value of the
velocity constant \(\kappa>0\); hence \(\sum_{k\le K_r}B_k\ge c_\Phi\) and the
right side of \eqref{eq:c3v26-conditional} is never positive.
\end{theorem}

\begin{proof}
Suppose \(B_0<c_\Phi\).  By Proposition~\ref{prop:c3v26-necessary},
\(\kappa<1/(2\sqrt3)<\frac12\), so the mean constraint of Theorem
C2-V36-constraint applies: with \(m_\kappa=\min\{1-\kappa,(2-4\kappa)/3\}=(2-4\kappa)/3\)
(the second is the smaller for every \(\kappa\ge-1\)),
\(c_\Phi\le\langle\varphi\rangle\le K_P^2/(3m_\kappa^2)\), i.e.\
\(c_\Phi^{1/2}\le K_P/(\sqrt3\,m_\kappa)\).  Combined with
\eqref{eq:c3v26-necessary},
\[
 \frac{\sqrt3}2\,\frac{K_P}{1-2\sqrt3\kappa}\ <\ \frac{K_P}{\sqrt3\,m_\kappa}
 \quad\Longleftrightarrow\quad
 \frac32m_\kappa\ <\ 1-2\sqrt3\kappa
 \quad\Longleftrightarrow\quad
 1-2\kappa\ <\ 1-2\sqrt3\kappa
 \quad\Longleftrightarrow\quad
 \kappa(2\sqrt3-2)<0 ,
\]
which is false for \(\kappa>0\).  A singular point has
\(\kappa\ge\varepsilon_{\rm reg}>0\) (Firewall C2-V36-AC2), so \(B_0\ge c_\Phi\).
\end{proof}

\begin{firewall}[Item (AC2): negative record N-C3-2, no mean floor from the mean spectrum]
\label{fw:c3v26-negative}
The route (a) of Assessment~\ref{ass:c3v24-direction} is closed, and
more sharply than expected: the condition \(\sum_{k\le K_r}B_k<c_\Phi\) is
not merely unverified, it is incompatible with the second cycle's
own mean constraint for every velocity constant.  The reason is
structural.  The mean spectrum bounds \(\langle a_k\rangle\) through the
flux bound, whose pressure term is \(K_P\) times a square root of the
mean energy; and the mean constraint bounds the mean energy itself
by \(K_P^2\) times a function of \(\kappa\).  Both come from the same
inequality, and the \(k=0\) bound is, up to the factor computed
above, the mean energy bound itself: the constant mode is allowed
by the mean spectrum to carry all the energy the class allows.  The
mean spectrum therefore cannot show that the energy is spread over
several levels on average; it shows only that no level carries more
than an explicit amount, an amount that is not below the pinching
constant.  What would be needed is a bound on the constant mode
alone that is strictly better than the flux bound, i.e.\ a
statement about the Gaussian momentum of a homogeneous blow-up that
does not pass through \(|n_0|\le a_0^{1/2}[\dots]\); the momentum identity
(C2-V17, V18) is the only other row of the ledger involving it, and
it is neutral.
\end{firewall}

\begin{assessment}[What remains of the variance floor]
\label{ass:c3v26-remains}
After V21--V26 the series has exactly one floor on the variance, the
pointwise explicit \(v_{\rm exp}\) of V23, exponentially small, and it has
shown that neither concentration (V22) nor the mean spectrum (V26)
improves it.  The remaining route is (b) of
Assessment~\ref{ass:c3v24-direction}: a nearly constant core has a
large Gaussian momentum, and the momentum identity decays it at rate
\(\frac12\) unless the Gaussian mean of the pressure gradient sustains
it.  V27 computes that mean exactly, as an integral of the
quadratic form \(\partial_iV_j\partial_jV_i\) against the gradient of the
Newtonian potential of the Gaussian, and bounds it on a nearly
constant core; V28 draws the consequence for the persistence of
nearly constant states.  No global regularity claim is made.
\end{assessment}

% =============================================================================
\section{Integrated V26 conclusion and next acquisition order}
\label{sec:c3v26-conclusion}
% =============================================================================

\begin{theorem}[What V26 proves]
\label{thm:c3v26-conclusion}
Relative to the two legacy cycles and V1--V25: the conditional mean
floor (Theorem~\ref{thm:c3v26-conditional}), the necessary conditions
(Proposition~\ref{prop:c3v26-necessary}), the incompatibility
(Theorem~\ref{thm:c3v26-incompatible}) and the negative record N-C3-2.
The full Navier--Stokes stability/global-regularity theorem remains
unproved.
\end{theorem}

\begin{proof}
The cited results.
\end{proof}

\begin{programme}[V27 acquisition order]
\label{prog:c3v27}
\begin{enumerate}[label=\textup{(AC\arabic*)},leftmargin=4.8em]
\item The exact identity: with \(q:=\partial_iV_j\partial_jV_i\) (so that
      \(-\Delta P_V=q\)) and \(g:=\Gamma_N*G\) the Newtonian potential of the
      Gaussian, \(\int\nabla P_V\,G=-\int q\,\nabla g\), with \(\nabla g\) radial,
      \(|\nabla g(y)|=M(|y|)/(4\pi|y|^2)\), \(M(s)=\int_{|z|<s}G\); hence the
      Gaussian mean of the pressure gradient is bounded by
      \(\frac1{4\pi}\int|q(y)|\,M(|y|)|y|^{-2}\dd y\).
\item On a nearly constant core (\(k_1=0\), \(\eta=16\,\mathrm{Var}\)): split at
      radius \(\rho\); inside, \(\int_{B_\rho}|q|\le\int_{B_\rho}|\nabla V|^2\le e^{\rho^2/(4\nu)}\varphi\,\|\nabla\widehat V\|_G^2\)
      with Lemma~\ref{lem:c3v24-h1}; outside, \(\int_{|y|>\rho}|q|\,|y|^{-2}\le\rho^{-2}\|\nabla V\|_{L^2}^2\le C\rho^{-2}\);
      state the resulting bound on \(|\overline{N(V)}|\) with the advective
      part from the dipole moments of \(V-\bar V\).
\item V28: the persistence alternative; V30 compulsory companion audit.
\end{enumerate}
\end{programme}

\begin{assessment}[Position after V26]
\label{ass:c3v26-position}
The conditional mean floor is void; the pointwise explicit floor is
the only one; the block turns to the nearly constant core.  No full
stability or global-regularity claim is made.
\end{assessment}

% =============================================================================
\section*{Version V26 (third cycle) append-only record}
\addcontentsline{toc}{section}{Version V26 (third cycle) append-only record}
% =============================================================================

The complete V25 source of the third cycle was retained before this
continuation.  Its SHA--256 digest was
\begin{center}\scriptsize\ttfamily
0afecaa7c07d3d57491ce39a5d82b63ff8c37783f498e020a18c81935b09665f.
\end{center}
V26 proved the conditional mean floor and showed that its condition
is incompatible with the class.  No full regularity claim is made.

% =============================================================================
\section{V27: the Gaussian mean of the pressure gradient, and the nonlinear mean on a nearly constant core}
\label{sec:c3v27-entry}
% =============================================================================

Items (AC1) and (AC2) of Programme~\ref{prog:c3v27} are carried out.
(AC1): the Gaussian mean of the pressure gradient of a divergence-free
field of the class is the integral of the quadratic form
\(\partial_iV_j\partial_jV_i\) against the gradient of the Newtonian potential
of the Gaussian, a radial field of size the Gaussian mass inside the
radius over four pi times the radius squared; it is bounded by the
enstrophy in any ball at the price of the ball's radius, and by the
class enstrophy outside the ball at the price of the inverse square
of the radius.  (AC2): on a nearly constant core the advective part
of the nonlinear mean is bounded by the square root of the variance
and the pressure part by the two-radius bound, the inside part being
exponentially expensive in the radius and linear in the variance.
The nonlinear mean is therefore small whenever the variance is
exponentially small in a radius beyond an explicit threshold, and
the momentum identity then decays the Gaussian momentum at rate at
least one quarter.  No global regularity claim is made.  The next
compulsory companion audit is V30.

% =============================================================================
\section{The Gaussian mean of the pressure gradient}
\label{sec:c3v27-pressure}
% =============================================================================

Throughout, \(V\) is the profile of an element of the family at a
homogeneous type-I point at some \(\sigma\): smooth, divergence-free,
\(|V|\le C_\infty\), \(\|V\|_{L^6(\R^3)}\le K_6\), \(\|\nabla V\|_{L^2(\R^3)}^2\le C\) (the
class bounds, C2-V4 and C2-V36), and \(P_V=\sum_{ij}R_iR_j(V_iV_j)\) the
pressure of C2-V36, with \(-\Delta P_V=\partial_i\partial_j(V_iV_j)=\partial_iV_j\,\partial_jV_i=:q\)
by \(\operatorname{div}V=0\).  Note \(|q|\le|\nabla V|^2\) (Cauchy--Schwarz for the
trace of \((\nabla V)^2\)), so \(q\in L^1(\R^3)\) with \(\|q\|_{L^1}\le C\), and
\(q\) is bounded and smooth.  Let \(\Gamma_N(y)=1/(4\pi|y|)\) be the
Newtonian kernel and
\begin{equation}
 g:=\Gamma_N*G,\qquad -\Delta g=G,\qquad
 \nabla g(y)=-\frac{M(|y|)}{4\pi|y|^2}\,\frac{y}{|y|},\qquad M(s):=\int_{|z|<s}G ,
 \label{eq:c3v27-g}
\end{equation}
so that \(|\nabla g(y)|=M(|y|)/(4\pi|y|^2)\le\min\bigl(|y|/3,\ |G|/(4\pi|y|^2)\bigr)\)
(Gauss's law for the radial potential; \(M(s)\le\frac{4\pi}3s^3\) since
\(G\le1\), and \(M\le|G|\)).

\begin{proposition}[Exact identity for the mean pressure gradient]
\label{prop:c3v27-identity}
\begin{equation}
 \boxed{\quad
 \int\nabla P_V\,G\,\dd z\ =\ -\int q\,\nabla g\,\dd y\ =\ \frac1{4\pi}\int\partial_iV_j\partial_jV_i(y)\,\frac{M(|y|)}{|y|^2}\,\frac{y}{|y|}\,\dd y ,
 \quad}
 \label{eq:c3v27-identity}
\end{equation}
and consequently, for every \(\rho>0\),
\begin{equation}
 \Bigl|\int\nabla P_V\,G\Bigr|\ \le\ \frac\rho3\int_{B_\rho}|\nabla V|^2\,\dd y+\frac{|G|}{4\pi\rho^2}\int_{|y|>\rho}|\nabla V|^2\,\dd y
 \ \le\ \frac\rho3\int_{B_\rho}|\nabla V|^2\,\dd y+\frac{|G|\,C}{4\pi\rho^2} .
 \label{eq:c3v27-two-radius}
\end{equation}
\end{proposition}

\begin{proof}
For \(q\in L^1\cap L^\infty\) the Newtonian potential \(\Gamma_N*q\) is
well defined, satisfies \(-\Delta(\Gamma_N*q)=q\), and lies in weak
\(L^3\); \(P_V\in L^3\) (C2-V36) satisfies the same equation, so the
difference is a harmonic function in \(L^3+L^{3,\infty}\), hence zero:
\(P_V=\Gamma_N*q\).  Integrating by parts (\(P_V\in L^3\), \(\nabla G\in L^{3/2}\)
with Gaussian decay), \(\int\nabla P_VG=-\int P_V\nabla G=-\int(\Gamma_N*q)\nabla G=-\int q\,(\Gamma_N*\nabla G)\)
by Fubini (\(q\in L^1\), \(\Gamma_N*\nabla G\) bounded), and
\(\Gamma_N*\nabla G=\nabla(\Gamma_N*G)=\nabla g\).  The formula for \(\nabla g\) is
Gauss's law for the radial solution of \(-\Delta g=G\): the flux of
\(\nabla g\) through the sphere of radius \(s\) is \(-M(s)\).  The bound
\eqref{eq:c3v27-two-radius} splits the integral at \(|y|=\rho\) and uses
\(|q|\le|\nabla V|^2\) with \(|\nabla g|\le|y|/3\le\rho/3\) inside and
\(|\nabla g|\le|G|/(4\pi\rho^2)\) outside.
\end{proof}

\begin{remark}[What the identity says]
\label{rem:c3v27-meaning}
The Gaussian mean of the pressure gradient is not a Gaussian
average of anything: it sees the strain of the profile everywhere,
weighted by an inverse-square kernel, and the enstrophy of the class
outside a ball of radius \(\rho\) contributes at most
\(|G|C/(4\pi\rho^2)\) however it is arranged.  Inside the ball the
kernel is small near the centre (\(|y|/3\)) and the strain is what the
core allows.  For the nearly constant core, the strain inside is
small in the Gaussian sense (Lemma~\ref{lem:c3v24-h1}), which converts
into an unweighted bound only at the exponential price
\(e^{\rho^2/(4\nu)}\); the two radii of \eqref{eq:c3v27-two-radius} are
in competition, and the competition is what the next section
quantifies.
\end{remark}

% =============================================================================
\section{The nonlinear mean on a nearly constant core}
\label{sec:c3v27-core}
% =============================================================================

Recall (C2-V17-mean-profile) \(\overline{N(V)}=\frac1{|G|}\int N(V)G\) with
\(N(V)=(V\cdot\nabla)V+\nabla P_V\) and \(\int(V\cdot\nabla)V\,G=\frac1{2\nu}\int(z\cdot V)V\,G\).
Put \(c_r:=(\frac18+2C_r^2)^{1/2}\) (Lemma~\ref{lem:c3v24-h1}).

\begin{proposition}[The nonlinear mean on a nearly constant core]
\label{prop:c3v27-core}
Suppose \(\mathrm{Var}<\frac1{32}\), \(\eta=16\,\mathrm{Var}\), and the dominant level
is \(k_1=0\).  Then, for every \(\rho>0\),
\begin{equation}
 \boxed{\quad
 |\overline{N(V)}|\ \le\ \frac{\varphi}{\nu^{1/2}|G|}\Bigl[\sqrt6\,\eta^{1/2}+(3+4c_r)^{1/2}\eta\Bigr]
 +\frac{\rho\,c_r\,\varphi}{3\nu|G|}\,e^{\rho^2/(4\nu)}\,\eta+\frac{C}{4\pi\rho^2} .
 \quad}
 \label{eq:c3v27-core}
\end{equation}
\end{proposition}

\begin{proof}
Write \(V=\bar V+\widetilde V\) with \(\bar V\) the Gaussian mean, a constant
vector; \(\widetilde V=V-\Pi_0V\) is divergence-free with
\(\|\widetilde V\|_G^2=\varphi(1-\hat a_0)\le\varphi\eta\) and \(|\bar V|\le(\varphi/|G|)^{1/2}\).

\emph{Advective part.}  \(\int(z\cdot V)V\,G=\int(z\cdot\widetilde V)\bar V\,G+\int(z\cdot\bar V)\widetilde V\,G+\int(z\cdot\widetilde V)\widetilde V\,G\),
the term \(\int(z\cdot\bar V)\bar VG\) vanishing by oddness.  The first
two are bounded by \(|\bar V|\,\|z\|_G\|\widetilde V\|_G\) each, with
\(\|z\|_G^2=\int|z|^2G=6\nu|G|\), so together by
\(2(\varphi/|G|)^{1/2}(6\nu|G|)^{1/2}(\varphi\eta)^{1/2}=2\sqrt6\,\nu^{1/2}\varphi\,\eta^{1/2}\).
The third is bounded by \(\||z|\widetilde V\|_G\|\widetilde V\|_G\), and the Gaussian
Hardy inequality (Lemma C2-V36-hardy, whose proof uses only
\(\operatorname{div}\widetilde V=0\)) gives
\(\||z|\widetilde V\|_G^2\le12\nu\|\widetilde V\|_G^2+16\nu^2\|\nabla\widetilde V\|_G^2\le12\nu\varphi\eta+16\nu\varphi c_r\eta=4\nu\varphi\eta(3+4c_r)\)
by Lemma~\ref{lem:c3v24-h1} (\(\nu\|\nabla\widetilde V\|_G^2=\nu\varphi\|\nabla(\widehat V-P)\|_G^2\le\varphi c_r\eta\)).
So the third term is at most \(2(3+4c_r)^{1/2}\nu^{1/2}\varphi\eta\), and
\(\frac1{2\nu|G|}|\int(z\cdot V)VG|\le\frac{\varphi}{\nu^{1/2}|G|}[\sqrt6\eta^{1/2}+(3+4c_r)^{1/2}\eta]\).

\emph{Pressure part.}  By \eqref{eq:c3v27-two-radius} divided by \(|G|\),
with \(\int_{B_\rho}|\nabla V|^2\le e^{\rho^2/(4\nu)}\|\nabla V\|_G^2=e^{\rho^2/(4\nu)}\|\nabla\widetilde V\|_G^2\le e^{\rho^2/(4\nu)}\varphi c_r\eta/\nu\)
(the gradient of the constant vanishes).
\end{proof}

\begin{corollary}[Decay of the Gaussian momentum on a nearly constant core]
\label{cor:c3v27-decay}
In the situation of Proposition~\ref{prop:c3v27-core}, let
\begin{equation}
 \rho_{\rm c}:=\Bigl(\frac{8\,C\,|G|^{1/2}}{\pi\,c_\Phi^{1/2}}\Bigr)^{1/2},\qquad
 \eta_{\rm c}:=\min\Bigl\{\frac{\nu|G|}{6144\,C_\Phi},\ \frac{1}{32(3+4c_r)^{1/2}}\Bigl(\frac{\nu|G|}{C_\Phi}\Bigr)^{1/2},\ \frac{3\nu}{32\,\rho_{\rm c}c_r}\Bigl(\frac{|G|}{C_\Phi}\Bigr)^{1/2}e^{-\rho_{\rm c}^2/(4\nu)}\Bigr\}.
 \label{eq:c3v27-etac}
\end{equation}
If \(\eta\le\eta_{\rm c}\) then \(|\overline{N(V)}|\le\frac18(\varphi/|G|)^{1/2}\le\frac14|\bar V|\), and hence
\begin{equation}
 \boxed{\quad
 \frac{\dd}{\dd\sigma}|\bar V|\ \le\ -\frac14|\bar V| .
 \quad}
 \label{eq:c3v27-decay}
\end{equation}
\end{corollary}

\begin{proof}
Since \(a_0=|G||\bar V|^2\ge(1-\eta)\varphi\ge\frac12\varphi\) for \(\eta\le\frac12\),
\(|\bar V|\ge(\varphi/(2|G|))^{1/2}\), so \(\frac18(\varphi/|G|)^{1/2}\le\frac{\sqrt2}8|\bar V|\le\frac14|\bar V|\).
It remains to bound each of the four terms of \eqref{eq:c3v27-core}
at \(\rho=\rho_{\rm c}\) by \(\frac1{32}(\varphi/|G|)^{1/2}\), using \(c_\Phi\le\varphi\le C_\Phi\).
First term: \(\sqrt6\varphi\eta^{1/2}/(\nu^{1/2}|G|)\le\frac1{32}(\varphi/|G|)^{1/2}\) iff
\(\eta\le\nu|G|/(6144\varphi)\), implied by the first entry of
\eqref{eq:c3v27-etac}.  Second term:
\((3+4c_r)^{1/2}\varphi\eta/(\nu^{1/2}|G|)\le\frac1{32}(\varphi/|G|)^{1/2}\) iff
\(\eta\le(\nu|G|/\varphi)^{1/2}/(32(3+4c_r)^{1/2})\), implied by the second
entry.  Third term:
\(\rho_{\rm c}c_r\varphi e^{\rho_{\rm c}^2/(4\nu)}\eta/(3\nu|G|)\le\frac1{32}(\varphi/|G|)^{1/2}\) iff
\(\eta\le3\nu(|G|/\varphi)^{1/2}e^{-\rho_{\rm c}^2/(4\nu)}/(32\rho_{\rm c}c_r)\), implied by
the third entry.  Fourth term: \(C/(4\pi\rho_{\rm c}^2)=\frac1{32}(c_\Phi/|G|)^{1/2}\le\frac1{32}(\varphi/|G|)^{1/2}\)
by the choice of \(\rho_{\rm c}\) and the pinching.  Summing,
\(|\overline{N(V)}|\le\frac18(\varphi/|G|)^{1/2}\).  Finally, where \(\bar V\ne0\),
\(\frac{\dd}{\dd\sigma}|\bar V|=\bar V\cdot\frac{\dd}{\dd\sigma}\bar V/|\bar V|\le-\frac12|\bar V|+|\overline{N(V)}|\le-\frac14|\bar V|\)
by the momentum identity \(\frac{\dd}{\dd\sigma}\bar V=-\frac12\bar V-\overline{N(V)}\)
(C2-V17-mean-profile).
\end{proof}

\begin{remark}[The threshold]
\label{rem:c3v27-threshold}
The threshold \(\eta_{\rm c}\) is exponentially small in
\(\rho_{\rm c}^2/(4\nu)=2C|G|^{1/2}/(\pi\nu c_\Phi^{1/2})\), through the inside term
of the pressure bound; the other entries are polynomial.  Its
exponent is polynomial in the class constants over the pinching
constant, while the exponent of \(v_{\rm exp}\) (Remark~\ref{rem:c3v23-size})
is quadratic in the ratio of the \(L^6\) bound to the pinching.  Which
is larger depends on the constants; both are far below any
order-one quantity.  The corollary is nonetheless of a different
kind from V23: it uses the equation (the momentum identity), and it
converts smallness of the variance into decay of a conserved-looking
quantity, which V28 turns into a statement about windows.
\end{remark}

% =============================================================================
\section{Integrated V27 conclusion and next acquisition order}
\label{sec:c3v27-conclusion}
% =============================================================================

\begin{theorem}[What V27 proves]
\label{thm:c3v27-conclusion}
Relative to the two legacy cycles and V1--V26: the exact identity for
the mean pressure gradient (Proposition~\ref{prop:c3v27-identity}), the
bound on the nonlinear mean on a nearly constant core
(Proposition~\ref{prop:c3v27-core}), and the decay corollary
(Corollary~\ref{cor:c3v27-decay}).  The full Navier--Stokes
stability/global-regularity theorem remains unproved.
\end{theorem}

\begin{proof}
The cited results.
\end{proof}

\begin{programme}[V28 acquisition order]
\label{prog:c3v28}
\begin{enumerate}[label=\textup{(AC\arabic*)},leftmargin=4.8em]
\item The locked level: while \(\mathrm{Var}<\frac1{32}\) on an interval of
      \(\sigma\), the dominant level \(k_1\) is constant (a change of level
      forces \(r\) through a midpoint \(\mu_k+\frac14\), where
      \(\mathrm{Var}\ge\frac1{16}\)).
\item The persistence alternative: on any interval of length
      \(\tau_*:=2\log(2C_\Phi/c_\Phi)\) on which \(k_1=0\), either
      \(\sup\mathrm{Var}\ge\eta_{\rm c}/16\), or the decay
      \eqref{eq:c3v27-decay} drives \(\hat a_0\) below \(\frac12\), contradicting
      \(\hat a_0\ge1-\eta\); hence \(\sup\mathrm{Var}\ge\min(\eta_{\rm c}/16,\frac1{32})\)
      on every such interval.  Combine with the pointwise floor for
      the levels \(1\le k_1\le K_r\) into a window floor.
\item V29: consolidation and direction for V31--V35; V30 compulsory
      companion audit.
\end{enumerate}
\end{programme}

\begin{assessment}[Position after V27]
\label{ass:c3v27-position}
The mean pressure gradient is an explicit strain integral; on a
nearly constant core the nonlinear mean is small below an explicit
threshold on the variance, and the momentum then decays.  No full
stability or global-regularity claim is made.
\end{assessment}

% =============================================================================
\section*{Version V27 (third cycle) append-only record}
\addcontentsline{toc}{section}{Version V27 (third cycle) append-only record}
% =============================================================================

The complete V26 source of the third cycle was retained before this
continuation.  Its SHA--256 digest was
\begin{center}\scriptsize\ttfamily
46609cccfcebf4f79cebdbf5076490038cd415447a0a247e545bb024eb3450ba.
\end{center}
V27 proved the identity for the mean pressure gradient and the
bound on the nonlinear mean on a nearly constant core.  No full
regularity claim is made.

% =============================================================================
\section{V28: the locked level, the persistence alternative, and the window floor}
\label{sec:c3v28-entry}
% =============================================================================

Items (AC1) and (AC2) of Programme~\ref{prog:c3v28} are carried out.
(AC1): while the variance stays below one thirty-second, the dominant
level cannot change, because a change forces the Rayleigh quotient
through a midpoint between two eigenvalues, where the variance is at
least one sixteenth.  (AC2): on a window of explicit log-length, a
nearly constant state cannot persist below the threshold of V27,
because its Gaussian momentum would decay by a factor that the
pinching forbids; combined with the pointwise floor for the other
levels this gives a window floor on the variance, and, through a
Lipschitz bound on the variance in log-time, a floor on its mean
against invariant measures.  The window floor's exponent is linear
in the class enstrophy over the square root of the pinching
constant, against the quadratic exponent of V23; which is larger
depends on the class constants, and both are far below order one.
No global regularity claim is made.  The next compulsory companion
audit is V30.

% =============================================================================
\section{The locked level}
\label{sec:c3v28-locked}
% =============================================================================

\begin{lemma}[The dominant level is locked while the variance is small]
\label{lem:c3v28-locked}
Let \(I\) be an interval of \(\sigma\) on which \(\mathrm{Var}<\frac1{16}\).  Then
the dominant level \(k_1\) of Definition~\ref{def:c3v21-dominant} is
constant on \(I\).
\end{lemma}

\begin{proof}
\(r=2\mathcal L/\varphi\) is continuous in \(\sigma\) (both are \(C^1\) on the
family, C2-V12).  If \(k_1\) took two different values on \(I\), the
intermediate value theorem would give a \(\sigma\in I\) with
\(r(\sigma)=\mu_k+\frac14\) for some \(k\) (the level nearest \(r\) changes only
when \(r\) crosses a midpoint between consecutive eigenvalues).  At
such \(\sigma\) every eigenvalue is at distance at least \(\frac14\) from \(r\),
so \(\mathrm{Var}=\sum_k(\mu_k-r)^2\hat a_k\ge\frac1{16}\), contradicting the
hypothesis.
\end{proof}

% =============================================================================
\section{The persistence alternative and the window floor}
\label{sec:c3v28-window}
% =============================================================================

With \(\eta_{\rm c}\) of \eqref{eq:c3v27-etac}, \(v^{(k)}\) of
\eqref{eq:c3v23-constants}, and \(K_r\) of Proposition~\ref{prop:c3v21-bounded},
define
\begin{equation}
 \tau_*:=2\log\frac{4C_\Phi}{c_\Phi},\qquad
 v_{\rm win}:=\min\Bigl(\frac{\eta_{\rm c}}{16},\ \frac1{32},\ \min_{1\le k\le K_r}v^{(k)}\Bigr)
 \label{eq:c3v28-vwin}
\end{equation}
(the last minimum is void when \(K_r=0\)).

\begin{theorem}[Window floor on the variance]
\label{thm:c3v28-window}
Let \(x_0\) be a homogeneous type-I point.  For every \(W\in\mathcal F(x_0)\)
and every \(\sigma_0\),
\begin{equation}
 \boxed{\quad
 \sup_{\sigma\in[\sigma_0,\sigma_0+\tau_*]}\mathrm{Var}(\sigma)\ \ge\ v_{\rm win} .
 \quad}
 \label{eq:c3v28-window}
\end{equation}
\end{theorem}

\begin{proof}
Suppose \(\mathrm{Var}<v_{\rm win}\) on \(I=[\sigma_0,\sigma_0+\tau_*]\).  Then
\(\mathrm{Var}<\frac1{32}\) on \(I\), so \(k_1\) is constant on \(I\)
(Lemma~\ref{lem:c3v28-locked}) and \(k_1\le K_r\).  If \(k_1\ge1\), Theorem
\ref{thm:c3v23-floor}'s proof gives \(\mathrm{Var}\ge v^{(k_1)}\ge v_{\rm win}\) at
every \(\sigma\in I\), a contradiction.  If \(k_1=0\), then \(\eta=16\,\mathrm{Var}<\eta_{\rm c}\)
on \(I\), and Corollary~\ref{cor:c3v27-decay} gives
\(\frac{\dd}{\dd\sigma}|\bar V|\le-\frac14|\bar V|\) on \(I\) (with \(\bar V\ne0\), as
\(a_0\ge\frac12\varphi>0\)); hence \(|\bar V(\sigma_0+\tau_*)|\le|\bar V(\sigma_0)|e^{-\tau_*/4}\)
and \(a_0(\sigma_0+\tau_*)\le a_0(\sigma_0)e^{-\tau_*/2}\le C_\Phi\cdot\frac{c_\Phi}{4C_\Phi}=\frac14c_\Phi\).
But \(a_0(\sigma_0+\tau_*)\ge(1-\eta)\varphi\ge\frac12c_\Phi\), a contradiction.
\end{proof}

\begin{corollary}[Mean floor from the window floor]
\label{cor:c3v28-mean}
Let \(\Lambda_{\rm Var}\) be a bound on \(|\frac{\dd}{\dd\sigma}\mathrm{Var}|\) on the
family (finite: \(\mathrm{Var}=\|L_-\widehat V\|_G^2-r^2\) is \(C^1\) in \(\sigma\) with
derivative bounded by the class derivative bounds up to order four,
(E5) and C2-V4).  Then for every invariant probability measure
\(\mathfrak m\) on \(\mathcal F(x_0)\),
\begin{equation}
 \langle\mathrm{Var}\rangle_{\mathfrak m}\ \ge\ \max\Bigl(v_{\rm exp},\ \frac{v_{\rm win}^2}{4\Lambda_{\rm Var}\tau_*}\Bigr)
 \qquad\text{if }v_{\rm win}\le2\Lambda_{\rm Var}\tau_* .
 \label{eq:c3v28-mean}
\end{equation}
\end{corollary}

\begin{proof}
The first bound is Theorem~\ref{thm:c3v23-floor}.  For the second, by
the invariance of \(\mathfrak m\), \(\langle\mathrm{Var}\rangle=\langle\frac1{\tau_*}\int_0^{\tau_*}\mathrm{Var}(\Theta_\sigma W)\dd\sigma\rangle\),
and on each window of length \(\tau_*\) there is a \(\sigma_1\) with
\(\mathrm{Var}(\sigma_1)\ge v_{\rm win}\); by the Lipschitz bound,
\(\mathrm{Var}\ge\frac12v_{\rm win}\) on an interval of length
\(v_{\rm win}/(2\Lambda_{\rm Var})\le\tau_*\) around \(\sigma_1\) (intersected with the
window, at least half of it lies inside), so the window average is
at least \(\frac12v_{\rm win}\cdot\frac{v_{\rm win}}{4\Lambda_{\rm Var}}/\tau_*\).
\end{proof}

\begin{firewall}[Status of \(\Lambda_{\rm Var}\)]
\label{fw:c3v28-lipschitz}
\(\Lambda_{\rm Var}\) is explicit in principle: \(\frac{\dd}{\dd\sigma}\mathrm{Var}=2\langle L_-\widehat V,L_-\partial_\sigma\widehat V\rangle_G-2r\frac{\dd r}{\dd\sigma}\)
with \(\partial_\sigma\widehat V=\varphi^{-1/2}(Y-\frac12(\varphi'/\varphi)V)\), and
\(L_-Y=L_-(L_-V-N(V))\) involves derivatives of \(V\) up to order four,
bounded on the family by the scale-invariant derivative bounds of
(E5) as transported in C2-V4; the series has not written these
bounds beyond order two, so \(\Lambda_{\rm Var}\) is not given a closed
form here.  Its role is only to convert the window floor into a mean
floor, and it enters polynomially.
\end{firewall}

% =============================================================================
\section{Comparison of the two floors}
\label{sec:c3v28-comparison}
% =============================================================================

\begin{assessment}[Two exponentially small floors of different origin]
\label{ass:c3v28-comparison}
The series now has two explicit floors on the variance at homogeneous
points, and neither is of order one.
\begin{enumerate}[label=\textup{(\roman*)},leftmargin=3.2em]
\item The pointwise floor \(v_{\rm exp}\) (V23) has exponent
      \((\rho_0^*)^2/(4\nu)=16(4\pi/3)^{-2/3}|G|^2K_6^4/(c_\Phi^2\nu)\) in the
      closed-form case: quadratic in the ratio of the \(L^6\) bound to
      the pinching.  It uses no property of the equation.
\item The window floor \(v_{\rm win}\) (V28) has, in its constant-level
      entry \(\eta_{\rm c}/16\), exponent
      \(\rho_{\rm c}^2/(4\nu)=2C|G|^{1/2}/(\pi\nu c_\Phi^{1/2})\): linear in the
      ratio of the class enstrophy to the square root of the
      pinching.  It uses the momentum identity and the pinching over
      a window of log-length \(\tau_*\).
\end{enumerate}
The ratio of the exponents is \(8\pi(4\pi/3)^{-2/3}|G|^{3/2}K_6^4/(Cc_\Phi^{3/2})\),
which with \(K_6=C_6C^{1/2}\), \(C=C_I^2\nu^{3/2}\) and the upper bound
\(c_\Phi\le K_6^2\cdot8\pi\nu/3\) is at least a numerical multiple of
\(C_6/C_I\): for a type-I constant \(C_I\) large compared with the
Sobolev constant, the window floor may be the larger; for \(C_I\) of
order one the pointwise floor may be.  In neither case is the
threshold \(4v/C_r\) of the push budget approached by any bound of the
series.  The conclusion of the block is therefore that the variance
floor is explicit, by two independent mechanisms, and small; the
mechanism that would make it of order one is not the \(L^6\)
exclusion of polynomials nor the momentum decay, but a reason why
a profile of the class cannot have a nearly constant Gaussian core
of radius of order \((4\nu)^{1/2}\), and no such reason is in the series.
No global regularity claim is made.
\end{assessment}

% =============================================================================
\section{Integrated V28 conclusion and next acquisition order}
\label{sec:c3v28-conclusion}
% =============================================================================

\begin{theorem}[What V28 proves]
\label{thm:c3v28-conclusion}
Relative to the two legacy cycles and V1--V27: the locked level
(Lemma~\ref{lem:c3v28-locked}), the window floor
(Theorem~\ref{thm:c3v28-window}) and the mean floor
(Corollary~\ref{cor:c3v28-mean}).  The full Navier--Stokes
stability/global-regularity theorem remains unproved.
\end{theorem}

\begin{proof}
The cited results.
\end{proof}

\begin{programme}[V29 acquisition order]
\label{prog:c3v29}
\begin{enumerate}[label=\textup{(AC\arabic*)},leftmargin=4.8em]
\item Consolidate V26--V28 with V21--V24 as the variance chapter of
      the third cycle, with the two negative records N-C3-1, N-C3-2
      and the two floors.
\item Fix the direction for V31--V35 by going upstream: the other
      compactness constant of the gate at homogeneous points is the
      defect rate \(\underline\delta\) (C2-V6, C2-V9-B), the long-window
      minimum of \(\int\|\partial_\sigma V\|_G^2\) over the family, positive
      because no element is self-similar (E9).  Determine whether it
      can be made explicit by the method of V21--V23: a quantitative
      form of the exclusion of self-similar profiles (the NRS--Tsai
      import (E9)), i.e.\ a lower bound on \(\|L_-V-N(V)\|_G\) for every
      field of the class in terms of the class constants.
\item V30 compulsory companion audit and delivery.
\end{enumerate}
\end{programme}

\begin{assessment}[Position after V28]
\label{ass:c3v28-position}
The window floor is proved; the variance chapter has two explicit
small floors and no route to an order-one one.  No full stability or
global-regularity claim is made.
\end{assessment}

% =============================================================================
\section*{Version V28 (third cycle) append-only record}
\addcontentsline{toc}{section}{Version V28 (third cycle) append-only record}
% =============================================================================

The complete V27 source of the third cycle was retained before this
continuation.  Its SHA--256 digest was
\begin{center}\scriptsize\ttfamily
21bd2408341a6a392c39b584eead3a1c3f629076c63732867b6bf36c23e60b6c.
\end{center}
V28 proved the locked level, the window floor and the mean floor.
No full regularity claim is made.

% =============================================================================
\section{V29: the variance chapter consolidated, and the direction for V31--V35}
\label{sec:c3v29-entry}
% =============================================================================

Items (AC1) and (AC2) of Programme~\ref{prog:c3v29} are carried out.
The variance chapter of the third cycle, V21--V28, is consolidated:
its two explicit floors, its two negative records, and its
conclusion that the variance route to the push budget is exhausted
at the level of explicit constants.  The direction for V31--V35 goes
upstream to the other compactness constant of the homogeneous gate,
the defect rate, and asks for its explicit form through a
quantitative version of the exclusion of self-similar profiles.  No
global regularity claim is made.  The next compulsory companion
audit is V30.

% =============================================================================
\section{The variance chapter}
\label{sec:c3v29-chapter}
% =============================================================================

\begin{theorem}[The variance chapter, V21--V28]
\label{thm:c3v29-chapter}
At a homogeneous type-I point, for every profile of the family:
\begin{enumerate}[label=\textup{(\arabic*)},leftmargin=3.2em]
\item (Spectral criteria, V21.)  \(\mathrm{Var}\ge\frac1{16}(1-\hat a_{k_1})\) with
      \(k_1\le K_r=\lfloor2C_r-\frac12\rfloor\) the dominant level, pointwise;
      a profile of variance below \(\frac1{32}\) is within \(\eta=16\,\mathrm{Var}\)
      of a polynomial field of exact level \(k_1\), in \(L^2(G)\) and,
      with the deviation \(c_r\eta/\nu\), in \(H^1(G)\) (V24); its
      Rayleigh quotient is within \(\eta/2\) of \(\mu_{k_1}\) (V24); its
      dominant level is locked while the variance stays below
      \(\frac1{16}\) (V28).
\item (Negative N-C3-1, V22.)  Concentration in a ball gives no
      floor: the divergence-free cut-off constant has variance at
      most \(96A^2\rho(4\pi\nu)^{-1/2}e^{-\rho^2/(4\nu)}\).
\item (Pointwise floor, V23.)  \(\mathrm{Var}\ge v_{\rm exp}\), explicit in \(\nu\),
      \(K_6\), \(c_\Phi\), \(C_r\) and, for \(2\le k\le K_r\), the finite-dimensional
      constants \(\gamma_k^{\rm top},\Lambda_k\); in closed form when
      \(C_r<\frac54\); exponent quadratic in \(K_6^2|G|/c_\Phi\); uses no
      property of the equation.
\item (Push budget, V24.)  The mean geometric push is positive when
      the mean outward normalized flux is below \(4v_{\rm exp}/C_r\), a
      threshold below \(8e^{-170}\); the route through the floor alone
      is closed for practical purposes.
\item (Negative N-C3-2, V26.)  The conditional mean floor
      \(\langle\mathrm{Var}\rangle\ge\frac1{16}(1-\sum_{k\le K_r}B_k/c_\Phi)\) is void:
      \(B_0\ge c_\Phi\) for every velocity constant, by the second
      cycle's mean constraint.
\item (Mean pressure gradient, V27.)  \(\int\nabla P_VG=-\int q\nabla g\) with
      \(q=\partial_iV_j\partial_jV_i\) and \(g\) the Newtonian potential of \(G\);
      on a nearly constant core the nonlinear mean is at most
      \(\frac18(\varphi/|G|)^{1/2}\) once \(\eta\le\eta_{\rm c}\), and the Gaussian
      momentum then decays at rate \(\frac14\).
\item (Window floor, V28.)  On every window of log-length
      \(\tau_*=2\log(4C_\Phi/c_\Phi)\), \(\sup\mathrm{Var}\ge v_{\rm win}=\min(\eta_{\rm c}/16,\frac1{32},\min_{1\le k\le K_r}v^{(k)})\),
      with exponent linear in \(C|G|^{1/2}/(\nu c_\Phi^{1/2})\); a mean floor
      follows through the Lipschitz constant of the variance.
\end{enumerate}
\end{theorem}

\begin{proof}
Theorems~\ref{thm:c3v24-consolidated}, \ref{thm:c3v26-conditional},
\ref{thm:c3v26-incompatible}, \ref{thm:c3v28-window};
Propositions~\ref{prop:c3v27-identity}, \ref{prop:c3v27-core};
Corollaries~\ref{cor:c3v27-decay}, \ref{cor:c3v28-mean};
Lemmas~\ref{lem:c3v24-h1}, \ref{lem:c3v28-locked}.
\end{proof}

\begin{assessment}[What the chapter settles]
\label{ass:c3v29-settles}
The chapter began from the second cycle's one signed quantity with a
compactness constant, and ends with that constant explicit by two
mechanisms and a proof that the mechanisms available to the series
cannot make it of order one: the \(L^6\) exclusion of polynomials
charges only the Gaussian tail beyond a capped radius; the momentum
decay charges only the time a nearly constant core can survive; the
mean spectrum is no better than the mean energy bound.  An order-one
floor would need a reason why a profile of the class cannot carry a
nearly constant Gaussian core of radius of order \((4\nu)^{1/2}\) for a
log-time of order one, and the series has no such reason.  The
chapter is closed; it is not a dead end in the sense of the negative
records, since its floors are usable wherever the variance appears
with a sign, but the push budget is not decided by it.
\end{assessment}

% =============================================================================
\section{Direction for V31--V35: the explicit defect rate}
\label{sec:c3v29-direction}
% =============================================================================

\begin{assessment}[Going upstream: the defect rate]
\label{ass:c3v29-direction}
The gate at homogeneous points (G1 localized, V13; two-case form,
V17) has one constant from compactness, the defect rate
\(\underline\delta=\lim m(\ell)/\ell\), the long-window minimum of
\(\int\|\partial_\sigma V\|_G^2\) over the family (C2-V6-m), positive because
no element of the family is self-similar and self-similar profiles
in \(L^q\) vanish (E9; C2-V9-B).  This is exactly the situation of the
variance before V21: a positive quantity, positive because an exact
degenerate case (an eigenfunction; a self-similar profile) is
excluded by the class, made uniform by compactness.  The method of
V21--V23 replaced the exclusion by a quantitative exclusion.  The
block V31--V35 asks whether the same can be done for the defect: a
lower bound on \(\|L_-V-N(V)\|_G\), for every divergence-free field of
the class with the pinching, explicit in the class constants.  The
exclusion (E9) rests on the Ne\v cas--R\accent23u\v zi\v cka--\v Sver\'ak
identity: for a self-similar profile the function
\(\Pi_{\rm NRS}:=P_V+\frac12|V|^2+\frac12z\cdot V\) satisfies an elliptic
equation with drift, \(\nu\Delta\Pi_{\rm NRS}-(V+\frac12z)\cdot\nabla\Pi_{\rm NRS}=\nu|\Omega|^2\),
whose maximum principle and decay force \(\Pi_{\rm NRS}\le0\), and a
further integral identity forces \(\Omega\equiv0\), hence \(V\equiv0\).  With a
defect \(Y=L_-V-N(V)\ne0\) the equation acquires a source linear in \(Y\)
and \(\nabla Y\); the plan is: V31, the exact \(\Pi_{\rm NRS}\)-equation with
defect source, for every field of the class; V32, the quantitative
maximum principle: the positive part of \(\Pi_{\rm NRS}\) is bounded by a
weighted norm of the source, with explicit constants from the
decay of the class; V33, the quantitative integral identity: the
Gaussian enstrophy \(\int|\Omega|^2G\) is bounded by the source and by the
positive part of \(\Pi_{\rm NRS}\); V34, from small Gaussian enstrophy to
small Gaussian energy against the pinching (Gaussian Poincar\'e for
divergence-free fields, C2-V8-spectrum(iii)), giving the explicit
defect floor \(\underline\delta_{\rm exp}\) if every step closes, and its
consequence for the gate G1 with explicit constants; V35 audits.  It
is expected that the floor, if obtained, is again exponentially
small in the class constants, and the value of the block is to
locate, as the variance chapter did, the mechanism that would be
needed for an order-one constant.  No global regularity claim is
made.
\end{assessment}

% =============================================================================
\section{Integrated V29 conclusion and next acquisition order}
\label{sec:c3v29-conclusion}
% =============================================================================

\begin{theorem}[What V29 establishes]
\label{thm:c3v29-conclusion}
Relative to the two legacy cycles and V1--V28: the consolidated
variance chapter (Theorem~\ref{thm:c3v29-chapter}) and the direction.
The full Navier--Stokes stability/global-regularity theorem remains
unproved.
\end{theorem}

\begin{proof}
The cited results.
\end{proof}

\begin{programme}[V30 acquisition order]
\label{prog:c3v30}
\begin{enumerate}[label=\textup{(AC\arabic*)},leftmargin=4.8em]
\item The compulsory companion audit: read the current SM and YM
      notes and record their digests and decision.
\item Record the position after thirty versions; delivery.
\end{enumerate}
\end{programme}

\begin{assessment}[Position after V29]
\label{ass:c3v29-position}
The variance chapter is closed with two explicit floors and two
negative records; the direction is the explicit defect rate.  No
full stability or global-regularity claim is made.
\end{assessment}

% =============================================================================
\section*{Version V29 (third cycle) append-only record}
\addcontentsline{toc}{section}{Version V29 (third cycle) append-only record}
% =============================================================================

The complete V28 source of the third cycle was retained before this
continuation.  Its SHA--256 digest was
\begin{center}\scriptsize\ttfamily
ba05d9d0ad6b7c0f22c8243e448c7d252ca24916c82651255e58e2b589044c32.
\end{center}
V29 consolidated the variance chapter and fixed the direction.  No
full regularity claim is made.

% =============================================================================
\section{V30: compulsory companion audit and the position after thirty versions}
\label{sec:c3v30-entry}
% =============================================================================

This is the thirtieth version of the third cycle.  The compulsory
review of the two companion programmes is performed first, and the
position after thirty versions is recorded with the plan for the
block V31--V35.  No global regularity claim is made.  The next
compulsory companion audit is V35.

% =============================================================================
\section{Compulsory V30 review of the current companion notes}
\label{sec:c3v30-companion-audit}
% =============================================================================

The current companion files in \texttt{latex\_notes} were inspected
read-only.  Both programmes have moved since the V25 audit: the SM
programme from V27 to V29 of its seventh series, the YM programme
from V4 to V6 of its sixth series.  The frozen checkpoints are
\begin{align}
 &\texttt{Renewal\_SM\_Einstein\_Continuum\_Research\_Notes\_V29.tex},
 \notag\\[-2pt]
 &\qquad
 \texttt{67eb65610f9ad52bb799df830646a3ad9f76a4d2800bbad8eb9264c9e85797d1},
 \label{eq:c3v30-SM-checkpoint}\\
 &\texttt{Renewal\_Yang\_Mills\_Mass\_Gap\_Research\_Notes\_V6.tex},
 \notag\\[-2pt]
 &\qquad
 \texttt{6f6cbd7532fa75c668c3b9d1e2a96a47018542b1ea59a7dd09d11415f609e88a}.
 \label{eq:c3v30-YM-checkpoint}
\end{align}
(SM V29 has 12096 lines; YM V6 has 2464 lines.  As at V25, the
folder also holds a YM file of the preceding number, \texttt{V5},
byte-identical to \texttt{V6}; the audit records the newer name.)

\emph{SM V28--V29.}  The SM programme treated the low physical branch
of its massless metric condition (the low symbol of a midpoint
Wilson kernel, a vertex passed through the overlap sign, the low
heat response) and then a kinematic separation of the high branch
(a unique flat overlap zero, a cut-off buffer removing a low--high
bridge, a complete high-response bound), closing its condition MMAC
with a mixed smooth owner.

\emph{YM V5--V6.}  The YM programme completed the order statistics of
its BKAR tree, closed a singly rooted hierarchy allocation with a
correction to its own V3 exhaustion statement, performed its
companion audit, and developed a one-layer projection at fixed final
vertex set with a conditional Cayley law on forest components, a
measure-preserving orientation factorization, and a one-node
vertex-accretion kernel.

\begin{theorem}[V30 companion-audit decision]
\label{thm:c3v30-companion-decision}
Nothing is imported from SM V29 or YM V6.  Neither bears on the
third cycle's chapters or on the gates.
\end{theorem}

\begin{proof}
The SM results concern symbols of lattice Dirac kernels and heat
responses; the YM results concern tree statistics and forest
projections in a cluster expansion.  None is a Navier--Stokes
statement, and none concerns the Gaussian spectral theory, the
class bounds, the momentum identity or the defect of this series.
\end{proof}

% =============================================================================
\section{Position after thirty versions, and the block V31--V35}
\label{sec:c3v30-position}
% =============================================================================

\begin{assessment}[Position after thirty versions of the third cycle]
\label{ass:c3v30-position-thirty}
Five chapters are written: cluster energy and the band (V2--V10),
localization (V11--V14), the anchor (V16--V19), the explicit floor
(V21--V24), and its continuation to the variance chapter (V26--V29).
The third cycle has so far produced two kinds of result: structural
(the type-I gate at homogeneous points localized to a bounded
neighbourhood of an anchor group, with a two-case form), and
quantitative (the variance floor explicit by two mechanisms, with
negative records N-C3-1 and N-C3-2 marking the mechanisms that do
not improve it).  The block V31--V35 applies the quantitative method
to the second compactness constant of the gate, the defect rate,
through the Ne\v cas--R\accent23u\v zi\v cka--\v Sver\'ak identity with a
defect source, as planned in Assessment~\ref{ass:c3v29-direction}:
V31 the exact equation, V32 the quantitative maximum principle, V33
the quantitative integral identity, V34 the floor and its
consequence for G1, V35 the audit.  No global regularity claim is
made.
\end{assessment}

% =============================================================================
\section{Integrated V30 conclusion and next acquisition order}
\label{sec:c3v30-conclusion}
% =============================================================================

\begin{theorem}[What V30 establishes]
\label{thm:c3v30-conclusion}
Relative to the two legacy cycles and V1--V29: the companion-audit
decision (Theorem~\ref{thm:c3v30-companion-decision}) and the position.
The full Navier--Stokes stability/global-regularity theorem remains
unproved.
\end{theorem}

\begin{proof}
The cited records.
\end{proof}

\begin{programme}[V31 acquisition order]
\label{prog:c3v31}
\begin{enumerate}[label=\textup{(AC\arabic*)},leftmargin=4.8em]
\item Derive, for every smooth divergence-free field \(V\) of the
      class with pressure \(P_V\) and defect \(Y:=L_-V-N(V)\), the exact
      equation satisfied by \(\Pi_{\rm NRS}=P_V+\frac12|V|^2+\frac12z\cdot V\):
      the drift-elliptic operator \(\nu\Delta-(V+\frac12z)\cdot\nabla\) applied
      to \(\Pi_{\rm NRS}\) equals \(\nu|\Omega|^2\) plus a source built from \(Y\),
      \(\nabla Y\), \(V\) and \(z\); write the source explicitly and verify
      that it vanishes when \(Y=0\), recovering the classical identity.
\item Record the decay of \(\Pi_{\rm NRS}\) at infinity available from the
      class bounds (\(P_V\in L^3\), \(V\in L^6\), \(|V|\le C_\infty\)) and note
      that \(\frac12z\cdot V\) need not decay: the classical argument
      uses \(V\in L^q\) with a decay that the class gives only in the
      integral sense; state what replaces pointwise decay.
\item The next compulsory companion audit is V35.
\end{enumerate}
\end{programme}

\begin{assessment}[Position after V30]
\label{ass:c3v30-position}
The companion audit is recorded; the defect-rate block begins.  No
full stability or global-regularity claim is made.
\end{assessment}

% =============================================================================
\section*{Version V30 (third cycle) append-only record}
\addcontentsline{toc}{section}{Version V30 (third cycle) append-only record}
% =============================================================================

The complete V29 source of the third cycle was retained before this
continuation.  Its SHA--256 digest was
\begin{center}\scriptsize\ttfamily
5e2b2567a002f8916fba850ef03254c8eb92e9632e2b7e983d9c69199a7e34c2.
\end{center}
V30 performed the compulsory companion audit and recorded the
position.  No full regularity claim is made.

% =============================================================================
\section{V31: the Ne\v cas--R\accent23u\v zi\v cka--\v Sver\'ak identity with a defect, and why it cannot be quantified in the class}
\label{sec:c3v31-entry}
% =============================================================================

Items (AC1) and (AC2) of Programme~\ref{prog:c3v31} are carried out.
(AC1): for every smooth divergence-free field of the class, the
Bernoulli-type function of Ne\v cas, R\accent23u\v zi\v cka and \v Sver\'ak
satisfies a drift-elliptic equation whose right side is the squared
vorticity plus the inner product of the defect with the drift; the
source is exactly linear in the defect, with no derivative of it,
and the classical identity is the case of zero defect.  (AC2): the
two integrations that could exploit the sign of the squared
vorticity are examined.  Against the Gaussian weight the identity is
the energy identity combined with the closure identity of the
second cycle, so it adds nothing to the ledger; against the constant
weight it requires the decay of the field at infinity that the exact
profile equation supplies in the classical theorem and that the
class bounds do not, and the same decay is needed by the maximum
principle.  The route to an explicit defect rate through a
quantitative form of (E9) is therefore closed within the class, and
recorded as a negative; the block is redirected to the mean-square
form of the second identity, whose lower side is explicit.  No
global regularity claim is made.  The next compulsory companion
audit is V35.

% =============================================================================
\section{The identity}
\label{sec:c3v31-identity}
% =============================================================================

Throughout, \(V\) is a smooth divergence-free field with the class
bounds (\(|V|\le C_\infty\), \(V\in L^6\), \(\nabla V\in L^2\), derivatives bounded),
\(P_V\) its pressure (C2-V36), \(\Omega=\operatorname{curl}V\),
\(N(V)=(V\cdot\nabla)V+\nabla P_V\), \(L_-V=\nu\Delta V-\frac12z\cdot\nabla V-\frac12V\), and
\begin{equation}
 Y:=L_-V-N(V),\qquad
 W:=V+\tfrac12z,\qquad
 \Pi:=P_V+\tfrac12|V|^2+\tfrac12z\cdot V .
 \label{eq:c3v31-notation}
\end{equation}
\(Y\) is divergence-free (\(\operatorname{div}L_-V=0\) and \(\operatorname{div}N(V)=0\) by
the definition of \(P_V\)); for the profile of an element of the
family, \(Y=\partial_\sigma V\).

\begin{lemma}[Two vector identities]
\label{lem:c3v31-vector}
For smooth \(V\) with \(\operatorname{div}V=0\):
\((V\cdot\nabla)V=\nabla(\frac12|V|^2)+\Omega\times V\) and
\((z\cdot\nabla)V=\nabla(z\cdot V)-V-z\times\Omega\).
\end{lemma}

\begin{proof}
\(\nabla(a\cdot b)=(a\cdot\nabla)b+(b\cdot\nabla)a+a\times\operatorname{curl}b+b\times\operatorname{curl}a\)
with \(a=b=V\) gives the first; with \(a=z\), \(b=V\), \(\operatorname{curl}z=0\)
and \((V\cdot\nabla)z=V\) gives the second.
\end{proof}

\begin{proposition}[Gradient form of the defect equation]
\label{prop:c3v31-gradient}
\begin{equation}
 \nabla\Pi\ =\ \nu\Delta V-\Omega\times W-Y .
 \label{eq:c3v31-gradient}
\end{equation}
\end{proposition}

\begin{proof}
By Lemma~\ref{lem:c3v31-vector},
\(-\frac12(z\cdot\nabla)V-\frac12V=-\frac12\nabla(z\cdot V)+\frac12z\times\Omega\) and
\(-(V\cdot\nabla)V=-\nabla(\frac12|V|^2)-\Omega\times V\), so
\(Y=L_-V-N(V)=\nu\Delta V-\nabla\Pi+\frac12z\times\Omega-\Omega\times V=\nu\Delta V-\nabla\Pi-\Omega\times(V+\frac12z)\).
\end{proof}

\begin{theorem}[The NRS identity with a defect]
\label{thm:c3v31-NRS}
For every smooth divergence-free field of the class,
\begin{equation}
 \boxed{\quad
 \nu\Delta\Pi-W\cdot\nabla\Pi\ =\ \nu|\Omega|^2+W\cdot Y .
 \quad}
 \label{eq:c3v31-NRS}
\end{equation}
In particular, for the profile of an element of the family at every
\(\sigma\), with \(Y=\partial_\sigma V\); and for \(Y=0\) the classical identity
\(\nu\Delta\Pi-W\cdot\nabla\Pi=\nu|\Omega|^2\) of (E9) is recovered.
\end{theorem}

\begin{proof}
Take the divergence of \eqref{eq:c3v31-gradient}: \(\Delta\Pi=\nu\Delta\operatorname{div}V-\operatorname{div}(\Omega\times W)-\operatorname{div}Y=-\operatorname{div}(\Omega\times W)\).
Now \(\operatorname{div}(\Omega\times W)=W\cdot\operatorname{curl}\Omega-\Omega\cdot\operatorname{curl}W\), with
\(\operatorname{curl}W=\Omega\) and \(\operatorname{curl}\Omega=\operatorname{curl}\operatorname{curl}V=-\Delta V\); hence
\(\Delta\Pi=W\cdot\Delta V+|\Omega|^2\).  By \eqref{eq:c3v31-gradient},
\(\nu\,W\cdot\Delta V=W\cdot\nabla\Pi+W\cdot(\Omega\times W)+W\cdot Y=W\cdot\nabla\Pi+W\cdot Y\).
Multiplying \(\Delta\Pi=W\cdot\Delta V+|\Omega|^2\) by \(\nu\) and substituting gives
\eqref{eq:c3v31-NRS}.
\end{proof}

\begin{remark}[The form of the source]
\label{rem:c3v31-source}
The programme anticipated a source involving \(\nabla Y\); there is
none, because \(\operatorname{div}Y=0\).  The source \(W\cdot Y=V\cdot Y+\frac12z\cdot Y\)
is the pointwise density of the two inner products that the ledger
already averages: \(\langle V,Y\rangle_G=\frac12\varphi'\) and
\(\langle z,Y\rangle_G=0\) (the second because \(\int(z\cdot Y)G=-2\nu\int Y\cdot\nabla G=2\nu\int(\operatorname{div}Y)G=0\)
for every divergence-free \(Y\)).
\end{remark}

% =============================================================================
\section{The Gaussian integral is the ledger}
\label{sec:c3v31-gaussian}
% =============================================================================

For a smooth weight \(f\) let \(L:=\nu\Delta-W\cdot\nabla\) and
\(L^*f:=\nu\Delta f+\operatorname{div}(fW)=\nu\Delta f+\frac32f+W\cdot\nabla f\)
(\(\operatorname{div}W=\frac32\)), so that \(\int(L\Pi)f=\int\Pi\,L^*f\) whenever the
boundary terms vanish.

\begin{proposition}[The Gaussian weight]
\label{prop:c3v31-gaussian}
\(L^*G=-\frac{z\cdot V}{2\nu}\,G\), and the Gaussian integral of
\eqref{eq:c3v31-NRS} reads
\begin{equation}
 \nu\,\mathcal E+\langle W,Y\rangle_G\ =\ -\frac1{2\nu}\int\Pi\,(z\cdot V)\,G ,
 \qquad \mathcal E=\int|\Omega|^2G .
 \label{eq:c3v31-gaussian}
\end{equation}
This identity is equivalent to the energy identity
\(\frac12\varphi'=\langle V,Y\rangle_G=-2\mathcal L-\frac14j\) combined with the closure
identity \(R:=\int(z\cdot V)^2G=8\nu\mathcal L-4\nu^2\mathcal E\) (Theorem
C2-V21-closure); it adds nothing to the ledger.
\end{proposition}

\begin{proof}
\(\nabla G=-\frac z{2\nu}G\), \(\nu\Delta G=(\frac{|z|^2}{4\nu}-\frac32)G\), and
\(W\cdot\nabla G=-\frac{V\cdot z+\frac12|z|^2}{2\nu}G\); summing with \(\frac32G\) gives
\(L^*G=-\frac{z\cdot V}{2\nu}G\).  The boundary terms vanish by the
Gaussian decay against the polynomial growth of \(\Pi\) and \(W\).  For
the equivalence, expand \(\Pi\):
\(-\frac1{2\nu}\int\Pi(z\cdot V)G=-\frac1{2\nu}\int P_V(z\cdot V)G-\frac1{4\nu}\int|V|^2(z\cdot V)G-\frac1{4\nu}R=-\frac14j-\frac1{4\nu}R\),
by the flux formula \(j=\frac1\nu\int|V|^2(z\cdot V)G+\frac2\nu\int P_V(z\cdot V)G\)
(proof of Theorem C2-V36-flux-bound).  With the closure identity,
\(-\frac1{4\nu}R=-2\mathcal L+\nu\mathcal E\), so \eqref{eq:c3v31-gaussian} reads
\(\langle W,Y\rangle_G=-2\mathcal L-\frac14j\); and \(\langle W,Y\rangle_G=\langle V,Y\rangle_G+\frac12\langle z,Y\rangle_G=\langle V,Y\rangle_G\)
(Remark~\ref{rem:c3v31-source}), while
\(\langle V,Y\rangle_G=\langle V,L_-V\rangle_G-\langle V,N\rangle_G=-2\mathcal L-\frac14j\) is the
energy identity (C2-V1, C2-V40: \(\langle V,N\rangle_G=\frac14j\)).
\end{proof}

\begin{remark}[The closure identity is the Gaussian NRS structure]
\label{rem:c3v31-closure}
Conversely, the closure identity of C2-V21, found there as a
relation between the radial kinetic energy, the Dirichlet form and
the Gaussian enstrophy, is the Gaussian integral of the classical
identity \eqref{eq:c3v31-NRS}: the second cycle's ledger contained
the NRS structure in weighted form without naming it.  This is
consistent with the second cycle's closure of the ledger search
(C2-V22): the NRS structure produces no row that the ledger did not
already have.
\end{remark}

% =============================================================================
\section{The constant weight needs decay the class does not give}
\label{sec:c3v31-constant}
% =============================================================================

\begin{proposition}[The constant weight and the maximum principle]
\label{prop:c3v31-constant}
\begin{enumerate}[label=\textup{(\roman*)},leftmargin=3.2em]
\item \(L^*1=\frac32\), and for the cut-off \(\chi_R=\chi(\cdot/R)\) of V22,
      \begin{equation}
       \frac32\int\Pi\chi_R=\nu\int|\Omega|^2\chi_R+\int(W\cdot Y)\chi_R-\nu\int\Pi\Delta\chi_R-\int\Pi\,W\cdot\nabla\chi_R ,
       \label{eq:c3v31-cutoff}
      \end{equation}
      with boundary terms bounded by
      \(A(1+C_\infty/R)\int_{S_R}(|P_V|+\frac12|V|^2)+A(1+C_\infty/R)\,R\int_{S_R}|V|\)
      on the shell \(S_R=\{R\le|z|\le2R\}\) (for \(R\ge\nu\)).
\item The class bounds give only \(\int_{S_R}|P_V|\le\|P_V\|_{L^3(S_R)}|S_R|^{2/3}\)
      and \(R\int_{S_R}|V|\le R\|V\|_{L^6(S_R)}|S_R|^{5/6}\), i.e.\ \(o(R^2)\) and
      \(o(R^{7/2})\), and \(\int\Pi\chi_R\) itself contains \(\frac12\int(z\cdot V)\chi_R=-\frac14\int|z|^2V\cdot\nabla\chi_R\),
      of the same size: no term of \eqref{eq:c3v31-cutoff} has a
      limit as \(R\to\infty\) under the class bounds alone.
\item The maximum principle for \(L\Pi\ge W\cdot Y\) yields \(\Pi\le\sup_{\partial B_R}\Pi+(\text{source})\)
      and needs \(\Pi\to0\) at infinity, i.e.\ \(z\cdot V\to0\); the class
      gives \(|V|\le C_\infty\) and integral decay, not \(|V|=o(|z|^{-1})\).
\item The weights \(f=G^\alpha\) satisfy
      \(\nu\Delta f+\frac12z\cdot\nabla f=(\alpha(\alpha-1)\frac{|z|^2}{4\nu}-\frac32\alpha)f\), a
      multiple of \(f\) only for \(\alpha\in\{0,1\}\): no third radial Gaussian
      weight is adjoint-invariant for the drift \(\frac12z\).
\end{enumerate}
\end{proposition}

\begin{proof}
(i) \(\int(L\Pi)\chi_R=\int\Pi L^*\chi_R\) with \(L^*\chi_R=\nu\Delta\chi_R+\frac32\chi_R+W\cdot\nabla\chi_R\);
the bounds use \(|\nabla\chi_R|\le A/R\), \(|\Delta\chi_R|\le A/R^2\), \(|W|\le C_\infty+R\)
on \(S_R\), and \(|\Pi|\le|P_V|+\frac12|V|^2+R|V|\) there.  (ii) H\"older on the
shell, \(|S_R|=\frac{28\pi}3R^3\), and the tails of \(\|P_V\|_{L^3}\), \(\|V\|_{L^6}\)
tend to zero; the identity \(z\cdot V=\operatorname{div}(\frac12|z|^2V)\).  (iii) The
weak maximum principle for the operator \(L\) with bounded
coefficients on \(B_R\).  (iv) Direct computation, as in the proof of
Proposition~\ref{prop:c3v31-gaussian}.
\end{proof}

\begin{firewall}[Item (AC2): negative record N-C3-3, no quantitative NRS in the class]
\label{fw:c3v31-negative}
The exclusion (E9) uses, besides the identity \eqref{eq:c3v31-NRS}
with \(Y=0\), the decay of the profile at infinity, which
Ne\v cas--R\accent23u\v zi\v cka--\v Sver\'ak and Tsai derive from the exact
profile equation (an elliptic equation with the linear part
\(\nu\Delta-\frac12z\cdot\nabla-\frac12\), whose decaying solutions decay
fast).  A field with a small defect is not a solution of that
equation, and the class bounds give no pointwise decay
(Proposition~\ref{prop:c3v31-constant}(ii),(iii)); the only weight
under which the identity integrates within the class is the
Gaussian, and there it is the ledger
(Proposition~\ref{prop:c3v31-gaussian}).  Hence the defect rate
\(\underline\delta\) cannot be made explicit by quantifying (E9) with the
class bounds alone.  What would be needed is a decay statement for
approximate profiles, i.e.\ for fields of the class whose defect is
small in \(L^2(G)\): a quantitative unique-continuation-type result
of the kind already found missing on the lower side of the
constant-chasing gate (C2-V38).  The plan of
Assessment~\ref{ass:c3v29-direction} for V32--V34 is withdrawn.
\end{firewall}

% =============================================================================
\section{Redirection: the mean-square form of the second identity}
\label{sec:c3v31-redirect}
% =============================================================================

\begin{assessment}[Direction for V32--V34]
\label{ass:c3v31-direction}
The mean identities of C2-V12 give, by expanding \(\|Y\|_G^2\) and
\(\langle N,Y\rangle_G\) in \(L_-V\) and \(N\),
\(\langle\|N(V)\|_G^2\rangle_{\mathfrak m}=\langle\|L_-V\|_G^2\rangle_{\mathfrak m}+\langle\|Y\|_G^2\rangle_{\mathfrak m}\):
the nonlinearity exceeds the linear part in mean square by the mean
defect.  The lower side of this identity is explicit and pointwise:
\(\|L_-V\|_G^2=\sum_k\mu_k^2a_k\ge r^2\varphi\ge\frac14c_\Phi\); its upper side is
bounded by the class: \(\|N\|_G^2\le2C_\infty^2\|\nabla V\|_G^2+2\|\nabla P_V\|_G^2\),
and \(\|\nabla P_V\|_{L^2(\R^3)}\le2C_\infty\|\nabla V\|_{L^2(\R^3)}\) with no
Calder\'on--Zygmund constant, since the symbol of \(R_iR_j\) is a
projection.  V32 records the identity and the explicit constraint
that results, \(\frac14c_\Phi+\underline\delta\le\kappa^2(\frac12\langle d\rangle+8\nu C)\), and
compares it with the constant-chasing constraint of C2-V36; V33
examines what the constraint excludes (a joint bound on the
pinching, the velocity constant and the class enstrophy), and
whether the pressure part can be taken in Gaussian form with a
smaller constant; V34 consolidates the defect block V31--V33 with
the negative record N-C3-3 and fixes the direction for V36--V40;
V35 audits.  No global regularity claim is made.
\end{assessment}

% =============================================================================
\section{Integrated V31 conclusion and next acquisition order}
\label{sec:c3v31-conclusion}
% =============================================================================

\begin{theorem}[What V31 proves]
\label{thm:c3v31-conclusion}
Relative to the two legacy cycles and V1--V30: the NRS identity with
a defect (Theorem~\ref{thm:c3v31-NRS}), its Gaussian integral as the
ledger (Proposition~\ref{prop:c3v31-gaussian}), the constant-weight
and maximum-principle obstructions
(Proposition~\ref{prop:c3v31-constant}) and the negative record
N-C3-3.  The full Navier--Stokes stability/global-regularity theorem
remains unproved.
\end{theorem}

\begin{proof}
The cited results.
\end{proof}

\begin{programme}[V32 acquisition order]
\label{prog:c3v32}
\begin{enumerate}[label=\textup{(AC\arabic*)},leftmargin=4.8em]
\item Prove the mean-square identity
      \(\langle\|N\|_G^2\rangle=\langle\|L_-V\|_G^2\rangle+\langle\|Y\|_G^2\rangle\) from
      C2-V12-mean-identities, and the pointwise lower bound
      \(\|L_-V\|_G^2\ge(r^2+\mathrm{Var})\varphi\ge\frac14c_\Phi\).
\item Prove \(\|\nabla P_V\|_{L^2(\R^3)}\le2C_\infty\|\nabla V\|_{L^2(\R^3)}\) by
      Plancherel, and \(\|(V\cdot\nabla)V\|_G^2\le C_\infty^2\|\nabla V\|_G^2=\frac{\kappa^2}4d\);
      deduce \(\|N\|_G^2\le\frac{\kappa^2}2d+8\kappa^2\nu C\) and the explicit
      constraint at every homogeneous point.
\item The next compulsory companion audit is V35.
\end{enumerate}
\end{programme}

\begin{assessment}[Position after V31]
\label{ass:c3v31-position}
The NRS identity with a defect is exact and adds nothing to the
ledger in Gaussian form; the quantitative-NRS route is closed in the
class; the block turns to the mean-square identity.  No full
stability or global-regularity claim is made.
\end{assessment}

% =============================================================================
\section*{Version V31 (third cycle) append-only record}
\addcontentsline{toc}{section}{Version V31 (third cycle) append-only record}
% =============================================================================

The complete V30 source of the third cycle was retained before this
continuation.  Its SHA--256 digest was
\begin{center}\scriptsize\ttfamily
8b0be78ce8838926551e39ff4856a78b0364a7b7481e1414f46dab3806772681.
\end{center}
V31 proved the NRS identity with a defect and recorded the negative
record N-C3-3.  No full regularity claim is made.

% =============================================================================
\section{V32: the mean-square identity and its explicit constraint}
\label{sec:c3v32-entry}
% =============================================================================

Items (AC1) and (AC2) of Programme~\ref{prog:c3v32} are carried out.
(AC1): against every invariant measure at a homogeneous type-I
point, the mean square of the nonlinearity equals the mean square of
the linear part plus the mean square of the defect; the linear part
is bounded below pointwise by the squared Rayleigh quotient times
the Gaussian energy, hence by a quarter of the pinching constant.
(AC2): the nonlinearity is bounded above by the velocity constant
squared times a Gaussian dissipation and an unweighted enstrophy
term, the pressure gradient being bounded in \(L^2\) by twice the
velocity bound times the enstrophy with no singular-integral
constant.  The two sides give an explicit constraint on the
pinching, the defect rate, the velocity constant and the class
enstrophy, which is stronger than the constant-chasing constraint of
the second cycle for small velocity constants; it excludes nothing
by itself, because the pinching and the defect rate have no lower
bounds in the series.  No global regularity claim is made.  The
next compulsory companion audit is V35.

% =============================================================================
\section{The mean-square identity}
\label{sec:c3v32-identity}
% =============================================================================

\begin{theorem}[Mean-square form of the second identity]
\label{thm:c3v32-mean-square}
Let \(x_0\) be a homogeneous type-I point and \(\mathfrak m\) an invariant
probability measure on \(\mathcal F(x_0)\).  Then
\begin{equation}
 \boxed{\quad
 \langle\|N(V)\|_G^2\rangle_{\mathfrak m}\ =\ \langle\|L_-V\|_G^2\rangle_{\mathfrak m}+\langle\|Y\|_G^2\rangle_{\mathfrak m},
 \qquad
 \langle\langle L_-V,N(V)\rangle_G\rangle_{\mathfrak m}\ =\ \langle\|L_-V\|_G^2\rangle_{\mathfrak m} ,
 \quad}
 \label{eq:c3v32-mean-square}
\end{equation}
and \(\langle\|Y\|_G^2\rangle_{\mathfrak m}\ge\underline\delta\).
\end{theorem}

\begin{proof}
Pointwise, \(\|Y\|_G^2=\|L_-V\|_G^2-2\langle L_-V,N\rangle_G+\|N\|_G^2\) and
\(\langle N,Y\rangle_G=\langle L_-V,N\rangle_G-\|N\|_G^2\); their sum is
\(\|L_-V\|_G^2-\langle L_-V,N\rangle_G\).  The mean identity
\(\langle\|Y\|_G^2\rangle+\langle\langle N,Y\rangle_G\rangle=0\) (C2-V12-mean-identities,
the second line) gives the second equality of
\eqref{eq:c3v32-mean-square}; substituting it into the mean of the
first expansion gives the first.  The lower bound on the mean defect
is the definition of \(\underline\delta\) with C2-V12.
\end{proof}

\begin{proposition}[Explicit lower side]
\label{prop:c3v32-lower}
For every profile of the family and every \(\sigma\),
\begin{equation}
 \|L_-V\|_G^2=\sum_k\mu_k^2a_k=(r^2+\mathrm{Var})\,\varphi\ \ge\ \Bigl(\frac14+v_{\rm exp}\Bigr)c_\Phi .
 \label{eq:c3v32-lower}
\end{equation}
\end{proposition}

\begin{proof}
Spectral theorem (C2-V31-degree): \(\|L_-V\|_G^2=\sum_k\mu_k^2a_k=\varphi\sum_k\mu_k^2\hat a_k\),
and \(\sum_k\mu_k^2\hat a_k=\mathrm{Var}+r^2\) by \eqref{eq:c3v21-spectral}; then
\(r\ge\frac12\), \(\mathrm{Var}\ge v_{\rm exp}\) (Theorem~\ref{thm:c3v23-floor}) and
\(\varphi\ge c_\Phi\).
\end{proof}

% =============================================================================
\section{The explicit upper side}
\label{sec:c3v32-upper}
% =============================================================================

\begin{lemma}[The pressure gradient in \(L^2\)]
\label{lem:c3v32-pressure}
For every smooth divergence-free \(V\) with \(|V|\le C_\infty\) and
\(\nabla V\in L^2(\R^3)\),
\begin{equation}
 \|\nabla P_V\|_{L^2(\R^3)}\ \le\ \|\nabla(V\otimes V)\|_{L^2(\R^3)}\ \le\ 2C_\infty\|\nabla V\|_{L^2(\R^3)} .
 \label{eq:c3v32-pressure}
\end{equation}
\end{lemma}

\begin{proof}
\(-\Delta P_V=\partial_i\partial_j(V_iV_j)\) gives, in Fourier variables,
\(\widehat{\partial_lP_V}(\xi)=-\sum_{ij}\frac{\xi_i\xi_j}{|\xi|^2}\,\widehat{\partial_l(V_iV_j)}(\xi)\).
The matrix \(\xi\otimes\xi/|\xi|^2\) has Frobenius norm one, so
\(|\widehat{\partial_lP_V}(\xi)|\le|\widehat{\partial_l(V\otimes V)}(\xi)|\) (Frobenius norm of
the matrix-valued transform) for every \(\xi\), and Plancherel gives
\(\|\partial_lP_V\|_2\le\|\partial_l(V\otimes V)\|_2\).  Pointwise
\(|\partial_l(V\otimes V)|\le2|V||\partial_lV|\le2C_\infty|\partial_lV|\); sum over \(l\).
\end{proof}

\begin{proposition}[Explicit upper side]
\label{prop:c3v32-upper}
For every profile of the family and every \(\sigma\), with
\(\kappa=C_\infty\nu^{-1/2}\), \(d=4\nu\|\nabla V\|_G^2\) and \(C\) the class
enstrophy bound \(\|\nabla V\|_{L^2(\R^3)}^2\le C\),
\begin{equation}
 \|N(V)\|_G^2\ \le\ 2\|(V\cdot\nabla)V\|_G^2+2\|\nabla P_V\|_G^2\ \le\ \frac{\kappa^2}2\,d+8\kappa^2\nu C .
 \label{eq:c3v32-upper}
\end{equation}
\end{proposition}

\begin{proof}
\(|(V\cdot\nabla)V|\le C_\infty|\nabla V|\) gives \(\|(V\cdot\nabla)V\|_G^2\le C_\infty^2\|\nabla V\|_G^2=C_\infty^2d/(4\nu)=\kappa^2d/4\).
Since \(G\le1\), \(\|\nabla P_V\|_G^2\le\|\nabla P_V\|_{L^2}^2\le4C_\infty^2\|\nabla V\|_{L^2}^2\le4C_\infty^2C=4\kappa^2\nu C\)
by Lemma~\ref{lem:c3v32-pressure}.
\end{proof}

% =============================================================================
\section{The constraint}
\label{sec:c3v32-constraint}
% =============================================================================

\begin{theorem}[Explicit mean-square constraint at a homogeneous point]
\label{thm:c3v32-constraint}
At every homogeneous type-I point, for every invariant probability
measure \(\mathfrak m\),
\begin{equation}
 \boxed{\quad
 \Bigl(\frac14+v_{\rm exp}\Bigr)c_\Phi+\underline\delta\ \le\ \langle(r^2+\mathrm{Var})\varphi\rangle+\langle\|Y\|_G^2\rangle\ \le\ \kappa^2\Bigl(\frac12\langle d\rangle+8\nu C\Bigr) ,
 \quad}
 \label{eq:c3v32-constraint}
\end{equation}
and, for \(\kappa<\frac12\), with \(m_\kappa=(2-4\kappa)/3\) and the mean
constraint of C2-V36,
\begin{equation}
 \Bigl(\frac14+v_{\rm exp}\Bigr)c_\Phi+\underline\delta\ \le\ \kappa^2\Bigl(\frac{K_P^2}{2m_\kappa^2}+8\nu C\Bigr) .
 \label{eq:c3v32-constraint-small}
\end{equation}
\end{theorem}

\begin{proof}
Theorem~\ref{thm:c3v32-mean-square} with Propositions~\ref{prop:c3v32-lower}
and \ref{prop:c3v32-upper} averaged; then \(\langle d\rangle\le K_P^2/m_\kappa^2\)
(Theorem C2-V36-constraint) for \(\kappa<\frac12\).
\end{proof}

\begin{assessment}[Comparison with the constant-chasing constraint]
\label{ass:c3v32-comparison}
The second cycle's constraint (C2-V36) bounds the pinching by
\(c_\Phi\le K_P^2/(3m_\kappa^2)\), with no factor of \(\kappa\); the present
one bounds it by \(4\kappa^2(K_P^2/(2m_\kappa^2)+8\nu C)\), which is smaller
when \(\kappa^2<\frac16\cdot\frac{K_P^2/m_\kappa^2}{K_P^2/m_\kappa^2+16\nu C}\).  It also bounds
the defect rate: \(\underline\delta\le\kappa^2(\frac12\langle d\rangle+8\nu C)\), so that
a homogeneous blow-up with a small velocity constant is nearly
self-similar in the mean-square sense, its mean defect being at most
\(\kappa^2\) times an explicit class quantity, while its pinching is at
most \(\kappa^2\) times the same.  Both readings are consistent with the
class: a small velocity constant means a small profile, hence a
small nonlinearity, hence a small defect and a small Gaussian energy.
The constraint is a new explicit row of the ledger; it is not an
exclusion, because neither \(c_\Phi\) nor \(\underline\delta\) has a lower bound
in the series beyond positivity, and the \(\varepsilon\)-regularity floor
\(\kappa\ge\varepsilon_{\rm reg}\) bounds \(\kappa\) below, not above.
\end{assessment}

\begin{firewall}[What the constraint does not do]
\label{fw:c3v32-no-exclusion}
An exclusion would follow from \eqref{eq:c3v32-constraint-small} if
the pinching or the defect rate were bounded below explicitly in
terms of the class constants by more than the right side.  The
pinching \(c_\Phi\) is produced by compactness from the homogeneity
constant (C2-V12) and the defect rate \(\underline\delta\) by compactness
from (E9) (C2-V6); V31 recorded that the latter cannot be made
explicit within the class, and the former is subject to the same
obstruction (its lower side is the lower side of the constant-chasing
gate, C2-V37--V38).  The constraint therefore relates three
non-explicit or unbounded constants to explicit ones, and its
content is the reading of Assessment~\ref{ass:c3v32-comparison}.
\end{firewall}

% =============================================================================
\section{Integrated V32 conclusion and next acquisition order}
\label{sec:c3v32-conclusion}
% =============================================================================

\begin{theorem}[What V32 proves]
\label{thm:c3v32-conclusion}
Relative to the two legacy cycles and V1--V31: the mean-square
identity (Theorem~\ref{thm:c3v32-mean-square}), the explicit lower and
upper sides (Propositions~\ref{prop:c3v32-lower}, \ref{prop:c3v32-upper},
Lemma~\ref{lem:c3v32-pressure}) and the constraint
(Theorem~\ref{thm:c3v32-constraint}).  The full Navier--Stokes
stability/global-regularity theorem remains unproved.
\end{theorem}

\begin{proof}
The cited results.
\end{proof}

\begin{programme}[V33 acquisition order]
\label{prog:c3v33}
\begin{enumerate}[label=\textup{(AC\arabic*)},leftmargin=4.8em]
\item Sharpen the upper side: the advective part in the form
      \(\|(V\cdot\nabla)V\|_G^2\le\|V\|_{L^\infty(\mathrm{supp})}^2\|\nabla V\|_G^2\) is already
      Gaussian; examine whether the pressure part admits a Gaussian
      bound \(\|\nabla P_V\|_G^2\le c\,C_\infty^2(\|\nabla V\|_G^2+\text{tail})\) by
      splitting the source \(q=\partial_iV_j\partial_jV_i\) at a radius \(\rho\) as in
      V27 (near part by the \(L^2\) bound of Lemma~\ref{lem:c3v32-pressure}
      applied to a cut-off, far part by the Newtonian kernel), so
      that the constraint reads \(\frac14c_\Phi+\underline\delta\le\kappa^2(c_1\langle d\rangle+c_2(\rho)\nu C)\)
      with \(c_2(\rho)\to0\); record the best form.
\item The pointwise version: \(\|Y\|_G^2=\|N\|_G^2-\|L_-V\|_G^2+2\langle Y,L_-V\rangle_G\)
      with \(\langle Y,L_-V\rangle_G=-\frac{\dd}{\dd\sigma}\mathcal L\) (C2-V2-second);
      determine what the explicit sides say pointwise about
      \(\frac{\dd}{\dd\sigma}\mathcal L\) at times when the defect is small.
\item The next compulsory companion audit is V35.
\end{enumerate}
\end{programme}

\begin{assessment}[Position after V32]
\label{ass:c3v32-position}
The mean-square identity gives an explicit constraint relating the
pinching, the defect rate, the velocity constant and the class
enstrophy; it excludes nothing by itself.  No full stability or
global-regularity claim is made.
\end{assessment}

% =============================================================================
\section*{Version V32 (third cycle) append-only record}
\addcontentsline{toc}{section}{Version V32 (third cycle) append-only record}
% =============================================================================

The complete V31 source of the third cycle was retained before this
continuation.  Its SHA--256 digest was
\begin{center}\scriptsize\ttfamily
f2c18ed269b0ed5d04593f915de25556b6e09a4bdb68743f44b99e549fad3af2.
\end{center}
V32 proved the mean-square identity and its explicit constraint.
No full regularity claim is made.

% =============================================================================
\section{V33: the pointwise and window forms of the mean-square identity}
\label{sec:c3v33-entry}
% =============================================================================

Items (AC1) and (AC2) of Programme~\ref{prog:c3v33} are carried out.
(AC1): the pressure part of the upper side cannot be put in Gaussian
form without an exponential price, because the pressure gradient on
the Gaussian bulk is generated by the strain everywhere and the
class controls the strain only in the unweighted sense; the form of
V32 is kept and the obstruction recorded.  (AC2): the mean-square
identity is the mean of an exact pointwise identity, the derivative
of the Dirichlet form being half the excess of the squared
nonlinearity over the squared linear part and the squared defect;
its explicit sides give a pointwise criterion for the descent of
the Dirichlet form and, integrated over a window, an explicit bound
on the window-averaged defect of every element of the family, not
only of invariant measures.  No global regularity claim is made.
The next compulsory companion audit is V35.

% =============================================================================
\section{The pressure part stays unweighted}
\label{sec:c3v33-pressure}
% =============================================================================

\begin{firewall}[Item (AC1): no Gaussian form for the pressure part]
\label{fw:c3v33-pressure}
Let \(\rho>0\).  Splitting the Gaussian integral at \(|z|=\rho\),
\(\|\nabla P_V\|_G^2\le\int_{B_\rho}|\nabla P_V|^2G+e^{-\rho^2/(4\nu)}\|\nabla P_V\|_{L^2}^2\), and
the second term is at most \(4\kappa^2\nu Ce^{-\rho^2/(4\nu)}\) by
Lemma~\ref{lem:c3v32-pressure}.  On \(B_\rho\), \(\nabla P_V=\nabla\Gamma_N*q\) splits
into the part generated by the strain in \(B_{2\rho}\) and the part
generated outside; the outside part is at most \(C/(4\pi\rho^2)\)
pointwise on \(B_\rho\) (as in Proposition~\ref{prop:c3v27-identity}), and
the inside part is bounded in \(L^2\) by the \(L^2\) norm of the strain
in \(B_{4\rho}\) times \(2C_\infty\) plus cut-off terms, by the argument of
Lemma~\ref{lem:c3v32-pressure} applied to \(\chi_{2\rho}V\otimes V\).  The strain in
\(B_{4\rho}\) is controlled by the Gaussian dissipation only at the
price \(e^{4\rho^2/\nu}\): \(\int_{B_{4\rho}}|\nabla V|^2\le e^{4\rho^2/\nu}\|\nabla V\|_G^2\).  Hence
every Gaussian form of the pressure part has the shape
\(\|\nabla P_V\|_G^2\le c_1e^{4\rho^2/\nu}\kappa^2d+c_2(\rho)\kappa^2\nu C\) with \(c_2(\rho)\to0\), and
the exponential factor makes it weaker than the unweighted form of
Proposition~\ref{prop:c3v32-upper} for every \(\rho\) of order
\((4\nu)^{1/2}\) or larger, while for smaller \(\rho\) the tail term is not
small.  The unweighted form is kept.  The obstruction is the same
as in V27: the Gaussian weight and the nonlocal pressure do not
commute, and the class bounds are unweighted.
\end{firewall}

% =============================================================================
\section{The pointwise identity}
\label{sec:c3v33-pointwise}
% =============================================================================

\begin{proposition}[Pointwise form]
\label{prop:c3v33-pointwise}
For every profile of the family and every \(\sigma\),
\begin{equation}
 \boxed{\quad
 2\,\frac{\dd\mathcal L}{\dd\sigma}\ =\ \|N(V)\|_G^2-\|L_-V\|_G^2-\|Y\|_G^2 ,
 \quad}
 \label{eq:c3v33-pointwise}
\end{equation}
and therefore
\begin{equation}
 2\,\frac{\dd\mathcal L}{\dd\sigma}\ \le\ \kappa^2\Bigl(\frac{d}2+8\nu C\Bigr)-(r^2+\mathrm{Var})\varphi-\|Y\|_G^2 .
 \label{eq:c3v33-pointwise-bound}
\end{equation}
In particular the Dirichlet form decreases at every \(\sigma\) at which
\(\kappa^2(\frac d2+8\nu C)<(r^2+\mathrm{Var})\varphi\), and a fortiori at every \(\sigma\) at
which \(\kappa^2(\frac d2+8\nu C)<(\frac14+v_{\rm exp})c_\Phi\).
\end{proposition}

\begin{proof}
\(\frac{\dd}{\dd\sigma}\mathcal L=-\|Y\|_G^2-\langle N,Y\rangle_G=-\langle Y,L_-V\rangle_G\) (Theorem
C2-V2-second, with \(Y+N=L_-V\)).  Expanding \(N=L_-V-Y\):
\(\|N\|_G^2=\|L_-V\|_G^2-2\langle L_-V,Y\rangle_G+\|Y\|_G^2\), so
\(-2\langle Y,L_-V\rangle_G=\|N\|_G^2-\|L_-V\|_G^2-\|Y\|_G^2\).  The bound is
Propositions~\ref{prop:c3v32-lower} and \ref{prop:c3v32-upper}.
\end{proof}

\begin{remark}[Reading]
\label{rem:c3v33-reading}
The Dirichlet form of the profile grows exactly when the
nonlinearity, in Gaussian mean square, exceeds the linear part and
the defect together.  Since the mean of \(\frac{\dd}{\dd\sigma}\mathcal L\)
vanishes against invariant measures, the excess averages to zero,
which is Theorem~\ref{thm:c3v32-mean-square}.  A small velocity
constant makes the excess negative at every \(\sigma\) at which the
Gaussian dissipation is not large, and the Dirichlet form then
decreases; it cannot decrease forever on a family on which it is
pinched below by \(\frac14c_\Phi\), so the Gaussian dissipation must be
large often enough: this is the pointwise mechanism behind the mean
constraint of V32.
\end{remark}

% =============================================================================
\section{The window form}
\label{sec:c3v33-window}
% =============================================================================

Let \(C_{\mathcal L}:=(2\nu C+C_\Phi)/4\), so that \(0<\frac14c_\Phi\le\mathcal L\le C_{\mathcal L}\) on the
family (\(\mathcal L=\frac18(d+2\varphi)\), \(d\le4\nu C\), \(\varphi\le C_\Phi\)).

\begin{theorem}[Window form of the mean-square identity]
\label{thm:c3v33-window}
For every \(W\in\mathcal F(x_0)\) at a homogeneous type-I point, every
\(\sigma_0\) and every \(L>0\), writing \(\langle f\rangle_L:=\frac1L\int_{\sigma_0}^{\sigma_0+L}f\,\dd\sigma\),
\begin{equation}
 \boxed{\quad
 \Bigl|\langle\|N\|_G^2\rangle_L-\langle\|L_-V\|_G^2\rangle_L-\langle\|Y\|_G^2\rangle_L\Bigr|\ \le\ \frac{2C_{\mathcal L}}{L} ,
 \quad}
 \label{eq:c3v33-window}
\end{equation}
and hence
\begin{equation}
 \langle\|Y\|_G^2\rangle_L\ \le\ \kappa^2\Bigl(\frac12\langle d\rangle_L+8\nu C\Bigr)-\Bigl(\frac14+v_{\rm exp}\Bigr)c_\Phi+\frac{2C_{\mathcal L}}L .
 \label{eq:c3v33-window-defect}
\end{equation}
\end{theorem}

\begin{proof}
Integrate \eqref{eq:c3v33-pointwise} over the window:
\(2(\mathcal L(\sigma_0+L)-\mathcal L(\sigma_0))=\int(\|N\|^2-\|L_-V\|^2-\|Y\|^2)\), and
\(|\mathcal L(\sigma_0+L)-\mathcal L(\sigma_0)|\le C_{\mathcal L}\) (both values lie in
\([\frac14c_\Phi,C_{\mathcal L}]\)).  Divide by \(L\); then the explicit sides.
\end{proof}

\begin{corollary}[Explicit window bound on the defect of every element]
\label{cor:c3v33-defect-window}
For every element of the family, every window of log-length \(L\),
the window-averaged defect satisfies
\(\langle\|Y\|_G^2\rangle_L\le\kappa^2(\frac12\langle d\rangle_L+8\nu C)+2C_{\mathcal L}/L\); and since
\(\langle d\rangle_L\le4\nu C\) always, \(\langle\|Y\|_G^2\rangle_L\le10\kappa^2\nu C+2C_{\mathcal L}/L\)
for every element, without invariant measures.  In particular the
windowed defect of C2-V6, \(m(L)\le L\,\langle\|Y\|_G^2\rangle_L\), satisfies
\(m(L)\le10\kappa^2\nu CL+2C_{\mathcal L}\) and \(\underline\delta\le10\kappa^2\nu C\).
\end{corollary}

\begin{proof}
\eqref{eq:c3v33-window-defect} with \(\frac12\langle d\rangle_L\le2\nu C\); the
windowed defect is a minimum over the family, hence at most the
value on any element; \(\underline\delta=\lim m(L)/L\).
\end{proof}

\begin{remark}[The defect rate is at most \(\kappa^2\) times the class enstrophy]
\label{rem:c3v33-rate}
\(\underline\delta\le10\kappa^2\nu C=10C_\infty^2C\) is a pointwise-derived, explicit
upper bound on the compactness constant of the gate, complementing
the trivial \(\underline\delta\le C_N\) of C2-V6.  It says that the
anti-alignment rate available to the gate G1 is at most the square
of the velocity bound times the class enstrophy: any argument that
closes G1 by exhibiting a positive lower bound on the mean of
\(\langle N,Y\rangle_G+\|Y\|_G^2\) must produce a quantity smaller than this.
\end{remark}

% =============================================================================
\section{Integrated V33 conclusion and next acquisition order}
\label{sec:c3v33-conclusion}
% =============================================================================

\begin{theorem}[What V33 proves]
\label{thm:c3v33-conclusion}
Relative to the two legacy cycles and V1--V32: the pointwise form
(Proposition~\ref{prop:c3v33-pointwise}), the window form
(Theorem~\ref{thm:c3v33-window}) and the explicit window bound on
the defect of every element (Corollary~\ref{cor:c3v33-defect-window}).
The full Navier--Stokes stability/global-regularity theorem remains
unproved.
\end{theorem}

\begin{proof}
The cited results.
\end{proof}

\begin{programme}[V34 acquisition order]
\label{prog:c3v34}
\begin{enumerate}[label=\textup{(AC\arabic*)},leftmargin=4.8em]
\item Consolidate V31--V33 as the defect block: the NRS identity with
      a defect, the negative record N-C3-3, the mean-square identity
      with its explicit sides, the pointwise and window forms, and
      the explicit upper bound on the defect rate.
\item Fix the direction for V36--V40.  Two candidates: (a) the
      pointwise descent criterion of Proposition~\ref{prop:c3v33-pointwise}
      combined with the trap chapter of the second cycle (the
      Gaussian dissipation \(d\) is bounded below on trapped windows
      by the energetic-cell dissipation of C2-V48), to determine
      whether the Dirichlet form must oscillate with an explicit
      amplitude; (b) the constant-chasing chapter revisited with the
      new row: the joint constraint on \((c_\Phi,\underline\delta,\kappa,C)\)
      together with C2-V36--V38, to determine whether any region of
      parameters is excluded once the \(\varepsilon\)-regularity floor
      \(\varepsilon_{\rm reg}\) is written explicitly from (E5).
\item V35 compulsory companion audit and delivery.
\end{enumerate}
\end{programme}

\begin{assessment}[Position after V33]
\label{ass:c3v33-position}
The mean-square identity has pointwise and window forms; the defect
rate is at most \(10\kappa^2\nu C\).  No full stability or
global-regularity claim is made.
\end{assessment}

% =============================================================================
\section*{Version V33 (third cycle) append-only record}
\addcontentsline{toc}{section}{Version V33 (third cycle) append-only record}
% =============================================================================

The complete V32 source of the third cycle was retained before this
continuation.  Its SHA--256 digest was
\begin{center}\scriptsize\ttfamily
74470da1f51f5fbb8d3aeb1619d9d23992a5b7d0107a56de256eeeb40f68eaf8.
\end{center}
V33 proved the pointwise and window forms and the explicit bound on
the defect rate.  No full regularity claim is made.

% =============================================================================
\section{V34: the defect block consolidated, and extremal configurations}
\label{sec:c3v34-entry}
% =============================================================================

Items (AC1) and (AC2) of Programme~\ref{prog:c3v34} are carried out.
The defect block V31--V33 is consolidated.  For the direction, the
first of the two candidates of V33 is sharpened into a method: on
the compact invariant family every bounded continuously
differentiable observable attains its supremum and infimum, and at
the attaining configuration its log-time derivative vanishes, so
that every identity of the ledger holds there pointwise with the
derivative term removed.  Applied to the Dirichlet form with the
explicit sides of V32 this gives an explicit upper bound on the
supremum of the Gaussian dissipation over the family, and, for a
small velocity constant, an explicit positive lower bound on its
infimum, the lower side that the constant-chasing chapter of the
second cycle had recorded as non-explicit.  The block V36--V40 is
fixed as the systematic use of extremal configurations.  A
superfluous hypothesis in the corollary of V33 is removed in place.
No global regularity claim is made.  The next compulsory companion
audit is V35.

% =============================================================================
\section{The defect block}
\label{sec:c3v34-block}
% =============================================================================

\begin{theorem}[The defect block, V31--V33]
\label{thm:c3v34-block}
At a homogeneous type-I point:
\begin{enumerate}[label=\textup{(\arabic*)},leftmargin=3.2em]
\item (NRS with a defect, V31.)  \(\nu\Delta\Pi-W\cdot\nabla\Pi=\nu|\Omega|^2+W\cdot Y\)
      for every field of the class, with \(\Pi=P_V+\frac12|V|^2+\frac12z\cdot V\),
      \(W=V+\frac12z\), \(Y=L_-V-N(V)\); its Gaussian integral is the energy
      identity with the closure identity of C2-V21; its constant
      integral and the maximum principle need decay the class does
      not give (negative record N-C3-3: the defect rate cannot be
      made explicit by quantifying (E9) in the class).
\item (Mean-square identity, V32.)  \(\langle\|N\|_G^2\rangle=\langle\|L_-V\|_G^2\rangle+\langle\|Y\|_G^2\rangle\);
      \(\|L_-V\|_G^2=(r^2+\mathrm{Var})\varphi\ge(\frac14+v_{\rm exp})c_\Phi\);
      \(\|N\|_G^2\le\frac{\kappa^2}2d+8\kappa^2\nu C\) with
      \(\|\nabla P_V\|_{L^2}\le2C_\infty\|\nabla V\|_{L^2}\); the constraint
      \((\frac14+v_{\rm exp})c_\Phi+\underline\delta\le\kappa^2(\frac12\langle d\rangle+8\nu C)\).
\item (Pointwise and window forms, V33.)
      \(2\mathcal L'=\|N\|_G^2-\|L_-V\|_G^2-\|Y\|_G^2\); the window average of
      the excess is at most \(2C_{\mathcal L}/L\) in absolute value for every
      element; \(\underline\delta\le10\kappa^2\nu C\); the pressure part has no
      Gaussian form without an exponential price.
\end{enumerate}
\end{theorem}

\begin{proof}
Theorems~\ref{thm:c3v31-NRS}, \ref{thm:c3v32-mean-square},
\ref{thm:c3v32-constraint}, \ref{thm:c3v33-window};
Propositions~\ref{prop:c3v31-gaussian}, \ref{prop:c3v31-constant},
\ref{prop:c3v32-lower}, \ref{prop:c3v32-upper}, \ref{prop:c3v33-pointwise};
Corollary~\ref{cor:c3v33-defect-window}.
\end{proof}

\begin{remark}[In-place correction in V33]
\label{rem:c3v34-correction}
Corollary~\ref{cor:c3v33-defect-window} was stated with the
hypothesis \(\kappa<\frac12\), which none of its bounds uses (the bound
\(\langle d\rangle_L\le4\nu C\) replaces the mean constraint of C2-V36 there);
the hypothesis has been removed in the V33 text, in the draft and
in the file, before this version was appended.  The V33 digest
recorded below is that of the corrected file.
\end{remark}

% =============================================================================
\section{Extremal configurations}
\label{sec:c3v34-extremal}
% =============================================================================

\begin{lemma}[Extremal configurations]
\label{lem:c3v34-extremal}
Let \(x_0\) be a homogeneous type-I point and \(F:\mathcal F(x_0)\to\R\) a
continuous observable such that \(\sigma\mapsto F(\Theta_\sigma W)\) is \(C^1\) for
every \(W\).  Then \(F\) attains its supremum and its infimum on
\(\mathcal F(x_0)\), and at an attaining element \(W^*\),
\(\frac{\dd}{\dd\sigma}F(\Theta_\sigma W^*)|_{\sigma=0}=0\).
\end{lemma}

\begin{proof}
\((\mathcal F(x_0),\rho)\) is compact (C2-V4) and \(F\) continuous, so the
extrema are attained.  The family is \(\Theta\)-invariant for all
\(\sigma\in\R\) (its elements are ancient solutions of the class, and
their rescalings by every factor lie in the family, C2-V3--V4), so
\(\sigma\mapsto F(\Theta_\sigma W^*)\) is a \(C^1\) function on \(\R\) with an
interior extremum at \(\sigma=0\); its derivative vanishes there.
\end{proof}

\begin{theorem}[Explicit bounds on the extremal Dirichlet form]
\label{thm:c3v34-dirichlet}
Let \(\mathcal L_{\max}:=\sup_{\mathcal F(x_0)}\mathcal L\) and \(\mathcal L_{\min}:=\inf_{\mathcal F(x_0)}\mathcal L\).
Then
\begin{align}
 \mathcal L_{\max}&\ \le\ \mathcal L_+:=\frac{\kappa^2C_\Phi}2+\frac12\bigl(\kappa^4C_\Phi^2+8\kappa^2\nu CC_\Phi\bigr)^{1/2},
 \label{eq:c3v34-Lmax}\\
 \mathcal L_{\min}&\ \ge\ \mathcal L_-:=\frac{c_\Phi}4\Bigl(1+\frac1{4\kappa^2}\Bigr)-2\nu C ,
 \label{eq:c3v34-Lmin}
\end{align}
and consequently, for the Gaussian dissipation \(d=8\mathcal L-2\varphi\) on the
whole family,
\begin{equation}
 \boxed{\quad
 8\mathcal L_--2C_\Phi\ \le\ d\ \le\ 8\mathcal L_+-2c_\Phi .
 \quad}
 \label{eq:c3v34-d}
\end{equation}
The lower bound is positive whenever
\(\kappa^2<c_\Phi/(4(8\nu C+C_\Phi-c_\Phi))\).
\end{theorem}

\begin{proof}
Apply Lemma~\ref{lem:c3v34-extremal} to \(F=\mathcal L\) (\(C^1\) in \(\sigma\) by
C2-V2).  At a maximizer, \(0=2\mathcal L'=\|N\|_G^2-\|L_-V\|_G^2-\|Y\|_G^2\le\|N\|_G^2-\|L_-V\|_G^2\)
(Proposition~\ref{prop:c3v33-pointwise}), and with the explicit sides
(Propositions~\ref{prop:c3v32-lower}, \ref{prop:c3v32-upper}),
\(d=8\mathcal L-2\varphi\) and \(r^2\varphi=4\mathcal L^2/\varphi\):
\(0\le\kappa^2(4\mathcal L-\varphi+8\nu C)-4\mathcal L^2/\varphi\), i.e.\
\(4\mathcal L^2-4\kappa^2\varphi\mathcal L-\kappa^2\varphi(8\nu C-\varphi)\le0\), whence
\(\mathcal L\le\frac{\varphi}2\bigl[\kappa^2+(\kappa^4+\kappa^2(8\nu C-\varphi)/\varphi)^{1/2}\bigr]\le\frac{\kappa^2\varphi}2+\frac12(\kappa^4\varphi^2+8\kappa^2\nu C\varphi)^{1/2}\),
and \(\varphi\le C_\Phi\) gives \eqref{eq:c3v34-Lmax}.  At a minimizer,
\(0=2\mathcal L'\) gives \(\|N\|_G^2=\|L_-V\|_G^2+\|Y\|_G^2\ge\|L_-V\|_G^2\ge\frac14\varphi\), so
\(\frac14\varphi\le\kappa^2(4\mathcal L-\varphi+8\nu C)\), i.e.\
\(\mathcal L\ge\frac\varphi4(1+\frac1{4\kappa^2})-2\nu C\), and \(\varphi\ge c_\Phi\) gives
\eqref{eq:c3v34-Lmin} (the coefficient of \(\varphi\) is positive).  The
bounds on \(d\) follow from \(c_\Phi\le\varphi\le C_\Phi\) and
\(\mathcal L_-\le\mathcal L\le\mathcal L_+\) on the family.  Positivity of the lower
bound: \(2c_\Phi(1+\frac1{4\kappa^2})>16\nu C+2C_\Phi\).
\end{proof}

\begin{assessment}[What the extremal bounds add]
\label{ass:c3v34-extremal}
The second cycle's constant-chasing chapter had an explicit upper
side on the mean dissipation (C2-V36) and no explicit lower side
(C2-V37--V38).  Theorem~\ref{thm:c3v34-dirichlet} gives both sides
for the infimum and supremum over the whole family, not only for
means, at the price of the velocity constant: the upper bound
\eqref{eq:c3v34-Lmax} improves the class bound \(d\le4\nu C\) when
\(\kappa\) is small, and the lower bound \eqref{eq:c3v34-Lmin} is
positive when \(\kappa^2\) is below an explicit ratio of the pinching to
the class constants.  Neither excludes a homogeneous point: the
lower bound on \(d\) is consistent with the upper bound as long as
\(8\mathcal L_--2C_\Phi\le8\mathcal L_+-2c_\Phi\), which is an inequality between
\(c_\Phi/\kappa^2\) and the class constants that the series cannot decide,
the pinching having no explicit lower bound.  The method is what
matters: an extremal configuration turns every mean identity of the
ledger into a pointwise identity at one configuration, and the
pointwise explicit sides of V21--V33 can then be applied to it.
\end{assessment}

% =============================================================================
\section{Direction for V36--V40: the extremal ledger}
\label{sec:c3v34-direction}
% =============================================================================

\begin{assessment}[The extremal ledger]
\label{ass:c3v34-direction}
The block V36--V40 applies Lemma~\ref{lem:c3v34-extremal} to the
observables of the ledger in turn, and records at each extremal
configuration the pointwise identity with its explicit sides.
V36: the Gaussian energy \(\varphi\) (at extremal energy, \(j=-8\mathcal L\)
exactly, the inward flux equals eight times the Dirichlet form;
combined with the flux bound of C2-V36 this bounds \(\mathcal L\) at the
extremal energy, hence \(d\) there) and the Rayleigh quotient \(r\) (at
extremal \(r\), \(\mathrm{Var}=\varphi^{1/2}\langle\mathcal R,\widehat N\rangle_G\), the variance
equals the nonlinear correlation, an explicit upper bound on the
variance at extremal Rayleigh quotient against the floor \(v_{\rm exp}\)).
V37: the constant mode \(a_0\) and the momentum (at extremal Gaussian
momentum, \(\bar V\cdot\overline{N(V)}=-\frac12|\bar V|^2\), with the bound of V27
on \(\overline{N(V)}\) on a nearly constant core).  V38: the defect
\(\|Y\|_G^2\) itself and the cross term \(\langle N,Y\rangle_G\) (at extremal
defect the third identity of C2-V8 holds pointwise), with the
question whether the extremal defect configuration is constrained
to be far from self-similar by the explicit sides.  V39
consolidates and asks whether two extremal identities can be
combined into an exclusion; V40 audits.  No global regularity
claim is made.
\end{assessment}

% =============================================================================
\section{Integrated V34 conclusion and next acquisition order}
\label{sec:c3v34-conclusion}
% =============================================================================

\begin{theorem}[What V34 establishes]
\label{thm:c3v34-conclusion}
Relative to the two legacy cycles and V1--V33: the consolidated
defect block (Theorem~\ref{thm:c3v34-block}), the extremal
configurations lemma (Lemma~\ref{lem:c3v34-extremal}), the explicit
bounds on the extremal Dirichlet form
(Theorem~\ref{thm:c3v34-dirichlet}) and the direction.  The full
Navier--Stokes stability/global-regularity theorem remains unproved.
\end{theorem}

\begin{proof}
The cited results.
\end{proof}

\begin{programme}[V35 acquisition order]
\label{prog:c3v35}
\begin{enumerate}[label=\textup{(AC\arabic*)},leftmargin=4.8em]
\item The compulsory companion audit: read the current SM and YM
      notes and record their digests and decision.
\item Record the position after thirty-five versions; delivery.
\end{enumerate}
\end{programme}

\begin{assessment}[Position after V34]
\label{ass:c3v34-position}
The defect block is consolidated; extremal configurations give
explicit two-sided bounds on the Dirichlet form over the family; the
direction is the extremal ledger.  No full stability or
global-regularity claim is made.
\end{assessment}

% =============================================================================
\section*{Version V34 (third cycle) append-only record}
\addcontentsline{toc}{section}{Version V34 (third cycle) append-only record}
% =============================================================================

The complete V33 source of the third cycle, after the in-place
correction of Remark~\ref{rem:c3v34-correction}, was retained before
this continuation.  Its SHA--256 digest was
\begin{center}\scriptsize\ttfamily
10f6b3bd54b64d63856ffad91405dffbba4d0b04ba70318c461dc40a5d7c513d.
\end{center}
V34 consolidated the defect block and proved the extremal bounds.
No full regularity claim is made.

% =============================================================================
\section{V35: compulsory companion audit and the position after thirty-five versions}
\label{sec:c3v35-entry}
% =============================================================================

This is the thirty-fifth version of the third cycle.  The compulsory
review of the two companion programmes is performed first, and the
position after thirty-five versions is recorded with the plan for
the block V36--V40.  No global regularity claim is made.  The next
compulsory companion audit is V40.

% =============================================================================
\section{Compulsory V35 review of the current companion notes}
\label{sec:c3v35-companion-audit}
% =============================================================================

The current companion files in \texttt{latex\_notes} were inspected
read-only.  Both programmes have moved since the V30 audit: the SM
programme from V29 to V31 of its seventh series, the YM programme
from V6 to V8 of its sixth series.  The frozen checkpoints are
\begin{align}
 &\texttt{Renewal\_SM\_Einstein\_Continuum\_Research\_Notes\_V31.tex},
 \notag\\[-2pt]
 &\qquad
 \texttt{b7815fc841bb4768d6d84d5c01645249d15bb0c00bd851fe296b5e290fc0a775},
 \label{eq:c3v35-SM-checkpoint}\\
 &\texttt{Renewal\_Yang\_Mills\_Mass\_Gap\_Research\_Notes\_V8.tex},
 \notag\\[-2pt]
 &\qquad
 \texttt{26a0cbbbcc83b8261f65a7411f518a24442a015cb4462d6738bda7ed43ece6bb}.
 \label{eq:c3v35-YM-checkpoint}
\end{align}
(SM V31 has 12764 lines; YM V8 has 3339 lines.  As at V25 and V30,
the folder also holds a YM file of the preceding number, \texttt{V7},
byte-identical to \texttt{V8}; the audit records the newer name.)

\emph{SM V30--V31.}  The SM programme performed its own companion
audit, acquired intrinsic half-edge data on a spin manifold and a
global spin-mesh card, explained why a separate global gap proof is
necessary, and then proved a global Wilson commutator and IMS gap
through uniform spin-graph geometry, a quadratic lattice partition,
an exact collar replacement and a whole-operator gap.

\emph{YM V6--V8.}  The YM programme corrected a bottleneck-owner
statement into an affine-layer owner, derived an exact
layer-allocation law from complete spacings, studied quotient
histories on a common final vertex set and an exact nonlaminar
defect, then, after an archived freshness no-go changed its route,
computed an exact rooted-forest marginal, showed endpoint lifts
preserve it, closed its condition VAFPC, and isolated the component
defect that remains.

\begin{theorem}[V35 companion-audit decision]
\label{thm:c3v35-companion-decision}
Nothing is imported from SM V31 or YM V8.  Neither bears on the
third cycle's chapters or on the gates.
\end{theorem}

\begin{proof}
The SM results concern spectral gaps of lattice Dirac operators on
spin meshes; the YM results concern forest marginals in a cluster
expansion.  None is a Navier--Stokes statement, and none concerns
the Gaussian spectral theory, the class bounds, the ledger
identities or the extremal method of this series.
\end{proof}

% =============================================================================
\section{Position after thirty-five versions, and the block V36--V40}
\label{sec:c3v35-position}
% =============================================================================

\begin{assessment}[Position after thirty-five versions of the third cycle]
\label{ass:c3v35-position-thirty-five}
Six chapters are written: cluster energy and the band (V2--V10),
localization (V11--V14), the anchor (V16--V19), the explicit floor
(V21--V24), the variance chapter (V26--V29), and the defect block
(V31--V34).  The last two blocks made two compactness constants of
the homogeneous gate explicit or bounded them: the variance floor
explicit and exponentially small (V23, V28), the defect rate bounded
above by \(10\kappa^2\nu C\) (V33) and shown not to be explicitly boundable
below by quantifying (E9) in the class (N-C3-3, V31).  They also
produced two new explicit rows of the ledger, the mean-square
constraint (V32) and the two-sided extremal bounds on the Dirichlet
form (V34), and a method, extremal configurations, that turns mean
identities into pointwise ones.  The three negative records N-C3-1
(concentration gives no floor), N-C3-2 (the mean spectrum gives no
mean floor) and N-C3-3 (no quantitative NRS in the class) mark the
routes closed.  The block V36--V40 follows
Assessment~\ref{ass:c3v34-direction}: V36 the extremal energy and the
extremal Rayleigh quotient, V37 the extremal constant mode and
momentum, V38 the extremal defect and cross term, V39 consolidation
and the question of combining extremal identities, V40 the audit.
No global regularity claim is made.
\end{assessment}

% =============================================================================
\section{Integrated V35 conclusion and next acquisition order}
\label{sec:c3v35-conclusion}
% =============================================================================

\begin{theorem}[What V35 establishes]
\label{thm:c3v35-conclusion}
Relative to the two legacy cycles and V1--V34: the companion-audit
decision (Theorem~\ref{thm:c3v35-companion-decision}) and the position.
The full Navier--Stokes stability/global-regularity theorem remains
unproved.
\end{theorem}

\begin{proof}
The cited records.
\end{proof}

\begin{programme}[V36 acquisition order]
\label{prog:c3v36}
\begin{enumerate}[label=\textup{(AC\arabic*)},leftmargin=4.8em]
\item Extremal energy: at a maximizer or minimizer of \(\varphi\) over the
      family, \(\varphi'=0\) gives \(j=-8\mathcal L\); combine with the flux
      bound of C2-V36 (\(|j|\le\kappa(d+4\varphi)+K_P(d+3\varphi)^{1/2}\)) and
      \(d=8\mathcal L-2\varphi\) to bound \(\mathcal L\), hence \(d\), at the extremal
      energy; record the explicit statement and compare with
      Theorem~\ref{thm:c3v34-dirichlet}.
\item Extremal Rayleigh quotient: at an extremum of \(r\),
      \(\mathrm{Var}=\varphi^{1/2}\langle\mathcal R,\widehat N\rangle_G\) (C2-V31); bound the right
      side by the class and compare with the floor \(v_{\rm exp}\); record
      whether the extremal Rayleigh quotient has an explicit two-sided
      bound.
\item The next compulsory companion audit is V40.
\end{enumerate}
\end{programme}

\begin{assessment}[Position after V35]
\label{ass:c3v35-position}
The companion audit is recorded; the extremal-ledger block begins.
No full stability or global-regularity claim is made.
\end{assessment}

% =============================================================================
\section*{Version V35 (third cycle) append-only record}
\addcontentsline{toc}{section}{Version V35 (third cycle) append-only record}
% =============================================================================

The complete V34 source of the third cycle was retained before this
continuation.  Its SHA--256 digest was
\begin{center}\scriptsize\ttfamily
7654acc67cefe80caf7e83bfd14aefdc3c400618b0774e44795efda8029ba26b.
\end{center}
V35 performed the compulsory companion audit and recorded the
position.  No full regularity claim is made.

% =============================================================================
\section{V36: the extremal energy and the extremal Rayleigh quotient}
\label{sec:c3v36-entry}
% =============================================================================

Items (AC1) and (AC2) of Programme~\ref{prog:c3v36} are carried out.
(AC1): at a configuration of extremal Gaussian energy the flux
equals minus the dissipation plus twice the energy, exactly and
pointwise, so that the flux bound of the second cycle gives, for a
velocity constant below one half, an explicit bound on the
dissipation plus three times the energy at that configuration; in
particular the supremum of the Gaussian energy over the family is
explicitly bounded, which makes the constant \(C_\Phi\) explicit for
small velocity constants.  (AC2): at a configuration of extremal
Rayleigh quotient the variance equals the normalized correlation of
the Rayleigh residual with the nonlinearity, hence is at most the
squared nonlinearity over the energy; with the explicit sides this
is an explicit upper bound, quadratic in the velocity constant,
against the explicit floor, and for a small velocity constant the
extremal-Rayleigh profile is nearly a single Hermite level.  No
global regularity claim is made.  The next compulsory companion
audit is V40.

% =============================================================================
\section{The extremal energy}
\label{sec:c3v36-energy}
% =============================================================================

\begin{theorem}[Extremal Gaussian energy]
\label{thm:c3v36-energy}
Let \(W^*\) be a maximizer or a minimizer of \(\varphi\) over \(\mathcal F(x_0)\)
(Lemma~\ref{lem:c3v34-extremal}), and let \(d,j,\mathcal L\) be evaluated at
\(W^*\) at \(\sigma=0\).  Then
\begin{equation}
 \boxed{\quad
 j=-8\mathcal L=-(d+2\varphi)\ <0 ,
 \quad}
 \label{eq:c3v36-flux}
\end{equation}
and if \(\kappa<\frac12\), with \(m_\kappa=(2-4\kappa)/3\),
\begin{equation}
 d+3\varphi\ \le\ \frac{K_P^2}{m_\kappa^2}\qquad\text{at }W^* .
 \label{eq:c3v36-constraint}
\end{equation}
Consequently, for \(\kappa<\frac12\),
\begin{equation}
 \boxed{\quad
 \sup_{\mathcal F(x_0)}\varphi\ \le\ \frac{K_P^2}{3m_\kappa^2},
 \qquad
 d\ \le\ \frac{K_P^2}{m_\kappa^2}-3c_\Phi\ \text{ at a minimizer of }\varphi .
 \quad}
 \label{eq:c3v36-sup}
\end{equation}
\end{theorem}

\begin{proof}
\(\frac{\dd}{\dd\sigma}\varphi=2\langle V,Y\rangle_G=-4\mathcal L-\frac12j\) (the energy identity,
Proposition~\ref{prop:c3v31-gaussian}'s proof), and \(\varphi'=0\) at \(W^*\)
gives \(j=-8\mathcal L=-(d+2\varphi)\), negative since \(\varphi\ge c_\Phi>0\).  The
flux bound (Theorem C2-V36-flux-bound) then reads
\(d+2\varphi\le\kappa(d+4\varphi)+K_P(d+3\varphi)^{1/2}\), i.e.\
\((1-\kappa)d+(2-4\kappa)\varphi\le K_P(d+3\varphi)^{1/2}\); since \(m_\kappa\le1-\kappa\) and
\(3m_\kappa=2-4\kappa\), the left side is at least \(m_\kappa(d+3\varphi)\), and
dividing by \((d+3\varphi)^{1/2}>0\) gives \eqref{eq:c3v36-constraint}.  At a
maximizer, \(3\varphi\le d+3\varphi\) gives the first bound of
\eqref{eq:c3v36-sup}; at a minimizer, \(\varphi\ge c_\Phi\) gives the second.
\end{proof}

\begin{remark}[Comparison]
\label{rem:c3v36-comparison}
The second cycle proved \eqref{eq:c3v36-constraint} in the mean
(Theorem C2-V36-constraint); the extremal configuration gives it
pointwise at the extremal energy, and the consequence for the
supremum of \(\varphi\) is new: for \(\kappa<\frac12\) the class constant \(C_\Phi\)
of the family may be taken to be \(K_P^2/(3m_\kappa^2)\), explicit in \(\kappa\),
\(\nu\), the Calder\'on--Zygmund constant and \(K_6\), in place of the
\(L^6\) bound \(K_6^2\cdot8\pi\nu/3\); which of the two is smaller depends
on the constants.  The sign of the flux at extremal energy is the
pointwise form of the mean inward flux of the first cycle: when the
Gaussian energy is at its largest or its smallest over the family,
the flux is inward and equals eight times the Dirichlet form.
Theorem~\ref{thm:c3v34-dirichlet} bounds \(d\) on the whole family in
terms of \(\kappa\), \(c_\Phi\), \(C_\Phi\), \(C\); \eqref{eq:c3v36-sup} bounds it at the
minimal-energy configuration in terms of \(K_P\), \(m_\kappa\), \(c_\Phi\); the two
are independent rows.
\end{remark}

% =============================================================================
\section{The extremal Rayleigh quotient}
\label{sec:c3v36-rayleigh}
% =============================================================================

\begin{theorem}[Extremal Rayleigh quotient]
\label{thm:c3v36-rayleigh}
Let \(W^*\) be a maximizer or a minimizer of \(r\) over \(\mathcal F(x_0)\), and
let all quantities be evaluated at \(W^*\) at \(\sigma=0\).  Then
\begin{equation}
 \mathrm{Var}=\frac1\varphi\bigl\langle L_-V+rV,\,N(V)\bigr\rangle_G
 =\varphi^{1/2}\langle\mathcal R,\widehat N\rangle_G ,
 \label{eq:c3v36-var-identity}
\end{equation}
and therefore
\begin{equation}
 \boxed{\quad
 v_{\rm exp}\ \le\ \mathrm{Var}\ \le\ \frac{\|N(V)\|_G^2}{\varphi}\ \le\ \frac{\kappa^2}{c_\Phi}\Bigl(\frac d2+8\nu C\Bigr)\ \le\ \frac{10\kappa^2\nu C}{c_\Phi}
 \qquad\text{at }W^* .
 \quad}
 \label{eq:c3v36-var-bound}
\end{equation}
If moreover \(10\kappa^2\nu C/c_\Phi<\frac1{32}\), the profile of \(W^*\) at \(\sigma=0\)
is within \(\eta=160\kappa^2\nu C/c_\Phi\) in \(L^2(G)\) of a divergence-free
polynomial field of a single Hermite level \(k_1\le K_r\)
(Corollary~\ref{cor:c3v21-nearly-polynomial}).
\end{theorem}

\begin{proof}
The Rayleigh identity (Theorem C2-V31-identity)
\(\frac{\dd r}{\dd\sigma}=-2\mathrm{Var}+\frac2\varphi\langle L_-V+rV,N(V)\rangle_G\) with \(r'=0\)
gives \eqref{eq:c3v36-var-identity}.  By Cauchy--Schwarz and
\(\|L_-V+rV\|_G^2=\varphi\,\mathrm{Var}\) (equation \eqref{eq:c3v22-var-lambda} with
\(\lambda=r\), multiplied by \(\varphi\)),
\(\varphi\,\mathrm{Var}\le(\varphi\,\mathrm{Var})^{1/2}\|N\|_G\), so \(\mathrm{Var}\le\|N\|_G^2/\varphi\); then
Proposition~\ref{prop:c3v32-upper}, \(\varphi\ge c_\Phi\) and \(d\le4\nu C\).  The
lower bound is Theorem~\ref{thm:c3v23-floor}, and the last statement
is Corollary~\ref{cor:c3v21-nearly-polynomial}.
\end{proof}

\begin{assessment}[What the extremal Rayleigh quotient says]
\label{ass:c3v36-rayleigh}
The variance of the second cycle's Rayleigh chapter, bounded below
by compactness (C2-V34) and explicitly (V23), is now also bounded
above at the extremal-Rayleigh configuration, by the squared
velocity constant times the class enstrophy over the pinching.  For
a small velocity constant this forces the extremal-Rayleigh profile
to be nearly a single Hermite level, i.e.\ nearly a polynomial field
in the Gaussian sense, while V22--V23 showed that such a profile has
a Gaussian core of capped radius: the blow-up profile at its
extremal Rayleigh quotient is, for small \(\kappa\), a nearly polynomial
core of level at most \(K_r\).  The two bounds in
\eqref{eq:c3v36-var-bound} are compatible for every \(\kappa>0\), because
\(v_{\rm exp}\) is exponentially small; an exclusion would require
\(v_{\rm exp}>10\kappa^2\nu C/c_\Phi\), i.e.\ a velocity constant below
\((c_\Phi v_{\rm exp}/(10\nu C))^{1/2}\), a threshold far below any plausible
\(\varepsilon\)-regularity floor.  The identity
\eqref{eq:c3v36-var-identity} is the pointwise form of the mean push
budget of C2-V33 at one configuration.
\end{assessment}

% =============================================================================
\section{Integrated V36 conclusion and next acquisition order}
\label{sec:c3v36-conclusion}
% =============================================================================

\begin{theorem}[What V36 proves]
\label{thm:c3v36-conclusion}
Relative to the two legacy cycles and V1--V35: the extremal energy
(Theorem~\ref{thm:c3v36-energy}), with the explicit bound on the
supremum of the Gaussian energy for \(\kappa<\frac12\), and the extremal
Rayleigh quotient (Theorem~\ref{thm:c3v36-rayleigh}).  The full
Navier--Stokes stability/global-regularity theorem remains unproved.
\end{theorem}

\begin{proof}
The cited results.
\end{proof}

\begin{programme}[V37 acquisition order]
\label{prog:c3v37}
\begin{enumerate}[label=\textup{(AC\arabic*)},leftmargin=4.8em]
\item The extremal spectrum: at an extremum of \(a_k=\|\Pi_kV\|_G^2\) over
      the family, the modewise identity (C2-V40) gives \(n_k=-\mu_ka_k\)
      pointwise; with the modewise flux bound (C2-V41) deduce an
      explicit bound on \(\sup_{\mathcal F(x_0)}a_k\) for every \(k\), the
      pointwise counterpart of the mean spectrum; record whether the
      obstruction of V26 (the constant mode may carry all the energy
      the class allows) persists for suprema.
\item The extremal momentum: at an extremum of \(|\bar V|^2\),
      \(\bar V\cdot\overline{N(V)}=-\frac12|\bar V|^2\); combine with the bound of V27
      on \(\overline{N(V)}\) on a nearly constant core and with
      \eqref{eq:c3v36-var-bound}, and record the explicit statement.
\item The next compulsory companion audit is V40.
\end{enumerate}
\end{programme}

\begin{assessment}[Position after V36]
\label{ass:c3v36-position}
Extremal energy and extremal Rayleigh quotient give explicit
pointwise rows: an explicit \(C_\Phi\) for small \(\kappa\), and a two-sided
bound on the variance at the extremal Rayleigh quotient.  No full
stability or global-regularity claim is made.
\end{assessment}

% =============================================================================
\section*{Version V36 (third cycle) append-only record}
\addcontentsline{toc}{section}{Version V36 (third cycle) append-only record}
% =============================================================================

The complete V35 source of the third cycle was retained before this
continuation.  Its SHA--256 digest was
\begin{center}\scriptsize\ttfamily
24c26eb13ee9d6a439c1187479acc5d60ba995d32e36e27b609c49207c4eed21.
\end{center}
V36 proved the extremal-energy and extremal-Rayleigh statements.
No full regularity claim is made.

% =============================================================================
\section{V37: the extremal spectrum and the extremal momentum}
\label{sec:c3v37-entry}
% =============================================================================

Items (AC1) and (AC2) of Programme~\ref{prog:c3v37} are carried out.
(AC1): at a configuration of extremal energy in a given Hermite
level the modewise identity holds pointwise, and the modewise flux
bound then bounds the supremum of that level's energy over the
family explicitly, decaying like the inverse of the degree: the
pointwise spectrum that V21 recorded as unavailable from means is
available from extremal configurations.  The obstruction of V26
persists for the constant level, so no unconditional pointwise
variance floor of order one follows; a conditional one follows for
profiles whose dominant level is positive.  (AC2): at a
configuration of extremal Gaussian momentum the mean nonlinearity
is anti-parallel to the momentum with half its size, which
contradicts the smallness of the mean nonlinearity on a nearly
constant core below the threshold of V27; hence either the Gaussian
momentum never exceeds half the energy, or at its maximum the
variance is at least the explicit threshold.  No global regularity
claim is made.  The next compulsory companion audit is V40.

% =============================================================================
\section{The extremal spectrum}
\label{sec:c3v37-spectrum}
% =============================================================================

For \(k\ge0\) put
\begin{equation}
 \beta_k:=\kappa\Bigl(\frac k2\Bigr)^{1/2}C_\Phi^{1/2}+\kappa\,(4\nu C+3C_\Phi)^{1/2}+\frac{K_P}4(3+2k)^{1/2},
 \qquad
 B_k^{\sup}:=\frac{\beta_k^2}{\mu_k^2}=\frac{4\beta_k^2}{(1+k)^2} .
 \label{eq:c3v37-Bsup}
\end{equation}

\begin{theorem}[Extremal spectrum]
\label{thm:c3v37-spectrum}
For every \(k\ge0\),
\begin{equation}
 \boxed{\quad
 \sup_{\mathcal F(x_0)}a_k\ \le\ B_k^{\sup},
 \qquad\text{so that}\qquad
 \hat a_k\le\frac{B_k^{\sup}}{c_\Phi}\ \text{ for every profile at every }\sigma ,
 \quad}
 \label{eq:c3v37-sup}
\end{equation}
and \(B_k^{\sup}\le C_{\rm sp}^{\sup}/k\) for \(k\ge1\) with
\(C_{\rm sp}^{\sup}:=20[\kappa C_\Phi^{1/2}+\kappa(4\nu C+3C_\Phi)^{1/2}+\frac14K_P]^2\).
\end{theorem}

\begin{proof}
\(a_k=\|\Pi_kV\|_G^2\) is continuous on the family and \(C^1\) in \(\sigma\) with
\(a_k'=-2\mu_ka_k-2n_k\) (Theorem C2-V40-modewise).  At a maximizer \(W_k\)
of \(a_k\) (Lemma~\ref{lem:c3v34-extremal}), \(a_k'=0\) gives \(n_k=-\mu_ka_k\),
and the modewise flux bound (Theorem C2-V41-bound) gives
\(\mu_ka_k=|n_k|\le a_k^{1/2}[\kappa(k/2)^{1/2}\varphi^{1/2}+\kappa(d+3\varphi)^{1/2}+\frac{K_P}4(3+2k)^{1/2}]\le a_k^{1/2}\beta_k\)
at \(W_k\), using \(\varphi\le C_\Phi\) and \(d\le4\nu C\).  If \(a_k(W_k)>0\), divide
by \(a_k^{1/2}\) and square; the supremum is the value at \(W_k\).  The
normalized form uses \(\varphi\ge c_\Phi\).  For \(k\ge1\), \(\mu_k\ge k/2\) and
\((k/2)^{1/2}\le(5k)^{1/2}\), \(1\le(5k)^{1/2}\), \((3+2k)^{1/2}\le(5k)^{1/2}\) give
\(\beta_k\le(5k)^{1/2}[\kappa C_\Phi^{1/2}+\kappa(4\nu C+3C_\Phi)^{1/2}+\frac14K_P]\) and the \(1/k\)
form with the constant \(20[\dots]^2\).
\end{proof}

\begin{remark}[Against Firewall V21]
\label{rem:c3v37-pointwise}
Firewall~\ref{fw:c3v21-mean-not-pointwise} recorded that the mean
spectrum of C2-V41 gives no pointwise bound.  The extremal method
does: \eqref{eq:c3v37-sup} is a bound at every configuration, because
the supremum of each \(a_k\) is itself attained at a configuration
where the modewise identity holds with the derivative removed.  The
constants \(B_k^{\sup}\) differ from the \(B_k\) of V26 only in the
replacement of the means \(\langle\varphi\rangle\), \(\langle d+3\varphi\rangle\) by the class
bounds; with Theorem~\ref{thm:c3v36-energy}, for \(\kappa<\frac12\), \(C_\Phi\) may
itself be taken explicit.
\end{remark}

\begin{proposition}[Conditional pointwise floor for a positive dominant level]
\label{prop:c3v37-conditional}
For every profile at every \(\sigma\) whose dominant level \(k_1\) is
positive,
\begin{equation}
 \mathrm{Var}\ \ge\ \frac1{16}\Bigl(1-\frac{\max_{1\le k\le K_r}B_k^{\sup}}{c_\Phi}\Bigr),
 \label{eq:c3v37-conditional}
\end{equation}
a floor of order one whenever \(\max_{1\le k\le K_r}B_k^{\sup}<c_\Phi\).
For the constant level the analogous condition \(B_0^{\sup}<c_\Phi\) is
incompatible with the class for every \(\kappa>0\), exactly as in
Theorem~\ref{thm:c3v26-incompatible}.
\end{proposition}

\begin{proof}
Proposition~\ref{prop:c3v21-criteria}(ii) with
\(\hat a_{k_1}\le B_{k_1}^{\sup}/c_\Phi\) and \(1\le k_1\le K_r\).  For \(k=0\):
\(B_0^{\sup}=4[\kappa(4\nu C+3C_\Phi)^{1/2}+\frac{\sqrt3}4K_P]^2\ge4[\kappa(3c_\Phi)^{1/2}+\frac{\sqrt3}4K_P]^2\),
and the proof of Theorem~\ref{thm:c3v26-incompatible} applies
verbatim to this lower bound.
\end{proof}

\begin{firewall}[Item (AC1): what persists and what changes]
\label{fw:c3v37-obstruction}
The obstruction of V26 persists for the constant level: the sup
bound on the constant mode is not below the pinching, for the same
structural reason (the constant-mode bound and the energy bound are
the same inequality).  What changes is the case of a positive
dominant level: whether \(\max_{1\le k\le K_r}B_k^{\sup}<c_\Phi\) can hold is
an explicit inequality between \(c_\Phi\) and the class constants
which the series cannot decide (\(B_1^{\sup}\ge\frac5{16}K_P^2\) while
\(c_\Phi\le K_P^2/(3m_\kappa^2)\le\frac34K_P^2\) for small \(\kappa\): the condition
is neither implied nor excluded).  The dichotomy that results is:
at every \(\sigma\), either the dominant level is the constant one, or
\eqref{eq:c3v37-conditional} holds.  The constant-level case is the
nearly constant core of V27--V28, which the momentum decay treats
over windows.
\end{firewall}

% =============================================================================
\section{The extremal momentum}
\label{sec:c3v37-momentum}
% =============================================================================

\begin{theorem}[Extremal Gaussian momentum]
\label{thm:c3v37-momentum}
Let \(W_0\) be a maximizer of \(a_0=|G||\bar V|^2\) over \(\mathcal F(x_0)\), with
\(a_0(W_0)>0\).  Then at \(W_0\), \(\sigma=0\),
\begin{equation}
 \bar V\cdot\overline{N(V)}=-\tfrac12|\bar V|^2,\qquad\text{hence}\qquad
 |\overline{N(V)}|\ \ge\ \tfrac12|\bar V| .
 \label{eq:c3v37-momentum}
\end{equation}
Consequently one of the following holds:
\begin{enumerate}[label=\textup{(\alph*)},leftmargin=3.2em]
\item \(\sup_{\mathcal F(x_0)}a_0<\frac12\sup_{\mathcal F(x_0)}\varphi\): the Gaussian momentum
      never carries half the Gaussian energy;
\item at \(W_0\) the variance satisfies \(\mathrm{Var}\ge\min(\eta_{\rm c}/16,\frac1{32})\)
      with \(\eta_{\rm c}\) of \eqref{eq:c3v27-etac}.
\end{enumerate}
\end{theorem}

\begin{proof}
\(a_0'=-a_0-2n_0\) with \(n_0=\langle\Pi_0V,N\rangle_G=|G|\,\bar V\cdot\overline{N(V)}\)
(Proposition C2-V40-low-modes); \(a_0'=0\) at \(W_0\) gives
\eqref{eq:c3v37-momentum}, and Cauchy--Schwarz the norm inequality.
Suppose (b) fails: \(\mathrm{Var}<\frac1{32}\) and \(\eta=16\,\mathrm{Var}<\eta_{\rm c}\) at
\(W_0\).  If the dominant level at \(W_0\) is \(0\), Corollary
\ref{cor:c3v27-decay} gives \(|\overline{N(V)}|\le\frac14|\bar V|\), contradicting
\eqref{eq:c3v37-momentum} (as \(\bar V\ne0\)).  So the dominant level is
positive, and then \(\hat a_0\le\sum_{k\ne k_1}\hat a_k\le\eta<\frac12\)
(Corollary~\ref{cor:c3v21-nearly-polynomial}), i.e.\
\(a_0(W_0)<\frac12\varphi(W_0)\le\frac12\sup\varphi\), which is (a).
\end{proof}

\begin{assessment}[Reading]
\label{ass:c3v37-momentum}
The momentum identity, neutral in the mean (V18), acquires a sign at
the extremal-momentum configuration: there the nonlinear mean
opposes the momentum with at least half its size.  On a nearly
constant core the nonlinear mean is small (V27), so the profile of
maximal momentum is either not a nearly constant core (its momentum
is below half its energy, alternative (a)) or has a variance above
the explicit threshold (alternative (b)).  Either way the nearly
constant core of V27--V28, the one configuration the variance
floors could not exclude, is not the configuration of maximal
Gaussian momentum unless its variance is at least \(\eta_{\rm c}/16\).
This does not raise the floor at other configurations.
\end{assessment}

% =============================================================================
\section{Integrated V37 conclusion and next acquisition order}
\label{sec:c3v37-conclusion}
% =============================================================================

\begin{theorem}[What V37 proves]
\label{thm:c3v37-conclusion}
Relative to the two legacy cycles and V1--V36: the extremal spectrum
(Theorem~\ref{thm:c3v37-spectrum}), the conditional pointwise floor
(Proposition~\ref{prop:c3v37-conditional}) and the extremal momentum
(Theorem~\ref{thm:c3v37-momentum}).  The full Navier--Stokes
stability/global-regularity theorem remains unproved.
\end{theorem}

\begin{proof}
The cited results.
\end{proof}

\begin{programme}[V38 acquisition order]
\label{prog:c3v38}
\begin{enumerate}[label=\textup{(AC\arabic*)},leftmargin=4.8em]
\item The extremal defect: at an extremum of \(\mathfrak y=\|Y\|_G^2\) over
      the family, the third identity of C2-V7--V8
      (\(\frac{\dd}{\dd\sigma}\frac12\|Y\|_G^2=\langle Y,L_-Y\rangle_G-\langle Y,N'(V)Y\rangle_G\))
      holds with zero left side; with \(\langle Y,L_-Y\rangle_G\le-\frac12\|Y\|_G^2\)
      (C2-V8-spectrum(ii)) deduce \(\langle Y,N'(V)Y\rangle_G\le-\frac12\|Y\|_G^2\) at
      the maximal defect and \(\ge\) at the minimal defect; bound
      \(|\langle Y,N'(V)Y\rangle_G|\) by the class and record the explicit
      two-sided consequence for \(\sup\|Y\|_G^2\) and \(\inf\|Y\|_G^2\).
\item The extremal cross term \(\mathfrak c=\langle N,Y\rangle_G\): its derivative
      from the ledger, if available in closed form (C2-V8), and the
      identity at its extremum; otherwise record that the cross term
      is not a \(C^1\) observable of the ledger.
\item The next compulsory companion audit is V40.
\end{enumerate}
\end{programme}

\begin{assessment}[Position after V37]
\label{ass:c3v37-position}
The extremal spectrum is explicit and pointwise; the extremal
momentum gives a sign to the momentum identity; the nearly constant
core is not the maximal-momentum configuration unless its variance
is above threshold.  No full stability or global-regularity claim
is made.
\end{assessment}

% =============================================================================
\section*{Version V37 (third cycle) append-only record}
\addcontentsline{toc}{section}{Version V37 (third cycle) append-only record}
% =============================================================================

The complete V36 source of the third cycle was retained before this
continuation.  Its SHA--256 digest was
\begin{center}\scriptsize\ttfamily
bc3658c09ae92c4482ba0724e50714bef749afc7c3c46bb157c65c50ffb2bc6f.
\end{center}
V37 proved the extremal spectrum and the extremal momentum.  No
full regularity claim is made.

% =============================================================================
\section{V38: the extremal defect and the extremal cross term}
\label{sec:c3v38-entry}
% =============================================================================

Items (AC1) and (AC2) of Programme~\ref{prog:c3v38} are carried out.
(AC1): at a configuration of extremal defect the Gaussian energy
identity of the defect holds with zero left side, so the linearized
nonlinearity is exactly minus the Dirichlet form of the defect
there, negative by at least half the defect: the third sign of the
second cycle, a mean statement, holds pointwise at the extremal
defect.  The explicit evaluation of this sign is carried out for the
advective parts of the linearized nonlinearity and is blocked by
the linearized pressure, whose Gaussian pairing with the defect
requires unweighted norms of the defect that the class does not
provide; the obstruction is the one of V27, V31 and V33 and is
recorded.  (AC2): at an extremum of the cross term the pairing of
the nonlinearity with the log-scale acceleration of the defect
equals minus the linearized quadratic form; the identity is
recorded without explicit content.  No global regularity claim is
made.  The next compulsory companion audit is V40.

% =============================================================================
\section{The extremal defect}
\label{sec:c3v38-defect}
% =============================================================================

On the family every element is smooth with uniform derivative
bounds of all orders (E5, C2-V4), so \(\mathfrak y:=\|Y\|_G^2\) is \(C^1\) in \(\sigma\)
and the identity of Theorem C2-V7-defect-energy holds for every \(\sigma\):
\begin{equation}
 \frac12\frac{\dd}{\dd\sigma}\|Y\|_G^2=\langle Y,L_-Y\rangle_G-\langle Y,N'(V)Y\rangle_G,
 \qquad
 \langle Y,L_-Y\rangle_G=-\nu\|\nabla Y\|_G^2-\frac12\|Y\|_G^2 ,
 \label{eq:c3v38-identity}
\end{equation}
with \(N'(V)Y=(Y\cdot\nabla)V+(V\cdot\nabla)Y+\nabla P'\), \(P'=\partial_\sigma P_V\), and
\begin{equation}
 \langle Y,N'(V)Y\rangle_G=\int Y\cdot(Y\cdot\nabla)V\,G+\frac1{4\nu}\int|Y|^2(z\cdot V)G+\frac1{2\nu}\int P'(z\cdot Y)G .
 \label{eq:c3v38-N'}
\end{equation}

\begin{theorem}[Extremal defect]
\label{thm:c3v38-defect}
Let \(W_Y\) be a maximizer or a minimizer of \(\mathfrak y=\|Y\|_G^2\) over
\(\mathcal F(x_0)\).  Then at \(W_Y\), \(\sigma=0\),
\begin{equation}
 \boxed{\quad
 \langle Y,N'(V)Y\rangle_G=-\nu\|\nabla Y\|_G^2-\frac12\|Y\|_G^2\ \le\ -\frac12\|Y\|_G^2 .
 \quad}
 \label{eq:c3v38-sign}
\end{equation}
In particular, at the configuration of maximal defect,
\(\langle Y,N'(V)Y\rangle_G\le-\frac12\sup_{\mathcal F(x_0)}\|Y\|_G^2\).
\end{theorem}

\begin{proof}
Lemma~\ref{lem:c3v34-extremal} with \(F=\mathfrak y\) and
\eqref{eq:c3v38-identity}.
\end{proof}

\begin{remark}[The third sign, pointwise]
\label{rem:c3v38-third-sign}
The second cycle's third mean sign,
\(\int\langle Y,N'(V)Y\rangle_G\dd\mathfrak m\le-\frac12\underline\delta\) (C2-V12), is the
average of \eqref{eq:c3v38-identity}; Theorem~\ref{thm:c3v38-defect} is
its pointwise form at one configuration, with the exact value.  It
reads: where the profile changes fastest in log-time, the
linearized nonlinearity damps the change at least as strongly as
the self-similar scaling does.
\end{remark}

% =============================================================================
\section{Explicit evaluation, and where it stops}
\label{sec:c3v38-explicit}
% =============================================================================

\begin{proposition}[The advective parts of the linearized form]
\label{prop:c3v38-advective}
For every profile of the family and every \(\sigma\), with
\(\kappa=C_\infty\nu^{-1/2}\) and \(C_\infty'\) the class bound on \(|\nabla V|\),
\begin{equation}
 \Bigl|\int Y\cdot(Y\cdot\nabla)V\,G\Bigr|\le C_\infty'\|Y\|_G^2,\qquad
 \Bigl|\frac1{4\nu}\int|Y|^2(z\cdot V)G\Bigr|\le\frac{\sqrt3}2\kappa\|Y\|_G^2+\frac\nu2\|\nabla Y\|_G^2+\frac{\kappa^2}2\|Y\|_G^2 ,
 \label{eq:c3v38-advective}
\end{equation}
and the pressure part equals \(\langle Y,\nabla P'\rangle_G\).  Consequently, at
a maximizer \(W_Y\) of the defect,
\begin{equation}
 \frac\nu2\|\nabla Y\|_G^2+\Bigl(\frac12-C_\infty'-\frac{\sqrt3}2\kappa-\frac{\kappa^2}2\Bigr)\|Y\|_G^2\ \le\ \bigl|\langle Y,\nabla P'\rangle_G\bigr|
 \qquad\text{at }W_Y .
 \label{eq:c3v38-at-max}
\end{equation}
\end{proposition}

\begin{proof}
The first bound is \(|\nabla V|\le C_\infty'\).  For the second,
\(|z\cdot V|\le C_\infty|z|\) and Cauchy--Schwarz give
\(\frac{C_\infty}{4\nu}\|Y\|_G\||z|Y\|_G\); the Gaussian Hardy inequality
(Lemma C2-V36-hardy, valid for the divergence-free \(Y\)) gives
\(\||z|Y\|_G\le(12\nu)^{1/2}\|Y\|_G+4\nu\|\nabla Y\|_G\), and
\(\frac{C_\infty}{4\nu}(12\nu)^{1/2}=\frac{\sqrt3}2\kappa\), while
\(C_\infty\|Y\|_G\|\nabla Y\|_G\le\frac\nu2\|\nabla Y\|_G^2+\frac{C_\infty^2}{2\nu}\|Y\|_G^2\).  The
pressure part: \((z\cdot Y)G=-2\nu\,Y\cdot\nabla G\), so
\(\frac1{2\nu}\int P'(z\cdot Y)G=-\int P'\,Y\cdot\nabla G=\int Y\cdot\nabla P'\,G\) by
\(\operatorname{div}Y=0\).  At \(W_Y\), \eqref{eq:c3v38-sign} gives
\(\nu\|\nabla Y\|^2+\frac12\|Y\|^2=-\langle Y,N'(V)Y\rangle_G\), bounded by the three
parts.
\end{proof}

\begin{firewall}[Item (AC1): the linearized pressure blocks the explicit evaluation]
\label{fw:c3v38-pressure}
If the pressure pairing were absent, \eqref{eq:c3v38-at-max} would
say that for \(C_\infty'+\frac{\sqrt3}2\kappa+\frac{\kappa^2}2<\frac12\) the maximal defect
vanishes, hence the defect vanishes on the family, hence the family
is self-similar, which (E9) excludes: an explicit exclusion of
homogeneous points with small derivative constants.  The pressure
pairing is not absent.  \(P'\) is the pressure of the linearized
nonlinearity, \(-\Delta P'=2\partial_iV_j\partial_jY_i\), and
\(\langle Y,\nabla P'\rangle_G\) is bounded, by the argument of
Lemma~\ref{lem:c3v32-pressure}, by \(\|Y\|_G\|\nabla P'\|_{L^2(\R^3)}\le\|Y\|_G(2C_\infty\|\nabla Y\|_{L^2(\R^3)}+2C_\infty'\|Y\|_{L^2(\R^3)})\),
which involves the unweighted norms of \(Y=\partial_\sigma V\); these are
not finite in general (\(\partial_\sigma V\) contains \(\frac12z\cdot\nabla V\), which
is not in \(L^2(\R^3)\) under the class bounds).  A Gaussian bound
would need the Gaussian-weighted Riesz transform, which is not
bounded, or a split of the source at a radius with the exponential
price of V27 and V33.  The explicit evaluation therefore stops at
\eqref{eq:c3v38-at-max}, with the pressure pairing as a named
non-explicit term; the obstruction is the fourth appearance of the
same one (V27, V31, V33, V38), and it is now recorded as a standing
firewall of the third cycle: \emph{the nonlocal pressure and the
Gaussian weight do not commute, and every route that needs the
pressure in Gaussian form pays an exponential in the radius.}
\end{firewall}

% =============================================================================
\section{The extremal cross term}
\label{sec:c3v38-cross}
% =============================================================================

\begin{proposition}[Extremal cross term]
\label{prop:c3v38-cross}
\(\mathfrak c:=\langle N(V),Y\rangle_G\) is \(C^1\) in \(\sigma\) on the family with
\(\mathfrak c'=\langle N'(V)Y,Y\rangle_G+\langle N(V),\partial_\sigma Y\rangle_G\) (Theorem C2-V8-third-sign),
and \(\partial_\sigma Y=L_-Y-N'(V)Y\).  At a maximizer or minimizer \(W_c\) of
\(\mathfrak c\),
\begin{equation}
 \langle N(V),\partial_\sigma Y\rangle_G=-\langle Y,N'(V)Y\rangle_G,
 \qquad\text{i.e.}\qquad
 \langle N,L_-Y\rangle_G=\langle N,N'(V)Y\rangle_G-\langle Y,N'(V)Y\rangle_G .
 \label{eq:c3v38-cross}
\end{equation}
\end{proposition}

\begin{proof}
Lemma~\ref{lem:c3v34-extremal} and the linearized equation
(Proposition C2-V7-defect-equation).
\end{proof}

\begin{firewall}[Item (AC2): no explicit content]
\label{fw:c3v38-cross}
Both sides of \eqref{eq:c3v38-cross} involve \(N'(V)Y\) with its pressure
part, and \(\langle N,L_-Y\rangle_G\) involves second derivatives of \(Y\); no
side has a sign or an explicit bound in the class.  The identity is
recorded as a row of the extremal ledger without consequence.
\end{firewall}

% =============================================================================
\section{Integrated V38 conclusion and next acquisition order}
\label{sec:c3v38-conclusion}
% =============================================================================

\begin{theorem}[What V38 proves]
\label{thm:c3v38-conclusion}
Relative to the two legacy cycles and V1--V37: the extremal defect
(Theorem~\ref{thm:c3v38-defect}), the advective evaluation
(Proposition~\ref{prop:c3v38-advective}) with the standing firewall on
the pressure, and the extremal cross term
(Proposition~\ref{prop:c3v38-cross}).  The full Navier--Stokes
stability/global-regularity theorem remains unproved.
\end{theorem}

\begin{proof}
The cited results.
\end{proof}

\begin{programme}[V39 acquisition order]
\label{prog:c3v39}
\begin{enumerate}[label=\textup{(AC\arabic*)},leftmargin=4.8em]
\item Consolidate V34--V38 as the extremal ledger: the lemma, and
      the rows at extremal Dirichlet form, energy, Rayleigh quotient,
      spectrum, momentum, defect, cross term, with their explicit
      sides and the standing pressure firewall.
\item Combining extremal identities: record that extrema of
      different observables are attained at different configurations
      in general, so that two rows combine only through a joint
      observable; examine the one-parameter family \(F_\lambda=\varphi-\lambda\mathcal L\)
      and \(F=\mathcal L-\frac{r_0}2\varphi\), and state what, if anything, their
      extremal rows add.  Fix the direction for V41--V45.
\item V40 compulsory companion audit and delivery.
\end{enumerate}
\end{programme}

\begin{assessment}[Position after V38]
\label{ass:c3v38-position}
The extremal defect gives the third sign pointwise; its explicit
evaluation stops at the linearized pressure, now a standing
firewall.  No full stability or global-regularity claim is made.
\end{assessment}

% =============================================================================
\section*{Version V38 (third cycle) append-only record}
\addcontentsline{toc}{section}{Version V38 (third cycle) append-only record}
% =============================================================================

The complete V37 source of the third cycle was retained before this
continuation.  Its SHA--256 digest was
\begin{center}\scriptsize\ttfamily
11a3b39df304e6406a8ab92f0c20c3186f60110baea8d354b655175d97a4c962.
\end{center}
V38 proved the extremal-defect and extremal-cross-term statements
and recorded the standing pressure firewall.  No full regularity
claim is made.

% =============================================================================
\section{V39: the extremal ledger consolidated, and the angle form of the gate}
\label{sec:c3v39-entry}
% =============================================================================

Items (AC1) and (AC2) of Programme~\ref{prog:c3v39} are carried out.
The extremal ledger V34--V38 is consolidated: seven rows, each a
mean identity of the second cycle turned pointwise at one extremal
configuration, with explicit sides where the class permits and the
standing pressure firewall where it does not.  Combining rows is
examined and found to add nothing beyond the rows of joint
observables.  The gate G1 is then reformulated with the mean-square
identity: the required anti-alignment of the nonlinearity with the
defect is an angle statement, and no comparison of norms can decide
it, because the nonlinearity exceeds the defect in mean square by
at least a quarter of the pinching.  The direction for V41--V45 is
the angle itself, in its Lamb--Bernoulli form.  No global
regularity claim is made.  The next compulsory companion audit is
V40.

% =============================================================================
\section{The extremal ledger}
\label{sec:c3v39-ledger}
% =============================================================================

\begin{theorem}[The extremal ledger, V34--V38]
\label{thm:c3v39-ledger}
At a homogeneous type-I point, each of the following observables
attains its supremum and infimum on \(\mathcal F(x_0)\), and at an
attaining configuration:
\begin{enumerate}[label=\textup{(\arabic*)},leftmargin=3.2em]
\item (Dirichlet form.)  \(\|N\|_G^2=\|L_-V\|_G^2+\|Y\|_G^2\);
      \(\mathcal L_-\le\mathcal L\le\mathcal L_+\) explicit (V34).
\item (Energy.)  \(j=-(d+2\varphi)\); for \(\kappa<\frac12\), \(d+3\varphi\le K_P^2/m_\kappa^2\)
      and \(\sup\varphi\le K_P^2/(3m_\kappa^2)\) (V36).
\item (Rayleigh quotient.)  \(\mathrm{Var}=\varphi^{-1}\langle L_-V+rV,N\rangle_G\le\|N\|_G^2/\varphi\le10\kappa^2\nu C/c_\Phi\) (V36).
\item (Level energies.)  \(n_k=-\mu_ka_k\); \(\sup a_k\le B_k^{\sup}\), pointwise
      spectrum \(\hat a_k\le B_k^{\sup}/c_\Phi\) (V37).
\item (Momentum.)  \(\bar V\cdot\overline{N(V)}=-\frac12|\bar V|^2\); either
      \(\sup a_0<\frac12\sup\varphi\) or \(\mathrm{Var}\ge\min(\eta_{\rm c}/16,\frac1{32})\) at the
      maximal momentum (V37).
\item (Defect.)  \(\langle Y,N'(V)Y\rangle_G=-\nu\|\nabla Y\|_G^2-\frac12\|Y\|_G^2\); the
      advective parts are explicit, the linearized pressure pairing
      is not (standing firewall) (V38).
\item (Cross term.)  \(\langle N,\partial_\sigma Y\rangle_G=-\langle Y,N'(V)Y\rangle_G\), without
      explicit content (V38).
\end{enumerate}
\end{theorem}

\begin{proof}
Lemma~\ref{lem:c3v34-extremal}; Theorems~\ref{thm:c3v34-dirichlet},
\ref{thm:c3v36-energy}, \ref{thm:c3v36-rayleigh}, \ref{thm:c3v37-spectrum},
\ref{thm:c3v37-momentum}, \ref{thm:c3v38-defect};
Proposition~\ref{prop:c3v38-cross}.
\end{proof}

\begin{proposition}[Combining rows]
\label{prop:c3v39-combining}
For \(C^1\) observables \(F_1,F_2\) and real \(\alpha,\beta\), the observable
\(\alpha F_1+\beta F_2\) attains its extrema, and there
\(\alpha F_1'+\beta F_2'=0\).  In particular, at an extremum of \(F_\lambda:=\varphi-\lambda\mathcal L\),
\begin{equation}
 -4\mathcal L-\frac12j\ =\ \frac\lambda2\bigl(\|N\|_G^2-\|L_-V\|_G^2-\|Y\|_G^2\bigr),
 \label{eq:c3v39-Flambda}
\end{equation}
and at an extremum of \(\mathcal L-\frac{r_0}2\varphi\) (\(r_0\) fixed),
\(\|N\|_G^2-\|L_-V\|_G^2-\|Y\|_G^2=-r_0(4\mathcal L+\frac12j)\).
\end{proposition}

\begin{proof}
Lemma~\ref{lem:c3v34-extremal} for the combined observable, with
\(\varphi'=-4\mathcal L-\frac12j\) and \(2\mathcal L'=\|N\|^2-\|L_-V\|^2-\|Y\|^2\).
\end{proof}

\begin{assessment}[Item (AC2): what combining adds]
\label{ass:c3v39-combining}
Extrema of different observables are attained at different
configurations in general, so rows (1)--(7) cannot be conjoined at
one configuration; the rows of joint observables
(Proposition~\ref{prop:c3v39-combining}) interpolate between them,
and each has the same structure as its parents: an exact relation
between an explicit quantity (\(\|N\|_G^2\), bounded above; \(\|L_-V\|_G^2\),
bounded below; \(j\), bounded by the flux bound) and non-explicit
ones (\(\|Y\|_G^2\), \(\varphi\) and \(\mathcal L\) at the configuration).  No
choice of \(\lambda\) or \(r_0\) removes the non-explicit quantities, and
none produces an inequality between explicit constants alone; the
extremal ledger produces explicit two-sided bounds on the ranges of
the observables over the family (V34, V36, V37), which is what it
can do, and no exclusion, because every range bound has the
pinching \(c_\Phi\) or the defect rate \(\underline\delta\) on its non-explicit
side.  The method is recorded as closed at this level; it will be
reused whenever a new observable with an explicit derivative
appears.
\end{assessment}

% =============================================================================
\section{The angle form of gate G1}
\label{sec:c3v39-angle}
% =============================================================================

\begin{theorem}[Gate G1 is a statement about an angle]
\label{thm:c3v39-angle}
Let \(x_0\) be a homogeneous type-I point and \(\mathfrak m\) an invariant
probability measure on \(\mathcal F(x_0)\).  Then
\begin{equation}
 \langle\|N\|_G^2\rangle_{\mathfrak m}-\langle\|Y\|_G^2\rangle_{\mathfrak m}=\langle\|L_-V\|_G^2\rangle_{\mathfrak m}\ \ge\ \Bigl(\frac14+v_{\rm exp}\Bigr)c_\Phi\ >0 ,
 \label{eq:c3v39-norms}
\end{equation}
and the mean identity \(\langle\langle N,Y\rangle_G\rangle=-\langle\|Y\|_G^2\rangle\) reads, with
\(\cos\theta:=\langle N,Y\rangle_G/(\|N\|_G\|Y\|_G)\) wherever both are nonzero,
\begin{equation}
 \boxed{\quad
 \bigl\langle\|N\|_G\|Y\|_G\cos\theta\bigr\rangle_{\mathfrak m}=-\langle\|Y\|_G^2\rangle_{\mathfrak m},
 \qquad
 \cos\theta=-\frac{\|Y\|_G^2+\mathcal L'}{\|N\|_G\|Y\|_G}\ \text{ pointwise}.
 \quad}
 \label{eq:c3v39-angle}
\end{equation}
The missing statement of gate G1 (Assessment~\ref{ass:c3v1-gates}(G1)),
\(\int\langle N,Y\rangle\,\dd\mathfrak m\ge-\theta_0\int\|Y\|^2\dd\mathfrak m\) with \(\theta_0<1\), is
therefore the statement that \eqref{eq:c3v39-angle} is impossible:
that the defect cannot be, on average and weighted by
\(\|N\|_G\|Y\|_G\), anti-aligned with the nonlinearity to the extent
\(-\langle\|Y\|_G^2\rangle\).
\end{theorem}

\begin{proof}
Theorem~\ref{thm:c3v32-mean-square} and Proposition~\ref{prop:c3v32-lower}
for \eqref{eq:c3v39-norms}; the mean identity of C2-V12 and the
pointwise form \(\langle N,Y\rangle_G=-\|Y\|_G^2-\mathcal L'\) (Proposition
\ref{prop:c3v33-pointwise}) for \eqref{eq:c3v39-angle}.
\end{proof}

\begin{firewall}[No norm comparison can close G1]
\label{fw:c3v39-no-norms}
By Cauchy--Schwarz, \(|\langle\langle N,Y\rangle_G\rangle|\le\langle\|N\|_G^2\rangle^{1/2}\langle\|Y\|_G^2\rangle^{1/2}\),
and by \eqref{eq:c3v39-norms} the right side exceeds \(\langle\|Y\|_G^2\rangle\)
strictly; so the exact anti-alignment demanded by
\eqref{eq:c3v39-angle} is compatible with every bound on the norms of
\(N\) and \(Y\), and corresponds to a mean cosine of size
\((\langle\|Y\|^2\rangle/\langle\|N\|^2\rangle)^{1/2}<1\), bounded away from \(-1\).  Any
proof of G1 must use the direction of \(Y\) relative to \(N\), not
their sizes.  This closes, for the third cycle, the family of
attempts that bound \(\|N\|_G\) or \(\|Y\|_G\) (the mean-square constraint
of V32, the window bounds of V33, the extremal bounds of V34--V38):
they constrain the ranges of the observables and cannot decide the
gate.  The direction of \(Y\) relative to \(N\) is what the
Lamb--Bernoulli form of the cross term describes:
\(\langle N,Y\rangle_G=\int\det(\Omega,V,Y)G+\frac1{2\nu}\int\Pi(z\cdot Y)G\)
(C2-V2, V1(S1)), the Gaussian mean of the defect against the Lamb
vector \(\Omega\times V\) plus the Bernoulli work against the radial
defect.
\end{firewall}

\begin{assessment}[Direction for V41--V45: the Lamb--defect angle]
\label{ass:c3v39-direction}
The block V41--V45 studies the angle directly.  V41: explicit bounds
on the two Lamb--Bernoulli pieces at every configuration
(\(|\int\det(\Omega,V,Y)G|\le C_\infty\mathcal E^{1/2}\|Y\|_G\) with \(\mathcal E\le\frac2\nu\mathcal L\) by the
closure identity; \(|\frac1{2\nu}\int(\Pi-c)(z\cdot Y)G|\le\frac1{2\nu}\||z|(\Pi-c)\|_G\|Y\|_G\)
with \(\||z|P_V\|_G\) explicitly bounded through the \(L^3\) bound and the
Gaussian moment \(\int|z|^6G^3\), since the pressure enters without a
derivative), hence an explicit \(M_{NY}\) with \(|\langle N,Y\rangle_G|\le M_{NY}\|Y\|_G\)
and, at extremal Dirichlet form, \(\|Y\|_G\le M_{NY}\).  V42: the pointwise
structure \(\det(\Omega,V,Y)=Y\cdot(\Omega\times V)\) on the anchor's neighbourhood
(V13, V17): the anti-alignment is the defect pointing against the
Lamb vector, Bernoulli-corrected, and the question is which
configurations of \(\Omega\), \(V\), \(Y\) on at most \(N_C\) cells realize a
Gaussian mean of \(-\|Y\|_G^2\); the stationary-band and cut-off
examples of V6 and V22 are tested against it.  V43: the extremal
Dirichlet form row (1) in Lamb--Bernoulli form, the one
configuration at which the anti-alignment is exact and pointwise.
V44 consolidates; V45 audits.  No global regularity claim is made.
\end{assessment}

% =============================================================================
\section{Integrated V39 conclusion and next acquisition order}
\label{sec:c3v39-conclusion}
% =============================================================================

\begin{theorem}[What V39 establishes]
\label{thm:c3v39-conclusion}
Relative to the two legacy cycles and V1--V38: the consolidated
extremal ledger (Theorem~\ref{thm:c3v39-ledger}), the combination
rows (Proposition~\ref{prop:c3v39-combining}), the angle form of the
gate (Theorem~\ref{thm:c3v39-angle}) and the firewall on norm
comparisons.  The full Navier--Stokes stability/global-regularity
theorem remains unproved.
\end{theorem}

\begin{proof}
The cited results.
\end{proof}

\begin{programme}[V40 acquisition order]
\label{prog:c3v40}
\begin{enumerate}[label=\textup{(AC\arabic*)},leftmargin=4.8em]
\item The compulsory companion audit: read the current SM and YM
      notes and record their digests and decision.
\item Record the position after forty versions; delivery.
\end{enumerate}
\end{programme}

\begin{assessment}[Position after V39]
\label{ass:c3v39-position}
The extremal ledger is consolidated and closed at its level; gate
G1 is an angle statement that norms cannot decide; the direction is
the Lamb--defect angle.  No full stability or global-regularity
claim is made.
\end{assessment}

% =============================================================================
\section*{Version V39 (third cycle) append-only record}
\addcontentsline{toc}{section}{Version V39 (third cycle) append-only record}
% =============================================================================

The complete V38 source of the third cycle was retained before this
continuation.  Its SHA--256 digest was
\begin{center}\scriptsize\ttfamily
59932a358fc764eb14f71cd6bf6d01a25268b114977677e22c170bfb066f6178.
\end{center}
V39 consolidated the extremal ledger and recorded the angle form of
the gate.  No full regularity claim is made.

% =============================================================================
\section{V40: compulsory companion audit and the position after forty versions}
\label{sec:c3v40-entry}
% =============================================================================

This is the fortieth version of the third cycle.  The compulsory
review of the two companion programmes is performed first, and the
position after forty versions is recorded with the plan for the
block V41--V45.  No global regularity claim is made.  The next
compulsory companion audit is V45.

% =============================================================================
\section{Compulsory V40 review of the current companion notes}
\label{sec:c3v40-companion-audit}
% =============================================================================

The current companion files in \texttt{latex\_notes} were inspected
read-only.  Both programmes have moved since the V35 audit: the SM
programme from V31 to V32 of its seventh series, the YM programme
from V8 to V9 of its sixth series.  The frozen checkpoints are
\begin{align}
 &\texttt{Renewal\_SM\_Einstein\_Continuum\_Research\_Notes\_V32.tex},
 \notag\\[-2pt]
 &\qquad
 \texttt{5509dd0ae6f3fb30997c68fabbf0af429f2e775a356420bd67d73f75932c2b5f},
 \label{eq:c3v40-SM-checkpoint}\\
 &\texttt{Renewal\_Yang\_Mills\_Mass\_Gap\_Research\_Notes\_V9.tex},
 \notag\\[-2pt]
 &\qquad
 \texttt{df86458225b0a5c6b4138066784dc82eaf837b0d03e192d7ddcab342565f3d5a}.
 \label{eq:c3v40-YM-checkpoint}
\end{align}
(SM V32 has 13098 lines; YM V9 has 4214 lines.  The YM file was
rewritten between the first inspection for this audit and the
assembly of this version; the digest recorded is that of the final
inspection, and the file of the preceding number, present at the
first inspection, had been removed.)

\emph{SM V32.}  The SM programme recorded the failure of a partition
interpolation and a reduction: what an algebra-preserving
interpolation can be, a frozen anisotropic coordinate gap, and one
mesh with a finite catalogue of star types.

\emph{YM V8--V9.}  The YM programme resolved owned components
(owned-edge colourings resolve to an incidence tree, the gate
coefficient contains an exact component-tree factor, the post-norm
grouping still has a superexponential count, a corrected form of
its condition SOCC), and then calibrated a Duhamel coefficient: an
exact distinct-type Duhamel coefficient, all forgotten groupings as
one forest, the lift of the grouping forest to label edges, an exact
tree--forest overlay and its lift allocation.

\begin{theorem}[V40 companion-audit decision]
\label{thm:c3v40-companion-decision}
Nothing is imported from SM V32 or YM V9.  Neither bears on the
third cycle's chapters or on the gates.
\end{theorem}

\begin{proof}
The SM results concern interpolations of lattice Dirac operators;
the YM results concern colourings, trees and Duhamel coefficients
in a cluster expansion.
None is a Navier--Stokes statement, and none concerns the ledger
identities, the extremal method or the Lamb--Bernoulli form of the
cross term.
\end{proof}

% =============================================================================
\section{Position after forty versions, and the block V41--V45}
\label{sec:c3v40-position}
% =============================================================================

\begin{assessment}[Position after forty versions of the third cycle]
\label{ass:c3v40-position-forty}
Seven chapters are written: cluster energy and the band (V2--V10),
localization (V11--V14), the anchor (V16--V19), the explicit floor
(V21--V24), the variance chapter (V26--V29), the defect block
(V31--V34), and the extremal ledger (V34--V39).  The cycle's
quantitative programme has reached its natural boundary: every
compactness constant of the homogeneous gate that can be made
explicit by the class bounds alone has been (the variance floor),
every one that cannot has its obstruction named (the defect rate,
N-C3-3; the pinching, through the pressure firewall), and the mean
identities of the second cycle have pointwise counterparts at
extremal configurations.  The gate G1 itself has been shown to be a
statement about the angle between the nonlinearity and the defect
that no comparison of norms can decide (V39).  The block V41--V45
studies that angle in its Lamb--Bernoulli form, as planned in
Assessment~\ref{ass:c3v39-direction}: V41 explicit bounds on the two
pieces and on the defect at extremal Dirichlet form; V42 the
pointwise structure of the Lamb--defect product on the anchor's
neighbourhood and the test configurations; V43 the extremal
Dirichlet form in Lamb--Bernoulli form; V44 consolidation; V45
audit.  No global regularity claim is made.
\end{assessment}

% =============================================================================
\section{Integrated V40 conclusion and next acquisition order}
\label{sec:c3v40-conclusion}
% =============================================================================

\begin{theorem}[What V40 establishes]
\label{thm:c3v40-conclusion}
Relative to the two legacy cycles and V1--V39: the companion-audit
decision (Theorem~\ref{thm:c3v40-companion-decision}) and the position.
The full Navier--Stokes stability/global-regularity theorem remains
unproved.
\end{theorem}

\begin{proof}
The cited records.
\end{proof}

\begin{programme}[V41 acquisition order]
\label{prog:c3v41}
\begin{enumerate}[label=\textup{(AC\arabic*)},leftmargin=4.8em]
\item Prove the explicit bounds on the Lamb piece,
      \(|\int\det(\Omega,V,Y)G|\le C_\infty\mathcal E^{1/2}\|Y\|_G\) with
      \(\mathcal E\le\frac2\nu\mathcal L\le\frac2\nu C_{\mathcal L}\), and on the Bernoulli piece,
      \(|\frac1{2\nu}\int(\Pi-c)(z\cdot Y)G|\le\frac1{2\nu}\||z|(\Pi-c)\|_G\|Y\|_G\), with
      \(\||z|P_V\|_G\le\|P_V\|_{L^3}(\int|z|^6G^3)^{1/6}\) and
      \(\||z||V|^2\|_G\le C_\infty\||z|V\|_G\le C_\infty\cdot2\nu^{1/2}(d+3\varphi)^{1/2}\); assemble
      the explicit constant \(M_{NY}\) with \(|\langle N,Y\rangle_G|\le M_{NY}\|Y\|_G\)
      at every configuration.
\item Deduce \(\|Y\|_G\le M_{NY}\) at every extremal-Dirichlet
      configuration and \(\langle\|Y\|_G^2\rangle\le\langle M_{NY}^2\rangle\) against every
      invariant measure; compare with the window bound
      \(\underline\delta\le10\kappa^2\nu C\) of V33 and record which is smaller.
\item The next compulsory companion audit is V45.
\end{enumerate}
\end{programme}

\begin{assessment}[Position after V40]
\label{ass:c3v40-position}
The companion audit is recorded; the Lamb--defect angle block
begins.  No full stability or global-regularity claim is made.
\end{assessment}

% =============================================================================
\section*{Version V40 (third cycle) append-only record}
\addcontentsline{toc}{section}{Version V40 (third cycle) append-only record}
% =============================================================================

The complete V39 source of the third cycle was retained before this
continuation.  Its SHA--256 digest was
\begin{center}\scriptsize\ttfamily
9a966cced810ae996c45e75e168a8fcd25c03e6e8ee4bbffd5e8dad24434fa13.
\end{center}
V40 performed the compulsory companion audit and recorded the
position.  No full regularity claim is made.

% =============================================================================
\section{V41: explicit bounds on the Lamb and Bernoulli pieces of the cross term}
\label{sec:c3v41-entry}
% =============================================================================

Items (AC1) and (AC2) of Programme~\ref{prog:c3v41} are carried out.
(AC1): the cross term of the gate, in its Lamb--Bernoulli form, is
bounded at every configuration by an explicit multiple of the
defect's Gaussian norm; the Lamb piece by the velocity constant
times the square root of twice the Dirichlet form, through the
closure identity, and the Bernoulli piece by the Gaussian sixth
moment of the pressure through the class \(L^3\) bound, the pressure
entering without a derivative, plus the Gaussian Hardy bound on the
kinetic part.  (AC2): at every extremal-Dirichlet configuration the
defect is at most this multiple, and against every invariant
measure the mean squared defect is at most its mean square; the
comparison with the window bound of V33 shows neither dominates,
the new bound having a part independent of the velocity constant.
No global regularity claim is made.  The next compulsory companion
audit is V45.

% =============================================================================
\section{The two pieces}
\label{sec:c3v41-pieces}
% =============================================================================

Throughout, the cross term is (Theorem C2-V3-cross)
\begin{equation}
 \langle N(V),Y\rangle_G=\int\det(\Omega,V,Y)\,G+\frac1{2\nu}\int\Pi\,(z\cdot Y)\,G,
 \qquad\Pi=P_V+\tfrac12|V|^2 ,
 \label{eq:c3v41-LB}
\end{equation}
and since \(\int(z\cdot Y)G=0\) (Remark~\ref{rem:c3v31-source}) the
constant \(\Pi\) may be shifted by any \(c\in\R\).  Put
\begin{equation}
 I_6:=\int_{\R^3}|z|^6G^3\,\dd z=\frac{105\,\pi^{3/2}}{8}\Bigl(\frac{4\nu}3\Bigr)^{9/2} .
 \label{eq:c3v41-I6}
\end{equation}

\begin{proposition}[The Lamb piece]
\label{prop:c3v41-lamb}
For every profile of the family and every \(\sigma\),
\begin{equation}
 \Bigl|\int\det(\Omega,V,Y)\,G\Bigr|\ \le\ C_\infty\,\mathcal E^{1/2}\|Y\|_G\ \le\ \kappa\,(2\mathcal L)^{1/2}\,\|Y\|_G\ \le\ \kappa\,(2C_{\mathcal L})^{1/2}\,\|Y\|_G .
 \label{eq:c3v41-lamb}
\end{equation}
\end{proposition}

\begin{proof}
\(|\det(\Omega,V,Y)|\le|\Omega||V||Y|\le C_\infty|\Omega||Y|\), Cauchy--Schwarz, and
\(\mathcal E=\|\Omega\|_G^2\le\frac2\nu\mathcal L\) by the closure identity (Theorem
C2-V21-closure); \(C_\infty(2\mathcal L/\nu)^{1/2}=\kappa(2\mathcal L)^{1/2}\); \(\mathcal L\le C_{\mathcal L}\).
\end{proof}

\begin{proposition}[The Bernoulli piece]
\label{prop:c3v41-bernoulli}
For every profile of the family and every \(\sigma\),
\begin{equation}
 \Bigl|\frac1{2\nu}\int\Pi\,(z\cdot Y)\,G\Bigr|\ \le\ \frac1{2\nu}\Bigl(\||z|P_V\|_G+\tfrac12\||z||V|^2\|_G\Bigr)\|Y\|_G ,
 \label{eq:c3v41-bernoulli}
\end{equation}
with
\begin{equation}
 \||z|P_V\|_G\le\|P_V\|_{L^3}\,I_6^{1/6}\le C_{\rm CZ}K_6^2\,I_6^{1/6},
 \qquad
 \||z||V|^2\|_G\le C_\infty\||z|V\|_G\le2\kappa\nu\,(d+3\varphi)^{1/2} .
 \label{eq:c3v41-moments}
\end{equation}
\end{proposition}

\begin{proof}
Cauchy--Schwarz in \(L^2(G)\) with \(|z|\) attached to \(\Pi\).  For the
pressure moment, H\"older with exponents \(3\) and \(\frac32\):
\(\int|z|^2P_V^2G\le(\int|P_V|^3)^{2/3}(\int|z|^6G^3)^{1/3}\), and \(\|P_V\|_3\le C_{\rm CZ}K_6^2\)
(C2-V36).  The Gaussian moment: \(\int|z|^6e^{-3|z|^2/(4\nu)}\dd z=4\pi\int_0^\infty s^8e^{-as^2}\dd s\)
with \(a=3/(4\nu)\), and \(\int_0^\infty s^8e^{-as^2}\dd s=\Gamma(\frac92)/(2a^{9/2})=\frac{105\sqrt\pi}{32}a^{-9/2}\),
giving \eqref{eq:c3v41-I6}.  For the kinetic moment, \(|V|\le C_\infty\) and
the Gaussian Hardy inequality \(\||z|V\|_G\le2\nu^{1/2}(d+3\varphi)^{1/2}\)
(Lemma C2-V36-hardy).
\end{proof}

% =============================================================================
\section{The explicit constant and its consequences}
\label{sec:c3v41-constant}
% =============================================================================

\begin{theorem}[Explicit bound on the cross term]
\label{thm:c3v41-MNY}
Define, at a configuration,
\begin{equation}
 M_{NY}:=\kappa(2\mathcal L)^{1/2}+\frac{C_{\rm CZ}K_6^2I_6^{1/6}}{2\nu}+\frac\kappa2(d+3\varphi)^{1/2},
 \qquad
 \overline M:=\kappa(2C_{\mathcal L})^{1/2}+\frac{C_{\rm CZ}K_6^2I_6^{1/6}}{2\nu}+\frac\kappa2(4\nu C+3C_\Phi)^{1/2} .
 \label{eq:c3v41-MNY}
\end{equation}
Then for every profile of the family and every \(\sigma\),
\begin{equation}
 \boxed{\quad
 |\langle N(V),Y\rangle_G|\ \le\ M_{NY}\,\|Y\|_G\ \le\ \overline M\,\|Y\|_G .
 \quad}
 \label{eq:c3v41-bound}
\end{equation}
Consequently: at every maximizer or minimizer of \(\mathcal L\) over the
family, \(\|Y\|_G\le M_{NY}\le\overline M\); and for every invariant
probability measure, \(\langle\|Y\|_G^2\rangle_{\mathfrak m}\le\langle M_{NY}^2\rangle_{\mathfrak m}\le\overline M^2\),
so that \(\underline\delta\le\overline M^2\).
\end{theorem}

\begin{proof}
\eqref{eq:c3v41-LB} with Propositions~\ref{prop:c3v41-lamb} and
\ref{prop:c3v41-bernoulli}; \(\frac1{2\nu}\cdot\frac12\cdot2\kappa\nu(d+3\varphi)^{1/2}=\frac\kappa2(d+3\varphi)^{1/2}\);
the class bounds for \(\overline M\).  At an extremum of \(\mathcal L\),
\(\langle N,Y\rangle_G=-\|Y\|_G^2\) (Proposition~\ref{prop:c3v33-pointwise} with
\(\mathcal L'=0\), or row (1) of Theorem~\ref{thm:c3v39-ledger}), so
\(\|Y\|_G^2\le M_{NY}\|Y\|_G\).  Against \(\mathfrak m\),
\(\langle\|Y\|_G^2\rangle=-\langle\langle N,Y\rangle_G\rangle\le\langle M_{NY}\|Y\|_G\rangle\le\langle M_{NY}^2\rangle^{1/2}\langle\|Y\|_G^2\rangle^{1/2}\).
\end{proof}

\begin{remark}[Comparison with the window bound of V33]
\label{rem:c3v41-comparison}
V33 gave \(\underline\delta\le10\kappa^2\nu C\), quadratic in \(\kappa\); the present
\(\overline M^2\) contains the term \((C_{\rm CZ}K_6^2I_6^{1/6}/(2\nu))^2\), independent of
\(\kappa\), which is a numerical multiple of \(K_P^2\) (with
\(K_P=4\nu^{-1/2}(4\pi\nu/3)^{1/4}C_{\rm CZ}K_6^2\), the ratio
\(C_{\rm CZ}K_6^2I_6^{1/6}/(2\nu)\big/(K_P/4)=(105\pi^{3/2}/8)^{1/6}(4/3)^{3/4}/(2(4\pi/3)^{1/4})\approx0.89\)).
For small \(\kappa\) the window bound is smaller; for \(\kappa\) of order one
the two are comparable.  The difference in structure is the point:
the V33 bound came from \(\|N\|_G\), where the pressure enters through
its gradient and needs the unweighted enstrophy; the present bound
comes from the Lamb--Bernoulli form, where the pressure enters
without a derivative and the Gaussian weight is harmless.  The
Bernoulli form is therefore the one in which the pressure can be
kept in Gaussian form, an exception to the standing firewall of V38
that the next versions exploit.
\end{remark}

% =============================================================================
\section{Integrated V41 conclusion and next acquisition order}
\label{sec:c3v41-conclusion}
% =============================================================================

\begin{theorem}[What V41 proves]
\label{thm:c3v41-conclusion}
Relative to the two legacy cycles and V1--V40: the explicit bounds
on the Lamb and Bernoulli pieces (Propositions~\ref{prop:c3v41-lamb},
\ref{prop:c3v41-bernoulli}) and the explicit bound on the cross term
with its consequences (Theorem~\ref{thm:c3v41-MNY}).  The full
Navier--Stokes stability/global-regularity theorem remains unproved.
\end{theorem}

\begin{proof}
The cited results.
\end{proof}

\begin{programme}[V42 acquisition order]
\label{prog:c3v42}
\begin{enumerate}[label=\textup{(AC\arabic*)},leftmargin=4.8em]
\item The pointwise structure: with the Lamb vector \(\ell:=\Omega\times V\),
      \(\det(\Omega,V,Y)=Y\cdot\ell\); write the localized gate of V13 in the
      form \(\int_U Y\cdot\ell\,G+\frac1{2\nu}\int_U(\Pi-c)(z\cdot Y)G\) and record
      that on the anchor's neighbourhood the anti-alignment is the
      defect pointing against the Lamb vector, corrected by the
      Bernoulli work against the radial defect.
\item The nearly constant core: for a profile with dominant level
      \(0\) and \(\eta=16\,\mathrm{Var}\), bound the Lamb piece by
      \(C_\infty\|\Omega\|_G\|Y\|_G\) with \(\|\Omega\|_G^2\le\|\nabla V\|_G^2\le\varphi c_r\eta/\nu\)
      (Lemma~\ref{lem:c3v24-h1}), so that the anti-alignment on a
      nearly constant core is carried by the Bernoulli piece up to
      \(O(\eta^{1/2})\); combine with the extremal-Dirichlet row to state
      what the Bernoulli work must be at an extremal-Dirichlet
      configuration that is a nearly constant core.
\item The next compulsory companion audit is V45.
\end{enumerate}
\end{programme}

\begin{assessment}[Position after V41]
\label{ass:c3v41-position}
The cross term is bounded explicitly by the defect's norm at every
configuration; the Bernoulli form keeps the pressure in Gaussian
form.  No full stability or global-regularity claim is made.
\end{assessment}

% =============================================================================
\section*{Version V41 (third cycle) append-only record}
\addcontentsline{toc}{section}{Version V41 (third cycle) append-only record}
% =============================================================================

The complete V40 source of the third cycle was retained before this
continuation.  Its SHA--256 digest was
\begin{center}\scriptsize\ttfamily
aada8e4c304a775a06deeffe485ea714a7ea11857cf0b2a1985011f3801010e9.
\end{center}
V41 proved the explicit bounds on the Lamb and Bernoulli pieces.
No full regularity claim is made.

% =============================================================================
\section{V42: the Lamb vector, the pointwise structure of the gate, and the nearly constant core}
\label{sec:c3v42-entry}
% =============================================================================

Items (AC1) and (AC2) of Programme~\ref{prog:c3v42} are carried out.
(AC1): the Lamb piece of the cross term is the Gaussian pairing of
the defect with the Lamb vector, which is orthogonal to velocity
and vorticity pointwise, so only the component of the defect normal
to the plane of velocity and vorticity enters; the localized gate is
restated in this form on the anchor's neighbourhood.  (AC2): on a
nearly constant core the Lamb vector is small in the Gaussian norm,
with an explicit bound linear in the square root of the variance,
and the kinetic part of the Bernoulli work is small in the same
way, so that the cross term is the pressure work against the
radial defect, equivalently the Gaussian pairing of the defect with
the pressure gradient, up to an explicit error; at an
extremal-Dirichlet configuration that is a nearly constant core the
pressure work is negative and equals minus the squared defect up to
that error, and by the identity of V27 the pressure gradient on the
core is generated by the far strain.  No global regularity claim is
made.  The next compulsory companion audit is V45.

% =============================================================================
\section{The Lamb vector and the pointwise structure}
\label{sec:c3v42-lamb}
% =============================================================================

\begin{definition}[Lamb vector]
\label{def:c3v42-lamb}
\(\ell:=\Omega\times V\), so that \((V\cdot\nabla)V=\ell+\nabla\frac12|V|^2\),
\(N(V)=\ell+\nabla\Pi\), and \(\det(\Omega,V,Y)=\ell\cdot Y\).
\end{definition}

\begin{proposition}[Pointwise structure of the Lamb piece]
\label{prop:c3v42-pointwise}
Pointwise, \(\ell\perp V\), \(\ell\perp\Omega\), and \(|\ell|\le|\Omega||V|\); hence
\(\ell\cdot Y=\ell\cdot Y^\perp\) with \(Y^\perp\) the component of \(Y\) orthogonal
to \(\operatorname{span}(V,\Omega)\), and
\begin{equation}
 \Bigl|\int\ell\cdot Y\,G\Bigr|\ \le\ \int|\Omega||V||Y^\perp|G\ \le\ C_\infty\|\Omega\|_G\|Y^\perp\|_G .
 \label{eq:c3v42-lamb-perp}
\end{equation}
The localized gate \eqref{eq:c3v13-gate} reads
\begin{equation}
 \int\Bigl[\int_U\ell\cdot Y^\perp\,G+\frac1{2\nu}\int_U(\Pi-c)(z\cdot Y)G\Bigr]\dd\mathfrak m\ \le\ -\underline\delta+O(\varepsilon_{\rm th}^{1/4}) .
 \label{eq:c3v42-gate}
\end{equation}
\end{proposition}

\begin{proof}
Properties of the cross product; \eqref{eq:c3v13-gate} with
\(\det(\Omega,V,Y)=\ell\cdot Y=\ell\cdot Y^\perp\).
\end{proof}

\begin{remark}[Reading]
\label{rem:c3v42-reading}
The gate asks that, on the anchor's neighbourhood and on average,
the defect point against the Lamb vector in the direction normal to
the velocity--vorticity plane, or that the Bernoulli pressure do
negative work against the radial defect, at a total rate at least
the defect rate.  The two pieces are of different nature: the Lamb
piece is local and vanishes where the flow is Beltrami
(\(\Omega\parallel V\)) or irrotational; the Bernoulli piece is nonlocal
through the pressure and vanishes where the defect is tangential
(\(z\cdot Y=0\)).  The second cycle's vanishing statements (C2-V3) are
the pointwise reading; the third cycle's localization (V13, V17)
restricts both to at most \(N_C(2r_2+1)^3\) cells.
\end{remark}

% =============================================================================
\section{The nearly constant core}
\label{sec:c3v42-core}
% =============================================================================

\begin{proposition}[On a nearly constant core the cross term is pressure work]
\label{prop:c3v42-core}
Suppose \(\mathrm{Var}<\frac1{32}\), \(\eta=16\,\mathrm{Var}\), and the dominant level is
\(k_1=0\).  Then, with \(c_r=(\frac18+2C_r^2)^{1/2}\),
\begin{align}
 \Bigl|\int\ell\cdot Y\,G\Bigr|&\ \le\ \kappa\,(2\varphi c_r\eta)^{1/2}\,\|Y\|_G ,
 \label{eq:c3v42-lamb-core}\\
 \Bigl|\frac1{4\nu}\int|V|^2(z\cdot Y)G\Bigr|&\ \le\ \Bigl[\frac{\varphi}{|G|^{1/2}}\Bigl(\frac{(3+4c_r)\eta}{\nu}\Bigr)^{1/2}+\kappa\bigl(\varphi\eta(3+4c_r)\bigr)^{1/2}\Bigr]\|Y\|_G ,
 \label{eq:c3v42-kinetic-core}
\end{align}
and therefore
\begin{equation}
 \boxed{\quad
 \langle N(V),Y\rangle_G=\frac1{2\nu}\int P_V(z\cdot Y)G+E_{\rm core}\|Y\|_G=\langle\nabla P_V,Y\rangle_G+E_{\rm core}\|Y\|_G,
 \qquad |E_{\rm core}|\le C_{\rm core}\,\eta^{1/2},
 \quad}
 \label{eq:c3v42-core}
\end{equation}
with \(C_{\rm core}:=\kappa(2C_\Phi c_r)^{1/2}+C_\Phi|G|^{-1/2}((3+4c_r)/\nu)^{1/2}+\kappa(C_\Phi(3+4c_r))^{1/2}\).
\end{proposition}

\begin{proof}
Write \(V=\bar V+\widetilde V\) as in Proposition~\ref{prop:c3v27-core}:
\(\|\widetilde V\|_G^2\le\varphi\eta\), \(\nu\|\nabla\widetilde V\|_G^2\le\varphi c_r\eta\)
(Lemma~\ref{lem:c3v24-h1}), \(\||z|\widetilde V\|_G^2\le4\nu\varphi\eta(3+4c_r)\), \(|\bar V|\le(\varphi/|G|)^{1/2}\).
Lamb piece: \(\Omega=\operatorname{curl}\widetilde V\), \(|\Omega|^2\le2|\nabla\widetilde V|^2\), so
\(\|\Omega\|_G\le(2\varphi c_r\eta/\nu)^{1/2}\) and \eqref{eq:c3v42-lamb-perp} with
\(C_\infty=\kappa\nu^{1/2}\) gives \eqref{eq:c3v42-lamb-core}.  Kinetic part:
\(|V|^2=|\bar V|^2+2\bar V\cdot\widetilde V+|\widetilde V|^2\); the constant contributes
\(\frac{|\bar V|^2}{4\nu}\int(z\cdot Y)G=0\); the cross part is bounded by
\(\frac1{2\nu}|\bar V|\||z|\widetilde V\|_G\|Y\|_G\le\frac1{2\nu}(\varphi/|G|)^{1/2}\cdot2(\nu\varphi\eta(3+4c_r))^{1/2}\|Y\|_G\),
which is the first term of \eqref{eq:c3v42-kinetic-core}; the quadratic
part, with \(|\widetilde V|\le|V|+|\bar V|\le2C_\infty\), by
\(\frac1{2\nu}C_\infty\||z|\widetilde V\|_G\|Y\|_G=\kappa(\varphi\eta(3+4c_r))^{1/2}\|Y\|_G\), the second
term.  Then
\eqref{eq:c3v41-LB} with \(\Pi=P_V+\frac12|V|^2\), and
\(\frac1{2\nu}\int P_V(z\cdot Y)G=\langle\nabla P_V,Y\rangle_G\) by \((z\cdot Y)G=-2\nu Y\cdot\nabla G\)
and \(\operatorname{div}Y=0\); \(\varphi\le C_\Phi\) for \(C_{\rm core}\).
\end{proof}

\begin{corollary}[Extremal-Dirichlet nearly constant core]
\label{cor:c3v42-extremal-core}
If a maximizer or minimizer of \(\mathcal L\) over the family is a nearly
constant core (\(k_1=0\), \(\eta<\frac12\)), then at it
\begin{equation}
 \langle\nabla P_V,Y\rangle_G=-\|Y\|_G^2-E_{\rm core}\|Y\|_G\ \le\ -\|Y\|_G\bigl(\|Y\|_G-C_{\rm core}\eta^{1/2}\bigr):
 \label{eq:c3v42-extremal-core}
\end{equation}
the pressure gradient does negative work against the defect, of
size the squared defect, up to the error; in particular
\(\|Y\|_G\le\|\nabla P_V\|_G+C_{\rm core}\eta^{1/2}\) there.
\end{corollary}

\begin{proof}
Row (1) of the extremal ledger (\(\langle N,Y\rangle_G=-\|Y\|_G^2\) at extremal
\(\mathcal L\)) with \eqref{eq:c3v42-core}; Cauchy--Schwarz for the last bound.
\end{proof}

\begin{assessment}[The pressure gradient on the core is generated by the far strain]
\label{ass:c3v42-far}
By Proposition~\ref{prop:c3v27-identity}, \(\nabla P_V=-\nabla\Gamma_N*q\) with
\(q=\partial_iV_j\partial_jV_i\), and on a nearly constant core the strain is
small in the Gaussian sense (\(\int_{B_\rho}|q|\le e^{\rho^2/(4\nu)}\varphi c_r\eta/\nu\))
while the class strain outside a ball of radius \(\rho\) contributes
to \(\nabla P_V\) on the core a field of size at most \(C/(4\pi\rho^2)\)
pointwise, nearly constant across the Gaussian bulk when \(\rho\) is
large compared with \((4\nu)^{1/2}\).  The gate on a nearly constant
core is therefore, up to the explicit errors, the statement that
the far field's pressure gradient on the core is anti-aligned with
the defect at the rate of the defect rate; and the momentum identity
(V18, V27) says that the same pressure gradient is what sustains or
decays the core's Gaussian momentum.  The nearly constant core is
thus the configuration on which the gate is entirely nonlocal: its
Lamb piece is negligible, its kinetic Bernoulli piece is negligible,
and what remains is the pairing of the core's defect with a
pressure gradient produced elsewhere.  No sign follows: the far
strain is unconstrained by the class beyond its total enstrophy.
This is recorded as the structural reason why the nearly constant
core resisted every floor of the variance chapter.
\end{assessment}

% =============================================================================
\section{Integrated V42 conclusion and next acquisition order}
\label{sec:c3v42-conclusion}
% =============================================================================

\begin{theorem}[What V42 proves]
\label{thm:c3v42-conclusion}
Relative to the two legacy cycles and V1--V41: the pointwise
structure of the Lamb piece (Proposition~\ref{prop:c3v42-pointwise}),
the pressure-work form of the cross term on a nearly constant core
(Proposition~\ref{prop:c3v42-core}) and its extremal-Dirichlet
corollary (Corollary~\ref{cor:c3v42-extremal-core}).  The full
Navier--Stokes stability/global-regularity theorem remains unproved.
\end{theorem}

\begin{proof}
The cited results.
\end{proof}

\begin{programme}[V43 acquisition order]
\label{prog:c3v43}
\begin{enumerate}[label=\textup{(AC\arabic*)},leftmargin=4.8em]
\item The extremal Dirichlet form in Lamb--Bernoulli form: at an
      extremum of \(\mathcal L\), \(\int\ell\cdot Y^\perp G+\frac1{2\nu}\int(\Pi-c)(z\cdot Y)G=-\|Y\|_G^2\);
      with Proposition~\ref{prop:c3v41-lamb}, the Bernoulli work is at
      most \(-\|Y\|_G(\|Y\|_G-\kappa(2\mathcal L)^{1/2})\), negative whenever
      \(\|Y\|_G>\kappa(2\mathcal L)^{1/2}\); record the explicit alternative at the
      extremal Dirichlet form: either the defect is at most
      \(\kappa(2\mathcal L)^{1/2}\) or the Bernoulli work is negative of size at
      least \(\|Y\|_G(\|Y\|_G-\kappa(2\mathcal L)^{1/2})\).
\item The Beltrami and tangential cases: at a configuration where
      \(Y^\perp=0\) (defect in the velocity--vorticity plane) the gate is
      pure Bernoulli work; at one where \(z\cdot Y=0\) it is pure Lamb
      pairing; state the two reduced gates and what the explicit
      bounds say for each.
\item The next compulsory companion audit is V45.
\end{enumerate}
\end{programme}

\begin{assessment}[Position after V42]
\label{ass:c3v42-position}
The gate is a Lamb pairing plus a Bernoulli work; on a nearly
constant core it is pure pressure work, generated by the far
strain.  No full stability or global-regularity claim is made.
\end{assessment}

% =============================================================================
\section*{Version V42 (third cycle) append-only record}
\addcontentsline{toc}{section}{Version V42 (third cycle) append-only record}
% =============================================================================

The complete V41 source of the third cycle was retained before this
continuation.  Its SHA--256 digest was
\begin{center}\scriptsize\ttfamily
6d7839ad6e62981af33cc63f1320363db9095de25b623b4d316f030ec8846039.
\end{center}
V42 recorded the pointwise structure of the gate and the pressure
form on a nearly constant core.  No full regularity claim is made.

% =============================================================================
\section{V43: the extremal Dirichlet form in Lamb--Bernoulli form, and the two reduced gates}
\label{sec:c3v43-entry}
% =============================================================================

Items (AC1) and (AC2) of Programme~\ref{prog:c3v43} are carried out.
(AC1): at a configuration of extremal Dirichlet form the Lamb
pairing and the Bernoulli work sum exactly to minus the squared
defect, and the explicit bound on the Lamb pairing gives an
alternative: either the defect is at most the velocity constant
times the square root of twice the Dirichlet form, or the Bernoulli
work is negative with an explicit lower bound on its size.  (AC2):
the two degenerate cases of the pointwise structure, defect in the
velocity--vorticity plane and tangential defect, reduce the gate to
its Bernoulli and its Lamb piece respectively; the reduced gates
are stated with what the explicit bounds say, the tangential case
giving an explicit bound on the defect rate of the same form as the
window bound.  No global regularity claim is made.  The next
compulsory companion audit is V45.

% =============================================================================
\section{The extremal Dirichlet form in Lamb--Bernoulli form}
\label{sec:c3v43-extremal}
% =============================================================================

Write \(\mathcal B:=\frac1{2\nu}\int(\Pi-c)(z\cdot Y)G\) for the Bernoulli work
(independent of \(c\)) and \(\mathcal A:=\int\ell\cdot Y^\perp G\) for the Lamb pairing, so
that \(\langle N,Y\rangle_G=\mathcal A+\mathcal B\) (Theorem C2-V3-cross,
Proposition~\ref{prop:c3v42-pointwise}).

\begin{theorem}[Extremal Dirichlet form, Lamb--Bernoulli form]
\label{thm:c3v43-extremal}
Let \(W^*\) be a maximizer or minimizer of \(\mathcal L\) over \(\mathcal F(x_0)\).  At
\(W^*\), \(\sigma=0\):
\begin{equation}
 \boxed{\quad
 \mathcal A+\mathcal B=-\|Y\|_G^2,
 \qquad
 -\|Y\|_G\bigl(\|Y\|_G+\kappa(2\mathcal L)^{1/2}\bigr)\ \le\ \mathcal B\ \le\ -\|Y\|_G\bigl(\|Y\|_G-\kappa(2\mathcal L)^{1/2}\bigr).
 \quad}
 \label{eq:c3v43-extremal}
\end{equation}
Hence one of the following holds at \(W^*\):
\begin{enumerate}[label=\textup{(\alph*)},leftmargin=3.2em]
\item \(\|Y\|_G\le\kappa(2\mathcal L)^{1/2}\): the defect at the extremal Dirichlet
      form is at most the velocity constant times the square root of
      twice the Dirichlet form;
\item \(\|Y\|_G>\kappa(2\mathcal L)^{1/2}\) and the Bernoulli work is negative,
      \(\mathcal B\le-\|Y\|_G(\|Y\|_G-\kappa(2\mathcal L)^{1/2})<0\).
\end{enumerate}
In both cases \(\|Y\|_G\le M_{NY}\) (Theorem~\ref{thm:c3v41-MNY}).
\end{theorem}

\begin{proof}
Row (1) of the extremal ledger, \(\langle N,Y\rangle_G=-\|Y\|_G^2\) at \(W^*\), and
\(|\mathcal A|\le\kappa(2\mathcal L)^{1/2}\|Y\|_G\) (Proposition~\ref{prop:c3v41-lamb}).
\end{proof}

\begin{remark}[Reading]
\label{rem:c3v43-reading}
The Dirichlet form is extremal exactly when the nonlinearity and the
defect are anti-aligned to the full extent of the defect; the Lamb
pairing can supply at most \(\kappa(2\mathcal L)^{1/2}\|Y\|_G\) of it, so unless
the defect is small compared with the velocity constant, the
Bernoulli pressure must do negative work against the radial defect
at the extremal Dirichlet form.  With the explicit lower bound
\(\mathcal L\ge\mathcal L_-\) (V34, positive for small \(\kappa\)), case (a) reads
\(\|Y\|_G\le\kappa(2\mathcal L)^{1/2}\), which for small \(\kappa\) is a small defect at
the extremal configuration; the two cases are the pointwise
counterpart of the tension recorded in V32--V33 between the defect
rate and \(\kappa^2\).  No exclusion: both cases are compatible with the
class.
\end{remark}

% =============================================================================
\section{The two reduced gates}
\label{sec:c3v43-reduced}
% =============================================================================

\begin{proposition}[The Beltrami-type reduction]
\label{prop:c3v43-beltrami}
Suppose that on the support of a blow-up measure \(\mathfrak m\) the
defect lies pointwise in the velocity--vorticity plane,
\(Y^\perp\equiv0\) (in particular wherever the flow is Beltrami,
\(\Omega\parallel V\), or irrotational).  Then the gate \eqref{eq:c3v42-gate}
reads \(\int\mathcal B_U\,\dd\mathfrak m\le-\underline\delta+O(\varepsilon_{\rm th}^{1/4})\) with
\(\mathcal B_U\) the Bernoulli work on the configuration neighbourhood, and
the explicit bound gives
\begin{equation}
 \underline\delta\ \le\ \Bigl\langle\frac1{2\nu}\||z|(\Pi-c)\|_{L^2(G;U)}\|Y\|_G\Bigr\rangle+O(\varepsilon_{\rm th}^{1/4})
 \ \le\ \Bigl(\frac{C_{\rm CZ}K_6^2I_6^{1/6}}{2\nu}+\frac\kappa2(4\nu C+3C_\Phi)^{1/2}\Bigr)\langle\|Y\|_G^2\rangle^{1/2}+O(\varepsilon_{\rm th}^{1/4}) .
 \label{eq:c3v43-beltrami}
\end{equation}
\end{proposition}

\begin{proof}
\(\mathcal A=0\) pointwise; Proposition~\ref{prop:c3v41-bernoulli} restricted to
\(U\) (the restriction only decreases the \(L^2(G;U)\) norm) and
Cauchy--Schwarz against \(\mathfrak m\).
\end{proof}

\begin{proposition}[The tangential reduction]
\label{prop:c3v43-tangential}
Suppose that on the support of \(\mathfrak m\) the defect is tangential,
\(z\cdot Y\equiv0\).  Then the gate reads
\(\int\mathcal A_U\,\dd\mathfrak m\le-\underline\delta+O(\varepsilon_{\rm th}^{1/4})\), and
\begin{equation}
 \underline\delta\ \le\ \kappa\,\langle(2\mathcal L)^{1/2}\|Y\|_G\rangle+O(\varepsilon_{\rm th}^{1/4})
 \ \le\ \kappa(2C_{\mathcal L})^{1/2}\langle\|Y\|_G^2\rangle^{1/2}+O(\varepsilon_{\rm th}^{1/4}) ;
 \label{eq:c3v43-tangential}
\end{equation}
moreover, against every invariant measure,
\(\langle\|Y\|_G^2\rangle=-\langle\mathcal A\rangle\le\kappa(2C_{\mathcal L})^{1/2}\langle\|Y\|_G^2\rangle^{1/2}\), so
\(\langle\|Y\|_G^2\rangle\le2\kappa^2C_{\mathcal L}=\kappa^2(\nu C+\tfrac12C_\Phi)\) and
\(\underline\delta\le\kappa^2(\nu C+\frac12C_\Phi)\).
\end{proposition}

\begin{proof}
\(\mathcal B=0\) pointwise; Proposition~\ref{prop:c3v41-lamb}; the mean
identity \(\langle\langle N,Y\rangle_G\rangle=-\langle\|Y\|_G^2\rangle\) with \(\langle N,Y\rangle_G=\mathcal A\);
\(C_{\mathcal L}=(2\nu C+C_\Phi)/4\).
\end{proof}

\begin{assessment}[Item (AC2): what the reductions show]
\label{ass:c3v43-reductions}
The two reductions are illustrations, not cases that occur: a
blow-up measure whose support has \(Y^\perp\equiv0\) or \(z\cdot Y\equiv0\)
throughout is not known to exist, and the class does not force
either.  Their value is comparative.  In the tangential case the
gate is purely local and the explicit bound on the defect rate,
\(\kappa^2(\nu C+\frac12C_\Phi)\), is of the same form as the window bound
\(10\kappa^2\nu C\) of V33, confirming that the Lamb pairing alone cannot
carry an anti-alignment beyond \(\kappa^2\) times the class enstrophy.
In the Beltrami-type case the gate is purely nonlocal, and the
bound has the \(\kappa\)-independent pressure term: the Bernoulli work
is the piece that survives a small velocity constant.  Taken with
V42, the picture of the gate is: its local part is bounded by
\(\kappa\) times the Dirichlet form and vanishes on Beltrami flows and
on nearly constant cores; its nonlocal part is the pressure work
against the radial defect, bounded by the Gaussian sixth moment of
the pressure, and it is where the anti-alignment must come from
whenever the local part is small.  Every route to closing G1 must
therefore bound the pressure work against the radial defect from
below by something better than \(-\|Y\|_G^2\), which is a statement
about the sign of \(\int P_V(z\cdot Y)G\), the correlation of the pressure
with the radial defect.  No global regularity claim is made.
\end{assessment}

% =============================================================================
\section{Integrated V43 conclusion and next acquisition order}
\label{sec:c3v43-conclusion}
% =============================================================================

\begin{theorem}[What V43 proves]
\label{thm:c3v43-conclusion}
Relative to the two legacy cycles and V1--V42: the extremal
Dirichlet form in Lamb--Bernoulli form (Theorem~\ref{thm:c3v43-extremal})
and the two reduced gates (Propositions~\ref{prop:c3v43-beltrami},
\ref{prop:c3v43-tangential}).  The full Navier--Stokes
stability/global-regularity theorem remains unproved.
\end{theorem}

\begin{proof}
The cited results.
\end{proof}

\begin{programme}[V44 acquisition order]
\label{prog:c3v44}
\begin{enumerate}[label=\textup{(AC\arabic*)},leftmargin=4.8em]
\item Consolidate V41--V43 as the angle chapter: the explicit bound
      \(M_{NY}\), the pointwise structure, the nearly constant core as
      pure pressure work, the extremal Dirichlet form in
      Lamb--Bernoulli form, and the reduced gates.
\item Fix the direction for V46--V50: the correlation of the pressure
      with the radial defect, \(\int P_V(z\cdot Y)G\); its exact expression
      through the Newtonian kernel (\(P_V=\Gamma_N*q\)) as a bilinear form
      in the strain and the radial defect; the sign question on
      the anchor's neighbourhood; and whether the extremal method
      applies to an observable whose derivative is the Bernoulli
      work.
\item V45 compulsory companion audit and delivery.
\end{enumerate}
\end{programme}

\begin{assessment}[Position after V43]
\label{ass:c3v43-position}
At the extremal Dirichlet form the Bernoulli work is negative unless
the defect is small; the gate's nonlocal part is the pressure work
against the radial defect.  No full stability or global-regularity
claim is made.
\end{assessment}

% =============================================================================
\section*{Version V43 (third cycle) append-only record}
\addcontentsline{toc}{section}{Version V43 (third cycle) append-only record}
% =============================================================================

The complete V42 source of the third cycle was retained before this
continuation.  Its SHA--256 digest was
\begin{center}\scriptsize\ttfamily
fd52c5fd55e4bd4f727d2787945fbad09e81c819898b79a5357943ba717de01e.
\end{center}
V43 proved the extremal Dirichlet form in Lamb--Bernoulli form and
the reduced gates.  No full regularity claim is made.

% =============================================================================
\section{V44: the angle chapter consolidated, and the pressure--radial-defect correlation}
\label{sec:c3v44-entry}
% =============================================================================

Items (AC1) and (AC2) of Programme~\ref{prog:c3v44} are carried out.
The angle chapter V41--V43 is consolidated.  For the direction, the
Bernoulli work is written as a bilinear form in the strain invariant
and the Newtonian potential of the radial-defect density, a
potential of dipole type since the radial defect has zero Gaussian
mass; and the Bernoulli work is shown to be the log-time derivative
of a quarter of the flux up to the pairing of the linearized
Bernoulli pressure with the radial velocity, so that the extremal
method applies at the extremal flux.  The block V46--V50 is fixed on
the sign of this correlation.  No global regularity claim is made.
The next compulsory companion audit is V45.

% =============================================================================
\section{The angle chapter}
\label{sec:c3v44-chapter}
% =============================================================================

\begin{theorem}[The angle chapter, V39 and V41--V43]
\label{thm:c3v44-chapter}
At a homogeneous type-I point:
\begin{enumerate}[label=\textup{(\arabic*)},leftmargin=3.2em]
\item (Angle form of G1, V39.)  \(\langle\|N\|^2\rangle-\langle\|Y\|^2\rangle=\langle\|L_-V\|^2\rangle\ge(\frac14+v_{\rm exp})c_\Phi\),
      and the gate is the impossibility of
      \(\langle\|N\|\|Y\|\cos\theta\rangle=-\langle\|Y\|^2\rangle\); no norm comparison decides
      it.
\item (Explicit bound, V41.)  \(|\langle N,Y\rangle_G|\le M_{NY}\|Y\|_G\) with
      \(M_{NY}=\kappa(2\mathcal L)^{1/2}+C_{\rm CZ}K_6^2I_6^{1/6}/(2\nu)+\frac\kappa2(d+3\varphi)^{1/2}\),
      the pressure entering in Gaussian form through its \(L^3\) bound;
      \(\|Y\|_G\le M_{NY}\) at extremal Dirichlet form and
      \(\underline\delta\le\overline M^2\).
\item (Structure, V42.)  \(\langle N,Y\rangle_G=\mathcal A+\mathcal B\), the Lamb pairing
      \(\mathcal A=\int\ell\cdot Y^\perp G\) seeing only the defect normal to the
      velocity--vorticity plane, and the Bernoulli work
      \(\mathcal B=\frac1{2\nu}\int(\Pi-c)(z\cdot Y)G\); on a nearly constant core
      \(\langle N,Y\rangle_G=\langle\nabla P_V,Y\rangle_G+O(\eta^{1/2})\|Y\|_G\) with the
      pressure gradient generated by the far strain.
\item (Extremal Dirichlet form, V43.)  \(\mathcal A+\mathcal B=-\|Y\|_G^2\) there;
      either \(\|Y\|_G\le\kappa(2\mathcal L)^{1/2}\) or \(\mathcal B<0\) with an explicit size;
      the tangential reduction gives \(\underline\delta\le\kappa^2(\nu C+\frac12C_\Phi)\),
      the Beltrami-type reduction leaves the \(\kappa\)-independent
      pressure term.
\end{enumerate}
\end{theorem}

\begin{proof}
Theorems~\ref{thm:c3v39-angle}, \ref{thm:c3v41-MNY}, \ref{thm:c3v43-extremal};
Propositions~\ref{prop:c3v42-pointwise}, \ref{prop:c3v42-core},
\ref{prop:c3v43-beltrami}, \ref{prop:c3v43-tangential};
Firewall~\ref{fw:c3v39-no-norms}.
\end{proof}

% =============================================================================
\section{The pressure--radial-defect correlation}
\label{sec:c3v44-correlation}
% =============================================================================

\begin{proposition}[Bilinear form of the pressure work]
\label{prop:c3v44-bilinear}
With \(q=\partial_iV_j\partial_jV_i=|S|^2-\frac12|\Omega|^2\) (\(S\) the strain-rate tensor),
\(P_V=\Gamma_N*q\), and
\begin{equation}
 \psi_Y:=\Gamma_N*\bigl((z\cdot Y)G\bigr),\qquad -\Delta\psi_Y=(z\cdot Y)G,\qquad \int(z\cdot Y)G=0 ,
 \label{eq:c3v44-psi}
\end{equation}
the pressure work is
\begin{equation}
 \boxed{\quad
 \frac1{2\nu}\int P_V\,(z\cdot Y)\,G\ =\ \frac1{2\nu}\int q\,\psi_Y\,\dd y
 \ =\ \frac1{2\nu}\int\Bigl(|S|^2-\frac12|\Omega|^2\Bigr)\psi_Y\,\dd y ,
 \quad}
 \label{eq:c3v44-bilinear}
\end{equation}
and \(\psi_Y\) is of dipole type: \(|\psi_Y(y)|\le\frac{D_Y}{4\pi|y|^2}(1+o(1))\) as
\(|y|\to\infty\) with dipole moment \(D_Y=|\int z\,(z\cdot Y)G|\le\||z|^2\|_G\|Y\|_G=(60\nu^2|G|)^{1/2}\|Y\|_G\);
in particular \(\psi_Y\in L^p\) for every \(p>\frac32\) and the bilinear form
converges absolutely for \(q\in L^1\cap L^\infty\).
\end{proposition}

\begin{proof}
\(P_V=\Gamma_N*q\) (Proposition~\ref{prop:c3v27-identity}) and Fubini,
\(\int(\Gamma_N*q)f=\int q(\Gamma_N*f)\) for \(f=(z\cdot Y)G\), which is bounded with
Gaussian decay; \(q=\partial_iV_j\partial_jV_i=\operatorname{tr}((\nabla V)^2)=|S|^2-|A|^2\) with
\(A\) the antisymmetric part, \(|A|^2=\frac12|\Omega|^2\).  Zero mass of \(f\):
Remark~\ref{rem:c3v31-source}.  A zero-mass density with Gaussian
decay has a Newtonian potential whose leading term at infinity is
the dipole term, \(\frac1{4\pi}\frac{y\cdot\int zf}{|y|^3}\); the dipole moment is
bounded by Cauchy--Schwarz with \(\int|z|^4G=60\nu^2|G|\).
\end{proof}

\begin{proposition}[The Bernoulli work is the derivative of the flux, up to a pairing]
\label{prop:c3v44-flux}
For every profile of the family and every \(\sigma\),
\begin{equation}
 \mathcal B=\frac1{2\nu}\int\Pi\,(z\cdot Y)G=\frac14\frac{\dd j}{\dd\sigma}-\frac1{2\nu}\int\partial_\sigma\Pi\,(z\cdot V)G,
 \qquad \partial_\sigma\Pi=P'+V\cdot Y ,
 \label{eq:c3v44-flux}
\end{equation}
and at a maximizer or minimizer of the flux \(j\) over the family,
\begin{equation}
 \mathcal B=-\frac1{2\nu}\int(P'+V\cdot Y)(z\cdot V)\,G .
 \label{eq:c3v44-extremal-flux}
\end{equation}
\end{proposition}

\begin{proof}
\(j=\frac2\nu\int\Pi(z\cdot V)G\) (C2-V32), so \(\frac14j=\frac1{2\nu}\int\Pi(z\cdot V)G\) and
differentiating in \(\sigma\) with \(\partial_\sigma V=Y\) gives \eqref{eq:c3v44-flux};
Lemma~\ref{lem:c3v34-extremal} for \(F=j\) (continuous, \(C^1\) in \(\sigma\) on
the family by the derivative bounds).
\end{proof}

\begin{assessment}[Direction for V46--V50: the sign of the correlation]
\label{ass:c3v44-direction}
The gate's nonlocal part is the correlation of the strain invariant
\(|S|^2-\frac12|\Omega|^2\) with the dipole potential of the radial defect.
Its sign is the question.  V46: the potential \(\psi_Y\) on the anchor's
neighbourhood, where the radial defect lives up to
\(O(\varepsilon_{\rm th}^{1/4})\) (V11--V13): \(\psi_Y\) is positive where the radial
defect is outward-weighted and the strain invariant is positive
where strain dominates rotation; write the correlation on the
configuration as a sum over cells of the cell's strain invariant
against the potential of the other cells' radial defects and its
own, and record the self-term's sign.  V47: the extremal flux row
\eqref{eq:c3v44-extremal-flux} with the explicit bounds on the
pairing (the \(V\cdot Y\) part by the Hardy bound, the \(P'\) part again
by the linearized pressure, now paired with \(z\cdot V\) rather than
\(z\cdot Y\), i.e.\ with the radial velocity, whose Gaussian moments the
class does control through \(\||z|V\|_G\le2\nu^{1/2}(d+3\varphi)^{1/2}\)); determine
whether the pairing \(\int P'(z\cdot V)G=\int q'\psi_V\) with \(q'=2\partial_iV_j\partial_jY_i\)
and \(\psi_V\) the potential of \((z\cdot V)G\) admits an explicit bound
through the dipole decay of \(\psi_V\) and the \(L^1\) bound on \(q'\), i.e.\
whether the standing firewall is bypassed when the linearized
pressure is paired with the radial velocity.  V48: if V47 closes,
the explicit two-sided bound on the Bernoulli work at the extremal
flux, and its combination with the extremal-Dirichlet alternative
of V43.  V49 consolidates; V50 audits.  No global regularity claim
is made.
\end{assessment}

% =============================================================================
\section{Integrated V44 conclusion and next acquisition order}
\label{sec:c3v44-conclusion}
% =============================================================================

\begin{theorem}[What V44 establishes]
\label{thm:c3v44-conclusion}
Relative to the two legacy cycles and V1--V43: the consolidated
angle chapter (Theorem~\ref{thm:c3v44-chapter}), the bilinear form
(Proposition~\ref{prop:c3v44-bilinear}), the flux-derivative form of
the Bernoulli work (Proposition~\ref{prop:c3v44-flux}) and the
direction.  The full Navier--Stokes stability/global-regularity
theorem remains unproved.
\end{theorem}

\begin{proof}
The cited results.
\end{proof}

\begin{programme}[V45 acquisition order]
\label{prog:c3v45}
\begin{enumerate}[label=\textup{(AC\arabic*)},leftmargin=4.8em]
\item The compulsory companion audit: read the current SM and YM
      notes and record their digests and decision.
\item Record the position after forty-five versions; delivery.
\end{enumerate}
\end{programme}

\begin{assessment}[Position after V44]
\label{ass:c3v44-position}
The angle chapter is consolidated; the Bernoulli work is a bilinear
form in the strain invariant and a dipole potential, and the
derivative of the flux up to a pairing.  No full stability or
global-regularity claim is made.
\end{assessment}

% =============================================================================
\section*{Version V44 (third cycle) append-only record}
\addcontentsline{toc}{section}{Version V44 (third cycle) append-only record}
% =============================================================================

The complete V43 source of the third cycle was retained before this
continuation.  Its SHA--256 digest was
\begin{center}\scriptsize\ttfamily
976dfc38ee564683d51c36fd14659b255c3fcd53d3dc110877e8fb84bdeb3ab4.
\end{center}
V44 consolidated the angle chapter and fixed the direction.  No
full regularity claim is made.

% =============================================================================
\section{V45: compulsory companion audit and the position after forty-five versions}
\label{sec:c3v45-entry}
% =============================================================================

This is the forty-fifth version of the third cycle.  The compulsory
review of the two companion programmes is performed first, and the
position after forty-five versions is recorded with the plan for
the block V46--V50.  No global regularity claim is made.  The next
compulsory companion audit is V50.

% =============================================================================
\section{Compulsory V45 review of the current companion notes}
\label{sec:c3v45-companion-audit}
% =============================================================================

The current companion files in \texttt{latex\_notes} were inspected
read-only.  Both programmes have moved since the V40 audit: the SM
programme from V32 to V33 of its seventh series, the YM programme
from V9 to V11 of its sixth series.  The frozen checkpoints are
\begin{align}
 &\texttt{Renewal\_SM\_Einstein\_Continuum\_Research\_Notes\_V33.tex},
 \notag\\[-2pt]
 &\qquad
 \texttt{2120d65f6edca90cbeb18d28aa329eed912c2528940a422b98f473b33c522252},
 \label{eq:c3v45-SM-checkpoint}\\
 &\texttt{Renewal\_Yang\_Mills\_Mass\_Gap\_Research\_Notes\_V11.tex},
 \notag\\[-2pt]
 &\qquad
 \texttt{223bed98d69854fa8ae2aa8351edee32cb7b6cca14ac8d51a82c8ef9e8f52fcb}.
 \label{eq:c3v45-YM-checkpoint}
\end{align}
(SM V33 has 13590 lines; YM V11 has 5027 lines.  As at V40, the YM
file was rewritten between the first inspection for this audit and
the assembly of this version: at the first inspection \texttt{V11}
was byte-identical to \texttt{V10} (4605 lines, digest
\texttt{0a7418e3\ldots}); the digest recorded is that of the final
inspection.)

\emph{SM V33.}  The SM programme constructed a concrete global
simplicial spin mesh: a face-compatible edgewise subdivision, an
intrinsic nodal spin interpolation, global Galerkin Dirac and
Wilson forms with moment consistency, and a remaining stratified
frozen gap problem.

\emph{YM V10--V11.}  The YM programme performed its companion audit
and developed a weighted overlay (one grouping edge, simultaneous
exchange on an arbitrary overlay, a high-excess rejected-forest
card), then exact grouping-forest determinants: the weighted
Laplacian of resolved components, two exact forest identities, a
diffuse-component closure, why the component mass is not the
original total mark, and a concentrated component-Laplacian card.

\begin{theorem}[V45 companion-audit decision]
\label{thm:c3v45-companion-decision}
Nothing is imported from SM V33 or YM V11.  Neither bears on the
third cycle's chapters or on the gates.
\end{theorem}

\begin{proof}
The SM results concern simplicial meshes and Galerkin forms for
lattice Dirac operators; the YM results concern overlays and
determinants of forests in a cluster expansion.  None is a Navier--Stokes statement, and
none concerns the Lamb--Bernoulli form, the pressure--radial-defect
correlation or the extremal method.
\end{proof}

% =============================================================================
\section{Position after forty-five versions, and the block V46--V50}
\label{sec:c3v45-position}
% =============================================================================

\begin{assessment}[Position after forty-five versions of the third cycle]
\label{ass:c3v45-position-forty-five}
Eight chapters are written: cluster energy and the band (V2--V10),
localization (V11--V14), the anchor (V16--V19), the explicit floor
(V21--V24), the variance chapter (V26--V29), the defect block
(V31--V34), the extremal ledger (V34--V39), and the angle chapter
(V41--V44).  The gate G1 has been reduced, step by step, to the sign
of one bilinear form: the correlation of the strain invariant
\(|S|^2-\frac12|\Omega|^2\) with the dipole potential of the radial defect
on the anchor's neighbourhood, everything else in the cross term
being bounded explicitly by the velocity constant times the
Dirichlet form (the Lamb pairing) or being small on the
configurations the floors could not exclude (the nearly constant
core).  The standing pressure firewall is the reason the reduction
stops there, and V44 identified the one pairing, linearized
pressure against radial velocity, where the firewall may be
bypassed because the radial velocity has controlled Gaussian
moments.  The block V46--V50 follows Assessment~\ref{ass:c3v44-direction}:
V46 the potential of the radial defect on the configuration and the
cell decomposition of the correlation; V47 the extremal flux row
and the bypass question; V48 the two-sided bound at the extremal
flux if V47 closes; V49 consolidation; V50 audit.  No global
regularity claim is made.
\end{assessment}

% =============================================================================
\section{Integrated V45 conclusion and next acquisition order}
\label{sec:c3v45-conclusion}
% =============================================================================

\begin{theorem}[What V45 establishes]
\label{thm:c3v45-conclusion}
Relative to the two legacy cycles and V1--V44: the companion-audit
decision (Theorem~\ref{thm:c3v45-companion-decision}) and the position.
The full Navier--Stokes stability/global-regularity theorem remains
unproved.
\end{theorem}

\begin{proof}
The cited records.
\end{proof}

\begin{programme}[V46 acquisition order]
\label{prog:c3v46}
\begin{enumerate}[label=\textup{(AC\arabic*)},leftmargin=4.8em]
\item The potential of the radial defect: for a density
      \(f=(z\cdot Y)G\) of zero mass supported, up to \(O(\varepsilon_{\rm th}^{1/4})\),
      on the configuration neighbourhood \(U\) of at most \(N_C(2r_2+1)^3\)
      cells, write \(\psi_Y=\sum_\beta\psi_\beta\) with \(\psi_\beta=\Gamma_N*(f\phi_\beta)\) the cell
      potentials, and the correlation
      \(\int q\psi_Y=\sum_{\alpha,\beta}\int_\alpha q\,\psi_\beta\); bound the far-cell
      contributions by the dipole decay and the class enstrophy,
      and isolate the self-terms \(\int_\alpha q\psi_\alpha\).
\item The self-term: for one cell, \(\int_\alpha q\,\Gamma_N*(f\phi_\alpha)\) with
      \(q=|S|^2-\frac12|\Omega|^2\) and \(f=(z\cdot Y)G\); record that it has no
      sign in general, and the two model cases (a cell in which
      \(z\cdot Y\) has one sign; a cell in which the flow is
      strain-dominated), with what each gives.
\item The next compulsory companion audit is V50.
\end{enumerate}
\end{programme}

\begin{assessment}[Position after V45]
\label{ass:c3v45-position}
The companion audit is recorded; the correlation block begins.  No
full stability or global-regularity claim is made.
\end{assessment}

% =============================================================================
\section*{Version V45 (third cycle) append-only record}
\addcontentsline{toc}{section}{Version V45 (third cycle) append-only record}
% =============================================================================

The complete V44 source of the third cycle was retained before this
continuation.  Its SHA--256 digest was
\begin{center}\scriptsize\ttfamily
c0dc701ce35f6a4fcb0259ca2a9210c74d91509d9439635fe5cf2c6dcee6154e.
\end{center}
V45 performed the compulsory companion audit and recorded the
position.  No full regularity claim is made.

% =============================================================================
\section{V46: the pressure work on the configuration as an energy--Hessian pairing, and its sign structure}
\label{sec:c3v46-entry}
% =============================================================================

Items (AC1) and (AC2) of Programme~\ref{prog:c3v46} are carried out,
with one change of form.  The strain form of the pressure work
(V44) is not the right one for localization, because the far quiet
cells have small energy but not small enstrophy; the energy form,
in which the pressure is generated by the velocity tensor, is the
one in which the far field contributes to the pressure on the
configuration at order the threshold (V11, C2-V58).  In that form
the pressure work on the configuration is, up to explicit far
corrections, the pairing of the velocity tensor on a halo of the
configuration with the Hessian of the Newtonian potential of the
radial-defect density on the configuration; it decomposes over
cells, and the cross terms between distinct cells have the
multipole form of a point mass, whose sign is fixed by the
orientation of the velocity relative to the radial-defect cell:
radial velocity gives the sign of the cell's radial-defect mass,
azimuthal velocity the opposite.  The gate's nonlocal part is
therefore a geometric statement about swirl around cells of
outward radial defect and radial flow around cells of inward radial
defect.  No global regularity claim is made.  The next compulsory
companion audit is V50.

% =============================================================================
\section{Localization of the pressure work in the energy form}
\label{sec:c3v46-localization}
% =============================================================================

Notation of V11--V13 and V16--V17: on the trapped sub-windows at a
homogeneous type-I point, \(U\) is the configuration neighbourhood
(at most \(N_C(2r_2+1)^3\) parabolic cells), \(F\) its complement of quiet
cells (energy \(e_\beta<\varepsilon_{\rm th}\), trapped enstrophy), and \(U'\) the halo
of \(U\), the cells within lattice distance \(2\) of \(U\).  Let
\(\chi_{U'}\) be a smooth cut-off equal to \(1\) on \(U'\) and vanishing
outside the cells within distance \(1\) of \(U'\), and split the pressure
by its source:
\begin{equation}
 P_V=P^{\rm loc}+P^{\rm far},\qquad
 P^{\rm loc}:=\Gamma_N*\partial_i\partial_j\bigl(\chi_{U'}V_iV_j\bigr),\qquad
 P^{\rm far}:=\Gamma_N*\partial_i\partial_j\bigl((1-\chi_{U'})V_iV_j\bigr).
 \label{eq:c3v46-split}
\end{equation}

\begin{proposition}[The far pressure is flat on the configuration]
\label{prop:c3v46-far}
There is \(C_{10}\), depending on \(\nu\), \(C_\infty\), \(C\) and the counts, such
that on the trapped sub-windows \(|\nabla P^{\rm far}|\le C_{10}\varepsilon_{\rm th}\) on \(U\), and
hence, with \(c\) the value of \(P^{\rm far}\) at a fixed point of \(U\),
\begin{equation}
 \Bigl|\frac1{2\nu}\int_U(P^{\rm far}-c)(z\cdot Y)G\Bigr|\ \le\ \frac{C_{10}\varepsilon_{\rm th}}{2\nu}\,\operatorname{diam}(U)\,\||z|\|_{L^2(G;U)}\,\|Y\|_G\ =:\ C_{11}\varepsilon_{\rm th}\|Y\|_G .
 \label{eq:c3v46-far}
\end{equation}
\end{proposition}

\begin{proof}
The source of \(P^{\rm far}\) is supported on quiet cells, and its gradient at
\(z\in U\) is \(\int\nabla\partial_i\partial_j\Gamma_N(z-y)(1-\chi_{U'})V_iV_j(y)\dd y\), with a
kernel bounded by \(C|z-y|^{-4}\) and a source bounded on each quiet
cell by its energy, at most \(\varepsilon_{\rm th}\); the lattice sum
\(\sum_\gamma(1+\operatorname{dist}(z,Q_\gamma))^{-4}\) converges (this is the mechanism of
Lemma C2-V58-far, used in the proof of Theorem~\ref{thm:c3v11-far}).
The bound on the work is the mean value theorem along a path in \(U\)
and Cauchy--Schwarz.
\end{proof}

\begin{theorem}[The pressure work on the configuration as an energy--Hessian pairing]
\label{thm:c3v46-pairing}
Let \(f_U:=(z\cdot Y)G\,\mathbf 1_U\) and \(\psi_U:=\Gamma_N*f_U\).  On the trapped
sub-windows,
\begin{equation}
 \boxed{\quad
 \frac1{2\nu}\int_U(P_V-c)(z\cdot Y)G\ =\ \frac1{2\nu}\int\chi_{U'}\,(V\otimes V):\nabla^2\psi_U\,\dd y\ +\ E_{\rm far},
 \qquad |E_{\rm far}|\le C_{11}\varepsilon_{\rm th}\|Y\|_G ,
 \quad}
 \label{eq:c3v46-pairing}
\end{equation}
and the pairing decomposes over the cells of \(U\): with
\(f_\beta:=(z\cdot Y)G\phi_\beta\) and \(\psi_\beta:=\Gamma_N*f_\beta\),
\begin{equation}
 \int\chi_{U'}(V\otimes V):\nabla^2\psi_U=\sum_{\beta\in U}\int\chi_{U'}(V\otimes V):\nabla^2\psi_\beta .
 \label{eq:c3v46-cells}
\end{equation}
\end{theorem}

\begin{proof}
By \eqref{eq:c3v46-split} and Proposition~\ref{prop:c3v46-far}, the
work of \(P^{\rm far}-c\) is \(E_{\rm far}\).  For the local part,
\(\int P^{\rm loc}f_U=\int\bigl(\Gamma_N*\partial_i\partial_j(\chi_{U'}V_iV_j)\bigr)f_U=\int\chi_{U'}V_iV_j\,\partial_i\partial_j(\Gamma_N*f_U)\)
by two integrations by parts (the source is compactly supported and
smooth, \(f_U\) bounded with Gaussian decay, \(\Gamma_N*f_U\) smooth away
from the support of \(f_U\) and \(C^{1,\alpha}\) across it; the boundary
terms vanish at infinity by the decay of \(\Gamma_N*f_U\) and its
derivatives).  Linearity in \(f_U=\sum_{\beta\in U}f_\beta\) (the partition of
unity on \(U\)) gives \eqref{eq:c3v46-cells}.
\end{proof}

\begin{remark}[Why the energy form]
\label{rem:c3v46-why}
In the strain form of V44 the far cells contribute to the pressure
on \(U\) through their enstrophy with a kernel \(|z-y|^{-2}\), a
contribution the class bounds only by the total enstrophy \(C\), not
by the threshold; in the energy form they contribute through their
energy with a kernel \(|z-y|^{-4}\), at order \(\varepsilon_{\rm th}\).  The two
forms are exact representations of the same pressure, the
difference being terms supported on the cut; the energy form is the
one adapted to the quiet far field, and it is the form of the
second cycle's far-pressure lemma.  The Hessian of the Newtonian
potential is the singular-integral operator \(-R_iR_j\) with its
local part \(-\frac13\delta_{ij}f\), so \eqref{eq:c3v46-pairing} is the
statement that the pressure work is \(-\frac1{2\nu}\int f_U\,R_iR_j(\chi_{U'}V_iV_j)\),
the localized pressure paired with the radial defect, as it should
be; the value of writing it with the Hessian is the cell
decomposition and the multipole structure below.
\end{remark}

% =============================================================================
\section{The multipole form of the cross terms, and the sign structure}
\label{sec:c3v46-multipole}
% =============================================================================

\begin{proposition}[Point-mass form of a cell's Hessian]
\label{prop:c3v46-multipole}
Let \(\beta\in U\), \(z_\beta\) its centre, \(m_\beta:=\int f_\beta\) the cell's
radial-defect mass, and \(\mu_\beta:=\int|y-z_\beta||f_\beta(y)|\dd y\) its absolute
first moment.  For \(|z-z_\beta|\ge2\), with \(\rho=|z-z_\beta|\) and \(\hat e=(z-z_\beta)/\rho\),
\begin{equation}
 \nabla^2\psi_\beta(z)=\frac{m_\beta}{4\pi\rho^3}\bigl(3\hat e\otimes\hat e-I\bigr)+R_\beta(z),\qquad
 |R_\beta(z)|\le\frac{C_{12}\,\mu_\beta}{\rho^4},
 \label{eq:c3v46-hessian}
\end{equation}
so that
\begin{equation}
 (V\otimes V):\nabla^2\psi_\beta(z)=\frac{m_\beta}{4\pi\rho^3}\bigl(3(V\cdot\hat e)^2-|V|^2\bigr)+O\Bigl(\frac{\mu_\beta|V|^2}{\rho^4}\Bigr).
 \label{eq:c3v46-sign}
\end{equation}
In particular, at a point at distance at least \(2\) from the cell:
if \(m_\beta>0\) (outward radial defect on the cell), a velocity radial
with respect to the cell (\(V\parallel\hat e\)) contributes positively and
an azimuthal one (\(V\perp\hat e\)) negatively; if \(m_\beta<0\), the
reverse.
\end{proposition}

\begin{proof}
Taylor expansion of \(\nabla^2\Gamma_N(z-y)\) in \(y\) about \(z_\beta\) over the
support of \(f_\beta\) (a cell of diameter \(\sqrt3\), at distance at least
\(2-\frac{\sqrt3}2\) from \(z\)): the zeroth-order term is
\(m_\beta\nabla^2\Gamma_N(z-z_\beta)=\frac{m_\beta}{4\pi\rho^3}(3\hat e\otimes\hat e-I)\), and the
remainder is bounded by \(\sup|\nabla^3\Gamma_N|\) on the segment times the
first moment, with \(|\nabla^3\Gamma_N(w)|\le C_{12}'|w|^{-4}\).  The sign
statement is \(3(V\cdot\hat e)^2-|V|^2=3|V|^2\cos^2\alpha-|V|^2\), positive for
\(\cos^2\alpha>\frac13\) and negative for \(\cos^2\alpha<\frac13\), with \(\alpha\) the
angle between \(V\) and the direction to the cell.
\end{proof}

\begin{proposition}[Self-term]
\label{prop:c3v46-self}
For a cell \(\alpha\in U\), the self-term \(\int\chi_{U'}\phi_\alpha(V\otimes V):\nabla^2\psi_\alpha\) has
no sign in general.  If \(f_\alpha\ge0\) on the cell (outward radial defect),
then \(\operatorname{tr}\nabla^2\psi_\alpha=-f_\alpha\le0\) on the cell, so an isotropic
velocity tensor on the cell (\(V\otimes V\) replaced by its trace part
\(\frac13|V|^2I\)) gives the self-term \(-\frac13\int\phi_\alpha|V|^2f_\alpha\le0\); the
anisotropic part \((V\otimes V)^\circ\) pairs with the traceless part of the
Hessian and has either sign.
\end{proposition}

\begin{proof}
\(-\Delta\psi_\alpha=f_\alpha\); the decomposition of \(V\otimes V\) into its trace and
traceless parts; the traceless part of the Hessian of a potential is
the Riesz-transform part, of either sign.
\end{proof}

\begin{assessment}[Item (AC2): the sign structure of the gate's nonlocal part]
\label{ass:c3v46-structure}
On the configuration, up to \(O(\varepsilon_{\rm th})\), the pressure work is a
sum over cells of the configuration of the velocity tensor on the
halo paired with the Hessian of each cell's radial-defect potential.
Its sign structure, from Propositions~\ref{prop:c3v46-multipole} and
\ref{prop:c3v46-self}, is:
\begin{enumerate}[label=\textup{(\roman*)},leftmargin=3.2em]
\item cross terms: a cell of outward radial defect (\(m_\beta>0\))
      contributes negative work where the surrounding velocity is
      azimuthal about it (swirl) and positive where it is radial
      (inflow or outflow); a cell of inward radial defect
      (\(m_\beta<0\)) the reverse;
\item self-terms: the isotropic part of the velocity tensor on a
      cell of outward radial defect gives negative work; the
      anisotropic part has no sign.
\end{enumerate}
The gate G1 (\(\mathcal B\le-\|Y\|_G^2\) at extremal Dirichlet form, and
\(\int\mathcal B\,\dd\mathfrak m\le-\underline\delta\) on average when the Lamb pairing is
small) therefore asks, in geometric terms, that the configuration
swirl about the cells where the defect pushes outward and flow
radially about the cells where it pushes inward, on average and
weighted by the radial-defect masses, at a rate at least the defect
rate.  This is the first form of the gate in which every term is a
local geometric quantity of the configuration (velocity orientation
relative to radial-defect cells), with the nonlocality reduced to
the multipole kernels; whether the class forces the opposite
orientation somewhere is the question the next versions ask.  No
global regularity claim is made.
\end{assessment}

% =============================================================================
\section{Integrated V46 conclusion and next acquisition order}
\label{sec:c3v46-conclusion}
% =============================================================================

\begin{theorem}[What V46 proves]
\label{thm:c3v46-conclusion}
Relative to the two legacy cycles and V1--V45: the flatness of the far
pressure on the configuration (Proposition~\ref{prop:c3v46-far}), the
energy--Hessian pairing with its cell decomposition
(Theorem~\ref{thm:c3v46-pairing}), the multipole form of the cross
terms (Proposition~\ref{prop:c3v46-multipole}) and the self-term
(Proposition~\ref{prop:c3v46-self}).  The full Navier--Stokes
stability/global-regularity theorem remains unproved.
\end{theorem}

\begin{proof}
The cited results.
\end{proof}

\begin{programme}[V47 acquisition order]
\label{prog:c3v47}
\begin{enumerate}[label=\textup{(AC\arabic*)},leftmargin=4.8em]
\item The bypass at the extremal flux: write the pairing of
      \eqref{eq:c3v44-extremal-flux} in the energy form,
      \(\int P'(z\cdot V)G=\int(V\otimes Y+Y\otimes V):\nabla^2\psi_V\) with
      \(\psi_V=\Gamma_N*((z\cdot V)G)\), whose mass vanishes and whose dipole
      moment is \(2\nu|G|\bar V\); with the pointwise growth bound
      \(|Y(z)|\le C_Y'(1+|z|)\) on the family (to be derived from the
      derivative bounds (E5) and the pressure gradient's local
      boundedness, or recorded as a hypothesis), bound the pairing
      explicitly through \(|\nabla^2\psi_V(y)|\le C(1+|y|)^{-3}\) and the class
      bounds; record whether the standing firewall is bypassed in
      this pairing.
\item If (AC1) closes: the explicit two-sided bound on the Bernoulli
      work at the extremal flux, and its combination with
      Theorem~\ref{thm:c3v43-extremal}.
\item The next compulsory companion audit is V50.
\end{enumerate}
\end{programme}

\begin{assessment}[Position after V46]
\label{ass:c3v46-position}
The gate's nonlocal part is an energy--Hessian pairing over the
configuration's cells with an explicit multipole sign structure.
No full stability or global-regularity claim is made.
\end{assessment}

% =============================================================================
\section*{Version V46 (third cycle) append-only record}
\addcontentsline{toc}{section}{Version V46 (third cycle) append-only record}
% =============================================================================

The complete V45 source of the third cycle was retained before this
continuation.  Its SHA--256 digest was
\begin{center}\scriptsize\ttfamily
6fbbd2bc6dd988fcbbfa7533e4a938e6168549c5fecb1648e85029c5d624671b.
\end{center}
V46 proved the energy--Hessian form of the pressure work on the
configuration and its sign structure.  No full regularity claim is
made.

% =============================================================================
\section{V47: the pairing at the extremal flux is explicit without a bypass, and the pressure firewall is revised}
\label{sec:c3v47-entry}
% =============================================================================

Items (AC1) and (AC2) of Programme~\ref{prog:c3v47} are carried out,
with a different outcome from the one planned.  The double
integration by parts proposed as a bypass is not justified in the
class: the linearized pressure grows linearly at infinity, through
the radial derivative of the pressure that the rescaling produces,
and the boundary terms of the second integration do not vanish.  It
is also unnecessary.  The linearized pressure decomposes exactly
into minus the pressure, minus half its radial derivative, and the
rescaled pressure of the time derivative; each piece is bounded
pointwise by the class, the growing piece linearly, and the Gaussian
absorbs the growth.  The pairing with the radial velocity is
therefore explicitly bounded, and the same observation revises the
standing firewall of V38: every Gaussian pairing with a pressure or
a linearized pressure in the series is explicitly bounded by
pointwise class constants of supremum type, because the class
bounds all derivatives pointwise; what the firewall correctly
records is that such bounds do not reduce to the Gaussian
dissipation, the enstrophy and the velocity constant, so they carry
no smallness.  The Bernoulli work at the extremal flux is bounded
two-sidedly, and the extremal-defect row of V38 is completed with
an explicit supremum bound on the defect for small derivative
constants.  No global regularity claim is made.  The next
compulsory companion audit is V50.

% =============================================================================
\section{Pointwise bounds}
\label{sec:c3v47-pointwise}
% =============================================================================

The class bounds of C2-V4 are used with the pointwise derivative
bounds that (E5) transports to the family, denoted
\(|\nabla^kV|\le C_\infty^{(k)}\) for \(k\le4\) (\(C_\infty^{(0)}=C_\infty\), \(C_\infty^{(1)}=C_\infty'\),
\(C_\infty^{(2)}=C_\infty''\)), together with \(\|V\|_{L^6}\le K_6\) and \(\|\nabla V\|_{L^2}^2\le C\).
The kernel of \(R_iR_j\) is \(K_{ij}(w)=\frac1{4\pi}\bigl(3w_iw_j/|w|^2-\delta_{ij}\bigr)|w|^{-3}\)
(plus the local term \(-\frac13\delta_{ij}\delta\)), and
\begin{equation}
 c_1:=\int_{|w|<1}|K(w)|_F|w|\,\dd w,\qquad
 c_p:=\Bigl(\int_{|w|>1}|K|_F^{p'}\Bigr)^{1/p'}\ (p'=p/(p-1)),\quad p\in\{2,3,6\},
 \label{eq:c3v47-c}
\end{equation}
are numerical constants (\(|K|_F\) the Frobenius norm of the matrix
kernel, at most \(|w|^{-3}/\pi\)).

\begin{lemma}[Pointwise bounds on Riesz transforms]
\label{lem:c3v47-riesz}
For a matrix-valued \(g=(g_{ij})\), bounded, Lipschitz and in \(L^p(\R^3)\),
\(p\in\{2,3,6\}\) (Frobenius norms),
\(\|\sum_{ij}R_iR_jg_{ij}\|_{L^\infty}\le\frac13\|\operatorname{tr}g\|_\infty+c_1[g]_{\rm Lip}+c_p\|g\|_{L^p}\).
\end{lemma}

\begin{proof}
The local term \(-\frac13\operatorname{tr}g\); the near part
\(\int_{|w|<1}K(w):(g(x-w)-g(x))\dd w\) (the
kernel has zero mean on spheres) bounded by \(c_1[g]_{\rm Lip}\); the far
part by H\"older with \(c_p\).
\end{proof}

\begin{lemma}[Pressure, time derivative, defect, linearized pressure]
\label{lem:c3v47-decomposition}
On the family, at every \(\sigma\), with \(V_t:=\nu\Delta V-(V\cdot\nabla)V-\nabla P_V\) (the
rescaled time derivative of the solution) and \(P_t\) the pressure of
the tensor \(V\otimes V_t+V_t\otimes V\) (the rescaled pressure of the time
derivative):
\begin{enumerate}[label=\textup{(\roman*)},leftmargin=3.2em]
\item \(|P_V|\le C_{P,0}\), \(|\nabla P_V|\le C_P\), \(|\nabla^2P_V|\le C_P'\), with
      \(C_{P,0}=\frac13C_\infty^2+2c_1C_\infty C_\infty'+c_3K_6^2\),
      \(C_P=\sqrt3\bigl[\frac23C_\infty C_\infty'+c_1(2C_\infty'^2+2C_\infty C_\infty'')+2c_2C_\infty C^{1/2}\bigr]\) (the
      bracket bounds each component \(\partial_lP_V\), the factor \(\sqrt3\) the
      Euclidean norm), and
      \(C_P'\) explicit in \(C_\infty^{(k)}\), \(k\le3\), \(K_6\), \(C\);
\item \(Y=-\frac12V-\frac12z\cdot\nabla V+V_t\) with \(|V_t|\le C_t:=\nu C_\infty''+C_\infty C_\infty'+C_P\);
      hence \(|Y(z)|\le C_Y'+\frac12C_\infty'|z|\), \(C_Y':=\frac12C_\infty+C_t\);
\item \(P'=\partial_\sigma P_V=-P_V-\frac12z\cdot\nabla P_V+P_t\) with \(|P_t|\le C_{P_t}\) and
      \(|\nabla P_t|\le C_{P_t}'\) explicit in \(C_\infty^{(k)}\), \(k\le4\), \(K_6\), \(C\); hence
      \(|P'(z)|\le C_{P,0}+C_{P_t}+\frac12C_P|z|\) and
      \(|\nabla P'(z)|\le\frac32C_P+C_{P_t}'+\frac12C_P'|z|\).
\end{enumerate}
\end{lemma}

\begin{proof}
(i) Lemma~\ref{lem:c3v47-riesz} with \(g=V_iV_j\) (\(p=3\), \(\|g\|_3\le K_6^2\),
\([g]_{\rm Lip}\le2C_\infty C_\infty'\)), \(g=\partial_l(V_iV_j)\) (\(p=2\), \(\|g\|_2\le2C_\infty C^{1/2}\),
\([g]_{\rm Lip}\le2C_\infty'^2+2C_\infty C_\infty''\)), and \(g=\partial_l\partial_m(V_iV_j)\) (split into
\(\partial V\partial V\in L^2\) and \(V\partial^2V\in L^6\)).  (ii) With \(V(z,\sigma)=\lambda u(x_0+\lambda z,T_*-\lambda^2)\),
\(\lambda=e^{-\sigma/2}\): \(\partial_\sigma V=-\frac12V-\frac12z\cdot\nabla V+\lambda^3\partial_tu\), and
\(\lambda^3\partial_tu\) is the rescaled Navier--Stokes right side
\(\nu\Delta V-(V\cdot\nabla)V-\nabla P_V\); the bound is the class.  (iii) Likewise
\(P_V=\lambda^2p\) gives \(\partial_\sigma P_V=-P_V-\frac12z\cdot\nabla P_V+\lambda^4\partial_tp\), and
\(\lambda^4\partial_tp\) is the pressure of \(\partial_t(u\otimes u)\) rescaled, i.e.\ of
\(V\otimes V_t+V_t\otimes V\); Lemma~\ref{lem:c3v47-riesz} with \(p=6\)
(\(\|V\otimes V_t\|_6\le K_6C_t\)) and the derivative bounds for the Lipschitz
seminorms (\(|\nabla V_t|\le\nu C_\infty'''+C_\infty'^2+C_\infty C_\infty''+C_P'\), and one order
more for \(\nabla P_t\)).
\end{proof}

\begin{firewall}[The bypass is not justified]
\label{fw:c3v47-no-bypass}
The identity \(\int P'(z\cdot V)G=\int(V\otimes Y+Y\otimes V):\nabla^2\psi_V\) proposed in
V46 requires two integrations by parts against a source of linear
growth (\(Y\), Lemma~\ref{lem:c3v47-decomposition}(ii)) and a potential
\(\psi_V\) of dipole type; the boundary term of the second integration
on the sphere of radius \(\rho\) is of size \(\rho^2\cdot\rho\cdot\rho^{-3}=O(1)\) and does
not vanish, and the alternative order of integration meets
\(|\nabla P'|\) of linear growth against \(|\psi_V|\sim\rho^{-2}\), a boundary term
of size \(O(\rho)\).  The identity is therefore not established in the
class (it holds when \(\bar V=0\), where \(\psi_V\) is of quadrupole type),
and it is not used.
\end{firewall}

% =============================================================================
\section{The pairing at the extremal flux}
\label{sec:c3v47-pairing}
% =============================================================================

\begin{theorem}[Explicit bound on the pairing of the linearized pressure with the radial velocity]
\label{thm:c3v47-pairing}
On the family, at every \(\sigma\), with \(X=\||z|V\|_G\le2\nu^{1/2}(d+3\varphi)^{1/2}\),
\begin{equation}
 \boxed{\quad
 \Bigl|\int P'(z\cdot V)G\Bigr|\ \le\ \bar J:=\Bigl(C_{P,0}+C_{P_t}+\frac{(6\nu)^{1/2}}2C_P\Bigr)|G|^{1/2}\cdot2\nu^{1/2}(d+3\varphi)^{1/2} .
 \quad}
 \label{eq:c3v47-J}
\end{equation}
\end{theorem}

\begin{proof}
By Lemma~\ref{lem:c3v47-decomposition}(iii), \(|P'|\le C_{P,0}+C_{P_t}+\frac12C_P|z|\),
so \(|\int P'(z\cdot V)G|\le(C_{P,0}+C_{P_t})\int|z||V|G+\frac12C_P\int|z|^2|V|G\); by
Cauchy--Schwarz, \(\int|z||V|G\le X\|1\|_G=X|G|^{1/2}\) and
\(\int|z|^2|V|G\le X\||z|\|_G=X(6\nu|G|)^{1/2}\); then the Hardy bound on \(X\)
(Lemma C2-V36-hardy).
\end{proof}

\begin{corollary}[The Bernoulli work at the extremal flux]
\label{cor:c3v47-extremal-flux}
At every maximizer or minimizer of the flux \(j\) over the family,
\begin{equation}
 |\mathcal B|\ \le\ \frac{\bar J}{2\nu}+\kappa\,(d+3\varphi)^{1/2}\,\|Y\|_G .
 \label{eq:c3v47-B}
\end{equation}
\end{corollary}

\begin{proof}
Proposition~\ref{prop:c3v44-flux}: \(\mathcal B=-\frac1{2\nu}\int(P'+V\cdot Y)(z\cdot V)G\) at
extremal \(j\); Theorem~\ref{thm:c3v47-pairing} for the first term, and
\(|\int(V\cdot Y)(z\cdot V)G|\le C_\infty X\|Y\|_G\le2\kappa\nu(d+3\varphi)^{1/2}\|Y\|_G\) for the
second.
\end{proof}

\begin{remark}[The two extremal rows do not combine]
\label{rem:c3v47-no-combination}
Theorem~\ref{thm:c3v43-extremal} bounds \(\mathcal B\) at the extremal
Dirichlet form, Corollary~\ref{cor:c3v47-extremal-flux} at the extremal
flux; the configurations differ in general and the two bounds do
not conjoin (Assessment~\ref{ass:c3v39-combining}).  Item (AC2) of
Programme~\ref{prog:c3v47} is closed with this record.
\end{remark}

% =============================================================================
\section{The pressure firewall revised}
\label{sec:c3v47-firewall}
% =============================================================================

\begin{theorem}[Every Gaussian pressure pairing is explicit at supremum level]
\label{thm:c3v47-revised}
Let \(m_1:=\int|z|G=32\pi\nu^2\) and \(m_2:=\int|z|^2G=6\nu|G|\).  On the family, at
every \(\sigma\):
\begin{align}
 \|\nabla P_V\|_G^2&\le C_P^2|G| ,
 \label{eq:c3v47-gradP}\\
 |\langle Y,\nabla P'\rangle_G|&\le C_{YP'}:=C_Y'\Bigl(\frac32C_P+C_{P_t}'\Bigr)|G|
 +\Bigl[\frac12C_\infty'\Bigl(\frac32C_P+C_{P_t}'\Bigr)+\frac12C_Y'C_P'\Bigr]m_1+\frac14C_\infty'C_P'\,m_2 .
 \label{eq:c3v47-YP'}
\end{align}
Consequently, at a maximizer \(W_Y\) of the defect over the family,
\begin{equation}
 \frac\nu2\|\nabla Y\|_G^2+\Bigl(\frac12-C_\infty'-\frac{\sqrt3}2\kappa-\frac{\kappa^2}2\Bigr)\|Y\|_G^2\ \le\ C_{YP'},
 \qquad
 \sup_{\mathcal F(x_0)}\|Y\|_G^2\ \le\ \frac{C_{YP'}}{\frac12-C_\infty'-\frac{\sqrt3}2\kappa-\frac{\kappa^2}2}
 \ \text{ if the denominator is positive.}
 \label{eq:c3v47-supY}
\end{equation}
\end{theorem}

\begin{proof}
\eqref{eq:c3v47-gradP}: \(|\nabla P_V|\le C_P\) and \(\int G=|G|\).  \eqref{eq:c3v47-YP'}:
\(|Y||\nabla P'|\le(C_Y'+\frac12C_\infty'|z|)(\frac32C_P+C_{P_t}'+\frac12C_P'|z|)\) integrated
against \(G\), with \(\int|z|G=4\pi\int_0^\infty s^3e^{-s^2/(4\nu)}\dd s=4\pi\cdot8\nu^2\).
\eqref{eq:c3v47-supY}: Proposition~\ref{prop:c3v38-advective},
\eqref{eq:c3v38-at-max}, with the pressure pairing bounded by \(C_{YP'}\).
\end{proof}

\begin{firewall}[The standing firewall, revised form]
\label{fw:c3v47-revised}
Firewall~\ref{fw:c3v38-pressure} stated that the Gaussian pairings
involving the linearized pressure are not explicitly bounded in the
class.  That was too strong: the class bounds all derivatives of
\(V\) pointwise (E5), the Riesz transforms of Lipschitz functions with
an \(L^p\) tail are pointwise bounded (Lemma~\ref{lem:c3v47-riesz}), and
the Gaussian absorbs polynomial growth; hence every such pairing is
bounded by an explicit constant of supremum type, as
\eqref{eq:c3v47-gradP}--\eqref{eq:c3v47-YP'} show.  What remains true,
and is the revised firewall, is that these bounds are in terms of
the pointwise constants \(C_\infty^{(k)}\), \(K_6\), \(C\) and the Gaussian
moments, not in terms of \(d\), \(\varphi\), \(C\) and \(\kappa\) alone: they carry no
factor of \(\kappa\) and no Gaussian dissipation, so they give no
smallness in the regimes (small velocity constant, small variance,
nearly constant core) in which the other terms of the identities
are small.  In particular \eqref{eq:c3v47-supY} is an explicit
supremum bound on the defect, valid when the derivative constants
satisfy \(C_\infty'+\frac{\sqrt3}2\kappa+\frac{\kappa^2}2<\frac12\), and not an exclusion; the
exclusion that Firewall~\ref{fw:c3v38-pressure} described as blocked
would have needed the pairing to be small, and it is not.  The
in-place status of V38 is unchanged (its statements are correct);
this record supersedes the wording ``not explicitly bounded'' by
``bounded only at supremum level''.
\end{firewall}

% =============================================================================
\section{Integrated V47 conclusion and next acquisition order}
\label{sec:c3v47-conclusion}
% =============================================================================

\begin{theorem}[What V47 proves]
\label{thm:c3v47-conclusion}
Relative to the two legacy cycles and V1--V46: the pointwise bounds
and the decomposition of the linearized pressure
(Lemma~\ref{lem:c3v47-decomposition}), the explicit bound on the
pairing (Theorem~\ref{thm:c3v47-pairing}) and on the Bernoulli work at
the extremal flux (Corollary~\ref{cor:c3v47-extremal-flux}), and the
revised firewall with the explicit supremum bound on the defect
(Theorem~\ref{thm:c3v47-revised}).  The full Navier--Stokes
stability/global-regularity theorem remains unproved.
\end{theorem}

\begin{proof}
The cited results.
\end{proof}

\begin{programme}[V48 acquisition order]
\label{prog:c3v48}
\begin{enumerate}[label=\textup{(AC\arabic*)},leftmargin=4.8em]
\item Apply the revised firewall back through the series: list every
      pairing recorded as non-explicit (V27 core bound, V31 constant
      weight, V33 Gaussian pressure form, V38 pairing) and state for
      each whether the supremum bound changes its conclusion; the
      expected answer is that only V38's is completed
      (\eqref{eq:c3v47-supY}) and the others are unaffected, the
      obstruction of V31 being decay, not pressure.
\item The extremal ledger with explicit constants: restate rows (6)
      and (7) of Theorem~\ref{thm:c3v39-ledger} and the flux row with
      the constants of this version, and record the supremum bound
      \eqref{eq:c3v47-supY} against the window bound of V33 and the
      Lamb--Bernoulli bound of V41.
\item The next compulsory companion audit is V50.
\end{enumerate}
\end{programme}

\begin{assessment}[Position after V47]
\label{ass:c3v47-position}
The pairing at the extremal flux is explicit; the pressure firewall
is revised to supremum level; the defect has an explicit supremum
bound for small derivative constants.  No full stability or
global-regularity claim is made.
\end{assessment}

% =============================================================================
\section*{Version V47 (third cycle) append-only record}
\addcontentsline{toc}{section}{Version V47 (third cycle) append-only record}
% =============================================================================

The complete V46 source of the third cycle was retained before this
continuation.  Its SHA--256 digest was
\begin{center}\scriptsize\ttfamily
06bc83cb7707fbb6731fc528b5d63c741f9bfa894e0a3f582ba57c2de7f2075f.
\end{center}
V47 proved the explicit pairing at the extremal flux and revised the
pressure firewall.  No full regularity claim is made.

% =============================================================================
\section{V48: the revised firewall applied through the series, and the extremal ledger with explicit constants}
\label{sec:c3v48-entry}
% =============================================================================

Items (AC1) and (AC2) of Programme~\ref{prog:c3v48} are carried out.
(AC1): the four places where the series recorded a pressure pairing
as non-explicit are re-examined under the revised firewall; only
the extremal-defect row of V38 changes, into the explicit supremum
bound of V47; the core bound of V27 and the Gaussian pressure form
of V33 were already explicit and are only made cruder by supremum
constants, and the obstruction of V31 is the absence of decay, not
the pressure.  (AC2): the extremal ledger's rows for the defect, the
cross term and the flux are restated with the constants of V47, and
the three bounds on the defect rate now in the series, the window
bound, the Lamb--Bernoulli bound and the supremum bound, are
compared: they are incomparable in general, each smallest in a
different regime of the class constants.  No global regularity
claim is made.  The next compulsory companion audit is V50.

% =============================================================================
\section{The revised firewall through the series}
\label{sec:c3v48-back}
% =============================================================================

\begin{proposition}[Item (AC1)]
\label{prop:c3v48-back}
Under Theorem~\ref{thm:c3v47-revised}:
\begin{enumerate}[label=\textup{(\roman*)},leftmargin=3.2em]
\item (V27.)  The bound on the mean pressure gradient on a nearly
      constant core, \eqref{eq:c3v27-two-radius}, is explicit in \(d\), \(C\)
      and \(\rho\); the supremum bound \(|\int\nabla P_VG|\le C_P|G|\) is
      cruder and carries no smallness in \(\eta\); the conclusion of V27
      (Corollary~\ref{cor:c3v27-decay}) is unchanged.
\item (V31.)  The obstruction to the constant weight and the maximum
      principle is the absence of pointwise decay of \(z\cdot V\) and of
      the boundary terms of \eqref{eq:c3v31-cutoff}; no pressure pairing
      is involved, and the negative record N-C3-3 stands.
\item (V33.)  \(\|\nabla P_V\|_G^2\le C_P^2|G|\) is a Gaussian form of the
      pressure part with a supremum constant; it does not depend on
      \(d\) and so does not serve the purpose of Firewall~\ref{fw:c3v33-pressure},
      which was a bound in terms of \(d\) and \(\kappa\); that firewall
      stands with its wording.
\item (V38.)  The pairing \(\langle Y,\nabla P'\rangle_G\) is bounded by \(C_{YP'}\),
      and \eqref{eq:c3v38-at-max} becomes the explicit supremum bound
      \eqref{eq:c3v47-supY} on the defect; the wording of
      Firewall~\ref{fw:c3v38-pressure} is superseded by
      Firewall~\ref{fw:c3v47-revised}.
\end{enumerate}
\end{proposition}

\begin{proof}
(i) \(|\int\nabla P_VG|\le\int|\nabla P_V|G\le C_P|G|\).  (ii) Proposition
\ref{prop:c3v31-constant}.  (iii) \eqref{eq:c3v47-gradP}.  (iv)
Theorem~\ref{thm:c3v47-revised}.
\end{proof}

% =============================================================================
\section{The extremal ledger with explicit constants}
\label{sec:c3v48-ledger}
% =============================================================================

\begin{theorem}[Rows (6), (7) and the flux row, explicit]
\label{thm:c3v48-rows}
At a homogeneous type-I point, with the constants of
Lemma~\ref{lem:c3v47-decomposition} and Theorem~\ref{thm:c3v47-revised}:
\begin{enumerate}[label=\textup{(\arabic*)},leftmargin=3.2em]
\item (Defect.)  At a maximizer of \(\|Y\|_G^2\):
      \(\nu\|\nabla Y\|_G^2+\frac12\|Y\|_G^2=-\langle Y,N'(V)Y\rangle_G\le(C_\infty'+\frac{\sqrt3}2\kappa+\frac{\kappa^2}2)\|Y\|_G^2+\frac\nu2\|\nabla Y\|_G^2+C_{YP'}\);
      hence \(\sup\|Y\|_G^2\le C_{YP'}/(\frac12-C_\infty'-\frac{\sqrt3}2\kappa-\frac{\kappa^2}2)\) when the
      denominator is positive, and \(\sup\nu\|\nabla Y\|_G^2\le2C_{YP'}+2(C_\infty'+\frac{\sqrt3}2\kappa+\frac{\kappa^2}2)\sup\|Y\|_G^2\)
      at that configuration.
\item (Cross term.)  At an extremum of \(\langle N,Y\rangle_G\):
      \(\langle N,\partial_\sigma Y\rangle_G=-\langle Y,N'(V)Y\rangle_G\), and both sides are bounded
      in absolute value by \((C_\infty'+\frac{\sqrt3}2\kappa+\frac{\kappa^2}2)\|Y\|_G^2+\frac\nu2\|\nabla Y\|_G^2+C_{YP'}\).
\item (Flux.)  At an extremum of \(j\): \(|\mathcal B|\le\bar J/(2\nu)+\kappa(d+3\varphi)^{1/2}\|Y\|_G\),
      with \(\bar J\) of \eqref{eq:c3v47-J}; and the flux itself satisfies
      \(|j|\le\kappa(d+4\varphi)+K_P(d+3\varphi)^{1/2}\) there as everywhere.
\end{enumerate}
\end{theorem}

\begin{proof}
(1) Theorem~\ref{thm:c3v38-defect}, Proposition~\ref{prop:c3v38-advective}
and \eqref{eq:c3v47-YP'}; the gradient bound by moving the
\(\|Y\|_G^2\) terms to the right.  (2) Proposition~\ref{prop:c3v38-cross}
with the same bound on \(\langle Y,N'(V)Y\rangle_G\).  (3)
Corollary~\ref{cor:c3v47-extremal-flux} and Theorem C2-V36-flux-bound.
\end{proof}

\begin{assessment}[Three bounds on the defect rate]
\label{ass:c3v48-three}
The series now bounds the defect rate \(\underline\delta\), the compactness
constant of the gate, in three explicit ways:
\begin{enumerate}[label=\textup{(\roman*)},leftmargin=3.2em]
\item the window bound \(\underline\delta\le10\kappa^2\nu C\) (V33), from
      \(\|N\|_G^2\) with the pressure gradient in unweighted \(L^2\);
\item the Lamb--Bernoulli bound \(\underline\delta\le\overline M^2\) (V41), with
      \(\overline M=\kappa(2C_{\mathcal L})^{1/2}+C_{\rm CZ}K_6^2I_6^{1/6}/(2\nu)+\frac\kappa2(4\nu C+3C_\Phi)^{1/2}\),
      from the pressure in Gaussian \(L^2\) through its \(L^3\) bound;
\item the supremum bound \(\underline\delta\le\sup\|Y\|_G^2\le C_{YP'}/(\frac12-C_\infty'-\frac{\sqrt3}2\kappa-\frac{\kappa^2}2)\)
      (V47), valid for small derivative constants, from the
      linearized pressure at supremum level.
\end{enumerate}
The first is quadratic in \(\kappa\) and vanishes with it; the second has
a \(\kappa\)-independent pressure term; the third depends on the
pointwise derivative constants up to order four and blows up as
the denominator vanishes.  None is comparable to another in
general, and each is the smallest in a regime: (i) for small \(\kappa\)
with \(C\) of order one, (ii) when the enstrophy \(C\) is large compared
with \(K_6^4/\nu^2\), (iii) when the derivative constants are small
compared with \(\frac12\).  The bounds are on the non-explicit side of
the gate, the defect rate; no lower bound on it is explicit (V31),
and the gate remains as stated in V39 and V46.
\end{assessment}

% =============================================================================
\section{Integrated V48 conclusion and next acquisition order}
\label{sec:c3v48-conclusion}
% =============================================================================

\begin{theorem}[What V48 establishes]
\label{thm:c3v48-conclusion}
Relative to the two legacy cycles and V1--V47: the review of the
pairings (Proposition~\ref{prop:c3v48-back}), the explicit rows
(Theorem~\ref{thm:c3v48-rows}) and the comparison of the three bounds
on the defect rate.  The full Navier--Stokes
stability/global-regularity theorem remains unproved.
\end{theorem}

\begin{proof}
The cited results.
\end{proof}

\begin{programme}[V49 acquisition order]
\label{prog:c3v49}
\begin{enumerate}[label=\textup{(AC\arabic*)},leftmargin=4.8em]
\item Consolidate V46--V48 as the correlation block: the
      energy--Hessian form of the pressure work on the configuration
      with its multipole sign structure, the explicit pairing at the
      extremal flux, the revised firewall, and the three bounds on
      the defect rate.
\item Fix the direction for V51--V55.  Two candidates: (a) the
      geometric gate of V46 tested on the configurations of the
      anchor chapter (a single energetic group with a diffuse far
      field): whether the multipole sign structure, combined with
      the count and the localization, bounds the pressure work on
      the anchor-alone event by an explicit multiple of the
      radial-defect masses, and what the momentum identity says
      about those masses; (b) the internal audit of the third cycle
      at its midpoint (V51--V55 as a re-derivation of the chapters
      V21--V48 in the manner of C2-V80--V87), since the cycle has
      corrected its own direction three times (V21, V22, V47) and a
      systematic check is due before V60.  Choose (b) first if any
      statement of V41--V48 is found to depend on a superseded one.
\item V50 compulsory companion audit and delivery.
\end{enumerate}
\end{programme}

\begin{assessment}[Position after V48]
\label{ass:c3v48-position}
The revised firewall changes only V38; the extremal rows are
explicit; the defect rate has three incomparable explicit upper
bounds.  No full stability or global-regularity claim is made.
\end{assessment}

% =============================================================================
\section*{Version V48 (third cycle) append-only record}
\addcontentsline{toc}{section}{Version V48 (third cycle) append-only record}
% =============================================================================

The complete V47 source of the third cycle was retained before this
continuation.  Its SHA--256 digest was
\begin{center}\scriptsize\ttfamily
d99f60d9162264836c514320d35aa9cccb76c044557525625ab4d892d4c96f20.
\end{center}
V48 applied the revised firewall through the series and made the
extremal rows explicit.  No full regularity claim is made.

% =============================================================================
\section{V49: the correlation block consolidated, and the midpoint audit ordered}
\label{sec:c3v49-entry}
% =============================================================================

Items (AC1) and (AC2) of Programme~\ref{prog:c3v49} are carried out.
The correlation block V46--V48 is consolidated.  For the direction,
the second candidate of V48 is chosen: the third cycle has corrected
its own direction three times (V21, V22, V47) and revised one
firewall (V47), and its chapters V21--V48 depend on one another
through explicit constants that were defined in one version and
used in others; before the cycle passes its midpoint a systematic
re-derivation is due, in the manner of the second cycle's internal
audit (C2-V80--V87).  The block V51--V55 is fixed as that audit,
with V55 the companion audit and delivery.  The geometric gate of
V46 on the anchor-alone event, the first candidate, is deferred to
V56--V60.  No global regularity claim is made.  The next compulsory
companion audit is V50.

% =============================================================================
\section{The correlation block}
\label{sec:c3v49-block}
% =============================================================================

\begin{theorem}[The correlation block, V44 and V46--V48]
\label{thm:c3v49-block}
At a homogeneous type-I point:
\begin{enumerate}[label=\textup{(\arabic*)},leftmargin=3.2em]
\item (Bilinear form, V44.)  The pressure work is
      \(\frac1{2\nu}\int q\psi_Y\) with \(q=|S|^2-\frac12|\Omega|^2\) and \(\psi_Y\) the dipole-type
      potential of the radial-defect density; and
      \(\mathcal B=\frac14j'-\frac1{2\nu}\int(P'+V\cdot Y)(z\cdot V)G\).
\item (Energy--Hessian form on the configuration, V46.)  On the
      trapped sub-windows,
      \(\frac1{2\nu}\int_U(P_V-c)(z\cdot Y)G=\frac1{2\nu}\int\chi_{U'}(V\otimes V):\nabla^2\psi_U+E_{\rm far}\),
      \(|E_{\rm far}|\le C_{11}\varepsilon_{\rm th}\|Y\|_G\), decomposing over the cells of \(U\);
      cross terms have the point-mass form
      \(\frac{m_\beta}{4\pi\rho^3}(3(V\cdot\hat e)^2-|V|^2)+O(\mu_\beta|V|^2\rho^{-4})\): negative
      for swirl about cells of outward radial defect and for radial
      flow about cells of inward radial defect; self-terms have the
      isotropic part \(-\frac13\int\phi_\alpha|V|^2f_\alpha\) and an unsigned
      anisotropic part.
\item (Explicit pairing, V47.)  \(|\int P'(z\cdot V)G|\le\bar J\) with
      \(P'=-P_V-\frac12z\cdot\nabla P_V+P_t\); at extremal flux
      \(|\mathcal B|\le\bar J/(2\nu)+\kappa(d+3\varphi)^{1/2}\|Y\|_G\); the bypass identity is
      not established in the class.
\item (Revised firewall, V47--V48.)  Every Gaussian pressure pairing
      is explicit at supremum level; \(\sup\|Y\|_G^2\le C_{YP'}/(\frac12-C_\infty'-\frac{\sqrt3}2\kappa-\frac{\kappa^2}2)\)
      when the denominator is positive; only V38's row changes; the
      defect rate has three incomparable explicit upper bounds.
\end{enumerate}
\end{theorem}

\begin{proof}
Propositions~\ref{prop:c3v44-bilinear}, \ref{prop:c3v44-flux},
\ref{prop:c3v46-far}, \ref{prop:c3v46-multipole}, \ref{prop:c3v46-self},
\ref{prop:c3v48-back}; Theorems~\ref{thm:c3v46-pairing},
\ref{thm:c3v47-pairing}, \ref{thm:c3v47-revised}, \ref{thm:c3v48-rows};
Corollary~\ref{cor:c3v47-extremal-flux}; Firewall~\ref{fw:c3v47-no-bypass}.
\end{proof}

\begin{assessment}[What the block settles]
\label{ass:c3v49-settles}
The gate's nonlocal part is a finite bilinear form over the
configuration's cells, with an explicit geometric sign structure and
explicit far corrections, and every pairing that enters it or the
extremal rows is explicitly bounded.  What the block does not do is
give the form a sign: the multipole structure fixes the sign of each
cross term in terms of the orientation of the velocity relative to a
radial-defect cell, and the class does not fix that orientation.  The
block also corrected the series twice: the strain form of V44 was
replaced by the energy form for localization (V46), and the standing
firewall of V38 was revised (V47).  These corrections are the reason
for the audit ordered below.
\end{assessment}

% =============================================================================
\section{Direction for V51--V55: the midpoint internal audit}
\label{sec:c3v49-direction}
% =============================================================================

\begin{assessment}[The audit]
\label{ass:c3v49-audit}
The block V51--V55 re-derives, in order, the results of V21--V48 that
carry explicit constants or that later versions use, checks each
against the statements it depends on, and records any correction in
place with a note in the version that makes it, as V34 did for V33.
The order follows the chapters.  V51: the explicit floor
(V21--V24) and the variance chapter (V26--V28): the one-level
criterion, the growth lemma, the floor \(v_{\rm exp}\) and its closed form,
the \(H^1\) closeness, the nearly constant core and the window floor;
in particular the constants \(\rho_k^{(6)}\), \(v^{(k)}\), \(\eta_{\rm c}\), \(\tau_*\) are
recomputed.  V52: the defect block (V31--V34): the NRS identity with
a defect, the mean-square identity and its explicit sides, the
window form, the extremal lemma and the bounds \(\mathcal L_\pm\).  V53: the
extremal ledger (V36--V38) and the angle chapter (V41--V43): the rows
and their constants, \(B_k^{\sup}\), \(M_{NY}\), \(I_6\), the reduced gates.
V54: the correlation block (V44, V46--V48): the energy--Hessian
form, the multipole expansion, the pointwise constants of V47 and
the three bounds on the defect rate; and the list of every
statement of V21--V48 that cites a superseded wording (V38's
firewall) with its current reading.  V55: the compulsory companion
audit, the audit's summary and delivery.  No global regularity
claim is made.
\end{assessment}

% =============================================================================
\section{Integrated V49 conclusion and next acquisition order}
\label{sec:c3v49-conclusion}
% =============================================================================

\begin{theorem}[What V49 establishes]
\label{thm:c3v49-conclusion}
Relative to the two legacy cycles and V1--V48: the consolidated
correlation block (Theorem~\ref{thm:c3v49-block}) and the audit
order.  The full Navier--Stokes stability/global-regularity theorem
remains unproved.
\end{theorem}

\begin{proof}
The cited results.
\end{proof}

\begin{programme}[V50 acquisition order]
\label{prog:c3v50}
\begin{enumerate}[label=\textup{(AC\arabic*)},leftmargin=4.8em]
\item The compulsory companion audit: read the current SM and YM
      notes and record their digests and decision.
\item Record the position after fifty versions; delivery.
\end{enumerate}
\end{programme}

\begin{assessment}[Position after V49]
\label{ass:c3v49-position}
The correlation block is consolidated; the midpoint internal audit
is ordered for V51--V55.  No full stability or global-regularity
claim is made.
\end{assessment}

% =============================================================================
\section*{Version V49 (third cycle) append-only record}
\addcontentsline{toc}{section}{Version V49 (third cycle) append-only record}
% =============================================================================

The complete V48 source of the third cycle was retained before this
continuation.  Its SHA--256 digest was
\begin{center}\scriptsize\ttfamily
46514d751156a6bc329ab5385713660e326e20559366bc158e2bfffce441a517.
\end{center}
V49 consolidated the correlation block and ordered the audit.  No
full regularity claim is made.

% =============================================================================
\section{V50: compulsory companion audit and the position after fifty versions}
\label{sec:c3v50-entry}
% =============================================================================

This is the fiftieth version of the third cycle.  The compulsory
review of the two companion programmes is performed first, and the
position at the midpoint of the cycle is recorded with the plan for
the internal audit V51--V55.  No global regularity claim is made.
The next compulsory companion audit is V55.

% =============================================================================
\section{Compulsory V50 review of the current companion notes}
\label{sec:c3v50-companion-audit}
% =============================================================================

The current companion files in \texttt{latex\_notes} were inspected
read-only.  Both programmes have moved since the V45 audit: the SM
programme from V33 to V34 of its seventh series, the YM programme
from V11 to V12 of its sixth series.  The frozen checkpoints are
\begin{align}
 &\texttt{Renewal\_SM\_Einstein\_Continuum\_Research\_Notes\_V34.tex},
 \notag\\[-2pt]
 &\qquad
 \texttt{488c5745c8042ea4c3bd6a85cad1a4a0a843a95e49b1f68469bb11c34cf6a1ea},
 \label{eq:c3v50-SM-checkpoint}\\
 &\texttt{Renewal\_Yang\_Mills\_Mass\_Gap\_Research\_Notes\_V12.tex},
 \notag\\[-2pt]
 &\qquad
 \texttt{d5d429e3e963e6ea92321b9fb53e5ad339247ca04091292d8ac02350c5e07d06}.
 \label{eq:c3v50-YM-checkpoint}
\end{align}
(SM V34 has 13938 lines; YM V12 has 5452 lines.)

\emph{SM V34.}  The SM programme treated the regular simplicial bulk
gap: regular edgewise bulk symbols, the removal of every regular
bulk doubler, and a codimension-one face operator.

\emph{YM V12.}  The YM programme obtained an exact Cayley--Pr\"ufer
degree law: an outer weighted Cayley law, why its proposed tail
route cannot be uniform, the factorial that survives finite-family
aggregation, and a conditional closure of its rank-balanced row.

\begin{theorem}[V50 companion-audit decision]
\label{thm:c3v50-companion-decision}
Nothing is imported from SM V34 or YM V12.  Neither bears on the
third cycle's chapters or on the gates.
\end{theorem}

\begin{proof}
The SM results concern symbols and doublers of lattice Dirac
operators; the YM results concern degree laws of trees in a cluster
expansion.  None is a Navier--Stokes statement, and none concerns
the ledger, the extremal method, the Lamb--Bernoulli form or the
correlation block.
\end{proof}

% =============================================================================
\section{Position at the midpoint of the third cycle}
\label{sec:c3v50-position}
% =============================================================================

\begin{assessment}[Position after fifty versions of the third cycle]
\label{ass:c3v50-position-fifty}
Nine chapters are written: cluster energy and the band (V2--V10),
localization (V11--V14), the anchor (V16--V19), the explicit floor
(V21--V24), the variance chapter (V26--V29), the defect block
(V31--V34), the extremal ledger (V34--V39), the angle chapter
(V41--V44), and the correlation block (V46--V49).  The cycle's
results are of three kinds.  Structural: the type-I gate at
homogeneous points is a mean inequality on a bounded neighbourhood
of an anchor group (V13, V17), an angle statement that norms cannot
decide (V39), and, in its nonlocal part, a finite bilinear form over
the configuration's cells with an explicit multipole sign structure
(V46).  Quantitative: the variance floor is explicit by two
mechanisms (V23, V28); the defect rate has three explicit upper
bounds (V33, V41, V47) and no explicit lower bound (N-C3-3); the
mean identities of the second cycle have pointwise forms at
extremal configurations with explicit two-sided bounds (V34--V38,
V47); every Gaussian pressure pairing is explicit at supremum level
(V47).  Negative: N-C3-1 (concentration gives no floor), N-C3-2
(the mean spectrum gives no mean floor), N-C3-3 (no quantitative
NRS in the class), and the unjustified bypass (V47).  The cycle has
corrected its own direction at V21, V22 and V47 and revised one
firewall; the internal audit V51--V55 (Assessment~\ref{ass:c3v49-audit})
re-derives V21--V48 before the second half of the cycle, which is
planned to begin with the geometric gate on the anchor-alone event
(V56--V60).  No global regularity claim is made.
\end{assessment}

% =============================================================================
\section{Integrated V50 conclusion and next acquisition order}
\label{sec:c3v50-conclusion}
% =============================================================================

\begin{theorem}[What V50 establishes]
\label{thm:c3v50-conclusion}
Relative to the two legacy cycles and V1--V49: the companion-audit
decision (Theorem~\ref{thm:c3v50-companion-decision}) and the
midpoint position.  The full Navier--Stokes stability/global-regularity
theorem remains unproved.
\end{theorem}

\begin{proof}
The cited records.
\end{proof}

\begin{programme}[V51 acquisition order]
\label{prog:c3v51}
\begin{enumerate}[label=\textup{(AC\arabic*)},leftmargin=4.8em]
\item Audit of the explicit floor and the variance chapter
      (V21--V24, V26--V28): re-derive the one-level and two-level
      criteria, the dominant-level bound, the cut-off constant's
      variance, the growth lemma with its constants, the floor
      \(v_{\rm exp}\) with the closed-form case, the \(H^1\) closeness, the
      nonlinear mean on the nearly constant core with \(\eta_{\rm c}\), the
      locked level and the window floor with \(\tau_*\); record every
      constant recomputed and every discrepancy.
\item Record the dependency list of V21--V28 on the second cycle's
      statements (C2-V31, C2-V34, C2-V36, C2-V40--V41) and confirm
      each citation's content against the legacy file.
\item The next compulsory companion audit is V55.
\end{enumerate}
\end{programme}

\begin{assessment}[Position after V50]
\label{ass:c3v50-position}
The companion audit is recorded; the internal audit begins.  No
full stability or global-regularity claim is made.
\end{assessment}

% =============================================================================
\section*{Version V50 (third cycle) append-only record}
\addcontentsline{toc}{section}{Version V50 (third cycle) append-only record}
% =============================================================================

The complete V49 source of the third cycle was retained before this
continuation.  Its SHA--256 digest was
\begin{center}\scriptsize\ttfamily
bc655bf8463cfa738cd57e412552468893696c7922cde74300d3807651cbb7e7.
\end{center}
V50 performed the compulsory companion audit and recorded the
midpoint position.  No full regularity claim is made.

% =============================================================================
\section{V51: internal audit, part 1: the explicit floor and the variance chapter}
\label{sec:c3v51-entry}
% =============================================================================

Items (AC1) and (AC2) of Programme~\ref{prog:c3v51} are carried out.
The statements of V21--V24 and V26--V28 that carry explicit constants
or are used later were re-derived from their hypotheses and their
constants recomputed.  No correction is required.  Two points are
recorded for precision: the size remark of V23 rests on the
numerical value \(144(4\pi/3)^{-2/3}\pi\approx174.1\), and the bound
``below \(e^{-170}\)'' uses the four units of margin; and the
window-floor theorem of V28 uses the three hypotheses of the decay
corollary of V27 together, all supplied by the window hypothesis.
The dependency list on the second cycle is recorded and each
citation confirmed against the legacy file.  No global regularity
claim is made.  The next compulsory companion audit is V55.

% =============================================================================
\section{Re-derivations}
\label{sec:c3v51-rederivations}
% =============================================================================

\begin{theorem}[Audit of V21--V24]
\label{thm:c3v51-floor}
The following were re-derived and hold as stated.
\begin{enumerate}[label=\textup{(\arabic*)},leftmargin=3.2em]
\item (V21.)  With \(r\in[\mu_{k_0},\mu_{k_0+1})\) and spacing \(\frac12\): the nearest
      level is within \(\frac14\), all others at least \(\frac14\), those outside
      \(\{k_0,k_0+1\}\) at least \(\frac12\); \(\mathrm{Var}\ge(\mu_{k_1}-r)^2\hat a_{k_1}+\frac1{16}(1-\hat a_{k_1})\)
      and \(\mathrm{Var}\ge\frac14(1-\hat a_{k_0}-\hat a_{k_0+1})\).  \(\mu_{k_1}\le r+\frac14\le C_r+\frac14\)
      gives \(k_1\le2C_r-\frac12\), and Chebyshev gives
      \(\sum_{k\ge K}\hat a_k\le r/\mu_K=2r/(1+K)\).  \(C_r=(4\nu C+2C_\Phi)/(4c_\Phi)\) is
      \(2C_{\mathcal L}/c_\Phi\) with \(C_{\mathcal L}=(2\nu C+C_\Phi)/4\) of V33, consistent.
\item (V22.)  For \(P_\rho=\operatorname{curl}(\frac12\chi_\rho(c_0\times z))\) on the annulus:
      \(|P_\rho|\le(1+A)|c_0|\), \(|\nabla P_\rho|\le\frac52A|c_0|/\rho\le3A|c_0|/\rho\),
      \(|\nabla^2P_\rho|\le3A|c_0|/\rho^2\le4A|c_0|/\rho^2\) (the products expanded
      term by term); \((L_-+\frac12)P_\rho\) bounded by \(A|c_0|(3+4\nu/\rho^2)\);
      tail \(\int_{|z|>\rho}G\le8\pi\nu(\rho+2\nu/\rho)e^{-\rho^2/(4\nu)}\) by the
      integration by parts and \(e^{-s^2/(4\nu)}\le(s/\rho)e^{-s^2/(4\nu)}\) for
      \(s\ge\rho\); the numerical reduction for \(\rho^2\ge12\nu\)
      (\((3+\frac13)^2\le12\), \(\rho+2\nu/\rho\le2\rho\), \(384\pi\nu/(4\pi\nu)^{3/2}=96(4\pi\nu)^{-1/2}\))
      gives \eqref{eq:c3v22-cutoff-var}.  Growth lemma: \(\gamma_0=4\pi/(3|G|)\);
      for \(Az\), \(\|Az\|_G^2=2\nu|G||A|_F^2\) and \(\int_{B_\rho}|Az|^2=\frac{4\pi}{15}\rho^5|A|_F^2\),
      so \(\gamma_1=4\pi/(30\nu|G|)\).
\item (V23.)  H\"older gives \(\int_{B_\rho}|\widehat V|^2\le(4\pi/3)^{2/3}\rho^2K_6^2/c_\Phi\); the
      comparison at \(\rho=\rho_k^*\) gives \(\eta\ge\frac18\gamma_k\rho^{2k+3}e^{-\rho^2/(4\nu)}\)
      and \(v^{(k)}=\gamma_k\rho^{2k+3}e^{-\rho^2/(4\nu)}/128\).  Closed forms:
      \(\rho_0^{(6)}=8(4\pi/3)^{-1/3}|G|K_6^2/c_\Phi\), \(\gamma_0/128=\pi/(96|G|)\);
      \((\rho_1^{(6)})^3=120(4\pi/3)^{2/3}\cdot2\nu|G|K_6^2/(4\pi c_\Phi)\), \(\gamma_1/128=\pi/(960\nu|G|)\).
      Size: with \(x:=(\rho_0^*)^2/(4\nu)\ge144(4\pi/3)^{-2/3}\pi\approx174.1\),
      \(v^{(0)}=x^{3/2}e^{-x}/(96\pi^{1/2})\), decreasing in \(x\) on this range,
      at most \(13.6\,e^{-174}<e^{-170}\); the remark's conclusion holds with
      this computation.
\item (V24.)  \(\nu\|\nabla\Pi_k\widehat V\|_G^2=\frac k2\hat a_k\) and the \(G\)-orthogonality of
      the gradients of distinct eigenfunctions of the symmetric
      operator give \(\nu\|\nabla(\widehat V-P)\|_G^2=\sum_{k\ne k_1}\frac k2\hat a_k\le(\frac18+2C_r^2)^{1/2}\eta\);
      \(|r-\mu_{k_1}|\le\mathrm{Var}^{1/2}\eta^{1/2}+\frac14\eta=\frac12\eta\).
\end{enumerate}
\end{theorem}

\begin{proof}
Each item is the re-derivation described; the computations were
carried out independently of the original text and agree with it.
\end{proof}

\begin{theorem}[Audit of V26--V28]
\label{thm:c3v51-variance}
The following were re-derived and hold as stated.
\begin{enumerate}[label=\textup{(\arabic*)},leftmargin=3.2em]
\item (V26.)  \(B_0^{1/2}\ge2\sqrt3\kappa c_\Phi^{1/2}+\frac{\sqrt3}2K_P\); \(B_0<c_\Phi\) forces
      \(\kappa<1/(2\sqrt3)\) and \(c_\Phi^{1/2}(1-2\sqrt3\kappa)>\frac{\sqrt3}2K_P\); with
      \(c_\Phi^{1/2}\le K_P/(\sqrt3m_\kappa)\), \(m_\kappa=(2-4\kappa)/3\), the two are compatible
      only if \(1-2\kappa<1-2\sqrt3\kappa\), false for \(\kappa>0\).
\item (V27.)  \(\int\nabla P_VG=-\int q\nabla g\) with \(|\nabla g|=M(|y|)/(4\pi|y|^2)\le\min(|y|/3,|G|/(4\pi|y|^2))\);
      on the core, with \(\|\widetilde V\|_G^2\le\varphi\eta\), \(\nu\|\nabla\widetilde V\|_G^2\le\varphi c_r\eta\),
      \(\||z|\widetilde V\|_G^2\le4\nu\varphi\eta(3+4c_r)\): the advective mean is at most
      \(\varphi\nu^{-1/2}|G|^{-1}[\sqrt6\eta^{1/2}+(3+4c_r)^{1/2}\eta]\) and the pressure
      mean at most \(\rho c_r\varphi e^{\rho^2/(4\nu)}\eta/(3\nu|G|)+C/(4\pi\rho^2)\); the
      four entries of \eqref{eq:c3v27-etac} make each term at most
      \(\frac1{32}(\varphi/|G|)^{1/2}\), and \(|\bar V|\ge(\varphi/(2|G|))^{1/2}\) gives
      \(|\overline{N(V)}|\le\frac18(\varphi/|G|)^{1/2}\le\frac14|\bar V|\).
\item (V28.)  A change of dominant level forces \(r\) through a
      midpoint, where \(\mathrm{Var}\ge\frac1{16}\).  On a window with
      \(\mathrm{Var}<v_{\rm win}\): \(\mathrm{Var}<\frac1{32}\), \(\eta<\eta_{\rm c}\), and \(k_1\) locked, the
      three hypotheses of Corollary~\ref{cor:c3v27-decay}; with
      \(\tau_*=2\log(4C_\Phi/c_\Phi)\), \(a_0\) falls to at most \(c_\Phi/4<\frac12c_\Phi\le(1-\eta)\varphi\).
\end{enumerate}
\end{theorem}

\begin{proof}
As above.
\end{proof}

% =============================================================================
\section{Dependencies on the second cycle}
\label{sec:c3v51-dependencies}
% =============================================================================

\begin{proposition}[Citations of V21--V28 confirmed]
\label{prop:c3v51-citations}
The following statements of the second cycle are cited in V21--V28
and were confirmed, by inspection of the legacy file, to contain
what is cited: C2-V8-spectrum (spectrum of \(L_-\), invariance of the
divergence-free subspace); C2-V12 (\(\varphi\), \(\mathcal L\) continuous, \(C^1\) in
\(\sigma\); mean identities); C2-V17-mean-profile (the momentum identity
\(\frac{\dd}{\dd\sigma}\bar V=-\frac12\bar V-\overline{N(V)}\) with
\(\overline{N(V)}=\frac1{2\nu|G|}\int[(z\cdot V)V+P_Vz]G\)); C2-V21-closure
(\(R=8\nu\mathcal L-4\nu^2\mathcal E\)); C2-V31-degree (spectral form of the variance);
C2-V33 (\(\frac12\le r\le C_r\)); C2-V34-eigen and C2-V34-variance
(polynomial eigenfunctions, compactness floor \(v_0\), mean push);
C2-V36 (\(K_6\), the Gaussian Hardy inequality \(X^2\le4\nu(d+3\varphi)\), the flux
bound, the mean constraint for \(\kappa<\frac12\)); C2-V40-modewise and
C2-V40-low-modes (\(a_k'=-2\mu_ka_k-2n_k\), \(\langle n_k\rangle=-\mu_k\langle a_k\rangle\),
\(V_0=\bar V\), \(a_0=|G||\bar V|^2\), \(n_0=|G|\bar V\cdot\overline N\)); C2-V41-modes and
C2-V41-bound and C2-V41-spectrum (\(\nu\|\nabla V_k\|_G^2=\frac k2a_k\),
\(\int|z|^2|V_k|^2G\le4\nu(3+2k)a_k\), the modewise flux bound, \(\langle a_k\rangle\le B_k\)).
\end{proposition}

\begin{proof}
Inspection of the file \texttt{Renewal\_NS\_Stability\_Research\_Notes\_V100\_f2.tex}
at the labelled statements (digest recorded in V1 of this cycle).
\end{proof}

% =============================================================================
\section{Integrated V51 conclusion and next acquisition order}
\label{sec:c3v51-conclusion}
% =============================================================================

\begin{theorem}[What V51 establishes]
\label{thm:c3v51-conclusion}
The re-derivations (Theorems~\ref{thm:c3v51-floor}, \ref{thm:c3v51-variance})
and the confirmed citations (Proposition~\ref{prop:c3v51-citations}):
V21--V24 and V26--V28 stand without correction.  The full
Navier--Stokes stability/global-regularity theorem remains unproved.
\end{theorem}

\begin{proof}
The cited records.
\end{proof}

\begin{programme}[V52 acquisition order]
\label{prog:c3v52}
\begin{enumerate}[label=\textup{(AC\arabic*)},leftmargin=4.8em]
\item Audit of the defect block (V31--V34): re-derive the vector
      identities and the NRS identity with a defect, the Gaussian
      adjoint computation, the mean-square identity, the Plancherel
      bound on the pressure gradient, the pointwise and window forms,
      the extremal lemma and the bounds \(\mathcal L_\pm\).
\item Confirm the citations C2-V2-second, C2-V3-cross, C2-V4 (compact
      family, derivative bounds), C2-V6-m, C2-V7-defect-energy,
      C2-V9-B, C2-V12-mean-identities.
\item The next compulsory companion audit is V55.
\end{enumerate}
\end{programme}

\begin{assessment}[Position after V51]
\label{ass:c3v51-position}
The floor and variance chapters pass the audit unchanged.  No full
stability or global-regularity claim is made.
\end{assessment}

% =============================================================================
\section*{Version V51 (third cycle) append-only record}
\addcontentsline{toc}{section}{Version V51 (third cycle) append-only record}
% =============================================================================

The complete V50 source of the third cycle was retained before this
continuation.  Its SHA--256 digest was
\begin{center}\scriptsize\ttfamily
9ef8d1759abaac883a4f004b9899c7a81859c031418f06788cf5d63ef58320ab.
\end{center}
V51 audited V21--V24 and V26--V28.  No full regularity claim is made.

% =============================================================================
\section{V52: internal audit, part 2: the defect block}
\label{sec:c3v52-entry}
% =============================================================================

Items (AC1) and (AC2) of Programme~\ref{prog:c3v52} are carried out.
The statements of V31--V34 were re-derived from their hypotheses and
their constants recomputed.  No correction is required.  One point
is recorded for precision: the extremal lemma of V34 uses that the
orbit of every element of the family is two-sided in log-time, which
the second cycle's construction of the family supplies (rescalings
by every positive factor lie in the family), and the derivative of
the observable at the extremal configuration vanishes for that
reason and not merely because the extremum is attained.  The
citations of the second cycle are confirmed.  No global regularity
claim is made.  The next compulsory companion audit is V55.

% =============================================================================
\section{Re-derivations}
\label{sec:c3v52-rederivations}
% =============================================================================

\begin{theorem}[Audit of V31--V34]
\label{thm:c3v52-defect}
The following were re-derived and hold as stated.
\begin{enumerate}[label=\textup{(\arabic*)},leftmargin=3.2em]
\item (V31.)  \(\nabla(a\cdot b)=(a\cdot\nabla)b+(b\cdot\nabla)a+a\times\operatorname{curl}b+b\times\operatorname{curl}a\)
      gives the two vector identities; substitution gives
      \(\nabla\Pi=\nu\Delta V-\Omega\times W-Y\); the divergence, with
      \(\operatorname{div}(\Omega\times W)=W\cdot\operatorname{curl}\Omega-\Omega\cdot\operatorname{curl}W=-W\cdot\Delta V-|\Omega|^2\)
      and \(\operatorname{div}Y=0\), gives \(\Delta\Pi=W\cdot\Delta V+|\Omega|^2\); eliminating \(\Delta V\)
      through the gradient form gives \(\nu\Delta\Pi-W\cdot\nabla\Pi=\nu|\Omega|^2+W\cdot Y\).
      The adjoint on the Gaussian: \(\nu\Delta G=(\frac{|z|^2}{4\nu}-\frac32)G\),
      \(W\cdot\nabla G=-\frac{V\cdot z+\frac12|z|^2}{2\nu}G\), \(\operatorname{div}W=\frac32\), sum
      \(-\frac{z\cdot V}{2\nu}G\); the expansion of \(-\frac1{2\nu}\int\Pi(z\cdot V)G\) as
      \(-\frac14j-\frac R{4\nu}\) with the flux formula, and \(R=8\nu\mathcal L-4\nu^2\mathcal E\),
      reduce the identity to \(\langle V,Y\rangle_G=-2\mathcal L-\frac14j\), the energy
      identity.  The constant weight: \(L^*1=\frac32\), and the cut-off
      boundary terms are \(o(R^2)\) and \(o(R^{7/2})\) under the class bounds.
\item (V32.)  \(\|Y\|^2+\langle N,Y\rangle=\|L_-V\|^2-\langle L_-V,N\rangle\) pointwise and the
      mean identity give \(\langle\langle L_-V,N\rangle\rangle=\langle\|L_-V\|^2\rangle\), hence
      \(\langle\|N\|^2\rangle=\langle\|L_-V\|^2\rangle+\langle\|Y\|^2\rangle\).  \(\|L_-V\|_G^2=\varphi\sum\mu_k^2\hat a_k=(r^2+\mathrm{Var})\varphi\).
      Plancherel: \(\widehat{\partial_lP_V}=-(\xi\otimes\xi/|\xi|^2):\widehat{\partial_l(V\otimes V)}\), and the
      projection has Frobenius norm one, so \(\|\partial_lP_V\|_2\le\|\partial_l(V\otimes V)\|_2\le2C_\infty\|\partial_lV\|_2\).
      \(\|N\|_G^2\le2C_\infty^2\|\nabla V\|_G^2+2\|\nabla P_V\|_{L^2}^2\le\frac{\kappa^2}2d+8\kappa^2\nu C\).
\item (V33.)  \(2\mathcal L'=-2\langle Y,L_-V\rangle=\|N\|^2-\|L_-V\|^2-\|Y\|^2\); integration over
      a window and \(\frac14c_\Phi\le\mathcal L\le C_{\mathcal L}=(2\nu C+C_\Phi)/4\) give the window
      form with \(2C_{\mathcal L}/L\); dropping the negative term and using
      \(\frac12\langle d\rangle_L\le2\nu C\) gives \(\langle\|Y\|^2\rangle_L\le10\kappa^2\nu C+2C_{\mathcal L}/L\) and
      \(\underline\delta\le10\kappa^2\nu C\).
\item (V34.)  At a maximizer of \(\mathcal L\): \(0\le\kappa^2(4\mathcal L-\varphi+8\nu C)-4\mathcal L^2/\varphi\),
      whose larger root is \(\frac\varphi2[\kappa^2+(\kappa^4+\kappa^2(8\nu C-\varphi)/\varphi)^{1/2}]\le\frac{\kappa^2\varphi}2+\frac12(\kappa^4\varphi^2+8\kappa^2\nu C\varphi)^{1/2}\),
      increasing in \(\varphi\); at a minimizer, \(\frac14\varphi\le\kappa^2(4\mathcal L-\varphi+8\nu C)\)
      gives \(\mathcal L\ge\frac\varphi4(1+\frac1{4\kappa^2})-2\nu C\), increasing in \(\varphi\).  The
      positivity condition \(8\mathcal L_--2C_\Phi>0\) is \(\kappa^2<c_\Phi/(4(8\nu C+C_\Phi-c_\Phi))\).
\end{enumerate}
\end{theorem}

\begin{proof}
Each item is the re-derivation described; the computations agree
with the original text.
\end{proof}

\begin{remark}[The extremal lemma's hypothesis]
\label{rem:c3v52-two-sided}
Lemma~\ref{lem:c3v34-extremal} concludes \(F'=0\) at an attaining
element from the two-sidedness of the orbit: for \(W^*\in\mathcal F(x_0)\)
and every \(\tau\in\R\), \(\Theta_\tau W^*\in\mathcal F(x_0)\) (C2-V3--V4: the family is the
set of local limits of rescalings, closed under rescaling by every
positive factor).  If the family were invariant only under \(\tau\ge0\),
an extremum at \(\sigma=0\) would give only a one-sided sign of the
derivative.  The audit confirms the two-sidedness from the second
cycle's definition, and records that every extremal row of
V34--V38 and V43, V47 depends on it.
\end{remark}

% =============================================================================
\section{Dependencies on the second cycle}
\label{sec:c3v52-dependencies}
% =============================================================================

\begin{proposition}[Citations of V31--V34 confirmed]
\label{prop:c3v52-citations}
Confirmed by inspection of the legacy file: C2-V2-second
(\(\frac{\dd}{\dd\sigma}\mathcal L=-\|Y\|_G^2-\langle N,Y\rangle_G\), \(\mathcal L\) \(C^1\) in \(\sigma\));
C2-V3-cross (Lamb--Bernoulli form); C2-V4 (compactness of the family,
the uniform derivative bounds, closure under rescaling); C2-V6-m
(the windowed defect and the rate \(\underline\delta\)); C2-V7-defect-energy
(the defect's Gaussian energy identity and the form of \(N'(V)Y\));
C2-V9-B (alternative (B) at homogeneous points); C2-V12-mean-identities
(the three mean identities and the exact signs); C2-V21-closure;
C2-V32 (\(j=\frac2\nu\int\Pi(z\cdot V)G\)); C2-V36 (flux bound, \(K_6\), \(K_P\), mean
constraint, \(\varepsilon_{\rm reg}\) in Firewall C2-V36-AC2).
\end{proposition}

\begin{proof}
Inspection at the labelled statements.
\end{proof}

% =============================================================================
\section{Integrated V52 conclusion and next acquisition order}
\label{sec:c3v52-conclusion}
% =============================================================================

\begin{theorem}[What V52 establishes]
\label{thm:c3v52-conclusion}
The re-derivations (Theorem~\ref{thm:c3v52-defect}) and the confirmed
citations (Proposition~\ref{prop:c3v52-citations}): V31--V34 stand
without correction.  The full Navier--Stokes stability/global-regularity
theorem remains unproved.
\end{theorem}

\begin{proof}
The cited records.
\end{proof}

\begin{programme}[V53 acquisition order]
\label{prog:c3v53}
\begin{enumerate}[label=\textup{(AC\arabic*)},leftmargin=4.8em]
\item Audit of the extremal ledger (V36--V38) and the angle chapter
      (V41--V43): re-derive the rows and their constants, in
      particular \(B_k^{\sup}\) and its \(1/k\) form, \(M_{NY}\), \(I_6\), the
      core bounds of V42 and the reduced gates of V43; apply any
      correction in place with a note.
\item Confirm the citations C2-V17, C2-V31-identity, C2-V32-push,
      C2-V36-hardy, C2-V40, C2-V41 (including the step of C2-V41
      giving the \(1/k\) form of the mean spectrum).
\item The next compulsory companion audit is V55.
\end{enumerate}
\end{programme}

\begin{assessment}[Position after V52]
\label{ass:c3v52-position}
The defect block passes the audit unchanged.  No full stability or
global-regularity claim is made.
\end{assessment}

% =============================================================================
\section*{Version V52 (third cycle) append-only record}
\addcontentsline{toc}{section}{Version V52 (third cycle) append-only record}
% =============================================================================

The complete V51 source of the third cycle was retained before this
continuation.  Its SHA--256 digest was
\begin{center}\scriptsize\ttfamily
fd510fcc5a118c96be86ce97aebc82c23c8cc23dd2fa89b8c659b7f61e23920c.
\end{center}
V52 audited V31--V34.  No full regularity claim is made.

% =============================================================================
\section{V53: internal audit, part 3: the extremal ledger and the angle chapter}
\label{sec:c3v53-entry}
% =============================================================================

Items (AC1) and (AC2) of Programme~\ref{prog:c3v53} are carried out.
The statements of V36--V38 and V41--V43 were re-derived and their
constants recomputed.  One correction was found and applied in
place: the constant of the \(1/k\) form of the extremal spectrum in
V37 was justified by the inequality \((3+2k)^{1/2}\le(3k)^{1/2}\), which
fails for \(k=1,2\); the correct reduction uses \((5k)^{1/2}\) and gives
the constant \(20[\kappa C_\Phi^{1/2}+\kappa(4\nu C+3C_\Phi)^{1/2}+\frac14K_P]^2\) in place of
\(8[\dots+K_P]^2\).  The same step occurs in the second cycle's mean
spectrum (Theorem C2-V41-spectrum), whose constant \(C_{\rm sp}\) should
read \(20[\kappa\langle\varphi\rangle^{1/2}+\kappa\langle d+3\varphi\rangle^{1/2}+\frac14K_P]^2\); the legacy
file is frozen and the correction is recorded here.  No statement of
this cycle depends on the value of either constant beyond its
\(O(1/k)\) form, and none is affected.  Everything else stands.  No
global regularity claim is made.  The next compulsory companion
audit is V55.

% =============================================================================
\section{The correction}
\label{sec:c3v53-correction}
% =============================================================================

\begin{proposition}[Corrected \(1/k\) form of the extremal spectrum]
\label{prop:c3v53-correction}
For \(k\ge1\), with \(\beta_k\) of \eqref{eq:c3v37-Bsup},
\(\beta_k\le(5k)^{1/2}[\kappa C_\Phi^{1/2}+\kappa(4\nu C+3C_\Phi)^{1/2}+\frac14K_P]\) and \(\mu_k\ge k/2\), so
\begin{equation}
 B_k^{\sup}=\frac{\beta_k^2}{\mu_k^2}\ \le\ \frac{C_{\rm sp}^{\sup}}{k},
 \qquad
 C_{\rm sp}^{\sup}:=20\Bigl[\kappa C_\Phi^{1/2}+\kappa(4\nu C+3C_\Phi)^{1/2}+\frac14K_P\Bigr]^2 .
 \label{eq:c3v53-Csp}
\end{equation}
The text of Theorem~\ref{thm:c3v37-spectrum} and its proof have been
corrected in place accordingly, in the draft and in the file, before
this version was appended; the V52 digest recorded below is that of
the corrected file.
\end{proposition}

\begin{proof}
For \(k\ge1\): \((k/2)^{1/2}\le(5k)^{1/2}\), \(1\le(5k)^{1/2}\), and \(3+2k\le5k\); hence
\(\beta_k\le(5k)^{1/2}[\dots]\), \(\beta_k^2\le5k[\dots]^2\), and \(\mu_k^{-2}\le4/k^2\).
\end{proof}

\begin{remark}[The same step in the second cycle]
\label{rem:c3v53-legacy}
Theorem C2-V41-spectrum states \(\langle a_k\rangle\le C_{\rm sp}/k\) with
\(C_{\rm sp}=8[\kappa\langle\varphi\rangle^{1/2}+\kappa\langle d+3\varphi\rangle^{1/2}+K_P]^2\), by the same
reduction; the reduction is valid for \(k\ge3\) and fails for \(k=1,2\)
(\(3+2k=5,7\) against \(3k=3,6\)).  The \(O(1/k)\) statement is correct with
\(C_{\rm sp}\) replaced by \(20[\kappa\langle\varphi\rangle^{1/2}+\kappa\langle d+3\varphi\rangle^{1/2}+\frac14K_P]^2\),
which also covers \(k=1,2\).  The legacy file is not edited; the
in-cycle uses of C2-V41 (Firewall~\ref{fw:c3v21-mean-not-pointwise},
V26, V37) use only the \(O(1/k)\) form or the exact bracket
\eqref{eq:c3v26-Bk}, and are unaffected.
\end{remark}

% =============================================================================
\section{Re-derivations}
\label{sec:c3v53-rederivations}
% =============================================================================

\begin{theorem}[Audit of V36--V38]
\label{thm:c3v53-ledger}
The following were re-derived and hold as stated.
\begin{enumerate}[label=\textup{(\arabic*)},leftmargin=3.2em]
\item (V36, energy.)  \(\varphi'=2\langle V,Y\rangle_G=-4\mathcal L-\frac12j\); at \(\varphi'=0\),
      \(j=-8\mathcal L=-(d+2\varphi)\); the flux bound rearranged is
      \((1-\kappa)d+(2-4\kappa)\varphi\le K_P(d+3\varphi)^{1/2}\), the left side at least
      \(m_\kappa(d+3\varphi)\) since \(m_\kappa\le1-\kappa\) and \(3m_\kappa=2-4\kappa\); hence
      \(d+3\varphi\le K_P^2/m_\kappa^2\), \(\sup\varphi\le K_P^2/(3m_\kappa^2)\), and
      \(d\le K_P^2/m_\kappa^2-3c_\Phi\) at minimal energy.
\item (V36, Rayleigh.)  \(r'=-2\mathrm{Var}+\frac2\varphi\langle L_-V+rV,N\rangle_G\) (C2-V31-identity);
      \(\|L_-V+rV\|_G^2=\varphi\,\mathrm{Var}\); Cauchy--Schwarz gives \(\mathrm{Var}\le\|N\|_G^2/\varphi\)
      at \(r'=0\), then \(\kappa^2(\frac d2+8\nu C)/c_\Phi\le10\kappa^2\nu C/c_\Phi\).
\item (V37, spectrum.)  \(a_k'=-2\mu_ka_k-2n_k\), so \(n_k=-\mu_ka_k\) at an
      extremum, and \(\mu_ka_k\le a_k^{1/2}\beta_k\) with the class bounds;
      the \(1/k\) form as corrected (Proposition~\ref{prop:c3v53-correction}).
      The conditional floor and the incompatibility for \(k=0\) follow
      as in V26.
\item (V37, momentum.)  \(a_0'=-a_0-2n_0\), \(n_0=|G|\bar V\cdot\overline{N(V)}\); at an
      extremum \(\bar V\cdot\overline{N(V)}=-\frac12|\bar V|^2\); if the dominant level
      is \(0\) with \(\eta<\eta_{\rm c}\), V27 gives \(|\overline{N(V)}|\le\frac14|\bar V|\), a
      contradiction, so either the level is positive and
      \(\hat a_0\le\eta<\frac12\), or \(\mathrm{Var}\ge\min(\eta_{\rm c}/16,\frac1{32})\).
\item (V38.)  \(\frac12\mathfrak y'=\langle Y,L_-Y\rangle_G-\langle Y,N'(V)Y\rangle_G\) with
      \(\langle Y,L_-Y\rangle_G=-\nu\|\nabla Y\|_G^2-\frac12\|Y\|_G^2\); at an extremum the sign
      statement; the advective bounds \(C_\infty'\|Y\|^2\) and, by Hardy for
      \(Y\), \(\frac{C_\infty}{4\nu}\|Y\|((12\nu)^{1/2}\|Y\|+4\nu\|\nabla Y\|)\le\frac{\sqrt3}2\kappa\|Y\|^2+\frac\nu2\|\nabla Y\|^2+\frac{\kappa^2}2\|Y\|^2\);
      the pressure part \(\langle Y,\nabla P'\rangle_G\) by \((z\cdot Y)G=-2\nu Y\cdot\nabla G\).  The
      firewall's wording is superseded by V47 (recorded in V48).
\end{enumerate}
\end{theorem}

\begin{proof}
As described.
\end{proof}

\begin{theorem}[Audit of V41--V43]
\label{thm:c3v53-angle}
The following were re-derived and hold as stated.
\begin{enumerate}[label=\textup{(\arabic*)},leftmargin=3.2em]
\item (V41.)  \(|\det(\Omega,V,Y)|\le|\Omega||V||Y|\), \(\mathcal E\le\frac2\nu\mathcal L\) (closure),
      \(C_\infty(2\mathcal L/\nu)^{1/2}=\kappa(2\mathcal L)^{1/2}\).  H\"older (\(3,\frac32\)):
      \(\int|z|^2P_V^2G\le\|P_V\|_3^2(\int|z|^6G^3)^{1/3}\); \(\int|z|^6e^{-3|z|^2/(4\nu)}=4\pi\Gamma(\frac92)/(2a^{9/2})\)
      with \(a=3/(4\nu)\), \(\Gamma(\frac92)=\frac{105}{16}\sqrt\pi\), giving
      \(I_6=\frac{105\pi^{3/2}}8(\frac{4\nu}3)^{9/2}\).  The kinetic moment
      \(C_\infty\||z|V\|_G\le2\kappa\nu(d+3\varphi)^{1/2}\); dividing by \(2\nu\) and halving
      gives \(\frac\kappa2(d+3\varphi)^{1/2}\).  The ratio of the pressure term to
      \(K_P/4\) is \((105\pi^{3/2}/8)^{1/6}(4/3)^{3/4}/(2(4\pi/3)^{1/4})\approx2.045\cdot1.241/2.861\approx0.89\).
\item (V42.)  \(\ell=\Omega\times V\perp V,\Omega\); on the core, \(\|\Omega\|_G^2\le2\|\nabla\widetilde V\|_G^2\le2\varphi c_r\eta/\nu\)
      gives the Lamb bound \(\kappa(2\varphi c_r\eta)^{1/2}\|Y\|\); the kinetic Bernoulli
      part: constant term zero, cross term
      \(\frac1{2\nu}|\bar V|\||z|\widetilde V\|\|Y\|\le\varphi|G|^{-1/2}((3+4c_r)\eta/\nu)^{1/2}\|Y\|\),
      quadratic term with \(|\widetilde V|\le2C_\infty\): \(\frac1{2\nu}C_\infty\||z|\widetilde V\|\|Y\|=\kappa(\varphi\eta(3+4c_r))^{1/2}\|Y\|\)
      (as corrected before the V42 assembly).
\item (V43.)  \(\mathcal A+\mathcal B=-\|Y\|^2\) at extremal \(\mathcal L\) with \(|\mathcal A|\le\kappa(2\mathcal L)^{1/2}\|Y\|\);
      the tangential reduction: \(\langle\|Y\|^2\rangle=-\langle\mathcal A\rangle\le\kappa(2C_{\mathcal L})^{1/2}\langle\|Y\|^2\rangle^{1/2}\),
      \(2C_{\mathcal L}=\nu C+\frac12C_\Phi\).
\end{enumerate}
\end{theorem}

\begin{proof}
As described.
\end{proof}

\begin{proposition}[Citations of V36--V43 confirmed]
\label{prop:c3v53-citations}
Confirmed by inspection of the legacy file: C2-V17-mean-profile;
C2-V31-identity (the Rayleigh identity in both forms, with
\(\mathcal R=P^\perp L_-\widehat V\)); C2-V32-push (\(j=4\langle V,N\rangle_G=\frac2\nu\int\Pi(z\cdot V)G\));
C2-V36-hardy (\(X^2\le4\nu(d+3\varphi)\)); C2-V36-flux-bound; C2-V40-modewise
and C2-V40-low-modes; C2-V41-modes, C2-V41-bound (the bracket of
\eqref{eq:c3v37-Bsup} with \(\varphi,d\) in place of the class bounds),
C2-V41-spectrum (with the correction of Remark~\ref{rem:c3v53-legacy}).
\end{proposition}

\begin{proof}
Inspection at the labelled statements.
\end{proof}

% =============================================================================
\section{Integrated V53 conclusion and next acquisition order}
\label{sec:c3v53-conclusion}
% =============================================================================

\begin{theorem}[What V53 establishes]
\label{thm:c3v53-conclusion}
The correction (Proposition~\ref{prop:c3v53-correction}), the
re-derivations (Theorems~\ref{thm:c3v53-ledger}, \ref{thm:c3v53-angle})
and the confirmed citations (Proposition~\ref{prop:c3v53-citations}):
V36--V38 and V41--V43 stand, with the constant of the extremal
spectrum's \(1/k\) form corrected.  The full Navier--Stokes
stability/global-regularity theorem remains unproved.
\end{theorem}

\begin{proof}
The cited records.
\end{proof}

\begin{programme}[V54 acquisition order]
\label{prog:c3v54}
\begin{enumerate}[label=\textup{(AC\arabic*)},leftmargin=4.8em]
\item Audit of the correlation block (V44, V46--V48): re-derive the
      bilinear form and the dipole moment, the flux-derivative form,
      the far-pressure flatness and the energy--Hessian pairing with
      its cell decomposition, the multipole expansion, the pointwise
      constants of V47 (Riesz lemma, \(C_{P,0}\), \(C_P\), \(C_t\), \(C_Y'\), the
      decomposition of \(P'\)), \(\bar J\), \(C_{YP'}\) and the three bounds on
      the defect rate.
\item The list of every statement of V21--V48 whose wording is
      superseded (V38's firewall; V37's constant) with its current
      reading, and the list of in-place corrections of the cycle
      (V33, V37) with the versions recording them (V34, V53).
\item The next compulsory companion audit is V55.
\end{enumerate}
\end{programme}

\begin{assessment}[Position after V53]
\label{ass:c3v53-position}
The extremal ledger and the angle chapter pass the audit with one
constant corrected.  No full stability or global-regularity claim is
made.
\end{assessment}

% =============================================================================
\section*{Version V53 (third cycle) append-only record}
\addcontentsline{toc}{section}{Version V53 (third cycle) append-only record}
% =============================================================================

The complete V52 source of the third cycle, after the in-place
correction of Proposition~\ref{prop:c3v53-correction}, was retained
before this continuation.  Its SHA--256 digest was
\begin{center}\scriptsize\ttfamily
8d60fe591eccb90e2596b6fe2b5d623d0d7edf92dc50756768b55c62edbdcec9.
\end{center}
V53 audited V36--V38 and V41--V43 and corrected one constant.  No
full regularity claim is made.

% =============================================================================
\section{V54: internal audit, part 4: the correlation block, and the lists}
\label{sec:c3v54-entry}
% =============================================================================

Items (AC1) and (AC2) of Programme~\ref{prog:c3v54} are carried out.
The statements of V44 and V46--V48 were re-derived and their
constants recomputed.  One correction was found and applied in
place: the constant \(C_P\) of V47 was defined as a bound on each
component of the pressure gradient and then used as a bound on its
Euclidean norm; the definition now carries the factor \(\sqrt3\), and
every use of \(C_P\) (in \(C_t\), \(C_Y'\), \(\bar J\), \(C_{YP'}\)) is unchanged in form
with the corrected value.  The lists of superseded wordings and of
in-place corrections of the cycle are recorded.  No global
regularity claim is made.  The next compulsory companion audit is
V55.

% =============================================================================
\section{The correction}
\label{sec:c3v54-correction}
% =============================================================================

\begin{proposition}[Corrected definition of \(C_P\)]
\label{prop:c3v54-correction}
Lemma~\ref{lem:c3v47-riesz} bounds each \(\partial_lP_V=-\sum_{ij}R_iR_j\partial_l(V_iV_j)\)
by \(\frac23C_\infty C_\infty'+c_1(2C_\infty'^2+2C_\infty C_\infty'')+2c_2C_\infty C^{1/2}\); the Euclidean
norm \(|\nabla P_V|\) is at most \(\sqrt3\) times the largest component.
Lemma~\ref{lem:c3v47-decomposition}(i) now defines
\(C_P:=\sqrt3[\frac23C_\infty C_\infty'+c_1(2C_\infty'^2+2C_\infty C_\infty'')+2c_2C_\infty C^{1/2}]\), so that
\(|\nabla P_V|\le C_P\) as used in (ii), (iii), \eqref{eq:c3v47-J} and
\eqref{eq:c3v47-YP'}.  The correction was applied in the draft and in
the file before this version was appended; the V53 digest recorded
below is that of the corrected file.  \(C_P'\) was never given in
closed form and is understood as a bound on \(|\nabla^2P_V|\) (operator
norm on \(z\)); \(C_{P,0}\) bounds the scalar \(|P_V|\) and needs no factor.
\end{proposition}

\begin{proof}
\(|\nabla P_V|^2=\sum_l|\partial_lP_V|^2\le3\max_l|\partial_lP_V|^2\).
\end{proof}

% =============================================================================
\section{Re-derivations}
\label{sec:c3v54-rederivations}
% =============================================================================

\begin{theorem}[Audit of V44 and V46--V48]
\label{thm:c3v54-correlation}
The following were re-derived and hold as stated (with
Proposition~\ref{prop:c3v54-correction}).
\begin{enumerate}[label=\textup{(\arabic*)},leftmargin=3.2em]
\item (V44.)  \(\int P_Vf=\int q(\Gamma_N*f)\) by Fubini for \(f=(z\cdot Y)G\);
      \(q=\operatorname{tr}((\nabla V)^2)=|S|^2-|A|^2\), \(|A|^2=\frac12|\Omega|^2\); zero mass of \(f\);
      dipole moment \(|\int zf|\le(\int|z|^4G)^{1/2}\|Y\|_G\) with \(\int|z|^4G=60\nu^2|G|\)
      (\(4\pi\int s^6e^{-s^2/(4\nu)}\dd s=4\pi\cdot\frac{15\sqrt\pi}{16}(4\nu)^{7/2}=\frac{15}4(4\nu)^2|G|\)).
      \(\frac14j=\frac1{2\nu}\int\Pi(z\cdot V)G\) differentiated in \(\sigma\) gives the
      flux-derivative form; \(\partial_\sigma\Pi=P'+V\cdot Y\).
\item (V46.)  The far pressure's gradient on \(U\) is bounded through
      the \(|z-y|^{-4}\) kernel and the quiet energies, the mechanism of
      Lemma C2-V58-far; the two integrations by parts against a
      compactly supported smooth source and a Newtonian potential
      with decay are justified; the multipole expansion of
      \(\nabla^2\Gamma_N(z-y)\) about the cell centre gives the point-mass term
      \(\frac{m_\beta}{4\pi\rho^3}(3\hat e\otimes\hat e-I)\) (the Hessian of \(1/(4\pi|w|)\)) with
      remainder bounded by \(\sup|\nabla^3\Gamma_N|\) times the first moment;
      \(3(V\cdot\hat e)^2-|V|^2=|V|^2(3\cos^2\alpha-1)\); \(\operatorname{tr}\nabla^2\psi_\alpha=-f_\alpha\).
\item (V47.)  The rescaling \(V=\lambda u(x_0+\lambda z,T_*-\lambda^2)\), \(\lambda=e^{-\sigma/2}\), gives
      \(\partial_\sigma V=-\frac12V-\frac12z\cdot\nabla V+\lambda^3\partial_tu\) and
      \(\partial_\sigma P_V=-P_V-\frac12z\cdot\nabla P_V+\lambda^4\partial_tp\); \(\lambda^3\partial_tu=\nu\Delta V-(V\cdot\nabla)V-\nabla P_V\)
      is the rescaled equation; \(\lambda^4\partial_tp\) is the pressure of the
      rescaled \(\partial_t(u\otimes u)\).  The Riesz lemma's near part uses the
      zero spherical mean of the kernel, its far part H\"older.
      \(\bar J\): \(|P'|\le C_{P,0}+C_{P_t}+\frac12C_P|z|\), \(\int|z||V|G\le X|G|^{1/2}\),
      \(\int|z|^2|V|G\le X(6\nu|G|)^{1/2}\), \(X\le2\nu^{1/2}(d+3\varphi)^{1/2}\).  \(C_{YP'}\): the
      product of the two linear bounds integrated against \(G\) with
      \(\int|z|G=32\pi\nu^2\), \(\int|z|^2G=6\nu|G|\).  The boundary terms of the
      proposed bypass are \(O(1)\) and \(O(\rho)\) as stated.
\item (V48.)  The review of the four pairings and the three bounds
      on the defect rate are records; their content was checked
      against V27, V31, V33, V38, V41 and V47.
\end{enumerate}
\end{theorem}

\begin{proof}
As described.
\end{proof}

% =============================================================================
\section{The lists}
\label{sec:c3v54-lists}
% =============================================================================

\begin{proposition}[Superseded wordings and in-place corrections of the third cycle, V1--V54]
\label{prop:c3v54-lists}
\begin{enumerate}[label=\textup{(\roman*)},leftmargin=3.2em]
\item \emph{Superseded wordings (text unchanged, reading changed by a
      later version).}  Firewall~\ref{fw:c3v38-pressure} (``not
      explicitly bounded'') is read as Firewall~\ref{fw:c3v47-revised}
      (``bounded only at supremum level''), recorded in V47--V48.
      The direction statements of V19--V20 for V22 (the Hermite-mass
      computation) are superseded by V21; the direction of V19 for
      an explicit floor from concentration is closed by V22
      (N-C3-1); the plan of V29 for V32--V34 (quantitative NRS) is
      withdrawn in V31 (N-C3-3); the bypass proposed in V46 is
      withdrawn in V47.
\item \emph{In-place corrections (text changed, recorded in a later
      version).}  V33, Corollary~\ref{cor:c3v33-defect-window}: the
      hypothesis \(\kappa<\frac12\) removed (recorded in V34).  V37,
      Theorem~\ref{thm:c3v37-spectrum}: the constant of the \(1/k\) form
      corrected to \(20[\dots+\frac14K_P]^2\) (recorded in V53).  V47,
      Lemma~\ref{lem:c3v47-decomposition}(i): the factor \(\sqrt3\) in \(C_P\)
      (recorded in V54).  Corrections made before a version was first
      appended (V6, V13, V16, V24, V27, V32, V40, V42, V45, V46,
      V47) are part of the appended text and not listed.
\item \emph{Legacy statements found imprecise (not edited).}
      Theorem C2-V41-spectrum's constant \(C_{\rm sp}\), Remark~\ref{rem:c3v53-legacy}.
\end{enumerate}
\end{proposition}

\begin{proof}
The records of the cited versions.
\end{proof}

% =============================================================================
\section{Integrated V54 conclusion and next acquisition order}
\label{sec:c3v54-conclusion}
% =============================================================================

\begin{theorem}[What V54 establishes]
\label{thm:c3v54-conclusion}
The correction (Proposition~\ref{prop:c3v54-correction}), the
re-derivations (Theorem~\ref{thm:c3v54-correlation}) and the lists
(Proposition~\ref{prop:c3v54-lists}): V44 and V46--V48 stand, with
the definition of \(C_P\) corrected.  The full Navier--Stokes
stability/global-regularity theorem remains unproved.
\end{theorem}

\begin{proof}
The cited records.
\end{proof}

\begin{programme}[V55 acquisition order]
\label{prog:c3v55}
\begin{enumerate}[label=\textup{(AC\arabic*)},leftmargin=4.8em]
\item The compulsory companion audit: read the current SM and YM
      notes and record their digests and decision.
\item The audit's summary: the statement that V21--V48 stand with
      the three in-place corrections, and the position after
      fifty-five versions; delivery.
\end{enumerate}
\end{programme}

\begin{assessment}[Position after V54]
\label{ass:c3v54-position}
The correlation block passes the audit with one definition
corrected; the lists are recorded.  No full stability or
global-regularity claim is made.
\end{assessment}

% =============================================================================
\section*{Version V54 (third cycle) append-only record}
\addcontentsline{toc}{section}{Version V54 (third cycle) append-only record}
% =============================================================================

The complete V53 source of the third cycle, after the in-place
correction of Proposition~\ref{prop:c3v54-correction}, was retained
before this continuation.  Its SHA--256 digest was
\begin{center}\scriptsize\ttfamily
aca8cfeeaa4fd647cd0fc1603b6f1d1e817a2cb56e9bae31a190f449760d1182.
\end{center}
V54 audited V44 and V46--V48 and corrected one definition.  No full
regularity claim is made.

% =============================================================================
\section{V55: compulsory companion audit, the internal audit's summary, and the position after fifty-five versions}
\label{sec:c3v55-entry}
% =============================================================================

This is the fifty-fifth version of the third cycle.  The compulsory
review of the two companion programmes is performed first; the
internal audit V51--V54 is summarized; and the position after
fifty-five versions is recorded with the plan for V56--V60.  No
global regularity claim is made.  The next compulsory companion
audit is V60.

% =============================================================================
\section{Compulsory V55 review of the current companion notes}
\label{sec:c3v55-companion-audit}
% =============================================================================

The current companion files in \texttt{latex\_notes} were inspected
read-only.  Both programmes have moved since the V50 audit: the SM
programme from V34 to V36 of its seventh series, the YM programme
from V12 to V13 of its sixth series.  The frozen checkpoints are
\begin{align}
 &\texttt{Renewal\_SM\_Einstein\_Continuum\_Research\_Notes\_V36.tex},
 \notag\\[-2pt]
 &\qquad
 \texttt{ee42aa4eaa7a21b4d5860aaef80a25e2e257829e28557a0d7ad4d000241978f9},
 \label{eq:c3v55-SM-checkpoint}\\
 &\texttt{Renewal\_Yang\_Mills\_Mass\_Gap\_Research\_Notes\_V13.tex},
 \notag\\[-2pt]
 &\qquad
 \texttt{659aac50cf47c3df5aae9ca655669834797f059ffb0ee7ff1bbdd12adf9b43ba}.
 \label{eq:c3v55-YM-checkpoint}
\end{align}
(SM V36 has 14682 lines; YM V13 has 5846 lines.)

\emph{SM V35--V36.}  The SM programme performed its own companion
audit and a face-interface acquisition, then a finite
Birman--Schwinger reduction of a face: a reference-interface
hypothesis, an exact zero-mode equivalence, its relation to the
Evans determinant, a resolvent formula with a quantitative card, a
self-adjoint finite matrix when the defect has fixed rank, and a
finite acquisition target.

\emph{YM V13.}  The YM programme studied sparse physical grouping
graphs: a physical compatibility graph, a star-cover criterion, a
strictly local primitive row, and what remains outside its local
catalogue.

\begin{theorem}[V55 companion-audit decision]
\label{thm:c3v55-companion-decision}
Nothing is imported from SM V36 or YM V13.  Neither bears on the
third cycle's chapters or on the gates.
\end{theorem}

\begin{proof}
The SM results concern Birman--Schwinger reductions of lattice Dirac
interfaces; the YM results concern compatibility graphs in a cluster
expansion.  None is a Navier--Stokes statement, and none concerns
the ledger, the extremal method, the Lamb--Bernoulli form or the
correlation block.
\end{proof}

% =============================================================================
\section{Summary of the internal audit}
\label{sec:c3v55-summary}
% =============================================================================

\begin{theorem}[Result of the midpoint internal audit, V51--V54]
\label{thm:c3v55-audit}
The statements of V21--V48 that carry explicit constants or are used
by later versions were re-derived independently of their text and
their constants recomputed (Theorems~\ref{thm:c3v51-floor},
\ref{thm:c3v51-variance}, \ref{thm:c3v52-defect}, \ref{thm:c3v53-ledger},
\ref{thm:c3v53-angle}, \ref{thm:c3v54-correlation}), and every citation
of the second cycle was confirmed against the legacy file
(Propositions~\ref{prop:c3v51-citations}, \ref{prop:c3v52-citations},
\ref{prop:c3v53-citations}).  Three in-place corrections were made,
none affecting a conclusion: the hypothesis \(\kappa<\frac12\) removed from
Corollary~\ref{cor:c3v33-defect-window} (V34); the constant of the
\(1/k\) form in Theorem~\ref{thm:c3v37-spectrum} corrected (V53); the
factor \(\sqrt3\) in the definition of \(C_P\) in
Lemma~\ref{lem:c3v47-decomposition} (V54).  One legacy constant was
found imprecise and its correction recorded without editing the
frozen file (Remark~\ref{rem:c3v53-legacy}).  The lists of superseded
wordings and corrections are Proposition~\ref{prop:c3v54-lists}.
With these, V21--V48 stand.
\end{theorem}

\begin{proof}
The cited records.
\end{proof}

% =============================================================================
\section{Position after fifty-five versions, and the block V56--V60}
\label{sec:c3v55-position}
% =============================================================================

\begin{assessment}[Position after fifty-five versions of the third cycle]
\label{ass:c3v55-position-fifty-five}
The audited state of the cycle is that of V50
(Assessment~\ref{ass:c3v50-position-fifty}) with the three corrections
of Theorem~\ref{thm:c3v55-audit}.  The second half of the cycle begins
with the first candidate of V48, deferred at V49: the geometric gate
of V46 on the anchor-alone event.  On that event (V17) the
configuration is a single group of at most \(N_C\) cells at the centre
with a diffuse far field, and the gate's nonlocal part is the
energy--Hessian pairing over those cells with the multipole sign
structure of V46.  The plan: V56, the radial-defect masses \(m_\beta\) of
the anchor's cells and their sum (zero up to the far correction),
with the momentum identity of V18 read as a statement about the
first moment of the radial defect; V57, the pairing on the anchor
group in the two-cell and one-cell cases (a single cell has only the
self-term, whose isotropic part has the sign of minus the cell's
radial-defect mass), and whether the self-term's anisotropic part is
bounded by the cell's energy and enstrophy explicitly; V58, the
count of V17 and the localization of V16 applied to bound the
number of cross terms and their multipole constants, giving an
explicit bound on the pressure work on the anchor-alone event in
terms of the radial-defect masses, the cell energies and the
orientation angles; V59, consolidation and the statement of the
gate on the anchor-alone event as an inequality between explicit
geometric quantities of a bounded configuration, with the direction
for V61--V65; V60, the compulsory companion audit and delivery.  No
global regularity claim is made.
\end{assessment}

% =============================================================================
\section{Integrated V55 conclusion and next acquisition order}
\label{sec:c3v55-conclusion}
% =============================================================================

\begin{theorem}[What V55 establishes]
\label{thm:c3v55-conclusion}
Relative to the two legacy cycles and V1--V54: the companion-audit
decision (Theorem~\ref{thm:c3v55-companion-decision}), the audit's
summary (Theorem~\ref{thm:c3v55-audit}) and the position.  The full
Navier--Stokes stability/global-regularity theorem remains unproved.
\end{theorem}

\begin{proof}
The cited records.
\end{proof}

\begin{programme}[V56 acquisition order]
\label{prog:c3v56}
\begin{enumerate}[label=\textup{(AC\arabic*)},leftmargin=4.8em]
\item On the anchor-alone event \(\mathcal E_0\) (Theorem~\ref{thm:c3v17-cases}),
      define the radial-defect masses \(m_\beta=\int(z\cdot Y)G\phi_\beta\) of the cells
      \(\beta\) of the anchor's neighbourhood \(U_0\); prove
      \(|\sum_{\beta\in U_0}m_\beta|\le C_{13}\varepsilon_{\rm th}^{1/2}\|Y\|_G\) (the far cells' radial
      defect is small by the quiet energies and the Gaussian
      weight), and bound each \(|m_\beta|\) by \(\||z|\phi_\beta^{1/2}\|_G\|Y\|_{L^2(G;\beta)}\).
\item Read the momentum identity of V18 on \(\mathcal E_0\) in terms of the
      first moment of the radial defect: \(\int zf_{U_0}\) and
      \(\frac{\dd}{\dd\sigma}\bar V\); record what the two rows (momentum,
      radial-defect masses) say jointly about the dipole moment of
      \(\psi_{U_0}\).
\item The next compulsory companion audit is V60.
\end{enumerate}
\end{programme}

\begin{assessment}[Position after V55]
\label{ass:c3v55-position}
The companion audit and the internal audit are recorded; the
geometric gate on the anchor-alone event begins.  No full stability
or global-regularity claim is made.
\end{assessment}

% =============================================================================
\section*{Version V55 (third cycle) append-only record}
\addcontentsline{toc}{section}{Version V55 (third cycle) append-only record}
% =============================================================================

The complete V54 source of the third cycle was retained before this
continuation.  Its SHA--256 digest was
\begin{center}\scriptsize\ttfamily
d22de20205525c228099796768fa47eedbcc213fb71acbcdb0fa682e96aa5857.
\end{center}
V55 performed the compulsory companion audit, summarized the
internal audit and recorded the position.  No full regularity claim
is made.

% =============================================================================
\section{V56: the radial-defect masses of the anchor, and the momentum identity as a dipole statement}
\label{sec:c3v56-entry}
% =============================================================================

Items (AC1) and (AC2) of Programme~\ref{prog:c3v56} are carried out,
the first with a corrected estimate.  The total radial-defect mass
of the anchor's neighbourhood equals minus that of the far field,
which is bounded not by a power of the threshold, as the programme
anticipated, but by the Gaussian tail at the anchor radius, since the
defect on the far cells is not small in the class; the tail is
explicit and exponentially small in the anchor radius.  Hence the
anchor group carries outward and inward radial defect in nearly
equal amounts.  The first moment of the radial-defect density is
exactly twice the viscosity times the Gaussian volume times the
Gaussian mean of the defect, so the momentum identity of the second
cycle is the statement that the dipole moment of the radial defect
equals minus the viscosity times the Gaussian volume times the
momentum, minus twice that times the nonlinear mean; at the
configuration of extremal momentum the dipole moment is orthogonal
to the momentum, and on a nearly constant core it is nearly
anti-parallel to it.  No global regularity claim is made.  The next
compulsory companion audit is V60.

% =============================================================================
\section{Gaussian tails}
\label{sec:c3v56-tails}
% =============================================================================

\begin{lemma}[Explicit Gaussian tails]
\label{lem:c3v56-tails}
For \(R>0\), with \(T_p(R):=\int_{|z|>R}|z|^pG\,\dd z\),
\begin{align}
 T_2(R)&\le4\pi e^{-R^2/(4\nu)}\Bigl(2\nu R^3+12\nu^2R+\frac{24\nu^3}R\Bigr),
 \label{eq:c3v56-T2}\\
 T_4(R)&\le4\pi e^{-R^2/(4\nu)}\Bigl(2\nu R^5+20\nu^2R^3+120\nu^3R+\frac{240\nu^4}R\Bigr).
 \label{eq:c3v56-T4}
\end{align}
\end{lemma}

\begin{proof}
With \(a=1/(4\nu)\): \(\int_R^\infty s^4e^{-as^2}\dd s=\frac{R^3e^{-aR^2}}{2a}+\frac3{2a}\int_R^\infty s^2e^{-as^2}\dd s\)
and \(\int_R^\infty s^2e^{-as^2}\dd s=\frac{Re^{-aR^2}}{2a}+\frac1{2a}\int_R^\infty e^{-as^2}\dd s\le e^{-aR^2}\bigl(\frac R{2a}+\frac1{4a^2R}\bigr)\)
(integrations by parts and \(\int_R^\infty e^{-as^2}\le e^{-aR^2}/(2aR)\)); with
\(\frac1{2a}=2\nu\) this gives \(\int_R^\infty s^4e^{-as^2}\le e^{-aR^2}(2\nu R^3+12\nu^2R+24\nu^3/R)\)
and \eqref{eq:c3v56-T2} after the factor \(4\pi\).  Likewise
\(\int_R^\infty s^6e^{-as^2}=\frac{R^5e^{-aR^2}}{2a}+\frac5{2a}\int_R^\infty s^4e^{-as^2}\), giving
\eqref{eq:c3v56-T4}.
\end{proof}

% =============================================================================
\section{The radial-defect masses of the anchor}
\label{sec:c3v56-masses}
% =============================================================================

On the anchor-alone event \(\mathcal E_0\) (Theorem~\ref{thm:c3v17-cases}), with
\(U_0\) the anchor's neighbourhood (the cells within lattice distance
\(r_2+R_{c_\Phi'}\) of the origin), put
\begin{equation}
 f:=(z\cdot Y)G,\qquad m_\beta:=\int f\phi_\beta\ (\beta\in U_0),\qquad
 R_0:=r_2+R_{c_\Phi'}-1 ,
 \label{eq:c3v56-notation}
\end{equation}
so that the partition functions of the cells not in \(U_0\) are
supported in \(\{|z|\ge R_0\}\).

\begin{proposition}[Total radial-defect mass of the anchor]
\label{prop:c3v56-total}
\begin{equation}
 \boxed{\quad
 \sum_{\beta\in U_0}m_\beta=-\int_{\{|z|\ge R_0\}}(z\cdot Y)G\,\bigl(1-\textstyle\sum_{\beta\in U_0}\phi_\beta\bigr),
 \qquad
 \Bigl|\sum_{\beta\in U_0}m_\beta\Bigr|\ \le\ T_2(R_0)^{1/2}\,\|Y\|_G ,
 \quad}
 \label{eq:c3v56-total}
\end{equation}
and for each cell \(|m_\beta|\le\bigl(\int|z|^2G\phi_\beta\bigr)^{1/2}\|Y\|_{L^2(G;\beta)}\).
\end{proposition}

\begin{proof}
\(\int f=\int(z\cdot Y)G=0\) for divergence-free \(Y\) (Remark~\ref{rem:c3v31-source})
and \(\sum_{\text{all }\beta}\phi_\beta=1\); the far part by Cauchy--Schwarz with
\(|z|G^{1/2}\) and \(|Y|G^{1/2}\) on \(\{|z|\ge R_0\}\).
\end{proof}

\begin{remark}[The anchor carries both signs of radial defect]
\label{rem:c3v56-both-signs}
\(R_{c_\Phi'}=\lceil(4\nu\log(2S_G/c_\Phi'))^{1/2}\rceil\) (Proposition~\ref{prop:c3v16-anchor}),
so \(e^{-R_0^2/(4\nu)}\) is of the order of \(c_\Phi'/S_G\) up to the shift by
\(r_2-1\), and \(T_2(R_0)\) is explicit and small when the pinching
constant is small compared with the lattice sum.  The programme's
expectation of a bound by a power of \(\varepsilon_{\rm th}\) was wrong: the
far cells are quiet in energy and enstrophy, but the defect
\(Y=-\frac12V-\frac12z\cdot\nabla V+V_t\) on them is not small (its second
derivatives and pressure gradient are only class-bounded), and only
the Gaussian weight controls their radial defect.  What
\eqref{eq:c3v56-total} says is that the anchor's cells cannot all have
radial defect of one sign beyond the tail: outward and inward
radial-defect cells coexist in the anchor group, their masses
cancelling up to \(T_2(R_0)^{1/2}\|Y\|_G\).  In the energy--Hessian
pairing of V46 this means the cross terms carry masses of both
signs, and the geometric reading of Assessment~\ref{ass:c3v46-structure}
applies to both kinds of cell within one group.
\end{remark}

% =============================================================================
\section{The first moment, and the momentum identity}
\label{sec:c3v56-dipole}
% =============================================================================

\begin{theorem}[The dipole moment of the radial defect is the Gaussian mean of the defect]
\label{thm:c3v56-dipole}
For every profile of the family and every \(\sigma\),
\begin{equation}
 \boxed{\quad
 D_Y:=\int z\,(z\cdot Y)\,G\,\dd z\ =\ 2\nu|G|\,\bar Y,\qquad \bar Y=\frac1{|G|}\int YG=\frac{\dd}{\dd\sigma}\bar V=-\frac12\bar V-\overline{N(V)} ,
 \quad}
 \label{eq:c3v56-dipole}
\end{equation}
so the momentum identity (C2-V17, V18) reads
\(D_Y=-\nu|G|\bar V-2\nu|G|\overline{N(V)}\).  Moreover:
\begin{enumerate}[label=\textup{(\roman*)},leftmargin=3.2em]
\item at a maximizer or minimizer of the Gaussian momentum \(a_0\)
      (Theorem~\ref{thm:c3v37-momentum}), \(D_Y\cdot\bar V=0\);
\item on a nearly constant core with \(\eta\le\eta_{\rm c}\)
      (Corollary~\ref{cor:c3v27-decay}), \(|D_Y+\nu|G|\bar V|\le\frac\nu4|G|^{1/2}\varphi^{1/2}\);
\item on \(\mathcal E_0\), the localized dipole moment
      \(D_{U_0}:=\int z\,f\sum_{\beta\in U_0}\phi_\beta\) satisfies \(|D_Y-D_{U_0}|\le T_4(R_0)^{1/2}\|Y\|_G\).
\end{enumerate}
\end{theorem}

\begin{proof}
\(z_jG=-2\nu\partial_jG\) gives \(\int z_iz_jY_jG=-2\nu\int z_iY\cdot\nabla G=2\nu\int\operatorname{div}(z_iY)G=2\nu\int Y_iG\)
by \(\operatorname{div}Y=0\); and \(\int YG=\frac{\dd}{\dd\sigma}\int VG\).  (i) \(\bar V\cdot\overline{N(V)}=-\frac12|\bar V|^2\)
there, so \(\bar V\cdot D_Y=-\nu|G||\bar V|^2+\nu|G||\bar V|^2=0\).  (ii)
\(|\overline{N(V)}|\le\frac18(\varphi/|G|)^{1/2}\) there.  (iii) Cauchy--Schwarz on
\(\{|z|\ge R_0\}\) with \(|z|^2G^{1/2}\) and \(|Y|G^{1/2}\).
\end{proof}

\begin{assessment}[Item (AC2): the two rows jointly]
\label{ass:c3v56-joint}
The potential \(\psi_{U_0}\) of the anchor's radial-defect density has
mass at most \(T_2(R_0)^{1/2}\|Y\|_G\) and dipole moment
\(-\nu|G|\bar V-2\nu|G|\overline{N(V)}\) up to \(T_4(R_0)^{1/2}\|Y\|_G\).  Its far field
is therefore, to leading order, the dipole field of the vector
\(-\nu|G|(\bar V+2\overline{N(V)})\): on a nearly constant core, the dipole field
of \(-\nu|G|\bar V\), opposite to the momentum; at extremal momentum, a
dipole orthogonal to the momentum.  Inside the configuration the
multipole expansion does not apply and the pairing of V46 is the
finite sum over cells.  The joint reading is that the radial defect
of the anchor group is a zero-mass, momentum-sized dipole
distribution over at most \(N_C(2r_2+1)^3\) cells, and the gate's
nonlocal part is the pairing of the group's velocity tensor with the
Hessian of that distribution's potential.  No global regularity
claim is made.
\end{assessment}

% =============================================================================
\section{Integrated V56 conclusion and next acquisition order}
\label{sec:c3v56-conclusion}
% =============================================================================

\begin{theorem}[What V56 proves]
\label{thm:c3v56-conclusion}
Relative to the two legacy cycles and V1--V55: the explicit tails
(Lemma~\ref{lem:c3v56-tails}), the total radial-defect mass of the
anchor (Proposition~\ref{prop:c3v56-total}) and the dipole identity
with its consequences (Theorem~\ref{thm:c3v56-dipole}).  The full
Navier--Stokes stability/global-regularity theorem remains unproved.
\end{theorem}

\begin{proof}
The cited results.
\end{proof}

\begin{programme}[V57 acquisition order]
\label{prog:c3v57}
\begin{enumerate}[label=\textup{(AC\arabic*)},leftmargin=4.8em]
\item The one-cell pairing: for a configuration whose radial defect
      is supported in one cell \(\alpha\) (up to the far correction), the
      pressure work is the self-term \(\frac1{2\nu}\int\chi_{U'}(V\otimes V):\nabla^2\psi_\alpha\);
      split \(\nabla^2\psi_\alpha\) into its trace part \(-\frac13f_\alpha I\) and its
      traceless part \(H_\alpha\); bound the traceless part's contribution by
      \(\frac1{2\nu}\|V\otimes V\|_{L^2(U')}\|H_\alpha\|_{L^2}\) with \(\|H_\alpha\|_{L^2}\le\|f_\alpha\|_{L^2}\)
      (the traceless part of the Hessian of a Newtonian potential is
      a Riesz transform of the density), and the trace part exactly;
      state the resulting two-sided bound on the one-cell pressure
      work in terms of \(\int\phi_\alpha|V|^2f_\alpha\), \(\|f_\alpha\|_{L^2}\) and the energy
      of \(U'\).
\item The two-cell pairing: cells \(\alpha,\beta\) at lattice distance
      \(\ge2\), masses \(m_\alpha,m_\beta\) with \(m_\alpha+m_\beta\) small
      (Proposition~\ref{prop:c3v56-total}): the cross terms by
      Proposition~\ref{prop:c3v46-multipole}, and the sign structure
      when \(m_\alpha\approx-m_\beta\) (a radial-defect dipole across the group).
\item The next compulsory companion audit is V60.
\end{enumerate}
\end{programme}

\begin{assessment}[Position after V56]
\label{ass:c3v56-position}
The anchor's radial defect is a zero-mass dipole distribution whose
dipole moment is the momentum identity.  No full stability or
global-regularity claim is made.
\end{assessment}

% =============================================================================
\section*{Version V56 (third cycle) append-only record}
\addcontentsline{toc}{section}{Version V56 (third cycle) append-only record}
% =============================================================================

The complete V55 source of the third cycle was retained before this
continuation.  Its SHA--256 digest was
\begin{center}\scriptsize\ttfamily
1267d7e9650965680c0c91b827629239731ab7f8716d83c7abcaf4e03161db6a.
\end{center}
V56 proved the radial-defect mass and dipole statements.  No full
regularity claim is made.

% =============================================================================
\section{V57: the one-cell and two-cell pairings}
\label{sec:c3v57-entry}
% =============================================================================

Items (AC1) and (AC2) of Programme~\ref{prog:c3v57} are carried out.
(AC1): when the radial defect is carried by one cell, the pressure
work is the self-term, whose trace part is exactly minus one sixth
of the viscosity's inverse times the cell's radial defect weighted by
the squared velocity, and whose traceless part is bounded by the
\(L^2\) norm of the cell's radial-defect density times the \(L^2\) norm of
the traceless velocity tensor on the halo, with no singular-integral
constant since the Riesz transforms are isometric on \(L^2\).  (AC2):
when the radial defect is carried by two cells with opposite masses,
the cross terms are, to leading multipole order, the difference of
the axial anisotropies of the velocity tensor on the two cells,
times the mass and the inverse cube of the distance: the cross
pressure work is negative when the flow is more azimuthal about the
outward-defect cell than about the inward-defect cell.  No global
regularity claim is made.  The next compulsory companion audit is
V60.

% =============================================================================
\section{The Hessian of a cell potential}
\label{sec:c3v57-hessian}
% =============================================================================

For a bounded compactly supported \(g\) let \(\psi=\Gamma_N*g\).  In Fourier
variables \(\widehat{\partial_i\partial_j\psi}=-(\xi_i\xi_j/|\xi|^2)\hat g\), so
\(\nabla^2\psi=-\mathsf R g\) with \(\mathsf R:=(R_iR_j)_{ij}\) the matrix of second Riesz
transforms, \(\operatorname{tr}\nabla^2\psi=-g\), and
\begin{equation}
 \nabla^2\psi=-\tfrac13g\,I+H,\qquad H:=\nabla^2\psi+\tfrac13g\,I\ \text{ traceless},
 \qquad \|H\|_{L^2}\le\|\mathsf Rg\|_{L^2}=\|g\|_{L^2}
 \label{eq:c3v57-H}
\end{equation}
(Frobenius norms: \(\sum_{ij}\|R_iR_jg\|_2^2=\int\sum_{ij}\xi_i^2\xi_j^2|\xi|^{-4}|\hat g|^2=\|g\|_2^2\), and
removing the trace part does not increase the Frobenius norm).

% =============================================================================
\section{The one-cell pairing}
\label{sec:c3v57-one-cell}
% =============================================================================

\begin{theorem}[One cell]
\label{thm:c3v57-one-cell}
Let the radial-defect density on the configuration be
\(f_\alpha=(z\cdot Y)G\phi_\alpha\) for one cell \(\alpha\) (the other cells of \(U\)
contributing nothing), \(\psi_\alpha=\Gamma_N*f_\alpha\), and \(U''\) the support of
\(\chi_{U'}\).  Then the pressure work of Theorem~\ref{thm:c3v46-pairing}
is
\begin{equation}
 \boxed{\quad
 \frac1{2\nu}\int\chi_{U'}(V\otimes V):\nabla^2\psi_\alpha
 =-\frac1{6\nu}\int\chi_{U'}|V|^2f_\alpha+\frac1{2\nu}\int\chi_{U'}(V\otimes V)^\circ:H_\alpha ,
 \qquad
 \Bigl|\frac1{2\nu}\int\chi_{U'}(V\otimes V)^\circ:H_\alpha\Bigr|\le\frac1{2\nu}\Bigl(\frac23\Bigr)^{1/2}\|V\|_{L^4(U'')}^2\,\|f_\alpha\|_{L^2} ,
 \quad}
 \label{eq:c3v57-one-cell}
\end{equation}
with \((V\otimes V)^\circ=V\otimes V-\frac13|V|^2I\) and
\(\|f_\alpha\|_{L^2}\le\sup_{Q_\alpha}(|z|G^{1/2})\,\|Y\|_{L^2(G;\alpha)}\).  In particular, if
\(z\cdot Y\ge0\) on the cell (outward radial defect), the trace part is
\(\le0\), and the one-cell pressure work is negative whenever
\(\frac13\int\chi_{U'}|V|^2f_\alpha>(\frac23)^{1/2}\|V\|_{L^4(U'')}^2\|f_\alpha\|_{L^2}\); if \(z\cdot Y\le0\)
on the cell, the trace part is \(\ge0\) and the reverse holds.
\end{theorem}

\begin{proof}
\eqref{eq:c3v57-H} with \(g=f_\alpha\); the trace part of \(V\otimes V\) pairs only
with the trace part of the Hessian and the traceless part only with
\(H_\alpha\) (both decompositions are orthogonal in the Frobenius inner
product); \(|(V\otimes V)^\circ|_F^2=|V|^4-\frac13|V|^4=\frac23|V|^4\), so
\(\|\chi_{U'}(V\otimes V)^\circ\|_{L^2}\le(\frac23)^{1/2}\|V\|_{L^4(U'')}^2\); Cauchy--Schwarz and
\eqref{eq:c3v57-H}.  The bound on \(\|f_\alpha\|_2\): \(f_\alpha^2=(z\cdot Y)^2G^2\phi_\alpha^2\le|z|^2G\cdot|Y|^2G\phi_\alpha\)
on the cell.
\end{proof}

\begin{remark}[Reading of the one-cell case]
\label{rem:c3v57-one-cell}
A single cell of outward radial defect does negative pressure work
through its trace part, of size one sixth of the inverse viscosity
times the squared velocity weighted by the radial defect, and the
traceless part can overturn it only if the velocity tensor on the
halo is strongly anisotropic relative to the Hessian of the cell's
potential.  For a single cell of inward radial defect the trace part
is positive.  Since the anchor's total mass is nearly zero
(Proposition~\ref{prop:c3v56-total}), a single cell can carry the
anchor's radial defect only if that mass is nearly zero, i.e.\ only
if \(z\cdot Y\) changes sign within the cell; the one-cell case with a
signed radial defect is therefore not the anchor-alone case but a
model of one cell within the group, and the trace parts of the
group's cells sum to \(-\frac1{6\nu}\sum_\alpha\int\chi_{U'}|V|^2f_\alpha=-\frac1{6\nu}\int\chi_{U'}|V|^2f_{U_0}\),
the radial defect weighted by the squared velocity, with no sign
from the masses alone.
\end{remark}

% =============================================================================
\section{The two-cell pairing}
\label{sec:c3v57-two-cell}
% =============================================================================

\begin{theorem}[Two cells of opposite mass]
\label{thm:c3v57-two-cell}
Let the radial defect on the configuration be carried by two cells
\(\alpha,\beta\) with centres at distance \(\rho\ge2\), unit vector
\(\hat e=(z_\beta-z_\alpha)/\rho\) along the axis, masses \(m_\alpha,m_\beta\), and absolute
first moments \(\mu_\alpha,\mu_\beta\) about their centres.  Let
\(Q_\gamma:=\int\chi_{U'}\phi_\gamma\bigl(3(V\cdot\hat e)^2-|V|^2\bigr)\) be the axial anisotropy of the
velocity tensor on cell \(\gamma\).  Then the cross part of the pressure
work,
\(\frac1{2\nu}\bigl[\int\chi_{U'}\phi_\alpha(V\otimes V):\nabla^2\psi_\beta+\int\chi_{U'}\phi_\beta(V\otimes V):\nabla^2\psi_\alpha\bigr]\),
equals
\begin{equation}
 \boxed{\quad
 \frac{1}{8\pi\nu\rho^3}\bigl(m_\beta Q_\alpha+m_\alpha Q_\beta\bigr)+E_{\rm mult},
 \qquad
 |E_{\rm mult}|\le\frac{C_{12}}{2\nu\rho^4}\Bigl(\mu_\beta\int\phi_\alpha|V|^2+\mu_\alpha\int\phi_\beta|V|^2\Bigr)+\frac{C_{14}}{2\nu\rho^4}\bigl(|m_\alpha|+|m_\beta|\bigr)\|V\|_{L^\infty}^2 ,
 \quad}
 \label{eq:c3v57-two-cell}
\end{equation}
the second error accounting for the variation of \(\rho^{-3}\) and \(\hat e\)
across a unit cell.  If \(m_\alpha=-m_\beta=:-m\) with \(m>0\) (outward radial
defect at \(\beta\), inward at \(\alpha\)),
\begin{equation}
 \frac{m}{8\pi\nu\rho^3}\bigl(Q_\alpha-Q_\beta\bigr)+E_{\rm mult}:
 \label{eq:c3v57-dipole}
\end{equation}
the cross work is negative, to leading order, exactly when the flow
on the inward cell is more azimuthal about the axis than the flow on
the outward cell (\(Q_\alpha<Q_\beta\)).
\end{theorem}

\begin{proof}
Proposition~\ref{prop:c3v46-multipole} applied to \(\nabla^2\psi_\beta\) on the cell
\(\alpha\) and to \(\nabla^2\psi_\alpha\) on the cell \(\beta\), with \(\hat e\) replaced by
\(-\hat e\) for the second (which leaves \((V\cdot\hat e)^2\) unchanged); the
point-mass term is evaluated at the cell centres, and its variation
across a cell (\(|\nabla_z(\rho^{-3}(3\hat e\otimes\hat e-I))|\le C_{14}'\rho^{-4}\) for \(\rho\ge2\))
is absorbed into \(E_{\rm mult}\) with the cell energies bounded by
\(\|V\|_\infty^2\) times the cell volume.  With \(m_\alpha=-m_\beta\),
\(m_\beta Q_\alpha+m_\alpha Q_\beta=m(Q_\alpha-Q_\beta)\).
\end{proof}

\begin{assessment}[Item (AC2): the radial-defect dipole across a group]
\label{ass:c3v57-dipole}
Proposition~\ref{prop:c3v56-total} makes the two-cell case with
\(m_\alpha\approx-m_\beta\) the generic two-cell anchor: a radial-defect dipole
across the group.  Its cross pressure work has the sign of
\(Q_\alpha-Q_\beta\), a comparison of the axial anisotropy of the velocity
tensor on the two cells, weighted by \(m/(8\pi\nu\rho^3)\); its self-terms
are the two one-cell terms of Theorem~\ref{thm:c3v57-one-cell}, whose
trace parts have opposite signs and whose sum is
\(-\frac1{6\nu}\int\chi_{U'}|V|^2(f_\alpha+f_\beta)\).  The gate on this configuration
asks that the sum of these four pieces be at most \(-\|Y\|_G^2\) at
extremal Dirichlet form, or \(-\underline\delta\) on average, up to the Lamb
pairing and the far corrections.  Every piece is an explicit
functional of the velocity and the radial defect on at most two
cells and their halo; none has a sign forced by the class.  The
class's constraints on the pieces are: the masses are bounded by
the cell defects, the anisotropies by the cell energies, and the
distance \(\rho\) by the group's diameter \(2r_2+4\); so the cross work is
at most \(\frac{|m|}{8\pi\nu\rho^3}(|Q_\alpha|+|Q_\beta|)\le\frac{|m|}{4\pi\nu\rho^3}(\int\phi_\alpha|V|^2+\int\phi_\beta|V|^2)\)
in absolute value, an explicit bound in the cell energies and the
radial-defect masses, which V58 assembles for the whole group.
\end{assessment}

% =============================================================================
\section{Integrated V57 conclusion and next acquisition order}
\label{sec:c3v57-conclusion}
% =============================================================================

\begin{theorem}[What V57 proves]
\label{thm:c3v57-conclusion}
Relative to the two legacy cycles and V1--V56: the Hessian
decomposition \eqref{eq:c3v57-H}, the one-cell pairing
(Theorem~\ref{thm:c3v57-one-cell}) and the two-cell pairing
(Theorem~\ref{thm:c3v57-two-cell}).  The full Navier--Stokes
stability/global-regularity theorem remains unproved.
\end{theorem}

\begin{proof}
The cited results.
\end{proof}

\begin{programme}[V58 acquisition order]
\label{prog:c3v58}
\begin{enumerate}[label=\textup{(AC\arabic*)},leftmargin=4.8em]
\item The whole group: with the count of V17 (at most
      \(N_C(2r_2+1)^3\) cells in \(U_0\)) and the diameter bound, assemble
      the pressure work on \(\mathcal E_0\) as the sum of the trace parts
      (\(-\frac1{6\nu}\int\chi_{U'}|V|^2f_{U_0}\)), the traceless self-terms (bounded
      by \(\frac1{2\nu}(\frac23)^{1/2}\|V\|_{L^4(U'')}^2\sum_\alpha\|f_\alpha\|_2\)), the cross terms
      (bounded by the multipole form with distances \(\ge2\), and by
      the direct bound \(\frac1{2\nu}\|V\|_{L^4(U'')}^2\|f_\beta\|_2\) for neighbouring
      cells), and the far correction \(E_{\rm far}\); state the explicit
      bound \(|\mathcal B_{U_0}|\le\mathcal M_{U_0}\) in terms of \(\|V\|_{L^4(U'')}\), the cell
      radial-defect norms and the counts.
\item The gate on \(\mathcal E_0\) in geometric form: combine with the Lamb
      bound on \(U_0\) (Proposition~\ref{prop:c3v41-lamb} localized) and
      Theorem~\ref{thm:c3v43-extremal}: at an extremal-Dirichlet
      configuration on \(\mathcal E_0\), \(\|Y\|_G^2\le\kappa(2\mathcal L)^{1/2}\|Y\|_G+\mathcal M_{U_0}+|E_{\rm far}|\);
      record what it bounds.
\item The next compulsory companion audit is V60.
\end{enumerate}
\end{programme}

\begin{assessment}[Position after V57]
\label{ass:c3v57-position}
The one-cell and two-cell pairings are explicit, with the trace
part signed by the radial defect and the cross terms signed by the
axial anisotropy difference.  No full stability or global-regularity
claim is made.
\end{assessment}

% =============================================================================
\section*{Version V57 (third cycle) append-only record}
\addcontentsline{toc}{section}{Version V57 (third cycle) append-only record}
% =============================================================================

The complete V56 source of the third cycle was retained before this
continuation.  Its SHA--256 digest was
\begin{center}\scriptsize\ttfamily
7fbf7c6c6474e18f25c8fb46172602d492be85e6855bd70a190e920762cda558.
\end{center}
V57 proved the one-cell and two-cell pairings.  No full regularity
claim is made.

% =============================================================================
\section{V58: the pressure work on the anchor-alone event, and the gate in geometric form}
\label{sec:c3v58-entry}
% =============================================================================

Items (AC1) and (AC2) of Programme~\ref{prog:c3v58} are carried out.
(AC1): the whole-group pressure work on the anchor-alone event is
decomposed, without summing over cells, into its trace part, the
radial defect of the neighbourhood weighted by the squared velocity
over six times the viscosity, and its traceless part, bounded by the
\(L^2\) norm of the neighbourhood's radial-defect density times the
\(L^4\) energy of the halo, with the far correction of V46; the
resulting bound is local: it involves the halo's energy, the
neighbourhood's radial defect and the anchor radius, and no global
constant.  (AC2): on the anchor-alone event the gate becomes an
inequality between the defect, the velocity constant times the
Dirichlet form, and the halo's \(L^4\) energy times the anchor radius
over the viscosity, at an extremal-Dirichlet configuration lying in
the event, and a bound on the defect rate by the same local
quantities in the mean over the event.  No global regularity claim
is made.  The next compulsory companion audit is V60.

% =============================================================================
\section{The whole-group bound}
\label{sec:c3v58-group}
% =============================================================================

On \(\mathcal E_0\), with \(U_0\) the anchor's neighbourhood, \(U'\) its halo and
\(U''\) the support of \(\chi_{U'}\) (cells within lattice distance \(3\) of
\(U_0\)), \(f_{U_0}=(z\cdot Y)G\mathbf 1_{U_0}\), \(\psi_{U_0}=\Gamma_N*f_{U_0}\), and
\(\rho_0:=r_2+R_{c_\Phi'}+1\) (so that \(|z|\le\rho_0\) on \(U_0\)).

\begin{theorem}[Pressure work on the anchor-alone event]
\label{thm:c3v58-group}
On the trapped sub-windows of \(\mathcal E_0\),
\begin{equation}
 \frac1{2\nu}\int_{U_0}(P_V-c)(z\cdot Y)G
 =-\frac1{6\nu}\int\chi_{U'}|V|^2f_{U_0}+\frac1{2\nu}\int\chi_{U'}(V\otimes V)^\circ:H_{U_0}+E_{\rm far},
 \label{eq:c3v58-decomposition}
\end{equation}
with \(H_{U_0}=\nabla^2\psi_{U_0}+\frac13f_{U_0}I\) traceless, \(|E_{\rm far}|\le C_{11}\varepsilon_{\rm th}\|Y\|_G\), and
\begin{equation}
 \boxed{\quad
 \Bigl|\frac1{2\nu}\int_{U_0}(P_V-c)(z\cdot Y)G\Bigr|\ \le\ \mathcal M_{U_0}+C_{11}\varepsilon_{\rm th}\|Y\|_G,
 \qquad
 \mathcal M_{U_0}:=\frac1{2\nu}\Bigl(\frac13+\Bigl(\frac23\Bigr)^{1/2}\Bigr)\|V\|_{L^4(U'')}^2\,\rho_0\,\|Y\|_{L^2(G;U_0)} .
 \quad}
 \label{eq:c3v58-bound}
\end{equation}
Moreover \(\|V\|_{L^4(U'')}^2\le C_\infty\,e_{U''}^{1/2}\) with \(e_{U''}=\int_{U''}|V|^2\) the
unweighted energy of \(U''\), at most \(K_6^2|U''|^{2/3}\), and
\(|U''|\le(2(r_2+R_{c_\Phi'}+3)+1)^3\).
\end{theorem}

\begin{proof}
Theorem~\ref{thm:c3v46-pairing} with \(U=U_0\), and the decomposition
\eqref{eq:c3v57-H} with \(g=f_{U_0}\).  The trace part is bounded by
\(\frac1{6\nu}\|\chi_{U'}|V|^2\|_{L^2}\|f_{U_0}\|_{L^2}\le\frac1{6\nu}\|V\|_{L^4(U'')}^2\|f_{U_0}\|_2\); the
traceless part by \(\frac1{2\nu}(\frac23)^{1/2}\|V\|_{L^4(U'')}^2\|H_{U_0}\|_2\) with
\(\|H_{U_0}\|_2\le\|f_{U_0}\|_2\); and \(\|f_{U_0}\|_2^2=\int_{U_0}(z\cdot Y)^2G^2\le\rho_0^2\int_{U_0}|Y|^2G\)
since \(G\le1\).  \(\|V\|_{L^4}^4\le\|V\|_\infty^2\|V\|_{L^2}^2\); \(\|V\|_{L^2(U'')}^2\le\|V\|_{L^6}^2|U''|^{2/3}\).
\end{proof}

\begin{remark}[Local against global]
\label{rem:c3v58-local}
Theorem~\ref{thm:c3v41-MNY} bounded the Bernoulli work on the whole
space by \(M_{NY}\|Y\|_G\) with the global pressure constant
\(C_{\rm CZ}K_6^2I_6^{1/6}/(2\nu)\); on the anchor-alone event the pressure
work on the neighbourhood is bounded by \eqref{eq:c3v58-bound}, in
which the pressure has disappeared: the energy form has replaced
the pressure by the velocity tensor on the halo, and the only
constants are the anchor radius and the halo's energy.  The
traceless part of the Hessian of the neighbourhood's potential is
controlled with no singular-integral constant because the Riesz
transforms are isometric on \(L^2\); the price is the \(L^4\) energy of
the halo in place of the Gaussian quantities.
\end{remark}

% =============================================================================
\section{The gate on the anchor-alone event in geometric form}
\label{sec:c3v58-gate}
% =============================================================================

\begin{theorem}[The gate on \(\mathcal E_0\)]
\label{thm:c3v58-gate}
On the trapped sub-windows of \(\mathcal E_0\):
\begin{enumerate}[label=\textup{(\roman*)},leftmargin=3.2em]
\item (Pointwise.)  If a maximizer or minimizer \(W^*\) of \(\mathcal L\) over
      the family lies in \(\mathcal E_0\), then at \(W^*\)
      \begin{equation}
       \|Y\|_G^2\ \le\ \kappa(2\mathcal L)^{1/2}\|Y\|_G+\mathcal M_{U_0}+C_{11}\varepsilon_{\rm th}\|Y\|_G+C_7\varepsilon_{\rm th}^{1/4}C_Y+C_{\rm class}(1+R_G^4)e^{-R_G^2/(8\nu)},
       \label{eq:c3v58-pointwise}
      \end{equation}
      hence
      \(\|Y\|_G\le\kappa(2\mathcal L)^{1/2}+\frac{\rho_0}{2\nu}(\frac13+(\frac23)^{1/2})\|V\|_{L^4(U'')}^2+C_{11}\varepsilon_{\rm th}+o(1)\)
      with \(o(1)\) the last two terms divided by \(\|Y\|_G\).
\item (Mean.)  In case (a) of Theorem~\ref{thm:c3v17-cases},
      \begin{equation}
       \frac12\underline\delta\,\mathfrak m(\mathcal E_0)\ \le\ \int_{\mathcal E_0}\Bigl[\kappa(2\mathcal L)^{1/2}\|Y\|_G+\mathcal M_{U_0}\Bigr]\dd\mathfrak m+O(\varepsilon_{\rm th}^{1/4})+O\bigl((1+R_G^4)e^{-R_G^2/(8\nu)}\bigr).
       \label{eq:c3v58-mean}
      \end{equation}
\end{enumerate}
\end{theorem}

\begin{proof}
(i) At extremal \(\mathcal L\), \(\langle N,Y\rangle_G=-\|Y\|_G^2\); on \(\mathcal E_0\) the cross term
is \(X_{U_0}\) up to the far cells' contribution \(O(\varepsilon_{\rm th}^{1/4}C_Y)\)
(Theorem~\ref{thm:c3v11-far}) and the tail beyond \(R_G\)
(Theorem~\ref{thm:c3v16-localized}); \(X_{U_0}=\mathcal A_{U_0}+\mathcal B_{U_0}\) with
\(|\mathcal A_{U_0}|\le\kappa(2\mathcal L)^{1/2}\|Y\|_G\) (Proposition~\ref{prop:c3v41-lamb}
restricted to \(U_0\)) and \(\mathcal B_{U_0}\) bounded by \eqref{eq:c3v58-bound}.
(ii) Theorem~\ref{thm:c3v17-cases}(a) with the same bounds integrated
against \(\mathfrak m\) on \(\mathcal E_0\).
\end{proof}

\begin{assessment}[Item (AC2): what the geometric form bounds]
\label{ass:c3v58-reading}
On the anchor-alone event the gate's demand, anti-alignment at the
rate of the defect, is met at most by two local sources: the Lamb
pairing, at most the velocity constant times the square root of
twice the Dirichlet form per unit defect, and the pressure work, at
most the halo's \(L^4\) energy times the anchor radius over the
viscosity per unit defect on the neighbourhood.  Both are bounded
by the configuration's own energy; neither has a sign.  The
inequality \eqref{eq:c3v58-mean} bounds the defect rate, in the mean
over the anchor-alone event, by these local quantities: it is a
fourth bound on \(\underline\delta\), local in the anchor.  It also says
where an exclusion would have to come from on this event: a proof
that the halo's velocity tensor cannot be oriented, relative to the
Hessian of the neighbourhood's radial-defect potential, so as to
make \eqref{eq:c3v58-decomposition} as negative as \(-\|Y\|_G^2\) plus the
Lamb pairing.  By V56 the potential is that of a zero-mass dipole
distribution of moment \(-\nu|G|(\bar V+2\overline{N(V)})\); by V57 the
orientation rule is explicit for one and two cells.  The class does
not fix the orientation, and no exclusion follows.  No global
regularity claim is made.
\end{assessment}

% =============================================================================
\section{Integrated V58 conclusion and next acquisition order}
\label{sec:c3v58-conclusion}
% =============================================================================

\begin{theorem}[What V58 proves]
\label{thm:c3v58-conclusion}
Relative to the two legacy cycles and V1--V57: the whole-group bound
(Theorem~\ref{thm:c3v58-group}) and the gate on the anchor-alone event
in geometric form (Theorem~\ref{thm:c3v58-gate}).  The full
Navier--Stokes stability/global-regularity theorem remains unproved.
\end{theorem}

\begin{proof}
The cited results.
\end{proof}

\begin{programme}[V59 acquisition order]
\label{prog:c3v59}
\begin{enumerate}[label=\textup{(AC\arabic*)},leftmargin=4.8em]
\item Consolidate V56--V58 as the anchor-gate block: the radial-defect
      masses and dipole, the one- and two-cell pairings, the
      whole-group bound and the gate in geometric form; record the
      fourth bound on the defect rate beside the three of V48.
\item Fix the direction for V61--V65.  With the gate on \(\mathcal E_0\) reduced
      to the orientation of the halo's velocity tensor relative to a
      dipole distribution's potential, the candidates are: (a) the
      orientation under the momentum identity, i.e.\ whether the
      radial defect's dipole (V56, anti-parallel to the momentum on a
      nearly constant core) and the velocity tensor on the core
      (nearly \(\bar V\otimes\bar V\) there) give the trace part a sign in the
      nearly constant case; (b) the second case of V17 (several groups
      within the Gaussian radius) treated with the same energy--Hessian
      form, the cross terms between groups at distances up to \(R_G\)
      having the multipole form with masses of each group nearly
      zero; (c) the internal consistency check of the four bounds on
      the defect rate against the explicit lower side of V34 for
      small velocity constants.  Choose (a) first.
\item V60 compulsory companion audit and delivery.
\end{enumerate}
\end{programme}

\begin{assessment}[Position after V58]
\label{ass:c3v58-position}
The pressure work on the anchor-alone event is bounded locally by
the halo's energy and the anchor radius; the gate on the event is an
inequality between local geometric quantities.  No full stability or
global-regularity claim is made.
\end{assessment}

% =============================================================================
\section*{Version V58 (third cycle) append-only record}
\addcontentsline{toc}{section}{Version V58 (third cycle) append-only record}
% =============================================================================

The complete V57 source of the third cycle was retained before this
continuation.  Its SHA--256 digest was
\begin{center}\scriptsize\ttfamily
6e3ebea8c9508fbf1e6dab5b2abc74f7d0d4fe2d288297de02e571b6d8ad4b1b.
\end{center}
V58 proved the whole-group bound and the gate in geometric form.
No full regularity claim is made.

% =============================================================================
\section{V59: the anchor-gate block consolidated, and the direction for V61--V65}
\label{sec:c3v59-entry}
% =============================================================================

Items (AC1) and (AC2) of Programme~\ref{prog:c3v59} are carried out.
The anchor-gate block V56--V58 is consolidated, and the fourth bound
on the defect rate is recorded beside the three of V48.  For the
direction, candidate (a) of V58 is chosen: the nearly constant core
on the anchor-alone event, where the radial defect's dipole is
nearly anti-parallel to the momentum and the velocity tensor is
nearly the momentum's tensor, so that the local pressure work has
an explicit leading form whose dipole part vanishes for a radial
cut-off and whose trace part is the anchor's total radial-defect
mass, itself exponentially small; the question for V61--V62 is the
quadrupole order.  No global regularity claim is made.  The next
compulsory companion audit is V60.

% =============================================================================
\section{The anchor-gate block}
\label{sec:c3v59-block}
% =============================================================================

\begin{theorem}[The anchor-gate block, V56--V58]
\label{thm:c3v59-block}
On the anchor-alone event \(\mathcal E_0\) at a homogeneous type-I point,
on the trapped sub-windows:
\begin{enumerate}[label=\textup{(\arabic*)},leftmargin=3.2em]
\item (Masses and dipole, V56.)  The radial-defect masses of the
      anchor's neighbourhood sum to at most \(T_2(R_0)^{1/2}\|Y\|_G\) in
      absolute value, an explicit Gaussian tail at the anchor radius;
      the dipole moment of the radial defect is \(2\nu|G|\bar Y=-\nu|G|\bar V-2\nu|G|\overline{N(V)}\),
      orthogonal to the momentum at extremal momentum and within
      \(\frac\nu4|G|^{1/2}\varphi^{1/2}\) of \(-\nu|G|\bar V\) on a nearly constant core.
\item (One and two cells, V57.)  A cell's self-term is
      \(-\frac1{6\nu}\int\chi|V|^2f_\alpha\) plus a traceless part bounded by
      \(\frac1{2\nu}(\frac23)^{1/2}\|V\|_{L^4(U'')}^2\|f_\alpha\|_2\); two cells of opposite mass
      \(\pm m\) at distance \(\rho\) have cross work \(\frac m{8\pi\nu\rho^3}(Q_\alpha-Q_\beta)\)
      to leading order, negative when the flow is more azimuthal
      about the inward cell.
\item (Whole group, V58.)  \(|\mathcal B_{U_0}|\le\mathcal M_{U_0}+C_{11}\varepsilon_{\rm th}\|Y\|_G\) with
      \(\mathcal M_{U_0}=\frac{\rho_0}{2\nu}(\frac13+(\frac23)^{1/2})\|V\|_{L^4(U'')}^2\|Y\|_{L^2(G;U_0)}\); the
      gate on \(\mathcal E_0\) is \eqref{eq:c3v58-pointwise} at an
      extremal-Dirichlet configuration in \(\mathcal E_0\) and
      \eqref{eq:c3v58-mean} in the mean.
\end{enumerate}
\end{theorem}

\begin{proof}
Proposition~\ref{prop:c3v56-total}; Theorems~\ref{thm:c3v56-dipole},
\ref{thm:c3v57-one-cell}, \ref{thm:c3v57-two-cell}, \ref{thm:c3v58-group},
\ref{thm:c3v58-gate}.
\end{proof}

\begin{proposition}[Fourth bound on the defect rate]
\label{prop:c3v59-fourth}
In case (a) of Theorem~\ref{thm:c3v17-cases}, with \(\mathfrak m(\mathcal E_0)>0\),
\begin{equation}
 \underline\delta\ \le\ \frac{2}{\mathfrak m(\mathcal E_0)}\int_{\mathcal E_0}\Bigl[\kappa(2\mathcal L)^{1/2}\|Y\|_G+\mathcal M_{U_0}\Bigr]\dd\mathfrak m+O(\varepsilon_{\rm th}^{1/4})+O\bigl((1+R_G^4)e^{-R_G^2/(8\nu)}\bigr)
 \ \le\ 2C_Y\Bigl[\kappa(2C_{\mathcal L})^{1/2}+\frac{\rho_0}{2\nu}\Bigl(\frac13+\Bigl(\frac23\Bigr)^{1/2}\Bigr)C_\infty K_6|U''|^{1/3}\Bigr]+o(1),
 \label{eq:c3v59-fourth}
\end{equation}
a bound local in the anchor (the halo's energy and the anchor
radius), beside the window bound (V33), the Lamb--Bernoulli bound
(V41) and the supremum bound (V47).
\end{proposition}

\begin{proof}
\eqref{eq:c3v58-mean} divided by \(\frac12\mathfrak m(\mathcal E_0)\); then \(\|Y\|_G\le C_Y\),
\(\mathcal L\le C_{\mathcal L}\), \(\|V\|_{L^4(U'')}^2\le C_\infty e_{U''}^{1/2}\le C_\infty K_6|U''|^{1/3}\), \(\|Y\|_{L^2(G;U_0)}\le C_Y\).
\end{proof}

\begin{assessment}[What the block settles]
\label{ass:c3v59-settles}
On the anchor-alone event the gate is now an inequality between
local geometric quantities of a bounded configuration, with every
constant explicit and the pressure eliminated in favour of the
velocity tensor on the halo; its sign structure is explicit for one
and two cells.  The block does not close the gate, because the
class does not fix the orientation of the velocity tensor relative
to the radial-defect potential.  It does isolate the one case in
which the orientation is nearly fixed: the nearly constant core,
where the velocity tensor is nearly \(\bar V\otimes\bar V\) and the radial
defect's dipole is nearly \(-\nu|G|\bar V\).  That case is the direction.
\end{assessment}

% =============================================================================
\section{Direction for V61--V65: the nearly constant core on the anchor-alone event}
\label{sec:c3v59-direction}
% =============================================================================

\begin{assessment}[The nearly constant anchor]
\label{ass:c3v59-direction}
On \(\mathcal E_0\) with dominant level \(0\) and \(\eta=16\,\mathrm{Var}\) small, write
\(V=\bar V+\widetilde V\) and \(V\otimes V=\bar V\otimes\bar V+(\bar V\otimes\widetilde V+\widetilde V\otimes\bar V)+\widetilde V\otimes\widetilde V\).  In
\eqref{eq:c3v58-decomposition} the trace part is
\(-\frac1{6\nu}|\bar V|^2\sum_{\beta\in U_0}m_\beta\) plus terms linear and quadratic in
\(\widetilde V\), the first at most \(\frac1{6\nu}|\bar V|^2T_2(R_0)^{1/2}\|Y\|_G\) by V56 and
the others small in \(\eta\) at the exponential price \(e^{\rho_0^2/(8\nu)}\) of
converting \(L^2(G)\) on \(U''\) to \(L^2\); the traceless part is
\((\bar V\otimes\bar V)^\circ:\int\chi_{U'}H_{U_0}\) plus the same corrections.  V61
computes \(\int\chi_{U'}\nabla^2\psi_{U_0}=\int\psi_{U_0}\nabla^2\chi_{U'}\) for a radial
cut-off \(\chi_{U'}(z)=\chi_0(|z|)\) supported in a shell around \(U''\): the
dipole term of \(\psi_{U_0}\) integrates to zero against \(\nabla^2\chi_{U'}\) by
oddness, and the leading term is the quadrupole of \(f_{U_0}\), of
size at most the second moment \(\int_{U_0}|z|^2|f_{U_0}|\) times the shell's
\(\sup|\nabla^2\chi_0|\) times the inverse fourth power of its radius.  V62
states the resulting explicit form of the local pressure work on
the nearly constant anchor: \(-\frac1{6\nu}|\bar V|^2\sum m_\beta+\frac1{2\nu}(\bar V\otimes\bar V)^\circ:\mathsf Q_{U_0}+O(\eta^{1/2}e^{\rho_0^2/(8\nu)})+O(\varepsilon_{\rm th})\),
with \(\mathsf Q_{U_0}\) the quadrupole tensor of the radial defect against
the cut-off; V63 asks whether the quadrupole tensor of the radial
defect on a nearly constant core has a form fixed by the profile
equation (through \(Y\approx-\frac12\bar V\) on the core, which makes
\(f\approx-\frac12(z\cdot\bar V)G\) there and its quadrupole explicit), and
computes the model value; V64 consolidates; V65 audits.  No global
regularity claim is made.
\end{assessment}

% =============================================================================
\section{Integrated V59 conclusion and next acquisition order}
\label{sec:c3v59-conclusion}
% =============================================================================

\begin{theorem}[What V59 establishes]
\label{thm:c3v59-conclusion}
Relative to the two legacy cycles and V1--V58: the consolidated
anchor-gate block (Theorem~\ref{thm:c3v59-block}), the fourth bound
on the defect rate (Proposition~\ref{prop:c3v59-fourth}) and the
direction.  The full Navier--Stokes stability/global-regularity
theorem remains unproved.
\end{theorem}

\begin{proof}
The cited results.
\end{proof}

\begin{programme}[V60 acquisition order]
\label{prog:c3v60}
\begin{enumerate}[label=\textup{(AC\arabic*)},leftmargin=4.8em]
\item The compulsory companion audit: read the current SM and YM
      notes and record their digests and decision.
\item Record the position after sixty versions; delivery.
\end{enumerate}
\end{programme}

\begin{assessment}[Position after V59]
\label{ass:c3v59-position}
The anchor-gate block is consolidated; the direction is the nearly
constant anchor and the quadrupole of its radial defect.  No full
stability or global-regularity claim is made.
\end{assessment}

% =============================================================================
\section*{Version V59 (third cycle) append-only record}
\addcontentsline{toc}{section}{Version V59 (third cycle) append-only record}
% =============================================================================

The complete V58 source of the third cycle was retained before this
continuation.  Its SHA--256 digest was
\begin{center}\scriptsize\ttfamily
791b1ec36d7c55f765d0a0cca18567a8858b4d90ee4ef4741dc8b59d3615f58c.
\end{center}
V59 consolidated the anchor-gate block and fixed the direction.  No
full regularity claim is made.

% =============================================================================
\section{V60: compulsory companion audit and the position after sixty versions}
\label{sec:c3v60-entry}
% =============================================================================

This is the sixtieth version of the third cycle.  The compulsory
review of the two companion programmes is performed first, and the
position after sixty versions is recorded with the plan for the
block V61--V65.  No global regularity claim is made.  The next
compulsory companion audit is V65.

% =============================================================================
\section{Compulsory V60 review of the current companion notes}
\label{sec:c3v60-companion-audit}
% =============================================================================

The current companion files in \texttt{latex\_notes} were inspected
read-only.  Both programmes have moved since the V55 audit: the SM
programme from V36 to V37 of its seventh series, the YM programme
from V13 to V14 of its sixth series.  The frozen checkpoints are
\begin{align}
 &\texttt{Renewal\_SM\_Einstein\_Continuum\_Research\_Notes\_V37.tex},
 \notag\\[-2pt]
 &\qquad
 \texttt{2a65fcc77aeb026758ca8f8e6be1ea419ffc787926c2ec15239aa32a1e8210d5},
 \label{eq:c3v60-SM-checkpoint}\\
 &\texttt{Renewal\_Yang\_Mills\_Mass\_Gap\_Research\_Notes\_V14.tex},
 \notag\\[-2pt]
 &\qquad
 \texttt{b08312b2b8e58f5b91ba81b504899df6478cb0377873c596423ca0feed7f35bc}.
 \label{eq:c3v60-YM-checkpoint}
\end{align}
(SM V37 has 15031 lines; YM V14 has 6277 lines.)

\emph{SM V37.}  The SM programme made an adiabatic acquisition of a
two-tail reference: why the construction is upstream of its
condition FBSC, a discrete localization estimate, an adiabatic
two-tail gap, the bulk paths already supplied by its V34, and the
obstruction to a bulk path.

\emph{YM V14.}  The YM programme audited its two attached programmes
and studied an attached-programme incidence: a graph-animal
two-anchor kernel, the passage from an anchor kernel to a
grouping-forest bound, and the missing incidence bridge.

\begin{theorem}[V60 companion-audit decision]
\label{thm:c3v60-companion-decision}
Nothing is imported from SM V37 or YM V14.  Neither bears on the
third cycle's chapters or on the gates.
\end{theorem}

\begin{proof}
The SM results concern adiabatic references for lattice Dirac
interfaces; the YM results concern kernels of graph animals in a
cluster expansion.  None is a Navier--Stokes statement, and none
concerns the ledger, the extremal method, the energy--Hessian form
or the anchor-gate block.
\end{proof}

% =============================================================================
\section{Position after sixty versions, and the block V61--V65}
\label{sec:c3v60-position}
% =============================================================================

\begin{assessment}[Position after sixty versions of the third cycle]
\label{ass:c3v60-position-sixty}
Ten chapters are written: cluster energy and the band (V2--V10),
localization (V11--V14), the anchor (V16--V19), the explicit floor
(V21--V24), the variance chapter (V26--V29), the defect block
(V31--V34), the extremal ledger (V34--V39), the angle chapter
(V41--V44), the correlation block (V46--V49), and the anchor-gate
block (V56--V59), with the midpoint internal audit (V51--V55) between
the last two.  Since V50 the cycle has audited itself, correcting
three constants in place, and has reduced the gate on the
anchor-alone event to an inequality between local geometric
quantities of a bounded configuration, with the pressure eliminated
in favour of the velocity tensor on the halo and every constant
explicit (V58); the radial defect of the anchor is a zero-mass dipole
distribution whose dipole moment is the momentum identity (V56), and
the sign structure of the pressure work is explicit for one and two
cells (V57).  The defect rate now has four explicit upper bounds
(V33, V41, V47, V59) and no explicit lower bound (N-C3-3).  The block
V61--V65 follows Assessment~\ref{ass:c3v59-direction}: the nearly
constant anchor, where the orientation of the velocity tensor
relative to the radial-defect potential is nearly fixed, and the
quadrupole of the radial defect against the cut-off decides the
local pressure work; V65 audits.  No global regularity claim is
made.
\end{assessment}

% =============================================================================
\section{Integrated V60 conclusion and next acquisition order}
\label{sec:c3v60-conclusion}
% =============================================================================

\begin{theorem}[What V60 establishes]
\label{thm:c3v60-conclusion}
Relative to the two legacy cycles and V1--V59: the companion-audit
decision (Theorem~\ref{thm:c3v60-companion-decision}) and the position.
The full Navier--Stokes stability/global-regularity theorem remains
unproved.
\end{theorem}

\begin{proof}
The cited records.
\end{proof}

\begin{programme}[V61 acquisition order]
\label{prog:c3v61}
\begin{enumerate}[label=\textup{(AC\arabic*)},leftmargin=4.8em]
\item For a radial cut-off \(\chi_{U'}(z)=\chi_0(|z|)\) with \(\chi_0=1\) on
      \([0,\rho_1]\) and \(0\) beyond \(\rho_2\) (\(\rho_1\ge\rho_0+2\), \(\rho_2=\rho_1+1\)),
      prove \(\int\chi_{U'}\nabla^2\psi_{U_0}=\int\psi_{U_0}\nabla^2\chi_{U'}\) and expand
      \(\psi_{U_0}\) on the shell \(\rho_1\le|z|\le\rho_2\) in multipoles: mass
      term (at most \(T_2(R_0)^{1/2}\|Y\|_G\) times the shell integral of
      \(|\nabla^2\chi_{U'}|/(4\pi|z|)\)), dipole term (zero by oddness against
      the radial cut-off's Hessian), quadrupole term (explicit in the
      second moment tensor of \(f_{U_0}\)), remainder (third moment over
      the fourth power of the shell radius).
\item State the resulting identity for \((\bar V\otimes\bar V)^\circ:\int\chi_{U'}H_{U_0}\)
      and its explicit bound.
\item The next compulsory companion audit is V65.
\end{enumerate}
\end{programme}

\begin{assessment}[Position after V60]
\label{ass:c3v60-position}
The companion audit is recorded; the nearly-constant-anchor block
begins.  No full stability or global-regularity claim is made.
\end{assessment}

% =============================================================================
\section*{Version V60 (third cycle) append-only record}
\addcontentsline{toc}{section}{Version V60 (third cycle) append-only record}
% =============================================================================

The complete V59 source of the third cycle was retained before this
continuation.  Its SHA--256 digest was
\begin{center}\scriptsize\ttfamily
e22705fc08a63e7a4335e7cb244803d7bf6a5e203179c11963b9899dd0d4b597.
\end{center}
V60 performed the compulsory companion audit and recorded the
position.  No full regularity claim is made.

% =============================================================================
\section{V61: the radial cut-off integrates the Hessian of the anchor's potential to its mass alone}
\label{sec:c3v61-entry}
% =============================================================================

Items (AC1) and (AC2) of Programme~\ref{prog:c3v61} are carried out,
with a stronger outcome than planned.  For a radial cut-off equal to
one on a ball containing the anchor's neighbourhood, the integral of
the Hessian of the neighbourhood's radial-defect potential against
the cut-off is exactly minus one third of the neighbourhood's total
radial-defect mass times the identity: every multipole of order at
least one contributes nothing, not only the dipole, because the
Hessian of a harmonic multipole has zero spherical mean on every
sphere outside the source and zero flux through every such sphere.
Consequently the traceless part of the Hessian integrates to zero
against the cut-off, and on a nearly constant anchor the local
pressure work is the total mass term, exponentially small by V56,
plus terms linear and quadratic in the deviation from the constant,
explicitly small in the square root of the variance.  The quadrupole
analysis planned for V62--V63 is unnecessary.  No global regularity
claim is made.  The next compulsory companion audit is V65.

% =============================================================================
\section{The radial cut-off}
\label{sec:c3v61-cutoff}
% =============================================================================

On \(\mathcal E_0\), with \(U_0\) the anchor's neighbourhood contained in
\(B_{\rho_0}\) (\(\rho_0=r_2+R_{c_\Phi'}+1\)), fix \(\rho_1:=\rho_0+2\), \(\rho_2:=\rho_0+3\), and a
radial cut-off \(\chi_{U'}(z)=\chi_0(|z|)\) with \(\chi_0\in C^\infty\), \(\chi_0=1\) on
\([0,\rho_1]\), \(\chi_0=0\) on \([\rho_2,\infty)\), \(0\le\chi_0\le1\).  This cut-off is one
on the halo \(U'\) (cells within lattice distance \(2\) of \(U_0\)) and
supported in the cells within distance \(3\) of \(U_0\), as
Theorem~\ref{thm:c3v46-pairing} and Theorem~\ref{thm:c3v58-group}
require, and \(U''=B_{\rho_2}\).  Let \(f=f_{U_0}=(z\cdot Y)G\mathbf 1_{U_0}\),
\(M:=\int f=\sum_{\beta\in U_0}m_\beta\), \(\psi=\Gamma_N*f\).

\begin{lemma}[Spheres outside the source see only the mass]
\label{lem:c3v61-spheres}
Let \(f\) be bounded, supported in \(B_{\rho_0}\), and \(\psi=\Gamma_N*f\).  For
every \(R>\rho_0\):
\begin{enumerate}[label=\textup{(\roman*)},leftmargin=3.2em]
\item \(\int_{\partial B_R}\partial_j\psi\,n_i\,\dd S=-\frac13M\delta_{ij}\);
\item \(\int_{\partial B_R}\partial_i\partial_j\psi\,\dd S=0\).
\end{enumerate}
\end{lemma}

\begin{proof}
On \(|z|=R>\rho_0\) the potential is its convergent multipole series
\(\psi=\sum_{\ell\ge0}\psi_\ell\), \(\psi_\ell(z)=\frac1{4\pi}\int f(y)\frac{|y|^\ell}{|z|^{\ell+1}}P_\ell(\hat z\cdot\hat y)\dd y\),
each \(\psi_\ell\) harmonic outside \(B_{\rho_0}\) and of the form
\(r^{-\ell-1}S_\ell(\hat z)\) with \(S_\ell\) a spherical harmonic of degree \(\ell\); the
series and its derivatives converge uniformly on \(\partial B_R\), so both
integrals may be taken term by term.  (i) \(\nabla\psi_\ell\) has components
of harmonic degrees \(\ell-1\) and \(\ell+1\) on the sphere, and \(n_i=\hat z_i\)
has degree \(1\); the spherical integral of a product vanishes unless
the degrees agree, so only \(\ell=0\) (\(\nabla\psi_0\otimes\hat z=-\frac M{4\pi R^2}\hat z\otimes\hat z\),
integral \(-\frac13M\delta_{ij}\)) and \(\ell=2\) can contribute; for \(\ell=2\), with
\(\psi_2=\frac{1}{8\pi r^3}\hat z\cdot\mathsf Q\hat z\) and \(\mathsf Q\) traceless symmetric,
\(\nabla\psi_2=\frac1{8\pi}\bigl(2\mathsf Qz/r^5-5(z\cdot\mathsf Qz)z/r^7\bigr)\), and with
\(\int_{S^2}\hat z_a\hat z_b\hat z_i\hat z_j=\frac{4\pi}{15}(\delta_{ab}\delta_{ij}+\delta_{ai}\delta_{bj}+\delta_{aj}\delta_{bi})\),
\(\int_{\partial B_R}\partial_j\psi_2\,\hat z_i\dd S=\frac1{8\pi R^2}\bigl(2\cdot\frac{4\pi}3\mathsf Q_{ij}-5\cdot\frac{8\pi}{15}\mathsf Q_{ij}\bigr)=0\).
(ii) \(\partial_i\partial_j\psi_\ell\) has components of degrees \(\ell-2\), \(\ell\), \(\ell+2\), so
its spherical integral vanishes for \(\ell\ne2\), and for \(\ell=0\) directly
(\(\int_{S^2}(3\hat z_i\hat z_j-\delta_{ij})=0\)); for \(\ell=2\), \(\partial_i\partial_j\psi_2=r^{-5}\Phi_{ij}(\hat z)\), and by (i) on two
spheres \(\varepsilon<R\) with the divergence theorem on the shell between
them, \(\int_\varepsilon^Rr^{-3}\dd r\int_{S^2}\Phi_{ij}\dd\Omega=0\), so \(\int_{S^2}\Phi_{ij}=0\).
\end{proof}

\begin{theorem}[The cut-off integral of the Hessian]
\label{thm:c3v61-hessian}
With the radial cut-off above and \(f\) supported in \(B_{\rho_0}\),
\begin{equation}
 \boxed{\quad
 \int\chi_{U'}\,\nabla^2\psi\,\dd z\ =\ -\frac13M\,I ,
 \qquad
 \int\chi_{U'}\,H\,\dd z\ =\ 0\quad(H=\nabla^2\psi+\tfrac13fI) .
 \quad}
 \label{eq:c3v61-hessian}
\end{equation}
\end{theorem}

\begin{proof}
\(\int\chi_{U'}\nabla^2\psi=\int_{B_{\rho_1}}\nabla^2\psi+\int_{\rho_1\le|z|\le\rho_2}\chi_0(|z|)\nabla^2\psi\).  The first
is \(\int_{\partial B_{\rho_1}}\nabla\psi\otimes n=-\frac13MI\) by the divergence theorem
(\(\psi\) is smooth) and Lemma~\ref{lem:c3v61-spheres}(i).  The second is
\(\int_{\rho_1}^{\rho_2}\chi_0(r)\bigl(\int_{\partial B_r}\nabla^2\psi\,\dd S\bigr)\dd r=0\) by
Lemma~\ref{lem:c3v61-spheres}(ii).  Since \(\chi_{U'}=1\) on the support of
\(f\), \(\int\chi_{U'}f=M\), and the second identity follows.
\end{proof}

\begin{remark}[Against the plan]
\label{rem:c3v61-plan}
The plan of V59--V60 expected the dipole term to vanish by oddness
and the quadrupole term to survive; the quadrupole also vanishes,
and so does every higher multipole, for the reason of
Lemma~\ref{lem:c3v61-spheres}: a radial cut-off weights spheres, and
on spheres outside the source the Hessian of the potential has
neither mean nor flux beyond the mass term.  The choice of a radial
cut-off is therefore not a convenience but the reason the anchor's
local pressure work has the simple form below.
\end{remark}

% =============================================================================
\section{The local pressure work on the nearly constant anchor}
\label{sec:c3v61-core}
% =============================================================================

\begin{theorem}[Local pressure work on the anchor, exact form]
\label{thm:c3v61-work}
On the trapped sub-windows of \(\mathcal E_0\), with \(\bar V\) the Gaussian mean
of \(V\) and \(\widetilde V=V-\bar V\),
\begin{equation}
 \boxed{\quad
 \frac1{2\nu}\int_{U_0}(P_V-c)(z\cdot Y)G
 =-\frac{|\bar V|^2}{6\nu}\,M
 +\frac1{2\nu}\int\chi_{U'}\bigl(\bar V\otimes\widetilde V+\widetilde V\otimes\bar V+\widetilde V\otimes\widetilde V\bigr):\nabla^2\psi
 +E_{\rm far} ,
 \quad}
 \label{eq:c3v61-work}
\end{equation}
with \(|M|\le T_2(R_0)^{1/2}\|Y\|_G\) and \(|E_{\rm far}|\le C_{11}\varepsilon_{\rm th}\|Y\|_G\).  If moreover
the dominant level is \(0\) with \(\eta=16\,\mathrm{Var}<\frac12\), then
\begin{equation}
 \Bigl|\frac1{2\nu}\int\chi_{U'}\bigl(\bar V\otimes\widetilde V+\widetilde V\otimes\bar V+\widetilde V\otimes\widetilde V\bigr):\nabla^2\psi\Bigr|
 \ \le\ \frac{\rho_0}{2\nu}\Bigl[2\Bigl(\frac\varphi{|G|}\Bigr)^{1/2}+2C_\infty\Bigr]e^{\rho_2^2/(8\nu)}(\varphi\eta)^{1/2}\,\|Y\|_{L^2(G;U_0)} ,
 \label{eq:c3v61-core}
\end{equation}
so that the local pressure work on a nearly constant anchor is at
most
\begin{equation}
 \frac{\varphi}{6\nu|G|}T_2(R_0)^{1/2}\|Y\|_G+\frac{\rho_0}{2\nu}\Bigl[2\Bigl(\frac\varphi{|G|}\Bigr)^{1/2}+2C_\infty\Bigr]e^{\rho_2^2/(8\nu)}(\varphi\eta)^{1/2}\|Y\|_G+C_{11}\varepsilon_{\rm th}\|Y\|_G
 =:\ \mathcal S_{\rm core}\,\|Y\|_G
 \label{eq:c3v61-small}
\end{equation}
in absolute value.
\end{theorem}

\begin{proof}
Theorem~\ref{thm:c3v46-pairing} with \(V\otimes V=\bar V\otimes\bar V+(\bar V\otimes\widetilde V+\widetilde V\otimes\bar V)+\widetilde V\otimes\widetilde V\)
and \((\bar V\otimes\bar V):\int\chi_{U'}\nabla^2\psi=-\frac13M|\bar V|^2\) by
Theorem~\ref{thm:c3v61-hessian}; \(|M|\) by Proposition~\ref{prop:c3v56-total}.
For \eqref{eq:c3v61-core}: Cauchy--Schwarz with
\(\|\nabla^2\psi\|_{L^2}=\|f\|_{L^2}\le\rho_0\|Y\|_{L^2(G;U_0)}\) (the Riesz isometry,
\eqref{eq:c3v57-H}, and \(G\le1\)); \(\|\chi_{U'}(\bar V\otimes\widetilde V+\widetilde V\otimes\bar V)\|_{L^2}\le2|\bar V|\|\widetilde V\|_{L^2(B_{\rho_2})}\)
and \(\|\chi_{U'}\widetilde V\otimes\widetilde V\|_{L^2}\le\|\widetilde V\|_{L^\infty}\|\widetilde V\|_{L^2(B_{\rho_2})}\le2C_\infty\|\widetilde V\|_{L^2(B_{\rho_2})}\);
\(\|\widetilde V\|_{L^2(B_{\rho_2})}\le e^{\rho_2^2/(8\nu)}\|\widetilde V\|_G\le e^{\rho_2^2/(8\nu)}(\varphi\eta)^{1/2}\)
(\(G\ge e^{-\rho_2^2/(4\nu)}\) on the ball) and \(|\bar V|\le(\varphi/|G|)^{1/2}\); then
\(|\bar V|^2\le\varphi/|G|\) in the mass term.
\end{proof}

\begin{assessment}[Item (AC2): the nearly constant anchor does no pressure work]
\label{ass:c3v61-none}
On a nearly constant anchor-alone configuration the pressure work on
the neighbourhood is at most \(\mathcal S_{\rm core}\|Y\|_G\), with \(\mathcal S_{\rm core}\)
explicit and small: exponentially small in the anchor radius through
\(T_2(R_0)^{1/2}\), of order \(\eta^{1/2}\) through the deviation (at the
explicit price \(e^{\rho_2^2/(8\nu)}\)), and of order \(\varepsilon_{\rm th}\) through the
far field.  With the Lamb piece already of order \(\eta^{1/2}\) on such a
core (Proposition~\ref{prop:c3v42-core}) and the far cells' cross
term of order \(\varepsilon_{\rm th}^{1/4}\) (V11) and the tail beyond \(R_G\)
(V16), the whole cross term \(\langle N,Y\rangle_G\) on a nearly constant
anchor-alone configuration is small: it cannot supply the
anti-alignment \(-\|Y\|_G^2\) demanded at an extremal-Dirichlet
configuration, nor its share of \(-\underline\delta\) in the mean.  V62 states
the two consequences: the defect at an extremal-Dirichlet nearly
constant anchor is explicitly small, and the blow-up measure cannot
be concentrated on nearly constant anchor-alone configurations.  No
global regularity claim is made.
\end{assessment}

% =============================================================================
\section{Integrated V61 conclusion and next acquisition order}
\label{sec:c3v61-conclusion}
% =============================================================================

\begin{theorem}[What V61 proves]
\label{thm:c3v61-conclusion}
Relative to the two legacy cycles and V1--V60: the spheres lemma
(Lemma~\ref{lem:c3v61-spheres}), the cut-off integral of the Hessian
(Theorem~\ref{thm:c3v61-hessian}) and the exact form and smallness of
the local pressure work on the nearly constant anchor
(Theorem~\ref{thm:c3v61-work}).  The full Navier--Stokes
stability/global-regularity theorem remains unproved.
\end{theorem}

\begin{proof}
The cited results.
\end{proof}

\begin{programme}[V62 acquisition order]
\label{prog:c3v62}
\begin{enumerate}[label=\textup{(AC\arabic*)},leftmargin=4.8em]
\item Assemble the full bound on \(|\langle N,Y\rangle_G|\) on a nearly constant
      anchor-alone configuration: Lamb piece
      (Proposition~\ref{prop:c3v42-core}), kinetic Bernoulli piece
      (Proposition~\ref{prop:c3v42-core}, restricted to \(U_0\) with the far
      part in \(X_F\)), local pressure work (\eqref{eq:c3v61-small}),
      far cells (Theorem~\ref{thm:c3v11-far}), tail beyond \(R_G\)
      (Theorem~\ref{thm:c3v16-localized}); state
      \(|\langle N,Y\rangle_G|\le\mathcal S\,\|Y\|_G+\mathcal S'\) with explicit \(\mathcal S,\mathcal S'\).
\item The two consequences: (a) at a maximizer or minimizer of \(\mathcal L\)
      that is a nearly constant anchor-alone configuration,
      \(\|Y\|_G\le\mathcal S+\mathcal S'/\|Y\|_G\), i.e.\ \(\|Y\|_G\le\mathcal S+\mathcal S'^{1/2}\); (b) for every
      blow-up measure \(\mathfrak m\), with \(A\) the set of nearly constant
      anchor-alone configurations (\(\eta\le\eta_1\) on \(\mathcal E_0\)),
      \(\mathfrak m(A^c)\ge(\underline\delta-\mathcal S C_Y-\mathcal S')/(\overline M C_Y)\) whenever positive.
\item The next compulsory companion audit is V65.
\end{enumerate}
\end{programme}

\begin{assessment}[Position after V61]
\label{ass:c3v61-position}
A radial cut-off integrates the Hessian of the anchor's potential to
its mass alone; the nearly constant anchor does no pressure work
beyond explicit small terms.  No full stability or global-regularity
claim is made.
\end{assessment}

% =============================================================================
\section*{Version V61 (third cycle) append-only record}
\addcontentsline{toc}{section}{Version V61 (third cycle) append-only record}
% =============================================================================

The complete V60 source of the third cycle was retained before this
continuation.  Its SHA--256 digest was
\begin{center}\scriptsize\ttfamily
4b4cfe9f8b58c4e2ac1ed4359fe9388798dc498cf6cbd9dbebf14935ec3f9cfd.
\end{center}
V61 proved the cut-off integral of the Hessian and the smallness of
the local pressure work on the nearly constant anchor.  No full
regularity claim is made.

% =============================================================================
\section{V62: the cross term on a nearly constant anchor is small, and two consequences}
\label{sec:c3v62-entry}
% =============================================================================

Items (AC1) and (AC2) of Programme~\ref{prog:c3v62} are carried out.
(AC1): on a nearly constant anchor-alone configuration the whole
cross term is bounded by an explicit small multiple of the defect
plus an explicit small constant, the pieces being the Lamb pairing,
the kinetic Bernoulli work, the local pressure work of V61, the far
cells and the tail beyond the Gaussian radius.  (AC2): at a
configuration of extremal Dirichlet form that is a nearly constant
anchor, the defect is explicitly small; and for every blow-up
measure, the set of nearly constant anchor-alone configurations
cannot carry all the measure: its complement has explicit positive
measure whenever the defect rate exceeds the explicit small
quantities.  These are the first signed statements of the cycle on
the nearly constant core, the configuration that the variance
chapter could not exclude.  No global regularity claim is made.
The next compulsory companion audit is V65.

% =============================================================================
\section{The cross term on a nearly constant anchor}
\label{sec:c3v62-cross}
% =============================================================================

Throughout, the configuration is on the trapped sub-windows of the
anchor-alone event \(\mathcal E_0\), with dominant level \(0\) and
\(\eta=16\,\mathrm{Var}\le\eta_1<\frac12\); \(\bar V\), \(\widetilde V\), \(M\), \(\rho_0\), \(\rho_2\), \(T_2(R_0)\) are as
in V56 and V61, \(c_r=(\frac18+2C_r^2)^{1/2}\).  Put
\begin{align}
 \mathcal S(\eta)&:=\kappa(2\varphi c_r\eta)^{1/2}
 +\frac{\varphi}{|G|^{1/2}}\Bigl(\frac{(3+4c_r)\eta}\nu\Bigr)^{1/2}+\kappa\bigl(\varphi\eta(3+4c_r)\bigr)^{1/2}
 +\frac{\varphi}{4\nu|G|}T_2(R_0)^{1/2}+\mathcal S_{\rm core}(\eta),
 \label{eq:c3v62-S}\\
 \mathcal S'&:=C_7\varepsilon_{\rm th}^{1/4}C_Y+C_{\rm class}(1+R_G^4)e^{-R_G^2/(8\nu)} ,
 \label{eq:c3v62-Sprime}
\end{align}
with \(\mathcal S_{\rm core}\) of \eqref{eq:c3v61-small}; \(\mathcal S(\eta)\to\frac{\varphi}{4\nu|G|}T_2(R_0)^{1/2}+\frac{\varphi}{6\nu|G|}T_2(R_0)^{1/2}+C_{11}\varepsilon_{\rm th}\)
as \(\eta\to0\), and \(\mathcal S'\) is independent of the configuration.

\begin{theorem}[The cross term on a nearly constant anchor]
\label{thm:c3v62-cross}
\begin{equation}
 \boxed{\quad
 |\langle N(V),Y\rangle_G|\ \le\ \mathcal S(\eta)\,\|Y\|_G+\mathcal S' .
 \quad}
 \label{eq:c3v62-cross}
\end{equation}
\end{theorem}

\begin{proof}
On \(\mathcal E_0\), \(\langle N,Y\rangle_G=X_{U_0}+X_F+(\text{tail beyond }R_G)\) with
\(|X_F|\le C_7\varepsilon_{\rm th}^{1/4}C_Y\) (Theorem~\ref{thm:c3v11-far}) and the tail at
most \(C_{\rm class}(1+R_G^4)e^{-R_G^2/(8\nu)}\) (Theorem~\ref{thm:c3v16-localized}),
which give \(\mathcal S'\).  \(X_{U_0}=\int_{U_0}\ell\cdot Y\,G+\frac1{4\nu}\int_{U_0}|V|^2(z\cdot Y)G+\frac1{2\nu}\int_{U_0}(P_V-c)(z\cdot Y)G\).
Lamb piece: \(|\int_{U_0}\ell\cdot YG|\le C_\infty\|\Omega\|_G\|Y\|_G\le\kappa(2\varphi c_r\eta)^{1/2}\|Y\|_G\)
(Proposition~\ref{prop:c3v42-core}, the restriction to \(U_0\) only
lowering the bound).  Kinetic piece: with \(|V|^2=|\bar V|^2+2\bar V\cdot\widetilde V+|\widetilde V|^2\),
the constant part is \(\frac{|\bar V|^2}{4\nu}\int_{U_0}(z\cdot Y)G=\frac{|\bar V|^2}{4\nu}M\), at most
\(\frac{\varphi}{4\nu|G|}T_2(R_0)^{1/2}\|Y\|_G\) (Proposition~\ref{prop:c3v56-total}); the
cross and quadratic parts are bounded as in
Proposition~\ref{prop:c3v42-core} (restricted to \(U_0\)), giving the
second and third terms of \eqref{eq:c3v62-S}.  Pressure piece:
\eqref{eq:c3v61-small}.
\end{proof}

% =============================================================================
\section{Two consequences}
\label{sec:c3v62-consequences}
% =============================================================================

\begin{theorem}[The defect at an extremal-Dirichlet nearly constant anchor is small]
\label{thm:c3v62-extremal}
If a maximizer or minimizer of \(\mathcal L\) over the family is a nearly
constant anchor-alone configuration (on \(\mathcal E_0\), dominant level \(0\),
\(\eta\le\eta_1\)), then at it
\begin{equation}
 \boxed{\quad
 \|Y\|_G\ \le\ \mathcal S(\eta)+\mathcal S'^{1/2} .
 \quad}
 \label{eq:c3v62-extremal}
\end{equation}
\end{theorem}

\begin{proof}
At extremal \(\mathcal L\), \(\|Y\|_G^2=-\langle N,Y\rangle_G\le\mathcal S\|Y\|_G+\mathcal S'\) by
\eqref{eq:c3v62-cross}; \(x^2\le ax+b\) with \(a,b\ge0\) gives \(x\le a+b^{1/2}\).
\end{proof}

\begin{theorem}[The blow-up measure is not concentrated on nearly constant anchors]
\label{thm:c3v62-measure}
Let \(\mathfrak m\) be a blow-up measure (alternative (B)) at a homogeneous
type-I point, and let \(A\subset\mathcal F(x_0)\) be the set of configurations that
are on the trapped sub-windows of \(\mathcal E_0\) with dominant level \(0\) and
\(\eta\le\eta_1\).  Then
\begin{equation}
 \boxed{\quad
 \mathfrak m(A^c)\ \ge\ \frac{\underline\delta-\mathcal S(\eta_1)C_Y-\mathcal S'}{\overline M\,C_Y}
 \qquad\text{whenever the numerator is positive,}
 \quad}
 \label{eq:c3v62-measure}
\end{equation}
with \(\overline M\) of \eqref{eq:c3v41-MNY}.
\end{theorem}

\begin{proof}
\(-\underline\delta\ge\int\langle N,Y\rangle_G\dd\mathfrak m=\int_A+\int_{A^c}\).  On \(A\),
\(\langle N,Y\rangle_G\ge-\mathcal S(\eta)\|Y\|_G-\mathcal S'\ge-\mathcal S(\eta_1)C_Y-\mathcal S'\) (the terms of
\(\mathcal S\) are nondecreasing in \(\eta\)); on \(A^c\), \(\langle N,Y\rangle_G\ge-\overline M\|Y\|_G\ge-\overline MC_Y\)
(Theorem~\ref{thm:c3v41-MNY}).  Hence
\(-\underline\delta\ge-(\mathcal S(\eta_1)C_Y+\mathcal S')\mathfrak m(A)-\overline MC_Y\mathfrak m(A^c)\ge-(\mathcal S(\eta_1)C_Y+\mathcal S')-\overline MC_Y\mathfrak m(A^c)\).
\end{proof}

\begin{assessment}[What the two consequences say]
\label{ass:c3v62-reading}
The variance chapter left one configuration untouched by every
floor: the nearly constant Gaussian core, whose variance the class
allows to be exponentially small (V22--V23) and whose only known
constraint was transience over windows (V28).  On the anchor-alone
event the anchor-gate block now says more.  First, such a core
cannot be a configuration of extremal Dirichlet form unless its
defect is explicitly small, \eqref{eq:c3v62-extremal}: the
configurations at which the Dirichlet form is largest or smallest
are, if nearly constant anchors, nearly self-similar.  Second, the
blow-up statistics cannot be carried by nearly constant anchors
alone, \eqref{eq:c3v62-measure}: the anti-alignment that the defect
rate demands must come from configurations that are either not
anchor-alone or not nearly constant, and their measure is bounded
below explicitly.  Neither statement is an exclusion: the first is
compatible with a small defect at the extremal configuration (the
class does not bound the defect below pointwise), and the second is
compatible with any positive measure on \(A^c\).  What they close is
the possibility that a homogeneous type-I blow-up is, statistically,
a nearly constant anchor: it is not.  The explicit quantities
\(\mathcal S(\eta_1)C_Y+\mathcal S'\) are small in the threshold, the anchor tail and
\(\eta_1^{1/2}\), and the condition of positivity of the numerator in
\eqref{eq:c3v62-measure} is an explicit inequality between the
defect rate and these; the defect rate has no explicit lower bound
(N-C3-3), so whether the condition holds is not decided by the
series, and the statement is recorded as conditional on it.
\end{assessment}

% =============================================================================
\section{Integrated V62 conclusion and next acquisition order}
\label{sec:c3v62-conclusion}
% =============================================================================

\begin{theorem}[What V62 proves]
\label{thm:c3v62-conclusion}
Relative to the two legacy cycles and V1--V61: the cross term on a
nearly constant anchor (Theorem~\ref{thm:c3v62-cross}), the defect at
an extremal-Dirichlet nearly constant anchor
(Theorem~\ref{thm:c3v62-extremal}) and the measure statement
(Theorem~\ref{thm:c3v62-measure}).  The full Navier--Stokes
stability/global-regularity theorem remains unproved.
\end{theorem}

\begin{proof}
The cited results.
\end{proof}

\begin{programme}[V63 acquisition order]
\label{prog:c3v63}
\begin{enumerate}[label=\textup{(AC\arabic*)},leftmargin=4.8em]
\item Remove the restriction to the anchor-alone event: in case (b) of
      Theorem~\ref{thm:c3v17-cases} (several groups within \(R_G\)), apply
      the radial cut-off and Theorem~\ref{thm:c3v61-hessian} to
      \(U_{R_G}\) (all cells within \(R_G+r_2\)) in place of \(U_0\); the total
      mass of \(f_{U_{R_G}}\) is bounded by \(T_2(R_G-1)^{1/2}\|Y\|_G\), smaller
      still; state the analogue of Theorem~\ref{thm:c3v62-cross} for a
      nearly constant configuration with several groups, and of the
      two consequences.
\item Remove the restriction to the trapped sub-windows if the
      constants permit (the far-cell bounds of V11 hold there; record
      what fails elsewhere), and record the resulting statement for
      every nearly constant configuration of the family.
\item The next compulsory companion audit is V65.
\end{enumerate}
\end{programme}

\begin{assessment}[Position after V62]
\label{ass:c3v62-position}
The cross term on a nearly constant anchor is explicitly small; the
defect at an extremal-Dirichlet nearly constant anchor is small; the
blow-up measure is not concentrated on nearly constant anchors.  No
full stability or global-regularity claim is made.
\end{assessment}

% =============================================================================
\section*{Version V62 (third cycle) append-only record}
\addcontentsline{toc}{section}{Version V62 (third cycle) append-only record}
% =============================================================================

The complete V61 source of the third cycle was retained before this
continuation.  Its SHA--256 digest was
\begin{center}\scriptsize\ttfamily
993aea13747b02b79043a9adf7a18fc45573278707912f6b283c8a819f301551.
\end{center}
V62 proved the smallness of the cross term on a nearly constant
anchor and its two consequences.  No full regularity claim is made.

% =============================================================================
\section{V63: the nearly constant configuration with several groups, and the trapped-window restriction}
\label{sec:c3v63-entry}
% =============================================================================

Items (AC1) and (AC2) of Programme~\ref{prog:c3v63} are carried out.
(AC1): the restriction to the anchor-alone event is removed, at a
price.  With several groups the configuration neighbourhood is the
set of cells within the Gaussian radius, and the radial cut-off of
V61 applies to any ball containing it; but the far-pressure flatness
of V46 needs the cells near the cut-off to be quiet and the
energetic groups beyond it to be far, and groups beyond the Gaussian
radius are unpinned (V16).  A quiet annulus at distance at least an
explicit \(D_{\rm far}\) from every group exists by pigeonhole within an
explicit range of radii, and the cut-off is placed there; the
deviation terms then carry the factor \(e^{\rho^2/(8\nu)}\) at that radius,
which is explicit and of the order of \(\exp(c\,\varepsilon_{\rm th}^{-1/2})\).  The
two consequences of V62 hold for nearly constant multi-group
configurations whose variance is below the inverse of that factor:
an explicit but extremely small range.  The anchor-alone case
remains the meaningful one.  (AC2): outside the trapped sub-windows
the far cells' enstrophy is bounded only by the class and the
far-cell cross term is not small; the statements of V61--V63 are
recorded on the trapped sub-windows, and the obstruction elsewhere
is named.  No global regularity claim is made.  The next compulsory
companion audit is V65.

% =============================================================================
\section{Several groups: the quiet annulus and the cut-off}
\label{sec:c3v63-multi}
% =============================================================================

In case (b) of Theorem~\ref{thm:c3v17-cases} the groups within the
Gaussian radius are contained in \(U_{R_G}\), the cells within lattice
distance \(R_G+r_2\) of the origin; the other groups (at most \(N_C\) in
all, each of lattice diameter at most \(2r_2+4\)) lie at unpinned
positions beyond.  Recall \(R_G\) is fixed by \(e^{-R_G^2/(8\nu)}=\varepsilon_{\rm th}\)
(V16), so \(e^{-R_G^2/(4\nu)}=\varepsilon_{\rm th}^2\).

\begin{lemma}[A quiet annulus]
\label{lem:c3v63-annulus}
Let \(D_{\rm far}:=\bigl(C_{15}N_Ca_C(2r_2+1)^3/\varepsilon_{\rm th}\bigr)^{1/4}\), with \(C_{15}\) the
constant of the far-pressure kernel bound (\(|\nabla\partial_i\partial_j\Gamma_N(w)|\le C_{15}|w|^{-4}\)).
There is \(\rho_1'\in[R_G+r_2+2,\ R_G+r_2+2+(N_C+1)(2D_{\rm far}+2r_2+6)]\) such that
no energetic cell lies at distance less than \(D_{\rm far}\) from the sphere
\(\{|z|=\rho_1'\}\).
\end{lemma}

\begin{proof}
Each group's cells occupy radii in an interval of length at most
\(2r_2+4\); the set of radii within \(D_{\rm far}\) of some energetic cell is
a union of at most \(N_C\) intervals of length at most \(2D_{\rm far}+2r_2+4\);
its total length is less than the length of the range, so a radius
outside it exists.
\end{proof}

\begin{proposition}[The multi-group neighbourhood]
\label{prop:c3v63-multi}
With \(\rho_1'\) of Lemma~\ref{lem:c3v63-annulus}, \(\rho_2':=\rho_1'+1\), the radial
cut-off \(\chi'\) equal to one on \(B_{\rho_1'}\) and vanishing outside \(B_{\rho_2'}\),
\(U':=\) the cells within \(B_{\rho_1'}\), \(f':=(z\cdot Y)G\mathbf 1_{U'}\), \(M':=\int f'\),
\(\psi':=\Gamma_N*f'\), on the trapped sub-windows:
\begin{enumerate}[label=\textup{(\roman*)},leftmargin=3.2em]
\item \(|M'|\le T_2(\rho_1'-1)^{1/2}\|Y\|_G\), and \(T_2(\rho_1'-1)\le T_2(R_G)\le4\pi\varepsilon_{\rm th}^2(2\nu R_G^3+12\nu^2R_G+24\nu^3/R_G)\);
\item the far pressure generated by \((1-\chi')V\otimes V\) has gradient at
      most \(C_{10}'\varepsilon_{\rm th}\) on \(U'\), with \(C_{10}'\) explicit (the quiet cells
      as in Proposition~\ref{prop:c3v46-far}, the energetic groups
      beyond the annulus through \(C_{15}N_Ca_C(2r_2+1)^3D_{\rm far}^{-4}=\varepsilon_{\rm th}\));
\item \(\int\chi'\nabla^2\psi'=-\frac13M'I\);
\item the cross term equals \(X_{U'}\) up to \(O(\varepsilon_{\rm th}^{1/4}C_Y)\) from the
      quiet cells within \(B_{\rho_1'}\) not in the groups and the tail
      beyond \(R_G\) (the cells of \(U'\) beyond \(R_G\) are within the tail's
      bound, and their inclusion in \(U'\) only moves them from the
      tail to the local part, where the same bound applies).
\end{enumerate}
\end{proposition}

\begin{proof}
(i) Proposition~\ref{prop:c3v56-total} with \(U'\) in place of \(U_0\), and
Lemma~\ref{lem:c3v56-tails}.  (ii) Proposition~\ref{prop:c3v46-far}'s
mechanism for the quiet cells; for an energetic cell at distance
\(d\ge D_{\rm far}\) from the cut-off's support, its source is bounded by
its energy \(\le a_C\) and the kernel by \(C_{15}d^{-4}\); the groups beyond
number at most \(N_C\) with at most \((2r_2+1)^3\) cells each.  (iii)
Theorem~\ref{thm:c3v61-hessian} with \(\rho_0\) replaced by \(\rho_1'-1\).  (iv)
Theorems~\ref{thm:c3v11-far} and \ref{thm:c3v16-localized}.
\end{proof}

\begin{theorem}[The nearly constant multi-group configuration]
\label{thm:c3v63-multi}
On the trapped sub-windows, for a configuration with dominant level
\(0\) and \(\eta=16\,\mathrm{Var}<\frac12\), in case (b),
\begin{equation}
 |\langle N(V),Y\rangle_G|\ \le\ \mathcal S_{(b)}(\eta)\|Y\|_G+\mathcal S' ,
 \label{eq:c3v63-cross}
\end{equation}
with \(\mathcal S'\) of \eqref{eq:c3v62-Sprime} and \(\mathcal S_{(b)}(\eta)\) obtained from
\(\mathcal S(\eta)\) of \eqref{eq:c3v62-S} by replacing \(T_2(R_0)\) by \(T_2(\rho_1'-1)\),
\(\rho_0\) by \(\rho_1'-1\), \(\rho_2\) by \(\rho_2'\), and \(C_{11}\) by the constant
\(C_{11}'\) of (ii).  Consequently Theorems~\ref{thm:c3v62-extremal} and
\ref{thm:c3v62-measure} hold with \(\mathcal E_0\) replaced by the whole family
and \(\mathcal S\) by \(\mathcal S_{(b)}\).  The deviation term of \(\mathcal S_{(b)}\) carries
\(e^{\rho_2'^2/(8\nu)}\) with \(\rho_2'\le R_G+r_2+3+(N_C+1)(2D_{\rm far}+2r_2+6)\), of the order
of \(\exp(c\,\varepsilon_{\rm th}^{-1/2})\) for an explicit \(c\), so \(\mathcal S_{(b)}(\eta)\) is small
only for \(\eta\) below \(\exp(-2c\,\varepsilon_{\rm th}^{-1/2})\) times explicit factors.
\end{theorem}

\begin{proof}
The proof of Theorem~\ref{thm:c3v62-cross} with
Proposition~\ref{prop:c3v63-multi} in place of Propositions
\ref{prop:c3v56-total}, \ref{prop:c3v46-far} and Theorem
\ref{thm:c3v61-hessian}; the consequences as in V62.  \(D_{\rm far}\sim\varepsilon_{\rm th}^{-1/4}\)
gives \(\rho_2'^2/(8\nu)\sim c\,\varepsilon_{\rm th}^{-1/2}\).
\end{proof}

\begin{assessment}[Item (AC1): the extension is explicit and weak]
\label{ass:c3v63-weak}
The anchor-alone case is the meaningful one: there the cut-off sits
at the anchor radius, the deviation factor is \(e^{(\rho_0+3)^2/(8\nu)}\) with
\(\rho_0\) of the order of \((4\nu\log(2S_G/c_\Phi'))^{1/2}\), i.e.\ a power of
\(S_G/c_\Phi'\), and \(\mathcal S(\eta)\) is small for \(\eta\) below a power of
\(c_\Phi'/S_G\).  With several groups the unpinned positions force the
cut-off outward by an amount of order \(\varepsilon_{\rm th}^{-1/4}\) to keep the
far pressure flat, and the range of variances covered collapses to
\(\exp(-c\varepsilon_{\rm th}^{-1/2})\).  This is not an artefact of the method but
of the class: an energetic group near the cut-off produces a
pressure gradient of order one on the configuration, and its
velocity tensor must then be part of the local pairing, which the
Gaussian weight does not control at that distance.  The multi-group
statement is recorded for completeness; the results of V61--V62 are
stated for the anchor-alone event.
\end{assessment}

% =============================================================================
\section{The trapped-window restriction}
\label{sec:c3v63-trapped}
% =============================================================================

\begin{firewall}[Item (AC2): outside the trapped sub-windows]
\label{fw:c3v63-trapped}
The bounds of V61--V63 use, through Theorem~\ref{thm:c3v11-far}, that
the far cells have trapped enstrophy (\(\int_{Q_\beta}|\nabla V|^2\le C_1\varepsilon_{\rm th}^{1/2}\)
on quiet cells), which the trap chapter of the second cycle
(C2-V57, C2-V62) supplies on the trapped sub-windows, the parts of
every window after the capture time \(\tau_{\rm cap}\).  Outside them the
quiet cells' enstrophy is bounded only by the class (\(C_\infty'^2\) per
cell), the far Lamb piece is bounded by \(C_\infty(2S_GC_\infty'^2)^{1/2}\|Y\|_G\)
with no smallness, and the far-cell cross term \(X_F\) is of order
one.  The statements of V61--V63 are therefore statements on the
trapped sub-windows; since these cover all of every window beyond a
log-time \(\tau_{\rm cap}\le\ell_M/2\), and the blow-up measure is
invariant, the measure statement \eqref{eq:c3v62-measure} applies to
the restriction of \(\mathfrak m\) to the trapped parts, which carry at least
\(1-\tau_{\rm cap}/\ell_M\ge\frac12\) of the mass of every window; the constants
are adjusted by this factor.  The obstruction elsewhere is the
untrapped enstrophy of quiet cells, and it is not removable by the
class bounds alone.
\end{firewall}

% =============================================================================
\section{Integrated V63 conclusion and next acquisition order}
\label{sec:c3v63-conclusion}
% =============================================================================

\begin{theorem}[What V63 proves]
\label{thm:c3v63-conclusion}
Relative to the two legacy cycles and V1--V62: the quiet annulus
(Lemma~\ref{lem:c3v63-annulus}), the multi-group neighbourhood
(Proposition~\ref{prop:c3v63-multi}) and the nearly constant
multi-group configuration (Theorem~\ref{thm:c3v63-multi}), with the
trapped-window firewall.  The full Navier--Stokes
stability/global-regularity theorem remains unproved.
\end{theorem}

\begin{proof}
The cited results.
\end{proof}

\begin{programme}[V64 acquisition order]
\label{prog:c3v64}
\begin{enumerate}[label=\textup{(AC\arabic*)},leftmargin=4.8em]
\item Consolidate V61--V63 as the nearly-constant-anchor block: the
      cut-off integral of the Hessian, the smallness of the cross
      term on a nearly constant anchor, the two consequences, the
      multi-group extension and the trapped-window restriction.
\item Fix the direction for V66--V70: the complement \(A^c\) of the
      nearly constant anchors carries the anti-alignment; on \(A^c\) the
      configuration is either multi-group within the Gaussian radius
      or spectrally spread (\(\mathrm{Var}\ge\eta_1/16\)); for the spread case,
      determine what the explicit floors and the extremal rows say
      about the cross term, and whether the pointwise spectrum of
      V37 bounds the Lamb pairing on a spread anchor.
\item V65 compulsory companion audit and delivery.
\end{enumerate}
\end{programme}

\begin{assessment}[Position after V63]
\label{ass:c3v63-position}
The nearly constant configuration is treated with several groups at
an explicit but weak level; the statements hold on the trapped
sub-windows.  No full stability or global-regularity claim is made.
\end{assessment}

% =============================================================================
\section*{Version V63 (third cycle) append-only record}
\addcontentsline{toc}{section}{Version V63 (third cycle) append-only record}
% =============================================================================

The complete V62 source of the third cycle was retained before this
continuation.  Its SHA--256 digest was
\begin{center}\scriptsize\ttfamily
c0aca4d330de68579fedd179f3cc286da5124997b58267d2b48f810bbd59b491.
\end{center}
V63 extended the nearly constant statements to several groups and
recorded the trapped-window restriction.  No full regularity claim
is made.

% =============================================================================
\section{V64: the nearly-constant-anchor block consolidated, and the direction for V66--V70}
\label{sec:c3v64-entry}
% =============================================================================

Items (AC1) and (AC2) of Programme~\ref{prog:c3v64} are carried out.
The nearly-constant-anchor block V61--V63 is consolidated.  For the
direction, the complement of the nearly constant anchors is examined:
by the measure statement of V62 it carries the anti-alignment, and
its configurations are either multi-group within the Gaussian
radius or spectrally spread; the spread anchor-alone case is the
next target, with the pointwise spectrum of V37 and the extremal
rows as the tools.  No global regularity claim is made.  The next
compulsory companion audit is V65.

% =============================================================================
\section{The nearly-constant-anchor block}
\label{sec:c3v64-block}
% =============================================================================

\begin{theorem}[The nearly-constant-anchor block, V61--V63]
\label{thm:c3v64-block}
At a homogeneous type-I point, on the trapped sub-windows:
\begin{enumerate}[label=\textup{(\arabic*)},leftmargin=3.2em]
\item (Radial cut-off, V61.)  For a bounded density \(f\) supported in
      \(B_{\rho_0}\) and a radial cut-off equal to one on \(B_{\rho_1}\), \(\rho_1>\rho_0\),
      \(\int\chi\nabla^2(\Gamma_N*f)=-\frac13(\int f)I\): every multipole of order at
      least one contributes nothing.
\item (Nearly constant anchor, V61--V62.)  On the anchor-alone event
      with dominant level \(0\) and \(\eta=16\,\mathrm{Var}\le\eta_1<\frac12\), the local
      pressure work is \(-\frac{|\bar V|^2}{6\nu}M\) plus deviation terms,
      at most \(\mathcal S_{\rm core}(\eta)\|Y\|_G\), and the whole cross term satisfies
      \(|\langle N,Y\rangle_G|\le\mathcal S(\eta)\|Y\|_G+\mathcal S'\) with \(\mathcal S,\mathcal S'\) explicit and small
      (in \(\eta^{1/2}\), the anchor tail \(T_2(R_0)^{1/2}\), and the threshold).
\item (Consequences, V62.)  At a maximizer or minimizer of \(\mathcal L\)
      that is a nearly constant anchor, \(\|Y\|_G\le\mathcal S(\eta)+\mathcal S'^{1/2}\); for
      every blow-up measure, \(\mathfrak m(A^c)\ge(\underline\delta-\mathcal S(\eta_1)C_Y-\mathcal S')/(\overline MC_Y)\)
      whenever positive, \(A\) the nearly constant anchors.
\item (Several groups, V63.)  The same with the cut-off in a quiet
      annulus at distance \(D_{\rm far}\sim\varepsilon_{\rm th}^{-1/4}\) from every group,
      explicit but with the deviation factor \(\exp(c\varepsilon_{\rm th}^{-1/2})\).
\item (Restriction, V63.)  On the trapped sub-windows only; elsewhere
      the far cells' untrapped enstrophy blocks the far-cell bound.
\end{enumerate}
\end{theorem}

\begin{proof}
Theorems~\ref{thm:c3v61-hessian}, \ref{thm:c3v61-work}, \ref{thm:c3v62-cross},
\ref{thm:c3v62-extremal}, \ref{thm:c3v62-measure}, \ref{thm:c3v63-multi};
Firewall~\ref{fw:c3v63-trapped}.
\end{proof}

\begin{assessment}[What the block settles]
\label{ass:c3v64-settles}
The nearly constant Gaussian core, left untouched by the variance
floors (V22--V23) and constrained only in time by the momentum decay
(V28), is now constrained in the statistics: it does no pressure
work on its own neighbourhood beyond explicit small terms, because a
radial cut-off sees only the total radial-defect mass of the
neighbourhood, which is exponentially small in the anchor radius; so
it cannot carry the anti-alignment, cannot be an extremal-Dirichlet
configuration with a large defect, and cannot carry the blow-up
measure.  The mechanism is the harmonic structure of the Newtonian
potential, and it is the one point of the third cycle at which the
nonlocal pressure was not an obstruction but the source of a sign:
the pressure work of a nearly uniform flow against its own radial
defect vanishes to all multipole orders.  The block does not close
the gate: the anti-alignment lives on the complement, and the class
does not fix the orientation there.
\end{assessment}

% =============================================================================
\section{Direction for V66--V70: the spread anchor}
\label{sec:c3v64-direction}
% =============================================================================

\begin{assessment}[The complement of the nearly constant anchors]
\label{ass:c3v64-direction}
By Theorem~\ref{thm:c3v62-measure} the set \(A^c\) carries the
anti-alignment.  Its configurations are of two kinds: multi-group
within the Gaussian radius (case (b) of V17), or anchor-alone with
\(\mathrm{Var}>\eta_1/16\), the spread anchor.  The block V66--V70 takes the
spread anchor.  On it the one-level criterion (V21) gives
\(\hat a_{k_1}\le1-\eta_1\): the spectral mass is not concentrated on one
level, so the profile is not nearly polynomial, and the Lamb vector
is not small; the Lamb pairing on \(U_0\) is bounded by
\(\kappa(2\mathcal L)^{1/2}\|Y\|_G\) (V41) and the pressure work by \(\mathcal M_{U_0}\) (V58),
both of order one.  V66: the pointwise spectrum (V37) on the spread
anchor: \(\hat a_k\le B_k^{\sup}/c_\Phi\) for every \(k\), combined with
\(\mathrm{Var}>\eta_1/16\) and \(k_1\le K_r\), gives a lower bound on the spectral
mass at levels \(k\ge1\) and hence on the Gaussian dissipation
\(\nu\|\nabla\widehat V\|_G^2=\sum\frac k2\hat a_k\ge\frac12(1-\hat a_0)\); state the explicit
dissipation floor on the spread anchor and its consequence for the
extremal rows there.  V67: the Lamb pairing on the spread anchor
through the vorticity: \(\|\Omega\|_G^2=\mathcal E\le\frac2\nu\mathcal L\) above and, from the
dissipation floor, below (\(\mathcal E\ge\) the Gaussian enstrophy floor of
C2-V19 if available, else through \(\|\nabla V\|_G\) and the divergence-free
identity); determine whether the Lamb pairing on a spread anchor is
bounded below in absolute value in the mean, and record.  V68: the
gate on the spread anchor as the sum of a Lamb pairing of order one
and a pressure work of order one, both explicit, and the question of
whether the extremal-Dirichlet row can be realized on a spread
anchor with both pieces negative; the test configurations of the
second cycle (the columnar swirl, the plane wave) evaluated on the
anchor.  V69 consolidates and fixes the direction for V71--V75; V70
audits.  No global regularity claim is made.
\end{assessment}

% =============================================================================
\section{Integrated V64 conclusion and next acquisition order}
\label{sec:c3v64-conclusion}
% =============================================================================

\begin{theorem}[What V64 establishes]
\label{thm:c3v64-conclusion}
Relative to the two legacy cycles and V1--V63: the consolidated
nearly-constant-anchor block (Theorem~\ref{thm:c3v64-block}) and the
direction.  The full Navier--Stokes stability/global-regularity
theorem remains unproved.
\end{theorem}

\begin{proof}
The cited results.
\end{proof}

\begin{programme}[V65 acquisition order]
\label{prog:c3v65}
\begin{enumerate}[label=\textup{(AC\arabic*)},leftmargin=4.8em]
\item The compulsory companion audit: read the current SM and YM
      notes and record their digests and decision.
\item Record the position after sixty-five versions; delivery.
\end{enumerate}
\end{programme}

\begin{assessment}[Position after V64]
\label{ass:c3v64-position}
The nearly-constant-anchor block is consolidated; the direction is
the spread anchor.  No full stability or global-regularity claim is
made.
\end{assessment}

% =============================================================================
\section*{Version V64 (third cycle) append-only record}
\addcontentsline{toc}{section}{Version V64 (third cycle) append-only record}
% =============================================================================

The complete V63 source of the third cycle was retained before this
continuation.  Its SHA--256 digest was
\begin{center}\scriptsize\ttfamily
79fbdbecaee33d6b761798e65c322d4cdd34082fba059a40f113dbe683872389.
\end{center}
V64 consolidated the nearly-constant-anchor block and fixed the
direction.  No full regularity claim is made.

% =============================================================================
\section{V65: compulsory companion audit and the position after sixty-five versions}
\label{sec:c3v65-entry}
% =============================================================================

This is the sixty-fifth version of the third cycle.  The compulsory
review of the two companion programmes is performed first, and the
position after sixty-five versions is recorded with the plan for
the block V66--V70.  No global regularity claim is made.  The next
compulsory companion audit is V70.

% =============================================================================
\section{Compulsory V65 review of the current companion notes}
\label{sec:c3v65-companion-audit}
% =============================================================================

The current companion files in \texttt{latex\_notes} were inspected
read-only.  Both programmes have moved since the V60 audit: the SM
programme from V37 to V39 of its seventh series, the YM programme
from V14 to V16 of its sixth series.  The frozen checkpoints are
\begin{align}
 &\texttt{Renewal\_SM\_Einstein\_Continuum\_Research\_Notes\_V39.tex},
 \notag\\[-2pt]
 &\qquad
 \texttt{ef17e5f9b5c6f2877cafa777d01132055ae161b612c198cae2427d8a6471cf40},
 \label{eq:c3v65-SM-checkpoint}\\
 &\texttt{Renewal\_Yang\_Mills\_Mass\_Gap\_Research\_Notes\_V16.tex},
 \notag\\[-2pt]
 &\qquad
 \texttt{610909f687d0453d2a731a3bb1c255467c80f6e871e95672d6e8b3d522b50c8d}.
 \label{eq:c3v65-YM-checkpoint}
\end{align}
(SM V39 has 15755 lines; YM V16 has 7167 lines.)

\emph{SM V38--V39.}  The SM programme constructed a universal gapped
bulk path, then a certified finite-range approximation of its
condition FBSC: a correction of a naive denominator-clearing route,
an explicit polynomial inverse, a certified surrogate for the face
matrix, a finite interval certificate on parameter cells,
completeness of refinement when a gap exists, and the data still
needed from its V33 assembly.

\emph{YM V15--V16.}  The YM programme obtained exact hypergraph
activities for owned components and a stratified, moment-retaining
form of its condition COHAC, then canonical connected source rows: a
proper-partition short, a canonical analytic activity, a
connected-carrier sufficient condition, cofinal stratified carriers,
and the first missing Yang--Mills estimate.

\begin{theorem}[V65 companion-audit decision]
\label{thm:c3v65-companion-decision}
Nothing is imported from SM V39 or YM V16.  Neither bears on the
third cycle's chapters or on the gates.
\end{theorem}

\begin{proof}
The SM results concern certified approximations of lattice Dirac
face conditions; the YM results concern activities and carriers in a
cluster expansion.  None is a Navier--Stokes statement, and none
concerns the ledger, the extremal method, the energy--Hessian form
or the anchor blocks.
\end{proof}

% =============================================================================
\section{Position after sixty-five versions, and the block V66--V70}
\label{sec:c3v65-position}
% =============================================================================

\begin{assessment}[Position after sixty-five versions of the third cycle]
\label{ass:c3v65-position-sixty-five}
Eleven chapters are written: cluster energy and the band (V2--V10),
localization (V11--V14), the anchor (V16--V19), the explicit floor
(V21--V24), the variance chapter (V26--V29), the defect block
(V31--V34), the extremal ledger (V34--V39), the angle chapter
(V41--V44), the correlation block (V46--V49), the anchor-gate block
(V56--V59) and the nearly-constant-anchor block (V61--V64), with the
midpoint internal audit (V51--V55).  The last block produced the
cycle's first signed statements on the configuration the variance
floors could not exclude: a nearly constant anchor does no pressure
work on its neighbourhood to any multipole order beyond its
exponentially small total radial-defect mass (V61), so it cannot be
an extremal-Dirichlet configuration with a large defect and cannot
carry the blow-up measure (V62).  The complement, the spread anchors
and the multi-group configurations, carries the anti-alignment, and
the block V66--V70 follows Assessment~\ref{ass:c3v64-direction} onto the
spread anchor: the explicit dissipation floor from the pointwise
spectrum (V66), the Lamb pairing through the vorticity (V67), the
gate on the spread anchor with both pieces of order one and the
second cycle's test configurations (V68), consolidation (V69), audit
(V70).  No global regularity claim is made.
\end{assessment}

% =============================================================================
\section{Integrated V65 conclusion and next acquisition order}
\label{sec:c3v65-conclusion}
% =============================================================================

\begin{theorem}[What V65 establishes]
\label{thm:c3v65-conclusion}
Relative to the two legacy cycles and V1--V64: the companion-audit
decision (Theorem~\ref{thm:c3v65-companion-decision}) and the position.
The full Navier--Stokes stability/global-regularity theorem remains
unproved.
\end{theorem}

\begin{proof}
The cited records.
\end{proof}

\begin{programme}[V66 acquisition order]
\label{prog:c3v66}
\begin{enumerate}[label=\textup{(AC\arabic*)},leftmargin=4.8em]
\item The dissipation floor on the spread anchor: for a configuration
      with \(\mathrm{Var}\ge\eta_1/16\), prove \(1-\hat a_0\ge\) an explicit positive
      quantity (from \(\mathrm{Var}\le\sum_k(\mu_k-r)^2\hat a_k\) with \(\hat a_k\le B_k^{\sup}/c_\Phi\)
      and \(k_1\le K_r\), or directly from \(\mathrm{Var}\ge\eta_1/16\) and the
      spectral form), and deduce \(\nu\|\nabla\widehat V\|_G^2\ge\frac12(1-\hat a_0)\), i.e.\
      an explicit floor on the Gaussian dissipation \(d\) on the spread
      anchor.
\item Record the consequence for the extremal rows on the spread
      anchor (the Dirichlet form is bounded below there beyond
      \(\frac14\varphi\)) and for the Lamb bound (\(\mathcal E\) is bounded below).
\item The next compulsory companion audit is V70.
\end{enumerate}
\end{programme}

\begin{assessment}[Position after V65]
\label{ass:c3v65-position}
The companion audit is recorded; the spread-anchor block begins.  No
full stability or global-regularity claim is made.
\end{assessment}

% =============================================================================
\section*{Version V65 (third cycle) append-only record}
\addcontentsline{toc}{section}{Version V65 (third cycle) append-only record}
% =============================================================================

The complete V64 source of the third cycle was retained before this
continuation.  Its SHA--256 digest was
\begin{center}\scriptsize\ttfamily
fb32e6e78c46a7986c1acb8e798d4b8d06c535d6e17e7146c18e9b8b5d5e0e18.
\end{center}
V65 performed the compulsory companion audit and recorded the
position.  No full regularity claim is made.

% =============================================================================
\section{V66: the dissipation floor on spread configurations}
\label{sec:c3v66-entry}
% =============================================================================

Items (AC1) and (AC2) of Programme~\ref{prog:c3v66} are carried out.
(AC1): on every configuration whose variance is not small, the
spectral mass off the constant level is bounded below explicitly,
because the variance is bounded above by an explicit multiple of
that mass through the class bound on the self-similar operator
applied to the profile; hence the Gaussian dissipation of the
normalized profile, which is half the mean Hermite degree, is
bounded below, and so is the dissipation \(d\) and the Dirichlet form.
The floor holds for every spread configuration, anchor or not.
(AC2): the consequences for the extremal rows are recorded; for the
Lamb pairing there is none, because the Gaussian enstrophy is not
bounded below by the dissipation, the radial kinetic energy being
able to absorb it in the closure identity.  No global regularity
claim is made.  The next compulsory companion audit is V70.

% =============================================================================
% ==============================================================================
\section{The variance is bounded by the mass off the constant level}
\label{sec:c3v66-variance}
% =============================================================================

\begin{lemma}[The variance is bounded on the family]
\label{lem:c3v66-varmax}
For every profile of the family and every \(\sigma\),
\(\|L_-V\|_G\le\|Y\|_G+\|N(V)\|_G\le C_Y+C_N^{1/2}\) and
\begin{equation}
 \mathrm{Var}\ \le\ \mathrm{Var}_{\max}:=\frac{(C_Y+C_N^{1/2})^2}{c_\Phi} .
 \label{eq:c3v66-varmax}
\end{equation}
\end{lemma}

\begin{proof}
\(L_-V=Y+N(V)\) and the class bounds; \(\mathrm{Var}=\|L_-V\|_G^2/\varphi-r^2\le\|L_-V\|_G^2/c_\Phi\).
\end{proof}

\begin{theorem}[Mass off the constant level and the dissipation floor]
\label{thm:c3v66-floor}
For every profile of the family and every \(\sigma\),
\begin{equation}
 \boxed{\quad
 \mathrm{Var}\ \le\ \bigl(3\,\mathrm{Var}_{\max}+5C_r^2\bigr)\,(1-\hat a_0),
 \qquad\text{hence}\qquad
 1-\hat a_0\ \ge\ \frac{\mathrm{Var}}{3\,\mathrm{Var}_{\max}+5C_r^2} .
 \quad}
 \label{eq:c3v66-mass}
\end{equation}
Consequently, on a spread configuration (\(\mathrm{Var}\ge\eta_1/16\)),
\begin{equation}
 \boxed{\quad
 \nu\|\nabla\widehat V\|_G^2\ \ge\ \frac{\eta_1}{32(3\,\mathrm{Var}_{\max}+5C_r^2)},
 \qquad
 d\ \ge\ d_{\rm spread}:=\frac{c_\Phi\,\eta_1}{8(3\,\mathrm{Var}_{\max}+5C_r^2)},
 \qquad
 \mathcal L\ \ge\ \frac{d_{\rm spread}+2c_\Phi}8 .
 \quad}
 \label{eq:c3v66-floor}
\end{equation}
\end{theorem}

\begin{proof}
Write \(\varepsilon:=1-\hat a_0=\sum_{k\ge1}\hat a_k\).  By \eqref{eq:c3v21-spectral},
\(\mathrm{Var}=(\tfrac12-r)^2\hat a_0+\sum_{k\ge1}(\mu_k-r)^2\hat a_k\).  First,
\(r-\frac12=\sum_{k\ge1}(\mu_k-\frac12)\hat a_k\le\sum_{k\ge1}\mu_k\hat a_k\le\bigl(\sum_{k\ge1}\mu_k^2\hat a_k\bigr)^{1/2}\varepsilon^{1/2}\)
by Cauchy--Schwarz, so \((\frac12-r)^2\hat a_0\le\bigl(\sum_{k\ge1}\mu_k^2\hat a_k\bigr)\varepsilon\).  Second,
\(\sum_{k\ge1}(\mu_k-r)^2\hat a_k\le2\sum_{k\ge1}\mu_k^2\hat a_k+2r^2\varepsilon\).  Third,
\(\sum_{k\ge1}\mu_k^2\hat a_k\le\sum_k\mu_k^2\hat a_k=\mathrm{Var}+r^2\le\mathrm{Var}_{\max}+C_r^2\)
(Lemma~\ref{lem:c3v66-varmax}, \(r\le C_r\)).  Hence
\(\mathrm{Var}\le(\mathrm{Var}_{\max}+C_r^2)\varepsilon+2(\mathrm{Var}_{\max}+C_r^2)\varepsilon+2C_r^2\varepsilon=(3\mathrm{Var}_{\max}+5C_r^2)\varepsilon\).
For the floor: \(\nu\|\nabla\widehat V\|_G^2=\sum_k\frac k2\hat a_k\ge\frac12\varepsilon\) (Lemma C2-V41-modes),
\(d=4\nu\|\nabla V\|_G^2=4\varphi\,\nu\|\nabla\widehat V\|_G^2\ge4c_\Phi\cdot\frac12\varepsilon\), and \(\mathcal L=\frac18(d+2\varphi)\).
\end{proof}

\begin{remark}[Reading]
\label{rem:c3v66-reading}
A spread configuration has an explicit fraction of its Gaussian
energy above the constant level and therefore an explicit Gaussian
dissipation: the Dirichlet form of a spread configuration exceeds
\(\frac14\varphi\) by an explicit amount.  The bound is crude (the constant
\(3\mathrm{Var}_{\max}+5C_r^2\) is the price of allowing the mass off the constant
level to sit at high levels, where each unit of mass carries much
variance) but it is what the one-level criterion of V21 and the
class give.  With V62, the blow-up measure charges spread
configurations or multi-group ones; on the former the dissipation
is at least \(d_{\rm spread}\).
\end{remark}

% =============================================================================
\section{Consequences for the extremal rows and the Lamb bound}
\label{sec:c3v66-consequences}
% =============================================================================

\begin{proposition}[Extremal rows on spread configurations]
\label{prop:c3v66-rows}
\begin{enumerate}[label=\textup{(\roman*)},leftmargin=3.2em]
\item If a minimizer of \(\mathcal L\) over the family is a spread
      configuration, then \(\mathcal L_{\min}\ge(d_{\rm spread}+2c_\Phi)/8\), which
      improves \eqref{eq:c3v34-Lmin} whenever \(d_{\rm spread}+2c_\Phi>8\mathcal L_-\).
\item At a spread extremal-Rayleigh configuration
      (Theorem~\ref{thm:c3v36-rayleigh}), \(\eta_1/16\le\mathrm{Var}\le\|N\|_G^2/\varphi\), so
      \(\|N(V)\|_G^2\ge c_\Phi\eta_1/16\) there: the nonlinearity of a spread
      extremal-Rayleigh configuration is bounded below explicitly.
\item In the mean, with \(A\) the nearly constant anchors of V62 and
      \(B\) the spread configurations, \(\langle d\rangle\ge d_{\rm spread}\,\mathfrak m(B)\), and
      \(\mathfrak m(B)\ge\mathfrak m(A^c)-\mathfrak m(\text{multi-group, not spread})\).
\end{enumerate}
\end{proposition}

\begin{proof}
(i), (ii) direct; (iii) \(d\ge d_{\rm spread}\) on \(B\) and \(d\ge0\).
\end{proof}

\begin{firewall}[Item (AC2): no floor on the Gaussian enstrophy from the dissipation]
\label{fw:c3v66-enstrophy}
The Lamb bound \(|\int\ell\cdot YG|\le C_\infty\mathcal E^{1/2}\|Y\|_G\) is an upper bound; a
lower bound on the Lamb pairing on a spread configuration would
need \(\mathcal E=\|\Omega\|_G^2\) bounded below.  For a divergence-free field,
\(\int|\Omega|^2G=\int|\nabla V|^2G-\int qG\) with
\(\int qG=\int V_iV_j\partial_i\partial_jG=\frac{R}{4\nu^2}-\frac{\varphi}{2\nu}\), so
\(\mathcal E=\frac{d+2\varphi}{4\nu}-\frac R{4\nu^2}\), the closure identity of C2-V21; the
radial kinetic energy \(R=\int(z\cdot V)^2G\) is bounded above only by
\(X^2\le4\nu(d+3\varphi)\), which exceeds \(\nu(d+2\varphi)\), so the dissipation floor
gives no positive floor on \(\mathcal E\).  \(\mathcal E>0\) on the family (a
divergence-free irrotational field in \(L^6\) is zero), but its floor is
a compactness constant, like those of V34 and V6, and the method of
V21--V23 does not apply to it directly (the degenerate case,
\(\Omega=0\), is excluded by the same \(L^6\) argument as the eigenfunctions,
but the quantitative distance to it is measured by \(R\), not by a
spectral quantity).  The Lamb pairing on the spread anchor is
therefore bounded above only; V67 examines the Gaussian enstrophy
identity of C2-V19 for a floor in the mean or at extremal
configurations.
\end{firewall}

% =============================================================================
\section{Integrated V66 conclusion and next acquisition order}
\label{sec:c3v66-conclusion}
% =============================================================================

\begin{theorem}[What V66 proves]
\label{thm:c3v66-conclusion}
Relative to the two legacy cycles and V1--V65: the bound of the
variance on the family (Lemma~\ref{lem:c3v66-varmax}), the mass and
dissipation floors on spread configurations
(Theorem~\ref{thm:c3v66-floor}) and their consequences
(Proposition~\ref{prop:c3v66-rows}).  The full Navier--Stokes
stability/global-regularity theorem remains unproved.
\end{theorem}

\begin{proof}
The cited results.
\end{proof}

\begin{programme}[V67 acquisition order]
\label{prog:c3v67}
\begin{enumerate}[label=\textup{(AC\arabic*)},leftmargin=4.8em]
\item The Gaussian enstrophy identity (C2-V19) at extremal \(\mathcal E\): with
      the identity's exact form re-read from the legacy file, apply
      Lemma~\ref{lem:c3v34-extremal} to \(\mathcal E\) and bound the stretching
      \(\mathcal S=\int\Omega\cdot(\Omega\cdot\nabla)VG\) by \(C_\infty'\mathcal E\) and the transport
      \(\mathcal T=\frac1{4\nu}\int|\Omega|^2(z\cdot V)G\) by the Hardy inequality for the
      divergence-free \(\Omega\); state the resulting explicit two-sided
      bound on \(\sup\mathcal E\) and \(\inf\mathcal E\), and whether the lower one is
      positive.
\item If a positive explicit \(\inf\mathcal E\) results, the Lamb pairing on
      the spread anchor in the mean; otherwise the record.
\item The next compulsory companion audit is V70.
\end{enumerate}
\end{programme}

\begin{assessment}[Position after V66]
\label{ass:c3v66-position}
Spread configurations have an explicit dissipation floor; the
Gaussian enstrophy has none from it.  No full stability or
global-regularity claim is made.
\end{assessment}

% =============================================================================
\section*{Version V66 (third cycle) append-only record}
\addcontentsline{toc}{section}{Version V66 (third cycle) append-only record}
% =============================================================================

The complete V65 source of the third cycle was retained before this
continuation.  Its SHA--256 digest was
\begin{center}\scriptsize\ttfamily
c5479bedd2536262d5e649c4c1aee10b8a3a3272bcc586cab83bfab9c4a7a834.
\end{center}
V66 proved the dissipation floor on spread configurations.  No full
regularity claim is made.

% =============================================================================
\section{V67: the extremal Gaussian enstrophy, and a constraint on the velocity and gradient constants}
\label{sec:c3v67-entry}
% =============================================================================

Items (AC1) and (AC2) of Programme~\ref{prog:c3v67} are carried out.
(AC1): the Gaussian enstrophy identity of the second cycle is
re-read; at an extremum of the Gaussian enstrophy over the family
the vorticity dissipation plus the enstrophy equals the stretching
minus the transport, and both right-hand terms are bounded by the
class through the gradient constant, the velocity constant and the
Gaussian Hardy inequality for the divergence-free vorticity, with no
pressure term, since the vorticity equation has none.  The result
is an explicit inequality between class constants that every
homogeneous type-I point must satisfy: the gradient constant plus a
quadratic in the velocity constant is at least one.  It is the one
extremal row of the cycle in which the pressure does not appear and
the evaluation closes; it is of \(\varepsilon\)-regularity type and consistent
with (E5), and it is recorded as an explicit constraint, not as an
exclusion, since the class constants are not known to violate it.
(AC2): the identity gives no explicit floor on the Gaussian
enstrophy, in the mean or at extremal configurations, so the Lamb
pairing on the spread anchor stays bounded above only.  No global
regularity claim is made.  The next compulsory companion audit is
V70.

% =============================================================================
\section{The extremal Gaussian enstrophy}
\label{sec:c3v67-extremal}
% =============================================================================

The Gaussian enstrophy identity (Theorem C2-V19-enstrophy), for every
\(\sigma\),
\begin{equation}
 \frac12\frac{\dd}{\dd\sigma}\mathcal E=-\mathcal E_\nabla-\mathcal E+\mathcal S-\mathcal T,\qquad
 \mathcal E=\int|\Omega|^2G,\ \ \mathcal E_\nabla=\nu\int|\nabla\Omega|^2G,\ \ \mathcal S=\int\Omega\cdot(\Omega\cdot\nabla)V\,G,\ \ \mathcal T=\frac1{4\nu}\int|\Omega|^2(z\cdot V)G ,
 \label{eq:c3v67-identity}
\end{equation}
follows from the vorticity equation
\(\partial_\sigma\Omega=\nu\Delta\Omega-\frac12z\cdot\nabla\Omega-\Omega-(V\cdot\nabla)\Omega+(\Omega\cdot\nabla)V\) (Lemma
C2-V19-vorticity-equation) by the \(G\)-symmetry of \(\nu\Delta-\frac12z\cdot\nabla\); it
contains no pressure.  \(\mathcal E\) is continuous on the family, \(C^1\) in
\(\sigma\), and positive at every configuration (C2-V19).

\begin{lemma}[Bounds on stretching and transport]
\label{lem:c3v67-bounds}
For every profile of the family and every \(\sigma\),
\begin{equation}
 |\mathcal S|\le C_\infty'\,\mathcal E,\qquad
 |\mathcal T|\le\Bigl(\frac{\sqrt3}2\kappa+\frac{\kappa^2}2\Bigr)\mathcal E+\frac12\mathcal E_\nabla .
 \label{eq:c3v67-bounds}
\end{equation}
\end{lemma}

\begin{proof}
\(|(\Omega\cdot\nabla)V|\le|\Omega||\nabla V|\le C_\infty'|\Omega|\).  For the transport,
\(|\mathcal T|\le\frac{C_\infty}{4\nu}\||z|\Omega\|_G\|\Omega\|_G\), and the Gaussian Hardy inequality
for the divergence-free \(\Omega\) (Lemma C2-V36-hardy, whose proof uses
only the divergence-free condition) gives
\(\||z|\Omega\|_G\le(12\nu)^{1/2}\mathcal E^{1/2}+4\nu\|\nabla\Omega\|_G\); hence
\(|\mathcal T|\le\frac{C_\infty(12\nu)^{1/2}}{4\nu}\mathcal E+C_\infty\mathcal E^{1/2}\|\nabla\Omega\|_G\le\frac{\sqrt3}2\kappa\mathcal E+\frac\nu2\|\nabla\Omega\|_G^2+\frac{C_\infty^2}{2\nu}\mathcal E\),
with \(C_\infty(12\nu)^{1/2}/(4\nu)=\frac{\sqrt3}2\kappa\) and \(C_\infty^2/(2\nu)=\kappa^2/2\).
\end{proof}

\begin{theorem}[Extremal Gaussian enstrophy, and the constraint]
\label{thm:c3v67-extremal}
Let \(W_{\mathcal E}\) be a maximizer or minimizer of \(\mathcal E\) over \(\mathcal F(x_0)\).  At
\(W_{\mathcal E}\), \(\sigma=0\),
\begin{equation}
 \mathcal E_\nabla+\mathcal E=\mathcal S-\mathcal T,\qquad\text{hence}\qquad
 \frac12\mathcal E_\nabla+\Bigl(1-C_\infty'-\frac{\sqrt3}2\kappa-\frac{\kappa^2}2\Bigr)\mathcal E\ \le\ 0 .
 \label{eq:c3v67-row}
\end{equation}
Since \(\mathcal E>0\) and \(\mathcal E_\nabla\ge0\) there, every homogeneous type-I point
satisfies
\begin{equation}
 \boxed{\quad
 C_\infty'+\frac{\sqrt3}2\kappa+\frac{\kappa^2}2\ \ge\ 1 .
 \quad}
 \label{eq:c3v67-constraint}
\end{equation}
\end{theorem}

\begin{proof}
Lemma~\ref{lem:c3v34-extremal} with \(F=\mathcal E\) and \eqref{eq:c3v67-identity};
Lemma~\ref{lem:c3v67-bounds}; if the coefficient of \(\mathcal E\) in
\eqref{eq:c3v67-row} were positive, the left side would be positive.
\end{proof}

\begin{assessment}[Item (AC1): what the constraint is]
\label{ass:c3v67-constraint}
\eqref{eq:c3v67-constraint} says that at a homogeneous type-I point
the rescaled solution cannot have both a small velocity bound and a
small gradient bound: the vorticity's self-similar damping (the
\(-\mathcal E\) of the identity, one unit) must be balanced at the extremal
enstrophy by stretching and transport, which the class bounds by
the gradient constant and the velocity constant.  This is a
statement of \(\varepsilon\)-regularity type: the class (E5) already bounds the
velocity constant below by \(\varepsilon_{\rm reg}\) at a singular point (Firewall
C2-V36-AC2), and the parabolic derivative bounds give \(C_\infty'\) of the
order of \(\kappa\) for small \(\kappa\), so \eqref{eq:c3v67-constraint} is a
second, independent lower bound on the velocity constant, explicit
once the constant of the derivative bound is; it is consistent with
(E5) and is not an exclusion, the class constants not being known
to violate it.  Its interest for the series is structural: it is
the first extremal row whose explicit evaluation closes, and it
closes because the vorticity equation has no pressure; every row
with a pressure (V38, V47) stopped at a supremum bound.  The same
mechanism gives, at the extremal enstrophy, the explicit two-sided
bound \(\mathcal E_\nabla\le2(C_\infty'+\frac{\sqrt3}2\kappa+\frac{\kappa^2}2-1)\mathcal E\) on the vorticity
dissipation there.
\end{assessment}

\begin{remark}[Against the defect row of V38]
\label{rem:c3v67-against-V38}
Proposition~\ref{prop:c3v38-advective} gave, at the extremal defect,
\(\frac\nu2\|\nabla Y\|^2+(\frac12-C_\infty'-\frac{\sqrt3}2\kappa-\frac{\kappa^2}2)\|Y\|^2\le|\langle Y,\nabla P'\rangle_G|\),
the same structure with the constant \(\frac12\) in place of \(1\) and the
pressure pairing on the right; V47 bounded the pairing at supremum
level and no constraint followed.  For the vorticity the pairing is
absent, and the constraint \eqref{eq:c3v67-constraint} follows.  The
difference in the constants (\(1\) against \(\frac12\)) is the difference
between the scaling damping of the vorticity (\(-\Omega\)) and that of
the defect (\(-\frac12Y\)).
\end{remark}

% =============================================================================
\section{No explicit floor on the Gaussian enstrophy}
\label{sec:c3v67-floor}
% =============================================================================

\begin{firewall}[Item (AC2): the Lamb pairing stays bounded above only]
\label{fw:c3v67-no-floor}
At an extremum of \(\mathcal E\), \eqref{eq:c3v67-row} is an upper bound on
\(\mathcal E_\nabla\) in terms of \(\mathcal E\), not a lower bound on \(\mathcal E\); in the mean,
the mean form of Theorem C2-V19-enstrophy,
\(\langle\mathcal E_\nabla+\mathcal E\rangle=\langle\mathcal S-\mathcal T\rangle\), bounds \(\langle\mathcal E\rangle\) above by the same
constants times itself and gives no floor.  The positivity of
\(\inf_{\mathcal F(x_0)}\mathcal E\) remains a compactness fact (C2-V19), and the Lamb
pairing on the spread anchor is bounded above by
\(\kappa(2\mathcal L)^{1/2}\|Y\|_G\) only.  The dissipation floor of V66 does not
transfer to \(\mathcal E\) (Firewall~\ref{fw:c3v66-enstrophy}).  What would give
a floor is a quantitative form of the exclusion of irrotational
divergence-free fields, i.e.\ of harmonic gradients, in the class:
the analogue of V21--V23 with the distance to the space of harmonic
gradients in place of the distance to the eigenspaces; the
degenerate case is again excluded by the \(L^6\) bound, and the same
growth-versus-Gaussian-tail mechanism applies, so an exponentially
small explicit floor on \(\mathcal E\) is available by that method.  It is
not pursued: an exponentially small floor on the Lamb pairing
decides nothing, as V24 showed for the variance.
\end{firewall}

% =============================================================================
\section{Integrated V67 conclusion and next acquisition order}
\label{sec:c3v67-conclusion}
% =============================================================================

\begin{theorem}[What V67 proves]
\label{thm:c3v67-conclusion}
Relative to the two legacy cycles and V1--V66: the bounds on
stretching and transport (Lemma~\ref{lem:c3v67-bounds}), the extremal
enstrophy row and the constraint \eqref{eq:c3v67-constraint}
(Theorem~\ref{thm:c3v67-extremal}), and the record that the Gaussian
enstrophy has no explicit floor.  The full Navier--Stokes
stability/global-regularity theorem remains unproved.
\end{theorem}

\begin{proof}
The cited results.
\end{proof}

\begin{programme}[V68 acquisition order]
\label{prog:c3v68}
\begin{enumerate}[label=\textup{(AC\arabic*)},leftmargin=4.8em]
\item The gate on the spread anchor: on the trapped sub-windows of
      \(\mathcal E_0\) with \(\mathrm{Var}\ge\eta_1/16\), the cross term is
      \(\mathcal A_{U_0}+\mathcal B_{U_0}+O(\varepsilon_{\rm th}^{1/4})+O(\text{tail})\) with
      \(|\mathcal A_{U_0}|\le\kappa(2\mathcal L)^{1/2}\|Y\|_G\) and \(|\mathcal B_{U_0}|\le\mathcal M_{U_0}+C_{11}\varepsilon_{\rm th}\|Y\|_G\);
      combine with the dissipation floor (\(\mathcal L\ge(d_{\rm spread}+2c_\Phi)/8\)) and
      the extremal rows of V43 and V58 to state the gate on the
      spread anchor in explicit form, and record which of its pieces
      can be of order one.
\item The test configurations of the second cycle (the columnar
      swirl and the plane wave, C1-V88 as cited in V13) evaluated on
      an anchor: their radial-defect masses, Lamb pairing and
      pressure work, as models of the sign structure of V46 and V57.
\item The next compulsory companion audit is V70.
\end{enumerate}
\end{programme}

\begin{assessment}[Position after V67]
\label{ass:c3v67-position}
The extremal enstrophy row closes to an explicit constraint on the
class constants; the Gaussian enstrophy has no explicit floor.  No
full stability or global-regularity claim is made.
\end{assessment}

% =============================================================================
\section*{Version V67 (third cycle) append-only record}
\addcontentsline{toc}{section}{Version V67 (third cycle) append-only record}
% =============================================================================

The complete V66 source of the third cycle was retained before this
continuation.  Its SHA--256 digest was
\begin{center}\scriptsize\ttfamily
fb5cf1c832fce246743fb0d638f58ac6eb938e73d26bd2bd73a2788057ae6173.
\end{center}
V67 proved the extremal enstrophy row and its constraint.  No full
regularity claim is made.

% =============================================================================
\section{V68: the gate on the spread anchor, and the gate-neutral test configurations}
\label{sec:c3v68-entry}
% =============================================================================

Items (AC1) and (AC2) of Programme~\ref{prog:c3v68} are carried out.
(AC1): on a spread anchor the two pieces of the localized cross
term, the Lamb pairing and the pressure work, are each bounded by
explicit quantities of order one, the Dirichlet form is bounded
below beyond a quarter of the energy by the dissipation floor, and
the gate is an inequality between these; no piece has a sign, and
the dissipation floor bounds the Lamb piece's upper bound below
without bounding the piece itself.  (AC2): the two exact solutions
of the second cycle's checks, the columnar swirl and the decaying
plane wave, have vanishing nonlinearity, hence vanishing cross term:
they are gate-neutral, realizing neither the demanded
anti-alignment nor its opposite, and they show that the gate's
content lies entirely in configurations whose nonlinearity couples
to the defect, which neither exact solution does.  No global
regularity claim is made.  The next compulsory companion audit is
V70.

% =============================================================================
\section{The gate on the spread anchor}
\label{sec:c3v68-gate}
% =============================================================================

\begin{theorem}[The gate on the spread anchor, explicit form]
\label{thm:c3v68-gate}
On the trapped sub-windows of \(\mathcal E_0\), for a configuration with
\(\mathrm{Var}\ge\eta_1/16\):
\begin{enumerate}[label=\textup{(\roman*)},leftmargin=3.2em]
\item \(\langle N,Y\rangle_G=\mathcal A_{U_0}+\mathcal B_{U_0}+E\) with \(|E|\le\mathcal S'\) (of
      \eqref{eq:c3v62-Sprime}), \(|\mathcal A_{U_0}|\le\kappa(2\mathcal L)^{1/2}\|Y\|_G\), and
      \(|\mathcal B_{U_0}|\le\mathcal M_{U_0}+C_{11}\varepsilon_{\rm th}\|Y\|_G\) with \(\mathcal M_{U_0}\) of
      \eqref{eq:c3v58-bound}; and \(\mathcal L\ge(d_{\rm spread}+2c_\Phi)/8\)
      (Theorem~\ref{thm:c3v66-floor}).
\item If a maximizer or minimizer of \(\mathcal L\) is a spread anchor, then at
      it \(\|Y\|_G^2\le\kappa(2\mathcal L)^{1/2}\|Y\|_G+\mathcal M_{U_0}+C_{11}\varepsilon_{\rm th}\|Y\|_G+\mathcal S'\), and the
      anti-alignment \(-\|Y\|_G^2\) is realized by \(\mathcal A_{U_0}+\mathcal B_{U_0}\), each
      of which may be of order one.
\item In the mean over the spread anchors \(B_0:=A^c\cap\mathcal E_0\cap\{\text{trapped}\}\),
      \(\int_{B_0}\langle N,Y\rangle_G\dd\mathfrak m\ge-\int_{B_0}\bigl[\kappa(2\mathcal L)^{1/2}\|Y\|_G+\mathcal M_{U_0}\bigr]\dd\mathfrak m-(C_{11}\varepsilon_{\rm th}C_Y+\mathcal S')\mathfrak m(B_0)\).
\end{enumerate}
\end{theorem}

\begin{proof}
(i) Theorems~\ref{thm:c3v11-far}, \ref{thm:c3v16-localized},
Proposition~\ref{prop:c3v41-lamb}, Theorem~\ref{thm:c3v58-group},
Theorem~\ref{thm:c3v66-floor}.  (ii) Row (1) of the extremal ledger.
(iii) Integration of (i).
\end{proof}

\begin{assessment}[Item (AC1): what is of order one]
\label{ass:c3v68-order-one}
On the spread anchor nothing is small: the Lamb pairing's bound
\(\kappa(2\mathcal L)^{1/2}\|Y\|_G\) has \(\mathcal L\) bounded below beyond \(\frac14c_\Phi\) by the
dissipation floor, and the pressure work's bound \(\mathcal M_{U_0}\) has the
halo's \(L^4\) energy, bounded below by the anchor's energetic cell
(energy at least \(c'\), V16) through \(\|V\|_{L^4(U'')}^4\ge\|V\|_{L^2(Q)}^4/|Q|\) on
that cell.  Both bounds are upper bounds on absolute values; the
pieces themselves have the sign structure of V42 (Lamb: the defect
against the Lamb vector, normal to the velocity--vorticity plane) and
V46 (pressure: the velocity tensor against the multipole Hessians
of the radial-defect cells).  The gate on the spread anchor is thus
a statement about two orientations of order-one quantities, and the
class fixes neither; the spread anchor is where the anti-alignment
must live (V62), and it is where the series has no sign.  This is
the honest end of the anchor programme at the level of the class
bounds: the remaining content of G1 on the anchor-alone event is the
orientation of the defect relative to the Lamb vector and of the
velocity tensor relative to the radial-defect Hessians, on a
configuration of at most \(N_C(2r_2+1)^3\) cells with explicit energy,
dissipation and defect bounds.
\end{assessment}

% =============================================================================
\section{The gate-neutral test configurations}
\label{sec:c3v68-tests}
% =============================================================================

\begin{proposition}[The columnar swirl and the plane wave are gate-neutral]
\label{prop:c3v68-tests}
\begin{enumerate}[label=\textup{(\roman*)},leftmargin=3.2em]
\item For the columnar swirl \(V=w(\sigma,r)e_\theta\) (Proposition C2-V18-swirl,
      not in the class but an exact solution), \(N(V)\equiv0\), \(z\cdot V\equiv0\),
      \(Y=\partial_\sigma V\) is azimuthal, the Lamb vector \(\Omega\times V\) is poloidal
      and orthogonal to \(Y\), and the radial defect \(z\cdot Y\) vanishes:
      the Lamb pairing, the Bernoulli work and the cross term all
      vanish identically.
\item For the decaying plane wave \(V=a\cos(k\cdot z)\,b\,e^{-\nu|k|^2t(\sigma)}\) with
      \(b\perp k\), \((V\cdot\nabla)V=0\), \(P_V=0\), \(N(V)\equiv0\), the Lamb vector is the
      gradient \(-\nabla\frac12|V|^2\), and the Lamb pairing equals minus the
      kinetic Bernoulli work: the cross term vanishes identically.
\end{enumerate}
In both cases \(\langle N,Y\rangle_G=0\) and \(\|Y\|_G^2=-\frac{\dd}{\dd\sigma}\mathcal L\): the Dirichlet
form decays exactly at the rate of the squared defect, the mean
identity \(\langle\langle N,Y\rangle\rangle=-\langle\|Y\|^2\rangle\) cannot hold on an invariant
measure unless the defect vanishes, and neither flow is a rescaling
limit at a homogeneous point (consistently with the second cycle:
the swirl has infinite Dirichlet integral and the plane wave is
not in \(L^6\)).
\end{proposition}

\begin{proof}
(i) Proposition C2-V18-swirl: \((V\cdot\nabla)V=-(w^2/r)e_r\) is a gradient
cancelled by the pressure, so \(N\equiv0\); \(e_\theta\perp z\); \(\Omega=-\partial_zw\,e_r+r^{-1}\partial_r(rw)e_z\)
and \(\Omega\times e_\theta\) has no azimuthal component.  (ii) A plane wave
with \(b\perp k\) has \((V\cdot\nabla)V=(b\cdot k)(\dots)=0\) and zero pressure; then
\(\Omega\times V=(V\cdot\nabla)V-\nabla\frac12|V|^2=-\nabla\frac12|V|^2\), and
\(\langle-\nabla\frac12|V|^2,Y\rangle_G=-\frac1{2\nu}\int\frac12|V|^2(z\cdot Y)G\) by the
integration by parts of Theorem C2-V3-cross.  The last statements:
\(\frac{\dd}{\dd\sigma}\mathcal L=-\|Y\|^2-\langle N,Y\rangle\) with \(N=0\).
\end{proof}

\begin{assessment}[Item (AC2): what the test configurations show]
\label{ass:c3v68-tests}
The second cycle used these flows to check the vanishing of the
nonlinear terms in its identities; for the gate they show something
more specific.  A configuration with vanishing nonlinearity has a
Dirichlet form that decays at the full rate of its defect and can
never be stationary in the statistical sense; the gate's demand, a
cross term equal to minus the squared defect in the mean, is met
only by configurations in which the nonlinearity opposes the
defect.  The sign structures of V42 and V46 name the two ways: the
defect normal to the velocity--vorticity plane against the Lamb
vector, and the velocity tensor oriented against the radial-defect
Hessians.  Neither test flow has either coupling.  A test flow with
the couplings would be a nontrivial configuration of the type the
gate must exclude or admit, and the second cycle's example set does
not contain one; constructing one, in the class or in a model class,
is the natural check of the sign structures, and is recorded as a
candidate for a later block.
\end{assessment}

% =============================================================================
\section{Integrated V68 conclusion and next acquisition order}
\label{sec:c3v68-conclusion}
% =============================================================================

\begin{theorem}[What V68 proves]
\label{thm:c3v68-conclusion}
Relative to the two legacy cycles and V1--V67: the gate on the
spread anchor in explicit form (Theorem~\ref{thm:c3v68-gate}) and the
gate-neutrality of the test configurations
(Proposition~\ref{prop:c3v68-tests}).  The full Navier--Stokes
stability/global-regularity theorem remains unproved.
\end{theorem}

\begin{proof}
The cited results.
\end{proof}

\begin{programme}[V69 acquisition order]
\label{prog:c3v69}
\begin{enumerate}[label=\textup{(AC\arabic*)},leftmargin=4.8em]
\item Consolidate V66--V68 as the spread-anchor block: the
      dissipation floor, the extremal enstrophy row with its
      constraint, the gate on the spread anchor and the
      gate-neutral test configurations.
\item Fix the direction for V71--V75: the vorticity ledger.  Every
      functional of the vorticity alone has a pressure-free identity
      (the vorticity equation has no pressure), so its extremal row
      closes to an explicit constraint on class constants, as the
      Gaussian enstrophy did (V67); the candidates are the weighted
      enstrophies \(\int|z|^{2m}|\Omega|^2G\), the palinstrophy \(\int|\nabla\Omega|^2G\),
      and the helicity's vorticity part.  Determine the constraints
      and whether, combined with the structure of (E5) (derivative
      constants of order \(\kappa\) for small \(\kappa\)), they give an explicit
      \(\varepsilon\)-regularity threshold for homogeneous type-I points in terms
      of the constants of (E5).
\item V70 compulsory companion audit and delivery.
\end{enumerate}
\end{programme}

\begin{assessment}[Position after V68]
\label{ass:c3v68-position}
The gate on the spread anchor is explicit with no small piece; the
classical test flows are gate-neutral.  No full stability or
global-regularity claim is made.
\end{assessment}

% =============================================================================
\section*{Version V68 (third cycle) append-only record}
\addcontentsline{toc}{section}{Version V68 (third cycle) append-only record}
% =============================================================================

The complete V67 source of the third cycle was retained before this
continuation.  Its SHA--256 digest was
\begin{center}\scriptsize\ttfamily
fba554f85d8269fd36c190edfea5a7782718c87d5211eacba9968ef2d9a2eb91.
\end{center}
V68 stated the gate on the spread anchor and the gate-neutrality of
the test configurations.  No full regularity claim is made.

% =============================================================================
\section{V69: the spread-anchor block consolidated, and the vorticity ledger ordered}
\label{sec:c3v69-entry}
% =============================================================================

Items (AC1) and (AC2) of Programme~\ref{prog:c3v69} are carried out.
The spread-anchor block V66--V68 is consolidated: the dissipation
floor on spread configurations, the extremal Gaussian enstrophy row
closing to an explicit constraint on the class constants, the gate
on the spread anchor with no small piece, and the gate-neutrality
of the classical test flows.  For the direction, the observation of
V67 is made systematic: every functional of the vorticity alone has
a pressure-free identity, so its extremal row closes; the block
V71--V75 writes the vorticity ledger at extremal configurations and
determines what the resulting constraints say jointly.  No global
regularity claim is made.  The next compulsory companion audit is
V70.

% =============================================================================
\section{The spread-anchor block}
\label{sec:c3v69-block}
% =============================================================================

\begin{theorem}[The spread-anchor block, V66--V68]
\label{thm:c3v69-block}
At a homogeneous type-I point:
\begin{enumerate}[label=\textup{(\arabic*)},leftmargin=3.2em]
\item (Dissipation floor, V66.)  \(\mathrm{Var}\le(3\mathrm{Var}_{\max}+5C_r^2)(1-\hat a_0)\) with
      \(\mathrm{Var}_{\max}=(C_Y+C_N^{1/2})^2/c_\Phi\); on a spread configuration
      (\(\mathrm{Var}\ge\eta_1/16\)), \(d\ge d_{\rm spread}=c_\Phi\eta_1/(8(3\mathrm{Var}_{\max}+5C_r^2))\) and
      \(\mathcal L\ge(d_{\rm spread}+2c_\Phi)/8\); the Gaussian enstrophy has no floor
      from it.
\item (Extremal enstrophy, V67.)  At an extremum of \(\mathcal E\),
      \(\mathcal E_\nabla+\mathcal E=\mathcal S-\mathcal T\) with \(|\mathcal S|\le C_\infty'\mathcal E\) and
      \(|\mathcal T|\le(\frac{\sqrt3}2\kappa+\frac{\kappa^2}2)\mathcal E+\frac12\mathcal E_\nabla\); hence every
      homogeneous type-I point satisfies \(C_\infty'+\frac{\sqrt3}2\kappa+\frac{\kappa^2}2\ge1\),
      the first extremal row of the cycle whose explicit evaluation
      closes, because the vorticity equation has no pressure.
\item (Gate on the spread anchor, V68.)  The Lamb pairing and the
      pressure work on the anchor are bounded by explicit order-one
      quantities and have no sign; the classical test flows (columnar
      swirl, plane wave) have vanishing nonlinearity and are
      gate-neutral.
\end{enumerate}
\end{theorem}

\begin{proof}
Theorems~\ref{thm:c3v66-floor}, \ref{thm:c3v67-extremal}, \ref{thm:c3v68-gate};
Proposition~\ref{prop:c3v68-tests}; Firewalls~\ref{fw:c3v66-enstrophy},
\ref{fw:c3v67-no-floor}.
\end{proof}

\begin{assessment}[What the block settles]
\label{ass:c3v69-settles}
The anchor programme (V56--V68) has reached its natural end at the
level of the class bounds.  On the anchor-alone event the gate is an
inequality between explicit local quantities; the nearly constant
anchor is excluded from carrying the anti-alignment (V62); on the
spread anchor every quantity is explicit and of order one, and the
gate is a question about two orientations that the class does not
fix (V68).  What the block added beyond the anchor programme is the
constraint \eqref{eq:c3v67-constraint}: an inequality between the
class constants that every homogeneous type-I point must satisfy,
obtained from an identity without pressure at an extremal
configuration.  It is of \(\varepsilon\)-regularity type and does not exclude
anything the series knows, but it is the first time the extremal
method produced a relation between explicit constants alone, with
no pinching or defect rate on the non-explicit side; the reason is
that the observable, the Gaussian enstrophy, is positive on the
family by an argument that does not pass through compactness (an
irrotational divergence-free field in \(L^6\) is zero), and its
identity has no pressure.  Whether the vorticity ledger has more
such rows is the direction.
\end{assessment}

% =============================================================================
\section{Direction for V71--V75: the vorticity ledger}
\label{sec:c3v69-direction}
% =============================================================================

\begin{assessment}[The vorticity ledger at extremal configurations]
\label{ass:c3v69-direction}
For any smooth functional \(F(\Omega)\) of the vorticity alone, the
identity for \(\frac{\dd}{\dd\sigma}F\) follows from the vorticity equation and
contains no pressure; at an extremum of \(F\) over the family the
identity holds with zero left side and every term is bounded by
the class through derivative constants and the Gaussian Hardy
inequality.  V71: the weighted enstrophies \(\mathcal E_m:=\int|z|^{2m}|\Omega|^2G\)
(\(m=1,2\)): their identities (the weight \(|z|^{2m}\) commutes with the
scaling term up to \(m\) times the unweighted quantity and introduces
the moment \(\int|z|^{2m-2}z\cdot\Omega\,\Omega\cdot\nabla V\)-type couplings), the bounds on
stretching and transport with the weight, and the extremal rows;
record whether each gives a constraint of the form
\(a_mC_\infty'+b_m\kappa+c_m\kappa^2\ge m+1\) (the scaling damping of \(|z|^{2m}|\Omega|^2G\)
being \(m+1\) units).  V72: the palinstrophy \(\mathcal P:=\int|\nabla\Omega|^2G\): its
identity has damping \(\frac32\) units and stretching terms with
\(\nabla^2V\), bounded by \(C_\infty''\); the extremal row gives a constraint
involving \(C_\infty''\).  V73: the joint reading: the constraints
\(C_\infty'+\frac{\sqrt3}2\kappa+\frac{\kappa^2}2\ge1\), \(\dots\ge m+1\), \(\dots\ge\frac32\) with the
structure of (E5) (each \(C_\infty^{(k)}\) is at most \(c_k\kappa\) for small \(\kappa\), by
the parabolic derivative bounds for small solutions) give
\(\kappa\ge\kappa_k\) with \(\kappa_k\) explicit in \(c_k\); determine the strongest and
compare with \(\varepsilon_{\rm reg}\) of (E5) and with the constant-chasing
constraints of C2-V36--V38; record whether any constraint is not
implied by (E5) alone.  V74 consolidates and fixes the direction for
V76--V80; V75 audits.  No global regularity claim is made.
\end{assessment}

% =============================================================================
\section{Integrated V69 conclusion and next acquisition order}
\label{sec:c3v69-conclusion}
% =============================================================================

\begin{theorem}[What V69 establishes]
\label{thm:c3v69-conclusion}
Relative to the two legacy cycles and V1--V68: the consolidated
spread-anchor block (Theorem~\ref{thm:c3v69-block}) and the direction.
The full Navier--Stokes stability/global-regularity theorem remains
unproved.
\end{theorem}

\begin{proof}
The cited results.
\end{proof}

\begin{programme}[V70 acquisition order]
\label{prog:c3v70}
\begin{enumerate}[label=\textup{(AC\arabic*)},leftmargin=4.8em]
\item The compulsory companion audit: read the current SM and YM
      notes and record their digests and decision.
\item Record the position after seventy versions; delivery.
\end{enumerate}
\end{programme}

\begin{assessment}[Position after V69]
\label{ass:c3v69-position}
The spread-anchor block is consolidated; the direction is the
vorticity ledger at extremal configurations.  No full stability or
global-regularity claim is made.
\end{assessment}

% =============================================================================
\section*{Version V69 (third cycle) append-only record}
\addcontentsline{toc}{section}{Version V69 (third cycle) append-only record}
% =============================================================================

The complete V68 source of the third cycle was retained before this
continuation.  Its SHA--256 digest was
\begin{center}\scriptsize\ttfamily
ba094ef2ec4914c43f73beafa9ed94736ae6d9f8e2ed2321fbed8a66113c4c36.
\end{center}
V69 consolidated the spread-anchor block and ordered the vorticity
ledger.  No full regularity claim is made.

% =============================================================================
\section{V70: compulsory companion audit and the position after seventy versions}
\label{sec:c3v70-entry}
% =============================================================================

This is the seventieth version of the third cycle.  The compulsory
review of the two companion programmes is performed first, and the
position after seventy versions is recorded with the plan for the
block V71--V75.  No global regularity claim is made.  The next
compulsory companion audit is V75.

% =============================================================================
\section{Compulsory V70 review of the current companion notes}
\label{sec:c3v70-companion-audit}
% =============================================================================

The current companion files in \texttt{latex\_notes} were inspected
read-only.  Both programmes have moved since the V65 audit: the SM
programme from V39 to V41 of its seventh series, the YM programme
from V16 to V19 of its sixth series.  The frozen checkpoints are
\begin{align}
 &\texttt{Renewal\_SM\_Einstein\_Continuum\_Research\_Notes\_V41.tex},
 \notag\\[-2pt]
 &\qquad
 \texttt{c54883e08fdc32e403ebc2dd79785694c4844ba55ee129d7473f50f852313e25},
 \label{eq:c3v70-SM-checkpoint}\\
 &\texttt{Renewal\_Yang\_Mills\_Mass\_Gap\_Research\_Notes\_V19.tex},
 \notag\\[-2pt]
 &\qquad
 \texttt{b74dab144e3d275c230c4de1eeb2cd96d903c948895f08882e97b921d286ac37}.
 \label{eq:c3v70-YM-checkpoint}
\end{align}
(SM V41 has 16587 lines; YM V19 has 8208 lines.  At the first
inspection for this audit the SM file was V40 and the folder also
held YM V18, a distinct earlier state; both had been superseded by
the final inspection, whose digests are recorded.)

\emph{SM V40--V41.}  The SM programme performed its own companion
audit and set up a face-block ledger (normalized affine-simplex
data, tangential Fourier and normal-layer blocks, uniform row-sum
constants for its condition PFC, exact symmetry before absolute
bounds), then an incidence extraction and assembly validation:
global combinatorial records, geometry reconstructed from one
record, incidence validation rows, operator validation identities,
a machine-independent certificate manifest, and the remaining
concrete input.

\emph{YM V17--V19.}  The YM programme developed maximal-correlation
shell transport (one boundary channel, nested shell transport, a
spectral formulation and its circularity boundary) and then a
noncircular Dobrushin row: one-link conditional influence,
volume-uniform covariance decay, and the consequences with their
exact scope.

\begin{theorem}[V70 companion-audit decision]
\label{thm:c3v70-companion-decision}
Nothing is imported from SM V41 or YM V19.  Neither bears on the
third cycle's chapters or on the gates.
\end{theorem}

\begin{proof}
The SM results concern block ledgers and validation records of
lattice Dirac faces; the YM
results concern Dobrushin conditions and covariance decay in a
lattice gauge theory.  None is a Navier--Stokes statement, and none
concerns the ledger, the extremal method, the energy--Hessian form,
the anchor blocks or the vorticity identities.
\end{proof}

% =============================================================================
\section{Position after seventy versions, and the block V71--V75}
\label{sec:c3v70-position}
% =============================================================================

\begin{assessment}[Position after seventy versions of the third cycle]
\label{ass:c3v70-position-seventy}
Twelve chapters are written: cluster energy and the band (V2--V10),
localization (V11--V14), the anchor (V16--V19), the explicit floor
(V21--V24), the variance chapter (V26--V29), the defect block
(V31--V34), the extremal ledger (V34--V39), the angle chapter
(V41--V44), the correlation block (V46--V49), the anchor-gate block
(V56--V59), the nearly-constant-anchor block (V61--V64) and the
spread-anchor block (V66--V69), with the midpoint internal audit
(V51--V55).  The anchor programme is complete at the level of the
class bounds: the gate on the anchor-alone event is explicit
(V58), nearly constant anchors cannot carry the anti-alignment
(V62), and on spread anchors the gate is a question of two
orientations of order-one explicit quantities (V68).  The cycle's
one relation between explicit constants alone is the constraint
\(C_\infty'+\frac{\sqrt3}2\kappa+\frac{\kappa^2}2\ge1\) at every homogeneous type-I point
(V67), from the pressure-free enstrophy identity at an extremal
configuration.  The block V71--V75 follows
Assessment~\ref{ass:c3v69-direction}: the vorticity ledger at
extremal configurations, weighted enstrophies (V71), palinstrophy
(V72), the joint reading against (E5) (V73), consolidation (V74),
audit (V75).  Thirty versions remain in the cycle; the closing block
V96--V100 is reserved for the cycle's manifest, constants sheet,
negative records and the draft of the fourth cycle's V1, as the
second cycle did at V89--V100.  No global regularity claim is made.
\end{assessment}

% =============================================================================
\section{Integrated V70 conclusion and next acquisition order}
\label{sec:c3v70-conclusion}
% =============================================================================

\begin{theorem}[What V70 establishes]
\label{thm:c3v70-conclusion}
Relative to the two legacy cycles and V1--V69: the companion-audit
decision (Theorem~\ref{thm:c3v70-companion-decision}) and the position.
The full Navier--Stokes stability/global-regularity theorem remains
unproved.
\end{theorem}

\begin{proof}
The cited records.
\end{proof}

\begin{programme}[V71 acquisition order]
\label{prog:c3v71}
\begin{enumerate}[label=\textup{(AC\arabic*)},leftmargin=4.8em]
\item Derive the identity for \(\mathcal E_1:=\int|z|^2|\Omega|^2G\) from the vorticity
      equation: the scaling term gives \(-2\mathcal E_1\) (the damping of \(\Omega\)
      and the weight's degree), the diffusion gives
      \(-\nu\int|z|^2|\nabla\Omega|^2G\) plus a commutator with \(|z|^2\) (a multiple of
      \(\int\Omega\cdot(z\cdot\nabla)\Omega\,G\) and \(\mathcal E\)), the stretching \(\int|z|^2\Omega\cdot(\Omega\cdot\nabla)VG\)
      and the transport \(\frac1{4\nu}\int|z|^2|\Omega|^2(z\cdot V)G\) minus the term
      from the weight's gradient, \(\frac12\int|\Omega|^2(z\cdot V)\cdot2\,G\)-type; write
      it exactly and verify the signs on the columnar swirl.
\item The extremal row for \(\mathcal E_1\) and its constraint.
\item The next compulsory companion audit is V75.
\end{enumerate}
\end{programme}

\begin{assessment}[Position after V70]
\label{ass:c3v70-position}
The companion audit is recorded; the vorticity-ledger block begins.
No full stability or global-regularity claim is made.
\end{assessment}

% =============================================================================
\section*{Version V70 (third cycle) append-only record}
\addcontentsline{toc}{section}{Version V70 (third cycle) append-only record}
% =============================================================================

The complete V69 source of the third cycle was retained before this
continuation.  Its SHA--256 digest was
\begin{center}\scriptsize\ttfamily
4465cd41c7858874ef19473d6086eb5a5c4b0eed1edffc02c719810aedcc2a53.
\end{center}
V70 performed the compulsory companion audit and recorded the
position.  No full regularity claim is made.

% =============================================================================
\section{V71: the weighted Gaussian enstrophy identity and its extremal row}
\label{sec:c3v71-entry}
% =============================================================================

Items (AC1) and (AC2) of Programme~\ref{prog:c3v71} are carried out,
the first with a corrected damping constant.  The identity for the
second Gaussian moment of the vorticity is derived exactly from the
vorticity equation: the diffusion and scaling terms, commuted with
the weight, give a damping of three halves (not two), a positive
commutator equal to three times the viscosity times the unweighted
enstrophy, and cancelling fourth-moment terms; the transport
produces the weighted transport minus four times the viscosity
times the unweighted transport, and the stretching is weighted.
At an extremum of the weighted enstrophy the row closes, but not to
a constants-only constraint: the positive commutator carries the
unweighted enstrophy to the right side, and the row bounds the
second moment of the vorticity by an explicit multiple of the
enstrophy and the vorticity dissipation whenever the derivative and
velocity constants are below the damping.  The constants-only
constraint of V67 is thus special to the lowest row, whose
commutator vanishes; the higher rows give explicit concentration
bounds on the vorticity's Gaussian moments.  No global regularity
claim is made.  The next compulsory companion audit is V75.

% =============================================================================
\section{The identity}
\label{sec:c3v71-identity}
% =============================================================================

Put, for the vorticity \(\Omega\) of the profile,
\begin{equation}
 \mathcal E_1:=\int|z|^2|\Omega|^2G,\quad
 \mathcal E_{\nabla,1}:=\nu\int|z|^2|\nabla\Omega|^2G,\quad
 \mathcal S_1:=\int|z|^2\,\Omega\cdot(\Omega\cdot\nabla)V\,G,\quad
 \mathcal T_1:=\frac1{4\nu}\int|z|^2|\Omega|^2(z\cdot V)\,G ,
 \label{eq:c3v71-notation}
\end{equation}
beside \(\mathcal E,\mathcal E_\nabla,\mathcal S,\mathcal T\) of V67.

\begin{theorem}[Weighted Gaussian enstrophy identity]
\label{thm:c3v71-identity}
For every profile of the family and every \(\sigma\),
\begin{equation}
 \boxed{\quad
 \frac12\frac{\dd}{\dd\sigma}\mathcal E_1=-\mathcal E_{\nabla,1}-\frac32\mathcal E_1+3\nu\,\mathcal E+\mathcal S_1-\mathcal T_1+4\nu\,\mathcal T .
 \quad}
 \label{eq:c3v71-identity}
\end{equation}
\end{theorem}

\begin{proof}
Multiply the vorticity equation
\(\partial_\sigma\Omega=\nu\Delta\Omega-\frac12z\cdot\nabla\Omega-\Omega-(V\cdot\nabla)\Omega+(\Omega\cdot\nabla)V\) by \(|z|^2\Omega G\) and
integrate; all integrations by parts are justified by the class
bounds and the Gaussian decay.  Diffusion:
\(\nu\int|z|^2\Omega\cdot\Delta\Omega\,G=-\nu\int|z|^2|\nabla\Omega|^2G-\frac\nu2\int\nabla|\Omega|^2\cdot\nabla(|z|^2G)\),
and \(\nabla(|z|^2G)=(2-\frac{|z|^2}{2\nu})z\,G=:wG\); integrating by parts again,
\(-\frac\nu2\int\nabla|\Omega|^2\cdot wG=\frac\nu2\int|\Omega|^2(\operatorname{div}w-\frac{w\cdot z}{2\nu})G\) with
\(\operatorname{div}w=6-\frac{5|z|^2}{2\nu}\) and \(\frac{w\cdot z}{2\nu}=\frac{|z|^2}\nu-\frac{|z|^4}{4\nu^2}\), so the
diffusion term is
\(-\mathcal E_{\nabla,1}+3\nu\mathcal E-\frac74\mathcal E_1+\frac1{8\nu}\int|z|^4|\Omega|^2G\).
Scaling: \(-\frac12\int|z|^2\Omega\cdot(z\cdot\nabla\Omega)G=-\frac14\int|z|^2z\cdot\nabla|\Omega|^2G=\frac14\int|\Omega|^2\operatorname{div}(|z|^2zG)\)
and \(\operatorname{div}(|z|^2zG)=(5|z|^2-\frac{|z|^4}{2\nu})G\), giving
\(\frac54\mathcal E_1-\frac1{8\nu}\int|z|^4|\Omega|^2G\); the fourth moments cancel
against the diffusion term.  The damping \(-\Omega\) gives \(-\mathcal E_1\).
Transport: \(-\int|z|^2\Omega\cdot(V\cdot\nabla)\Omega G=-\frac12\int|z|^2V\cdot\nabla|\Omega|^2G=\frac12\int|\Omega|^2\operatorname{div}(|z|^2VG)\)
with \(\operatorname{div}(|z|^2VG)=(2z\cdot V-\frac{|z|^2(z\cdot V)}{2\nu})G\) (\(\operatorname{div}V=0\)), giving
\(\int|\Omega|^2(z\cdot V)G-\mathcal T_1=4\nu\mathcal T-\mathcal T_1\).  Stretching: \(\mathcal S_1\).  Summing:
\(-\mathcal E_{\nabla,1}+3\nu\mathcal E-\frac74\mathcal E_1+\frac54\mathcal E_1-\mathcal E_1+\mathcal S_1-\mathcal T_1+4\nu\mathcal T\).
\end{proof}

\begin{remark}[Checks]
\label{rem:c3v71-checks}
With the weight \(1\) in place of \(|z|^2\) the same computation gives
\(-\mathcal E_\nabla+(\frac1{8\nu}\mathcal E_1-\frac34\mathcal E)+(\frac34\mathcal E-\frac1{8\nu}\mathcal E_1)-\mathcal E+\mathcal S-\mathcal T\), the identity
of C2-V19, as it must.  On the columnar swirl (\(z\cdot V\equiv0\)) both
transport terms vanish.  The programme's guess of a damping \(2\) was
wrong: the scaling operator \(-\frac12z\cdot\nabla-1\) acting on \(|z|^2|\Omega|^2G\)
gives, after the integrations by parts, the net \(-\frac32\), the weight
contributing \(+\frac54\) against the \(-\frac74\) of the commuted diffusion
and the \(-1\) of the damping; the positive commutator \(3\nu\mathcal E\) is the
Gaussian Laplacian of the weight.
\end{remark}

% =============================================================================
\section{The extremal row}
\label{sec:c3v71-extremal}
% =============================================================================

\begin{lemma}[Bounds on the weighted stretching and transport]
\label{lem:c3v71-bounds}
For every profile of the family and every \(\sigma\),
\begin{equation}
 |\mathcal S_1|\le C_\infty'\mathcal E_1,\qquad
 |\mathcal T_1|\le\Bigl(\frac{\sqrt3}2\kappa+\frac{\kappa}{\sqrt2}+\kappa^2\Bigr)\mathcal E_1+\frac{\kappa}{\sqrt2}\nu\mathcal E+\frac12\mathcal E_{\nabla,1} .
 \label{eq:c3v71-bounds}
\end{equation}
\end{lemma}

\begin{proof}
\(|\mathcal T_1|\le\frac{C_\infty}{4\nu}\int|z|^3|\Omega|^2G\le\frac{C_\infty}{4\nu}\mathcal E_2^{1/2}\mathcal E_1^{1/2}\) with
\(\mathcal E_2:=\int|z|^4|\Omega|^2G\).  The Gaussian Hardy inequality of Lemma
C2-V36-hardy holds for every smooth field \(F\) (its proof uses
\(\operatorname{div}(|F|^2z)=3|F|^2+z\cdot\nabla|F|^2\), not the divergence-free condition):
\(\int|z|^2|F|^2G\le12\nu\|F\|_G^2+16\nu^2\|\nabla F\|_G^2\).  With \(F=|z|\Omega\),
\(\|\nabla F\|_G^2\le2\mathcal E+2\int|z|^2|\nabla\Omega|^2G=2\mathcal E+2\mathcal E_{\nabla,1}/\nu\), so
\(\mathcal E_2\le12\nu\mathcal E_1+32\nu^2\mathcal E+32\nu\mathcal E_{\nabla,1}\) and
\(\mathcal E_2^{1/2}\le(12\nu\mathcal E_1)^{1/2}+(32\nu^2\mathcal E)^{1/2}+(32\nu\mathcal E_{\nabla,1})^{1/2}\).  Then
\(\frac{C_\infty}{4\nu}(12\nu)^{1/2}=\frac{\sqrt3}2\kappa\),
\(\frac{C_\infty}{4\nu}(32\nu^2)^{1/2}(\mathcal E\mathcal E_1)^{1/2}=\sqrt2\kappa(\nu\mathcal E)^{1/2}\mathcal E_1^{1/2}\le\frac\kappa{\sqrt2}(\nu\mathcal E+\mathcal E_1)\),
\(\frac{C_\infty}{4\nu}(32\nu)^{1/2}(\mathcal E_{\nabla,1}\mathcal E_1)^{1/2}=\sqrt2\kappa(\mathcal E_{\nabla,1}\mathcal E_1)^{1/2}\le\frac12\mathcal E_{\nabla,1}+\kappa^2\mathcal E_1\).
The stretching by \(|\nabla V|\le C_\infty'\).
\end{proof}

\begin{theorem}[Extremal weighted enstrophy]
\label{thm:c3v71-extremal}
Let \(W_1\) be a maximizer or minimizer of \(\mathcal E_1\) over \(\mathcal F(x_0)\).  At
\(W_1\), \(\sigma=0\),
\begin{equation}
 \boxed{\quad
 \frac12\mathcal E_{\nabla,1}+\Bigl(\frac32-C_\infty'-\frac{\sqrt3}2\kappa-\frac\kappa{\sqrt2}-\kappa^2\Bigr)\mathcal E_1
 \ \le\ \Bigl(3+\frac\kappa{\sqrt2}+2\sqrt3\kappa+2\kappa^2\Bigr)\nu\,\mathcal E+2\nu\,\mathcal E_\nabla .
 \quad}
 \label{eq:c3v71-row}
\end{equation}
In particular, if \(\theta_1:=\frac32-C_\infty'-\frac{\sqrt3}2\kappa-\frac\kappa{\sqrt2}-\kappa^2>0\), then
\begin{equation}
 \sup_{\mathcal F(x_0)}\mathcal E_1\ \le\ \frac{(3+\frac\kappa{\sqrt2}+2\sqrt3\kappa+2\kappa^2)\nu\,\mathcal E(W_1)+2\nu\,\mathcal E_\nabla(W_1)}{\theta_1}
 \ \le\ \frac{(3+\frac\kappa{\sqrt2}+2\sqrt3\kappa+2\kappa^2)\,2\nu C+4\nu^2C_\infty''^2|G|}{\theta_1} .
 \label{eq:c3v71-sup}
\end{equation}
\end{theorem}

\begin{proof}
Lemma~\ref{lem:c3v34-extremal} with \(F=\mathcal E_1\) (continuous on the family,
\(C^1\) in \(\sigma\) by the derivative bounds) and \eqref{eq:c3v71-identity}:
\(\mathcal E_{\nabla,1}+\frac32\mathcal E_1=3\nu\mathcal E+\mathcal S_1-\mathcal T_1+4\nu\mathcal T\); Lemma~\ref{lem:c3v71-bounds} and
\(4\nu|\mathcal T|\le4\nu(\frac{\sqrt3}2\kappa+\frac{\kappa^2}2)\mathcal E+2\nu\mathcal E_\nabla\) (Lemma~\ref{lem:c3v67-bounds}).
For \eqref{eq:c3v71-sup}: \(\mathcal E\le2\|\nabla V\|_G^2\le2C\) and
\(\mathcal E_\nabla=\nu\int|\nabla\Omega|^2G\le2\nu C_\infty''^2|G|\).
\end{proof}

\begin{assessment}[Item (AC2): concentration, not a constraint]
\label{ass:c3v71-reading}
Unlike the unweighted row (V67), the weighted row does not close to
a relation between constants alone: the commutator \(3\nu\mathcal E\) is
positive and of the order of the damping term, so the row bounds the
second Gaussian moment of the vorticity by the enstrophy and the
vorticity dissipation, i.e.\ it says that at the configuration of
maximal second moment the vorticity is Gaussian-concentrated within
an explicit radius of order \(\nu^{1/2}\) times an explicit function of
the constants, provided \(\theta_1>0\).  This is a concentration bound
of the kind the anchor chapter obtained for the energy by other
means (V16); it is explicit and it is not an exclusion.  The
constants-only constraint of V67 is therefore special: it needs an
observable whose identity has a pure damping and no positive
commutator, and among the Gaussian moments of the vorticity only the
zeroth has that property.  V72 examines the palinstrophy, whose
identity has a pure damping \(\frac32\) but whose stretching brings in
the enstrophy through the second derivatives of \(V\).
\end{assessment}

% =============================================================================
\section{Integrated V71 conclusion and next acquisition order}
\label{sec:c3v71-conclusion}
% =============================================================================

\begin{theorem}[What V71 proves]
\label{thm:c3v71-conclusion}
Relative to the two legacy cycles and V1--V70: the weighted Gaussian
enstrophy identity (Theorem~\ref{thm:c3v71-identity}), the bounds on
its terms (Lemma~\ref{lem:c3v71-bounds}) and the extremal row with the
concentration bound (Theorem~\ref{thm:c3v71-extremal}).  The full
Navier--Stokes stability/global-regularity theorem remains unproved.
\end{theorem}

\begin{proof}
The cited results.
\end{proof}

\begin{programme}[V72 acquisition order]
\label{prog:c3v72}
\begin{enumerate}[label=\textup{(AC\arabic*)},leftmargin=4.8em]
\item The palinstrophy \(\mathcal P:=\int|\nabla\Omega|^2G\): derive its identity from the
      gradient of the vorticity equation (damping \(\frac32\) from
      \(-\frac12z\cdot\nabla-1\) acting on \(\nabla\Omega\), diffusion \(-\nu\int|\nabla^2\Omega|^2G\),
      transport \(-\frac1{4\nu}\int|\nabla\Omega|^2(z\cdot V)G-\int\nabla\Omega:(\nabla V\cdot\nabla)\Omega G\), stretching
      \(\int\nabla\Omega:\nabla((\Omega\cdot\nabla)V)G\)); bound the terms by \(C_\infty'\), \(C_\infty''\), \(\kappa\) and
      the Hardy inequality; state the extremal row.
\item Record whether the row is a constants-only constraint or a
      concentration bound (the stretching's second-derivative part
      brings \(\mathcal E^{1/2}\mathcal P^{1/2}\)); state the resulting explicit bound on
      \(\sup\mathcal P\) in terms of \(\mathcal E\) when the damping exceeds the
      first-derivative constants.
\item The next compulsory companion audit is V75.
\end{enumerate}
\end{programme}

\begin{assessment}[Position after V71]
\label{ass:c3v71-position}
The weighted enstrophy row gives an explicit concentration bound on
the vorticity's second Gaussian moment, not a constraint on
constants.  No full stability or global-regularity claim is made.
\end{assessment}

% =============================================================================
\section*{Version V71 (third cycle) append-only record}
\addcontentsline{toc}{section}{Version V71 (third cycle) append-only record}
% =============================================================================

The complete V70 source of the third cycle was retained before this
continuation.  Its SHA--256 digest was
\begin{center}\scriptsize\ttfamily
0f705a7e9149d7f7e89c053e20bfc9d31dbc660f0c0907bc981a334088997e7f.
\end{center}
V71 proved the weighted enstrophy identity and its extremal row.
No full regularity claim is made.

% =============================================================================
\section{V72: the palinstrophy identity and its extremal row}
\label{sec:c3v72-entry}
% =============================================================================

Items (AC1) and (AC2) of Programme~\ref{prog:c3v72} are carried out.
(AC1): the identity for the Gaussian palinstrophy is derived from
the gradient of the vorticity equation: a pure damping of three
halves, the diffusion of the second derivatives, a transport term
of the same form as before, and a stretching term with two parts,
one bounded by the gradient constant times the palinstrophy and one
by the second-derivative constant times the geometric mean of the
palinstrophy and the enstrophy.  (AC2): at an extremum of the
palinstrophy the row closes to a ratio bound, palinstrophy at most
an explicit multiple of the enstrophy, whenever the damping exceeds
twice the gradient constant plus the velocity terms; the Gaussian
Poincar\'e inequality turns the ratio bound into a bound on the
palinstrophy by the Gaussian angular momentum alone when the
multiple is small, but not into a constants-only constraint.  The
palinstrophy row is therefore of the same kind as the weighted row
of V71, and the joint reading of V73 is prepared.  No global
regularity claim is made.  The next compulsory companion audit is
V75.

% =============================================================================
\section{The identity}
\label{sec:c3v72-identity}
% =============================================================================

Put \(\mathcal P:=\int|\nabla\Omega|^2G\) (so that \(\mathcal E_\nabla=\nu\mathcal P\)), \(\mathcal P_\nabla:=\nu\int|\nabla^2\Omega|^2G\),
\begin{equation}
 \mathcal S_{\mathcal P}:=\int\partial_j\Omega\cdot\partial_j\bigl((\Omega\cdot\nabla)V\bigr)G-\int\partial_j\Omega\cdot(\partial_jV\cdot\nabla)\Omega\,G,\qquad
 \mathcal T_{\mathcal P}:=\frac1{4\nu}\int|\nabla\Omega|^2(z\cdot V)\,G
 \label{eq:c3v72-notation}
\end{equation}
(summation over \(j\)).

\begin{theorem}[Gaussian palinstrophy identity]
\label{thm:c3v72-identity}
For every profile of the family and every \(\sigma\),
\begin{equation}
 \boxed{\quad
 \frac12\frac{\dd}{\dd\sigma}\mathcal P=-\mathcal P_\nabla-\frac32\mathcal P+\mathcal S_{\mathcal P}-\mathcal T_{\mathcal P} .
 \quad}
 \label{eq:c3v72-identity}
\end{equation}
\end{theorem}

\begin{proof}
Differentiate the vorticity equation:
\(\partial_\sigma\partial_j\Omega=\nu\Delta\partial_j\Omega-\frac12z\cdot\nabla\partial_j\Omega-\frac32\partial_j\Omega-\partial_j((V\cdot\nabla)\Omega)+\partial_j((\Omega\cdot\nabla)V)\),
since \(\partial_j(z\cdot\nabla\Omega)=\partial_j\Omega+z\cdot\nabla\partial_j\Omega\).  Multiply by \(\partial_j\Omega G\), sum
over \(j\) and integrate.  The \(G\)-symmetry of \(\nu\Delta-\frac12z\cdot\nabla\) gives
\(\sum_j\langle\partial_j\Omega,(\nu\Delta-\frac12z\cdot\nabla)\partial_j\Omega\rangle_G=-\nu\int|\nabla^2\Omega|^2G\); the damping
gives \(-\frac32\mathcal P\); the transport splits as
\(\partial_j((V\cdot\nabla)\Omega)=(\partial_jV\cdot\nabla)\Omega+(V\cdot\nabla)\partial_j\Omega\), and
\(-\sum_j\int\partial_j\Omega\cdot(V\cdot\nabla)\partial_j\Omega\,G=-\frac12\int V\cdot\nabla|\nabla\Omega|^2G=\frac12\int|\nabla\Omega|^2\operatorname{div}(VG)=-\mathcal T_{\mathcal P}\)
by \(\operatorname{div}V=0\); the remaining transport part and the stretching form
\(\mathcal S_{\mathcal P}\).
\end{proof}

\begin{lemma}[Bounds]
\label{lem:c3v72-bounds}
For every profile of the family and every \(\sigma\),
\begin{equation}
 |\mathcal S_{\mathcal P}|\le2C_\infty'\mathcal P+C_\infty''\,\mathcal E^{1/2}\mathcal P^{1/2},\qquad
 |\mathcal T_{\mathcal P}|\le\Bigl(\frac{\sqrt3}2\kappa+\frac{\kappa^2}2\Bigr)\mathcal P+\frac12\mathcal P_\nabla .
 \label{eq:c3v72-bounds}
\end{equation}
\end{lemma}

\begin{proof}
\(\partial_j((\Omega\cdot\nabla)V)=(\partial_j\Omega\cdot\nabla)V+(\Omega\cdot\nabla)\partial_jV\): the first pairs with \(\partial_j\Omega\) at
most \(C_\infty'|\nabla\Omega|^2\), the second at most \(C_\infty''|\Omega||\nabla\Omega|\), integrated by
Cauchy--Schwarz; the transport remainder \(|\partial_j\Omega\cdot(\partial_jV\cdot\nabla)\Omega|\le C_\infty'|\nabla\Omega|^2\).
The transport \(\mathcal T_{\mathcal P}\) by the Gaussian Hardy inequality for the field
\(\nabla\Omega\) (valid for every field, Lemma~\ref{lem:c3v71-bounds}'s proof),
exactly as in Lemma~\ref{lem:c3v67-bounds} with \(\Omega\) replaced by \(\nabla\Omega\).
\end{proof}

% =============================================================================
\section{The extremal row}
\label{sec:c3v72-extremal}
% =============================================================================

\begin{theorem}[Extremal palinstrophy: a ratio bound]
\label{thm:c3v72-extremal}
Let \(W_{\mathcal P}\) be a maximizer or minimizer of \(\mathcal P\) over \(\mathcal F(x_0)\), and
\(\theta_{\mathcal P}:=\frac32-2C_\infty'-\frac{\sqrt3}2\kappa-\frac{\kappa^2}2\).  At \(W_{\mathcal P}\), \(\sigma=0\),
\begin{equation}
 \boxed{\quad
 \frac12\mathcal P_\nabla+\theta_{\mathcal P}\,\mathcal P\ \le\ C_\infty''\,\mathcal E^{1/2}\mathcal P^{1/2},
 \qquad\text{hence, if }\theta_{\mathcal P}>0,\qquad
 \sup_{\mathcal F(x_0)}\mathcal P\ \le\ \Bigl(\frac{C_\infty''}{\theta_{\mathcal P}}\Bigr)^2\mathcal E(W_{\mathcal P})\ \le\ 2\Bigl(\frac{C_\infty''}{\theta_{\mathcal P}}\Bigr)^2C .
 \quad}
 \label{eq:c3v72-row}
\end{equation}
\end{theorem}

\begin{proof}
Lemma~\ref{lem:c3v34-extremal} with \(F=\mathcal P\) and \eqref{eq:c3v72-identity}:
\(\mathcal P_\nabla+\frac32\mathcal P=\mathcal S_{\mathcal P}-\mathcal T_{\mathcal P}\); Lemma~\ref{lem:c3v72-bounds}; divide by
\(\mathcal P^{1/2}\) when \(\mathcal P>0\) (if \(\mathcal P(W_{\mathcal P})=0\) at a maximizer the bound is
trivial); \(\mathcal E\le2\|\nabla V\|_G^2\le2C\).
\end{proof}

\begin{proposition}[The Gaussian Poincar\'e link]
\label{prop:c3v72-poincare}
With \(\bar\Omega:=\frac1{|G|}\int\Omega G=\frac1{2\nu|G|}\int z\times V\,G\) (the Gaussian angular
momentum, up to the factor), for every profile
\begin{equation}
 \mathcal E\ \le\ 2\nu\,\mathcal P+|G|\,|\bar\Omega|^2 ,
 \label{eq:c3v72-poincare}
\end{equation}
so that at a maximizer of \(\mathcal P\) with \(\theta_{\mathcal P}>0\) and
\(q_{\mathcal P}:=2\nu(C_\infty''/\theta_{\mathcal P})^2<1\),
\begin{equation}
 \sup_{\mathcal F(x_0)}\mathcal P\ \le\ \frac{(C_\infty''/\theta_{\mathcal P})^2}{1-q_{\mathcal P}}\,|G|\,|\bar\Omega(W_{\mathcal P})|^2 .
 \label{eq:c3v72-angular}
\end{equation}
\end{proposition}

\begin{proof}
The Gaussian Poincar\'e inequality (C2-V8-spectrum(iii)):
\(\int|\Omega-\bar\Omega|^2G\le2\nu\int|\nabla\Omega|^2G\), and \(\int|\Omega|^2G=\int|\Omega-\bar\Omega|^2G+|G||\bar\Omega|^2\).
\(\int\Omega G=\int\operatorname{curl}V\,G=-\int\nabla G\times V=\frac1{2\nu}\int z\times V\,G\).  Then
\eqref{eq:c3v72-row} with \(\mathcal E\le2\nu\mathcal P+|G||\bar\Omega|^2\) and \(\mathcal P(1-q_{\mathcal P})\le\dots\).
\end{proof}

\begin{assessment}[Item (AC2): ratio bounds, not constraints]
\label{ass:c3v72-reading}
The palinstrophy row has a pure damping, like the enstrophy row, but
its stretching contains the second derivatives of the velocity,
which pair the palinstrophy with the enstrophy; the row therefore
bounds the palinstrophy by an explicit multiple of the enstrophy
(a regularity ratio at the extremal configuration), and through the
Gaussian Poincar\'e inequality by the Gaussian angular momentum when
the multiple is small.  A constants-only constraint would follow if
the angular momentum vanished at the maximizer of \(\mathcal P\), which the
class does not give.  The vorticity ledger's rows are thus of three
kinds: the enstrophy row (constants-only constraint, V67), the
weighted rows (concentration bounds, V71), and the palinstrophy row
(ratio bound).  Each is explicit; only the first relates explicit
constants alone.  The joint reading (V73) asks what the three say
together with the structure of (E5).
\end{assessment}

% =============================================================================
\section{Integrated V72 conclusion and next acquisition order}
\label{sec:c3v72-conclusion}
% =============================================================================

\begin{theorem}[What V72 proves]
\label{thm:c3v72-conclusion}
Relative to the two legacy cycles and V1--V71: the palinstrophy
identity (Theorem~\ref{thm:c3v72-identity}), its bounds
(Lemma~\ref{lem:c3v72-bounds}), the extremal ratio bound
(Theorem~\ref{thm:c3v72-extremal}) and the Poincar\'e link
(Proposition~\ref{prop:c3v72-poincare}).  The full Navier--Stokes
stability/global-regularity theorem remains unproved.
\end{theorem}

\begin{proof}
The cited results.
\end{proof}

\begin{programme}[V73 acquisition order]
\label{prog:c3v73}
\begin{enumerate}[label=\textup{(AC\arabic*)},leftmargin=4.8em]
\item The joint reading: with the structure of (E5) for small
      velocity constants (\(C_\infty^{(k)}\le c_k\kappa\) for \(\kappa\le\kappa_{\rm small}\), the
      parabolic derivative bounds for small solutions of the rescaled
      equation, recorded as the form in which (E5) enters), derive
      from \eqref{eq:c3v67-constraint} the explicit threshold
      \(\kappa\ge\kappa_1(c_1)\), and from the ratio and concentration rows the
      regimes in which they apply; compare \(\kappa_1\) with \(\varepsilon_{\rm reg}\) and
      with the constant-chasing thresholds of C2-V36--V38 (\(\kappa<\frac12\),
      \(\kappa<\sqrt2\)); record whether any row says something not implied
      by (E5).
\item Record the vorticity ledger's three rows as explicit results
      of the cycle, in the form of a table with the observable, the
      damping, the commutator, the bound and the kind of conclusion.
\item The next compulsory companion audit is V75.
\end{enumerate}
\end{programme}

\begin{assessment}[Position after V72]
\label{ass:c3v72-position}
The palinstrophy row gives an explicit ratio bound; the vorticity
ledger has one constraint and two kinds of bound.  No full stability
or global-regularity claim is made.
\end{assessment}

% =============================================================================
\section*{Version V72 (third cycle) append-only record}
\addcontentsline{toc}{section}{Version V72 (third cycle) append-only record}
% =============================================================================

The complete V71 source of the third cycle was retained before this
continuation.  Its SHA--256 digest was
\begin{center}\scriptsize\ttfamily
aff4ee6122b4c4a41f75be9e2304b3b1598be489548f10775e121fc4fabdc4f8.
\end{center}
V72 proved the palinstrophy identity and its extremal row.  No full
regularity claim is made.

% =============================================================================
\section{V73: the joint reading of the vorticity ledger against (E5)}
\label{sec:c3v73-entry}
% =============================================================================

Items (AC1) and (AC2) of Programme~\ref{prog:c3v73} are carried out.
(AC1): the form in which (E5) enters, derivative constants
proportional to the velocity constant for small velocity constants,
is recorded as an explicit hypothesis on the imports with unnamed
constants \(\gamma_k\); the constraint of V67 is first read without
any hypothesis as a lower bound on the gradient constant, and under
the hypothesis as an explicit lower bound \(\kappa_1(\gamma_1)\) on the
velocity constant, bracketed on both sides; the regimes of the
concentration and ratio rows are explicit upper thresholds, and the
three rows apply on a common nonempty window if and only if the
first-derivative constant \(\gamma_1\) lies in an explicit interval.
The comparison with \(\varepsilon_{\rm reg}\) and with the constant-chasing
thresholds is a criterion in the constants: the enstrophy floor is
the stronger of the two floors if and only if \(\gamma_1<1/\varepsilon_{\rm reg}-\varepsilon_{\rm reg}/2-\sqrt3/2\),
and it empties the regime \(\kappa<\frac12\) of C2-V36 if and only if
\(\gamma_1<\frac74-\frac{\sqrt3}2\).  No row excludes anything for any value of
the constants.  (AC2): the three rows are tabulated.  No global
regularity claim is made.  The next compulsory companion audit is
V75.

% =============================================================================
\section{The form in which (E5) enters}
\label{sec:c3v73-E5}
% =============================================================================

The derivative constants \(C_\infty^{(k)}\) of the family are the bounds
\(|\nabla^kV|\le C_\infty^{(k)}\) transported to the family by C2-V4 from the derivative
bounds of (E5); in the units of the profile equation (\(z\) measured
in units of \(\nu^{1/2}\), \(V\) in units of \(\nu^{1/2}\)) the first, \(C_\infty'\), is
dimensionless, like \(\kappa=C_\infty\nu^{-1/2}\), which is why \eqref{eq:c3v67-constraint}
mixes \(C_\infty'\) and \(\kappa\) without a factor of \(\nu\), while \(\nu^{1/2}C_\infty''\) is
the dimensionless second-derivative constant.  The series has not written the
constants of the derivative bounds (Firewall~\ref{fw:c3v28-lipschitz}
for the orders three and four; the orders one and two enter only
through \(C_\infty',C_\infty''\)).  The joint reading needs one structural
property of these bounds: for a solution of the rescaled equation
bounded by a small multiple of \(\nu^{1/2}\), the interior parabolic
estimates give derivative bounds proportional to the velocity bound,
the nonlinear terms being absorbed for a small velocity constant.
It is recorded as a hypothesis on the imports, not as a result of
the series.

\begin{definition}[Hypothesis (E5-small)]
\label{def:c3v73-E5small}
There are constants \(\kappa_{\rm small}>0\) and \(\gamma_1,\gamma_2>0\), independent of the
point and of the profile, such that at every homogeneous type-I
point with velocity constant \(\kappa\le\kappa_{\rm small}\) the derivative constants
of the family satisfy
\begin{equation}
 C_\infty'\le\gamma_1\kappa,\qquad C_\infty''\le\gamma_2\,\kappa\,\nu^{-1/2} .
 \label{eq:c3v73-E5small}
\end{equation}
\end{definition}

\begin{remark}[Status and use]
\label{rem:c3v73-status}
(E5-small) is the small-constant form of the derivative bounds of
(E5) as transported by C2-V4; the general form has the derivative
constants bounded by polynomials in \(\kappa\) without a linear leading
term being asserted, and (E5-small) is the statement that the
linear term is the leading one below \(\kappa_{\rm small}\).  The constants
\(\gamma_1,\gamma_2,\kappa_{\rm small}\) are not computed in the series (the \(c_1,c_p\) of V47 are
kernel constants, unrelated).  Every statement of this version that
uses (E5-small) is explicit in \(\gamma_1,\gamma_2\); every statement that does
not is marked as such.
\end{remark}

% =============================================================================
\section{The enstrophy row as a floor}
\label{sec:c3v73-floor}
% =============================================================================

\begin{proposition}[Gradient floor, without hypothesis]
\label{prop:c3v73-gradient-floor}
Let \(\kappa_0:=-\frac{\sqrt3}2+\bigl(\frac34+2\bigr)^{1/2}=\frac{\sqrt{11}-\sqrt3}2\) (\(\approx0.792\)).  At every
homogeneous type-I point with \(\kappa<\kappa_0\),
\begin{equation}
 C_\infty'\ \ge\ 1-\frac{\sqrt3}2\kappa-\frac{\kappa^2}2\ >\ 0 .
 \label{eq:c3v73-gradient-floor}
\end{equation}
\end{proposition}

\begin{proof}
\eqref{eq:c3v67-constraint}; the right side vanishes at \(\kappa=\kappa_0\), the
positive root of \(\kappa^2/2+\frac{\sqrt3}2\kappa=1\), and is decreasing in \(\kappa\).
\end{proof}

\begin{theorem}[Velocity floor under (E5-small)]
\label{thm:c3v73-velocity-floor}
Put \(a:=\gamma_1+\frac{\sqrt3}2\) and
\begin{equation}
 \kappa_1(\gamma_1):=-a+\bigl(a^2+2\bigr)^{1/2}=\frac2{a+(a^2+2)^{1/2}},\qquad
 \frac2{2a+\sqrt2}\ \le\ \kappa_1\ <\ \frac1a\ \le\ \frac2{\sqrt3} .
 \label{eq:c3v73-kappa1}
\end{equation}
Under (E5-small), every homogeneous type-I point satisfies
\(\kappa\ge\min\{\kappa_1(\gamma_1),\kappa_{\rm small}\}\); if \(\kappa_1(\gamma_1)>\kappa_{\rm small}\), every homogeneous
type-I point satisfies \(\kappa>\kappa_{\rm small}\).
\end{theorem}

\begin{proof}
If \(\kappa\le\kappa_{\rm small}\), \eqref{eq:c3v67-constraint} and \eqref{eq:c3v73-E5small} give
\(f(\kappa):=a\kappa+\kappa^2/2\ge1\); \(f\) is increasing on \([0,\infty)\) with \(f(\kappa_1)=1\),
so \(\kappa\ge\kappa_1\).  The bracket: \((a^2+2)^{1/2}\in(a,a+\sqrt2]\) since
\((a+\sqrt2)^2=a^2+2+2\sqrt2a\ge a^2+2\); and \(a\ge\frac{\sqrt3}2\).
\end{proof}

\begin{assessment}[What the floor says]
\label{ass:c3v73-floor}
Without any hypothesis, the enstrophy row says that a homogeneous
type-I point with velocity constant below \(\kappa_0\approx0.79\) has a
gradient constant of order one, at least \(1-\frac{\sqrt3}2\kappa-\frac{\kappa^2}2\); (E5) bounds
\(C_\infty'\) above, and this bounds it below, so the two together confine
the gradient constant to an explicit interval at every such point.
Under (E5-small) the two bounds are incompatible below \(\kappa_1(\gamma_1)\),
and the velocity constant is at least \(\kappa_1\), which lies between
\(2/(2\gamma_1+\sqrt3+\sqrt2)\) and \(1/(\gamma_1+\sqrt3/2)\): for \(\gamma_1=1\), between \(0.39\) and
\(0.54\) (the exact value \(0.475\)); for \(\gamma_1=3\), between \(0.22\) and \(0.26\).
The floor is of \(\varepsilon\)-regularity type and is not an exclusion.
\end{assessment}

% =============================================================================
\section{The regimes of the concentration and ratio rows}
\label{sec:c3v73-regimes}
% =============================================================================

\begin{proposition}[Explicit applicability thresholds]
\label{prop:c3v73-thresholds}
Under (E5-small), for \(\kappa\le\kappa_{\rm small}\),
\begin{align}
 \theta_1&\ \ge\ \theta_1^-(\kappa):=\frac32-b\kappa-\kappa^2,\qquad b:=\gamma_1+\frac{\sqrt3}2+\frac1{\sqrt2},\qquad
 \theta_1^->0\iff\kappa<\kappa_c(\gamma_1):=-\frac b2+\Bigl(\frac{b^2}4+\frac32\Bigr)^{1/2},
 \label{eq:c3v73-theta1}\\
 \theta_{\mathcal P}&\ \ge\ \theta_{\mathcal P}^-(\kappa):=\frac32-a'\kappa-\frac{\kappa^2}2,\qquad a':=2\gamma_1+\frac{\sqrt3}2,\qquad
 \theta_{\mathcal P}^->0\iff\kappa<\kappa_r(\gamma_1):=-a'+\bigl(a'^2+3\bigr)^{1/2},
 \label{eq:c3v73-thetaP}
\end{align}
with \(\theta_1,\theta_{\mathcal P}\) of Theorems~\ref{thm:c3v71-extremal} and
\ref{thm:c3v72-extremal}.  On \(\kappa<\kappa_c\) the concentration bound
\eqref{eq:c3v71-sup} holds with \(\theta_1^-\) in place of \(\theta_1\) and
\(\gamma_2\kappa\nu^{-1/2}\) in place of \(C_\infty''\); on \(\kappa<\kappa_r\) the ratio bound
\eqref{eq:c3v72-row} holds in the form
\(\sup\mathcal P\le2(\gamma_2\kappa/\theta_{\mathcal P}^-)^2C/\nu\).
\end{proposition}

\begin{proof}
Substitute \eqref{eq:c3v73-E5small} in the definitions of \(\theta_1,\theta_{\mathcal P}\);
both lower bounds are decreasing in \(\kappa\) and vanish at the stated
positive roots.  The bounds of V71--V72 are monotone in \(\theta\) and in
\(C_\infty''\).
\end{proof}

\begin{theorem}[The common window]
\label{thm:c3v73-window}
Under (E5-small) with \(\kappa_{\rm small}\ge\kappa_1(\gamma_1)\):
\begin{enumerate}[label=\textup{(\arabic*)},leftmargin=3.2em]
\item The ratio row applies at the floor, \(\theta_{\mathcal P}^-(\kappa_1)>0\), if and only
      if \(\gamma_1\kappa_1<\frac12\), if and only if
      \(\gamma_1<\frac{\sqrt3+\sqrt7}4\) (\(\approx1.094\)); then \(\kappa_1<\kappa_r\).
\item The concentration row applies at the floor, \(\theta_1^-(\kappa_1)>0\), if
      and only if \(\kappa_1<\kappa_*:=\frac{\sqrt6-\sqrt2}2\) (\(\approx0.518\)), if and only if
      \(\gamma_1>\frac{3\sqrt2+\sqrt6-2\sqrt3}4\) (\(\approx0.807\)); then \(\kappa_1<\kappa_c\).
\item The three rows apply on a common nonempty window
      \([\kappa_1,\min\{\kappa_c,\kappa_r,\kappa_{\rm small}\})\) if and only if
      \(\frac{3\sqrt2+\sqrt6-2\sqrt3}4<\gamma_1<\frac{\sqrt3+\sqrt7}4\).
\end{enumerate}
\end{theorem}

\begin{proof}
At \(\kappa=\kappa_1\), \(a\kappa_1+\kappa_1^2/2=1\).  (1)
\(\theta_{\mathcal P}^-(\kappa_1)=\frac32-\gamma_1\kappa_1-(a\kappa_1+\kappa_1^2/2)=\frac12-\gamma_1\kappa_1\).  With
\(\kappa_1=2/(a+(a^2+2)^{1/2})\), \(\gamma_1\kappa_1<\frac12\iff4\gamma_1-a<(a^2+2)^{1/2}\iff3\gamma_1-\frac{\sqrt3}2<(a^2+2)^{1/2}\);
this holds when the left side is nonpositive, and otherwise
squares to \(9\gamma_1^2-3\sqrt3\gamma_1+\frac34<\gamma_1^2+\sqrt3\gamma_1+\frac{11}4\), i.e.\
\(4\gamma_1^2-2\sqrt3\gamma_1-1<0\), whose positive root is \((\sqrt3+\sqrt7)/4\); the
nonpositive case \(\gamma_1\le1/(2\sqrt3)\) lies below that root.  Since
\(\theta_{\mathcal P}^-\) is decreasing, \(\theta_{\mathcal P}^-(\kappa_1)>0\iff\kappa_1<\kappa_r\).
(2) \(\theta_1^-(\kappa_1)=\frac32-(a\kappa_1+\kappa_1^2/2)-\frac{\kappa_1}{\sqrt2}-\frac{\kappa_1^2}2=\frac12-\frac{\kappa_1}{\sqrt2}-\frac{\kappa_1^2}2\),
positive iff \(\kappa_1^2+\sqrt2\kappa_1-1<0\) iff \(\kappa_1<\kappa_*\), the positive root.
Since \(f\) is increasing, \(\kappa_1<\kappa_*\iff f(\kappa_*)>1\); \(\kappa_*^2=2-\sqrt3\), so
\(f(\kappa_*)=a\kappa_*+1-\frac{\sqrt3}2>1\iff a>\frac{\sqrt3}{2\kappa_*}=\frac{\sqrt3(\sqrt6+\sqrt2)}4=\frac{3\sqrt2+\sqrt6}4\),
i.e.\ \(\gamma_1>\frac{3\sqrt2+\sqrt6}4-\frac{\sqrt3}2\).  (3) is (1) and (2) together,
the window being nonempty exactly when both thresholds exceed
\(\kappa_1\).
\end{proof}

\begin{assessment}[Item (AC1): what the rows say jointly]
\label{ass:c3v73-joint}
The three rows are governed by one unknown constant, \(\gamma_1\), the
constant of the first-derivative bound for small solutions, and by
\(\gamma_2\) only through the size of the bounds, not their
applicability.  The enstrophy row is a floor on \(\kappa\) that decreases
with \(\gamma_1\); the ratio row applies below a threshold that also
decreases with \(\gamma_1\), faster, so the two are compatible only for
\(\gamma_1<1.094\); the concentration row applies below a threshold that
decreases more slowly than the floor, so it is compatible with the
floor only for \(\gamma_1>0.807\).  In the interval \((0.807,1.094)\) all three
rows apply at every homogeneous type-I point with \(\kappa\) in the
window; outside it at least one row is void at the floor.  In no case
is a row violated: the rows are consistent with (E5) for every value
of \(\gamma_1,\gamma_2\), and the joint reading produces no exclusion.  What it
produces is the location of the decision: the value of \(\gamma_1\), a
constant of parabolic regularity that the series has not computed
and that is in principle computable.
\end{assessment}

% =============================================================================
\section{Comparison with the floors and thresholds of the second cycle}
\label{sec:c3v73-comparison}
% =============================================================================

\begin{proposition}[Comparison with \(\varepsilon_{\rm reg}\)]
\label{prop:c3v73-epsreg}
Let \(\varepsilon_{\rm reg}\) be the \(\varepsilon\)-regularity floor of Firewall C2-V36-AC2
(\(\kappa\ge\varepsilon_{\rm reg}\) at every singular point).  Under (E5-small) with
\(\kappa_{\rm small}\ge\kappa_1\), the enstrophy floor is the stronger, \(\kappa_1(\gamma_1)>\varepsilon_{\rm reg}\), if
and only if
\begin{equation}
 \gamma_1\ <\ \frac1{\varepsilon_{\rm reg}}-\frac{\varepsilon_{\rm reg}}2-\frac{\sqrt3}2 .
 \label{eq:c3v73-epsreg}
\end{equation}
\end{proposition}

\begin{proof}
\(\kappa_1>\varepsilon\iff f(\varepsilon)<1\iff a\varepsilon+\varepsilon^2/2<1\iff a<1/\varepsilon-\varepsilon/2\).
\end{proof}

\begin{proposition}[Comparison with the constant-chasing thresholds]
\label{prop:c3v73-chasing}
\begin{enumerate}[label=\textup{(\arabic*)},leftmargin=3.2em]
\item \(\kappa_1(\gamma_1)<2/\sqrt3<\sqrt2\) for every \(\gamma_1\): the floor never empties
      the regime \(\kappa<\sqrt2\) in which the mean ledger of C2-V41 closes
      on the upper side.
\item \(\kappa_1(\gamma_1)>\frac12\) if and only if \(\gamma_1<\frac74-\frac{\sqrt3}2\) (\(\approx0.884\)).  Under
      (E5-small) with \(\kappa_{\rm small}\ge\frac12\) and \(\gamma_1<\frac74-\frac{\sqrt3}2\), no homogeneous
      type-I point has \(\kappa<\frac12\): the regime of the mean constraint
      of Theorem C2-V36-constraint, of \eqref{eq:c3v32-constraint-small}
      and of Proposition~\ref{prop:c3v26-necessary} is empty.
\end{enumerate}
\end{proposition}

\begin{proof}
(1) \eqref{eq:c3v73-kappa1}.  (2) \(\kappa_1>\frac12\iff f(\frac12)<1\iff\frac a2+\frac18<1\iff a<\frac74\).
If \(\kappa<\frac12\le\kappa_{\rm small}\), Theorem~\ref{thm:c3v73-velocity-floor} gives
\(\kappa\ge\kappa_1>\frac12\).
\end{proof}

\begin{assessment}[Consistency, not exclusion]
\label{ass:c3v73-consistency}
Proposition~\ref{prop:c3v73-chasing}(2) is a consistency statement
between two chapters of the series: the second cycle's constraints
bite for a small velocity constant, and the enstrophy row says that
a small velocity constant forces a gradient constant of order one,
which (E5-small) forbids below \(\kappa_1\).  If \(\gamma_1<0.884\) the two
regimes are disjoint and the constant-chasing constraints of
C2-V36 and V32 hold vacuously at homogeneous type-I points; if
\(\gamma_1\ge0.884\) they hold on \([\kappa_1,\frac12)\).  Neither case is an exclusion,
since every point in the admissible range satisfies every row.  The
comparison with \(\varepsilon_{\rm reg}\) is of the same kind: the series has
neither \(\varepsilon_{\rm reg}\) nor \(\gamma_1\) numerically, and \eqref{eq:c3v73-epsreg} records
which of the two floors is the operative one as a criterion.
\end{assessment}

\begin{firewall}[What is and is not implied by (E5)]
\label{fw:c3v73-not-implied}
Every row of the vorticity ledger is a consequence of the class
bounds, the compactness of the family and the vorticity equation;
none is a consequence of (E5) as an inequality.  The gradient floor
\eqref{eq:c3v73-gradient-floor} is not implied by (E5), which bounds
\(C_\infty'\) above only; it is implied by (E5) together with
\(\varepsilon\)-regularity only in the trivial sense that \(\kappa\ge\varepsilon_{\rm reg}\) is a
floor too, and the two floors are compared by
\eqref{eq:c3v73-epsreg}.  The concentration and ratio rows are
bounds on Gaussian quantities of the vorticity that (E5) bounds
otherwise (by the derivative constants and the Gaussian integral of
the weight, \(\mathcal E_1\le C_\infty'^2\cdot\int|z|^2G\)-type); the rows are sharper when
the extremal configuration's enstrophy is small, which is not known.
No row excludes a homogeneous type-I point for any values of
\(\gamma_1,\gamma_2,\varepsilon_{\rm reg}\); the vorticity ledger alone cannot close the
programme, and this is recorded as the negative record of the block.
\end{firewall}

% =============================================================================
\section{The three rows}
\label{sec:c3v73-table}
% =============================================================================

\begin{theorem}[Item (AC2): the vorticity ledger at extremal configurations]
\label{thm:c3v73-table}
At a homogeneous type-I point, with \(\theta_1,\theta_{\mathcal P}\) as in V71--V72 and
\(\Xi_1:=3+\frac\kappa{\sqrt2}+2\sqrt3\kappa+2\kappa^2\):

\medskip
\begin{center}\small
\begin{tabular}{@{}p{0.13\textwidth}p{0.07\textwidth}p{0.14\textwidth}p{0.4\textwidth}p{0.16\textwidth}@{}}
\toprule
\textbf{Observable} & \textbf{Damping} & \textbf{Commutator} & \textbf{Bound at the extremal configuration} & \textbf{Kind}\\
\midrule
\(\mathcal E=\int|\Omega|^2G\) (V67) & \(1\) & none &
 \(C_\infty'+\frac{\sqrt3}2\kappa+\frac{\kappa^2}2\ge1\) & constants-only constraint; floor \(\kappa\ge\kappa_1(\gamma_1)\)\\[4pt]
\(\mathcal E_1=\int|z|^2|\Omega|^2G\) (V71) & \(\frac32\) & \(+3\nu\mathcal E+4\nu\mathcal T\) &
 \(\sup\mathcal E_1\le\bigl(2\nu C\,\Xi_1+4\nu^2C_\infty''^2|G|\bigr)/\theta_1\) if \(\theta_1>0\) & concentration; regime \(\kappa<\kappa_c\)\\[4pt]
\(\mathcal P=\int|\nabla\Omega|^2G\) (V72) & \(\frac32\) & none (stretching couples to \(\mathcal E\) by \(C_\infty''\)) &
 \(\sup\mathcal P\le2(C_\infty''/\theta_{\mathcal P})^2C\) if \(\theta_{\mathcal P}>0\); \(\sup\mathcal P\le\frac{(C_\infty''/\theta_{\mathcal P})^2}{1-q_{\mathcal P}}|G||\bar\Omega|^2\) if \(q_{\mathcal P}<1\) & ratio; regime \(\kappa<\kappa_r\)\\
\bottomrule
\end{tabular}
\end{center}
\medskip

\noindent
Each bound is explicit in the class constants; the regimes are those
of Proposition~\ref{prop:c3v73-thresholds} under (E5-small).
\end{theorem}

\begin{proof}
Theorems~\ref{thm:c3v67-extremal}, \ref{thm:c3v71-extremal},
\ref{thm:c3v72-extremal}, Proposition~\ref{prop:c3v72-poincare},
Theorem~\ref{thm:c3v73-velocity-floor} and
Proposition~\ref{prop:c3v73-thresholds}.
\end{proof}

% =============================================================================
\section{Integrated V73 conclusion and next acquisition order}
\label{sec:c3v73-conclusion}
% =============================================================================

\begin{theorem}[What V73 establishes]
\label{thm:c3v73-conclusion}
Relative to the two legacy cycles and V1--V72: the gradient floor
(Proposition~\ref{prop:c3v73-gradient-floor}), the velocity floor
\(\kappa_1(\gamma_1)\) with its bracket (Theorem~\ref{thm:c3v73-velocity-floor}), the
applicability thresholds and the common window
(Proposition~\ref{prop:c3v73-thresholds},
Theorem~\ref{thm:c3v73-window}), the comparison criteria
(Propositions~\ref{prop:c3v73-epsreg}, \ref{prop:c3v73-chasing}), the
negative record (Firewall~\ref{fw:c3v73-not-implied}) and the table
(Theorem~\ref{thm:c3v73-table}).  The full Navier--Stokes
stability/global-regularity theorem remains unproved.
\end{theorem}

\begin{proof}
The cited results.
\end{proof}

\begin{programme}[V74 acquisition order]
\label{prog:c3v74}
\begin{enumerate}[label=\textup{(AC\arabic*)},leftmargin=4.8em]
\item Consolidate the vorticity-ledger block V71--V73 as one theorem
      with its negative record.
\item Fix the direction for V76--V80.  The candidate: the linear
      functionals of the vorticity, beginning with the Gaussian mean
      \(\bar\Omega\) of V72, whose identity is pure damping (one unit) with a
      source that is the Gaussian moment of \(z\times(V\times\Omega)\), the Lamb
      vector of V42 in the vorticity form; derive the identity, check
      it on the columnar swirl, and record whether its extremal row
      closes the angular-momentum variant \eqref{eq:c3v72-angular} of
      the palinstrophy row.
\item The next compulsory companion audit is V75.
\end{enumerate}
\end{programme}

\begin{assessment}[Position after V73]
\label{ass:c3v73-position}
The vorticity ledger's three rows are read jointly; they are
governed by one uncomputed constant of parabolic regularity and
exclude nothing.  No full stability or global-regularity claim is
made.
\end{assessment}

% =============================================================================
\section*{Version V73 (third cycle) append-only record}
\addcontentsline{toc}{section}{Version V73 (third cycle) append-only record}
% =============================================================================

The complete V72 source of the third cycle was retained before this
continuation.  Its SHA--256 digest was
\begin{center}\scriptsize\ttfamily
d9e869279d0bfb5ad5b486026bce170643250990423aaa407461b22de260a4d0.
\end{center}
V73 read the vorticity ledger jointly against (E5) and tabulated its
rows.  No full regularity claim is made.

% =============================================================================
\section{V74: the vorticity-ledger block consolidated, and the mean-vorticity loop}
\label{sec:c3v74-entry}
% =============================================================================

Items (AC1) and (AC2) of Programme~\ref{prog:c3v74} are carried out.
(AC1): the block V71--V73 is consolidated as one theorem with its
negative record, and an in-place correction of V73 is recorded.
(AC2): the identity for the Gaussian mean vorticity is derived; it
has pure damping of one unit and a source that is the Gaussian
first moment of \(V\times\Omega\), vanishing on every columnar swirl.  Its
extremal row bounds the mean vorticity by the weighted enstrophy at
the maximizing configuration, and, combined with the palinstrophy
row of V72 through the Gaussian Poincar\'e inequality, it closes a
loop between two extremal configurations: the result is the
cycle's second constants-only constraint,
\(3\kappa^2+(1+2\kappa^2)\,q_{\mathcal P}\ge1\) wherever the palinstrophy row applies.
The direction for V76--V80 is the explicit reading of this
constraint and the two-configuration method.  No global regularity
claim is made.  The next compulsory companion audit is V75.

% =============================================================================
\section{Correction record}
\label{sec:c3v74-correction}
% =============================================================================

\begin{remark}[In-place correction of V73]
\label{rem:c3v74-correction}
V73 as first assembled stated that all derivative constants
\(C_\infty^{(k)}\) are dimensionless in the units of the profile equation.
Only \(C_\infty'\) is; \(C_\infty^{(k)}\) carries the factor \(\nu^{(1-k)/2}\), and
\(\nu^{1/2}C_\infty''\) is the dimensionless second-derivative constant.
Hypothesis (E5-small) was corrected in place to
\(C_\infty''\le\gamma_2\kappa\nu^{-1/2}\) with \(\gamma_2\) dimensionless, and the two
consequences in Proposition~\ref{prop:c3v73-thresholds} (the
substitution for \(C_\infty''\), and the ratio bound
\(\sup\mathcal P\le2(\gamma_2\kappa/\theta_{\mathcal P}^-)^2C/\nu\)) were corrected accordingly.
No threshold of V73 is affected: \(\gamma_2\) enters the sizes of the
bounds only, not the applicability conditions, which involve
\(\gamma_1\) alone.  The digest of V73 recorded below is that of the
corrected file.
\end{remark}

% =============================================================================
\section{The vorticity-ledger block}
\label{sec:c3v74-block}
% =============================================================================

\begin{theorem}[The vorticity-ledger block, V71--V73]
\label{thm:c3v74-block}
At a homogeneous type-I point, with the notation of V67, V71, V72 and
\(\Xi_1\) of Theorem~\ref{thm:c3v73-table}:
\begin{enumerate}[label=\textup{(\arabic*)},leftmargin=3.2em]
\item (Weighted enstrophy, V71.)
      \(\frac12\mathcal E_1'=-\mathcal E_{\nabla,1}-\frac32\mathcal E_1+3\nu\mathcal E+\mathcal S_1-\mathcal T_1+4\nu\mathcal T\); at an extremum
      of \(\mathcal E_1\) with \(\theta_1>0\),
      \(\sup\mathcal E_1\le(2\nu C\,\Xi_1+4\nu^2C_\infty''^2|G|)/\theta_1\): a concentration
      bound, not a constraint, because of the positive commutator.
\item (Palinstrophy, V72.)  \(\frac12\mathcal P'=-\mathcal P_\nabla-\frac32\mathcal P+\mathcal S_{\mathcal P}-\mathcal T_{\mathcal P}\); at an
      extremum of \(\mathcal P\) with \(\theta_{\mathcal P}>0\),
      \(\sup\mathcal P\le(C_\infty''/\theta_{\mathcal P})^2\mathcal E(W_{\mathcal P})\le2(C_\infty''/\theta_{\mathcal P})^2C\), and with
      \(q_{\mathcal P}=2\nu(C_\infty''/\theta_{\mathcal P})^2<1\),
      \(\sup\mathcal P\le(C_\infty''/\theta_{\mathcal P})^2(1-q_{\mathcal P})^{-1}|G||\bar\Omega(W_{\mathcal P})|^2\): a ratio
      bound.
\item (Joint reading, V73.)  Without hypothesis,
      \(C_\infty'\ge1-\frac{\sqrt3}2\kappa-\frac{\kappa^2}2\) for \(\kappa<\kappa_0\approx0.79\).  Under (E5-small),
      \(\kappa\ge\min\{\kappa_1(\gamma_1),\kappa_{\rm small}\}\) with \(\kappa_1\) explicit and bracketed; the
      three rows apply on a common nonempty window iff
      \(0.807<\gamma_1<1.094\); the floor \(\kappa_1\) exceeds \(\varepsilon_{\rm reg}\) iff
      \(\gamma_1<1/\varepsilon_{\rm reg}-\varepsilon_{\rm reg}/2-\sqrt3/2\), and exceeds \(\frac12\) iff \(\gamma_1<0.884\).
\item (Negative record.)  No row excludes a homogeneous type-I
      point for any values of \(\gamma_1,\gamma_2,\varepsilon_{\rm reg}\).
\end{enumerate}
\end{theorem}

\begin{proof}
Theorems~\ref{thm:c3v71-identity}, \ref{thm:c3v71-extremal},
\ref{thm:c3v72-identity}, \ref{thm:c3v72-extremal},
Proposition~\ref{prop:c3v72-poincare},
Proposition~\ref{prop:c3v73-gradient-floor},
Theorems~\ref{thm:c3v73-velocity-floor}, \ref{thm:c3v73-window},
Propositions~\ref{prop:c3v73-epsreg}, \ref{prop:c3v73-chasing},
Firewall~\ref{fw:c3v73-not-implied}.
\end{proof}

% =============================================================================
\section{The Gaussian mean vorticity}
\label{sec:c3v74-mean}
% =============================================================================

Recall \(\bar\Omega=\frac1{|G|}\int\Omega\,G=\frac1{2\nu|G|}\int z\times V\,G\)
(Proposition~\ref{prop:c3v72-poincare}), and put
\begin{equation}
 \mathcal B:=\frac1{2\nu|G|}\int z\times(V\times\Omega)\,G
 =\frac1{2\nu|G|}\int\bigl[V\,(z\cdot\Omega)-\Omega\,(z\cdot V)\bigr]G .
 \label{eq:c3v74-source}
\end{equation}

\begin{theorem}[Mean-vorticity identity]
\label{thm:c3v74-identity}
For every profile of the family and every \(\sigma\),
\begin{equation}
 \boxed{\quad
 \frac{\dd}{\dd\sigma}\bar\Omega=-\bar\Omega+\mathcal B .
 \quad}
 \label{eq:c3v74-identity}
\end{equation}
\end{theorem}

\begin{proof}
The vorticity equation of the profile is
\(\partial_\sigma\Omega=\nu\Delta\Omega-\frac12z\cdot\nabla\Omega-\Omega+\operatorname{curl}(V\times\Omega)\), since
\(\operatorname{curl}((z\cdot\nabla)V)=(z\cdot\nabla)\Omega+\Omega\) and
\(-(V\cdot\nabla)\Omega+(\Omega\cdot\nabla)V=\operatorname{curl}(V\times\Omega)\) for divergence-free fields.
Linear part: \(\nu\int\Omega\,\Delta G=\int\Omega\bigl(\frac{|z|^2}{4\nu}-\frac32\bigr)G\) and
\(-\frac12\int(z\cdot\nabla\Omega)G=\frac12\int\Omega\operatorname{div}(zG)=\int\Omega\bigl(\frac32-\frac{|z|^2}{4\nu}\bigr)G\), so the
linear part integrates to \(-\int\Omega G\).  Source:
\(\int\operatorname{curl}F\,G=-\int\nabla G\times F=\frac1{2\nu}\int z\times F\,G\) with \(F=V\times\Omega\), and
\(z\times(V\times\Omega)=V(z\cdot\Omega)-\Omega(z\cdot V)\).  Divide by \(|G|\).  All integrals
converge by the class bounds.
\end{proof}

\begin{remark}[Check against the velocity form, and on the columnar swirl]
\label{rem:c3v74-check}
From \(\bar\Omega=\frac1{2\nu|G|}\int z\times V\,G\) directly: for the linear part,
\(L_0=\nu\Delta-\frac12z\cdot\nabla\) is symmetric in \(L^2(G)\) and \(L_0(e\times z)=-\frac12e\times z\),
so \(\int(L_-V)\cdot(e\times z)G=-\int V\cdot(e\times z)G\), the damping of one unit;
for the nonlinear part, \(\int z\times\nabla P\,G=0\) (the integrand
\(\varepsilon_{ijk}z_j\partial_kP\,G\) integrates by parts to \(\varepsilon_{ijk}Pz_jz_k G/(2\nu)=0\)), and
\(-\int z\times(V\cdot\nabla)V\,G=-\int z\times(\Omega\times V)G=\int z\times(V\times\Omega)G\), the
gradient part of \((V\cdot\nabla)V\) vanishing as the pressure does.  The
two routes agree.  On the columnar swirl \(V=w(\sigma,r)e_\theta\) of
Proposition~\ref{prop:c3v68-tests}, \(V\times\Omega=\frac12\partial_z(w^2)e_z+\frac wr\partial_r(rw)e_r\)
is poloidal and \(z\times(V\times\Omega)\) is azimuthal, so \(\mathcal B=0\) and
\eqref{eq:c3v74-identity} reduces to \(\bar\Omega'=-\bar\Omega\), the exact decay of
the linear profile equation that the swirl solves (\(N\equiv0\)).
\end{remark}

\begin{lemma}[Bounds on the source]
\label{lem:c3v74-bounds}
For every profile of the family,
\begin{equation}
 |\mathcal B|\ \le\ \frac{\kappa}{2(\nu|G|)^{1/2}}\,\mathcal E_1^{1/2},\qquad
 \mathcal E_1\le12\nu\,\mathcal E+16\nu^2\,\mathcal P .
 \label{eq:c3v74-bounds}
\end{equation}
\end{lemma}

\begin{proof}
\(|z\times(V\times\Omega)|\le C_\infty|z||\Omega|\) and \(\int|z||\Omega|G\le\mathcal E_1^{1/2}|G|^{1/2}\) by
Cauchy--Schwarz; \(C_\infty|G|^{1/2}/(2\nu|G|)=\kappa/(2(\nu|G|)^{1/2})\).  The second
bound is the Gaussian Hardy inequality for the field \(\Omega\) (proof of
Lemma~\ref{lem:c3v71-bounds}).
\end{proof}

\begin{theorem}[Extremal mean vorticity]
\label{thm:c3v74-extremal}
Let \(W_{\bar\Omega}\) be a maximizer of \(|\bar\Omega|^2\) over \(\mathcal F(x_0)\).  Then
\begin{equation}
 \boxed{\quad
 \sup_{\mathcal F(x_0)}|G|\,|\bar\Omega|^2\ =\ |G|\,|\bar\Omega(W_{\bar\Omega})|^2\ \le\ \frac{\kappa^2}{4\nu}\mathcal E_1(W_{\bar\Omega})
 \ \le\ \kappa^2\bigl(3\,\mathcal E+4\nu\,\mathcal P\bigr)(W_{\bar\Omega}) .
 \quad}
 \label{eq:c3v74-row}
\end{equation}
\end{theorem}

\begin{proof}
\(|\bar\Omega|^2\) is continuous on the family and \(C^1\) in \(\sigma\); Lemma~\ref{lem:c3v34-extremal}
and \eqref{eq:c3v74-identity} give \(0=\bar\Omega\cdot\bar\Omega'=-|\bar\Omega|^2+\bar\Omega\cdot\mathcal B\) at
\(W_{\bar\Omega}\), \(\sigma=0\), hence \(|\bar\Omega|\le|\mathcal B|\) there; then \eqref{eq:c3v74-bounds}.
\end{proof}

% =============================================================================
\section{The two-configuration loop}
\label{sec:c3v74-loop}
% =============================================================================

\begin{theorem}[The mean-vorticity loop: a second constants-only constraint]
\label{thm:c3v74-loop}
Let \(x_0\) be a homogeneous type-I point at which the palinstrophy
row applies, \(\theta_{\mathcal P}=\frac32-2C_\infty'-\frac{\sqrt3}2\kappa-\frac{\kappa^2}2>0\), and let
\(q_{\mathcal P}=2\nu(C_\infty''/\theta_{\mathcal P})^2\).  Then
\begin{equation}
 \boxed{\quad
 3\kappa^2+(1+2\kappa^2)\,q_{\mathcal P}\ \ge\ 1 .
 \quad}
 \label{eq:c3v74-constraint}
\end{equation}
Equivalently: if \(q_{\mathcal P}<1\) and \(3\kappa^2<1\), then
\((1-3\kappa^2)(1-q_{\mathcal P})\le5\kappa^2q_{\mathcal P}\).
\end{theorem}

\begin{proof}
Put \(M:=\sup_{\mathcal F(x_0)}|G||\bar\Omega|^2\) and \(\Pi:=\sup_{\mathcal F(x_0)}\mathcal P\), both finite and
attained.  If \(q_{\mathcal P}\ge1\) or \(3\kappa^2\ge1\), \eqref{eq:c3v74-constraint} holds;
assume \(q_{\mathcal P}<1\) and \(3\kappa^2<1\).  By \eqref{eq:c3v74-row} and the Gaussian
Poincar\'e inequality \eqref{eq:c3v72-poincare} at \(W_{\bar\Omega}\),
\(M\le\kappa^2(3(2\nu\Pi+M)+4\nu\Pi)=\kappa^2(10\nu\Pi+3M)\), i.e.\
\((1-3\kappa^2)M\le10\kappa^2\nu\Pi\).  By \eqref{eq:c3v72-angular} at \(W_{\mathcal P}\),
\(\Pi\le\frac{(C_\infty''/\theta_{\mathcal P})^2}{1-q_{\mathcal P}}|G||\bar\Omega(W_{\mathcal P})|^2\le\frac{q_{\mathcal P}}{2\nu(1-q_{\mathcal P})}M\).
Hence \((1-3\kappa^2)M\le\frac{5\kappa^2q_{\mathcal P}}{1-q_{\mathcal P}}M\).  If
\((1-3\kappa^2)(1-q_{\mathcal P})>5\kappa^2q_{\mathcal P}\), then \(M=0\), so \(\bar\Omega\equiv0\) on the family,
so \(\Pi=0\), so \(\nabla\Omega\equiv0\) and \(\Omega\equiv\bar\Omega=0\) on the family; but \(\mathcal E>0\) on the
family (Theorem~\ref{thm:c3v67-extremal}: a divergence-free
irrotational field in \(L^6\) is zero, and \(\varphi\ge c_\Phi>0\)).  Therefore
\((1-3\kappa^2)(1-q_{\mathcal P})\le5\kappa^2q_{\mathcal P}\), which expands to
\(1\le3\kappa^2+q_{\mathcal P}+2\kappa^2q_{\mathcal P}\).
\end{proof}

\begin{assessment}[What the loop is]
\label{ass:c3v74-loop}
The mean row bounds the mean vorticity at its maximizer by the
enstrophy and the palinstrophy there; the palinstrophy row bounds
the palinstrophy at its maximizer by the enstrophy there; the
Gaussian Poincar\'e inequality bounds the enstrophy anywhere by the
palinstrophy and the mean vorticity.  The three together bound the
supremum of each of the two extremal observables by a multiple of
the supremum of the other, at two different configurations, and
when the product of the multiples is below one both suprema vanish,
which the class forbids.  This is a new mechanism in the series:
the extremal method so far produced one row per observable, each
evaluated at its own configuration; the loop links two rows through
a functional inequality that holds at every configuration, and the
sign structure (both rows are upper bounds on a positive quantity
by another positive quantity, and the closing inequality is an upper
bound on the third by the two) makes the composition close without
any knowledge of where the extremal configurations are.  The
constraint \eqref{eq:c3v74-constraint} is explicit, it involves the
second-derivative constant through \(q_{\mathcal P}\), and it holds wherever
the palinstrophy row applies; it is of \(\varepsilon\)-regularity type, like
\eqref{eq:c3v67-constraint}, since a small velocity constant with
(E5-small) makes both \(3\kappa^2\) and \(q_{\mathcal P}\) small; its explicit floor
\(\kappa_2(\gamma_1,\gamma_2)\) is the first task of V76.
\end{assessment}

\begin{firewall}[What the loop does not do]
\label{fw:c3v74-loop}
The constraint holds only where \(\theta_{\mathcal P}>0\); where the palinstrophy
row is void, the loop says nothing, and Theorem~\ref{thm:c3v73-window}
shows that under (E5-small) this happens at the floor \(\kappa_1\) when
\(\gamma_1\ge1.094\).  The alternative form ``\(\theta_{\mathcal P}\le0\) or
\eqref{eq:c3v74-constraint}'' is what holds at every homogeneous type-I
point.  It is not an exclusion: the class constants are not known
to violate it.
\end{firewall}

% =============================================================================
\section{Direction for V76--V80}
\label{sec:c3v74-direction}
% =============================================================================

\begin{assessment}[The two-configuration method]
\label{ass:c3v74-direction}
V76: the explicit reading of \eqref{eq:c3v74-constraint} under
(E5-small): \(q_{\mathcal P}\le2\gamma_2^2\kappa^2/\theta_{\mathcal P}^-(\kappa)^2\), so the constraint fails
for \(\kappa\) below an explicit \(\kappa_2(\gamma_1,\gamma_2)\), with the leading form
\(\kappa_2\approx(3+\frac89\gamma_2^2)^{-1/2}\) for small \(\gamma_2\); the comparison of \(\kappa_2\) with
\(\kappa_1(\gamma_1)\) and the table of the three floors (\(\varepsilon_{\rm reg}\), \(\kappa_1\), \(\kappa_2\)).
V77: the loop between the enstrophy row (V67) and the mean row
through Poincar\'e without the palinstrophy: at the extremum of
\(\mathcal E\), \(\frac12\mathcal E_\nabla+(1-C_\infty'-\dots)\mathcal E\le0\) bounds \(\mathcal P(W_{\mathcal E})\) by \(\mathcal E(W_{\mathcal E})\), and
Poincar\'e bounds \(\mathcal E\) by \(\mathcal P\) and \(|\bar\Omega|^2\); record whether a loop with
the mean row closes and what constraint it gives.  V78: the first
moments \(\int z_j\Omega_i\,G\) (the vorticity dipole): their identities have
damping \(\frac32\) and commutators with the weight; record whether any
closes into a loop with the rows above.  V79: consolidation of the
loop block; V80: audit.  The general form of the method is: a family
of observables \(F_1,\dots,F_n\) whose extremal rows give
\(\sup F_i\le\sum_jA_{ij}\sup F_j\) with an explicit nonnegative matrix \(A\); if
the spectral radius of \(A\) is below one every \(\sup F_i\) vanishes, and
the class's positivity statements turn this into a constraint
\(\rho(A)\ge1\).  The vorticity ledger gives such a matrix on
\((|G||\bar\Omega|^2,\mathcal P)\) with \(A=\bigl(\begin{smallmatrix}3\kappa^2&10\kappa^2\nu\\ q_{\mathcal P}/(2\nu(1-q_{\mathcal P}))&0\end{smallmatrix}\bigr)\),
whose spectral radius is below one exactly when
\eqref{eq:c3v74-constraint} fails; enlarging the family of observables
is the direction.  No global regularity claim is made.
\end{assessment}

% =============================================================================
\section{Integrated V74 conclusion and next acquisition order}
\label{sec:c3v74-conclusion}
% =============================================================================

\begin{theorem}[What V74 proves]
\label{thm:c3v74-conclusion}
Relative to the two legacy cycles and V1--V73: the consolidated
block (Theorem~\ref{thm:c3v74-block}), the mean-vorticity identity
(Theorem~\ref{thm:c3v74-identity}) with its checks
(Remark~\ref{rem:c3v74-check}) and bounds (Lemma~\ref{lem:c3v74-bounds}),
the extremal mean-vorticity row (Theorem~\ref{thm:c3v74-extremal}),
and the loop constraint (Theorem~\ref{thm:c3v74-loop}).  The full
Navier--Stokes stability/global-regularity theorem remains
unproved.
\end{theorem}

\begin{proof}
The cited results.
\end{proof}

\begin{programme}[V75 acquisition order]
\label{prog:c3v75}
\begin{enumerate}[label=\textup{(AC\arabic*)},leftmargin=4.8em]
\item The compulsory companion audit: read the current SM and YM
      notes and record their digests and decision.
\item Record the position after seventy-five versions; delivery.
\end{enumerate}
\end{programme}

\begin{assessment}[Position after V74]
\label{ass:c3v74-position}
The vorticity ledger has a second constants-only constraint, from a
loop of two extremal rows.  No full stability or global-regularity
claim is made.
\end{assessment}

% =============================================================================
\section*{Version V74 (third cycle) append-only record}
\addcontentsline{toc}{section}{Version V74 (third cycle) append-only record}
% =============================================================================

The complete V73 source of the third cycle (after the in-place
correction of Remark~\ref{rem:c3v74-correction}) was retained before
this continuation.  Its SHA--256 digest was
\begin{center}\scriptsize\ttfamily
1fe93b1bbaa561ecb9213b041ae781a803a0281bb97acd0135fd518f0905b006.
\end{center}
V74 consolidated the vorticity-ledger block, derived the
mean-vorticity identity and closed the two-configuration loop.  No
full regularity claim is made.

% =============================================================================
\section{V75: compulsory companion audit and the position after seventy-five versions}
\label{sec:c3v75-entry}
% =============================================================================

This is the seventy-fifth version of the third cycle.  The compulsory
review of the two companion programmes is performed first, and the
position after seventy-five versions is recorded with the plan for
the block V76--V80.  No global regularity claim is made.  The next
compulsory companion audit is V80.

% =============================================================================
\section{Compulsory V75 review of the current companion notes}
\label{sec:c3v75-companion-audit}
% =============================================================================

The current companion files in \texttt{latex\_notes} were inspected
read-only.  Both programmes have moved since the V70 audit: the SM
programme from V41 to V44 of its seventh series, the YM programme
from V19 to V21 of its sixth series.  The frozen checkpoints are
\begin{align}
 &\texttt{Renewal\_SM\_Einstein\_Continuum\_Research\_Notes\_V44.tex},
 \notag\\[-2pt]
 &\qquad
 \texttt{92c8b334e6a04b879b4b6db25af983c095f68453c0a6007d2a21ff589f7d36d3},
 \label{eq:c3v75-SM-checkpoint}\\
 &\texttt{Renewal\_Yang\_Mills\_Mass\_Gap\_Research\_Notes\_V21.tex},
 \notag\\[-2pt]
 &\qquad
 \texttt{4393ea3951991d7e1bcaa67210893033d54b9b272b3260a38d1f90a3989bf150}.
 \label{eq:c3v75-YM-checkpoint}
\end{align}
(SM V44 has 17679 lines; YM V21 has 9142 lines.  The folder also
holds frozen legacy files of both programmes, \texttt{\_f1} to
\texttt{\_f6}, not inspected.)

\emph{SM V42--V44.}  The SM programme replaced missing mesh data by
a uniform theorem on a normalized shape class (a universal
shape-regular face space with a continuous complete face
certificate and a universal face gap card, with a finite
alternative when the card fails), then an adiabatic face replacement
with its moment gate and smooth-continuum comparison, and a
conservative self-adjoint collar completion (finite normal fibre at
zero tangential momentum, gap preservation, first-moment control,
covariant curved lift).

\emph{YM V20--V21.}  The YM programme performed its own companion
audit (recording the NS file at V68 with the digest
\texttt{ba094ef2\dots}, which agrees with the V69 record of this
series) and then an exact tree-gauge disintegration: correction of
its target to an exact trace-shell law, incident-star
disintegration, common-incidence cancellation, removal of a spurious
Jacobian, a response-matched oscillation, and two negative results
(response matching does not imply its condition WCIC; perfect mixing
does not imply a positive static Hessian), with the canonical-root
static card restored as its next acquisition.

\begin{theorem}[V75 companion-audit decision]
\label{thm:c3v75-companion-decision}
Nothing is imported from SM V44 or YM V21.  Neither bears on the
third cycle's chapters, on the vorticity ledger or on the loop
method.
\end{theorem}

\begin{proof}
The SM results concern face certificates and collar completions of
lattice Dirac interfaces; the YM results concern gauge-fixed
conditional structures of a lattice gauge theory.  None is a
Navier--Stokes statement, and none concerns the ledger, the extremal
method, the vorticity identities or the two-configuration loop.  The
one point of contact is procedural: the YM audit's record of the NS
V68 digest agrees with this series' own record, a cross-check of the
append-only discipline and not a mathematical import.
\end{proof}

% =============================================================================
\section{Position after seventy-five versions, and the block V76--V80}
\label{sec:c3v75-position}
% =============================================================================

\begin{assessment}[Position after seventy-five versions of the third cycle]
\label{ass:c3v75-position-seventyfive}
Thirteen chapters are written: those listed in
Assessment~\ref{ass:c3v70-position-seventy} and the vorticity-ledger
block (V71--V74).  The cycle now has two relations between explicit
constants alone at every homogeneous type-I point:
\(C_\infty'+\frac{\sqrt3}2\kappa+\frac{\kappa^2}2\ge1\) (V67, one row of the vorticity ledger at
one extremal configuration) and, wherever the palinstrophy row
applies, \(3\kappa^2+(1+2\kappa^2)q_{\mathcal P}\ge1\) (V74, two rows at two extremal
configurations closed by the Gaussian Poincar\'e inequality).  Both
are of \(\varepsilon\)-regularity type; under (E5-small) both are floors on
the velocity constant, explicit in the parabolic constants \(\gamma_1,\gamma_2\)
that the series has not computed; neither is an exclusion.  What
the block changed is the method: a family of observables whose
extremal rows bound each supremum by the others gives a nonnegative
matrix whose spectral radius must be at least one, and the
vorticity ledger supplies such matrices with explicit entries.  The
block V76--V80 follows Assessment~\ref{ass:c3v74-direction}: the
explicit floor \(\kappa_2\) (V76), the enstrophy--mean loop (V77), the
vorticity dipole (V78), consolidation (V79), audit (V80).  Twenty-five
versions remain; the closing block V96--V100 is reserved as before.
No global regularity claim is made.
\end{assessment}

% =============================================================================
\section{Integrated V75 conclusion and next acquisition order}
\label{sec:c3v75-conclusion}
% =============================================================================

\begin{theorem}[What V75 establishes]
\label{thm:c3v75-conclusion}
Relative to the two legacy cycles and V1--V74: the companion-audit
decision (Theorem~\ref{thm:c3v75-companion-decision}) and the position.
The full Navier--Stokes stability/global-regularity theorem remains
unproved.
\end{theorem}

\begin{proof}
The cited records.
\end{proof}

\begin{programme}[V76 acquisition order]
\label{prog:c3v76}
\begin{enumerate}[label=\textup{(AC\arabic*)},leftmargin=4.8em]
\item Under (E5-small), write \eqref{eq:c3v74-constraint} with
      \(q_{\mathcal P}\le2\gamma_2^2\kappa^2/\theta_{\mathcal P}^-(\kappa)^2\) and derive the explicit floor
      \(\kappa_2(\gamma_1,\gamma_2)\): the least \(\kappa\) with
      \(3\kappa^2+(1+2\kappa^2)\cdot2\gamma_2^2\kappa^2/\theta_{\mathcal P}^-(\kappa)^2\ge1\) on \(\kappa<\kappa_r(\gamma_1)\); its
      leading form for small \(\gamma_2\); the comparison with \(\kappa_1(\gamma_1)\)
      (which is larger, for which \(\gamma_1,\gamma_2\)); the table of the three
      floors \(\varepsilon_{\rm reg},\kappa_1,\kappa_2\).
\item Record the general two-configuration lemma (the matrix
      form of Assessment~\ref{ass:c3v74-direction}) as a stated
      lemma with proof, for use in V77--V78.
\item The next compulsory companion audit is V80.
\end{enumerate}
\end{programme}

\begin{assessment}[Position after V75]
\label{ass:c3v75-position}
The companion audit is recorded; the loop block begins.  No full
stability or global-regularity claim is made.
\end{assessment}

% =============================================================================
\section*{Version V75 (third cycle) append-only record}
\addcontentsline{toc}{section}{Version V75 (third cycle) append-only record}
% =============================================================================

The complete V74 source of the third cycle was retained before this
continuation.  Its SHA--256 digest was
\begin{center}\scriptsize\ttfamily
d985c84d63e8917d464ce5e67ef6cb56f90d5b126db98e19c180c9cc4bc3d532.
\end{center}
V75 performed the compulsory companion audit and recorded the
position.  No full regularity claim is made.

% =============================================================================
\section{V76: the matrix lemma and the explicit loop floor}
\label{sec:c3v76-entry}
% =============================================================================

Items (AC1) and (AC2) of Programme~\ref{prog:c3v76} are carried out,
in the order (AC2) then (AC1).  The two-configuration argument of
V74 is stated as a lemma on nonnegative matrices of extremal-row
coefficients: if the suprema of nonnegative observables satisfy
\(s\le As\) with \(A\ge0\) and \(\rho(A)<1\), all suprema vanish; the loop
constraint of V74 is its instance.  Under (E5-small) the constraint
is read as an explicit floor \(\kappa_2(\gamma_1,\gamma_2)\) on the velocity constant,
bracketed on both sides, equal to \(\min\{1/\sqrt3,\kappa_r\}\) when the
second-derivative constant vanishes.  The comparison with the
enstrophy floor \(\kappa_1(\gamma_1)\) is explicit: the loop floor is the
stronger only for \(1/\sqrt3<\gamma_1<1.094\) and \(\gamma_2\) below an explicit
function of \(\gamma_1\) that never exceeds \(0.042\); the entry \(3\kappa^2\) of the
matrix, coming from the Hardy conversion of the mean row, caps the
loop floor at \(1/\sqrt3\).  The three floors are tabulated.  No global
regularity claim is made.  The next compulsory companion audit is
V80.

% =============================================================================
\section{The matrix lemma}
\label{sec:c3v76-matrix}
% =============================================================================

\begin{lemma}[Two-configuration lemma]
\label{lem:c3v76-matrix}
Let \(x_0\) be a homogeneous type-I point, \(F_1,\dots,F_n\) nonnegative
continuous observables on \(\mathcal F(x_0)\), \(s_i:=\sup_{\mathcal F(x_0)}F_i<\infty\), and
\(A\in\R^{n\times n}\) a matrix with nonnegative entries such that
\begin{equation}
 s_i\ \le\ \sum_{j=1}^nA_{ij}\,s_j\qquad(i=1,\dots,n).
 \label{eq:c3v76-system}
\end{equation}
If the spectral radius \(\rho(A)\) is below one, then \(s_1=\dots=s_n=0\), i.e.\
every \(F_i\) vanishes identically on the family.  Consequently, if some
\(F_i\) is positive somewhere on the family, \(\rho(A)\ge1\).
\end{lemma}

\begin{proof}
\eqref{eq:c3v76-system} says \(s\le As\) entrywise with \(s\ge0\).  Since \(A\ge0\)
entrywise, \(As\le A(As)=A^2s\), and by induction \(s\le A^ks\) for every \(k\).
If \(\rho(A)<1\), \(A^k\to0\), so \(s\le0\), hence \(s=0\).
\end{proof}

\begin{remark}[How the system \eqref{eq:c3v76-system} arises]
\label{rem:c3v76-arise}
Each row of \eqref{eq:c3v76-system} comes from the extremal row of
\(F_i\) (Lemma~\ref{lem:c3v34-extremal}: at a maximizer of \(F_i\) the
identity holds with zero left side and bounds \(F_i\) there by other
observables evaluated there), followed by functional inequalities
valid at every configuration (Hardy, Poincar\'e, the class bounds)
that convert those observables into multiples of the \(s_j\).  The
maximizers of the various \(F_i\) are different configurations, and
the lemma needs no information on where they are.  For \(n=2\),
\(A=\bigl(\begin{smallmatrix}\alpha&\beta\\ \gamma&\delta\end{smallmatrix}\bigr)\ge0\) has
\(\rho(A)<1\) if and only if \(\alpha<1\), \(\delta<1\) and \((1-\alpha)(1-\delta)>\beta\gamma\).
\end{remark}

\begin{corollary}[V74 recovered]
\label{cor:c3v76-v74}
With \(F_1=|G||\bar\Omega|^2\), \(F_2=\mathcal P\) and, wherever \(\theta_{\mathcal P}>0\) and \(q_{\mathcal P}<1\),
\(A=\bigl(\begin{smallmatrix}3\kappa^2&10\kappa^2\nu\\ q_{\mathcal P}/(2\nu(1-q_{\mathcal P}))&0\end{smallmatrix}\bigr)\),
the system \eqref{eq:c3v76-system} holds (proof of
Theorem~\ref{thm:c3v74-loop}), \(F_2\) is positive somewhere (else \(\Omega\) is
constant, hence zero, on the family), and \(\rho(A)\ge1\) is
\((1-3\kappa^2)(1-q_{\mathcal P})\le5\kappa^2q_{\mathcal P}\), i.e.\ \eqref{eq:c3v74-constraint}.
\end{corollary}

\begin{proof}
Remark~\ref{rem:c3v76-arise} with \(\alpha=3\kappa^2\), \(\beta=10\kappa^2\nu\),
\(\gamma=q_{\mathcal P}/(2\nu(1-q_{\mathcal P}))\), \(\delta=0\): \(\beta\gamma=5\kappa^2q_{\mathcal P}/(1-q_{\mathcal P})\).
\end{proof}

% =============================================================================
\section{The explicit loop floor}
\label{sec:c3v76-floor}
% =============================================================================

Under (E5-small) and for \(\kappa\le\kappa_{\rm small}\), \(\theta_{\mathcal P}\ge\theta_{\mathcal P}^-(\kappa)\) and
\(q_{\mathcal P}\le q^-(\kappa):=2\gamma_2^2\kappa^2/\theta_{\mathcal P}^-(\kappa)^2\) on \(\kappa<\kappa_r(\gamma_1)\)
(Proposition~\ref{prop:c3v73-thresholds}); put
\begin{equation}
 h(\kappa):=3\kappa^2+(1+2\kappa^2)\,\frac{2\gamma_2^2\kappa^2}{\theta_{\mathcal P}^-(\kappa)^2},\qquad
 \kappa_2(\gamma_1,\gamma_2):=\inf\{\kappa\in(0,\kappa_r):h(\kappa)\ge1\}\quad(\inf\emptyset:=\kappa_r).
 \label{eq:c3v76-h}
\end{equation}
\(h\) is increasing on \([0,\kappa_r)\) with \(h(0)=0\), so \(h<1\) on \([0,\kappa_2)\) and
\(h\ge1\) on \([\kappa_2,\kappa_r)\).

\begin{theorem}[Loop floor under (E5-small)]
\label{thm:c3v76-floor}
Under (E5-small), every homogeneous type-I point satisfies
\(\kappa\ge\min\{\kappa_2(\gamma_1,\gamma_2),\kappa_{\rm small}\}\).  Moreover, with
\(\kappa_2^+:=(3+\frac89\gamma_2^2)^{-1/2}\),
\begin{equation}
 \kappa_2\ \le\ \min\{\kappa_2^+,\kappa_r\},\qquad
 \kappa_2\ \ge\ \kappa_2^-:=\Bigl(3+\frac{10}3\,\frac{\gamma_2^2}{\theta_{\mathcal P}^-(\kappa_2^+)^2}\Bigr)^{-1/2}
 \quad\text{if }\kappa_2^+<\kappa_r ,
 \label{eq:c3v76-bracket}
\end{equation}
and \(\kappa_2=\min\{1/\sqrt3,\kappa_r\}\) when \(\gamma_2=0\).
\end{theorem}

\begin{proof}
If \(\kappa\le\kappa_{\rm small}\) and \(\kappa<\kappa_r\), then \(\theta_{\mathcal P}>0\) and
\eqref{eq:c3v74-constraint} with \(q_{\mathcal P}\le q^-(\kappa)\) gives \(h(\kappa)\ge1\), so
\(\kappa\ge\kappa_2\); if \(\kappa\ge\kappa_r\), then \(\kappa\ge\kappa_2\) as well.  Upper bound:
\(h(\kappa)\ge3\kappa^2+\frac89\gamma_2^2\kappa^2\) since \(1+2\kappa^2\ge1\) and \(\theta_{\mathcal P}^-\le\frac32\), so
\(h(\kappa_2^+)\ge1\) if \(\kappa_2^+<\kappa_r\).  Lower bound: for \(\kappa\le\kappa_2^+<\kappa_r\),
\(1+2\kappa^2\le1+2(\kappa_2^+)^2\le\frac53\) (as \(\kappa_2^+\le1/\sqrt3\)) and
\(\theta_{\mathcal P}^-(\kappa)\ge\theta_{\mathcal P}^-(\kappa_2^+)>0\), so \(h(\kappa)\le\kappa^2(3+\frac{10}3\gamma_2^2/\theta_{\mathcal P}^-(\kappa_2^+)^2)<1\)
for \(\kappa<\kappa_2^-\); and \(\kappa_2^-\le\kappa_2^+\) since \(\frac{10}{3\theta^2}\ge\frac{40}{27}>\frac89\) for
\(\theta\le\frac32\).  For \(\gamma_2=0\), \(h=3\kappa^2\).
\end{proof}

\begin{remark}[Numerical values]
\label{rem:c3v76-values}
\(\kappa_2^+\) is \(0.577\), \(0.507\), \(0.302\) for \(\gamma_2=0,1,3\), independently of
\(\gamma_1\); for \((\gamma_1,\gamma_2)=(0.7,0.5)\) the bracket is \(0.089\le\kappa_2\le0.557\)
with \(\kappa_2=0.412\); for \((\gamma_1,\gamma_2)=(1,\gamma_2)\), \(\kappa_r=0.483<\kappa_2^+\) for
\(\gamma_2\le1\), and \(\kappa_2\) is \(0.483\), \(0.447\), \(0.309\) for \(\gamma_2=0,0.1,1\).  The
lower bound \eqref{eq:c3v76-bracket} is weak when \(\kappa_2^+\) is close to
\(\kappa_r\); the exact value is the root of \(h=1\), a quartic equation in
\(\kappa^2\) after clearing \(\theta_{\mathcal P}^-(\kappa)^2\), and is given by the bisection
recorded in the working notes.
\end{remark}

% =============================================================================
\section{Comparison of the floors}
\label{sec:c3v76-comparison}
% =============================================================================

\begin{theorem}[Loop floor against enstrophy floor]
\label{thm:c3v76-comparison}
Under (E5-small):
\begin{enumerate}[label=\textup{(\arabic*)},leftmargin=3.2em]
\item If \(\gamma_1\ge\frac{\sqrt3+\sqrt7}4\) (\(\approx1.094\)), then \(\kappa_2\le\kappa_r\le\kappa_1\).
\item If \(\gamma_1<\frac{\sqrt3+\sqrt7}4\), then \(\kappa_2>\kappa_1\) if and only if \(h(\kappa_1)<1\), if
      and only if \(\gamma_1>1/\sqrt3\) and
      \begin{equation}
       \gamma_2\ <\ \gamma_2^*(\gamma_1):=\Bigl(\frac12-\gamma_1\kappa_1\Bigr)
       \Bigl(\frac{1-3\kappa_1^2}{2\kappa_1^2(1+2\kappa_1^2)}\Bigr)^{1/2} .
       \label{eq:c3v76-gamma2star}
      \end{equation}
\item \(\gamma_2^*(\gamma_1)<0.042\) on \((1/\sqrt3,1.094)\), with values \(0.023\), \(0.041\),
      \(0.041\), \(0.017\), \(0.001\) at \(\gamma_1=0.6,0.7,0.8,1,1.09\).
\end{enumerate}
\end{theorem}

\begin{proof}
(1) Theorem~\ref{thm:c3v73-window}(1): \(\kappa_1<\kappa_r\) iff \(\gamma_1<(\sqrt3+\sqrt7)/4\).
(2) With \(\kappa_1<\kappa_r\) and \(h\) increasing, \(\kappa_2>\kappa_1\) iff \(h(\kappa_1)<1\).  At
\(\kappa_1\), \(\theta_{\mathcal P}^-(\kappa_1)=\frac12-\gamma_1\kappa_1\) (proof of Theorem~\ref{thm:c3v73-window}),
so \(h(\kappa_1)<1\) iff \(3\kappa_1^2<1\) and
\(2\gamma_2^2\kappa_1^2(1+2\kappa_1^2)<(1-3\kappa_1^2)(\frac12-\gamma_1\kappa_1)^2\); and \(\kappa_1<1/\sqrt3\) iff
\(f(1/\sqrt3)>1\) iff \(a/\sqrt3+\frac16>1\) iff \(a>\frac{5\sqrt3}6\) iff \(\gamma_1>\frac{\sqrt3}3\).
(3) Evaluation of \eqref{eq:c3v76-gamma2star}.
\end{proof}

\begin{theorem}[The three floors]
\label{thm:c3v76-table}
At a homogeneous type-I point, under (E5-small) with
\(\kappa_{\rm small}\ge\max\{\kappa_1,\kappa_2\}\), \(\kappa\ge\max\{\varepsilon_{\rm reg},\kappa_1(\gamma_1),\kappa_2(\gamma_1,\gamma_2)\}\), where:

\medskip
\begin{center}\small
\begin{tabular}{@{}p{0.12\textwidth}p{0.3\textwidth}p{0.5\textwidth}@{}}
\toprule
\textbf{Floor} & \textbf{Source} & \textbf{Explicit form}\\
\midrule
\(\varepsilon_{\rm reg}\) & (E5), CKN \(\varepsilon\)-regularity (Firewall C2-V36-AC2) & universal; not computed in the series\\[3pt]
\(\kappa_1(\gamma_1)\) & enstrophy row at one configuration (V67, V73) & \(-a+(a^2+2)^{1/2}\), \(a=\gamma_1+\frac{\sqrt3}2\); \(\frac2{2a+\sqrt2}\le\kappa_1<\frac1a\)\\[3pt]
\(\kappa_2(\gamma_1,\gamma_2)\) & mean--palinstrophy loop (V74, V76) & root of \(h=1\) on \((0,\kappa_r)\), or \(\kappa_r\); \(\kappa_2^-\le\kappa_2\le\min\{\kappa_2^+,\kappa_r\}\)\\
\bottomrule
\end{tabular}
\end{center}
\medskip

\noindent
The operative floor is \(\kappa_1\) unless \(1/\sqrt3<\gamma_1<1.094\) and \(\gamma_2<\gamma_2^*(\gamma_1)\),
or \(\varepsilon_{\rm reg}\) exceeds both.
\end{theorem}

\begin{proof}
Firewall C2-V36-AC2, Theorems~\ref{thm:c3v73-velocity-floor},
\ref{thm:c3v76-floor}, \ref{thm:c3v76-comparison}.
\end{proof}

\begin{assessment}[Item (AC1): what caps the loop floor]
\label{ass:c3v76-cap}
The loop floor never exceeds \(1/\sqrt3\), whatever the derivative
constants, because of the entry \(3\kappa^2\) of the matrix: the mean row
bounds \(|G||\bar\Omega|^2\) by \(\frac{\kappa^2}{4\nu}\mathcal E_1\), and the Hardy conversion
\(\mathcal E_1\le12\nu\mathcal E+16\nu^2\mathcal P\) with Poincar\'e returns \(3\kappa^2\) of \(|G||\bar\Omega|^2\)
itself.  The enstrophy floor has no such cap (it reaches \(\kappa_0=0.79\)
as \(\gamma_1\to0\)), which is why it dominates except in the sliver of
Theorem~\ref{thm:c3v76-comparison}(2).  A loop that improves on the
enstrophy floor must therefore use the enstrophy row itself as one
of its rows rather than converting the enstrophy through Poincar\'e:
at the maximizer of \(\mathcal E\), \eqref{eq:c3v67-row} bounds the palinstrophy
there by \(\frac2\nu\delta_{67}\mathcal E\) with \(\delta_{67}:=C_\infty'+\frac{\sqrt3}2\kappa+\frac{\kappa^2}2-1\ge0\), which is
small exactly when the enstrophy constraint is nearly saturated;
this is the enstrophy--mean loop of V77.
\end{assessment}

% =============================================================================
\section{Integrated V76 conclusion and next acquisition order}
\label{sec:c3v76-conclusion}
% =============================================================================

\begin{theorem}[What V76 proves]
\label{thm:c3v76-conclusion}
Relative to the two legacy cycles and V1--V75: the matrix lemma
(Lemma~\ref{lem:c3v76-matrix}) with V74 as its instance
(Corollary~\ref{cor:c3v76-v74}), the explicit loop floor with its
bracket (Theorem~\ref{thm:c3v76-floor}), the comparison
(Theorem~\ref{thm:c3v76-comparison}) and the table of floors
(Theorem~\ref{thm:c3v76-table}).  The full Navier--Stokes
stability/global-regularity theorem remains unproved.
\end{theorem}

\begin{proof}
The cited results.
\end{proof}

\begin{programme}[V77 acquisition order]
\label{prog:c3v77}
\begin{enumerate}[label=\textup{(AC\arabic*)},leftmargin=4.8em]
\item The enstrophy--mean loop: observables \(F_1=\mathcal E\), \(F_2=|G||\bar\Omega|^2\),
      \(F_3=\mathcal P\); rows: at the maximizer of \(\mathcal E\), \eqref{eq:c3v67-row} gives
      \(\nu\mathcal P\le2\delta_{67}\mathcal E\), then Poincar\'e gives \((1-4\delta_{67})s_1\le s_2\); at
      the maximizer of \(|\bar\Omega|^2\), \eqref{eq:c3v74-row} gives \(s_2\le3\kappa^2s_1+4\kappa^2\nu s_3\);
      at the maximizer of \(\mathcal P\), \eqref{eq:c3v72-row} gives \(s_3\le\frac{q_{\mathcal P}}{2\nu}s_1\).
      Apply Lemma~\ref{lem:c3v76-matrix}; write the constraint
      \(\rho(A)\ge1\) explicitly and verify it dominates
      \eqref{eq:c3v67-constraint} and \eqref{eq:c3v74-constraint}.
\item Its floor under (E5-small), with the leading form, against
      \(\kappa_1\).
\item The next compulsory companion audit is V80.
\end{enumerate}
\end{programme}

\begin{assessment}[Position after V76]
\label{ass:c3v76-position}
The loop method is a lemma; its first instance gives a floor capped
at \(1/\sqrt3\), weaker than the enstrophy floor except in a sliver.  No
full stability or global-regularity claim is made.
\end{assessment}

% =============================================================================
\section*{Version V76 (third cycle) append-only record}
\addcontentsline{toc}{section}{Version V76 (third cycle) append-only record}
% =============================================================================

The complete V75 source of the third cycle was retained before this
continuation.  Its SHA--256 digest was
\begin{center}\scriptsize\ttfamily
654135f3d7c2203cb486a38c910dfbdb51f2e5efbec9adb0e38cc19a797e739b.
\end{center}
V76 stated the matrix lemma and read the loop constraint as an
explicit floor.  No full regularity claim is made.

% =============================================================================
\section{V77: the enstrophy--mean loop and its four-observable variant}
\label{sec:c3v77-entry}
% =============================================================================

Items (AC1) and (AC2) of Programme~\ref{prog:c3v77} are carried out.
The loop with the enstrophy as one of its rows gives the third
constants-only constraint of the cycle,
\(4\delta_{67}+3\kappa^2+2\kappa^2q_{\mathcal P}\ge1\) wherever the palinstrophy row applies,
with \(\delta_{67}=C_\infty'+\frac{\sqrt3}2\kappa+\frac{\kappa^2}2-1\ge0\) the slack of the V67
constraint; it dominates V67 exactly where \(3\kappa^2+2\kappa^2q_{\mathcal P}<1\) and is
incomparable with V74.  Under (E5-small) its floor exceeds \(\kappa_1\)
only for \(1/\sqrt3<\gamma_1<1.094\) and \(\gamma_2\) below an explicit function
bounded by \(0.069\).  A four-observable variant replaces the Hardy
conversion of the mean row by the concentration row of V71; it gives
a fourth constraint, \(4\delta_{67}+\kappa^2(\Xi_1+q_{\mathcal P})/(4\theta_1)\ge1\), whose floor
never exceeds \(\kappa_1\) in its window of applicability, because the
concentration threshold \(\theta_1\) is at most \(0.072\) at \(\kappa_1\) wherever
the window is nonempty.  The block's negative record is stated: the
loops of the vorticity ledger with the rows written so far cannot
raise the floor above \(\max\{\kappa_1,1/\sqrt3\}\).  No global regularity claim
is made.  The next compulsory companion audit is V80.

% =============================================================================
\section{The enstrophy--mean loop}
\label{sec:c3v77-loop}
% =============================================================================

Put \(\delta_{67}:=C_\infty'+\frac{\sqrt3}2\kappa+\frac{\kappa^2}2-1\), nonnegative at every homogeneous
type-I point by \eqref{eq:c3v67-constraint}, and
\(s_1:=\sup\mathcal E\), \(s_2:=\sup|G||\bar\Omega|^2\), \(s_3:=\sup\mathcal P\) over \(\mathcal F(x_0)\).

\begin{lemma}[The three rows]
\label{lem:c3v77-rows}
At a homogeneous type-I point with \(\theta_{\mathcal P}>0\),
\begin{equation}
 s_1\le4\delta_{67}\,s_1+s_2,\qquad
 s_2\le3\kappa^2s_1+4\kappa^2\nu\,s_3,\qquad
 s_3\le\frac{q_{\mathcal P}}{2\nu}\,s_1 .
 \label{eq:c3v77-rows}
\end{equation}
\end{lemma}

\begin{proof}
First row: at a maximizer \(W_{\mathcal E}\) of \(\mathcal E\), \eqref{eq:c3v67-row} reads
\(\frac12\mathcal E_\nabla\le\delta_{67}\mathcal E\), i.e.\ \(\nu\mathcal P(W_{\mathcal E})\le2\delta_{67}s_1\); the Gaussian
Poincar\'e inequality \eqref{eq:c3v72-poincare} at \(W_{\mathcal E}\) gives
\(s_1=\mathcal E(W_{\mathcal E})\le2\nu\mathcal P(W_{\mathcal E})+|G||\bar\Omega(W_{\mathcal E})|^2\le4\delta_{67}s_1+s_2\).  Second row:
\eqref{eq:c3v74-row} at the maximizer of \(|\bar\Omega|^2\), with \(\mathcal E\le s_1\) and
\(\mathcal P\le s_3\) there.  Third row: \eqref{eq:c3v72-row} at the maximizer of
\(\mathcal P\), \(\sup\mathcal P\le(C_\infty''/\theta_{\mathcal P})^2\mathcal E(W_{\mathcal P})\le\frac{q_{\mathcal P}}{2\nu}s_1\).
\end{proof}

\begin{theorem}[Enstrophy--mean loop: a third constants-only constraint]
\label{thm:c3v77-loop}
At every homogeneous type-I point with \(\theta_{\mathcal P}>0\),
\begin{equation}
 \boxed{\quad
 4\delta_{67}+3\kappa^2+2\kappa^2q_{\mathcal P}\ \ge\ 1,\qquad\text{i.e.}\qquad
 4C_\infty'+2\sqrt3\,\kappa+5\kappa^2+2\kappa^2q_{\mathcal P}\ \ge\ 5 .
 \quad}
 \label{eq:c3v77-constraint}
\end{equation}
\end{theorem}

\begin{proof}
\eqref{eq:c3v77-rows} is \(s\le As\) with
\(A=\begin{pmatrix}4\delta_{67}&1&0\\ 3\kappa^2&0&4\kappa^2\nu\\ q_{\mathcal P}/(2\nu)&0&0\end{pmatrix}\ge0\).
Its characteristic polynomial is
\(p(\lambda)=\lambda^3-4\delta_{67}\lambda^2-3\kappa^2\lambda-2\kappa^2q_{\mathcal P}\): expanding along the first
row, \((\lambda-4\delta_{67})\lambda^2+1\cdot(-3\kappa^2\lambda-4\kappa^2\nu\cdot q_{\mathcal P}/(2\nu))\).  By Descartes'
rule (signs \(+,-,-,-\) or \(+,0,-,-\)), \(p\) has exactly one positive root when
\(q_{\mathcal P}>0\), and \(p(\lambda)=\lambda(\lambda^2-4\delta_{67}\lambda-3\kappa^2)\) has one positive root when
\(q_{\mathcal P}=0\); the Perron root \(\rho(A)\) is that root, and \(p<0\) on \((0,\rho(A))\),
\(p>0\) on \((\rho(A),\infty)\).  Hence \(\rho(A)<1\) iff \(p(1)>0\) iff
\(1-4\delta_{67}-3\kappa^2-2\kappa^2q_{\mathcal P}>0\).  \(F_1=\mathcal E\) is positive on the family, so
Lemma~\ref{lem:c3v76-matrix} gives \(\rho(A)\ge1\).  (Directly: substituting
the third row into the second and the second into the first gives
\(s_1\le(4\delta_{67}+3\kappa^2+2\kappa^2q_{\mathcal P})s_1\) with \(s_1>0\).)
\end{proof}

\begin{proposition}[Comparison with V67 and V74]
\label{prop:c3v77-comparison}
\begin{enumerate}[label=\textup{(\arabic*)},leftmargin=3.2em]
\item \eqref{eq:c3v77-constraint} implies \eqref{eq:c3v67-constraint} and is
      strictly stronger exactly where \(3\kappa^2+2\kappa^2q_{\mathcal P}<1\): there it
      says \(\delta_{67}\ge\frac14(1-3\kappa^2-2\kappa^2q_{\mathcal P})>0\), a positive slack in the
      enstrophy constraint.
\item Neither of \eqref{eq:c3v77-constraint} and \eqref{eq:c3v74-constraint}
      implies the other: the first has \(4\delta_{67}\) where the second has
      \(q_{\mathcal P}\), and both \(4\delta_{67}<q_{\mathcal P}\) and \(4\delta_{67}>q_{\mathcal P}\) are compatible with
      the class.
\item Wherever \(\theta_{\mathcal P}>0\), the conjunction of the two is
      \(4\delta_{67}+3\kappa^2+2\kappa^2q_{\mathcal P}\ge1\) and \(3\kappa^2+(1+2\kappa^2)q_{\mathcal P}\ge1\), i.e.\
      \(3\kappa^2+2\kappa^2q_{\mathcal P}\ge1-\min\{4\delta_{67},q_{\mathcal P}\}\).
\end{enumerate}
\end{proposition}

\begin{proof}
(1) \(\delta_{67}\ge0\) is \eqref{eq:c3v67-constraint}; rearrange.  (2) The
series bounds neither \(\delta_{67}\) nor \(q_{\mathcal P}\) by the other.  (3) Rearrange.
\end{proof}

% =============================================================================
\section{The floor of the enstrophy--mean loop}
\label{sec:c3v77-floor}
% =============================================================================

Under (E5-small), \(\delta_{67}\le\delta^-(\kappa):=\gamma_1\kappa+\frac{\sqrt3}2\kappa+\frac{\kappa^2}2-1\) and
\(q_{\mathcal P}\le q^-(\kappa)\) (V76); put
\begin{equation}
 g_3(\kappa):=4\delta^-(\kappa)+3\kappa^2+2\kappa^2q^-(\kappa),\qquad
 \kappa_3(\gamma_1,\gamma_2):=\inf\{\kappa\in(0,\kappa_r):g_3(\kappa)\ge1\}\quad(\inf\emptyset:=\kappa_r),
 \label{eq:c3v77-g3}
\end{equation}
\(g_3\) increasing on \([0,\kappa_r)\).

\begin{theorem}[Floor of the enstrophy--mean loop]
\label{thm:c3v77-floor}
Under (E5-small), every homogeneous type-I point satisfies
\(\kappa\ge\min\{\kappa_3,\kappa_{\rm small}\}\), and:
\begin{enumerate}[label=\textup{(\arabic*)},leftmargin=3.2em]
\item \(\kappa_3\le\min\{\max\{\kappa_1,1/\sqrt3\},\kappa_r\}\).
\item For \(\gamma_1<\frac{\sqrt3+\sqrt7}4\): \(\kappa_3>\kappa_1\) iff \(g_3(\kappa_1)<1\) iff \(\gamma_1>1/\sqrt3\) and
      \begin{equation}
       \gamma_2<\gamma_3^*(\gamma_1):=\Bigl(\frac12-\gamma_1\kappa_1\Bigr)\frac{(1-3\kappa_1^2)^{1/2}}{2\kappa_1^2}
       =\gamma_2^*(\gamma_1)\Bigl(\frac{1+2\kappa_1^2}{2\kappa_1^2}\Bigr)^{1/2} ;
       \label{eq:c3v77-gamma3star}
      \end{equation}
      \(\gamma_3^*\) takes the values \(0.036\), \(0.067\), \(0.069\), \(0.055\), \(0.031\), \(0.002\)
      at \(\gamma_1=0.6,0.7,0.8,0.9,1,1.09\), and is below \(0.07\) throughout.
\item For \(\gamma_1\ge\frac{\sqrt3+\sqrt7}4\): \(\kappa_3\le\kappa_r\le\kappa_1\).
\end{enumerate}
\end{theorem}

\begin{proof}
The floor as in Theorem~\ref{thm:c3v76-floor}, from
\eqref{eq:c3v77-constraint}.  (1) For \(\kappa\ge\kappa_1\), \(\delta^-(\kappa)\ge0\), so
\(g_3(\kappa)\ge3\kappa^2\ge1\) once \(\kappa\ge1/\sqrt3\).  (2) \(g_3(\kappa_1)=3\kappa_1^2+2\kappa_1^2q^-(\kappa_1)\)
since \(\delta^-(\kappa_1)=0\), with \(q^-(\kappa_1)=2\gamma_2^2\kappa_1^2/(\frac12-\gamma_1\kappa_1)^2\); \(g_3(\kappa_1)<1\)
iff \(3\kappa_1^2<1\) and \(4\gamma_2^2\kappa_1^4<(1-3\kappa_1^2)(\frac12-\gamma_1\kappa_1)^2\); \(\kappa_1<1/\sqrt3\) iff
\(\gamma_1>1/\sqrt3\) (proof of Theorem~\ref{thm:c3v76-comparison}).  The
identity with \(\gamma_2^*\) is \eqref{eq:c3v76-gamma2star}.  (3)
Theorem~\ref{thm:c3v73-window}(1).
\end{proof}

\begin{remark}[Values]
\label{rem:c3v77-values}
For \((\gamma_1,\gamma_2)=(0.7,0)\): \(\kappa_3=0.554\) against \(\kappa_1=0.544\); for
\((1,0)\): \(\kappa_3=\kappa_r=0.483\) against \(\kappa_1=0.475\) (the loop empties the
whole window of the palinstrophy row); for \((0.7,0.1)\): \(\kappa_3=0.538<\kappa_1\);
for \((0.3,\gamma_2)\): \(\kappa_3\le0.637<\kappa_1=0.667\).  The leading form
\(\kappa\approx5/(4\gamma_1+2\sqrt3)\) of \eqref{eq:c3v77-constraint} for small \(\kappa\) is not
attained, the quadratic terms being of order one at the floor.
\end{remark}

% =============================================================================
\section{The four-observable variant}
\label{sec:c3v77-four}
% =============================================================================

The entry \(3\kappa^2\) of the loops comes from the Hardy conversion
\(\mathcal E_1\le12\nu\mathcal E+16\nu^2\mathcal P\) of the mean row.  The concentration row of V71
bounds \(\sup\mathcal E_1\) instead, at the cost of the threshold \(\theta_1\).

\begin{theorem}[The four-observable loop]
\label{thm:c3v77-four}
At every homogeneous type-I point with \(\theta_1>0\) and \(\theta_{\mathcal P}>0\),
with \(\Xi_1=3+\frac\kappa{\sqrt2}+2\sqrt3\kappa+2\kappa^2\),
\begin{equation}
 \boxed{\quad
 4\delta_{67}+\frac{\kappa^2\,(\Xi_1+q_{\mathcal P})}{4\theta_1}\ \ge\ 1 .
 \quad}
 \label{eq:c3v77-four}
\end{equation}
\end{theorem}

\begin{proof}
With \(s_4:=\sup\mathcal E_1\): the first form of \eqref{eq:c3v74-row} gives
\(s_2\le\frac{\kappa^2}{4\nu}s_4\); \eqref{eq:c3v71-row} at the maximizer of \(\mathcal E_1\) gives
\(\theta_1s_4\le\Xi_1\nu\,s_1+2\nu\mathcal E_\nabla\le\Xi_1\nu\,s_1+2\nu^2s_3\); with the first and
third rows of \eqref{eq:c3v77-rows},
\(s_2\le\frac{\kappa^2}{4\nu}\cdot\frac1{\theta_1}\bigl(\Xi_1\nu+2\nu^2\frac{q_{\mathcal P}}{2\nu}\bigr)s_1=\frac{\kappa^2(\Xi_1+q_{\mathcal P})}{4\theta_1}s_1\)
and \(s_1\le\bigl(4\delta_{67}+\frac{\kappa^2(\Xi_1+q_{\mathcal P})}{4\theta_1}\bigr)s_1\), \(s_1>0\).
\end{proof}

\begin{proposition}[The four-observable floor never exceeds \(\kappa_1\)]
\label{prop:c3v77-four-floor}
Under (E5-small), let \(\kappa_4\) be the floor of \eqref{eq:c3v77-four} on
\((0,\min\{\kappa_c,\kappa_r\})\), defined as in \eqref{eq:c3v77-g3} with
\(g_4(\kappa):=4\delta^-(\kappa)+\kappa^2(\Xi_1(\kappa)+q^-(\kappa))/(4\theta_1^-(\kappa))\).  Then \(\kappa_4\le\kappa_1\) for
every \(\gamma_1,\gamma_2\).
\end{proposition}

\begin{proof}
If \(\kappa_1\ge\min\{\kappa_c,\kappa_r\}\) there is nothing to prove.  Otherwise
\(\theta_1^-(\kappa_1)>0\) and \(\theta_{\mathcal P}^-(\kappa_1)>0\), which by Theorem~\ref{thm:c3v73-window}
means \(0.807<\gamma_1<1.094\), and there \(\theta_1^-(\kappa_1)=\frac12-\frac{\kappa_1}{\sqrt2}-\frac{\kappa_1^2}2\) with
\(\kappa_1\in(0.4577,0.5176)\), so \(\theta_1^-(\kappa_1)<0.072\) and \(\kappa_1^2\Xi_1(\kappa_1)>0.209\cdot5.3\);
hence \(g_4(\kappa_1)\ge\kappa_1^2\Xi_1(\kappa_1)/(4\theta_1^-(\kappa_1))>3.8>1\), so \(\kappa_4\le\kappa_1\).
\end{proof}

\begin{firewall}[Negative record of the loop block]
\label{fw:c3v77-negative}
With the rows written so far (enstrophy, weighted enstrophy,
palinstrophy, mean vorticity) and the functional inequalities used
(Hardy, Poincar\'e), every loop of the vorticity ledger gives a floor
at most \(\max\{\kappa_1,1/\sqrt3\}\) under (E5-small), and exceeds \(\kappa_1\) only in
the slivers of Theorems~\ref{thm:c3v76-comparison} and
\ref{thm:c3v77-floor}(2).  Two structural facts are responsible: the
Hardy constant returns \(3\kappa^2\) of the mean vorticity itself, capping
at \(1/\sqrt3\); and the thresholds \(\theta_1,\theta_{\mathcal P}\), which are the damping of
the higher rows minus their stretching, are small at \(\kappa_1\) in the
windows where the higher rows apply, because \(\kappa_1\) is itself the
point where stretching balances damping in the lowest row.  A loop
that raises the floor needs a row whose stretching is controlled by
something other than \(C_\infty'\) and \(\kappa\), or a functional inequality
sharper than Hardy for divergence-free fields with small mean.
\end{firewall}

% =============================================================================
\section{Integrated V77 conclusion and next acquisition order}
\label{sec:c3v77-conclusion}
% =============================================================================

\begin{theorem}[What V77 proves]
\label{thm:c3v77-conclusion}
Relative to the two legacy cycles and V1--V76: the enstrophy--mean
loop and its constraint (Lemma~\ref{lem:c3v77-rows},
Theorem~\ref{thm:c3v77-loop}, Proposition~\ref{prop:c3v77-comparison}),
its floor (Theorem~\ref{thm:c3v77-floor}), the four-observable
constraint (Theorem~\ref{thm:c3v77-four}) and the bound on its floor
(Proposition~\ref{prop:c3v77-four-floor}), and the negative record
(Firewall~\ref{fw:c3v77-negative}).  The full Navier--Stokes
stability/global-regularity theorem remains unproved.
\end{theorem}

\begin{proof}
The cited results.
\end{proof}

\begin{programme}[V78 acquisition order]
\label{prog:c3v78}
\begin{enumerate}[label=\textup{(AC\arabic*)},leftmargin=4.8em]
\item The vorticity dipole \(D_{ij}:=\int z_j\Omega_i\,G\): derive its identity
      (expected: pure damping \(\frac32\), since \(L_0z_j=-\frac12z_j\), and source
      \(-\int\varepsilon_{ijl}(V\times\Omega)_lG+\frac1{2\nu}\int z_j(z\times(V\times\Omega))_iG\)); bound the
      source by \(\mathcal E\) and \(\mathcal E_1\); the extremal row for \(|D|^2\).
\item The second-order Gaussian Poincar\'e inequality
      \(\int|\Omega-P_{\le1}\Omega|^2G\le2\nu^2\int|\nabla^2\Omega|^2G\) and what a loop through it
      needs (a row for \(\int|\nabla^2\Omega|^2G\), with the third-derivative
      constant); record the tower and whether it can pass the cap.
\item The next compulsory companion audit is V80.
\end{enumerate}
\end{programme}

\begin{assessment}[Position after V77]
\label{ass:c3v77-position}
The cycle has four constants-only constraints, all of
\(\varepsilon\)-regularity type, and a negative record on loops.  No full
stability or global-regularity claim is made.
\end{assessment}

% =============================================================================
\section*{Version V77 (third cycle) append-only record}
\addcontentsline{toc}{section}{Version V77 (third cycle) append-only record}
% =============================================================================

The complete V76 source of the third cycle was retained before this
continuation.  Its SHA--256 digest was
\begin{center}\scriptsize\ttfamily
64b7c404f310dd59f22f4b6c5e22f6a1b78d02949fcd53e70a0e18befefa783f.
\end{center}
V77 closed the enstrophy--mean loop and its four-observable variant
and recorded the negative record of the loop block.  No full
regularity claim is made.

% =============================================================================
\section{V78: the vorticity dipole, the second-order Poincar\'e inequality, and the tower}
\label{sec:c3v78-entry}
% =============================================================================

Items (AC1) and (AC2) of Programme~\ref{prog:c3v78} are carried out.
The identity for the vorticity dipole \(D=\int z\otimes\Omega\,G\) is derived: pure
damping \(\frac32\), no commutator, traceless, with an explicit source
that vanishes on the columnar swirl; its extremal row bounds the
dipole by the enstrophy and the weighted enstrophy.  The second-order
Gaussian Poincar\'e inequality expresses the enstrophy through the
mean, the dipole and the Gaussian integral of \(|\nabla^2\Omega|^2\).  A loop
through it needs a row for that last observable, which brings in
the third-derivative constant; but even before that row is written,
the dipole's Hardy-converted coefficient is \(\frac{80}9\kappa^2\), so the loop
through the dipole caps the floor at \(3/\sqrt{80}\approx0.335\), below the cap
\(1/\sqrt3\) of the mean-vorticity loops.  The tower of higher rows is
recorded with its damping and its derivative constants, and the
negative record of V77 is extended to it.  No global regularity
claim is made.  The next compulsory companion audit is V80.

% =============================================================================
\section{The dipole identity}
\label{sec:c3v78-identity}
% =============================================================================

Put \(D_{ij}:=\int z_j\Omega_i\,G\) (a \(3\times3\) matrix, \(|D|\) its Frobenius norm) and
\begin{equation}
 (S_D)_{ij}:=-\varepsilon_{ijl}\int(V\times\Omega)_l\,G+\frac1{2\nu}\int z_j\,\bigl(z\times(V\times\Omega)\bigr)_i\,G .
 \label{eq:c3v78-source}
\end{equation}

\begin{theorem}[Dipole identity]
\label{thm:c3v78-identity}
For every profile of the family and every \(\sigma\),
\begin{equation}
 \boxed{\quad
 \frac{\dd}{\dd\sigma}D=-\frac32D+S_D,\qquad \operatorname{tr}D=0 .
 \quad}
 \label{eq:c3v78-identity}
\end{equation}
\end{theorem}

\begin{proof}
With the vorticity equation in the form of the proof of
Theorem~\ref{thm:c3v74-identity}: the linear part is
\(\int z_j(L_0\Omega_i)G-\int z_j\Omega_iG=\int\Omega_i(L_0z_j)G-D_{ij}=-\frac12D_{ij}-D_{ij}\), by the
symmetry of \(L_0=\nu\Delta-\frac12z\cdot\nabla\) in \(L^2(G)\) and \(L_0z_j=-\frac12z_j\).  The
source: for \(F=V\times\Omega\),
\(\int z_j\varepsilon_{ikl}\partial_kF_l\,G=-\int\varepsilon_{ikl}F_l\,\partial_k(z_jG)=-\varepsilon_{ijl}\int F_lG+\frac1{2\nu}\int\varepsilon_{ikl}z_jz_kF_lG\),
and \(\varepsilon_{ikl}z_kF_l=(z\times F)_i\).  The trace:
\(\int z\cdot\operatorname{curl}V\,G=\int V\cdot\operatorname{curl}(zG)=\int V\cdot(\nabla G\times z)=0\) since
\(\nabla G\parallel z\).
\end{proof}

\begin{remark}[Check on the columnar swirl]
\label{rem:c3v78-swirl}
For \(V=w(r)e_\theta\), \(\Omega=r^{-1}\partial_r(rw)e_z\), the dipole vanishes (\(\Omega_3\) is
radial, so \(\int z_j\Omega_3G=0\) for every \(j\)), and the source vanishes:
\(V\times\Omega=\frac wr\partial_r(rw)e_r\) has zero Gaussian mean by angular symmetry,
and \(z\times(V\times\Omega)=z_3\frac wr\partial_r(rw)e_\theta\) is odd in \(z_3\).  The identity
reduces to \(0=0\), consistently.
\end{remark}

\begin{lemma}[Bounds on the source]
\label{lem:c3v78-bounds}
For every profile of the family,
\begin{equation}
 |S_D|\ \le\ \kappa\,|G|^{1/2}\Bigl((2\nu)^{1/2}\mathcal E^{1/2}+\bigl(\tfrac32\bigr)^{1/2}\mathcal E_1^{1/2}\Bigr).
 \label{eq:c3v78-bounds}
\end{equation}
\end{lemma}

\begin{proof}
The matrix \((\varepsilon_{ijl}F_l)_{ij}\) has Frobenius norm \(\sqrt2|F|\), so the first
term of \eqref{eq:c3v78-source} is at most
\(\sqrt2\int|V||\Omega|G\le\sqrt2C_\infty|G|^{1/2}\mathcal E^{1/2}\).  The second is at most
\(\frac1{2\nu}\int|z|^2|V||\Omega|G\le\frac{C_\infty}{2\nu}\bigl(\int|z|^2G\bigr)^{1/2}\mathcal E_1^{1/2}=\frac{C_\infty}{2\nu}(6\nu|G|)^{1/2}\mathcal E_1^{1/2}\),
using \(\int|z|^2G=6\nu|G|\).  With \(C_\infty=\kappa\nu^{1/2}\):
\(\sqrt2\kappa\nu^{1/2}=\kappa(2\nu)^{1/2}\) and \(\kappa\nu^{1/2}(6\nu)^{1/2}/(2\nu)=\kappa(3/2)^{1/2}\).
\end{proof}

\begin{theorem}[Extremal dipole]
\label{thm:c3v78-extremal}
Let \(W_D\) be a maximizer of \(|D|^2\) over \(\mathcal F(x_0)\).  Then
\begin{equation}
 \boxed{\quad
 \sup_{\mathcal F(x_0)}|D|^2\ \le\ \frac49|S_D(W_D)|^2\ \le\ \kappa^2|G|\Bigl(\frac{16}9\nu\,\mathcal E+\frac43\mathcal E_1\Bigr)(W_D) .
 \quad}
 \label{eq:c3v78-row}
\end{equation}
\end{theorem}

\begin{proof}
Lemma~\ref{lem:c3v34-extremal} with \(F=|D|^2\): \(0=D:D'=-\frac32|D|^2+D:S_D\) at
\(W_D\), so \(|D|\le\frac23|S_D|\); then \eqref{eq:c3v78-bounds} and \((x+y)^2\le2x^2+2y^2\).
\end{proof}

% =============================================================================
\section{The second-order Poincar\'e inequality and the loop through the dipole}
\label{sec:c3v78-poincare}
% =============================================================================

\begin{proposition}[Second-order Gaussian Poincar\'e inequality]
\label{prop:c3v78-poincare}
With \(\mathcal Q:=\int|\nabla^2\Omega|^2G\), for every profile,
\begin{equation}
 \mathcal E\ \le\ |G|\,|\bar\Omega|^2+\frac{|D|^2}{2\nu|G|}+2\nu^2\,\mathcal Q .
 \label{eq:c3v78-poincare}
\end{equation}
\end{proposition}

\begin{proof}
Hermite expansion in \(L^2(G)\) (C2-V8-spectrum): \(-L_0\) has eigenvalue
\(k/2\) on the level-\(k\) space \(H_k\), and \(\nu\int|\nabla f|^2G=-\int fL_0f\,G\), so
\(\int|\nabla f|^2G=\sum_k\frac k{2\nu}\|P_kf\|_G^2\) and, the gradient mapping \(H_k\) to
\(H_{k-1}\), \(\int|\nabla^2f|^2G=\sum_k\frac{k(k-1)}{4\nu^2}\|P_kf\|_G^2\ge\frac1{2\nu^2}\sum_{k\ge2}\|P_kf\|_G^2\).
The level-0 part of \(\Omega\) has norm \(|G||\bar\Omega|^2\); the level-1 part is
\(\sum_jc_jz_j\) with \(c_j=\int\Omega z_jG/(2\nu|G|)\), of norm
\(\sum_{ij}D_{ij}^2/(2\nu|G|)\).
\end{proof}

\begin{theorem}[The loop through the dipole, and its cap]
\label{thm:c3v78-cap}
At a homogeneous type-I point, with \(s_1=\sup\mathcal E\), \(s_2=\sup|G||\bar\Omega|^2\),
\(s_3=\sup\mathcal P\), \(s_D=\sup|D|^2\), \(s_{\mathcal Q}=\sup\mathcal Q\),
\begin{equation}
 s_1\le s_2+\frac{s_D}{2\nu|G|}+2\nu^2s_{\mathcal Q},\qquad
 \frac{s_D}{2\nu|G|}\le\frac{80}9\kappa^2s_1+\frac{32}3\kappa^2\nu\,s_3 ,
 \label{eq:c3v78-dipole-rows}
\end{equation}
so that any loop closing through \eqref{eq:c3v78-poincare} with the
dipole row has, in the row for \(s_1\), a diagonal coefficient at least
\(\frac{80}9\kappa^2\) (after substituting \(s_2\) and \(s_D\)), and under (E5-small)
its floor is at most \(3/\sqrt{80}\approx0.335\) whenever the coefficients of
\(s_3\) and \(s_{\mathcal Q}\) are nonnegative.
\end{theorem}

\begin{proof}
The first inequality is \eqref{eq:c3v78-poincare} at the maximizer of
\(\mathcal E\).  The second: \eqref{eq:c3v78-row} with Hardy,
\(\mathcal E_1\le12\nu\mathcal E+16\nu^2\mathcal P\), gives
\(s_D\le\kappa^2|G|(\frac{16}9\nu s_1+\frac43(12\nu s_1+16\nu^2s_3))=\kappa^2|G|\nu(\frac{160}9s_1+\frac{64}3\nu s_3)\);
divide by \(2\nu|G|\).  The cap: a constraint \(\rho(A)\ge1\) is implied by the
\(s_1\)-row coefficient alone being at least one, and \(\frac{80}9\kappa^2\ge1\) iff
\(\kappa\ge3/\sqrt{80}\); the floor is the least \(\kappa\) at which the constraint
holds, hence at most this value.
\end{proof}

\begin{assessment}[Item (AC2): the tower]
\label{ass:c3v78-tower}
The loop through \eqref{eq:c3v78-poincare} needs a row for \(\mathcal Q\).  Its
identity has the form of V72's: damping \(2\) (weight \(1\) of \(\Omega\) plus
\(\frac12\) per derivative), a dissipation \(\nu\int|\nabla^3\Omega|^2G\), stretching terms
pairing \(\nabla^2\Omega\) with \(\nabla^2((\Omega\cdot\nabla)V)\), bounded by
\(C_\infty'\mathcal Q+C_\infty''\mathcal P^{1/2}\mathcal Q^{1/2}+C_\infty'''\mathcal E^{1/2}\mathcal Q^{1/2}\) up to combinatorial factors, and a
transport with two commutator terms bounded by \(C_\infty'\mathcal Q+C_\infty''\mathcal P^{1/2}\mathcal Q^{1/2}\)
and the Hardy term \((\frac{\sqrt3}2\kappa+\frac{\kappa^2}2)\mathcal Q+\frac12\nu\int|\nabla^3\Omega|^2G\); its extremal
row bounds \(\sup\mathcal Q\) by \(\sup\mathcal P\) and \(\sup\mathcal E\) with a threshold
\(\theta_{\mathcal Q}=2-cC_\infty'-\frac{\sqrt3}2\kappa-\frac{\kappa^2}2\) and coefficients in \(C_\infty''\), \(C_\infty'''\).  The
\(m\)-th level \(\int|\nabla^m\Omega|^2G\) has damping \(1+\frac m2\), threshold
\(\theta_m=1+\frac m2-c_mC_\infty'-\frac{\sqrt3}2\kappa-\frac{\kappa^2}2\), and stretching coefficients up to
\(C_\infty^{(m+1)}\); under (E5-small) each level adds a constant \(\gamma_{m+1}\).  The
tower is finite at each depth and explicit, but by
Theorem~\ref{thm:c3v78-cap} no loop that enters through the dipole
can raise the floor above \(0.335\), and no loop that enters through
the mean alone above \(1/\sqrt3\) (Firewall~\ref{fw:c3v77-negative}); the
Poincar\'e route returns the low Hermite levels of \(\Omega\) with
coefficients \(c\kappa^2\), \(c\ge3\), whatever the depth.  The tower is
therefore not written beyond this record.
\end{assessment}

\begin{firewall}[Negative record extended]
\label{fw:c3v78-negative}
Loops of the vorticity ledger whose closing inequality is a Gaussian
Poincar\'e inequality of any order, with the rows of the tower and
the Hardy conversion of the low-level rows, have floors at most
\(\max\{\kappa_1,1/\sqrt3\}\) under (E5-small).  What could pass the cap is a
closing inequality that does not return the low levels with an
order-one multiple of \(\kappa^2\): for instance a bound on \(|\bar\Omega|\) or \(|D|\)
by the enstrophy with a small constant, which would need the mean
row's source \(\mathcal B\) to be small for reasons other than the class bounds
(a cancellation in \(\int z\times(V\times\Omega)G\)), or the divergence-free
structure of \(\Omega\) in the Hardy step.  Neither is available.
\end{firewall}

% =============================================================================
\section{Integrated V78 conclusion and next acquisition order}
\label{sec:c3v78-conclusion}
% =============================================================================

\begin{theorem}[What V78 proves]
\label{thm:c3v78-conclusion}
Relative to the two legacy cycles and V1--V77: the dipole identity
(Theorem~\ref{thm:c3v78-identity}), its bounds and extremal row
(Lemma~\ref{lem:c3v78-bounds}, Theorem~\ref{thm:c3v78-extremal}), the
second-order Poincar\'e inequality (Proposition~\ref{prop:c3v78-poincare}),
the cap of the dipole loop (Theorem~\ref{thm:c3v78-cap}), and the
tower record with the extended negative record
(Assessment~\ref{ass:c3v78-tower}, Firewall~\ref{fw:c3v78-negative}).  The
full Navier--Stokes stability/global-regularity theorem remains
unproved.
\end{theorem}

\begin{proof}
The cited results.
\end{proof}

\begin{programme}[V79 acquisition order]
\label{prog:c3v79}
\begin{enumerate}[label=\textup{(AC\arabic*)},leftmargin=4.8em]
\item Consolidate the loop block V74, V76--V78 as one theorem: the
      four constants-only constraints, the matrix lemma, the floors
      and the two negative records.
\item Fix the direction for V81--V95.  Candidates, to be decided by
      what can be made explicit: (a) computing \(\gamma_1\), the constant of
      the first-derivative bound for small solutions of the profile
      equation, from the class (the kernel constants of V47 and the
      pressure bounds), which would make \(\kappa_1\) a number; (b) the
      cancellation in the mean row's source
      \(\int z\times(V\times\Omega)G=\int[V(z\cdot\Omega)-\Omega(z\cdot V)]G\) for divergence-free fields;
      (c) closing the cycle's ledger early and reserving V81--V95 for
      the manifest.
\item The next compulsory companion audit is V80.
\end{enumerate}
\end{programme}

\begin{assessment}[Position after V78]
\label{ass:c3v78-position}
The dipole row is written and the loop method is bounded above by
explicit caps; the block's negative record is complete.  No full
stability or global-regularity claim is made.
\end{assessment}

% =============================================================================
\section*{Version V78 (third cycle) append-only record}
\addcontentsline{toc}{section}{Version V78 (third cycle) append-only record}
% =============================================================================

The complete V77 source of the third cycle was retained before this
continuation.  Its SHA--256 digest was
\begin{center}\scriptsize\ttfamily
a8c3557b95c48bdd5972962a8e40e97063da753c3589cdccd1f83c6bcdc72ad0.
\end{center}
V78 derived the dipole identity and recorded the caps of the loop
method.  No full regularity claim is made.

% =============================================================================
\section{V79: the loop block consolidated, and the direction for V81--V95}
\label{sec:c3v79-entry}
% =============================================================================

Items (AC1) and (AC2) of Programme~\ref{prog:c3v79} are carried out.
The loop block V74, V76--V78 is consolidated: the matrix lemma, the
four constants-only constraints, their floors under (E5-small) and
the two negative records.  For the direction, candidate (a) of
Programme~\ref{prog:c3v79} is chosen: the constant \(\gamma_1\) of the
first-derivative bound for small solutions is computable from the
mild formulation of the ancient solutions with two kernel constants
(the gradient of the heat kernel and the gradient of the Oseen
kernel in \(L^1\)), and with it the floor \(\kappa_1\) becomes a number;
V81--V95 are assigned to this computation, to the second-derivative
constant, and to the numerical table of the floors.  Candidates (b)
and (c) are recorded with the reasons for not choosing them.  No
global regularity claim is made.  The next compulsory companion
audit is V80.

% =============================================================================
\section{The loop block}
\label{sec:c3v79-block}
% =============================================================================

\begin{theorem}[The loop block, V74 and V76--V78]
\label{thm:c3v79-block}
At a homogeneous type-I point, with the notation of V67, V71, V72,
V74:
\begin{enumerate}[label=\textup{(\arabic*)},leftmargin=3.2em]
\item (Matrix lemma, V76.)  If nonnegative continuous observables
      have suprema \(s\) with \(s\le As\), \(A\ge0\) entrywise, and \(\rho(A)<1\), all
      vanish; hence \(\rho(A)\ge1\) whenever one is positive somewhere.
\item (Rows.)  Mean vorticity: \(\bar\Omega'=-\bar\Omega+\mathcal B\), \(|G||\bar\Omega|^2\le\frac{\kappa^2}{4\nu}\mathcal E_1\le\kappa^2(3\mathcal E+4\nu\mathcal P)\)
      at its maximizer (V74).  Dipole: \(D'=-\frac32D+S_D\), \(\operatorname{tr}D=0\),
      \(|D|^2\le\kappa^2|G|(\frac{16}9\nu\mathcal E+\frac43\mathcal E_1)\) at its maximizer (V78).  With the
      enstrophy, weighted-enstrophy and palinstrophy rows of
      V67, V71, V72.
\item (Constraints, wherever \(\theta_{\mathcal P}>0\).)
      \(3\kappa^2+(1+2\kappa^2)q_{\mathcal P}\ge1\) (V74);
      \(4\delta_{67}+3\kappa^2+2\kappa^2q_{\mathcal P}\ge1\) (V77), dominating V67 where
      \(3\kappa^2+2\kappa^2q_{\mathcal P}<1\); and, where also \(\theta_1>0\),
      \(4\delta_{67}+\kappa^2(\Xi_1+q_{\mathcal P})/(4\theta_1)\ge1\) (V77).
\item (Floors under (E5-small).)  \(\kappa_2\le\min\{\kappa_2^+,\kappa_r\}\) with
      \(\kappa_2^+=(3+\frac89\gamma_2^2)^{-1/2}\); \(\kappa_3\le\min\{\max\{\kappa_1,1/\sqrt3\},\kappa_r\}\); \(\kappa_4\le\kappa_1\);
      \(\kappa_2>\kappa_1\) iff \(1/\sqrt3<\gamma_1<1.094\) and \(\gamma_2<\gamma_2^*(\gamma_1)\le0.042\);
      \(\kappa_3>\kappa_1\) iff the same with \(\gamma_3^*(\gamma_1)\le0.069\).
\item (Negative records.)  Loops closing through Gaussian Poincar\'e
      inequalities of any order, with the tower's rows and the Hardy
      conversion, have floors at most \(\max\{\kappa_1,1/\sqrt3\}\); entering
      through the dipole, at most \(3/\sqrt{80}\).
\end{enumerate}
\end{theorem}

\begin{proof}
Lemma~\ref{lem:c3v76-matrix}; Theorems~\ref{thm:c3v74-identity},
\ref{thm:c3v74-extremal}, \ref{thm:c3v78-identity}, \ref{thm:c3v78-extremal};
Theorems~\ref{thm:c3v74-loop}, \ref{thm:c3v77-loop}, \ref{thm:c3v77-four};
Theorems~\ref{thm:c3v76-floor}, \ref{thm:c3v76-comparison},
\ref{thm:c3v77-floor}, Proposition~\ref{prop:c3v77-four-floor};
Firewalls~\ref{fw:c3v77-negative}, \ref{fw:c3v78-negative},
Theorem~\ref{thm:c3v78-cap}.
\end{proof}

\begin{assessment}[What the block settles]
\label{ass:c3v79-settles}
The vorticity ledger at extremal configurations yields relations
between explicit constants, and the loop method multiplies them,
but every relation so obtained is of \(\varepsilon\)-regularity type: it says
that the velocity constant of a homogeneous type-I point is not
small, in a form explicit in the derivative constants.  The
strongest such floor in the series is \(\kappa_1(\gamma_1)\), from the single
enstrophy row, and the block shows that the loops cannot improve it
beyond slivers of the parameter space.  The unknown that decides
the value of every floor is \(\gamma_1\); the block leaves it as the one
quantity of the vorticity chapter that the series has not made
explicit.
\end{assessment}

% =============================================================================
\section{Direction for V81--V95: the explicit gradient constant}
\label{sec:c3v79-direction}
% =============================================================================

\begin{assessment}[The mild formulation gives \(\gamma_1\) with two kernel constants]
\label{ass:c3v79-direction}
Let \(W\in\mathcal F(x_0)\) be an ancient solution on \((-\infty,0)\times\R^3\) with
\(|W(s)|\le C_\infty(-s)^{-1/2}\), \(W(s)\in L^6\), and pressure \(\sum R_iR_j(W_iW_j)\) (V47).
It satisfies the Duhamel formula from any \(s_0<s\),
\(W(s)=e^{\nu(s-s_0)\Delta}W(s_0)-\int_{s_0}^se^{\nu(s-\tau)\Delta}\PP\operatorname{div}(W\otimes W)(\tau)\dd\tau\)
(the class and (E8)).  Differentiating, with the gradient falling on
the heat kernel in the first term and on \(W\otimes W\) in the second,
\begin{equation}
 \|\nabla W(s)\|_\infty\le\frac{c_h\|W(s_0)\|_\infty}{(\nu(s-s_0))^{1/2}}
 +\int_{s_0}^s\frac{c_O\,\|\nabla(W\otimes W)(\tau)\|_\infty}{(\nu(s-\tau))^{1/2}}\dd\tau ,
 \label{eq:c3v79-duhamel}
\end{equation}
where \(c_h:=\sup_t(\nu t)^{1/2}\|\nabla K_{\nu t}\|_{L^1}=2/\sqrt\pi\) (the heat kernel, an
explicit Gaussian moment) and \(c_O:=\sup_t(\nu t)^{1/2}\|\nabla K^O_{\nu t}\|_{L^1}\) with
\(K^O_t\) the Oseen kernel of \(e^{t\Delta}\PP\), finite and explicit
(\(K^O_t=K_t\delta_{ij}+\partial_i\partial_j\Phi_t\), \(\Phi_t=(-\Delta)^{-1}K_t=\operatorname{erf}(|x|/(2t^{1/2}))/(4\pi|x|)\)).
With \(s_0=2s\) and \(N:=\sup_{\tau<0}(-\tau)\|\nabla W(\tau)\|_\infty\) (which is \(C_\infty'\) in the units
of the profile), \eqref{eq:c3v79-duhamel} gives
\(N\le\frac{c_h}{\sqrt2}\kappa+2\sqrt2\,c_O\,\kappa\,N\), using
\(\int_{2s}^s(s-\tau)^{-1/2}(-\tau)^{-3/2}\dd\tau=\sqrt2\,(-s)^{-1}\); hence
\begin{equation}
 C_\infty'\ \le\ \gamma_1(\kappa)\,\kappa,\qquad
 \gamma_1(\kappa)=\frac{c_h/\sqrt2}{1-2\sqrt2\,c_O\kappa}
 \qquad(\kappa<\tfrac1{2\sqrt2c_O}),
 \label{eq:c3v79-gamma1}
\end{equation}
which is (E5-small) with \(\kappa_{\rm small}\) any value below \(1/(2\sqrt2c_O)\) and
\(\gamma_1=\gamma_1(\kappa_{\rm small})\), all explicit once \(c_O\) is bounded.  The a priori
finiteness of \(N\) on each window holds because \(W\) is smooth with
the class derivative bounds; the estimate is then an absorption.
V81 sets up the mild formulation and computes \(c_h\); V82 bounds \(c_O\)
explicitly from the formula for \(\partial_i\partial_j\Phi_t\); V83 proves
\eqref{eq:c3v79-gamma1} in the class; V84 evaluates \(\kappa_1(\gamma_1(\kappa))\) as a
fixed-point floor, i.e.\ the least \(\kappa\) with \(\gamma_1(\kappa)\kappa+\frac{\sqrt3}2\kappa+\frac{\kappa^2}2\ge1\), a
number; V85 audits.  V86--V89 bound the second-derivative constant
by one more differentiation of the Duhamel formula, giving \(\gamma_2(\kappa)\)
and the explicit values of \(q_{\mathcal P}\) and of the loop floors; V90 audits;
V91--V94 write the numerical table of the floors and compare with the
literature's explicit \(\varepsilon\)-regularity constants as far as the series
records them; V95 audits; V96--V100 close the cycle.  No global
regularity claim is made: the outcome is an explicit lower bound on
the velocity constant of a homogeneous type-I point, a quantitative
\(\varepsilon\)-regularity statement, not an exclusion.
\end{assessment}

\begin{remark}[Candidates (b) and (c)]
\label{rem:c3v79-others}
(b) A cancellation in \(\int[V(z\cdot\Omega)-\Omega(z\cdot V)]G\) for divergence-free
fields would lower the coefficient \(3\kappa^2\); none is visible in the
class (both terms are of the size of their bounds on a Beltrami-type
configuration with \(\Omega\parallel V\) localized off the origin), and the
question is deferred.  (c) Closing early would leave \(\gamma_1\) unknown
while it is computable; the computation is the natural end of the
vorticity chapter.
\end{remark}

% =============================================================================
\section{Integrated V79 conclusion and next acquisition order}
\label{sec:c3v79-conclusion}
% =============================================================================

\begin{theorem}[What V79 establishes]
\label{thm:c3v79-conclusion}
Relative to the two legacy cycles and V1--V78: the consolidated loop
block (Theorem~\ref{thm:c3v79-block}) and the direction
(Assessment~\ref{ass:c3v79-direction}, whose formulas are a plan, to be
proved in V81--V83).  The full Navier--Stokes stability/global-regularity
theorem remains unproved.
\end{theorem}

\begin{proof}
The cited results.
\end{proof}

\begin{programme}[V80 acquisition order]
\label{prog:c3v80}
\begin{enumerate}[label=\textup{(AC\arabic*)},leftmargin=4.8em]
\item The compulsory companion audit: read the current SM and YM
      notes and record their digests and decision.
\item Record the position after eighty versions; delivery.
\end{enumerate}
\end{programme}

\begin{assessment}[Position after V79]
\label{ass:c3v79-position}
The loop block is consolidated; the direction is the explicit
gradient constant.  No full stability or global-regularity claim is
made.
\end{assessment}

% =============================================================================
\section*{Version V79 (third cycle) append-only record}
\addcontentsline{toc}{section}{Version V79 (third cycle) append-only record}
% =============================================================================

The complete V78 source of the third cycle was retained before this
continuation.  Its SHA--256 digest was
\begin{center}\scriptsize\ttfamily
f8ebd5428061146ca0ca07c0cf3a7a999fa9131103b306cd1721d6640ed42967.
\end{center}
V79 consolidated the loop block and fixed the direction for
V81--V95.  No full regularity claim is made.

% =============================================================================
\section{V80: compulsory companion audit and the position after eighty versions}
\label{sec:c3v80-entry}
% =============================================================================

This is the eightieth version of the third cycle.  The compulsory
review of the two companion programmes is performed first, and the
position after eighty versions is recorded with the plan for the
block V81--V85.  No global regularity claim is made.  The next
compulsory companion audit is V85.

% =============================================================================
\section{Compulsory V80 review of the current companion notes}
\label{sec:c3v80-companion-audit}
% =============================================================================

The current companion files in \texttt{latex\_notes} were inspected
read-only.  Both programmes have moved since the V75 audit: the SM
programme from V44 to V45 of its seventh series, the YM programme
from V21 to V22 of its sixth series.  The frozen checkpoints are
\begin{align}
 &\texttt{Renewal\_SM\_Einstein\_Continuum\_Research\_Notes\_V45.tex},
 \notag\\[-2pt]
 &\qquad
 \texttt{5552283b0d2ab6f97f486c89143c35a84a2ef50b4187ace0d775e83b5283c121},
 \label{eq:c3v80-SM-checkpoint}\\
 &\texttt{Renewal\_Yang\_Mills\_Mass\_Gap\_Research\_Notes\_V22.tex},
 \notag\\[-2pt]
 &\qquad
 \texttt{4fcc777a1932deaa2efc708be1d3ab6a7c13cd290b0c3ba17112a8dd3325c393}.
 \label{eq:c3v80-YM-checkpoint}
\end{align}
(SM V45 has 18015 lines; YM V22 has 9612 lines.  At the inspection
the folder also held a file named YM V23, byte-identical to YM V22,
a version in preparation; it is not a distinct state and is not
recorded separately.)

\emph{SM V45.}  The SM programme performed its own companion audit
and then treated the frozen edge sector fan, a Lipschitz filling of
the transverse fan, two-normal adiabatic localization and the joint
moment defect, continuing its collar-completion chapter.

\emph{YM V22.}  The YM programme audited a transfer interface with
the Grand Tensor manuscript, imported a wrapping-toron obstruction at
its proved scope, forced a response-fixed or local-vacuum
trajectory, proved that exponential Gram decay implies a cyclic
spectral gap and that a common exponent on a dense family gives a
full gap, derived a Dobrushin interaction-distance tail, proved that
its condition WCIC gives a finite-bank Osterwalder--Schrader gap with
an explicit mass rate, and stated an exact conditional continuum
mass handoff.

\begin{theorem}[V80 companion-audit decision]
\label{thm:c3v80-companion-decision}
Nothing is imported from SM V45 or YM V22.  Neither bears on the
third cycle's chapters, on the vorticity ledger, on the loop method
or on the explicit gradient constant.
\end{theorem}

\begin{proof}
The SM results concern edge fans and adiabatic localization of
lattice Dirac interfaces; the YM results concern spectral gaps of
lattice transfer matrices under conditional-influence hypotheses.
None is a Navier--Stokes statement, and none concerns the ledger, the
extremal method, the vorticity identities, the matrix lemma or the
mild formulation of the ancient solutions.
\end{proof}

% =============================================================================
\section{Position after eighty versions, and the block V81--V85}
\label{sec:c3v80-position}
% =============================================================================

\begin{assessment}[Position after eighty versions of the third cycle]
\label{ass:c3v80-position-eighty}
Fourteen chapters are written: those listed in
Assessment~\ref{ass:c3v75-position-seventyfive} and the loop block
(V74, V76--V79).  The cycle's explicit results at a homogeneous
type-I point are now: the enstrophy constraint
\(C_\infty'+\frac{\sqrt3}2\kappa+\frac{\kappa^2}2\ge1\) (V67); the three loop constraints (V74, V77)
wherever the palinstrophy row applies; the concentration and ratio
bounds (V71, V72); the mean-vorticity and dipole identities and rows
(V74, V78); and, under (E5-small), the floors \(\kappa_1,\kappa_2,\kappa_3,\kappa_4\) with
\(\kappa_1\) the operative one except in explicit slivers.  Every one of
them is of \(\varepsilon\)-regularity type and none is an exclusion; the
negative records of V73, V77 and V78 say that the ledger and its
loops cannot do more with the rows written.  The one uncomputed
quantity is \(\gamma_1\), and V79 showed that it is computable from the
mild formulation with two kernel constants.  The block V81--V85
follows Assessment~\ref{ass:c3v79-direction}: the mild formulation and
\(c_h\) (V81), the Oseen kernel and \(c_O\) (V82), the bound
\(C_\infty'\le\gamma_1(\kappa)\kappa\) in the class (V83), the numerical floor (V84), audit
(V85).  Twenty versions remain; the closing block V96--V100 is
reserved as before.  No global regularity claim is made.
\end{assessment}

% =============================================================================
\section{Integrated V80 conclusion and next acquisition order}
\label{sec:c3v80-conclusion}
% =============================================================================

\begin{theorem}[What V80 establishes]
\label{thm:c3v80-conclusion}
Relative to the two legacy cycles and V1--V79: the companion-audit
decision (Theorem~\ref{thm:c3v80-companion-decision}) and the position.
The full Navier--Stokes stability/global-regularity theorem remains
unproved.
\end{theorem}

\begin{proof}
The cited records.
\end{proof}

\begin{programme}[V81 acquisition order]
\label{prog:c3v81}
\begin{enumerate}[label=\textup{(AC\arabic*)},leftmargin=4.8em]
\item The mild formulation of the ancient solutions of the family:
      state and prove that every \(W\in\mathcal F(x_0)\) satisfies the Duhamel
      formula on \(\R^3\) from every \(s_0<s\), from the class (smooth,
      \(L^6\), the Riesz pressure) and (E8); record the a priori
      finiteness of \(\sup_{[s_0,s]}\|\nabla W\|_\infty\) from the class derivative
      bounds.
\item The heat-kernel constant: \(\|\nabla K_t\|_{L^1}=2/(\pi t)^{1/2}\) (Gaussian
      moment), hence \(c_h=2/\sqrt\pi\); the commutation of \(\nabla\) with
      \(e^{t\Delta}\PP\operatorname{div}\) on the tensor \(W\otimes W\).
\item The next compulsory companion audit is V85.
\end{enumerate}
\end{programme}

\begin{assessment}[Position after V80]
\label{ass:c3v80-position}
The companion audit is recorded; the explicit-constant block begins.
No full stability or global-regularity claim is made.
\end{assessment}

% =============================================================================
\section*{Version V80 (third cycle) append-only record}
\addcontentsline{toc}{section}{Version V80 (third cycle) append-only record}
% =============================================================================

The complete V79 source of the third cycle was retained before this
continuation.  Its SHA--256 digest was
\begin{center}\scriptsize\ttfamily
f15c8245123af0e2e67f0da947ac0cd318269293e7a0db3c0945537214aeca04.
\end{center}
V80 performed the compulsory companion audit and recorded the
position.  No full regularity claim is made.

% =============================================================================
\section{V81: the mild formulation of the ancient solutions, and the first explicit floors}
\label{sec:c3v81-entry}
% =============================================================================

Items (AC1) and (AC2) of Programme~\ref{prog:c3v81} are carried out.
Every element of the family satisfies the Duhamel formula on every
finite window, and, the window being sent to \(-\infty\), the ancient
form without initial term.  The heat-kernel constant is
\(c_h=2/\sqrt\pi\).  The Oseen kernel is written explicitly and its
gradient is integrable, which defines the constant \(c_O\) (bounded
numerically in V82).  Two absorptions follow at once: from the
ancient form, every nonzero element of the family has velocity
constant at least \(1/(\pi c_O)\), an explicit \(\varepsilon\)-regularity floor of
the kind the series had only from (E5) without a value; and from
the finite-window form, the gradient constant satisfies
\(C_\infty'\le\gamma_1(\kappa,\theta)\kappa\) with \(\gamma_1\) explicit in \(c_h\), \(c_O\), \(\kappa\) and a free
window parameter \(\theta\).  A notation clash is recorded.  No global
regularity claim is made.  The next compulsory companion audit is
V85.

\begin{remark}[Notation]
\label{rem:c3v81-notation}
The constants \(\gamma_k\) of Lemma~\ref{lem:c3v22-growth} (growth of Hermite
projections) are unrelated to the constants \(\gamma_1,\gamma_2\) of
Hypothesis (E5-small) (V73); the clash is recorded here and not
repaired, both notations being local to their chapters.
\end{remark}

% =============================================================================
\section{Setting}
\label{sec:c3v81-setting}
% =============================================================================

Throughout this block, \(W\in\mathcal F(x_0)\) is an element of the family at a
homogeneous type-I point: a smooth divergence-free ancient solution
of the Navier--Stokes equations on \((-\infty,0)\times\R^3\) with
\begin{equation}
 |W(y,s)|\le C_\infty(-s)^{-1/2},\qquad |\nabla W(y,s)|\le C_\infty'(-s)^{-1},\qquad
 \|W(s)\|_{L^6}\le K_6(-s)^{-1/4},\qquad \|\nabla W(s)\|_{L^2}^2\le C(-s)^{-1/2},
 \label{eq:c3v81-class}
\end{equation}
derivative bounds of every order on half-lines \(s\le s_1<0\)
(Theorem~\ref{thm:c3v1-ancient}, C2-V4), and pressure
\(p=\sum_{ij}R_iR_j(W_iW_j)\) (C2-V36, Section~\ref{sec:c3v27-pressure}).  The
profile is \(V(z,\sigma)=\lambda W(\lambda z,-\lambda^2)\), \(\lambda=e^{-\sigma/2}\), so that
\(\|\nabla V(\sigma)\|_\infty=(-s)\|\nabla W(s)\|_\infty\) at \(s=-\lambda^2\), and
\begin{equation}
 \kappa=\nu^{-1/2}\sup_{s<0}(-s)^{1/2}\|W(s)\|_\infty,\qquad
 C_\infty'=\sup_{s<0}(-s)\|\nabla W(s)\|_\infty
 \label{eq:c3v81-constants}
\end{equation}
are the class constants (the suprema over the family being the
constants of the class; the arguments below apply to each \(W\)).
\(K_t(x)=(4\pi t)^{-3/2}e^{-|x|^2/(4t)}\) is the heat kernel and \(\PP=I-\nabla\Delta^{-1}\operatorname{div}\)
the Leray projector.

% =============================================================================
\section{The heat kernel and the Oseen kernel}
\label{sec:c3v81-kernels}
% =============================================================================

\begin{lemma}[Heat-kernel gradient constant]
\label{lem:c3v81-heat}
\(\|\nabla K_t\|_{L^1(\R^3)}=2/(\pi t)^{1/2}\); hence, for bounded \(g\),
\(\|\nabla e^{\nu t\Delta}g\|_\infty\le c_h(\nu t)^{-1/2}\|g\|_\infty\) with \(c_h=2/\sqrt\pi\).
\end{lemma}

\begin{proof}
\(\nabla K_t=-\frac x{2t}K_t\), so \(\|\nabla K_t\|_1=\frac1{2t}\int|x|K_t=\frac1{2t}\,\mathbb E|X|\) for
\(X\sim N(0,2tI_3)\); \(\mathbb E|X|=(2t)^{1/2}\,\mathbb E|N(0,I_3)|=(2t)^{1/2}\cdot2\sqrt2/\sqrt\pi\), giving
\(\frac{2}{\sqrt{\pi t}}\).  The operator bound is Young's inequality with
\(e^{\nu t\Delta}g=K_{\nu t}*g\).
\end{proof}

\begin{lemma}[The Oseen kernel]
\label{lem:c3v81-oseen}
The operator \(e^{t\Delta}\PP\) on \(L^2(\R^3)^3\) is the convolution with the matrix
kernel
\begin{equation}
 K^O_t=K_t\,I+\nabla^2\Phi_t,\qquad
 \Phi_t(x)=\frac{\operatorname{erf}\bigl(|x|/(2t^{1/2})\bigr)}{4\pi|x|}=(-\Delta)^{-1}K_t ,
 \label{eq:c3v81-oseen}
\end{equation}
with \(K^O_t(x)=t^{-3/2}K^O_1(x/t^{1/2})\).  \(K^O_1\) is smooth, \(|\nabla^kK^O_1(x)|\le c_k(1+|x|)^{-3-k}\),
and \(\nabla K^O_t\in L^1\) with \(\|\nabla K^O_t\|_{L^1}=t^{-1/2}\|\nabla K^O_1\|_{L^1}\).  Consequently, for
a bounded matrix field \(f\) with \(\operatorname{div}f\in L^2\),
\begin{equation}
 \|e^{\nu t\Delta}\PP\operatorname{div}f\|_\infty\le c_O\,(\nu t)^{-1/2}\|f\|_\infty,\qquad
 c_O:=\sup\bigl\{\|e^{\Delta}\PP\operatorname{div}f\|_\infty:\ \|f\|_\infty\le1\bigr\}\le\|\nabla K^O_1\|_{L^1}<\infty ,
 \label{eq:c3v81-cO}
\end{equation}
where \(\|f\|_\infty=\sup_x|f(x)|\) in the Frobenius norm and the last
\(L^1\) norm is that of the pointwise operator norm of the tensor
\(\nabla K^O_1\) acting on matrices.
\end{lemma}

\begin{proof}
The symbol of \(e^{t\Delta}\PP\) is \(e^{-t|\xi|^2}(\delta_{ij}-\xi_i\xi_j/|\xi|^2)\); the second term is
the symbol of \(\partial_i\partial_j\) applied to the function with symbol \(e^{-t|\xi|^2}/|\xi|^2\),
i.e.\ to \((-\Delta)^{-1}K_t=\Gamma_N*K_t\), \(\Gamma_N=1/(4\pi|y|)\).  For the radial
density \(K_t\), \(\Gamma_N*K_t(x)=\frac{Q(r)}{4\pi r}+\int_{|y|>r}\frac{K_t(y)}{4\pi|y|}\dd y\) with \(r=|x|\)
and \(Q(r)=\int_{|y|<r}K_t\); here \(Q(r)=\operatorname{erf}(r/(2t^{1/2}))-\frac{r}{(\pi t)^{1/2}}e^{-r^2/(4t)}\) and
\(\int_{|y|>r}\frac{K_t}{4\pi|y|}=\frac1{4\pi}\int_r^\infty4\pi\rho\,K_t(\rho)\dd\rho=\frac{2t}{(4\pi t)^{3/2}}e^{-r^2/(4t)}=\frac{e^{-r^2/(4t)}}{4\pi(\pi t)^{1/2}}\),
which cancels the second term of \(Q(r)/(4\pi r)\); this gives
\eqref{eq:c3v81-oseen} (the sign: \(\nabla^2\Phi_t\) has symbol \(-\xi_i\xi_je^{-t|\xi|^2}/|\xi|^2\)).
Scaling is that of \(K_t\).  Smoothness: \(\operatorname{erf}(r/2)/r\) is an even entire
function of \(r\), so \(\Phi_1\) is smooth.  Decay: \(\Phi_1=\frac1{4\pi r}-\frac{\operatorname{erfc}(r/2)}{4\pi r}\)
with \(\operatorname{erfc}(r/2)\le e^{-r^2/4}\) and all its derivatives Gaussian, while
\(|\nabla^{k+2}(1/r)|\le c\,r^{-3-k}\); together with the Gaussian \(K_1\) this
gives the stated decay, and \(\nabla K^O_1\in L^1\).  Then
\(e^{t\Delta}\PP\operatorname{div}f=\partial_j(K^O_t)_{il}*f_{lj}\) (the multipliers commute; the identity
holds on \(L^2\) and extends to bounded \(f\) with \(\operatorname{div}f\in L^2\) by
approximation, both sides being defined) and Young's inequality with
the \(t\)-scaling give \eqref{eq:c3v81-cO}; \(\nu t\) replaces \(t\) for \(e^{\nu t\Delta}\).
\end{proof}

% =============================================================================
\section{The mild formulation}
\label{sec:c3v81-mild}
% =============================================================================

\begin{theorem}[Duhamel formula on finite windows]
\label{thm:c3v81-duhamel}
For every \(W\in\mathcal F(x_0)\) and \(s_0<s<0\),
\begin{equation}
 W(s)=e^{\nu(s-s_0)\Delta}W(s_0)-\int_{s_0}^se^{\nu(s-\tau)\Delta}\PP\operatorname{div}\bigl(W\otimes W\bigr)(\tau)\,\dd\tau ,
 \label{eq:c3v81-duhamel}
\end{equation}
the integral converging absolutely in \(L^\infty\).
\end{theorem}

\begin{proof}
On \([s_0,s]\), by \eqref{eq:c3v81-class} and the class derivative bounds,
\(W\), \(\nabla W\), \(\nabla^2W\), \(\nabla p\) and \(\partial_\tau W=\nu\Delta W-(W\cdot\nabla)W-\nabla p\) are bounded and
continuous in \(\tau\) with values in \(L^\infty\) (Lemma~\ref{lem:c3v47-decomposition}
gives \(|\nabla p|\le C_P(-\tau)^{-3/2}\) in the ancient variables).  Put
\(g(\tau):=e^{\nu(s-\tau)\Delta}W(\tau)\); by dominated convergence \(g\) is differentiable
on \([s_0,s)\) with \(g'(\tau)=e^{\nu(s-\tau)\Delta}(\partial_\tau W-\nu\Delta W)(\tau)=-e^{\nu(s-\tau)\Delta}f(\tau)\),
\(f:=(W\cdot\nabla)W+\nabla p\), and \(g\) is continuous at \(\tau=s\) with \(g(s)=W(s)\).
Now \(f=\PP\operatorname{div}(W\otimes W)\): \(\operatorname{div}(W\otimes W)=(W\cdot\nabla)W\in L^2\) (bounded times \(L^2\)),
and \(-\nabla\Delta^{-1}\operatorname{div}\operatorname{div}(W\otimes W)=\nabla\sum R_iR_j(W_iW_j)=\nabla p\) since
\(R_iR_j=-\partial_i\partial_j\Delta^{-1}\).  Integrating \(g'\) from \(s_0\) to \(s\) gives
\eqref{eq:c3v81-duhamel}; absolute convergence by \eqref{eq:c3v81-cO}
with \(\|W\otimes W(\tau)\|_\infty\le C_\infty^2(-\tau)^{-1}\).
\end{proof}

\begin{theorem}[Ancient form]
\label{thm:c3v81-ancient}
For every \(W\in\mathcal F(x_0)\) and \(s<0\),
\begin{equation}
 W(s)=-\int_{-\infty}^se^{\nu(s-\tau)\Delta}\PP\operatorname{div}\bigl(W\otimes W\bigr)(\tau)\,\dd\tau ,
 \label{eq:c3v81-ancient}
\end{equation}
the integral converging absolutely in \(L^\infty\).
\end{theorem}

\begin{proof}
In \eqref{eq:c3v81-duhamel} let \(s_0\to-\infty\): \(\|e^{\nu(s-s_0)\Delta}W(s_0)\|_\infty\le\|W(s_0)\|_\infty\le C_\infty(-s_0)^{-1/2}\to0\),
and the integrand is bounded in \(L^\infty\) by
\(c_OC_\infty^2\nu^{-1/2}(s-\tau)^{-1/2}(-\tau)^{-1}\), integrable on \((-\infty,s)\).
\end{proof}

% =============================================================================
\section{The first explicit floors}
\label{sec:c3v81-floors}
% =============================================================================

\begin{theorem}[Explicit \(\varepsilon\)-regularity floor]
\label{thm:c3v81-epsilon}
Every nonzero \(W\in\mathcal F(x_0)\) has \(\sup_s(-s)^{1/2}\nu^{-1/2}\|W(s)\|_\infty\ge1/(\pi c_O)\).  Hence
every homogeneous type-I point satisfies
\begin{equation}
 \boxed{\quad
 \kappa\ \ge\ \frac1{\pi c_O} .
 \quad}
 \label{eq:c3v81-epsilon}
\end{equation}
\end{theorem}

\begin{proof}
Put \(\kappa_W:=\nu^{-1/2}\sup_s(-s)^{1/2}\|W(s)\|_\infty\le\kappa\), finite.  From
\eqref{eq:c3v81-ancient} and \eqref{eq:c3v81-cO},
\(\|W(s)\|_\infty\le\int_{-\infty}^sc_O(\nu(s-\tau))^{-1/2}\kappa_W^2\nu(-\tau)^{-1}\dd\tau=c_O\kappa_W^2\nu^{1/2}\int_{-\infty}^s(s-\tau)^{-1/2}(-\tau)^{-1}\dd\tau\).
With \(\tau=-(-s)r\), \(r\in[1,\infty)\), the integral is
\((-s)^{-1/2}\int_1^\infty(r-1)^{-1/2}r^{-1}\dd r=(-s)^{-1/2}\int_0^\infty u^{-1/2}(1+u)^{-1}\dd u=\pi(-s)^{-1/2}\).
Hence \((-s)^{1/2}\nu^{-1/2}\|W(s)\|_\infty\le\pi c_O\kappa_W^2\) for every \(s\), i.e.\
\(\kappa_W\le\pi c_O\kappa_W^2\); as \(\kappa_W>0\) for \(W\ne0\), \(\kappa_W\ge1/(\pi c_O)\).  The family
at a homogeneous type-I point contains nonzero elements
(Theorem~\ref{thm:c3v1-ancient}), and \(\kappa\ge\kappa_W\).
\end{proof}

\begin{theorem}[Explicit gradient bound]
\label{thm:c3v81-gradient}
For every \(\theta>0\) and every homogeneous type-I point with
\(4c_O\kappa\,(\theta/(1+\theta))^{1/2}<1\),
\begin{equation}
 \boxed{\quad
 C_\infty'\ \le\ \gamma_1(\kappa,\theta)\,\kappa,\qquad
 \gamma_1(\kappa,\theta):=\frac{c_h\,(\theta(1+\theta))^{-1/2}}{1-4c_O\kappa\,(\theta/(1+\theta))^{1/2}} .
 \quad}
 \label{eq:c3v81-gamma1}
\end{equation}
In particular (\(\theta=1\)) \(C_\infty'\le\frac{c_h/\sqrt2}{1-2\sqrt2c_O\kappa}\kappa\) for \(\kappa<1/(2\sqrt2c_O)\), which is
\eqref{eq:c3v79-gamma1}.
\end{theorem}

\begin{proof}
Fix \(s<0\), \(L:=-s\), \(s_0:=s-\theta L\).  Differentiate \eqref{eq:c3v81-duhamel}:
the gradient of the first term is \((\nabla K_{\nu(s-s_0)})*W(s_0)\), bounded by
\(c_h(\nu\theta L)^{-1/2}\|W(s_0)\|_\infty\le c_h(\nu\theta L)^{-1/2}C_\infty((1+\theta)L)^{-1/2}=c_h\kappa\,(\theta(1+\theta))^{-1/2}L^{-1}\)
(Lemma~\ref{lem:c3v81-heat}).  In the integral, \(\nabla\) commutes with
\(e^{\nu(s-\tau)\Delta}\PP\operatorname{div}\) and falls on \(W\otimes W\) (all are Fourier multipliers,
and \(\nabla(W\otimes W)\) is bounded with \(\operatorname{div}\) in \(L^2\)); \(|\nabla(W\otimes W)|\le2|W||\nabla W|\) in
Frobenius norm, so by \eqref{eq:c3v81-cO} the integral's gradient is
bounded by
\(\int_{s_0}^sc_O(\nu(s-\tau))^{-1/2}\,2C_\infty(-\tau)^{-1/2}\,C_\infty'(-\tau)^{-1}\dd\tau=2c_O\kappa C_\infty'\int_{s_0}^s(s-\tau)^{-1/2}(-\tau)^{-3/2}\dd\tau\).
With \(\tau=-Lr\), \(r\in[1,1+\theta]\), the integral is
\(L^{-1}\int_1^{1+\theta}(r-1)^{-1/2}r^{-3/2}\dd r=L^{-1}\int_0^{\theta^{1/2}}2(1+w^2)^{-3/2}\dd w=2L^{-1}(\theta/(1+\theta))^{1/2}\)
(\(r=1+w^2\)).  Multiplying by \(L\) and taking the supremum over \(s\),
\(C_\infty'\le c_h\kappa(\theta(1+\theta))^{-1/2}+4c_O\kappa(\theta/(1+\theta))^{1/2}C_\infty'\), and \(C_\infty'<\infty\) by the
class, so the last term is absorbed.
\end{proof}

\begin{assessment}[What the two floors are]
\label{ass:c3v81-reading}
\eqref{eq:c3v81-epsilon} is the small-ancient-solution Liouville
argument with its constant written: it is the series' first
\(\varepsilon\)-regularity floor with a value, and it makes (E5)'s \(\varepsilon_{\rm reg}\)
explicit from below in the mild formulation.  \eqref{eq:c3v81-gamma1}
is (E5-small) with explicit constants, valid without a smallness
hypothesis on \(\kappa\) beyond the absorption condition; the free
parameter \(\theta\) trades the initial term against the absorption, and
V83 will optimize it inside the enstrophy floor
\(\kappa_1=\min\{\kappa:\gamma_1(\kappa,\theta)\kappa+\frac{\sqrt3}2\kappa+\frac{\kappa^2}2\ge1\}\).  Both floors are numbers
once \(c_O\) is: V82.  The comparison between \(1/(\pi c_O)\), \(\kappa_1\) and the
loop floors is the numerical table of V83.  Neither floor is an
exclusion, and no global regularity claim is made.
\end{assessment}

% =============================================================================
\section{Integrated V81 conclusion and next acquisition order}
\label{sec:c3v81-conclusion}
% =============================================================================

\begin{theorem}[What V81 proves]
\label{thm:c3v81-conclusion}
Relative to the two legacy cycles and V1--V80: the kernel lemmas
(Lemmas~\ref{lem:c3v81-heat}, \ref{lem:c3v81-oseen}), the mild
formulation on finite windows and in ancient form
(Theorems~\ref{thm:c3v81-duhamel}, \ref{thm:c3v81-ancient}), the explicit
\(\varepsilon\)-regularity floor (Theorem~\ref{thm:c3v81-epsilon}) and the explicit
gradient bound (Theorem~\ref{thm:c3v81-gradient}).  The full
Navier--Stokes stability/global-regularity theorem remains
unproved.
\end{theorem}

\begin{proof}
The cited results.
\end{proof}

\begin{programme}[V82 acquisition order]
\label{prog:c3v82}
\begin{enumerate}[label=\textup{(AC\arabic*)},leftmargin=4.8em]
\item Bound \(c_O\) explicitly: write \(\nabla K^O_1=\nabla K_1\,I+\nabla^3\Phi_1\) in radial
      form (\(\Phi_1=\phi(r)\), \(\partial_i\partial_j\phi=\phi''\hat x_i\hat x_j+\frac{\phi'}r(\delta_{ij}-\hat x_i\hat x_j)\), one more
      derivative), bound the pointwise operator norm of the tensor by
      explicit radial functions, and integrate; give a closed-form or
      rigorously bounded numerical value with the bounding steps
      recorded (elementary bounds on \(\operatorname{erf}\), \(\operatorname{erfc}\) and Gaussian
      moments).
\item Record the resulting numerical values of \(1/(\pi c_O)\) and of
      \(\gamma_1(\kappa,1)\).
\item The next compulsory companion audit is V85.
\end{enumerate}
\end{programme}

\begin{assessment}[Position after V81]
\label{ass:c3v81-position}
The mild formulation is in the class, with an explicit
\(\varepsilon\)-regularity floor and an explicit gradient bound in terms of one
kernel constant.  No full stability or global-regularity claim is
made.
\end{assessment}

% =============================================================================
\section*{Version V81 (third cycle) append-only record}
\addcontentsline{toc}{section}{Version V81 (third cycle) append-only record}
% =============================================================================

The complete V80 source of the third cycle was retained before this
continuation.  Its SHA--256 digest was
\begin{center}\scriptsize\ttfamily
254aabd220ff6a2da7b7de31df9676ecab1792bdfd96908c2ac66f5f4076805a.
\end{center}
V81 established the mild formulation and the first explicit floors.
No full regularity claim is made.

% =============================================================================
\section{V82: the Oseen constant, certified and computed}
\label{sec:c3v82-entry}
% =============================================================================

Items (AC1) and (AC2) of Programme~\ref{prog:c3v82} are carried out.
The gradient of the Oseen kernel is written in closed radial form,
its pointwise operator norm on matrices is computed exactly, its
tail is bounded analytically, and two values of \(c_O\) are recorded:
a certified upper bound \(c_O\le4.208\), proved by a weighted
Cauchy--Schwarz inequality and Plancherel with closed-form Gaussian
moments, and a computed value \(c_O\le2.834\), obtained by quadrature of
the explicit radial integrand with the analytic tail, recorded as a
computation and not as a hand-verified bound.  The resulting values
of the \(\varepsilon\)-regularity floor of V81 are \(\kappa\ge0.0756\) (certified) and
\(\kappa\ge0.1123\) (computed).  No global regularity claim is made.  The
next compulsory companion audit is V85.

% =============================================================================
\section{Radial form of the Oseen gradient}
\label{sec:c3v82-radial}
% =============================================================================

Write \(r=|x|\), \(n=x/r\), \(\phi(r):=\Phi_1(x)=\operatorname{erf}(r/2)/(4\pi r)\), \(E:=\operatorname{erf}(r/2)\),
\(g:=e^{-r^2/4}/\sqrt\pi\) (so \(E'=g\)), and \(h_K:=\frac r2K_1(r)=\frac{r}{2}(4\pi)^{-3/2}e^{-r^2/4}\).

\begin{lemma}[Radial coefficients]
\label{lem:c3v82-radial}
\begin{align}
 \phi'&=\frac{gr-E}{4\pi r^2},\qquad
 \phi''=\frac{2E-2gr-\frac12r^3g}{4\pi r^3},\qquad
 \phi'''=\frac{(\frac14r^5+r^3+6r)g-6E}{4\pi r^4},
 \label{eq:c3v82-phi}\\
 c&:=\frac{\phi''}r-\frac{\phi'}{r^2}=\frac{3(E-gr)-\frac12r^3g}{4\pi r^4},
 \label{eq:c3v82-c}
\end{align}
and, with \(Y_{ijk}:=\delta_{ij}n_k+\delta_{ik}n_j+\delta_{jk}n_i-3n_in_jn_k\) (fully symmetric,
\(|Y|_F^2=6\), \(Y\perp n\otimes n\otimes n\)),
\begin{equation}
 \partial_i\partial_j\partial_k\Phi_1=\phi'''\,n_in_jn_k+c\,Y_{ijk},\qquad
 \partial_k\bigl(K_1\delta_{ij}\bigr)=-h_K\,n_k\delta_{ij} .
 \label{eq:c3v82-tensor}
\end{equation}
All coefficients are smooth on \([0,\infty)\) and vanish at \(r=0\)
(\(\phi'''\), \(c\), \(h_K\) are odd, of order \(r\)).
\end{lemma}

\begin{proof}
Differentiate \(\phi=E/(4\pi r)\) with \(E'=g\), \(E''=-\frac r2g\), \(E'''=(\frac{r^2}4-\frac12)g\).
The third derivative of a radial function is
\(\phi'''n_in_jn_k+(\phi''/r-\phi'/r^2)Y_{ijk}\).  \(|Y|_F^2=15-18+9=6\) from
\(|S|_F^2=15\), \(S\cdot n^{\otimes3}=3\) for \(S=\delta_{ij}n_k+\delta_{ik}n_j+\delta_{jk}n_i\).  Oddness: \(E\) is
odd, \(g\) even.
\end{proof}

\begin{proposition}[Pointwise operator norm]
\label{prop:c3v82-opnorm}
Let \(M(x)\) be the linear map from \(3\times3\) matrices \(f\) to vectors,
\(M(x)f:=\bigl(\partial_j(K^O_1)_{il}(x)f_{lj}\bigr)_i\), so that \(e^{\Delta}\PP\operatorname{div}f=M*f\).  With
\(\alpha:=\phi'''-h_K\),
\begin{equation}
 \|M(x)\|_{F\to2}=\max\Bigl\{\bigl(\alpha^2+2c^2\bigr)^{1/2},\ \Bigl(\frac{(2c-h_K)^2+h_K^2}2\Bigr)^{1/2}\Bigr\},\qquad
 |M(x)|_F^2=\phi'''^2+6c^2+3h_K^2-2\phi'''h_K-4ch_K .
 \label{eq:c3v82-opnorm}
\end{equation}
\end{proposition}

\begin{proof}
By \eqref{eq:c3v82-tensor}, \(Mf=\phi'''\,n\,(n\cdot fn)+c\,(fn+f^{\rm T}n+n\operatorname{tr}f-3n\,(n\cdot fn))-h_Kfn\).
Write \(f=f_s+f_a\) (symmetric and antisymmetric parts), \(\mu:=n\cdot fn\),
\(\tau:=\operatorname{tr}f\), \(f_sn=\mu n+w_s\), \(f_an=w_a\) with \(w_s,w_a\perp n\).  Then
\(fn=\mu n+w_s+w_a\), \(f^{\rm T}n=\mu n+w_s-w_a\), and
\[
 Mf=n\bigl[(\phi'''-h_K)\mu+c\,(\tau-\mu)\bigr]+(2c-h_K)\,w_s-h_K\,w_a .
\]
In an orthonormal basis with \(n\) last, \(|f|_F^2=|f_s|_F^2+|f_a|_F^2\ge\mu^2+\frac12(\tau-\mu)^2+2|w_s|^2+2|w_a|^2\)
(the \(2\times2\) block of \(f_s\) has trace \(\tau-\mu\), and a \(2\times2\) symmetric matrix
has Frobenius norm at least \(|\operatorname{trace}|/\sqrt2\)).  Cauchy--Schwarz gives
\(|Mf|^2\le(\alpha^2+2c^2)(\mu^2+\frac12(\tau-\mu)^2)+\frac12((2c-h_K)^2+h_K^2)(2|w_s|^2+2|w_a|^2)\),
and the two budgets sum to at most \(|f|_F^2\), so the norm is at most
the stated maximum; each value is attained (take \(w_s=w_a=0\) and
\((\mu,\tau-\mu)\) proportional to \((\alpha,2c)\); or \(\mu=\tau=0\) and \(w_s,w_a\) parallel with
lengths proportional to \(|2c-h_K|,|h_K|\) and signs aligned).  The
Frobenius norm follows from \(|n^{\otimes3}|=1\), \(|Y|^2=6\), \(|n_k\delta_{ij}|^2=3\),
\(n^{\otimes3}\cdot Y=0\), \(n^{\otimes3}\cdot(n_k\delta_{ij})=1\), \(Y\cdot(n_k\delta_{ij})=2\).
\end{proof}

\begin{lemma}[Tail]
\label{lem:c3v82-tail}
For \(r\ge12\), \(\|M(x)\|_{F\to2}\le(\sqrt{54}+10^{-9})/(4\pi r^4)\); hence
\(\int_{|x|\ge R}\|M\|_{F\to2}\dd x\le(\sqrt{54}+10^{-9})/R\) for \(R\ge12\).
\end{lemma}

\begin{proof}
With \(E=1-\operatorname{erfc}(r/2)\) and \(\operatorname{erfc}(r/2)\le\frac2rg\) (for \(r>0\)),
\(4\pi r^4\phi'''=-6+\varepsilon_1\), \(0\le\varepsilon_1\le(\frac14r^5+r^3+6r+\frac{12}r)g\), and
\(4\pi r^4c=3-\varepsilon_2\), \(0\le\varepsilon_2\le(\frac{r^3}2+3r+\frac6r)g\); also \(4\pi r^4h_K=\frac{r^5}4g\).
At \(r=12\), \(g=e^{-36}/\sqrt\pi<1.4\cdot10^{-16}\) and the three bounds are below
\(10^{-11}\); they decrease for \(r\ge12\).  Then
\(4\pi r^4(\alpha^2+2c^2)^{1/2}\le((6+10^{-10})^2+2(3+10^{-10})^2)^{1/2}<\sqrt{54}+10^{-9}\), and the
second entry of the maximum is smaller.  Integrate
\(4\pi r^2\cdot(\sqrt{54}+10^{-9})/(4\pi r^4)\) from \(R\) to \(\infty\).
\end{proof}

% =============================================================================
\section{The certified bound}
\label{sec:c3v82-certified}
% =============================================================================

\begin{theorem}[Certified Oseen constant]
\label{thm:c3v82-certified}
\(c_O\le\|\,|M|_F\,\|_{L^1(\R^3)}\le4.208\).
\end{theorem}

\begin{proof}
Let \(T:=\nabla K^O_1\) (the tensor \(\partial_k(K^O_1)_{ij}\)), so \(\|M(x)\|_{F\to2}\le|T(x)|_F\).  For
\(\lambda>0\), Cauchy--Schwarz gives
\(\int|T|_F\le\|(1+\lambda r^2)^{-1}\|_{L^2}\,\|(1+\lambda r^2)T\|_{L^2}\le\pi\lambda^{-3/4}\bigl(\|T\|_{L^2}+\lambda\|r^2T\|_{L^2}\bigr)\),
using \(\int_{\R^3}(1+\lambda|x|^2)^{-2}\dd x=\pi^2\lambda^{-3/2}\).  By Plancherel
(\(\hat f(\xi)=\int fe^{-ix\cdot\xi}\), \(\|f\|_2^2=(2\pi)^{-3}\|\hat f\|_2^2\)) with
\(\hat T_{ijk}=i\xi_k\bigl(\delta_{ij}-\xi_i\xi_j/\rho^2\bigr)e^{-\rho^2}\), \(\rho=|\xi|\):
\(\sum_{ijk}|\hat T_{ijk}|^2=(3-2+1)\rho^2e^{-2\rho^2}=2\rho^2e^{-2\rho^2}\), and \(r^2T\) has transform
\(-\Delta_\xi\hat T\).  With \(F:=\rho^{-2}e^{-\rho^2}\), \(P_{ijk}:=\xi_i\xi_j\xi_k\), \(S_{ijk}:=\delta_{ij}\xi_k+\delta_{ik}\xi_j+\delta_{jk}\xi_i\):
\(\Delta_\xi(PF)=2SF+P(F''+8F'/\rho)\) (Euler's identity \(\xi\cdot\nabla P=3P\), \(\Delta P=2S\)),
\(F''+8F'/\rho=(-10\rho^{-4}-10\rho^{-2}+4)e^{-\rho^2}\), and
\(\Delta_\xi(\xi_ke^{-\rho^2})=\xi_k(4\rho^2-10)e^{-\rho^2}\); with \(|S|^2=15\rho^2\), \(S\cdot P=3\rho^4\),
\(|P|^2=\rho^6\), \(S_{ijk}\delta_{ij}\xi_k=5\rho^2\), \(P_{ijk}\delta_{ij}\xi_k=\rho^4\), a direct expansion
gives
\[
 \sum_{ijk}|\Delta_\xi\hat T_{ijk}|^2=\bigl(40\rho^{-2}+80+168\rho^2-160\rho^4+32\rho^6\bigr)e^{-2\rho^2}.
\]
With \(m_n:=\int_0^\infty\rho^{2n}e^{-2\rho^2}\dd\rho=(2n-1)!!\,(\pi/2)^{1/2}/2^{2n+1}\):
\(\|T\|_2^2=(2\pi)^{-3}4\pi\cdot2m_2=0.011905\), and
\(\|r^2T\|_2^2=(2\pi)^{-3}4\pi(40m_0+80m_1+168m_2-160m_3+32m_4)=2.13101\); hence
\(A:=\|T\|_2=0.10911\), \(B:=\|r^2T\|_2=1.45980\).  The bound
\(\pi\lambda^{-3/4}(A+\lambda B)\) is minimal at \(\lambda=3A/B=0.22423\), where it equals
\(\pi\cdot4\cdot3^{-3/4}A^{1/4}B^{3/4}=4.2078\).
\end{proof}

\begin{remark}[Verification]
\label{rem:c3v82-verification}
The two \(L^2\) norms were also evaluated in real space from
\eqref{eq:c3v82-opnorm} by quadrature on \([10^{-4},10^5]\); they agree with
the closed forms to five digits (\(0.0119051\) and \(2.13093\)), which checks
the Fourier algebra.  The certified bound loses a factor \(1.5\) against
the computed value because the weight \((1+\lambda r^2)^{-1}\) is not adapted
to the \(r^{-4}\) tail; a weight with an \(r^4\) term (transform \(\Delta_\xi^2\)) would
tighten it and is left open.
\end{remark}

% =============================================================================
\section{The computed value}
\label{sec:c3v82-computed}
% =============================================================================

\begin{proposition}[Computed Oseen constant]
\label{prop:c3v82-computed}
\(\int_{|x|\le12}\|M\|_{F\to2}\dd x=2.2214\) by quadrature of the explicit radial
integrand \(4\pi r^2\|M\|_{F\to2}(r)\) of \eqref{eq:c3v82-opnorm} (trapezoidal
rule, \(10^5\) and \(4\cdot10^5\) nodes on \([10^{-4},200]\), agreeing to \(10^{-6}\), the part
beyond \(12\) subtracted with Lemma~\ref{lem:c3v82-tail}); with the tail,
\begin{equation}
 c_O\ \le\ c_O^{\rm num}:=2.834 .
 \label{eq:c3v82-num}
\end{equation}
The pointwise formula \eqref{eq:c3v82-opnorm} was checked against the
numerically assembled \(3\times9\) matrix of the tensor at nine radii.  The
value is recorded as a computation of an explicit integral, not as a
hand-verified bound; the certified bound is
Theorem~\ref{thm:c3v82-certified}.
\end{proposition}

\begin{corollary}[Item (AC2): the floors with values]
\label{cor:c3v82-values}
Every homogeneous type-I point satisfies
\begin{equation}
 \kappa\ \ge\ \frac1{\pi c_O}\ \ge\ 0.0756\ \text{(certified)},\qquad
 \kappa\ \ge\ 0.1123\ \text{(with }c_O^{\rm num}\text{)} ,
 \label{eq:c3v82-floors}
\end{equation}
and \(\gamma_1(\kappa,1)=\frac{(2/\pi)^{1/2}}{1-2\sqrt2c_O\kappa}\) is finite for \(\kappa<0.0840\) (certified) and
\(\kappa<0.1248\) (computed).
\end{corollary}

\begin{proof}
Theorems~\ref{thm:c3v81-epsilon}, \ref{thm:c3v81-gradient} with the two
values; \(c_h/\sqrt2=(2/\pi)^{1/2}\).
\end{proof}

\begin{assessment}[Reading]
\label{ass:c3v82-reading}
The velocity constant of a homogeneous type-I point is at least
\(0.0756\), a proved number, and at least \(0.1123\) if the quadrature is
accepted.  The window in which the gradient bound with \(\theta=1\) applies
lies below these floors for the certified value and barely above
them for the computed one, so the enstrophy floor
\(\kappa_1=\min\{\kappa:\gamma_1(\kappa,\theta)\kappa+\frac{\sqrt3}2\kappa+\frac{\kappa^2}2\ge1\}\) must be evaluated with the
window parameter \(\theta\) optimized; that is V83.  Neither floor is an
exclusion.
\end{assessment}

% =============================================================================
\section{Integrated V82 conclusion and next acquisition order}
\label{sec:c3v82-conclusion}
% =============================================================================

\begin{theorem}[What V82 proves]
\label{thm:c3v82-conclusion}
Relative to the two legacy cycles and V1--V81: the radial form and
the exact pointwise operator norm of the Oseen gradient
(Lemma~\ref{lem:c3v82-radial}, Proposition~\ref{prop:c3v82-opnorm}), the tail
(Lemma~\ref{lem:c3v82-tail}), the certified bound \(c_O\le4.208\)
(Theorem~\ref{thm:c3v82-certified}), the computed value \(c_O\le2.834\)
(Proposition~\ref{prop:c3v82-computed}) and the floors with values
(Corollary~\ref{cor:c3v82-values}).  The full Navier--Stokes
stability/global-regularity theorem remains unproved.
\end{theorem}

\begin{proof}
The cited results.
\end{proof}

\begin{programme}[V83 acquisition order]
\label{prog:c3v83}
\begin{enumerate}[label=\textup{(AC\arabic*)},leftmargin=4.8em]
\item The explicit enstrophy floor: for each of the two values of
      \(c_O\), the least \(\kappa\) with \(\min_\theta\gamma_1(\kappa,\theta)\kappa+\frac{\sqrt3}2\kappa+\frac{\kappa^2}2\ge1\) (the
      minimum over admissible \(\theta\), i.e.\ \(4c_O\kappa(\theta/(1+\theta))^{1/2}<1\)), with
      the optimal \(\theta\) recorded; the comparison with \(1/(\pi c_O)\); the
      table of the two floors with both values.
\item Record whether the gradient bound can be improved by a
      second iteration of the Duhamel formula (inserting the ancient
      form of \(W\) into the nonlinear term), and what it would need.
\item The next compulsory companion audit is V85.
\end{enumerate}
\end{programme}

\begin{assessment}[Position after V82]
\label{ass:c3v82-position}
The Oseen constant is bounded, and the first explicit floor is a
number.  No full stability or global-regularity claim is made.
\end{assessment}

% =============================================================================
\section*{Version V82 (third cycle) append-only record}
\addcontentsline{toc}{section}{Version V82 (third cycle) append-only record}
% =============================================================================

The complete V81 source of the third cycle was retained before this
continuation.  Its SHA--256 digest was
\begin{center}\scriptsize\ttfamily
f3a83402a09a77f11fd259d1552e6790636238dcdd7931f17fd7c6e6e89447a1.
\end{center}
V82 bounded the Oseen constant and gave the first floor a value.  No
full regularity claim is made.

% =============================================================================
\section{V83: the explicit enstrophy floor, and the table of floors with values}
\label{sec:c3v83-entry}
% =============================================================================

Items (AC1) and (AC2) of Programme~\ref{prog:c3v83} are carried out.
The enstrophy constraint of V67 combined with the explicit gradient
bound of V81, with the window parameter optimized, gives the floor
\(\kappa\ge0.1129\) (certified, with \(c_O\le4.208\)) and \(\kappa\ge0.1386\) (with the
computed \(c_O\le2.834\)) at every homogeneous type-I point; both exceed
the Liouville floors of V81--V82, so the enstrophy row raises the
explicit \(\varepsilon\)-regularity floor of the mild formulation by a factor
\(1.49\) (certified) and \(1.23\) (computed).  A second Liouville floor from
the gradient bound alone, \(\kappa\ge1/(4c_O)\), is recorded and is the
weakest.  The iteration of the Duhamel formula does not improve the
gradient bound; what would is a smaller constant than \(c_O\) for the
specific tensor \(W\otimes W\).  The table of floors is written.  No global
regularity claim is made.  The next compulsory companion audit is
V85.

% =============================================================================
\section{The gradient Liouville floor}
\label{sec:c3v83-liouville}
% =============================================================================

\begin{proposition}[Gradient Liouville floor]
\label{prop:c3v83-gradient-liouville}
Every homogeneous type-I point satisfies \(\kappa\ge1/(4c_O)\), i.e.\ \(\kappa\ge0.0594\)
(certified) and \(\kappa\ge0.0882\) (computed).
\end{proposition}

\begin{proof}
If \(4c_O\kappa<1\), the absorption of Theorem~\ref{thm:c3v81-gradient} holds
for every \(\theta>0\), and letting \(\theta\to\infty\) in \eqref{eq:c3v81-gamma1} gives
\(C_\infty'\le\gamma_1(\kappa,\theta)\kappa\to0\): \(\nabla W\equiv0\) for every element, so each \(W(s)\) is
constant in space, hence zero (bounded by \(C_\infty(-s)^{-1/2}\) and in \(L^6\)),
contradicting the existence of nonzero elements.
\end{proof}

\begin{remark}[The two Liouville floors]
\label{rem:c3v83-two}
The velocity form (Theorem~\ref{thm:c3v81-epsilon}) gives \(1/(\pi c_O)\), the
gradient form \(1/(4c_O)\); the former is larger because the ancient
integral \(\int_{-\infty}^s(s-\tau)^{-1/2}(-\tau)^{-1}\dd\tau=\pi(-s)^{-1/2}\) has the constant \(\pi\)
against the gradient form's \(2\cdot2\).  Both are the classical
small-ancient-solution argument with explicit constants.
\end{remark}

% =============================================================================
\section{The explicit enstrophy floor}
\label{sec:c3v83-enstrophy}
% =============================================================================

For \(\kappa>0\) let \(\Theta(\kappa):=\{\theta>0:4c_O\kappa(\theta/(1+\theta))^{1/2}<1\}\) (all \(\theta>0\) if \(4c_O\kappa<1\), an
interval \((0,\theta_{\max}(\kappa))\) otherwise) and
\begin{equation}
 \Gamma(\kappa):=\inf_{\theta\in\Theta(\kappa)}\gamma_1(\kappa,\theta)\,\kappa+\frac{\sqrt3}2\kappa+\frac{\kappa^2}2 ,\qquad
 \gamma_1(\kappa,\theta)=\frac{c_h(\theta(1+\theta))^{-1/2}}{1-4c_O\kappa(\theta/(1+\theta))^{1/2}} .
 \label{eq:c3v83-Gamma}
\end{equation}

\begin{theorem}[Explicit enstrophy floor]
\label{thm:c3v83-floor}
\(\Gamma\) is increasing on \((0,\infty)\), and every homogeneous type-I point
satisfies \(\Gamma(\kappa)\ge1\), hence \(\kappa\ge\kappa_1^{\rm expl}:=\min\{\kappa:\Gamma(\kappa)\ge1\}\).  With
\(c_h=2/\sqrt\pi\):
\begin{equation}
 \boxed{\quad
 \kappa\ \ge\ 0.1129\ \ (c_O=4.208,\ \theta^*=0.088),\qquad
 \kappa\ \ge\ 0.1386\ \ (c_O=2.834,\ \theta^*=0.148),
 \quad}
 \label{eq:c3v83-values}
\end{equation}
where \(\theta^*\) is the minimizing window parameter at the floor.
\end{theorem}

\begin{proof}
\eqref{eq:c3v67-constraint} and Theorem~\ref{thm:c3v81-gradient}: for
every \(\theta\in\Theta(\kappa)\), \(1\le C_\infty'+\frac{\sqrt3}2\kappa+\frac{\kappa^2}2\le\gamma_1(\kappa,\theta)\kappa+\frac{\sqrt3}2\kappa+\frac{\kappa^2}2\); take the
infimum.  Monotonicity: for fixed \(\theta\), \(\gamma_1(\kappa,\theta)\kappa\) is increasing in \(\kappa\)
on its domain and \(\Theta(\kappa)\) shrinks as \(\kappa\) grows, so the infimum is
increasing.  The values: bisection on \(\kappa\) with the infimum over \(\theta\)
on a geometric grid of \(2\cdot10^4\) points in \([10^{-4},10^2]\); at the floor,
\(\gamma_1\kappa=0.896\) and \(0.870\) respectively, the remaining \(0.104\) and \(0.130\)
being \(\frac{\sqrt3}2\kappa+\frac{\kappa^2}2\).  (The floor is a computed root of an
explicit monotone function; its first three digits are stable under
refinement of both grids.)
\end{proof}

\begin{remark}[Structure of the optimum]
\label{rem:c3v83-optimum}
At the floor the optimal window is short (\(\theta^*\approx0.09\) to \(0.15\) of the
current time-to-blow-up), because the absorption coefficient
\(4c_O\kappa(\theta/(1+\theta))^{1/2}\) is close to one for long windows at these
values of \(\kappa\) (\(4c_O\kappa=1.90\) and \(1.57\) at the two floors), while the
initial term costs only \((\theta(1+\theta))^{-1/2}\).  The enstrophy constraint is
then satisfied with the gradient term contributing \(0.87\) to \(0.90\) of
the unit, i.e.\ at the floor \(C_\infty'\) is allowed to be of order \(6\) to \(8\)
times \(\kappa\).
\end{remark}

% =============================================================================
\section{Item (AC2): iteration does not help}
\label{sec:c3v83-iteration}
% =============================================================================

\begin{assessment}[Iterating the Duhamel formula]
\label{ass:c3v83-iteration}
Inserting the ancient form \eqref{eq:c3v81-ancient} of \(W\) into the
nonlinear term of \eqref{eq:c3v81-duhamel} and bounding again by the
class reproduces the same two constants: the velocity bound
\(\kappa_W\le\pi c_O\kappa_W^2\) and the gradient bound \eqref{eq:c3v81-gamma1} are
already the fixed points of the iteration, and a further insertion
only replaces \(\kappa_W\) by \(\pi c_O\kappa_W^2\), which is not smaller unless
\(\kappa_W<1/(\pi c_O)\), the excluded case.  The constants improve only if
the kernel action on the specific tensor \(W\otimes W\) is bounded better
than by \(c_O\|W\otimes W\|_\infty\): (i) using \(\operatorname{div}W=0\), for which the trace part
of \(W\otimes W\) (which the pressure removes) could be handled separately;
(ii) using the \(L^6\) bound \(K_6\) on the far part of the kernel, where
\(\|\nabla K^O_t\|_{L^{6/5}}\) replaces the \(L^1\) norm on \(|x-y|>R\); (iii) using the
Gaussian energy bound \(C\) on \(\nabla W\) through \(L^2\)--\(L^\infty\) splittings.  Each
gives an explicit constant in terms of the class constants and
would lower \(c_O\) in the floors; none is pursued here, the floors being
recorded at the level of the pure sup-norm argument.
\end{assessment}

% =============================================================================
\section{The table of floors with values}
\label{sec:c3v83-table}
% =============================================================================

\begin{theorem}[Explicit floors on the velocity constant]
\label{thm:c3v83-table}
Every homogeneous type-I point satisfies each of the following, with
the certified value proved and the computed value depending on the
quadrature of Proposition~\ref{prop:c3v82-computed}:

\medskip
\begin{center}\small
\begin{tabular}{@{}p{0.34\textwidth}p{0.2\textwidth}p{0.17\textwidth}p{0.17\textwidth}@{}}
\toprule
\textbf{Floor} & \textbf{Formula} & \textbf{Certified} & \textbf{Computed}\\
\midrule
Gradient Liouville (V83) & \(1/(4c_O)\) & \(0.0594\) & \(0.0882\)\\
Velocity Liouville (V81) & \(1/(\pi c_O)\) & \(0.0756\) & \(0.1123\)\\
Enstrophy row with explicit gradient (V67, V81, V83) & \(\min\{\kappa:\Gamma(\kappa)\ge1\}\) & \(0.1129\) & \(0.1386\)\\
\bottomrule
\end{tabular}
\end{center}
\medskip

\noindent
The operative floor is the enstrophy floor; the loop floors of
V76--V77 need the second-derivative constant (V84).
\end{theorem}

\begin{proof}
Proposition~\ref{prop:c3v83-gradient-liouville},
Corollary~\ref{cor:c3v82-values}, Theorem~\ref{thm:c3v83-floor}.
\end{proof}

\begin{assessment}[Reading]
\label{ass:c3v83-reading}
The series now states, with a proved number: at a homogeneous type-I
singular point of a smooth periodic solution, the rescaled velocity
bound \(C_\infty(T_*-t)^{-1/2}\) has \(C_\infty\ge0.1129\,\nu^{1/2}\).  It is a quantitative
\(\varepsilon\)-regularity statement in the mild formulation, of the kind known
in the literature with various constants, and its interest for the
series is twofold: the constant is explicit and small enough to be
compared with the second cycle's constant-chasing thresholds
(\(\kappa<\frac12\), \(\kappa<\sqrt2\)), which it does not empty; and the enstrophy row,
the cycle's own contribution, improves the Liouville floor by a
factor \(1.5\), showing that the extremal method adds to the mild
formulation rather than duplicating it.  It is not an exclusion.
\end{assessment}

% =============================================================================
\section{Integrated V83 conclusion and next acquisition order}
\label{sec:c3v83-conclusion}
% =============================================================================

\begin{theorem}[What V83 proves]
\label{thm:c3v83-conclusion}
Relative to the two legacy cycles and V1--V82: the gradient
Liouville floor (Proposition~\ref{prop:c3v83-gradient-liouville}), the
explicit enstrophy floor (Theorem~\ref{thm:c3v83-floor}), the iteration
record (Assessment~\ref{ass:c3v83-iteration}) and the table
(Theorem~\ref{thm:c3v83-table}).  The full Navier--Stokes
stability/global-regularity theorem remains unproved.
\end{theorem}

\begin{proof}
The cited results.
\end{proof}

\begin{programme}[V84 acquisition order]
\label{prog:c3v84}
\begin{enumerate}[label=\textup{(AC\arabic*)},leftmargin=4.8em]
\item The second-derivative constant: differentiate
      \eqref{eq:c3v81-duhamel} twice, with \(\|\nabla^2K_t\|_{L^1}=c_{h,2}/t\)
      (\(c_{h,2}\) a Gaussian moment, computed) and
      \(|\nabla^2(W\otimes W)|\le2|W||\nabla^2W|+2|\nabla W|^2\); with
      \(N_2:=\sup(-s)^{3/2}\|\nabla^2W(s)\|_\infty=\nu^{1/2}C_\infty''\), derive
      \(N_2\le c_{h,2}\kappa\,\theta^{-1}(1+\theta)^{-1/2}+4c_O\kappa(\theta/(1+\theta))^{1/2}N_2+c_O\,C_\infty'^2\,I_2(\theta)\) with
      \(I_2\) the explicit window integral, hence \(\gamma_2(\kappa,\theta)\) with
      \(\nu^{1/2}C_\infty''\le\gamma_2\kappa\).
\item Evaluate \(q_{\mathcal P}=2\nu(C_\infty''/\theta_{\mathcal P})^2\) at the enstrophy floor with the
      explicit constants, and record whether the loop constraints
      of V74 and V77 apply there (\(\theta_{\mathcal P}>0\)) and what floors they
      give; the table of V83 extended.
\item The next compulsory companion audit is V85.
\end{enumerate}
\end{programme}

\begin{assessment}[Position after V83]
\label{ass:c3v83-position}
The enstrophy floor is a number, \(0.1129\), certified.  No full
stability or global-regularity claim is made.
\end{assessment}

% =============================================================================
\section*{Version V83 (third cycle) append-only record}
\addcontentsline{toc}{section}{Version V83 (third cycle) append-only record}
% =============================================================================

The complete V82 source of the third cycle was retained before this
continuation.  Its SHA--256 digest was
\begin{center}\scriptsize\ttfamily
afe7bf8aa51cdebeb0538f9df1d4f55633d482377399f2f27e8319b8d4244356.
\end{center}
V83 gave the enstrophy floor its value and wrote the table of
floors.  No full regularity claim is made.

% =============================================================================
\section{V84: the second-derivative constant, and which rows apply at the floors}
\label{sec:c3v84-entry}
% =============================================================================

Items (AC1) and (AC2) of Programme~\ref{prog:c3v84} are carried out.
The second derivative of the Duhamel formula gives
\(\nu^{1/2}C_\infty''\le\gamma_2(\kappa)\kappa\) with \(\gamma_2\) explicit in \(c_O\), the Gaussian moment
\(c_{h,2}\) of the second derivative of the heat kernel (certified \(\le1.794\)
in closed form, computed \(1.415\)) and the explicit \(\gamma_1\); at the floors,
\(\nu^{1/2}C_\infty''\le11.65\) (certified) and \(\le8.08\) (computed).  Which rows of the
vorticity ledger apply at the floors is then decided: the
palinstrophy row is void at every homogeneous type-I point with
\(\kappa<0.4569\), unconditionally, by the enstrophy constraint and the
gradient floor of V73 together, so the loop constraints of V74 and
V77 contribute nothing to the explicit floors; the concentration
row applies at the floors with explicit constants.  The table of V83
is extended.  No global regularity claim is made.  The next
compulsory companion audit is V85.

% =============================================================================
\section{The second derivative of the Duhamel formula}
\label{sec:c3v84-second}
% =============================================================================

\begin{lemma}[Second-derivative heat constant]
\label{lem:c3v84-heat2}
\(\nabla^2K_t=\frac1{2t}(\zeta\zeta^{\rm T}-I)K_t\) with \(\zeta=x/(2t)^{1/2}\), so
\(\int|\nabla^2K_t|_F\dd x=c_{h,2}/t\) with
\(c_{h,2}=\frac12\,\mathbb E\bigl(|\zeta|^4-2|\zeta|^2+3\bigr)^{1/2}\), \(\zeta\sim N(0,I_3)\); and
\begin{equation}
 c_{h,2}\ \le\ \frac12\bigl(\mathbb E|\,|\zeta|^2-1\,|+\sqrt2\bigr)
 =\frac12\Bigl(2+2\Bigl(\frac2\pi\Bigr)^{1/2}\Bigl(3e^{-1/2}-(2\pi)^{1/2}\operatorname{erf}(2^{-1/2})\Bigr)+\sqrt2\Bigr)=1.7936 ,
 \label{eq:c3v84-ch2}
\end{equation}
while the quadrature of the exact expectation gives \(c_{h,2}=1.4149\).
Consequently \(\|\nabla^2e^{\nu t\Delta}g\|_\infty\le c_{h,2}(\nu t)^{-1}\|g\|_\infty\) for bounded \(g\), in
Frobenius norm.
\end{lemma}

\begin{proof}
\(\partial_i\partial_jK_t=(x_ix_j/(4t^2)-\delta_{ij}/(2t))K_t\); substitute \(x=(2t)^{1/2}\zeta\), and
\(K_t\) is the density of \((2t)^{1/2}\zeta\).  \(|\zeta\zeta^{\rm T}-I|_F^2=|\zeta|^4-2|\zeta|^2+3=(|\zeta|^2-1)^2+2\),
and \((a^2+2)^{1/2}\le|a|+\sqrt2\).  With \(\rho=|\zeta|\) of density \((2/\pi)^{1/2}\rho^2e^{-\rho^2/2}\):
\(\mathbb E|\rho^2-1|=\mathbb E(\rho^2-1)+2\mathbb E[(1-\rho^2)\mathbf 1_{\rho<1}]=2+2(2/\pi)^{1/2}\int_0^1(\rho^2-\rho^4)e^{-\rho^2/2}\dd\rho\),
and \(\int_0^1\rho^2e^{-\rho^2/2}=-e^{-1/2}+(\pi/2)^{1/2}\operatorname{erf}(2^{-1/2})\),
\(\int_0^1\rho^4e^{-\rho^2/2}=-4e^{-1/2}+3(\pi/2)^{1/2}\operatorname{erf}(2^{-1/2})\), whose difference is
\(3e^{-1/2}-2(\pi/2)^{1/2}\operatorname{erf}(2^{-1/2})=3e^{-1/2}-(2\pi)^{1/2}\operatorname{erf}(2^{-1/2})\).  Numerically
\(\mathbb E|\rho^2-1|=2.1729\).
\end{proof}

\begin{lemma}[Window integral]
\label{lem:c3v84-J2}
\(J_2(\theta):=\int_1^{1+\theta}(r-1)^{-1/2}r^{-2}\dd r=\frac{\theta^{1/2}}{1+\theta}+\arctan\theta^{1/2}\), increasing from \(0\)
to \(\pi/2\).
\end{lemma}

\begin{proof}
\(r=1+w^2\): \(\int_0^{\theta^{1/2}}2(1+w^2)^{-2}\dd w=[w/(1+w^2)+\arctan w]_0^{\theta^{1/2}}\).
\end{proof}

\begin{theorem}[Explicit second-derivative bound]
\label{thm:c3v84-gamma2}
Let \(\gamma_1(\kappa):=\inf_{\theta\in\Theta(\kappa)}\gamma_1(\kappa,\theta)\) (V83).  For every \(\theta>0\) with
\(2c_O\kappa J_2(\theta)<1\), every homogeneous type-I point satisfies
\begin{equation}
 \boxed{\quad
 \nu^{1/2}C_\infty''\ \le\ \gamma_2(\kappa,\theta)\,\kappa,\qquad
 \gamma_2(\kappa,\theta):=\frac{c_{h,2}\,\theta^{-1}(1+\theta)^{-1/2}+2c_O\,\gamma_1(\kappa)^2\,\kappa\,J_2(\theta)}{1-2c_O\kappa J_2(\theta)} ,
 \quad}
 \label{eq:c3v84-gamma2}
\end{equation}
and \(\gamma_2(\kappa):=\inf_\theta\gamma_2(\kappa,\theta)\).  At the floors of V83:
\begin{equation}
 \nu^{1/2}C_\infty''\le11.65\ \ (\text{certified: }c_O=4.208,\ c_{h,2}=1.794,\ \kappa=0.1129,\ \theta=0.054),\qquad
 \nu^{1/2}C_\infty''\le8.08\ \ (\text{computed: }c_O=2.834,\ c_{h,2}=1.415,\ \kappa=0.1386,\ \theta=0.077).
 \label{eq:c3v84-values}
\end{equation}
\end{theorem}

\begin{proof}
Differentiate \eqref{eq:c3v81-duhamel} twice; \(\nabla^2\) falls on the heat
kernel in the first term and on \(W\otimes W\) in the second.  With \(L=-s\),
\(s_0=s-\theta L\), \(N_2:=\sup_s(-s)^{3/2}\|\nabla^2W(s)\|_\infty=C_\infty''\) (profile units, dimension
\(\nu^{-1/2}\)): the first term is at most
\(L^{3/2}c_{h,2}(\nu\theta L)^{-1}C_\infty((1+\theta)L)^{-1/2}=c_{h,2}\kappa\nu^{-1/2}\theta^{-1}(1+\theta)^{-1/2}\).  For the
second, the operator \(e^{\nu(s-\tau)\Delta}\PP\operatorname{div}\) acts on the \((l,j)\) indices of
\(f^{(ab)}_{lj}:=\partial_a\partial_b(W_lW_j)\) for each \((a,b)\); by Cauchy--Schwarz in the
measure \(\|M(y)\|\dd y\), \(\sum_{ab}|(M_t*f^{(ab)})(x)|^2\le(\int\|M_t\|)^2\sup|\nabla^2(W\otimes W)|_F^2\), so
the bound \eqref{eq:c3v81-cO} holds for the Frobenius norm over all
indices; and \(|\nabla^2(W\otimes W)|_F\le2|W||\nabla^2W|+2|\nabla W|^2\).  Hence the second term
is at most
\(L^{3/2}\int_{s_0}^sc_O(\nu(s-\tau))^{-1/2}\bigl(2C_\infty(-\tau)^{-1/2}N_2(-\tau)^{-3/2}+2C_\infty'^2(-\tau)^{-2}\bigr)\dd\tau\),
and with \(\tau=-Lr\), \(\int_{s_0}^s(s-\tau)^{-1/2}(-\tau)^{-2}\dd\tau=L^{-3/2}J_2(\theta)\), this equals
\(2c_O\kappa J_2(\theta)N_2+2c_O\nu^{-1/2}C_\infty'^2J_2(\theta)\).  Multiply by \(\nu^{1/2}\), use
\(C_\infty'\le\gamma_1(\kappa)\kappa\), take the supremum over \(s\) and absorb (\(N_2<\infty\) by the
class).  The values: minimization over \(\theta\) on the grid of V83.
\end{proof}

\begin{remark}[Absorption range]
\label{rem:c3v84-range}
Since \(J_2\to\pi/2\), the absorption holds for every \(\theta\) iff \(\pi c_O\kappa<1\), the
regime excluded by Theorem~\ref{thm:c3v81-epsilon}; at the floors
\(2c_O\kappa\pi=2.98\) and \(2.47\), so the window must again be short
(\(\theta\approx0.05\) to \(0.08\)) and the second-derivative constant is large,
\(\gamma_2\approx100\) and \(58\), the first term \(c_{h,2}/(\theta(1+\theta)^{1/2})\) dominating.
\end{remark}

% =============================================================================
\section{Which rows apply at the floors}
\label{sec:c3v84-rows}
% =============================================================================

\begin{theorem}[The palinstrophy row is void below \(0.4569\)]
\label{thm:c3v84-void}
At every homogeneous type-I point,
\(\theta_{\mathcal P}=\frac32-2C_\infty'-\frac{\sqrt3}2\kappa-\frac{\kappa^2}2\le-\frac12+\frac{\sqrt3}2\kappa+\frac{\kappa^2}2\), which is negative for
\(\kappa<\kappa_0':=\frac{\sqrt7-\sqrt3}2=0.4569\).  Hence the ratio bound of V72, the loop
constraints of V74 and V77 and the floors \(\kappa_2,\kappa_3,\kappa_4\) of V76--V77 are
void at every homogeneous type-I point with \(\kappa<0.4569\), in particular
at the explicit floors of V83, whatever the derivative constants.
\end{theorem}

\begin{proof}
Proposition~\ref{prop:c3v73-gradient-floor}: \(C_\infty'\ge1-\frac{\sqrt3}2\kappa-\frac{\kappa^2}2\) for
\(\kappa<\kappa_0\); substitute.  The root of \(-\frac12+\frac{\sqrt3}2\kappa+\frac{\kappa^2}2\) is
\(-\frac{\sqrt3}2+(\frac34+1)^{1/2}\).
\end{proof}

\begin{remark}[Consistency with V73]
\label{rem:c3v84-consistency}
Theorem~\ref{thm:c3v73-window}(1) found that the ratio row applies at
\(\kappa_1(\gamma_1)\) iff \(\gamma_1<(\sqrt3+\sqrt7)/4\), i.e.\ iff \(\kappa_1>0.4577\); the present
statement is its form without (E5-small).  The loop block's floors
were therefore always confined to \(\kappa\ge0.457\), and the slivers of
V76--V77 lie there.
\end{remark}

\begin{proposition}[The concentration row at the floors]
\label{prop:c3v84-concentration}
At the floors, \(\theta_1\ge\frac32-\gamma_1(\kappa)\kappa-\frac{\sqrt3}2\kappa-\frac\kappa{\sqrt2}-\kappa^2\) equals \(0.414\) (certified) and
\(0.392\) (computed), so the concentration bound \eqref{eq:c3v71-sup}
applies with explicit constants:
\begin{equation}
 \sup_{\mathcal F(x_0)}\int|z|^2|\Omega|^2G\ \le\ 16.9\,\nu C+1313\,\nu|G|\ \ (\text{certified}),\qquad
 \le\ 18.5\,\nu C+665\,\nu|G|\ \ (\text{computed}),
 \label{eq:c3v84-concentration}
\end{equation}
with \(|G|=(4\pi\nu)^{3/2}\) and \(C\) the enstrophy bound of the class.  The bound
is a statement about the Gaussian second moment of the vorticity
along the family at points with \(\kappa\) at the floor; for larger \(\kappa\) it
holds with \(\theta_1\) and \(C_\infty''\) re-evaluated, as long as \(\theta_1>0\).
\end{proposition}

\begin{proof}
\eqref{eq:c3v71-sup} with \(\Xi_1=3.496\), \(3.617\), \(\theta_1\) as stated, and
\(4\nu^2C_\infty''^2|G|=4(\nu^{1/2}C_\infty'')^2\nu|G|\) with \eqref{eq:c3v84-values}.
\end{proof}

\begin{theorem}[Item (AC2): the extended table]
\label{thm:c3v84-table}
Every homogeneous type-I point satisfies:

\medskip
\begin{center}\small
\begin{tabular}{@{}p{0.4\textwidth}p{0.14\textwidth}p{0.14\textwidth}p{0.2\textwidth}@{}}
\toprule
\textbf{Statement} & \textbf{Certified} & \textbf{Computed} & \textbf{Source}\\
\midrule
\(\kappa\ge1/(4c_O)\) & \(0.0594\) & \(0.0882\) & V83\\
\(\kappa\ge1/(\pi c_O)\) & \(0.0756\) & \(0.1123\) & V81--V82\\
\(\kappa\ge\kappa_1^{\rm expl}\) (enstrophy row) & \(0.1129\) & \(0.1386\) & V67, V81, V83\\
\(C_\infty'\le\gamma_1(\kappa)\kappa\) at the floor & \(0.896\) & \(0.870\) & V81, V83\\
\(\nu^{1/2}C_\infty''\le\gamma_2(\kappa)\kappa\) at the floor & \(11.65\) & \(8.08\) & V84\\
palinstrophy row, loops (V72, V74, V77) & void for \(\kappa<0.4569\) & same & V84\\
concentration row (V71) at the floor & applies, \(\theta_1\ge0.414\) & \(\theta_1\ge0.392\) & V84\\
\bottomrule
\end{tabular}
\end{center}
\medskip
\end{theorem}

\begin{proof}
The cited results.
\end{proof}

\begin{assessment}[Reading]
\label{ass:c3v84-reading}
With the explicit constants, the vorticity ledger's contribution to
the floor is exactly the enstrophy row; the loops, which needed the
palinstrophy row, are structurally void below \(0.457\), a fact that
follows from the enstrophy row itself (a small velocity constant
forces a gradient constant near one, which kills the palinstrophy
damping).  The explicit-constant block has thus produced: a certified
floor \(0.1129\) on the velocity constant at a homogeneous type-I point,
explicit derivative bounds at the floor, and the concentration of
the vorticity's Gaussian second moment.  The remaining versions of
the block (V86--V89) may improve \(c_O\) for the specific tensor along
the lines of Assessment~\ref{ass:c3v83-iteration}, which would raise
every number in the table, or turn to the closing of the cycle.
\end{assessment}

% =============================================================================
\section{Integrated V84 conclusion and next acquisition order}
\label{sec:c3v84-conclusion}
% =============================================================================

\begin{theorem}[What V84 proves]
\label{thm:c3v84-conclusion}
Relative to the two legacy cycles and V1--V83: the second-derivative
heat constant and the window integral (Lemmas~\ref{lem:c3v84-heat2},
\ref{lem:c3v84-J2}), the explicit second-derivative bound
(Theorem~\ref{thm:c3v84-gamma2}), the voidness of the palinstrophy row
below \(0.4569\) (Theorem~\ref{thm:c3v84-void}), the concentration row at
the floors (Proposition~\ref{prop:c3v84-concentration}) and the extended
table (Theorem~\ref{thm:c3v84-table}).  The full Navier--Stokes
stability/global-regularity theorem remains unproved.
\end{theorem}

\begin{proof}
The cited results.
\end{proof}

\begin{programme}[V85 acquisition order]
\label{prog:c3v85}
\begin{enumerate}[label=\textup{(AC\arabic*)},leftmargin=4.8em]
\item The compulsory companion audit: read the current SM and YM
      notes and record their digests and decision.
\item Record the position after eighty-five versions; decide
      between improving \(c_O\) for the tensor \(W\otimes W\) (V86--V89) and
      beginning the closing of the cycle; delivery.
\end{enumerate}
\end{programme}

\begin{assessment}[Position after V84]
\label{ass:c3v84-position}
The explicit-constant block is complete at the level of the
sup-norm mild formulation.  No full stability or global-regularity
claim is made.
\end{assessment}

% =============================================================================
\section*{Version V84 (third cycle) append-only record}
\addcontentsline{toc}{section}{Version V84 (third cycle) append-only record}
% =============================================================================

The complete V83 source of the third cycle was retained before this
continuation.  Its SHA--256 digest was
\begin{center}\scriptsize\ttfamily
28dd8ff205a2d211a76496e833bc4dee4df49ed944c70865ae238b56fa098acb.
\end{center}
V84 bounded the second-derivative constant and decided which rows
apply at the floors.  No full regularity claim is made.

% =============================================================================
\section{V85: compulsory companion audit and the position after eighty-five versions}
\label{sec:c3v85-entry}
% =============================================================================

This is the eighty-fifth version of the third cycle.  The compulsory
review of the two companion programmes is performed first, and the
position after eighty-five versions is recorded with the decision
for the block V86--V89.  No global regularity claim is made.  The
next compulsory companion audit is V90.

% =============================================================================
\section{Compulsory V85 review of the current companion notes}
\label{sec:c3v85-companion-audit}
% =============================================================================

The current companion files in \texttt{latex\_notes} were inspected
read-only.  Both programmes have moved since the V80 audit: the SM
programme from V45 to V49 of its seventh series, the YM programme
from V22 to V25 of its sixth series.  The frozen checkpoints are
\begin{align}
 &\texttt{Renewal\_SM\_Einstein\_Continuum\_Research\_Notes\_V49.tex},
 \notag\\[-2pt]
 &\qquad
 \texttt{32cddcf742bb53f6aa195c9e57d5571b861ee7b370883dbd02ce5a5486dc92f3},
 \label{eq:c3v85-SM-checkpoint}\\
 &\texttt{Renewal\_Yang\_Mills\_Mass\_Gap\_Research\_Notes\_V25.tex},
 \notag\\[-2pt]
 &\qquad
 \texttt{8d2598532555a9239fea347533e57ae802f43d4dec214f065bb8f91539d8fb08}.
 \label{eq:c3v85-YM-checkpoint}
\end{align}
(SM V49 has 19476 lines; YM V25 has 10996 lines.  The folder also
held YM V24, a distinct earlier state of 10504 lines with digest
\texttt{f0852058\dots}, superseded by V25.)

\emph{SM V46--V49.}  The SM programme completed its collar
constructions stratum by stratum: a conservative completion of the
codimension-two fan (joint transverse constant subspace, joint
completion and gap, transverse first moment, covariant lift along a
curved hinge), the codimension-three edge completion (three-normal
coefficient field, matching of the lower strata, three-normal
adiabatic gap, first moment and continuum safety), the
original-vertex four-normal completion (final simplicial link,
attachment of boundary strata, moment scaling), and a global
multiscale stratified gap (one global collar hierarchy, a uniform
weighted row ledger, a partition at the intermediate scale, the
global dimensionless gap and smooth-core equivalence).

\emph{YM V23--V25.}  The YM programme proved a compact Gibbs
concentration inequality and a staple-bifurcation incidence card
ruling out a uniform single-link form of its condition WCIC, with an
exact one-plaquette Wilson no-go; embedded the Wilson
low-temperature influence obstruction in exact parity-tree
coordinates and upgraded an earlier bare-bound failure to a
worst-case failure; and, after its own audit (which recorded SM V47
and this series' V81), proved sector-weight and within-sector
total-variation bounds, a conditional fixed-path product bound, a
subcritical path estimate, covariance decay and an
Osterwalder--Schrader mass rate, and identified the complete
condition needed from a repairing block.

\begin{theorem}[V85 companion-audit decision]
\label{thm:c3v85-companion-decision}
Nothing is imported from SM V49 or YM V25.  Neither bears on the
third cycle's chapters, on the vorticity ledger, on the loop method
or on the explicit-constant block.
\end{theorem}

\begin{proof}
The SM results concern stratified collar completions and a global
gap of lattice Dirac interfaces; the YM results concern Gibbs
concentration, incidence obstructions and conditional path
estimates of a lattice gauge theory.  None is a Navier--Stokes
statement, and none concerns the mild formulation, the Oseen
constant, the enstrophy row or the explicit floors.  The YM audit's
record of this series' V81 is a procedural cross-check only.
\end{proof}

% =============================================================================
\section{Position after eighty-five versions, and the block V86--V89}
\label{sec:c3v85-position}
% =============================================================================

\begin{assessment}[Position after eighty-five versions of the third cycle]
\label{ass:c3v85-position-eightyfive}
Fifteen chapters are written: those of
Assessment~\ref{ass:c3v80-position-eighty} and the explicit-constant
block (V81--V84).  The block's outcome is a proved number: every
homogeneous type-I point of a smooth periodic solution has velocity
constant \(\kappa\ge0.1129\) (Theorem~\ref{thm:c3v83-floor}, certified), with the
Liouville floors \(0.0756\) and \(0.0594\) below it, explicit derivative
bounds at the floor, and the concentration of the vorticity's
Gaussian second moment; the loop rows are structurally void below
\(0.4569\) (Theorem~\ref{thm:c3v84-void}).  Every number rests on the
Oseen constant \(c_O\), certified at \(4.208\) and computed at \(2.834\).
\end{assessment}

\begin{assessment}[Decision for V86--V89: the traceless reduction]
\label{ass:c3v85-decision}
Of the three improvements of Assessment~\ref{ass:c3v83-iteration}, the
first is elementary and exact: \(\PP\) annihilates gradients, so
\(e^{t\Delta}\PP\operatorname{div}(W\otimes W)=e^{t\Delta}\PP\operatorname{div}(W\otimes W-\frac13|W|^2I)\), and the traceless part of
\(W\otimes W\) has Frobenius norm \((2/3)^{1/2}|W|^2\); likewise the trace part of
\(\partial_k(W\otimes W)\) is \(\frac13\partial_k|W|^2\,I\) and is annihilated, and the traceless
part of \(a\otimes w+w\otimes a\) has norm at most \((8/3)^{1/2}|a||w|\) against \(2|a||w|\).
Both absorptions of V81 therefore hold with \(c_O\) replaced by
\((2/3)^{1/2}c_O\) in the velocity form and by \((2/3)^{1/2}c_O\) in the gradient
form, and the restriction of the operator to symmetric traceless
inputs has its own, smaller, pointwise norm (in the notation of
Proposition~\ref{prop:c3v82-opnorm}, \(\max\{(2/3)^{1/2}|\alpha-c|,|2c-h_K|/\sqrt2\}\), with the
same \(r^{-4}\) tail constant \(\sqrt{54}\)).  V86 proves the reduction and
recomputes the two Liouville floors; V87 computes the restricted
operator norm and its integral; V88 recomputes the enstrophy floor
and the derivative bounds with the reduced constants and rewrites
the table; V89 consolidates the block and fixes the use of V91--V95.
The second improvement (the \(L^6\) bound on the far field) is
recorded as the next candidate.  No global regularity claim is
made.
\end{assessment}

% =============================================================================
\section{Integrated V85 conclusion and next acquisition order}
\label{sec:c3v85-conclusion}
% =============================================================================

\begin{theorem}[What V85 establishes]
\label{thm:c3v85-conclusion}
Relative to the two legacy cycles and V1--V84: the companion-audit
decision (Theorem~\ref{thm:c3v85-companion-decision}), the position and
the decision.  The full Navier--Stokes stability/global-regularity
theorem remains unproved.
\end{theorem}

\begin{proof}
The cited records.
\end{proof}

\begin{programme}[V86 acquisition order]
\label{prog:c3v86}
\begin{enumerate}[label=\textup{(AC\arabic*)},leftmargin=4.8em]
\item Prove the traceless reduction: for a symmetric matrix field
      \(f\) with \(\operatorname{div}f\in L^2\), \(e^{t\Delta}\PP\operatorname{div}f=e^{t\Delta}\PP\operatorname{div}f_0\) with
      \(f_0=f-\frac13(\operatorname{tr}f)I\); the norms \(|(w\otimes w)_0|=(2/3)^{1/2}|w|^2\) and
      \(|(a\otimes w+w\otimes a)_0|\le(8/3)^{1/2}|a||w|\); the reduced absorptions
      \(\kappa_W\le(2/3)^{1/2}\pi c_O\kappa_W^2\) and
      \(C_\infty'\le c_h(\theta(1+\theta))^{-1/2}\kappa+2(8/3)^{1/2}c_O\kappa(\theta/(1+\theta))^{1/2}C_\infty'\).
\item Recompute the Liouville floors: \(\kappa\ge(3/2)^{1/2}/(\pi c_O)\) and the
      gradient form; values with both \(c_O\).
\item The next compulsory companion audit is V90.
\end{enumerate}
\end{programme}

\begin{assessment}[Position after V85]
\label{ass:c3v85-position}
The companion audit is recorded; the traceless-reduction block
begins.  No full stability or global-regularity claim is made.
\end{assessment}

% =============================================================================
\section*{Version V85 (third cycle) append-only record}
\addcontentsline{toc}{section}{Version V85 (third cycle) append-only record}
% =============================================================================

The complete V84 source of the third cycle was retained before this
continuation.  Its SHA--256 digest was
\begin{center}\scriptsize\ttfamily
c8a3c995285318161798d209ef20ca611f49a36efc5ca63a42a52795bffb6467.
\end{center}
V85 performed the compulsory companion audit and recorded the
position and the decision.  No full regularity claim is made.

% =============================================================================
\section{V86: the traceless reduction and the reduced Liouville floors}
\label{sec:c3v86-entry}
% =============================================================================

Items (AC1) and (AC2) of Programme~\ref{prog:c3v86} are carried out.
The Leray projector annihilates the trace part of the tensor
\(W\otimes W\) and of its derivatives, so every absorption of V81 and V84
holds with the tensor norms reduced by explicit factors:
\((2/3)^{1/2}\) for \(W\otimes W\), \((8/3)^{1/2}\) in place of \(2\) for \(\nabla(W\otimes W)\) and for the
\(W\otimes\nabla^2W\) part of \(\nabla^2(W\otimes W)\).  The two Liouville floors become
\(\kappa\ge(3/2)^{1/2}/(\pi c_O)\) and \(\kappa\ge1/(2(8/3)^{1/2}c_O)\): \(0.0926\) and \(0.0728\)
(certified), \(0.1376\) and \(0.1080\) (computed).  No global regularity
claim is made.  The next compulsory companion audit is V90.

% =============================================================================
\section{The reduction}
\label{sec:c3v86-reduction}
% =============================================================================

For a matrix \(f\), \(f_0:=f-\frac13(\operatorname{tr}f)I\) is its traceless part.

\begin{lemma}[Annihilation of the trace part]
\label{lem:c3v86-annihilation}
Let \(f\) be a bounded matrix field with \(\operatorname{div}f\in L^2\) and \(\nabla\operatorname{tr}f\in L^2\).
Then \(\PP\operatorname{div}f=\PP\operatorname{div}f_0\), and \(e^{t\Delta}\PP\operatorname{div}f=e^{t\Delta}\PP\operatorname{div}f_0\).
\end{lemma}

\begin{proof}
\(\operatorname{div}(\frac13(\operatorname{tr}f)I)=\frac13\nabla\operatorname{tr}f\), and \(\PP\nabla g=\nabla g-\nabla\Delta^{-1}\Delta g=0\) for \(\nabla g\in L^2\).
\end{proof}

\begin{lemma}[Norms of traceless parts]
\label{lem:c3v86-norms}
For vectors \(w,a\) and matrices \(f\):
\begin{enumerate}[label=\textup{(\roman*)},leftmargin=3.2em]
\item \(|(w\otimes w)_0|_F=(2/3)^{1/2}|w|^2\);
\item \(|(a\otimes w+w\otimes a)_0|_F^2=2|a|^2|w|^2+\frac23(a\cdot w)^2\le\frac83|a|^2|w|^2\);
\item \(|f_0|_F\le|f|_F\);
\item for \(W\) smooth, in Frobenius norm over all indices,
      \(|(\nabla(W\otimes W))_0|\le(8/3)^{1/2}|W||\nabla W|\) and
      \(|(\nabla^2(W\otimes W))_0|\le(8/3)^{1/2}|W||\nabla^2W|+2|\nabla W|^2\), the traceless part
      taken in the tensor indices of \(W\otimes W\).
\end{enumerate}
\end{lemma}

\begin{proof}
\(|f_0|^2=|f|^2-\frac13(\operatorname{tr}f)^2\).  (i) \(|w\otimes w|^2=|w|^4\), \(\operatorname{tr}=|w|^2\).  (ii)
\(|a\otimes w+w\otimes a|^2=2|a|^2|w|^2+2(a\cdot w)^2\), \(\operatorname{tr}=2a\cdot w\); then \((a\cdot w)^2\le|a|^2|w|^2\).
(iv) \(\partial_k(W\otimes W)=\partial_kW\otimes W+W\otimes\partial_kW\), so by (ii)
\(\sum_k|(\partial_k(W\otimes W))_0|^2\le\frac83\sum_k|\partial_kW|^2|W|^2\).  For the second derivative,
\(\partial_a\partial_b(W\otimes W)=(\partial_a\partial_bW\otimes W+W\otimes\partial_a\partial_bW)+(\partial_aW\otimes\partial_bW+\partial_bW\otimes\partial_aW)\); the
traceless part is linear, (ii) bounds the first group summed over
\((a,b)\) by \(\frac83|W|^2|\nabla^2W|^2\), and (iii) with
\(|\partial_aW\otimes\partial_bW+\partial_bW\otimes\partial_aW|\le2|\partial_aW||\partial_bW|\) bounds the second by \(4|\nabla W|^4\)
after summation; then the triangle inequality.
\end{proof}

% =============================================================================
\section{The reduced absorptions}
\label{sec:c3v86-absorptions}
% =============================================================================

\begin{theorem}[Reduced absorptions]
\label{thm:c3v86-absorptions}
For every \(W\in\mathcal F(x_0)\), with \(c_O\) of \eqref{eq:c3v81-cO} or, equivalently,
its restriction \(c_O^{s0}\le c_O\) to symmetric traceless inputs (V87):
\begin{enumerate}[label=\textup{(\arabic*)},leftmargin=3.2em]
\item \(\kappa_W\le(2/3)^{1/2}\pi c_O\kappa_W^2\); hence every homogeneous type-I point
      satisfies \(\kappa\ge(3/2)^{1/2}/(\pi c_O)\).
\item For \(\theta>0\) with \(2(8/3)^{1/2}c_O\kappa(\theta/(1+\theta))^{1/2}<1\),
      \begin{equation}
       C_\infty'\le\gamma_1^{\rm red}(\kappa,\theta)\,\kappa,\qquad
       \gamma_1^{\rm red}(\kappa,\theta):=\frac{c_h(\theta(1+\theta))^{-1/2}}{1-2(8/3)^{1/2}c_O\kappa(\theta/(1+\theta))^{1/2}} ;
       \label{eq:c3v86-gamma1}
      \end{equation}
      hence every homogeneous type-I point satisfies \(\kappa\ge1/(2(8/3)^{1/2}c_O)\).
\item For \(\theta>0\) with \((8/3)^{1/2}c_O\kappa J_2(\theta)<1\),
      \begin{equation}
       \nu^{1/2}C_\infty''\le\gamma_2^{\rm red}(\kappa,\theta)\,\kappa,\qquad
       \gamma_2^{\rm red}(\kappa,\theta):=\frac{c_{h,2}\theta^{-1}(1+\theta)^{-1/2}+2c_O\gamma_1^{\rm red}(\kappa)^2\kappa J_2(\theta)}{1-(8/3)^{1/2}c_O\kappa J_2(\theta)} ,
       \label{eq:c3v86-gamma2}
      \end{equation}
      with \(\gamma_1^{\rm red}(\kappa)\) the infimum of \eqref{eq:c3v86-gamma1} over admissible
      \(\theta\).
\end{enumerate}
\end{theorem}

\begin{proof}
In the proofs of Theorems~\ref{thm:c3v81-epsilon}, \ref{thm:c3v81-gradient}
and \ref{thm:c3v84-gamma2}, replace the tensor by its traceless part
(Lemma~\ref{lem:c3v86-annihilation}: \(\nabla\operatorname{tr}(W\otimes W)=2(\nabla W)^{\rm T}W\in L^2\), and
likewise for the derivatives) and bound it by
Lemma~\ref{lem:c3v86-norms}: \(|(W\otimes W)_0|\le(2/3)^{1/2}C_\infty^2(-\tau)^{-1}\) in (1);
\((8/3)^{1/2}\) in place of \(2\) in the coefficient of \(N\) in (2), so \(4c_O\) becomes
\(2(8/3)^{1/2}c_O\); in (3) the coefficient of \(N_2\) becomes \((8/3)^{1/2}c_O\kappa J_2\) while
the \(|\nabla W|^2\) term keeps its coefficient \(2\).  The floors as in
Theorem~\ref{thm:c3v81-epsilon} and Proposition~\ref{prop:c3v83-gradient-liouville}.
\end{proof}

\begin{corollary}[Item (AC2): reduced Liouville floors]
\label{cor:c3v86-values}
Every homogeneous type-I point satisfies
\begin{equation}
 \boxed{\quad
 \kappa\ \ge\ \frac{(3/2)^{1/2}}{\pi c_O}=\begin{cases}0.0926&\text{(certified, }c_O=4.208)\\ 0.1376&\text{(computed, }c_O=2.834)\end{cases},\qquad
 \kappa\ \ge\ \frac1{2(8/3)^{1/2}c_O}=\begin{cases}0.0728\\ 0.1080\end{cases}.
 \quad}
 \label{eq:c3v86-values}
\end{equation}
\end{corollary}

\begin{proof}
Theorem~\ref{thm:c3v86-absorptions} with the values of V82;
\((3/2)^{1/2}=1.2247\), \(2(8/3)^{1/2}=3.266\).
\end{proof}

\begin{assessment}[Reading]
\label{ass:c3v86-reading}
The reduction is exact and costs nothing; it raises the velocity
Liouville floor by \(22\%\) and the gradient one by \(22\%\).  It uses the
divergence-free structure only through the pressure (the trace part
of \(W\otimes W\) is what the pressure removes), which is the first of the
three improvements of Assessment~\ref{ass:c3v83-iteration}.  The
restricted operator norm (V87) adds a little more; the enstrophy
floor with the reduced constants is V88.
\end{assessment}

% =============================================================================
\section{Integrated V86 conclusion and next acquisition order}
\label{sec:c3v86-conclusion}
% =============================================================================

\begin{theorem}[What V86 proves]
\label{thm:c3v86-conclusion}
Relative to the two legacy cycles and V1--V85: the annihilation and
norm lemmas (Lemmas~\ref{lem:c3v86-annihilation}, \ref{lem:c3v86-norms}),
the reduced absorptions (Theorem~\ref{thm:c3v86-absorptions}) and the
reduced Liouville floors (Corollary~\ref{cor:c3v86-values}).  The full
Navier--Stokes stability/global-regularity theorem remains
unproved.
\end{theorem}

\begin{proof}
The cited results.
\end{proof}

\begin{programme}[V87 acquisition order]
\label{prog:c3v87}
\begin{enumerate}[label=\textup{(AC\arabic*)},leftmargin=4.8em]
\item The restricted operator norm: prove
      \(\|M(x)\|_{s0}=\max\{(2/3)^{1/2}|\alpha-c|,|2c-h_K|/\sqrt2\}\) on symmetric traceless
      inputs, from the decomposition of Proposition~\ref{prop:c3v82-opnorm};
      its tail; its integral \(c_O^{s0}\) computed, and the certified bound
      unchanged.
\item Record the values of \(c_O^{s0}\) and the gain over \(c_O\).
\item The next compulsory companion audit is V90.
\end{enumerate}
\end{programme}

\begin{assessment}[Position after V86]
\label{ass:c3v86-position}
The traceless reduction is proved and the Liouville floors are
raised.  No full stability or global-regularity claim is made.
\end{assessment}

% =============================================================================
\section*{Version V86 (third cycle) append-only record}
\addcontentsline{toc}{section}{Version V86 (third cycle) append-only record}
% =============================================================================

The complete V85 source of the third cycle was retained before this
continuation.  Its SHA--256 digest was
\begin{center}\scriptsize\ttfamily
6040f292e16b0eafc8bc8d4483c09395dabbf49b181acfa6e62c321327519293.
\end{center}
V86 proved the traceless reduction and raised the Liouville floors.
No full regularity claim is made.

% =============================================================================
\section{V87: the restricted Oseen constant}
\label{sec:c3v87-entry}
% =============================================================================

Items (AC1) and (AC2) of Programme~\ref{prog:c3v87} are carried out.
The operator \(e^{\Delta}\PP\operatorname{div}\) restricted to symmetric traceless inputs
has the pointwise norm \(\max\{\sqrt6|c|,|2c-h_K|/\sqrt2\}\) (using the exact identity
\(\phi'''-h_K=-2c\), which also simplifies the unrestricted norm), never
larger than the unrestricted one, equal to it for \(r\ge3\) and with the
same tail; its integral is \(c_O^{s0}\le2.764\) (computed), a gain of \(2.5\%\) over
\(c_O\), the near region where the restriction helps contributing
little.  The certified bound stays at \(4.208\).  No global regularity
claim is made.  The next compulsory companion audit is V90.

% =============================================================================
\section{The restricted norm}
\label{sec:c3v87-norm}
% =============================================================================

\begin{proposition}[Restricted pointwise norm]
\label{prop:c3v87-opnorm}
In the notation of Proposition~\ref{prop:c3v82-opnorm}, the norm of
\(M(x)\) on the five-dimensional space of symmetric traceless matrices
is
\begin{equation}
 \|M(x)\|_{s0}=\max\Bigl\{\Bigl(\frac23\Bigr)^{1/2}|\alpha-c|,\ \frac{|2c-h_K|}{\sqrt2}\Bigr\}\ \le\ \|M(x)\|_{F\to2},
 \label{eq:c3v87-opnorm}
\end{equation}
and, with \(c_O^{s0}:=\sup\{\|e^{\Delta}\PP\operatorname{div}f\|_\infty:\|f\|_\infty\le1,\ f\text{ symmetric traceless}\}\),
\(c_O^{s0}\le\int\|M\|_{s0}\dd x\); Theorems~\ref{thm:c3v86-absorptions} and
\ref{thm:c3v84-gamma2} hold with \(c_O^{s0}\) in place of \(c_O\).
\end{proposition}

\begin{proof}
In the proof of Proposition~\ref{prop:c3v82-opnorm}, symmetric means
\(w_a=0\) and traceless means \(\tau=0\), so \(Mf=n(\alpha-c)\mu+(2c-h_K)w_s\) with the
budget \(|f|_F^2\ge\mu^2+\frac12\mu^2+2|w_s|^2=\frac32\mu^2+2|w_s|^2\).  Hence
\(|Mf|^2\le\frac23(\alpha-c)^2\cdot\frac32\mu^2+\frac12(2c-h_K)^2\cdot2|w_s|^2\le\max\{\frac23(\alpha-c)^2,\frac12(2c-h_K)^2\}|f|^2\),
with equality for \(f=\operatorname{diag}(-\frac12,-\frac12,1)\mu\)-type inputs (\(w_s=0\), block
\(-\frac\mu2I_2\), which realizes \(|f|^2=\frac32\mu^2\)) or for \(f=w_s\otimes n+n\otimes w_s\).  The
inputs of the Duhamel terms, \((W\otimes W)_0\), \((\partial_k(W\otimes W))_0\) and
\((\partial_a\partial_b(W\otimes W))_0\), are symmetric traceless, so the Cauchy--Schwarz
argument of Theorem~\ref{thm:c3v84-gamma2} applies with \(\|M\|_{s0}\).
\end{proof}

\begin{lemma}[The identity \(\alpha=-2c\), comparison and tail]
\label{lem:c3v87-tail}
\(\phi'''-h_K=-2\bigl(\phi''/r-\phi'/r^2\bigr)\), i.e.\ \(\alpha=-2c\) identically.  Hence
\begin{equation}
 \|M\|_{s0}=\max\Bigl\{\sqrt6\,|c|,\ \frac{|2c-h_K|}{\sqrt2}\Bigr\},\qquad
 \|M\|_{F\to2}=\max\Bigl\{\sqrt6\,|c|,\ \Bigl(\frac{(2c-h_K)^2+h_K^2}2\Bigr)^{1/2}\Bigr\},
 \label{eq:c3v87-simplified}
\end{equation}
the two norms differ only where the second entry of the unrestricted
norm is the larger, i.e.\ for \(r<2.34\) (the crossing radius); for \(r\ge2.34\)
they are equal, and for \(r\ge12\), \(\|M\|_{s0}\le(\sqrt{54}+10^{-9})/(4\pi r^4)\), since
\(\sqrt6\cdot3=\sqrt{54}\).
\end{lemma}

\begin{proof}
For radial \(\phi\), \(\partial_r(\Delta\phi)=\phi'''+2\phi''/r-2\phi'/r^2\), and \(\Delta\Phi_1=-K_1\)
(\(\Phi_1=(-\Delta)^{-1}K_1\)), so \(\partial_r(\Delta\phi)=-\partial_rK_1=\frac r2K_1=h_K\); this is
\(\alpha=-2c\).  Then \((\alpha^2+2c^2)^{1/2}=\sqrt6|c|=(2/3)^{1/2}|\alpha-c|\), so the first entries of
\eqref{eq:c3v82-opnorm} and \eqref{eq:c3v87-opnorm} coincide, and the
second entries differ by the term \(h_K^2\).  The crossing radius and
the tail: Lemma~\ref{lem:c3v82-tail} with \(4\pi r^4c\to3\).
\end{proof}

\begin{proposition}[Item (AC2): computed restricted constant]
\label{prop:c3v87-computed}
\(\int_{|x|\le12}\|M\|_{s0}\dd x=2.1516\) by the quadrature of
Proposition~\ref{prop:c3v82-computed}, and with the tail
\begin{equation}
 c_O^{s0}\ \le\ 2.764\qquad(\text{computed; against }c_O\le2.834),
 \label{eq:c3v87-value}
\end{equation}
while the certified bound of Theorem~\ref{thm:c3v82-certified} applies
unchanged, \(c_O^{s0}\le4.208\).  Formula \eqref{eq:c3v87-opnorm} was checked
against the numerically assembled \(3\times5\) matrix of \(M\) on an
orthonormal basis of symmetric traceless matrices at seven radii.
\end{proposition}

\begin{proof}
As in Proposition~\ref{prop:c3v82-computed}, with \eqref{eq:c3v87-opnorm} as
integrand; the tail is Lemma~\ref{lem:c3v87-tail}.
\end{proof}

\begin{assessment}[Reading]
\label{ass:c3v87-reading}
The gain from the restriction is small because the integral of the
Oseen gradient is dominated by \(r\ge3\), where the maximizing input
is already symmetric and traceless (\(\alpha=-2c\) asymptotically).  The
traceless reduction of V86, which acts on the tensor rather than on
the kernel, is where the gain is.  What remains between the
certified \(4.208\) and the computed \(2.764\) is the certification of the
quadrature, which V89 addresses by a Lipschitz bound on the
integrand obtained from the Fourier side.
\end{assessment}

% =============================================================================
\section{Integrated V87 conclusion and next acquisition order}
\label{sec:c3v87-conclusion}
% =============================================================================

\begin{theorem}[What V87 proves]
\label{thm:c3v87-conclusion}
Relative to the two legacy cycles and V1--V86: the restricted norm
(Proposition~\ref{prop:c3v87-opnorm}), its comparison and tail
(Lemma~\ref{lem:c3v87-tail}) and the computed value
(Proposition~\ref{prop:c3v87-computed}).  The full Navier--Stokes
stability/global-regularity theorem remains unproved.
\end{theorem}

\begin{proof}
The cited results.
\end{proof}

\begin{programme}[V88 acquisition order]
\label{prog:c3v88}
\begin{enumerate}[label=\textup{(AC\arabic*)},leftmargin=4.8em]
\item The enstrophy floor with the reduced constants:
      \(\Gamma^{\rm red}(\kappa)=\inf_\theta\gamma_1^{\rm red}(\kappa,\theta)\kappa+\frac{\sqrt3}2\kappa+\frac{\kappa^2}2\) and its root for
      \(c_O\in\{4.208,2.834,2.764\}\); \(\gamma_2^{\rm red}\) and \(\theta_1\) at the floors;
      the table of V84 rewritten with three columns.
\item Record which floor is operative and by how much the
      reduction moved it.
\item The next compulsory companion audit is V90.
\end{enumerate}
\end{programme}

\begin{assessment}[Position after V87]
\label{ass:c3v87-position}
The restricted constant is computed.  No full stability or
global-regularity claim is made.
\end{assessment}

% =============================================================================
\section*{Version V87 (third cycle) append-only record}
\addcontentsline{toc}{section}{Version V87 (third cycle) append-only record}
% =============================================================================

The complete V86 source of the third cycle was retained before this
continuation.  Its SHA--256 digest was
\begin{center}\scriptsize\ttfamily
a3eed5bdc0cd652e058263b8327336e228263ca9b94545a7822439479cdd8254.
\end{center}
V87 computed the restricted Oseen constant.  No full regularity
claim is made.

% =============================================================================
\section{V88: the enstrophy floor with the reduced constants, and the three-column table}
\label{sec:c3v88-entry}
% =============================================================================

Items (AC1) and (AC2) of Programme~\ref{prog:c3v88} are carried out.
With the traceless reduction the enstrophy floor becomes
\(\kappa\ge0.1253\) (certified, Plancherel bound on \(c_O\)), \(\kappa\ge0.1544\) (computed,
\(c_O\)) and \(\kappa\ge0.1566\) (computed, restricted \(c_O^{s0}\)); the derivative
bounds and the concentration threshold at the floors are
recomputed, and the table is rewritten with three columns.  The
reduction moved the certified floor from \(0.1129\) to \(0.1253\) (\(+11\%\)) and
the computed one from \(0.1386\) to \(0.1566\) (\(+13\%\)).  An in-place
correction of V87 is recorded.  No global regularity claim is made.
The next compulsory companion audit is V90.

\begin{remark}[In-place correction of V87]
\label{rem:c3v88-correction}
Lemma~\ref{lem:c3v87-tail} as first assembled located the crossing of
the two entries of the unrestricted norm at ``about \(2.9\)''; the
crossing radius is \(2.34\) (the first entry \(\sqrt6|c|\) dominates for
\(r\ge2.34\), and for the restricted norm already for \(r\ge1.5\)).  The
statement was corrected in place; the digest of V87 recorded below
is that of the corrected file.  Nothing else depends on the value.
\end{remark}

% =============================================================================
\section{The reduced enstrophy floor}
\label{sec:c3v88-floor}
% =============================================================================

Let \(\Gamma^{\rm red}(\kappa):=\inf_\theta\gamma_1^{\rm red}(\kappa,\theta)\kappa+\frac{\sqrt3}2\kappa+\frac{\kappa^2}2\), the infimum over the
admissible \(\theta\) of \eqref{eq:c3v86-gamma1}, increasing in \(\kappa\) as in
Theorem~\ref{thm:c3v83-floor}.

\begin{theorem}[Reduced enstrophy floor]
\label{thm:c3v88-floor}
Every homogeneous type-I point satisfies \(\Gamma^{\rm red}(\kappa)\ge1\), hence
\(\kappa\ge\kappa_1^{\rm red}:=\min\{\kappa:\Gamma^{\rm red}(\kappa)\ge1\}\), with
\begin{equation}
 \boxed{\quad
 \kappa_1^{\rm red}=\begin{cases}
 0.1253 & (c_O=4.208,\ \text{certified};\ \theta^*=0.114,\ \gamma_1^{\rm red}\kappa=0.884)\\
 0.1544 & (c_O=2.834,\ \text{computed};\ \theta^*=0.200,\ \gamma_1^{\rm red}\kappa=0.854)\\
 0.1566 & (c_O^{s0}=2.764,\ \text{computed};\ \theta^*=0.208,\ \gamma_1^{\rm red}\kappa=0.852).
 \end{cases}
 \quad}
 \label{eq:c3v88-values}
\end{equation}
\end{theorem}

\begin{proof}
\eqref{eq:c3v67-constraint} with Theorem~\ref{thm:c3v86-absorptions}(2)
(and Proposition~\ref{prop:c3v87-opnorm} for the third line); the
values by the bisection of Theorem~\ref{thm:c3v83-floor} with the
coefficient \(2(8/3)^{1/2}c_O\) in place of \(4c_O\).
\end{proof}

\begin{proposition}[Derivative bounds and the concentration threshold at the reduced floors]
\label{prop:c3v88-derivatives}
At the floors of \eqref{eq:c3v88-values}, with \(c_{h,2}=1.794\) (certified)
or \(1.415\) (computed):
\(\gamma_1^{\rm red}=7.05,\ 5.53,\ 5.44\); \(\nu^{1/2}C_\infty''\le11.33,\ 7.85,\ 7.73\) (\(\gamma_2^{\rm red}=90.4,\ 50.8,\ 49.3\) at
\(\theta=0.062,\ 0.089,\ 0.092\)); and \(\theta_1\ge0.404,\ 0.379,\ 0.377\), so the
concentration row applies at each floor, while the palinstrophy row
remains void (Theorem~\ref{thm:c3v84-void}).
\end{proposition}

\begin{proof}
Theorem~\ref{thm:c3v86-absorptions}(3) and the definition of \(\theta_1\), with
the minimization of V84.
\end{proof}

% =============================================================================
\section{The three-column table}
\label{sec:c3v88-table}
% =============================================================================

\begin{theorem}[Item (AC2): explicit floors after the reduction]
\label{thm:c3v88-table}
Every homogeneous type-I point satisfies:

\medskip
\begin{center}\small
\begin{tabular}{@{}p{0.36\textwidth}p{0.15\textwidth}p{0.15\textwidth}p{0.17\textwidth}@{}}
\toprule
\textbf{Statement} & \textbf{Certified} \((c_O{=}4.208)\) & \textbf{Computed} \((c_O{=}2.834)\) & \textbf{Computed} \((c_O^{s0}{=}2.764)\)\\
\midrule
\(\kappa\ge1/(2(8/3)^{1/2}c_O)\) & \(0.0728\) & \(0.1080\) & \(0.1108\)\\
\(\kappa\ge(3/2)^{1/2}/(\pi c_O)\) & \(0.0926\) & \(0.1376\) & \(0.1410\)\\
\(\kappa\ge\kappa_1^{\rm red}\) (enstrophy row) & \(0.1253\) & \(0.1544\) & \(0.1566\)\\
\(C_\infty'\le\gamma_1^{\rm red}\kappa\) at the floor & \(0.884\) & \(0.854\) & \(0.852\)\\
\(\nu^{1/2}C_\infty''\) at the floor & \(11.33\) & \(7.85\) & \(7.73\)\\
concentration row at the floor, \(\theta_1\ge\) & \(0.404\) & \(0.379\) & \(0.377\)\\
palinstrophy row and loops & void \((\kappa<0.4569)\) & void & void\\
\bottomrule
\end{tabular}
\end{center}
\medskip

\noindent
The operative floor is the enstrophy row in every column.
\end{theorem}

\begin{proof}
Corollary~\ref{cor:c3v86-values}, Proposition~\ref{prop:c3v87-computed},
Theorem~\ref{thm:c3v88-floor}, Proposition~\ref{prop:c3v88-derivatives},
Theorem~\ref{thm:c3v84-void}.
\end{proof}

\begin{assessment}[Reading]
\label{ass:c3v88-reading}
The reduction moved every floor by \(11\) to \(23\%\); the enstrophy row
stays the operative floor and its margin over the Liouville floor
shrinks (from \(1.49\) to \(1.35\) certified, from \(1.23\) to \(1.11\) computed),
because the gradient absorption gains \((8/3)^{1/2}/2\) while the velocity one
gains \((2/3)^{1/2}\), a larger factor.  The certified column lags the
computed ones by the Plancherel loss on \(c_O\) alone; certifying the
quadrature (V89) makes the third column the certified one.
\end{assessment}

% =============================================================================
\section{Integrated V88 conclusion and next acquisition order}
\label{sec:c3v88-conclusion}
% =============================================================================

\begin{theorem}[What V88 proves]
\label{thm:c3v88-conclusion}
Relative to the two legacy cycles and V1--V87: the reduced enstrophy
floor (Theorem~\ref{thm:c3v88-floor}), the derivative bounds and
thresholds at the reduced floors (Proposition~\ref{prop:c3v88-derivatives})
and the table (Theorem~\ref{thm:c3v88-table}).  The full Navier--Stokes
stability/global-regularity theorem remains unproved.
\end{theorem}

\begin{proof}
The cited results.
\end{proof}

\begin{programme}[V89 acquisition order]
\label{prog:c3v89}
\begin{enumerate}[label=\textup{(AC\arabic*)},leftmargin=4.8em]
\item Certify the quadrature: prove \(\sup_x|T|_F\le\sqrt2/(4\pi^2)\) and
      \(\sup_x|\nabla T|_F\le3\sqrt2\pi^{1/2}/(16\pi^2)\) from the Fourier side, hence a
      Lipschitz constant \(L\le97\) for \(G(r)=4\pi r^2\|M(r)\|\) on \([0,12]\); the
      midpoint-rule error bound \(12Lh/4\); with \(N=2\cdot10^6\) cells the
      certified values \(c_O\le2.836\), \(c_O^{s0}\le2.766\).
\item The final certified table (one column), the consolidation of
      the block V86--V89, and the use of V91--V95.
\item The next compulsory companion audit is V90.
\end{enumerate}
\end{programme}

\begin{assessment}[Position after V88]
\label{ass:c3v88-position}
The reduced floors are tabulated.  No full stability or
global-regularity claim is made.
\end{assessment}

% =============================================================================
\section*{Version V88 (third cycle) append-only record}
\addcontentsline{toc}{section}{Version V88 (third cycle) append-only record}
% =============================================================================

The complete V87 source of the third cycle (after the in-place
correction of Remark~\ref{rem:c3v88-correction}) was retained before
this continuation.  Its SHA--256 digest was
\begin{center}\scriptsize\ttfamily
7ffc9ec072e882d1dc4bc6bdcb8d5bd04e3aa8bc2570cf40687ffeea31ea2ab2.
\end{center}
V88 recomputed the enstrophy floor with the reduced constants and
rewrote the table.  No full regularity claim is made.

% =============================================================================
\section{V89: certifying the quadrature, the final table, and the block consolidated}
\label{sec:c3v89-entry}
% =============================================================================

Items (AC1) and (AC2) of Programme~\ref{prog:c3v89} are carried out.
The integrand of the Oseen constant is shown to be Lipschitz on
\([0,12]\) with an explicit constant obtained from the Fourier side, so
that a midpoint sum with \(2\cdot10^6\) cells carries an error below
\(0.00175\); with the analytic tail this certifies \(c_O\le2.836\) and
\(c_O^{s0}\le2.766\), and every floor of the block is then proved with the
values of the former ``computed'' column: at every homogeneous
type-I point, \(\kappa\ge0.1565\).  The block V86--V89 is consolidated and
the direction for V91--V95 is fixed: the strain form of the enstrophy
row.  No global regularity claim is made.  The next compulsory
companion audit is V90.

% =============================================================================
\section{Certification of the quadrature}
\label{sec:c3v89-certification}
% =============================================================================

\begin{lemma}[Sup-norm bounds from the Fourier side]
\label{lem:c3v89-sup}
With \(T=\nabla K^O_1\) (the tensor \(\partial_k(K^O_1)_{ij}\)),
\begin{equation}
 \sup_x|T(x)|_F\le\frac{\sqrt2}{4\pi^2}=0.03583,\qquad
 \sup_x|\nabla T(x)|_F\le\frac{3\sqrt2\,\pi^{1/2}}{16\pi^2}=0.04763 .
 \label{eq:c3v89-sup}
\end{equation}
\end{lemma}

\begin{proof}
\(T(x)=(2\pi)^{-3}\int\hat T(\xi)e^{ix\cdot\xi}\dd\xi\) with \(|\hat T|_F=\sqrt2\rho e^{-\rho^2}\)
(Theorem~\ref{thm:c3v82-certified}), so
\(|T|_F\le(2\pi)^{-3}4\pi\sqrt2\int_0^\infty\rho^3e^{-\rho^2}\dd\rho=(2\pi)^{-3}\cdot2\sqrt2\pi\).  For \(\nabla T\), the
transform is \(i\xi_m\hat T_{ijk}\), of Frobenius norm \(\sqrt2\rho^2e^{-\rho^2}\), and
\(\int_0^\infty\rho^4e^{-\rho^2}\dd\rho=3\pi^{1/2}/8\).
\end{proof}

\begin{lemma}[Lipschitz bound for the integrand]
\label{lem:c3v89-lipschitz}
For either norm \(\|M\|\in\{\|M\|_{F\to2},\|M\|_{s0}\}\), the function
\(G(r):=4\pi r^2\|M(rn)\|\) (independent of the unit vector \(n\)) is Lipschitz
on \([0,12]\) with constant
\begin{equation}
 L\ \le\ 8\pi\cdot12\cdot\frac{\sqrt2}{4\pi^2}+4\pi\cdot144\cdot\frac{3\sqrt2\pi^{1/2}}{16\pi^2}=10.80+86.18<97 .
 \label{eq:c3v89-L}
\end{equation}
\end{lemma}

\begin{proof}
\(\|M(x)\|\le|T(x)|_F\) and \(|\,\|M(x)\|-\|M(y)\|\,|\le\|M(x)-M(y)\|\le|T(x)-T(y)|_F\le\sup|\nabla T|_F|x-y|\),
so \(r\mapsto\|M(rn)\|\) is bounded by \(\sup|T|_F\) and Lipschitz with constant
\(\sup|\nabla T|_F\); the product rule with \(4\pi r^2\) on \([0,12]\) gives
\eqref{eq:c3v89-L} by Lemma~\ref{lem:c3v89-sup}.
\end{proof}

\begin{theorem}[Certified Oseen constants]
\label{thm:c3v89-certified}
\(c_O\le2.836\) and \(c_O^{s0}\le2.766\).
\end{theorem}

\begin{proof}
For a Lipschitz \(G\) with constant \(L\) on \([0,12]\), the midpoint sum with
\(N\) cells of width \(h=12/N\) satisfies \(|\int_0^{12}G-h\sum_iG(r_i)|\le N\cdot Lh^2/4=3Lh\).
With \(N=2\cdot10^6\), \(h=6\cdot10^{-6}\) and \(L\le97\): error \(\le0.00175\).  The sums,
evaluated from the exact formulas \eqref{eq:c3v87-simplified} in double
precision (accumulated rounding below \(10^{-9}\); in the first cells the
cancellation in \(E-gr\) costs a relative error below \(10^{-3}\) on
values below \(10^{-15}\), immaterial), are \(2.221433\) for \(\|M\|_{F\to2}\) and
\(2.151581\) for \(\|M\|_{s0}\).  The tails from \(12\) are at most
\((\sqrt{54}+10^{-9})/12=0.612373\) (Lemmas~\ref{lem:c3v82-tail},
\ref{lem:c3v87-tail}).  Hence \(c_O\le2.221433+0.00175+0.612373<2.836\) and
\(c_O^{s0}\le2.151581+0.00175+0.612373<2.766\).
\end{proof}

\begin{remark}[Status]
\label{rem:c3v89-status}
The certification rests on: exact closed-form integrands (V82, V87),
a Lipschitz constant proved from Gaussian moments, an elementary
quadrature error bound, an analytic tail, and \(2\cdot10^6\) evaluations of
elementary functions in floating point.  It is a computer-assisted
bound of the ordinary kind; the Plancherel bound \(4.208\) of V82
remains the hand-only bound.
\end{remark}

% =============================================================================
\section{The final table}
\label{sec:c3v89-table}
% =============================================================================

\begin{theorem}[Item (AC2): certified explicit floors]
\label{thm:c3v89-table}
Every homogeneous type-I point of a smooth periodic solution
satisfies, with \(c_O^{s0}\le2.766\) and \(c_{h,2}\le1.794\), all certified:

\medskip
\begin{center}\small
\begin{tabular}{@{}p{0.5\textwidth}p{0.16\textwidth}p{0.2\textwidth}@{}}
\toprule
\textbf{Statement} & \textbf{Value} & \textbf{Source}\\
\midrule
\(\kappa\ge1/(2(8/3)^{1/2}c_O^{s0})\) (gradient Liouville) & \(0.1107\) & V83, V86, V89\\
\(\kappa\ge(3/2)^{1/2}/(\pi c_O^{s0})\) (velocity Liouville) & \(0.1410\) & V81, V86, V89\\
\(\kappa\ge\kappa_1^{\rm red}\) (enstrophy row; \(\theta^*=0.207\)) & \(\mathbf{0.1565}\) & V67, V81, V86, V88, V89\\
\(C_\infty'\le\gamma_1^{\rm red}\kappa\) at the floor (\(\gamma_1^{\rm red}=5.45\)) & \(0.852\) & V86\\
\(\nu^{1/2}C_\infty''\le\gamma_2^{\rm red}\kappa\) at the floor (\(\gamma_2^{\rm red}=55.7\), \(\theta=0.104\)) & \(8.72\) & V84, V86\\
concentration row applies at the floor & \(\theta_1\ge0.377\) & V71, V84\\
palinstrophy row and loop constraints & void for \(\kappa<0.4569\) & V84\\
\bottomrule
\end{tabular}
\end{center}
\medskip

\noindent
In particular, the rescaled velocity bound at a homogeneous type-I
singular point satisfies \(C_\infty\ge0.1565\,\nu^{1/2}\).
\end{theorem}

\begin{proof}
Theorem~\ref{thm:c3v89-certified} in Corollary~\ref{cor:c3v86-values},
Theorem~\ref{thm:c3v88-floor} and Proposition~\ref{prop:c3v88-derivatives}
(the bisection and minimizations repeated with \(c_O^{s0}=2.766\) and
\(c_{h,2}=1.794\): floor \(0.1565\), \(\gamma_1^{\rm red}\kappa=0.852\), \(\gamma_2^{\rm red}=55.7\),
\(\theta_1\ge0.377\)).
\end{proof}

% =============================================================================
\section{The block consolidated, and the direction for V91--V95}
\label{sec:c3v89-block}
% =============================================================================

\begin{theorem}[The traceless-reduction block, V86--V89]
\label{thm:c3v89-block}
At a homogeneous type-I point: the Leray projector annihilates the
trace parts of \(W\otimes W\) and its derivatives, reducing the tensor norms
by \((2/3)^{1/2}\) and \((8/3)^{1/2}/2\) (V86); the operator restricted to symmetric
traceless inputs has pointwise norm \(\max\{\sqrt6|c|,|2c-h_K|/\sqrt2\}\), with the
exact identity \(\phi'''-h_K=-2c\) (V87); the reduced enstrophy floor is
\(0.1253\) by hand and \(0.1565\) with the certified quadrature (V88, V89); the
quadrature is certified by a Fourier-side Lipschitz constant and the
midpoint rule (V89).
\end{theorem}

\begin{proof}
Theorems~\ref{thm:c3v86-absorptions}, \ref{thm:c3v88-floor},
\ref{thm:c3v89-certified}, \ref{thm:c3v89-table};
Propositions~\ref{prop:c3v87-opnorm}, \ref{prop:c3v87-computed};
Lemmas~\ref{lem:c3v87-tail}, \ref{lem:c3v89-lipschitz}.
\end{proof}

\begin{assessment}[Direction for V91--V95: the strain form of the enstrophy row]
\label{ass:c3v89-direction}
At the floor the gradient term carries \(0.85\) of the unit in the
enstrophy constraint, so the constraint's own stretching bound is
where a gain lies.  The stretching integrand is
\(\Omega\cdot(\Omega\cdot\nabla)V=\Omega^{\rm T}(\nabla V)\Omega=\Omega^{\rm T}S\Omega\) with \(S=\frac12(\nabla V+\nabla V^{\rm T})\) the strain,
traceless because \(\operatorname{div}V=0\); a traceless symmetric \(3\times3\) matrix has
largest eigenvalue at most \((2/3)^{1/2}|S|_F\le(2/3)^{1/2}|\nabla V|_F\).  Hence
\(|\mathcal S|\le(2/3)^{1/2}C_\infty'\mathcal E\) and the enstrophy constraint becomes
\((2/3)^{1/2}C_\infty'+\frac{\sqrt3}2\kappa+\frac{\kappa^2}2\ge1\), which moves every downstream
statement: the gradient floor of V73, the threshold \(0.4569\) of V84,
the floor \(\kappa_1\).  V91 proves the strain form and its consequences
in the class; V92 recomputes the certified floor (expected near
\(0.18\)); V93 examines the transport bound \(|\mathcal T|\) (the Hardy step) for a
similar structural gain and records the negative if none; V94
writes the constants sheet of the explicit chapter; V95 audits.  No
global regularity claim is made.
\end{assessment}

% =============================================================================
\section{Integrated V89 conclusion and next acquisition order}
\label{sec:c3v89-conclusion}
% =============================================================================

\begin{theorem}[What V89 proves]
\label{thm:c3v89-conclusion}
Relative to the two legacy cycles and V1--V88: the sup-norm and
Lipschitz lemmas (Lemmas~\ref{lem:c3v89-sup}, \ref{lem:c3v89-lipschitz}),
the certified Oseen constants (Theorem~\ref{thm:c3v89-certified}), the
final certified table (Theorem~\ref{thm:c3v89-table}) and the
consolidated block (Theorem~\ref{thm:c3v89-block}).  The full
Navier--Stokes stability/global-regularity theorem remains
unproved.
\end{theorem}

\begin{proof}
The cited results.
\end{proof}

\begin{programme}[V90 acquisition order]
\label{prog:c3v90}
\begin{enumerate}[label=\textup{(AC\arabic*)},leftmargin=4.8em]
\item The compulsory companion audit: read the current SM and YM
      notes and record their digests and decision.
\item Record the position after ninety versions; delivery.
\end{enumerate}
\end{programme}

\begin{assessment}[Position after V89]
\label{ass:c3v89-position}
The explicit floor \(0.1565\) is certified.  No full stability or
global-regularity claim is made.
\end{assessment}

% =============================================================================
\section*{Version V89 (third cycle) append-only record}
\addcontentsline{toc}{section}{Version V89 (third cycle) append-only record}
% =============================================================================

The complete V88 source of the third cycle was retained before this
continuation.  Its SHA--256 digest was
\begin{center}\scriptsize\ttfamily
8bd4016edbad6ec3d5f75146265fd86cae7533289ed311450b0e6a31c95c541d.
\end{center}
V89 certified the quadrature and consolidated the block.  No full
regularity claim is made.

% =============================================================================
\section{V90: compulsory companion audit and the position after ninety versions}
\label{sec:c3v90-entry}
% =============================================================================

This is the ninetieth version of the third cycle.  The compulsory
review of the two companion programmes is performed first, and the
position after ninety versions is recorded with the plan for the
block V91--V95.  No global regularity claim is made.  The next
compulsory companion audit is V95.

% =============================================================================
\section{Compulsory V90 review of the current companion notes}
\label{sec:c3v90-companion-audit}
% =============================================================================

The current companion files in \texttt{latex\_notes} were inspected
read-only.  Both programmes have moved since the V85 audit: the SM
programme from V49 to V50 of its seventh series (at a first
inspection for this audit the file was still V49; it was superseded
before the final inspection, whose digest is recorded), the YM
programme from V25 to V26 of its sixth series.  The frozen
checkpoints are
\begin{align}
 &\texttt{Renewal\_SM\_Einstein\_Continuum\_Research\_Notes\_V50.tex},
 \notag\\[-2pt]
 &\qquad
 \texttt{5931c1a01125c21735ac326608ef856f01124ac9d2439155a5c4cd04871495f2},
 \label{eq:c3v90-SM-checkpoint}\\
 &\texttt{Renewal\_Yang\_Mills\_Mass\_Gap\_Research\_Notes\_V26.tex},
 \notag\\[-2pt]
 &\qquad
 \texttt{d51eb25593783101c72d9b897a112c3b9cd0f3c48a51ef378b4c3b2aec635a4e}.
 \label{eq:c3v90-YM-checkpoint}
\end{align}
(SM V50 has 19900 lines; YM V26 has 11523 lines.  At the inspection
the folder also held a file named YM V27, byte-identical to YM V26,
a version in preparation, not recorded separately.)

\emph{SM V50.}  The SM programme performed its own companion audit
and wrote a completion-locality ledger: leading and completed global
kernels, a certified resolvent and sign locality, global rank and
trace norm, a record of why this does not yet prove
heat-coefficient locality, and a heat-flat jet card.

\emph{YM V26.}  The YM programme separated its degree-eighteen link
graph from its degree-twenty plaquette graph, computed the
thresholds \(17(a+p)<1\) and \(19(a+p)<1\), proved an exact
plaquette-star shielding, derived a one-point bad-plaquette
probability and proved its insufficiency for paths, defined a
multi-defect suppression card giving its path product and a
conditional explicit \(\beta\) threshold, and defined a gauge-compatible
defect-erasure card with its multi-defect and connected-animal
consequences.

\begin{theorem}[V90 companion-audit decision]
\label{thm:c3v90-companion-decision}
Nothing is imported from SM V50 or YM V26.  Neither bears on the
third cycle's chapters, on the explicit-constant block or on the
traceless-reduction block.
\end{theorem}

\begin{proof}
The SM results concern global kernels and locality ledgers of
lattice Dirac interfaces; the YM results concern graph degrees, shielding and defect
suppression in a lattice gauge theory; none is a Navier--Stokes
statement, and none concerns the mild formulation, the Oseen
constant, the enstrophy row or the explicit floors.
\end{proof}

% =============================================================================
\section{Position after ninety versions, and the block V91--V95}
\label{sec:c3v90-position}
% =============================================================================

\begin{assessment}[Position after ninety versions of the third cycle]
\label{ass:c3v90-position-ninety}
Sixteen chapters are written: those of
Assessment~\ref{ass:c3v85-position-eightyfive} and the
traceless-reduction block (V86--V89).  The cycle's headline result
is now certified with the values of the computed column: every
homogeneous type-I point of a smooth periodic solution has velocity
constant \(\kappa\ge0.1565\) (Theorem~\ref{thm:c3v89-table}), the Liouville
floors \(0.1410\) and \(0.1107\) below it, explicit derivative bounds at the
floor (\(C_\infty'\le0.852\), \(\nu^{1/2}C_\infty''\le8.72\)), the concentration of the
vorticity's Gaussian second moment, and the structural voidness of
the palinstrophy row below \(0.4569\).  The certification chain is: exact
radial kernels, an exact pointwise operator norm, an analytic tail,
a Fourier-side Lipschitz constant, and a midpoint sum with an
elementary error bound.  The block V91--V95 follows
Assessment~\ref{ass:c3v89-direction}: the strain form of the enstrophy
row (V91), its certified floor (V92, expected near \(0.17\)), the
transport bound examined (V93), the constants sheet of the explicit
chapter (V94), audit (V95).  Ten versions remain; the closing block
V96--V100 is reserved for the manifest, the negative records and the
draft of the fourth cycle's V1.  No global regularity claim is made.
\end{assessment}

% =============================================================================
\section{Integrated V90 conclusion and next acquisition order}
\label{sec:c3v90-conclusion}
% =============================================================================

\begin{theorem}[What V90 establishes]
\label{thm:c3v90-conclusion}
Relative to the two legacy cycles and V1--V89: the companion-audit
decision (Theorem~\ref{thm:c3v90-companion-decision}) and the position.
The full Navier--Stokes stability/global-regularity theorem remains
unproved.
\end{theorem}

\begin{proof}
The cited records.
\end{proof}

\begin{programme}[V91 acquisition order]
\label{prog:c3v91}
\begin{enumerate}[label=\textup{(AC\arabic*)},leftmargin=4.8em]
\item The strain form: \(\Omega\cdot(\Omega\cdot\nabla)V=\Omega^{\rm T}S\Omega\), \(S\) the traceless strain;
      \(\lambda_{\max}(S)\le(2/3)^{1/2}|S|_F\le(2/3)^{1/2}|\nabla V|_F\); hence
      \(|\mathcal S|\le(2/3)^{1/2}C_\infty'\mathcal E\) and the constraint
      \((2/3)^{1/2}C_\infty'+\frac{\sqrt3}2\kappa+\frac{\kappa^2}2\ge1\); record that \(C_\infty'\) bounds the
      Frobenius norm of \(\nabla V\) in the class and in the Duhamel bound.
\item Consequences rewritten: the gradient floor
      \(C_\infty'\ge(3/2)^{1/2}(1-\frac{\sqrt3}2\kappa-\frac{\kappa^2}2)\), the threshold below which the
      palinstrophy row is void, the slack \(\delta_{67}\) and the loop
      constraints in strain form.
\item The next compulsory companion audit is V95.
\end{enumerate}
\end{programme}

\begin{assessment}[Position after V90]
\label{ass:c3v90-position}
The companion audit is recorded; the strain-form block begins.  No
full stability or global-regularity claim is made.
\end{assessment}

% =============================================================================
\section*{Version V90 (third cycle) append-only record}
\addcontentsline{toc}{section}{Version V90 (third cycle) append-only record}
% =============================================================================

The complete V89 source of the third cycle was retained before this
continuation.  Its SHA--256 digest was
\begin{center}\scriptsize\ttfamily
ca78661f1a12c9020f52a7b5ec5aeac59b7ec21def419b801ef8cbc793a3ea31.
\end{center}
V90 performed the compulsory companion audit and recorded the
position.  No full regularity claim is made.

% =============================================================================
\section{V91: the strain form of the vorticity rows}
\label{sec:c3v91-entry}
% =============================================================================

Items (AC1) and (AC2) of Programme~\ref{prog:c3v91} are carried out.
Every stretching term of the vorticity ledger that is a quadratic
form in a vector field is bounded through the strain, the symmetric
traceless part of the velocity gradient, whose largest eigenvalue
is at most \((2/3)^{1/2}\) of its Frobenius norm; the enstrophy constraint
becomes \((2/3)^{1/2}C_\infty'+\frac{\sqrt3}2\kappa+\frac{\kappa^2}2\ge1\), and the weighted-enstrophy and
palinstrophy rows acquire the same factor on their \(C_\infty'\) terms.  The
consequences are rewritten: the gradient floor gains the factor
\((3/2)^{1/2}\), the threshold below which the palinstrophy row is void
is unchanged (\(0.4569\)) because both rows scale together, and every
statement of V73--V77 under (E5-small) holds with \(\gamma_1\) replaced by
\((2/3)^{1/2}\gamma_1\).  No global regularity claim is made.  The next
compulsory companion audit is V95.

% =============================================================================
\section{The strain bound}
\label{sec:c3v91-strain}
% =============================================================================

Throughout, \(C_\infty'\) bounds the Frobenius norm \(|\nabla V|_F\) on the family (the
Duhamel bounds of V81 and V86 are Frobenius bounds, and the class
bound of C2-V4 is at least as strong as the operator-norm bound
used in Lemma~\ref{lem:c3v67-bounds}).  Let \(S:=\frac12(\nabla V+\nabla V^{\rm T})\) be the
strain.

\begin{lemma}[Strain bound for quadratic forms]
\label{lem:c3v91-strain}
For every vector \(a\), \(|a\cdot(a\cdot\nabla)V|=|a^{\rm T}Sa|\le(2/3)^{1/2}|S|_F|a|^2\le(2/3)^{1/2}|\nabla V|_F|a|^2\).
\end{lemma}

\begin{proof}
\(a\cdot(a\cdot\nabla)V=a_ia_j\partial_jV_i=a^{\rm T}Aa\) with \(A_{ij}=\partial_jV_i\), and a quadratic form sees
only the symmetric part \(S\) of \(A\); \(\operatorname{tr}S=\operatorname{div}V=0\).  For symmetric
traceless \(S\) with eigenvalues \(\lambda_1\ge\lambda_2\ge\lambda_3\): \(\lambda_2+\lambda_3=-\lambda_1\), so
\(\lambda_2^2+\lambda_3^2\ge\frac12\lambda_1^2\) and \(|S|_F^2\ge\frac32\lambda_1^2\); likewise \(|S|_F^2\ge\frac32\lambda_3^2\); hence
\(|a^{\rm T}Sa|\le\max\{\lambda_1,-\lambda_3\}|a|^2\le(2/3)^{1/2}|S|_F|a|^2\).  Finally
\(|S|_F^2=\frac12(|A|_F^2+\operatorname{tr}(A^2))\le|A|_F^2\).
\end{proof}

% =============================================================================
\section{The rows in strain form}
\label{sec:c3v91-rows}
% =============================================================================

\begin{theorem}[Strain form of the enstrophy constraint]
\label{thm:c3v91-enstrophy}
For every profile, \(|\mathcal S|\le(2/3)^{1/2}C_\infty'\,\mathcal E\); hence every homogeneous
type-I point satisfies
\begin{equation}
 \boxed{\quad
 \Bigl(\frac23\Bigr)^{1/2}C_\infty'+\frac{\sqrt3}2\kappa+\frac{\kappa^2}2\ \ge\ 1 ,
 \quad}
 \label{eq:c3v91-constraint}
\end{equation}
and at an extremum of \(\mathcal E\), \(\frac12\mathcal E_\nabla\le\delta^S\mathcal E\) with
\(\delta^S:=(2/3)^{1/2}C_\infty'+\frac{\sqrt3}2\kappa+\frac{\kappa^2}2-1\ge0\).
\end{theorem}

\begin{proof}
\(\mathcal S=\int\Omega\cdot(\Omega\cdot\nabla)V\,G\) and Lemma~\ref{lem:c3v91-strain} with \(a=\Omega\); then
the proof of Theorem~\ref{thm:c3v67-extremal} with \(C_\infty'\) replaced by
\((2/3)^{1/2}C_\infty'\).
\end{proof}

\begin{theorem}[Strain form of the weighted-enstrophy and palinstrophy rows]
\label{thm:c3v91-rows}
For every profile, \(|\mathcal S_1|\le(2/3)^{1/2}C_\infty'\mathcal E_1\) and
\(|\mathcal S_{\mathcal P}|\le2(2/3)^{1/2}C_\infty'\mathcal P+C_\infty''\mathcal E^{1/2}\mathcal P^{1/2}\).  Consequently
Theorems~\ref{thm:c3v71-extremal} and \ref{thm:c3v72-extremal} hold with
\begin{equation}
 \theta_1^S:=\frac32-\Bigl(\frac23\Bigr)^{1/2}C_\infty'-\frac{\sqrt3}2\kappa-\frac\kappa{\sqrt2}-\kappa^2,\qquad
 \theta_{\mathcal P}^S:=\frac32-2\Bigl(\frac23\Bigr)^{1/2}C_\infty'-\frac{\sqrt3}2\kappa-\frac{\kappa^2}2
 \label{eq:c3v91-thetas}
\end{equation}
in place of \(\theta_1,\theta_{\mathcal P}\), and \(q_{\mathcal P}^S:=2\nu(C_\infty''/\theta_{\mathcal P}^S)^2\).
\end{theorem}

\begin{proof}
\(\mathcal S_1=\int|z|^2\Omega\cdot(\Omega\cdot\nabla)V\,G\): Lemma~\ref{lem:c3v91-strain} pointwise.  In
\(\mathcal S_{\mathcal P}\) (V72), the two terms \(\int\partial_j\Omega\cdot(\partial_j\Omega\cdot\nabla)V\,G\) and
\(-\int\partial_j\Omega\cdot(\partial_jV\cdot\nabla)\Omega\,G\) are, for each \(j\) and each component
respectively, quadratic forms: the first is \((\partial_j\Omega)^{\rm T}A(\partial_j\Omega)\) with
\(A_{ik}=\partial_kV_i\), the second is \(\sum_i(\nabla\Omega_i)^{\rm T}A(\nabla\Omega_i)\) since
\(\partial_j\Omega_i\,\partial_jV_k\,\partial_k\Omega_i=\sum_i(\nabla\Omega_i)_jA_{kj}(\nabla\Omega_i)_k\); both are bounded by
\((2/3)^{1/2}C_\infty'|\nabla\Omega|^2\); the third term of \(\mathcal S_{\mathcal P}\), \(\int\partial_j\Omega\cdot(\Omega\cdot\nabla)\partial_jV\,G\), is
not a quadratic form and keeps its bound \(C_\infty''\mathcal E^{1/2}\mathcal P^{1/2}\).
\end{proof}

% =============================================================================
\section{Consequences rewritten}
\label{sec:c3v91-consequences}
% =============================================================================

\begin{proposition}[Item (AC2): consequences in strain form]
\label{prop:c3v91-consequences}
At every homogeneous type-I point:
\begin{enumerate}[label=\textup{(\arabic*)},leftmargin=3.2em]
\item (Gradient floor.)  \(C_\infty'\ge(3/2)^{1/2}\bigl(1-\frac{\sqrt3}2\kappa-\frac{\kappa^2}2\bigr)\) for \(\kappa<\kappa_0=0.792\).
\item (Voidness.)  \(\theta_{\mathcal P}^S\le-\frac12+\frac{\sqrt3}2\kappa+\frac{\kappa^2}2<0\) for \(\kappa<0.4569\): the
      threshold of Theorem~\ref{thm:c3v84-void} is unchanged, because the
      strain factor enters the enstrophy constraint and the
      palinstrophy threshold in the ratio \(1:2\), as \(C_\infty'\) did.
\item (Under (E5-small).)  Every statement of V73, V76 and V77 holds
      with \(\gamma_1\) replaced by \((2/3)^{1/2}\gamma_1\): the floor is
      \(\kappa_1^S=-a^S+((a^S)^2+2)^{1/2}\) with \(a^S=(2/3)^{1/2}\gamma_1+\frac{\sqrt3}2\), the window and
      comparison thresholds of Theorem~\ref{thm:c3v73-window} and
      Propositions~\ref{prop:c3v73-epsreg}, \ref{prop:c3v73-chasing} apply to
      \((2/3)^{1/2}\gamma_1\), and the loop constraints of V74 and V77 hold with
      \(\delta^S\), \(q_{\mathcal P}^S\) in place of \(\delta_{67}\), \(q_{\mathcal P}\).
\item (Explicit constants.)  With the certified \(\gamma_1^{\rm red}\) of V86--V89,
      the floor is the least \(\kappa\) with
      \((2/3)^{1/2}\gamma_1^{\rm red}(\kappa)\kappa+\frac{\sqrt3}2\kappa+\frac{\kappa^2}2\ge1\): V92.
\end{enumerate}
\end{proposition}

\begin{proof}
(1) \eqref{eq:c3v91-constraint}.  (2) Substitute (1) in \(\theta_{\mathcal P}^S\):
\(\frac32-2(1-u)-u=-\frac12+u\) with \(u=\frac{\sqrt3}2\kappa+\frac{\kappa^2}2\).  (3) Each proof used
\(C_\infty'\) only through \(|\mathcal S|\), \(|\mathcal S_1|\), the \(C_\infty'\mathcal P\) term of \(|\mathcal S_{\mathcal P}|\) and the
hypothesis \(C_\infty'\le\gamma_1\kappa\).  (4) Theorem~\ref{thm:c3v86-absorptions}(2) in
\eqref{eq:c3v91-constraint}.
\end{proof}

\begin{assessment}[Reading]
\label{ass:c3v91-reading}
The strain bound is the vorticity-ledger counterpart of the
traceless reduction: there the pressure removed the trace of \(W\otimes W\),
here incompressibility removes the trace of the strain, and in both
cases the gain is the factor \((2/3)^{1/2}\) on the dominant term.  At the
floor of V89 the gradient term carried \(0.85\) of the unit; the strain
form lowers it to \(0.70\) at the same \(\kappa\), so the floor rises (to about
\(0.17\), V92).  The rows' structure is otherwise unchanged, and the
loops remain void below \(0.4569\).
\end{assessment}

% =============================================================================
\section{Integrated V91 conclusion and next acquisition order}
\label{sec:c3v91-conclusion}
% =============================================================================

\begin{theorem}[What V91 proves]
\label{thm:c3v91-conclusion}
Relative to the two legacy cycles and V1--V90: the strain bound
(Lemma~\ref{lem:c3v91-strain}), the strain forms of the enstrophy
constraint and of the weighted-enstrophy and palinstrophy rows
(Theorems~\ref{thm:c3v91-enstrophy}, \ref{thm:c3v91-rows}) and the
rewritten consequences (Proposition~\ref{prop:c3v91-consequences}).  The
full Navier--Stokes stability/global-regularity theorem remains
unproved.
\end{theorem}

\begin{proof}
The cited results.
\end{proof}

\begin{programme}[V92 acquisition order]
\label{prog:c3v92}
\begin{enumerate}[label=\textup{(AC\arabic*)},leftmargin=4.8em]
\item The certified strain-form floor: the least \(\kappa\) with
      \((2/3)^{1/2}\gamma_1^{\rm red}(\kappa)\kappa+\frac{\sqrt3}2\kappa+\frac{\kappa^2}2\ge1\) for \(c_O^{s0}=2.766\), with
      \(\theta^*\); the derivative bounds \(\gamma_1^{\rm red},\gamma_2^{\rm red}\) and the threshold
      \(\theta_1^S\) at the new floor; the final table rewritten.
\item Record the share of the unit carried by each term at the
      floor (strain, transport, damping).
\item The next compulsory companion audit is V95.
\end{enumerate}
\end{programme}

\begin{assessment}[Position after V91]
\label{ass:c3v91-position}
The rows are in strain form.  No full stability or global-regularity
claim is made.
\end{assessment}

% =============================================================================
\section*{Version V91 (third cycle) append-only record}
\addcontentsline{toc}{section}{Version V91 (third cycle) append-only record}
% =============================================================================

The complete V90 source of the third cycle (after the in-place
update of its companion-audit paragraph to SM V50, made before
delivery) was retained before this continuation.  Its SHA--256
digest was
\begin{center}\scriptsize\ttfamily
50b4154e59e7129d09227a910caa7cd295c12a2ab6b58705070a71c13ab22ac0.
\end{center}
V91 put the vorticity rows in strain form.  No full regularity claim
is made.

% =============================================================================
\section{V92: the certified strain-form floor}
\label{sec:c3v92-entry}
% =============================================================================

Items (AC1) and (AC2) of Programme~\ref{prog:c3v92} are carried out.
With the strain form of the enstrophy constraint and the certified
reduced gradient bound, every homogeneous type-I point satisfies
\(\kappa\ge0.1694\); the derivative bounds and the concentration threshold
are recomputed at the new floor, the table is rewritten, and the
share of the unit carried by each term at the floor is recorded:
strain \(0.839\), transport \(0.161\) (of which \(0.147\) from the Hardy term
\(\frac{\sqrt3}2\kappa\)).  No global regularity claim is made.  The next
compulsory companion audit is V95.

% =============================================================================
\section{The floor}
\label{sec:c3v92-floor}
% =============================================================================

Let \(\Gamma^S(\kappa):=(2/3)^{1/2}\gamma_1^{\rm red}(\kappa)\kappa+\frac{\sqrt3}2\kappa+\frac{\kappa^2}2\) with \(\gamma_1^{\rm red}(\kappa)\) the
infimum of \eqref{eq:c3v86-gamma1} over admissible \(\theta\) at \(c_O^{s0}=2.766\);
\(\Gamma^S\) is increasing.

\begin{theorem}[Certified strain-form floor]
\label{thm:c3v92-floor}
Every homogeneous type-I point of a smooth periodic solution
satisfies \(\Gamma^S(\kappa)\ge1\), hence
\begin{equation}
 \boxed{\quad
 \kappa\ \ge\ \kappa_1^S:=\min\{\kappa:\Gamma^S(\kappa)\ge1\}=0.1694
 \qquad(\theta^*=0.161,\ \gamma_1^{\rm red}=6.07,\ \gamma_1^{\rm red}\kappa=1.028),
 \quad}
 \label{eq:c3v92-floor}
\end{equation}
and at the floor \(\nu^{1/2}C_\infty''\le11.53\) (\(\gamma_2^{\rm red}=68.1\), \(\theta=0.084\), \(c_{h,2}=1.794\)),
\(\theta_1^S\ge0.366\), and the palinstrophy row is void.
\end{theorem}

\begin{proof}
\eqref{eq:c3v91-constraint} with Theorem~\ref{thm:c3v86-absorptions}(2) at
\(c_O^{s0}\le2.766\) (Theorem~\ref{thm:c3v89-certified}); the bisection and
minimizations of V83--V84 with the factor \((2/3)^{1/2}\) on the gradient
term; \(\theta_1^S\) from \eqref{eq:c3v91-thetas}; voidness from
Proposition~\ref{prop:c3v91-consequences}(2).
\end{proof}

\begin{proposition}[Item (AC2): shares at the floor]
\label{prop:c3v92-shares}
At \(\kappa=\kappa_1^S\) the constraint \eqref{eq:c3v91-constraint} is an equality
with the shares
\begin{equation}
 \Bigl(\frac23\Bigr)^{1/2}\gamma_1^{\rm red}\kappa=0.839,\qquad \frac{\sqrt3}2\kappa=0.147,\qquad \frac{\kappa^2}2=0.014 .
 \label{eq:c3v92-shares}
\end{equation}
The first is the strain (stretching) term, the second and third the
transport term as bounded in Lemma~\ref{lem:c3v67-bounds} through the
Gaussian Hardy inequality; the damping is the unit.
\end{proposition}

\begin{proof}
Evaluation at the floor.
\end{proof}

\begin{theorem}[The table after the strain form]
\label{thm:c3v92-table}
Every homogeneous type-I point satisfies, all certified:

\medskip
\begin{center}\small
\begin{tabular}{@{}p{0.5\textwidth}p{0.16\textwidth}p{0.2\textwidth}@{}}
\toprule
\textbf{Statement} & \textbf{Value} & \textbf{Source}\\
\midrule
\(\kappa\ge1/(2(8/3)^{1/2}c_O^{s0})\) (gradient Liouville) & \(0.1107\) & V83, V86, V89\\
\(\kappa\ge(3/2)^{1/2}/(\pi c_O^{s0})\) (velocity Liouville) & \(0.1410\) & V81, V86, V89\\
\(\kappa\ge\kappa_1^{\rm red}\) (enstrophy row, gradient form) & \(0.1565\) & V89\\
\(\kappa\ge\kappa_1^S\) (enstrophy row, strain form; \(\theta^*=0.161\)) & \(\mathbf{0.1694}\) & V91, V92\\
\(C_\infty'\le\gamma_1^{\rm red}\kappa\) at the floor & \(1.028\) & V86\\
\(\nu^{1/2}C_\infty''\le\gamma_2^{\rm red}\kappa\) at the floor & \(11.53\) & V84, V86\\
concentration row applies at the floor & \(\theta_1^S\ge0.366\) & V71, V91\\
palinstrophy row and loop constraints & void for \(\kappa<0.4569\) & V84, V91\\
\bottomrule
\end{tabular}
\end{center}
\medskip

\noindent
In particular \(C_\infty\ge0.1694\,\nu^{1/2}\) at every homogeneous type-I singular
point; the floor exceeds the velocity Liouville floor by the factor
\(1.20\).
\end{theorem}

\begin{proof}
Theorems~\ref{thm:c3v89-table}, \ref{thm:c3v92-floor}.
\end{proof}

\begin{assessment}[Reading]
\label{ass:c3v92-reading}
The strain form raised the floor by \(8\%\).  At the floor the transport
term carries \(0.16\) of the unit, nine tenths of it from the term
\(\frac{\sqrt3}2\kappa\) produced by the Gaussian Hardy inequality in
Lemma~\ref{lem:c3v67-bounds}; V93 examines whether that route is the
right one for the transport term, the alternative being the
transport form \(\int\Omega\cdot(V\cdot\nabla)\Omega\,G\) bounded directly by Cauchy--Schwarz.
\end{assessment}

% =============================================================================
\section{Integrated V92 conclusion and next acquisition order}
\label{sec:c3v92-conclusion}
% =============================================================================

\begin{theorem}[What V92 proves]
\label{thm:c3v92-conclusion}
Relative to the two legacy cycles and V1--V91: the certified
strain-form floor (Theorem~\ref{thm:c3v92-floor}), the shares
(Proposition~\ref{prop:c3v92-shares}) and the table
(Theorem~\ref{thm:c3v92-table}).  The full Navier--Stokes
stability/global-regularity theorem remains unproved.
\end{theorem}

\begin{proof}
The cited results.
\end{proof}

\begin{programme}[V93 acquisition order]
\label{prog:c3v93}
\begin{enumerate}[label=\textup{(AC\arabic*)},leftmargin=4.8em]
\item The transport term in transport form: from \((z\cdot V)G=-2\nu V\cdot\nabla G\)
      and \(\operatorname{div}V=0\), \(\mathcal T=\int\Omega\cdot(V\cdot\nabla)\Omega\,G\), so \(|\mathcal T|\le C_\infty\mathcal E^{1/2}\mathcal P^{1/2}\le\frac12\mathcal E_\nabla+\frac{\kappa^2}2\mathcal E\),
      without the Hardy term; the same for \(\mathcal T_{\mathcal P}\), and for the weighted
      row, where \(-\mathcal T_1+4\nu\mathcal T=-\int|z|^2\Omega\cdot(V\cdot\nabla)\Omega\,G\) exactly, bounded by
      \(\frac12\mathcal E_{\nabla,1}+\frac{\kappa^2}2\mathcal E_1\).  Rewrite the three rows and their
      thresholds; record the new voidness threshold of the
      palinstrophy row.
\item The constraint \((2/3)^{1/2}C_\infty'+\frac{\kappa^2}2\ge1\) and its consequences under
      (E5-small) and with the explicit constants (the floor, V94).
\item The next compulsory companion audit is V95.
\end{enumerate}
\end{programme}

\begin{assessment}[Position after V92]
\label{ass:c3v92-position}
The certified floor is \(0.1694\).  No full stability or
global-regularity claim is made.
\end{assessment}

% =============================================================================
\section*{Version V92 (third cycle) append-only record}
\addcontentsline{toc}{section}{Version V92 (third cycle) append-only record}
% =============================================================================

The complete V91 source of the third cycle was retained before this
continuation.  Its SHA--256 digest was
\begin{center}\scriptsize\ttfamily
9f9fb7dce68cb3a6bbe4ebe134ba5f8f2d19af420ef043756dae7e24aa860bb7.
\end{center}
V92 certified the strain-form floor.  No full regularity claim is
made.

% =============================================================================
\section{V93: the transport form of the transport terms}
\label{sec:c3v93-entry}
% =============================================================================

Items (AC1) and (AC2) of Programme~\ref{prog:c3v93} are carried out.
The transport terms of the three vorticity rows, written since V67
in the form \(\frac1{4\nu}\int|\cdot|^2(z\cdot V)G\) and bounded through the Gaussian Hardy
inequality, are equal to \(\int(\cdot)\cdot(V\cdot\nabla)(\cdot)\,G\) and are bounded directly by
Cauchy--Schwarz with the constant \(\frac{\kappa^2}2\) on the observable and \(\frac12\)
on its dissipation, without the term \(\frac{\sqrt3}2\kappa\) (and, for the weighted
row, with the \(4\nu\mathcal T\) term cancelling exactly).  The enstrophy constraint
becomes \((2/3)^{1/2}C_\infty'+\frac{\kappa^2}2\ge1\); the concentration bound loses its
\(C_\infty''\) term; the palinstrophy row and the loop constraints are void
at every homogeneous type-I point with \(\kappa<1\).  The bounds of V67,
V71 and V72 are not wrong and are not edited; they are superseded.
No global regularity claim is made.  The next compulsory companion
audit is V95.

% =============================================================================
\section{The transport identities}
\label{sec:c3v93-identities}
% =============================================================================

\begin{lemma}[Transport form]
\label{lem:c3v93-transport}
For every profile, with \(\widetilde{\mathcal T}_1:=\int|z|^2\,\Omega\cdot(V\cdot\nabla)\Omega\,G\),
\begin{equation}
 \mathcal T=\int\Omega\cdot(V\cdot\nabla)\Omega\,G,\qquad
 \mathcal T_{\mathcal P}=\int\partial_j\Omega\cdot(V\cdot\nabla)\partial_j\Omega\,G,\qquad
 -\mathcal T_1+4\nu\mathcal T=-\widetilde{\mathcal T}_1 .
 \label{eq:c3v93-identities}
\end{equation}
\end{lemma}

\begin{proof}
\((z\cdot V)G=-2\nu\,V\cdot\nabla G\).  For a scalar \(F\) with \(\operatorname{div}V=0\):
\(\frac1{4\nu}\int F(z\cdot V)G=-\frac12\int F\,V\cdot\nabla G=\frac12\int G\operatorname{div}(FV)=\frac12\int G\,V\cdot\nabla F\).  With
\(F=|\Omega|^2\): \(\frac12V\cdot\nabla|\Omega|^2=\Omega\cdot(V\cdot\nabla)\Omega\).  With \(F=|\nabla\Omega|^2\): likewise.
With \(F=|z|^2|\Omega|^2\): \(\frac12V\cdot\nabla(|z|^2|\Omega|^2)=(z\cdot V)|\Omega|^2+|z|^2\Omega\cdot(V\cdot\nabla)\Omega\), and
\(\int(z\cdot V)|\Omega|^2G=4\nu\mathcal T\); so \(\mathcal T_1=4\nu\mathcal T+\widetilde{\mathcal T}_1\).
\end{proof}

\begin{lemma}[Transport bounds]
\label{lem:c3v93-bounds}
For every profile,
\begin{equation}
 |\mathcal T|\le\frac12\mathcal E_\nabla+\frac{\kappa^2}2\mathcal E,\qquad
 |\widetilde{\mathcal T}_1|\le\frac12\mathcal E_{\nabla,1}+\frac{\kappa^2}2\mathcal E_1,\qquad
 |\mathcal T_{\mathcal P}|\le\frac12\mathcal P_\nabla+\frac{\kappa^2}2\mathcal P .
 \label{eq:c3v93-bounds}
\end{equation}
\end{lemma}

\begin{proof}
\(|\Omega\cdot(V\cdot\nabla)\Omega|\le|V||\Omega||\nabla\Omega|\le C_\infty|\Omega||\nabla\Omega|\), so
\(|\mathcal T|\le C_\infty\|\Omega\|_G\|\nabla\Omega\|_G=\kappa\nu^{1/2}\mathcal E^{1/2}(\mathcal E_\nabla/\nu)^{1/2}\le\frac12\mathcal E_\nabla+\frac{\kappa^2}2\mathcal E\).  For
\(\widetilde{\mathcal T}_1\): \(C_\infty\||z|\Omega\|_G\||z|\nabla\Omega\|_G=\kappa\nu^{1/2}\mathcal E_1^{1/2}(\mathcal E_{\nabla,1}/\nu)^{1/2}\).  For \(\mathcal T_{\mathcal P}\): with
\(\nabla\Omega\) in place of \(\Omega\).
\end{proof}

% =============================================================================
\section{The three rows in strain--transport form}
\label{sec:c3v93-rows}
% =============================================================================

\begin{theorem}[The rows in strain--transport form]
\label{thm:c3v93-rows}
At a homogeneous type-I point, with
\begin{equation}
 \theta_0^{ST}:=1-\Bigl(\frac23\Bigr)^{1/2}C_\infty'-\frac{\kappa^2}2,\qquad
 \theta_1^{ST}:=\frac32-\Bigl(\frac23\Bigr)^{1/2}C_\infty'-\frac{\kappa^2}2,\qquad
 \theta_{\mathcal P}^{ST}:=\frac32-2\Bigl(\frac23\Bigr)^{1/2}C_\infty'-\frac{\kappa^2}2 :
 \label{eq:c3v93-thetas}
\end{equation}
\begin{enumerate}[label=\textup{(\arabic*)},leftmargin=3.2em]
\item (Enstrophy.)  At an extremum of \(\mathcal E\), \(\frac12\mathcal E_\nabla+\theta_0^{ST}\mathcal E\le0\); hence
      \begin{equation}
       \boxed{\quad
       \Bigl(\frac23\Bigr)^{1/2}C_\infty'+\frac{\kappa^2}2\ \ge\ 1 .
       \quad}
       \label{eq:c3v93-constraint}
      \end{equation}
\item (Weighted enstrophy.)  At an extremum of \(\mathcal E_1\),
      \(\frac12\mathcal E_{\nabla,1}+\theta_1^{ST}\mathcal E_1\le3\nu\mathcal E\); hence, if \(\theta_1^{ST}>0\),
      \(\sup\mathcal E_1\le3\nu\mathcal E(W_1)/\theta_1^{ST}\le6\nu C/\theta_1^{ST}\).
\item (Palinstrophy.)  At an extremum of \(\mathcal P\),
      \(\frac12\mathcal P_\nabla+\theta_{\mathcal P}^{ST}\mathcal P\le C_\infty''\mathcal E^{1/2}\mathcal P^{1/2}\); the ratio bound of V72 holds with
      \(\theta_{\mathcal P}^{ST}\), where positive.
\end{enumerate}
\end{theorem}

\begin{proof}
The identities of V67, V71 (\eqref{eq:c3v71-identity}) and V72, with
\(|\mathcal S|\), \(|\mathcal S_1|\), \(|\mathcal S_{\mathcal P}|\) bounded by Theorems~\ref{thm:c3v91-enstrophy},
\ref{thm:c3v91-rows} and the transport terms by
Lemmas~\ref{lem:c3v93-transport}, \ref{lem:c3v93-bounds}; in the weighted
row, \(-\mathcal T_1+4\nu\mathcal T=-\widetilde{\mathcal T}_1\) so that no bound on \(\mathcal T\) and no \(\mathcal E_\nabla\)
term enters; \(\mathcal E\le2C\) as in V71.
\end{proof}

\begin{proposition}[Item (AC2): consequences]
\label{prop:c3v93-consequences}
At every homogeneous type-I point:
\begin{enumerate}[label=\textup{(\arabic*)},leftmargin=3.2em]
\item (Gradient floor.)  \(C_\infty'\ge(3/2)^{1/2}(1-\frac{\kappa^2}2)\) for \(\kappa<\sqrt2\).
\item (Concentration.)  \(\theta_1^{ST}=\theta_0^{ST}+\frac12\le\frac12\), and wherever \(\theta_1^{ST}>0\),
      \(\sup\mathcal E_1\le6\nu C/\theta_1^{ST}\), without any derivative constant.
\item (Voidness.)  \(\theta_{\mathcal P}^{ST}\le-\frac12+\frac{\kappa^2}2<0\) for \(\kappa<1\): the palinstrophy
      row, the loop constraints of V74 and V77 and the floors of
      V76--V77 are void at every homogeneous type-I point with \(\kappa<1\).
\item (Under (E5-small).)  \(\kappa\ge\min\{\kappa_1^{ST},\kappa_{\rm small}\}\) with
      \(\kappa_1^{ST}=-a+(a^2+2)^{1/2}\), \(a=(2/3)^{1/2}\gamma_1\), i.e.\ \(\kappa_1^{ST}\in[2/(2a+\sqrt2),1/a)\); for
      \(\gamma_1\to0\) the floor tends to \(\sqrt2\) (against \(0.792\) in V73).
\end{enumerate}
\end{proposition}

\begin{proof}
(1) \eqref{eq:c3v93-constraint}.  (2) \(\theta_1^{ST}-\theta_0^{ST}=\frac12\) and \(\theta_0^{ST}\le0\)
by (1); Theorem~\ref{thm:c3v93-rows}(2).  (3) Substitute (1) in \(\theta_{\mathcal P}^{ST}\):
\(\frac32-2(1-\frac{\kappa^2}2)-\frac{\kappa^2}2=-\frac12+\frac{\kappa^2}2\).  (4) As in
Theorem~\ref{thm:c3v73-velocity-floor} with \(f(\kappa)=a\kappa+\kappa^2/2\).
\end{proof}

\begin{assessment}[Reading]
\label{ass:c3v93-reading}
The transport term of the enstrophy identity is the derivative of
\(|\Omega|^2\) along \(V\); bounding it by Cauchy--Schwarz against the
dissipation is the natural estimate, and the Hardy route of V67,
which first moved the derivative onto the Gaussian weight, cost the
term \(\frac{\sqrt3}2\kappa\).  With the strain form, the constraint
\eqref{eq:c3v93-constraint} says: at a homogeneous type-I point the
strain constant is at least \((3/2)^{1/2}(1-\kappa^2/2)\); for a small velocity
constant the strain is of order one.  Every downstream statement
improves: the concentration bound is now a pure ratio \(6\nu C/\theta_1^{ST}\)
with \(\theta_1^{ST}\ge\frac12-\)(explicit) and equal to \(\frac12\) at the explicit floor; the
palinstrophy row is void up to \(\kappa=1\), which confines the loop block
to \(\kappa\ge1\); the (E5-small) floor \(\kappa_1^{ST}\) is larger.  The explicit floor
is V94.
\end{assessment}

% =============================================================================
\section{Integrated V93 conclusion and next acquisition order}
\label{sec:c3v93-conclusion}
% =============================================================================

\begin{theorem}[What V93 proves]
\label{thm:c3v93-conclusion}
Relative to the two legacy cycles and V1--V92: the transport
identities and bounds (Lemmas~\ref{lem:c3v93-transport},
\ref{lem:c3v93-bounds}), the rows in strain--transport form
(Theorem~\ref{thm:c3v93-rows}) and the consequences
(Proposition~\ref{prop:c3v93-consequences}).  The full Navier--Stokes
stability/global-regularity theorem remains unproved.
\end{theorem}

\begin{proof}
The cited results.
\end{proof}

\begin{programme}[V94 acquisition order]
\label{prog:c3v94}
\begin{enumerate}[label=\textup{(AC\arabic*)},leftmargin=4.8em]
\item The certified floor of \eqref{eq:c3v93-constraint} with
      \(\gamma_1^{\rm red}\) at \(c_O^{s0}=2.766\), its \(\theta^*\), the derivative bounds and the
      concentration bound at the floor; the final table.
\item The constants sheet of the explicit chapter (V81--V94): every
      constant, its value, its status (closed form, certified
      quadrature, hand bound) and where it is proved; the shares at
      the floor; the record of what would move the floor next.
\item The next compulsory companion audit is V95.
\end{enumerate}
\end{programme}

\begin{assessment}[Position after V93]
\label{ass:c3v93-position}
The transport terms are in transport form; the constraint is
\((2/3)^{1/2}C_\infty'+\frac{\kappa^2}2\ge1\).  No full stability or global-regularity claim
is made.
\end{assessment}

% =============================================================================
\section*{Version V93 (third cycle) append-only record}
\addcontentsline{toc}{section}{Version V93 (third cycle) append-only record}
% =============================================================================

The complete V92 source of the third cycle was retained before this
continuation.  Its SHA--256 digest was
\begin{center}\scriptsize\ttfamily
f4949f7e9f81dfb427605b413690691ab4d201713089d527a80ca5b6dde6dd0e.
\end{center}
V93 put the transport terms in transport form.  No full regularity
claim is made.

% =============================================================================
\section{V94: the certified floor in strain--transport form, and the constants sheet}
\label{sec:c3v94-entry}
% =============================================================================

Items (AC1) and (AC2) of Programme~\ref{prog:c3v94} are carried out.
With the constraint \((2/3)^{1/2}C_\infty'+\frac{\kappa^2}2\ge1\) and the certified reduced
gradient bound, every homogeneous type-I point of a smooth periodic
solution satisfies \(\kappa\ge0.1815\); at the floor the strain term carries
\(0.98\) of the unit, the concentration bound reads
\(\sup\int|z|^2|\Omega|^2G\le12\nu C\), and the palinstrophy row is void.  The
constants sheet of the explicit chapter is written.  No global
regularity claim is made.  The next compulsory companion audit is
V95.

% =============================================================================
\section{The floor}
\label{sec:c3v94-floor}
% =============================================================================

Let \(\Gamma^{ST}(\kappa):=(2/3)^{1/2}\gamma_1^{\rm red}(\kappa)\kappa+\frac{\kappa^2}2\), \(\gamma_1^{\rm red}\) at \(c_O^{s0}=2.766\), increasing
in \(\kappa\).

\begin{theorem}[Certified floor in strain--transport form]
\label{thm:c3v94-floor}
Every homogeneous type-I point of a smooth periodic solution
satisfies \(\Gamma^{ST}(\kappa)\ge1\), hence
\begin{equation}
 \boxed{\quad
 \kappa\ \ge\ \kappa_1^{ST}:=\min\{\kappa:\Gamma^{ST}(\kappa)\ge1\}=0.1815
 \qquad(\theta^*=0.131,\ \gamma_1^{\rm red}=6.64,\ \gamma_1^{\rm red}\kappa=1.205),
 \quad}
 \label{eq:c3v94-floor}
\end{equation}
i.e.\ \(C_\infty\ge0.1815\,\nu^{1/2}\).  At the floor: \(\nu^{1/2}C_\infty''\le14.63\) (\(\gamma_2^{\rm red}=80.6\),
\(\theta=0.070\), \(c_{h,2}=1.794\)); \(\theta_1^{ST}\ge\frac12\), so \(\sup_{\mathcal F(x_0)}\int|z|^2|\Omega|^2G\le12\nu C\); the
palinstrophy row is void (\(\kappa<1\)).  The shares of the unit are
\((2/3)^{1/2}\gamma_1^{\rm red}\kappa=0.984\) (strain) and \(\frac{\kappa^2}2=0.016\) (transport).
\end{theorem}

\begin{proof}
\eqref{eq:c3v93-constraint} with Theorem~\ref{thm:c3v86-absorptions}(2) at
\(c_O^{s0}\le2.766\); the bisection of V83 with \(\Gamma^{ST}\); \(\gamma_2^{\rm red}\) from
\eqref{eq:c3v86-gamma2}; \(\theta_1^{ST}=\frac32-(\text{gradient term}+\frac{\kappa^2}2)\ge\frac32-1\) at the
floor, where \(\Gamma^{ST}=1\); Theorem~\ref{thm:c3v93-rows}(2) and
Proposition~\ref{prop:c3v93-consequences}(3).
\end{proof}

\begin{theorem}[The final table of the explicit chapter]
\label{thm:c3v94-table}
Every homogeneous type-I point satisfies, all certified:

\medskip
\begin{center}\small
\begin{tabular}{@{}p{0.52\textwidth}p{0.14\textwidth}p{0.2\textwidth}@{}}
\toprule
\textbf{Statement} & \textbf{Value} & \textbf{Source}\\
\midrule
\(\kappa\ge1/(2(8/3)^{1/2}c_O^{s0})\) (gradient Liouville) & \(0.1107\) & V83, V86, V89\\
\(\kappa\ge(3/2)^{1/2}/(\pi c_O^{s0})\) (velocity Liouville) & \(0.1410\) & V81, V86, V89\\
\(\kappa\ge\kappa_1^{\rm red}\) (enstrophy row, gradient form) & \(0.1565\) & V89\\
\(\kappa\ge\kappa_1^S\) (strain form) & \(0.1694\) & V92\\
\(\kappa\ge\kappa_1^{ST}\) (strain--transport form; \(\theta^*=0.131\)) & \(\mathbf{0.1815}\) & V93, V94\\
\(C_\infty'\le\gamma_1^{\rm red}\kappa\) at the floor & \(1.205\) & V86\\
\(\nu^{1/2}C_\infty''\le\gamma_2^{\rm red}\kappa\) at the floor & \(14.63\) & V84, V86\\
\(\sup\int|z|^2|\Omega|^2G\le12\nu C\) at the floor (\(\theta_1^{ST}\ge\frac12\)) & --- & V71, V93\\
palinstrophy row and loop constraints & void for \(\kappa<1\) & V93\\
\bottomrule
\end{tabular}
\end{center}
\medskip

\noindent
The floor exceeds the velocity Liouville floor by the factor \(1.29\).
\end{theorem}

\begin{proof}
Theorems~\ref{thm:c3v92-table}, \ref{thm:c3v94-floor}.
\end{proof}

% =============================================================================
\section{The constants sheet of the explicit chapter}
\label{sec:c3v94-sheet}
% =============================================================================

\begin{theorem}[Item (AC2): constants sheet, V81--V94]
\label{thm:c3v94-sheet}
\medskip
\begin{center}\small
\begin{tabular}{@{}p{0.2\textwidth}p{0.36\textwidth}p{0.12\textwidth}p{0.2\textwidth}@{}}
\toprule
\textbf{Constant} & \textbf{Definition} & \textbf{Value} & \textbf{Status; where}\\
\midrule
\(c_h\) & \(\sup_t t^{1/2}\|\nabla K_t\|_{L^1}=2/\sqrt\pi\) & \(1.1284\) & closed form; V81\\
\(c_{h,2}\) & \(\sup_t t\,\|\,|\nabla^2K_t|_F\|_{L^1}=\frac12\mathbb E|\zeta\zeta^{\rm T}-I|_F\) & \(\le1.7936\) & closed-form bound; V84\\
 & & \(1.4149\) & quadrature (not used)\\
\(c_O\) & \(\|e^{\Delta}\PP\operatorname{div}\|_{\infty\to\infty}\) & \(\le4.208\) & hand (Plancherel); V82\\
 & & \(\le2.836\) & certified quadrature; V89\\
\(c_O^{s0}\) & same, symmetric traceless inputs & \(\le2.766\) & certified quadrature; V87, V89\\
\(\|M\|_{s0}(r)\) & \(\max\{\sqrt6|c|,|2c-h_K|/\sqrt2\}\), \(\phi'''-h_K=-2c\) & exact & V82, V87\\
tail & \(\|M\|\le(\sqrt{54}+10^{-9})/(4\pi r^4)\), \(r\ge12\) & exact & V82\\
Lipschitz & \(|G'|\le97\) on \([0,12]\) & exact & V89\\
\(\gamma_1^{\rm red}(\kappa,\theta)\) & \(c_h(\theta(1+\theta))^{-1/2}/(1-2(8/3)^{1/2}c_O^{s0}\kappa(\theta/(1+\theta))^{1/2})\) & formula & V81, V86\\
\(\gamma_2^{\rm red}(\kappa,\theta)\) & \eqref{eq:c3v86-gamma2} & formula & V84, V86\\
\(J_2(\theta)\) & \(\theta^{1/2}/(1+\theta)+\arctan\theta^{1/2}\) & closed form & V84\\
strain factor & \(\lambda_{\max}(S)\le(2/3)^{1/2}|S|_F\) & exact & V91\\
transport & \(|\mathcal T|\le\frac12\mathcal E_\nabla+\frac{\kappa^2}2\mathcal E\) & exact & V93\\
constraint & \((2/3)^{1/2}C_\infty'+\frac{\kappa^2}2\ge1\) & exact & V93\\
floor & \(\kappa\ge0.1815\), \(\theta^*=0.131\) & certified & V94\\
void threshold & palinstrophy row void for \(\kappa<1\) & exact & V93\\
\bottomrule
\end{tabular}
\end{center}
\medskip
\end{theorem}

\begin{proof}
The cited results.
\end{proof}

\begin{assessment}[What would move the floor next]
\label{ass:c3v94-next}
At the floor the constraint is \((2/3)^{1/2}\gamma_1^{\rm red}\kappa=0.98\): the floor is
decided by the gradient bound alone, and the vorticity ledger's
constants (strain, transport, damping) are now exact.  The gradient
bound \eqref{eq:c3v86-gamma1} has three parts: the heat-kernel term
\(c_h(\theta(1+\theta))^{-1/2}\), the absorption coefficient \(2(8/3)^{1/2}c_O^{s0}\kappa(\theta/(1+\theta))^{1/2}\),
and the window optimization.  Each is at its sup-norm optimum; what
remains is information the sup-norm argument discards: (i) the \(L^6\)
bound on the far part of the kernel action (needs \(K_6\), a class
constant from the enstrophy rate, not numerical); (ii) the Gaussian
energy bound on \(\nabla W\) (needs \(C\), likewise); (iii) the ancient form of
\(W(s_0)\) in the initial term, which replaces \(\|W(s_0)\|_\infty\) by the Oseen
integral and trades \(c_h\) against \(c_O\).  These are recorded for the
fourth cycle.  The reading of the number itself: a homogeneous
type-I singularity of a smooth periodic solution, if one exists, has
rescaled velocity bound \(C_\infty\ge0.18\,\nu^{1/2}\), strain constant of order
one, and a Gaussian second vorticity moment bounded by twelve
times the enstrophy scale; it is not excluded.
\end{assessment}

% =============================================================================
\section{Integrated V94 conclusion and next acquisition order}
\label{sec:c3v94-conclusion}
% =============================================================================

\begin{theorem}[What V94 proves]
\label{thm:c3v94-conclusion}
Relative to the two legacy cycles and V1--V93: the certified floor
(Theorem~\ref{thm:c3v94-floor}), the final table
(Theorem~\ref{thm:c3v94-table}) and the constants sheet
(Theorem~\ref{thm:c3v94-sheet}).  The full Navier--Stokes
stability/global-regularity theorem remains unproved.
\end{theorem}

\begin{proof}
The cited results.
\end{proof}

\begin{programme}[V95 acquisition order]
\label{prog:c3v95}
\begin{enumerate}[label=\textup{(AC\arabic*)},leftmargin=4.8em]
\item The compulsory companion audit: read the current SM and YM
      notes and record their digests and decision.
\item Record the position after ninety-five versions and the plan
      of the closing block V96--V100; delivery.
\end{enumerate}
\end{programme}

\begin{assessment}[Position after V94]
\label{ass:c3v94-position}
The explicit chapter is complete at the level of the sup-norm mild
formulation: certified floor \(0.1815\).  No full stability or
global-regularity claim is made.
\end{assessment}

% =============================================================================
\section*{Version V94 (third cycle) append-only record}
\addcontentsline{toc}{section}{Version V94 (third cycle) append-only record}
% =============================================================================

The complete V93 source of the third cycle was retained before this
continuation.  Its SHA--256 digest was
\begin{center}\scriptsize\ttfamily
d4bbf57d3a9010d29039d3af7949d819b795ae30cb69451b93fd9a5a02265bb2.
\end{center}
V94 certified the floor in strain--transport form and wrote the
constants sheet.  No full regularity claim is made.

% =============================================================================
\section{V95: compulsory companion audit and the position after ninety-five versions}
\label{sec:c3v95-entry}
% =============================================================================

This is the ninety-fifth version of the third cycle.  The compulsory
review of the two companion programmes is performed first, and the
position after ninety-five versions is recorded with the plan of
the closing block V96--V100.  No global regularity claim is made.
The next compulsory companion audit is V100.

% =============================================================================
\section{Compulsory V95 review of the current companion notes}
\label{sec:c3v95-companion-audit}
% =============================================================================

The current companion files in \texttt{latex\_notes} were inspected
read-only.  Since the V90 audit the SM programme has not moved (its
file is still V50, with the digest of \eqref{eq:c3v90-SM-checkpoint},
verified at this audit); the YM programme has moved from V26 to V27
of its sixth series.  The frozen checkpoints are
\begin{align}
 &\texttt{Renewal\_SM\_Einstein\_Continuum\_Research\_Notes\_V50.tex},
 \notag\\[-2pt]
 &\qquad
 \texttt{5931c1a01125c21735ac326608ef856f01124ac9d2439155a5c4cd04871495f2},
 \label{eq:c3v95-SM-checkpoint}\\
 &\texttt{Renewal\_Yang\_Mills\_Mass\_Gap\_Research\_Notes\_V27.tex},
 \notag\\[-2pt]
 &\qquad
 \texttt{c6b950842a93b5614903bd0d4c4bff20415e3733eddb81a8fb492bf65eb46f49}.
 \label{eq:c3v95-YM-checkpoint}
\end{align}
(SM V50 has 19900 lines; YM V27 has 11985 lines.  The folder also
held YM V26, the state recorded at V90, superseded.)

\emph{YM V27.}  The YM programme treated a relative centre-flux
erasure and a Peierls sector: small-tube centre labels obeying a
Bianchi cocycle with dual closed surfaces, the relative exactness
required by its fixed tree and block map, a Haar-preserving erasure
with a per-plaquette action gain, a centre-sheet Peierls bound with
explicit thresholds, a centre-sheet connectivity tail, and its
combination with good-sector paths into an Osterwalder--Schrader
spectral rate.

\begin{theorem}[V95 companion-audit decision]
\label{thm:c3v95-companion-decision}
Nothing is imported from SM V50 or YM V27.  Neither bears on the
third cycle's chapters or on the explicit chapter.
\end{theorem}

\begin{proof}
The SM file is unchanged since Theorem~\ref{thm:c3v90-companion-decision}.
The YM results concern centre fluxes, Peierls bounds and spectral
rates of a lattice gauge theory; none is a Navier--Stokes statement,
and none concerns the mild formulation, the Oseen constant, the
vorticity rows or the explicit floors.
\end{proof}

% =============================================================================
\section{Position after ninety-five versions, and the closing block}
\label{sec:c3v95-position}
% =============================================================================

\begin{assessment}[Position after ninety-five versions of the third cycle]
\label{ass:c3v95-position-ninetyfive}
Seventeen chapters are written: those of
Assessment~\ref{ass:c3v90-position-ninety} and the strain--transport
block (V91--V94).  The explicit chapter (V81--V94) closes with the
certified statement: every homogeneous type-I point of a smooth
periodic solution has rescaled velocity bound \(C_\infty\ge0.1815\,\nu^{1/2}\),
strain constant at least \((3/2)^{1/2}(1-\kappa^2/2)\), and Gaussian second
vorticity moment at most \(12\nu C\) at the floor; the palinstrophy row
and the loops are void below \(\kappa=1\).  The chain of the number is
recorded in the constants sheet (Theorem~\ref{thm:c3v94-sheet}).  The
closing block V96--V100 is: the manifest of the third cycle's
results by chapter (V96); the cycle's negative records collected,
with the reasons (V97); the constants sheet of the whole cycle and
the corrections log (V98); the draft of the fourth cycle's V1, with
the summary of major results and the directions recorded for it
(V99); the compulsory audit and the closure of the cycle (V100),
after which the file is frozen as the third legacy file and the
fourth cycle begins from a new V1.  No global regularity claim is
made.
\end{assessment}

% =============================================================================
\section{Integrated V95 conclusion and next acquisition order}
\label{sec:c3v95-conclusion}
% =============================================================================

\begin{theorem}[What V95 establishes]
\label{thm:c3v95-conclusion}
Relative to the two legacy cycles and V1--V94: the companion-audit
decision (Theorem~\ref{thm:c3v95-companion-decision}) and the position.
The full Navier--Stokes stability/global-regularity theorem remains
unproved.
\end{theorem}

\begin{proof}
The cited records.
\end{proof}

\begin{programme}[V96 acquisition order]
\label{prog:c3v96}
\begin{enumerate}[label=\textup{(AC\arabic*)},leftmargin=4.8em]
\item The manifest of the third cycle: for each chapter (V2--V10,
      V11--V14, V16--V19, V21--V24, V26--V29, V31--V34, V34--V39,
      V41--V44, V46--V49, V51--V55, V56--V59, V61--V64, V66--V69,
      V71--V74, V76--V79, V81--V84, V86--V89, V91--V94), the results
      that stand, with their labels, in one table each; no result
      re-proved, no result re-stated beyond its name and label.
\item The status line of each chapter: explicit, conditional,
      negative.
\item The next compulsory companion audit is V100.
\end{enumerate}
\end{programme}

\begin{assessment}[Position after V95]
\label{ass:c3v95-position}
The companion audit is recorded; the closing block begins.  No full
stability or global-regularity claim is made.
\end{assessment}

% =============================================================================
\section*{Version V95 (third cycle) append-only record}
\addcontentsline{toc}{section}{Version V95 (third cycle) append-only record}
% =============================================================================

The complete V94 source of the third cycle was retained before this
continuation.  Its SHA--256 digest was
\begin{center}\scriptsize\ttfamily
82e75dca644b070ecf9bcfb73efaf82ba0168c9a30e6fd666dac1bb6fd19ff3f.
\end{center}
V95 performed the compulsory companion audit and recorded the
position.  No full regularity claim is made.

% =============================================================================
\section{V96: the manifest of the third cycle}
\label{sec:c3v96-entry}
% =============================================================================

Items (AC1) and (AC2) of Programme~\ref{prog:c3v96} are carried out.
For each chapter of the third cycle the results that stand are
listed by label and name, in the order of the file; nothing is
re-proved or re-stated, and every entry is a theorem, proposition
or corollary of the cycle (lemmas, firewalls and definitions are
reached through the entries that use them).  The audit decisions
(V5, V10, \dots, V95), the ``What VNN proves'' conclusions and the
version records are not listed, being records.  The legacy
summaries of V1 (Theorems~\ref{thm:c3v1-ledgers} to
\ref{thm:c3v1-spikes}) are the imports from the two legacy cycles and
are listed first.  Each chapter carries a status line.  No global
regularity claim is made.  The next compulsory companion audit is
V100.

% =============================================================================
\section{Imports from the legacy cycles (V1)}
\label{sec:c3v96-legacy}
% =============================================================================

\begin{center}\small
\begin{tabular}{@{}p{0.3\textwidth}p{0.64\textwidth}@{}}
\toprule
\textbf{Label} & \textbf{Name}\\
\midrule
\texttt{thm:c3v1-ledgers} & Ledgers and the two-regime classification, first cycle V32--V40\\
\texttt{thm:c3v1-ancient} & The ancient limit and the type-I gate, first cycle V41--V50, V56--V65, V71--V75, V93\\
\texttt{thm:c3v1-rate-free} & Rate-free results and the three gates, first cycle V51--V55, V66--V70\\
\texttt{thm:c3v1-shapes} & Shapes of singular points, first cycle V76--V85\\
\texttt{thm:c3v1-gaussian} & The Gaussian ledger, first cycle V86--V99\\
\texttt{thm:c3v1-profile} & Profile framework and statistics, second cycle V1--V25\\
\texttt{thm:c3v1-constants} & Constant-chasing gate, second cycle V36--V41\\
\texttt{thm:c3v1-windows} & The intermediate class at its windows, second cycle V42--V55\\
\texttt{thm:c3v1-trap} & The cell trap and the Morrey floors, second cycle V57, V59, V71, V72\\
\texttt{thm:c3v1-propagation} & Finite configurations and propagation, second cycle V58--V69\\
\texttt{thm:c3v1-typeI} & Type-I structure and singular points, second cycle V63, V66, V67\\
\texttt{thm:c3v1-spikes} & Spikes, second cycle V72--V78\\
\bottomrule
\end{tabular}
\end{center}

% =============================================================================
\section{The chapters}
\label{sec:c3v96-chapters}
% =============================================================================

\subsection*{Cluster energy and the band (V2--V10)}
\emph{Status:} conditional (windows and class constants; consolidated in V5, V8, V10).

\begin{center}\small
\begin{tabular}{@{}p{0.05\textwidth}p{0.33\textwidth}p{0.56\textwidth}@{}}
\toprule
\textbf{V} & \textbf{Label} & \textbf{Name}\\
\midrule
2 & \texttt{thm:c3v2-monotone} & An isolated cluster loses energy\\
2 & \texttt{cor:c3v2-two} & Loss per window and the two-cluster case\\
2 & \texttt{prop:c3v2-lifetime} & Conditional lifetime\\
3 & \texttt{thm:c3v3-group} & Monotonicity for a group\\
3 & \texttt{prop:c3v3-accounting} & What the accounting gives and does not give\\
4 & \texttt{thm:c3v4-nucleation} & Nucleation takes at least two windows\\
4 & \texttt{cor:c3v4-two-windows} & Activity in the exterior over two windows\\
5 & \texttt{thm:c3v5-consolidated} & Cluster energy on a window, V2--V4\\
6 & \texttt{cor:c3v6-band} & Band cells over a window\\
7 & \texttt{prop:c3v7-seeds} & Seeds per window and their union\\
7 & \texttt{thm:c3v7-alternative} & Isolation alternative over \(K\) windows\\
7 & \texttt{prop:c3v7-approach} & Who approaches the layer\\
8 & \texttt{prop:c3v8-gaussian} & Gaussian ledger of band and energetic centres\\
8 & \texttt{prop:c3v8-rise} & Rise of a band centre in log-scale\\
8 & \texttt{thm:c3v8-consolidated} & The band chapter, V6--V8\\
10 & \texttt{thm:c3v10-consolidated} & V2--V9\\
\bottomrule
\end{tabular}
\end{center}

\subsection*{Localization at homogeneous points (V11--V14)}
\emph{Status:} explicit at homogeneous points; no sign for the localized push (V13).

\begin{center}\small
\begin{tabular}{@{}p{0.05\textwidth}p{0.33\textwidth}p{0.56\textwidth}@{}}
\toprule
\textbf{V} & \textbf{Label} & \textbf{Name}\\
\midrule
11 & \texttt{thm:c3v11-far} & The far field carries no anti-alignment\\
11 & \texttt{thm:c3v11-localized} & The configuration carries the anti-alignment\\
12 & \texttt{thm:c3v12-lower} & The configuration carries the defect\\
12 & \texttt{prop:c3v12-tensor} & Tensor identity on the configuration neighbourhood\\
13 & \texttt{prop:c3v13-decomposition} & Decomposition of the configuration part\\
13 & \texttt{thm:c3v13-no-sign} & No sign for the localized push\\
14 & \texttt{thm:c3v14-consolidated} & Localization at homogeneous points, V11--V13\\
\bottomrule
\end{tabular}
\end{center}

\subsection*{The anchor (V16--V19)}
\emph{Status:} explicit structure; the cross term lives within the Gaussian radius.

\begin{center}\small
\begin{tabular}{@{}p{0.05\textwidth}p{0.33\textwidth}p{0.56\textwidth}@{}}
\toprule
\textbf{V} & \textbf{Label} & \textbf{Name}\\
\midrule
16 & \texttt{prop:c3v16-anchor} & The anchor group\\
16 & \texttt{prop:c3v16-not-pinned} & The other groups are not pinned\\
16 & \texttt{thm:c3v16-localized} & The cross term lives on the groups within the Gaussian radius\\
17 & \texttt{thm:c3v17-cases} & Two-case form of the localized gate\\
17 & \texttt{prop:c3v17-count} & Count at the anchor's level\\
18 & \texttt{prop:c3v18-mean} & Mean and momentum on the anchor\\
18 & \texttt{prop:c3v18-dipole} & Dipole matrix of the anchor\\
19 & \texttt{thm:c3v19-consolidated} & The anchor chapter, V16--V18\\
\bottomrule
\end{tabular}
\end{center}

\subsection*{The explicit variance floor (V21--V24)}
\emph{Status:} explicit (\(v_{\rm exp}\)); the push budget closed for practical purposes.

\begin{center}\small
\begin{tabular}{@{}p{0.05\textwidth}p{0.33\textwidth}p{0.56\textwidth}@{}}
\toprule
\textbf{V} & \textbf{Label} & \textbf{Name}\\
\midrule
21 & \texttt{prop:c3v21-criteria} & One-level and two-level criteria\\
21 & \texttt{prop:c3v21-bounded} & Pointwise bound on the dominant level and on the spectral tail\\
21 & \texttt{cor:c3v21-nearly-polynomial} & Small variance means nearly polynomial of bounded level\\
22 & \texttt{prop:c3v22-cutoff-level-k} & Cut-off eigenfunctions of every level\\
23 & \texttt{thm:c3v23-floor} & Explicit variance floor at homogeneous points\\
23 & \texttt{cor:c3v23-closed-form} & The closed-form case\\
23 & \texttt{prop:c3v23-descent} & Pointwise descent of the Rayleigh quotient\\
24 & \texttt{thm:c3v24-mean-push} & Mean geometric push with the explicit floor\\
24 & \texttt{thm:c3v24-consolidated} & The explicit-floor chapter, V21--V23\\
\bottomrule
\end{tabular}
\end{center}

\subsection*{The variance chapter (V26--V29)}
\emph{Status:} explicit window floor; negative record N-C3-2 (V26).

\begin{center}\small
\begin{tabular}{@{}p{0.05\textwidth}p{0.33\textwidth}p{0.56\textwidth}@{}}
\toprule
\textbf{V} & \textbf{Label} & \textbf{Name}\\
\midrule
26 & \texttt{thm:c3v26-conditional} & Conditional mean floor\\
26 & \texttt{prop:c3v26-necessary} & Necessary conditions from the \(k=0\) term\\
26 & \texttt{thm:c3v26-incompatible} & The condition is incompatible with the class\\
27 & \texttt{prop:c3v27-identity} & Exact identity for the mean pressure gradient\\
27 & \texttt{prop:c3v27-core} & The nonlinear mean on a nearly constant core\\
27 & \texttt{cor:c3v27-decay} & Decay of the Gaussian momentum on a nearly constant core\\
28 & \texttt{thm:c3v28-window} & Window floor on the variance\\
28 & \texttt{cor:c3v28-mean} & Mean floor from the window floor\\
29 & \texttt{thm:c3v29-chapter} & The variance chapter, V21--V28\\
\bottomrule
\end{tabular}
\end{center}

\subsection*{The defect block (V31--V34)}
\emph{Status:} explicit constraints (NRS identity, mean-square identity, extremal Dirichlet form).

\begin{center}\small
\begin{tabular}{@{}p{0.05\textwidth}p{0.33\textwidth}p{0.56\textwidth}@{}}
\toprule
\textbf{V} & \textbf{Label} & \textbf{Name}\\
\midrule
31 & \texttt{prop:c3v31-gradient} & Gradient form of the defect equation\\
31 & \texttt{thm:c3v31-NRS} & The NRS identity with a defect\\
31 & \texttt{prop:c3v31-gaussian} & The Gaussian weight\\
31 & \texttt{prop:c3v31-constant} & The constant weight and the maximum principle\\
32 & \texttt{thm:c3v32-mean-square} & Mean-square form of the second identity\\
32 & \texttt{prop:c3v32-lower} & Explicit lower side\\
32 & \texttt{prop:c3v32-upper} & Explicit upper side\\
32 & \texttt{thm:c3v32-constraint} & Explicit mean-square constraint at a homogeneous point\\
33 & \texttt{prop:c3v33-pointwise} & Pointwise form\\
33 & \texttt{thm:c3v33-window} & Window form of the mean-square identity\\
33 & \texttt{cor:c3v33-defect-window} & Explicit window bound on the defect of every element\\
34 & \texttt{thm:c3v34-block} & The defect block, V31--V33\\
34 & \texttt{thm:c3v34-dirichlet} & Explicit bounds on the extremal Dirichlet form\\
\bottomrule
\end{tabular}
\end{center}

\subsection*{The extremal ledger (V36--V39)}
\emph{Status:} explicit rows at extremal configurations; gate G1 as an angle.

\begin{center}\small
\begin{tabular}{@{}p{0.05\textwidth}p{0.33\textwidth}p{0.56\textwidth}@{}}
\toprule
\textbf{V} & \textbf{Label} & \textbf{Name}\\
\midrule
36 & \texttt{thm:c3v36-energy} & Extremal Gaussian energy\\
36 & \texttt{thm:c3v36-rayleigh} & Extremal Rayleigh quotient\\
37 & \texttt{thm:c3v37-spectrum} & Extremal spectrum\\
37 & \texttt{prop:c3v37-conditional} & Conditional pointwise floor for a positive dominant level\\
37 & \texttt{thm:c3v37-momentum} & Extremal Gaussian momentum\\
38 & \texttt{thm:c3v38-defect} & Extremal defect\\
38 & \texttt{prop:c3v38-advective} & The advective parts of the linearized form\\
38 & \texttt{prop:c3v38-cross} & Extremal cross term\\
39 & \texttt{thm:c3v39-ledger} & The extremal ledger, V34--V38\\
39 & \texttt{prop:c3v39-combining} & Combining rows\\
39 & \texttt{thm:c3v39-angle} & Gate G1 is a statement about an angle\\
\bottomrule
\end{tabular}
\end{center}

\subsection*{The angle chapter (V41--V44)}
\emph{Status:} structural (Lamb--Bernoulli form, bilinear and flux-derivative forms).

\begin{center}\small
\begin{tabular}{@{}p{0.05\textwidth}p{0.33\textwidth}p{0.56\textwidth}@{}}
\toprule
\textbf{V} & \textbf{Label} & \textbf{Name}\\
\midrule
41 & \texttt{prop:c3v41-lamb} & The Lamb piece\\
41 & \texttt{prop:c3v41-bernoulli} & The Bernoulli piece\\
41 & \texttt{thm:c3v41-MNY} & Explicit bound on the cross term\\
42 & \texttt{prop:c3v42-pointwise} & Pointwise structure of the Lamb piece\\
42 & \texttt{prop:c3v42-core} & On a nearly constant core the cross term is pressure work\\
42 & \texttt{cor:c3v42-extremal-core} & Extremal-Dirichlet nearly constant core\\
43 & \texttt{thm:c3v43-extremal} & Extremal Dirichlet form, Lamb--Bernoulli form\\
43 & \texttt{prop:c3v43-beltrami} & The Beltrami-type reduction\\
43 & \texttt{prop:c3v43-tangential} & The tangential reduction\\
44 & \texttt{thm:c3v44-chapter} & The angle chapter, V39 and V41--V43\\
44 & \texttt{prop:c3v44-bilinear} & Bilinear form of the pressure work\\
44 & \texttt{prop:c3v44-flux} & The Bernoulli work is the derivative of the flux, up to a pairing\\
\bottomrule
\end{tabular}
\end{center}

\subsection*{The correlation block (V46--V49)}
\emph{Status:} explicit bounds at supremum level (energy--Hessian form, pressure firewall).

\begin{center}\small
\begin{tabular}{@{}p{0.05\textwidth}p{0.33\textwidth}p{0.56\textwidth}@{}}
\toprule
\textbf{V} & \textbf{Label} & \textbf{Name}\\
\midrule
46 & \texttt{prop:c3v46-far} & The far pressure is flat on the configuration\\
46 & \texttt{thm:c3v46-pairing} & The pressure work on the configuration as an energy--Hessian pairing\\
46 & \texttt{prop:c3v46-multipole} & Point-mass form of a cell's Hessian\\
46 & \texttt{prop:c3v46-self} & Self-term\\
47 & \texttt{thm:c3v47-pairing} & Explicit bound on the pairing of the linearized pressure with the radial velocity\\
47 & \texttt{cor:c3v47-extremal-flux} & The Bernoulli work at the extremal flux\\
47 & \texttt{thm:c3v47-revised} & Every Gaussian pressure pairing is explicit at supremum level\\
48 & \texttt{prop:c3v48-back} & Item (AC1)\\
48 & \texttt{thm:c3v48-rows} & Rows (6), (7) and the flux row, explicit\\
49 & \texttt{thm:c3v49-block} & The correlation block, V44 and V46--V48\\
\bottomrule
\end{tabular}
\end{center}

\subsection*{Midpoint internal audit (V51--V55)}
\emph{Status:} corrections in place (V53, V54) and citations confirmed.

\begin{center}\small
\begin{tabular}{@{}p{0.05\textwidth}p{0.33\textwidth}p{0.56\textwidth}@{}}
\toprule
\textbf{V} & \textbf{Label} & \textbf{Name}\\
\midrule
51 & \texttt{thm:c3v51-floor} & Audit of V21--V24\\
51 & \texttt{thm:c3v51-variance} & Audit of V26--V28\\
51 & \texttt{prop:c3v51-citations} & Citations of V21--V28 confirmed\\
52 & \texttt{thm:c3v52-defect} & Audit of V31--V34\\
52 & \texttt{prop:c3v52-citations} & Citations of V31--V34 confirmed\\
53 & \texttt{prop:c3v53-correction} & Corrected \(1/k\) form of the extremal spectrum\\
53 & \texttt{thm:c3v53-ledger} & Audit of V36--V38\\
53 & \texttt{thm:c3v53-angle} & Audit of V41--V43\\
53 & \texttt{prop:c3v53-citations} & Citations of V36--V43 confirmed\\
54 & \texttt{prop:c3v54-correction} & Corrected definition of \(C_P\)\\
54 & \texttt{thm:c3v54-correlation} & Audit of V44 and V46--V48\\
54 & \texttt{prop:c3v54-lists} & Superseded wordings and in-place corrections of the third cycle, V1--V54\\
55 & \texttt{thm:c3v55-audit} & Result of the midpoint internal audit, V51--V54\\
\bottomrule
\end{tabular}
\end{center}

\subsection*{The anchor-gate block (V56--V59)}
\emph{Status:} explicit gate on the anchor-alone event; no sign.

\begin{center}\small
\begin{tabular}{@{}p{0.05\textwidth}p{0.33\textwidth}p{0.56\textwidth}@{}}
\toprule
\textbf{V} & \textbf{Label} & \textbf{Name}\\
\midrule
56 & \texttt{prop:c3v56-total} & Total radial-defect mass of the anchor\\
56 & \texttt{thm:c3v56-dipole} & The dipole moment of the radial defect is the Gaussian mean of the defect\\
57 & \texttt{thm:c3v57-one-cell} & One cell\\
57 & \texttt{thm:c3v57-two-cell} & Two cells of opposite mass\\
58 & \texttt{thm:c3v58-group} & Pressure work on the anchor-alone event\\
58 & \texttt{thm:c3v58-gate} & The gate on \(\mathcal E_0\)\\
59 & \texttt{thm:c3v59-block} & The anchor-gate block, V56--V58\\
59 & \texttt{prop:c3v59-fourth} & Fourth bound on the defect rate\\
\bottomrule
\end{tabular}
\end{center}

\subsection*{The nearly-constant-anchor block (V61--V64)}
\emph{Status:} explicit exclusion of the nearly constant anchor as carrier of the anti-alignment.

\begin{center}\small
\begin{tabular}{@{}p{0.05\textwidth}p{0.33\textwidth}p{0.56\textwidth}@{}}
\toprule
\textbf{V} & \textbf{Label} & \textbf{Name}\\
\midrule
61 & \texttt{thm:c3v61-hessian} & The cut-off integral of the Hessian\\
61 & \texttt{thm:c3v61-work} & Local pressure work on the anchor, exact form\\
62 & \texttt{thm:c3v62-cross} & The cross term on a nearly constant anchor\\
62 & \texttt{thm:c3v62-extremal} & The defect at an extremal-Dirichlet nearly constant anchor is small\\
62 & \texttt{thm:c3v62-measure} & The blow-up measure is not concentrated on nearly constant anchors\\
63 & \texttt{prop:c3v63-multi} & The multi-group neighbourhood\\
63 & \texttt{thm:c3v63-multi} & The nearly constant multi-group configuration\\
64 & \texttt{thm:c3v64-block} & The nearly-constant-anchor block, V61--V63\\
\bottomrule
\end{tabular}
\end{center}

\subsection*{The spread-anchor block (V66--V69)}
\emph{Status:} explicit; first constants-only constraint (V67); test flows gate-neutral.

\begin{center}\small
\begin{tabular}{@{}p{0.05\textwidth}p{0.33\textwidth}p{0.56\textwidth}@{}}
\toprule
\textbf{V} & \textbf{Label} & \textbf{Name}\\
\midrule
66 & \texttt{thm:c3v66-floor} & Mass off the constant level and the dissipation floor\\
66 & \texttt{prop:c3v66-rows} & Extremal rows on spread configurations\\
67 & \texttt{thm:c3v67-extremal} & Extremal Gaussian enstrophy, and the constraint\\
68 & \texttt{thm:c3v68-gate} & The gate on the spread anchor, explicit form\\
68 & \texttt{prop:c3v68-tests} & The columnar swirl and the plane wave are gate-neutral\\
69 & \texttt{thm:c3v69-block} & The spread-anchor block, V66--V68\\
\bottomrule
\end{tabular}
\end{center}

\subsection*{The vorticity ledger (V71--V74)}
\emph{Status:} constraints of \(\varepsilon\)-regularity type; negative record (V73); the loop (V74).

\begin{center}\small
\begin{tabular}{@{}p{0.05\textwidth}p{0.33\textwidth}p{0.56\textwidth}@{}}
\toprule
\textbf{V} & \textbf{Label} & \textbf{Name}\\
\midrule
71 & \texttt{thm:c3v71-identity} & Weighted Gaussian enstrophy identity\\
71 & \texttt{thm:c3v71-extremal} & Extremal weighted enstrophy\\
72 & \texttt{thm:c3v72-identity} & Gaussian palinstrophy identity\\
72 & \texttt{thm:c3v72-extremal} & Extremal palinstrophy: a ratio bound\\
72 & \texttt{prop:c3v72-poincare} & The Gaussian Poincar\'e link\\
73 & \texttt{prop:c3v73-gradient-floor} & Gradient floor, without hypothesis\\
73 & \texttt{thm:c3v73-velocity-floor} & Velocity floor under (E5-small)\\
73 & \texttt{prop:c3v73-thresholds} & Explicit applicability thresholds\\
73 & \texttt{thm:c3v73-window} & The common window\\
73 & \texttt{prop:c3v73-epsreg} & Comparison with \(\varepsilon_{\rm reg}\)\\
73 & \texttt{prop:c3v73-chasing} & Comparison with the constant-chasing thresholds\\
73 & \texttt{thm:c3v73-table} & Item (AC2): the vorticity ledger at extremal configurations\\
74 & \texttt{thm:c3v74-block} & The vorticity-ledger block, V71--V73\\
74 & \texttt{thm:c3v74-identity} & Mean-vorticity identity\\
74 & \texttt{thm:c3v74-extremal} & Extremal mean vorticity\\
74 & \texttt{thm:c3v74-loop} & The mean-vorticity loop: a second constants-only constraint\\
\bottomrule
\end{tabular}
\end{center}

\subsection*{The loop block (V76--V79)}
\emph{Status:} constraints and floors, capped; negative records (V77, V78).

\begin{center}\small
\begin{tabular}{@{}p{0.05\textwidth}p{0.33\textwidth}p{0.56\textwidth}@{}}
\toprule
\textbf{V} & \textbf{Label} & \textbf{Name}\\
\midrule
76 & \texttt{cor:c3v76-v74} & V74 recovered\\
76 & \texttt{thm:c3v76-floor} & Loop floor under (E5-small)\\
76 & \texttt{thm:c3v76-comparison} & Loop floor against enstrophy floor\\
76 & \texttt{thm:c3v76-table} & The three floors\\
77 & \texttt{thm:c3v77-loop} & Enstrophy--mean loop: a third constants-only constraint\\
77 & \texttt{prop:c3v77-comparison} & Comparison with V67 and V74\\
77 & \texttt{thm:c3v77-floor} & Floor of the enstrophy--mean loop\\
77 & \texttt{thm:c3v77-four} & The four-observable loop\\
77 & \texttt{prop:c3v77-four-floor} & The four-observable floor never exceeds \(\kappa_1\)\\
78 & \texttt{thm:c3v78-identity} & Dipole identity\\
78 & \texttt{thm:c3v78-extremal} & Extremal dipole\\
78 & \texttt{prop:c3v78-poincare} & Second-order Gaussian Poincar\'e inequality\\
78 & \texttt{thm:c3v78-cap} & The loop through the dipole, and its cap\\
79 & \texttt{thm:c3v79-block} & The loop block, V74 and V76--V78\\
\bottomrule
\end{tabular}
\end{center}

\subsection*{The explicit-constant block (V81--V84)}
\emph{Status:} certified numbers (Oseen constant, Liouville and enstrophy floors).

\begin{center}\small
\begin{tabular}{@{}p{0.05\textwidth}p{0.33\textwidth}p{0.56\textwidth}@{}}
\toprule
\textbf{V} & \textbf{Label} & \textbf{Name}\\
\midrule
81 & \texttt{thm:c3v81-duhamel} & Duhamel formula on finite windows\\
81 & \texttt{thm:c3v81-ancient} & Ancient form\\
81 & \texttt{thm:c3v81-epsilon} & Explicit \(\varepsilon\)-regularity floor\\
81 & \texttt{thm:c3v81-gradient} & Explicit gradient bound\\
82 & \texttt{prop:c3v82-opnorm} & Pointwise operator norm\\
82 & \texttt{thm:c3v82-certified} & Certified Oseen constant\\
82 & \texttt{prop:c3v82-computed} & Computed Oseen constant\\
82 & \texttt{cor:c3v82-values} & Item (AC2): the floors with values\\
83 & \texttt{prop:c3v83-gradient-liouville} & Gradient Liouville floor\\
83 & \texttt{thm:c3v83-floor} & Explicit enstrophy floor\\
83 & \texttt{thm:c3v83-table} & Explicit floors on the velocity constant\\
84 & \texttt{thm:c3v84-gamma2} & Explicit second-derivative bound\\
84 & \texttt{thm:c3v84-void} & The palinstrophy row is void below \(0.4569\)\\
84 & \texttt{prop:c3v84-concentration} & The concentration row at the floors\\
84 & \texttt{thm:c3v84-table} & Item (AC2): the extended table\\
\bottomrule
\end{tabular}
\end{center}

\subsection*{The traceless-reduction block (V86--V89)}
\emph{Status:} certified quadrature; floor \(0.1565\).

\begin{center}\small
\begin{tabular}{@{}p{0.05\textwidth}p{0.33\textwidth}p{0.56\textwidth}@{}}
\toprule
\textbf{V} & \textbf{Label} & \textbf{Name}\\
\midrule
86 & \texttt{thm:c3v86-absorptions} & Reduced absorptions\\
86 & \texttt{cor:c3v86-values} & Item (AC2): reduced Liouville floors\\
87 & \texttt{prop:c3v87-opnorm} & Restricted pointwise norm\\
87 & \texttt{prop:c3v87-computed} & Item (AC2): computed restricted constant\\
88 & \texttt{thm:c3v88-floor} & Reduced enstrophy floor\\
88 & \texttt{prop:c3v88-derivatives} & Derivative bounds and the concentration threshold at the reduced floors\\
88 & \texttt{thm:c3v88-table} & Item (AC2): explicit floors after the reduction\\
89 & \texttt{thm:c3v89-certified} & Certified Oseen constants\\
89 & \texttt{thm:c3v89-table} & Item (AC2): certified explicit floors\\
89 & \texttt{thm:c3v89-block} & The traceless-reduction block, V86--V89\\
\bottomrule
\end{tabular}
\end{center}

\subsection*{The strain--transport block (V91--V94)}
\emph{Status:} certified floor \(0.1815\); constants sheet.

\begin{center}\small
\begin{tabular}{@{}p{0.05\textwidth}p{0.33\textwidth}p{0.56\textwidth}@{}}
\toprule
\textbf{V} & \textbf{Label} & \textbf{Name}\\
\midrule
91 & \texttt{thm:c3v91-enstrophy} & Strain form of the enstrophy constraint\\
91 & \texttt{thm:c3v91-rows} & Strain form of the weighted-enstrophy and palinstrophy rows\\
91 & \texttt{prop:c3v91-consequences} & Item (AC2): consequences in strain form\\
92 & \texttt{thm:c3v92-floor} & Certified strain-form floor\\
92 & \texttt{prop:c3v92-shares} & Item (AC2): shares at the floor\\
92 & \texttt{thm:c3v92-table} & The table after the strain form\\
93 & \texttt{thm:c3v93-rows} & The rows in strain--transport form\\
93 & \texttt{prop:c3v93-consequences} & Item (AC2): consequences\\
94 & \texttt{thm:c3v94-floor} & Certified floor in strain--transport form\\
94 & \texttt{thm:c3v94-table} & The final table of the explicit chapter\\
94 & \texttt{thm:c3v94-sheet} & Item (AC2): constants sheet, V81--V94\\
\bottomrule
\end{tabular}
\end{center}

% =============================================================================
\section{Integrated V96 conclusion and next acquisition order}
\label{sec:c3v96-conclusion}
% =============================================================================

\begin{theorem}[What V96 establishes]
\label{thm:c3v96-conclusion}
Relative to the two legacy cycles and V1--V95: the manifest above,
196 entries in eighteen chapters, each an existing labelled
result of the cycle.  The full Navier--Stokes stability/global-regularity
theorem remains unproved.
\end{theorem}

\begin{proof}
The cited labels.
\end{proof}

\begin{programme}[V97 acquisition order]
\label{prog:c3v97}
\begin{enumerate}[label=\textup{(AC\arabic*)},leftmargin=4.8em]
\item The negative records of the third cycle collected: N-C3-1
      (V22), N-C3-2 (V26), N-C3-3 (V31), the firewalls of V32, V47, V63,
      V66, V67, V73, V74, V77, V78, and the voidness statements of V84
      and V93; for each, the statement, the reason, and what would
      lift it.
\item The status of the gates (G1 and the angle form) after the
      cycle.
\item The next compulsory companion audit is V100.
\end{enumerate}
\end{programme}

\begin{assessment}[Position after V96]
\label{ass:c3v96-position}
The manifest is written.  No full stability or global-regularity
claim is made.
\end{assessment}

% =============================================================================
\section*{Version V96 (third cycle) append-only record}
\addcontentsline{toc}{section}{Version V96 (third cycle) append-only record}
% =============================================================================

The complete V95 source of the third cycle was retained before this
continuation.  Its SHA--256 digest was
\begin{center}\scriptsize\ttfamily
646eac0396bed8d5ee4fc5d23718df1aaf3aecdd6ee68faf03a197f293ff9688.
\end{center}
V96 wrote the manifest of the third cycle.  No full regularity claim
is made.

% =============================================================================
\section{V97: the negative records of the third cycle}
\label{sec:c3v97-entry}
% =============================================================================

Items (AC1) and (AC2) of Programme~\ref{prog:c3v97} are carried out.
The cycle's negative records are collected: the three numbered
records N-C3-1 to N-C3-3, the standing firewalls, and the voidness
statements, each with its statement, its reason and what would lift
it, by reference to the labelled firewall.  The status of the gates
after the cycle is recorded.  No global regularity claim is made.
The next compulsory companion audit is V100.

% =============================================================================
\section{The numbered records}
\label{sec:c3v97-numbered}
% =============================================================================

\begin{theorem}[Numbered negative records]
\label{thm:c3v97-numbered}
\medskip
\begin{center}\small
\begin{tabular}{@{}p{0.11\textwidth}p{0.27\textwidth}p{0.28\textwidth}p{0.28\textwidth}@{}}
\toprule
\textbf{Record} & \textbf{Statement} & \textbf{Reason} & \textbf{What would lift it}\\
\midrule
N-C3-1 (Firewall~\ref{fw:c3v22-no-floor-from-concentration}) &
 concentration of the profile in a ball gives no floor on the variance &
 the cut-off eigenfunctions of every level are nearly polynomial in the ball; the growth lemma bounds the tail only exponentially &
 a pointwise use of the equation on the ball (Firewall~\ref{fw:c3v23-no-equation})\\[3pt]
N-C3-2 (Firewall~\ref{fw:c3v26-negative}) &
 the mean spectrum gives no mean floor on the variance &
 the conditional floor's condition \(B_0<c_\Phi\) is incompatible with the class for every \(\kappa\) (Theorem~\ref{thm:c3v26-incompatible}) &
 nothing: the route is closed\\[3pt]
N-C3-3 (Firewall~\ref{fw:c3v31-negative}) &
 no quantitative NRS (defect-rate) inequality holds in the class &
 the defect rate \(\underline\delta\) comes from (E9) by compactness and cannot be made explicit in the class; the NRS identity with defect has a signed but non-explicit source &
 an explicit Liouville constant for (E9) in the class, or a different positivity\\
\bottomrule
\end{tabular}
\end{center}
\medskip
\end{theorem}

\begin{proof}
The cited firewalls.
\end{proof}

% =============================================================================
\section{The standing firewalls}
\label{sec:c3v97-firewalls}
% =============================================================================

\begin{theorem}[Standing firewalls of the cycle]
\label{thm:c3v97-firewalls}
\medskip
\begin{center}\small
\begin{tabular}{@{}p{0.2\textwidth}p{0.36\textwidth}p{0.38\textwidth}@{}}
\toprule
\textbf{Firewall} & \textbf{Statement} & \textbf{Reason; what would lift it}\\
\midrule
\ref{fw:c3v2-multiwindow} & the multi-window hypothesis on cluster energy is not supplied by the class & one window's loss is explicit, its repetition is not; a rate on the window sequence would\\[2pt]
\ref{fw:c3v3-AC1}, \ref{fw:c3v7-AC}, \ref{fw:c3v8-AC1} & the accounting of groups and seeds gives counts, not a contradiction & the counts are consistent with the class; a bound on the number of energetic centres per scale below the count would\\[2pt]
\ref{fw:c3v6-band} & the band is a large-scale object and is not localized & its energy sits at the Gaussian radius; a localization at the anchor scale would\\[2pt]
\ref{fw:c3v12-AC2}, \ref{fw:c3v18-neutral} & the cross term of the gate has no signed row on the configuration & tensor identity and dipole matrix have no sign; a sign from the equation would\\[2pt]
\ref{fw:c3v21-mean-not-pointwise}, \ref{fw:c3v21-high-empty} & the mean spectrum is not a pointwise bound; the high case of the one-level criterion is empty & levels above \(K_r\) are excluded by the Rayleigh bound; nothing to lift\\[2pt]
\ref{fw:c3v24-threshold} & the push threshold is explicit and out of reach with \(v_{\rm exp}\) & \(v_{\rm exp}\) is exponentially small; a polynomial floor would\\[2pt]
\ref{fw:c3v28-lipschitz} & \(\Lambda_{\rm Var}\) has no closed form & derivative bounds of order three and four not written; V84's method extends to them\\[2pt]
\ref{fw:c3v32-no-exclusion} & the mean-square constraint excludes nothing & pinching and defect rate have no explicit lower bound (N-C3-3)\\[2pt]
\ref{fw:c3v33-pressure} & no Gaussian form for the pressure part of the defect & the pressure is nonlocal; the energy--Hessian form (V46) is the substitute\\[2pt]
\ref{fw:c3v37-obstruction}, \ref{fw:c3v38-pressure}, \ref{fw:c3v38-cross} & the linearized pressure blocks the explicit evaluation of the spectrum, momentum and cross rows & pointwise supremum bounds replace it (V47) at the cost of explicitness in the constants\\[2pt]
\ref{fw:c3v39-no-norms} & no norm comparison can close G1 & G1 is an angle statement; only a sign would\\[2pt]
\ref{fw:c3v47-no-bypass}, \ref{fw:c3v47-revised} & the double integration by parts of V46 is not justified; the pressure firewall in revised form & boundary terms of order one; a decay of \(Y\) would\\[2pt]
\ref{fw:c3v63-trapped} & the multi-group extension holds only inside trapped sub-windows & the quiet annulus is not guaranteed; a trap at every time would\\[2pt]
\ref{fw:c3v66-enstrophy}, \ref{fw:c3v67-no-floor} & no floor on the Gaussian enstrophy from the dissipation; the Lamb pairing bounded above only & the remainder \(R\) can absorb; a sign of the pairing would\\[2pt]
\ref{fw:c3v73-not-implied}, \ref{fw:c3v74-loop}, \ref{fw:c3v77-negative}, \ref{fw:c3v78-negative} & the vorticity constraints are of \(\varepsilon\)-regularity type; the loops cap at \(\max\{\kappa_1,1/\sqrt3\}\) and \(3/\sqrt{80}\) & the Hardy conversion returns the low Hermite levels with \(c\kappa^2\), \(c\ge3\); a closing inequality with a small multiple would\\
\bottomrule
\end{tabular}
\end{center}
\medskip
\end{theorem}

\begin{proof}
The cited firewalls.
\end{proof}

\begin{proposition}[Voidness statements]
\label{prop:c3v97-void}
The palinstrophy row of V72, the loop constraints of V74 and V77
and the loop floors of V76--V77 are void at every homogeneous
type-I point with \(\kappa<1\) (Proposition~\ref{prop:c3v93-consequences}(3);
earlier \(0.4569\), Theorem~\ref{thm:c3v84-void}); they can bear only on
points with velocity constant at least one, where the explicit
gradient bound is not absorbing.  They are recorded as results of
the cycle whose range of application is empty below \(\kappa=1\).
\end{proposition}

\begin{proof}
The cited statements.
\end{proof}

% =============================================================================
\section{Item (AC2): the gates after the cycle}
\label{sec:c3v97-gates}
% =============================================================================

\begin{assessment}[Status of the gates]
\label{ass:c3v97-gates}
Gate G1 (Theorem~\ref{thm:c3v1-ancient}: every nonzero ancient limit at a
homogeneous type-I point is excluded) is unproved.  The cycle
reduced it to an angle statement (Theorem~\ref{thm:c3v39-angle}),
localized its cross term on the anchor (V16, V58), excluded the
nearly constant anchor as its carrier (V62), and showed that on a
spread anchor it is a question of two orientations of explicit
order-one quantities that the class does not fix (V68).  What the
cycle proved instead are constraints that every homogeneous
type-I point must satisfy: the explicit variance floor (V23), the
mean-square constraint (V32), the extremal rows (V36--V38, V43,
V47--V48), the enstrophy constraint in its final strain--transport
form (V93), and the certified floor \(\kappa\ge0.1815\) (V94).  None is an
exclusion; together they describe the singular point, if it exists,
as one with a rescaled velocity bound at least \(0.18\nu^{1/2}\), strain
constant of order one, Gaussian vorticity moment bounded by twelve
times the enstrophy scale, and an anti-alignment carried by a
spread anchor within the Gaussian radius.  The other two gates of
the first cycle (Theorem~\ref{thm:c3v1-rate-free}) were not touched
by the cycle.
\end{assessment}

% =============================================================================
\section{Integrated V97 conclusion and next acquisition order}
\label{sec:c3v97-conclusion}
% =============================================================================

\begin{theorem}[What V97 establishes]
\label{thm:c3v97-conclusion}
Relative to the two legacy cycles and V1--V96: the collected negative
records (Theorems~\ref{thm:c3v97-numbered}, \ref{thm:c3v97-firewalls},
Proposition~\ref{prop:c3v97-void}) and the status of the gates.  The
full Navier--Stokes stability/global-regularity theorem remains
unproved.
\end{theorem}

\begin{proof}
The cited records.
\end{proof}

\begin{programme}[V98 acquisition order]
\label{prog:c3v98}
\begin{enumerate}[label=\textup{(AC\arabic*)},leftmargin=4.8em]
\item The constants sheet of the whole cycle: the class constants
      (\(\kappa,C_\infty^{(k)},K_6,C,K_P,c_\Phi,C_\Phi,C_Y,C_N,C_r,K_r\)) with their status
      (hypothesis, compactness, explicit), the cycle's derived
      constants (\(v_{\rm exp}\), \(\rho_c\), \(N\)-C3 constants, the extremal-row
      constants, \(\varepsilon_{\rm reg}\)) and the explicit chapter's sheet by
      reference to Theorem~\ref{thm:c3v94-sheet}.
\item The corrections log: every in-place correction of the cycle
      (V23, V27, V33, V37, V42, V47, V53, V54, V61, V72, V73, V87, V90)
      with the version that recorded it.
\item The next compulsory companion audit is V100.
\end{enumerate}
\end{programme}

\begin{assessment}[Position after V97]
\label{ass:c3v97-position}
The negative records are collected.  No full stability or
global-regularity claim is made.
\end{assessment}

% =============================================================================
\section*{Version V97 (third cycle) append-only record}
\addcontentsline{toc}{section}{Version V97 (third cycle) append-only record}
% =============================================================================

The complete V96 source of the third cycle was retained before this
continuation.  Its SHA--256 digest was
\begin{center}\scriptsize\ttfamily
5cbb34cc7e7b86c557ddab937f8fba3427fb2d0e78d8b31b48ab825c5183c348.
\end{center}
V97 collected the negative records.  No full regularity claim is
made.

% =============================================================================
\section{V98: the constants sheet of the cycle and the corrections log}
\label{sec:c3v98-entry}
% =============================================================================

Items (AC1) and (AC2) of Programme~\ref{prog:c3v98} are carried out:
the constants of the class and of the cycle are tabulated with
their status, and every in-place correction of the cycle is
logged with the version that recorded it.  No global regularity
claim is made.  The next compulsory companion audit is V100.

% =============================================================================
\section{The constants sheet}
\label{sec:c3v98-sheet}
% =============================================================================

\begin{theorem}[Class constants]
\label{thm:c3v98-class}
\medskip
\begin{center}\small
\begin{tabular}{@{}p{0.14\textwidth}p{0.44\textwidth}p{0.36\textwidth}@{}}
\toprule
\textbf{Constant} & \textbf{Meaning} & \textbf{Status}\\
\midrule
\(C_I\), \(C\) & type-I enstrophy rate (I); \(\|\nabla W(s)\|_2^2\le C(-s)^{-1/2}\), \(C=C_I^2\nu^{3/2}\) & hypothesis of the type-I class\\
\(\kappa=C_\infty\nu^{-1/2}\) & rescaled velocity bound \(|W|\le C_\infty(-s)^{-1/2}\) & from (E5) and the class; bounded below by \(0.1815\) (V94)\\
\(C_\infty'\) & \(|\nabla V|_F\le C_\infty'\) on the family & from (E5)/C2-V4; explicit bound \(\gamma_1^{\rm red}(\kappa)\kappa\) (V86); floor \((3/2)^{1/2}(1-\kappa^2/2)\) (V93)\\
\(C_\infty''\), \(C_\infty^{(k)}\) & higher derivative bounds & from (E5)/C2-V4; \(\nu^{1/2}C_\infty''\le\gamma_2^{\rm red}(\kappa)\kappa\) (V84, V86); \(k\ge3\) not written\\
\(K_6\) & \(\|V\|_{L^6}\le K_6\) & Sobolev from \(C\); not numerical in the series\\
\(K_P\) & pressure constant of C2-V37 & class; not numerical\\
\(c_\Phi\le\varphi\le C_\Phi\) & pinching of the Gaussian energy & compactness (C2-V12); no explicit lower value (N-C3-3 analogue)\\
\(C_Y\), \(C_N\) & bounds on \(\|Y\|_G\), \(\|N\|_G\) & class, explicit in the derivative constants (V47)\\
\(C_r=(4\nu C+2C_\Phi)/(4c_\Phi)\), \(K_r=\lfloor2C_r-\frac12\rfloor\) & Rayleigh cap and top level & explicit in the class constants (V21)\\
\(\underline\delta\) & defect rate & compactness from (E9); not explicit (N-C3-3)\\
\(\varepsilon_{\rm reg}\) & CKN \(\varepsilon\)-regularity floor & from (E5); replaced from below by the explicit \(0.1410\) (V81, V86, V89)\\
\(\varepsilon_*\), \(M\), \(N_*\) & window, intermediate-class and separation constants of the legacy cycles & legacy; not used in V21--V94\\
\bottomrule
\end{tabular}
\end{center}
\medskip
\end{theorem}

\begin{theorem}[Derived constants of the cycle]
\label{thm:c3v98-derived}
\medskip
\begin{center}\small
\begin{tabular}{@{}p{0.2\textwidth}p{0.42\textwidth}p{0.32\textwidth}@{}}
\toprule
\textbf{Constant} & \textbf{Definition} & \textbf{Status; where}\\
\midrule
\(v_{\rm exp}\) & explicit variance floor at homogeneous points & explicit, exponentially small in the class constants; V23\\
\(\rho_c^2=8C|G|^{1/2}/(\pi c_\Phi^{1/2})\) & core radius of the nearly constant core & explicit; V27\\
\(v_{\rm win}\), \(\tau_*\), \(\Lambda_{\rm Var}\) & window floor, window length, Lipschitz constant of the variance & first two explicit, third not closed (Firewall~\ref{fw:c3v28-lipschitz}); V28\\
\(\delta_{67}\), \(\delta^S\) & slack of the enstrophy constraint & explicit, nonnegative; V67, V91, V93\\
\(\theta_1,\theta_{\mathcal P}\) and their \(S\), \(ST\) forms & thresholds of the weighted and palinstrophy rows & explicit; V71, V72, V91, V93\\
\(q_{\mathcal P}=2\nu(C_\infty''/\theta_{\mathcal P})^2\) & ratio constant & explicit; V72\\
\(\kappa_1,\kappa_2,\kappa_3,\kappa_4\) & floors under (E5-small) & explicit in \(\gamma_1,\gamma_2\); V73, V76, V77; void below \(\kappa=1\) except \(\kappa_1\) (V93)\\
\(\gamma_1,\gamma_2\) & constants of (E5-small) & given values by V81--V86: \(\gamma_1^{\rm red}(\kappa)\), \(\gamma_2^{\rm red}(\kappa)\)\\
\(c_h,c_{h,2},c_O,c_O^{s0}\) & kernel constants & Theorem~\ref{thm:c3v94-sheet}\\
\(\kappa_1^{\rm expl},\kappa_1^{\rm red},\kappa_1^S,\kappa_1^{ST}\) & certified floors \(0.1129\) (Plancherel), \(0.1565\), \(0.1694\), \(0.1815\) & V83, V89, V92, V94\\
\(1/(\pi c_O)\) forms & Liouville floors \(0.0756\), \(0.1123\), \(0.1410\) & V82, V86, V89\\
\bottomrule
\end{tabular}
\end{center}
\medskip
\end{theorem}

\begin{proof}
The cited results.
\end{proof}

% =============================================================================
\section{The corrections log}
\label{sec:c3v98-log}
% =============================================================================

\begin{proposition}[Item (AC2): in-place corrections of the third cycle]
\label{prop:c3v98-log}
Text of an appended version changed after its first appending,
each recorded in a later version:
\begin{enumerate}[label=\textup{(\arabic*)},leftmargin=3.2em]
\item V33, Corollary~\ref{cor:c3v33-defect-window}: hypothesis \(\kappa<\frac12\)
      removed; recorded in V34.
\item V37, Theorem~\ref{thm:c3v37-spectrum}: constant of the \(1/k\) form
      corrected to \(20[\kappa C_\Phi^{1/2}+\kappa(4\nu C+3C_\Phi)^{1/2}+\frac14K_P]^2\); recorded in V53
      (Proposition~\ref{prop:c3v53-correction}).
\item V47, Lemma~\ref{lem:c3v47-decomposition}(i): factor \(\sqrt3\) in \(C_P\);
      recorded in V54 (Proposition~\ref{prop:c3v54-correction}).
\item V73, Hypothesis (E5-small) and Proposition~\ref{prop:c3v73-thresholds}:
      dimension of \(C_\infty''\) (\(\gamma_2\kappa\nu^{-1/2}\)); recorded in V74
      (Remark~\ref{rem:c3v74-correction}).
\item V87, Lemma~\ref{lem:c3v87-tail}: crossing radius \(2.34\) in place of
      ``about \(2.9\)''; recorded in V88 (Remark~\ref{rem:c3v88-correction}).
\item V90, companion-audit paragraph: SM V50 in place of V49 (the
      SM file moved between the first inspection and the assembly);
      corrected before delivery; recorded in the V91 record.
\end{enumerate}
Corrections made to a draft before its first appending are part of
the appended text and are not logged (Proposition~\ref{prop:c3v54-lists}
lists those of V1--V54; later ones: V56 (constant name \(C_1\), remarked in
V57), V61 (the \(\ell=2\) argument replaced by the flux argument), V63,
V70 (companion file moved during the audit), V72 (sign of \(\int\Omega G\)),
V81 (sign of a symbol sentence)).  Superseded but unedited: the
bounds of Lemma~\ref{lem:c3v67-bounds} and of V71--V72 (superseded by
V91 and V93); Theorem C2-V41-spectrum's constant (legacy, not
edited; Remark~\ref{rem:c3v53-legacy}).
\end{proposition}

\begin{proof}
The records of the cited versions.
\end{proof}

% =============================================================================
\section{Integrated V98 conclusion and next acquisition order}
\label{sec:c3v98-conclusion}
% =============================================================================

\begin{theorem}[What V98 establishes]
\label{thm:c3v98-conclusion}
Relative to the two legacy cycles and V1--V97: the constants sheet
(Theorems~\ref{thm:c3v98-class}, \ref{thm:c3v98-derived}) and the
corrections log (Proposition~\ref{prop:c3v98-log}).  The full
Navier--Stokes stability/global-regularity theorem remains
unproved.
\end{theorem}

\begin{proof}
The cited records.
\end{proof}

\begin{programme}[V99 acquisition order]
\label{prog:c3v99}
\begin{enumerate}[label=\textup{(AC\arabic*)},leftmargin=4.8em]
\item The draft of the fourth cycle's V1: the summary of the major
      results of the three cycles (the legacy summaries of V1 and
      the third cycle's manifest, negative records and certified
      floor), the imports, the rules, and the directions recorded
      for the fourth cycle (Assessment~\ref{ass:c3v94-next}; the
      strain-form reading of the gates; the loop method above \(\kappa=1\)).
\item The list of what the fourth cycle must not redo.
\item The next compulsory companion audit is V100.
\end{enumerate}
\end{programme}

\begin{assessment}[Position after V98]
\label{ass:c3v98-position}
The constants and corrections are logged.  No full stability or
global-regularity claim is made.
\end{assessment}

% =============================================================================
\section*{Version V98 (third cycle) append-only record}
\addcontentsline{toc}{section}{Version V98 (third cycle) append-only record}
% =============================================================================

The complete V97 source of the third cycle was retained before this
continuation.  Its SHA--256 digest was
\begin{center}\scriptsize\ttfamily
c4ab65856497059a61cc5152ccdcec9b90afb9e70334e842fca412ce50fa3a48.
\end{center}
V98 wrote the constants sheet and the corrections log.  No full
regularity claim is made.

% =============================================================================
\section{V99: the draft of the fourth cycle's opening summary and directions}
\label{sec:c3v99-entry}
% =============================================================================

Items (AC1) and (AC2) of Programme~\ref{prog:c3v99} are carried out.
The summary of the third cycle that will open the fourth cycle's V1
is written here in its final form (seven theorems and a remark,
each a digest of a consolidated block, citing the labels of this
file); the directions recorded for the fourth cycle are stated,
with the first acquisition of the fourth cycle formulated; and the
list of what the fourth cycle must not redo is recorded.  No global
regularity claim is made.  The next compulsory companion audit is
V100, which closes the cycle.

% =============================================================================
\section{Summary of the third cycle (audited by its manifest)}
\label{sec:c3v99-cycle3}
% =============================================================================

All statements hold at a homogeneous type-I point of a smooth
periodic solution with \(T_*<\infty\), modulo (E1)--(E9), with the profile
notation of the second cycle; the third cycle's in-place corrections
(its V33, V37, V47, V73, V87, V90) are incorporated, and its
superseded bounds (its V67, V71, V72 transport and stretching
bounds) are replaced by their final forms (its V91, V93).  The
labels of the third cycle's file are named as ``Theorem C3-V67-extremal''
for its label \texttt{thm:c3v67-extremal}.

\begin{theorem}[Cluster energy, band, localization and anchor, third cycle V2--V19]
\label{thm:c3v99-early}
On a window at the seed scale, an isolated cluster loses an explicit
fraction of its energy per window, nucleation takes at least two
windows, and the accounting of groups and seeds gives explicit
counts consistent with the class (Theorem~\ref{thm:c3v10-consolidated}).
The band cells over a window form a large-scale object at the
Gaussian radius (Theorem~\ref{thm:c3v8-consolidated}).  At a homogeneous
type-I point the far field carries neither the anti-alignment nor
the defect, both live on the configuration
(Theorem~\ref{thm:c3v14-consolidated}), and the cross term of the gate
lives on the groups within the Gaussian radius, on an anchor group
whose mean, momentum and dipole matrix are explicit
(Theorem~\ref{thm:c3v19-consolidated}); no signed row for the localized
push exists (Theorem~\ref{thm:c3v13-no-sign}).
\end{theorem}

\begin{theorem}[Explicit floor and the variance chapter, third cycle V21--V29]
\label{thm:c3v99-variance}
The variance of the normalized profile satisfies the one-level
criterion \(\mathrm{Var}\ge\frac1{16}(1-\hat a_{k_1})\) with \(k_1\le K_r=\lfloor2C_r-\frac12\rfloor\)
(Proposition~\ref{prop:c3v21-criteria}) and the explicit floor
\(\mathrm{Var}\ge v_{\rm exp}\) (Theorem~\ref{thm:c3v23-floor}), exponentially small in
the class constants, which closes the push budget for practical
purposes (Theorem~\ref{thm:c3v24-mean-push}); concentration gives no
floor (N-C3-1) and the mean spectrum no mean floor (N-C3-2); the mean
pressure gradient has an exact identity, the Gaussian momentum
decays on a nearly constant core, and the variance has a window
floor and a mean floor (Theorem~\ref{thm:c3v29-chapter}).
\end{theorem}

\begin{theorem}[Defect block and extremal ledger, third cycle V31--V39]
\label{thm:c3v99-defect}
The NRS identity holds with a defect, \(\nu\Delta\Pi-W\cdot\nabla\Pi=\nu|\Omega|^2+W\cdot Y\), whose
Gaussian integral is the ledger, but no quantitative NRS inequality
holds in the class (N-C3-3); the mean-square identity
\(\langle\|N\|^2\rangle=\langle\|L_-V\|^2\rangle+\langle\|Y\|^2\rangle\) gives the explicit constraint
\((\frac14+v_{\rm exp})c_\Phi+\underline\delta\le\kappa^2(\frac12\langle d\rangle+8\nu C)\) and, pointwise on windows,
\(\underline\delta\le10\kappa^2\nu C\) (Theorem~\ref{thm:c3v34-block}).  Every \(C^1\) observable
attains its extrema on the compact \(\Theta\)-invariant family with
vanishing derivative there (Lemma~\ref{lem:c3v34-extremal}), which gives
explicit rows for the Dirichlet form, the energy, the Rayleigh
quotient, the spectrum, the momentum and the defect
(Theorem~\ref{thm:c3v39-ledger}); Gate G1 is a statement about the angle
between the nonlinearity and the defect (Theorem~\ref{thm:c3v39-angle}).
\end{theorem}

\begin{theorem}[Angle and correlation chapters, third cycle V41--V49]
\label{thm:c3v99-angle}
The cross term \(\langle N,Y\rangle_G\) splits into a Lamb piece and a Bernoulli
piece with the explicit bound of Theorem~\ref{thm:c3v41-MNY}; on a nearly
constant core the cross term is pressure work; the extremal
Dirichlet form has a Lamb--Bernoulli form with Beltrami-type and
tangential reductions; the Bernoulli work is the derivative of the
flux up to a pairing (Theorem~\ref{thm:c3v44-chapter}).  The pressure work
on the configuration is an energy--Hessian pairing with a multipole
sign structure, every Gaussian pressure pairing is explicit at
supremum level (the pressure firewall in revised form), and the
defect rate has three explicit upper bounds
(Theorem~\ref{thm:c3v49-block}).  The bypass of V46 is not justified
(Firewall~\ref{fw:c3v47-no-bypass}).
\end{theorem}

\begin{theorem}[The anchor gates, third cycle V56--V69]
\label{thm:c3v99-anchor}
The radial defect of the anchor has explicit masses and its dipole
moment is the Gaussian mean of the defect; the pressure work on the
anchor-alone event is an explicit gate on \(\mathcal E_0\) with no sign
(Theorem~\ref{thm:c3v59-block}).  The cut-off integral of the Hessian is
\(-\frac13MI\), a nearly constant anchor does no pressure work, its cross
term is small, and the blow-up measure is not concentrated on nearly
constant anchors (Theorem~\ref{thm:c3v64-block}).  On spread
configurations the dissipation has an explicit floor; the extremal
Gaussian enstrophy row is pressure-free and closes to a constraint
between explicit constants; the gate on the spread anchor is a
question of two orientations of explicit order-one quantities, and
the columnar swirl and the plane wave are gate-neutral
(Theorem~\ref{thm:c3v69-block}).
\end{theorem}

\begin{theorem}[The vorticity ledger and the loops, third cycle V71--V79]
\label{thm:c3v99-vorticity}
The weighted enstrophy, the palinstrophy, the mean vorticity and
the vorticity dipole have pressure-free identities with pure
dampings \(\frac32,\frac32,1,\frac32\) (the first with the commutator \(3\nu\mathcal E\)), and
extremal rows (Theorems~\ref{thm:c3v74-block}, \ref{thm:c3v79-block}).  If
nonnegative observables have suprema \(s\le As\) with \(A\ge0\) and \(\rho(A)<1\),
all vanish (Lemma~\ref{lem:c3v76-matrix}); the loops so closed give
constraints between explicit constants wherever the palinstrophy
row applies, but their floors are capped at \(\max\{\kappa_1,1/\sqrt3\}\) and
\(3/\sqrt{80}\) (Firewalls~\ref{fw:c3v77-negative}, \ref{fw:c3v78-negative}), and
the palinstrophy row is void at every homogeneous type-I point
with \(\kappa<1\) (Proposition~\ref{prop:c3v93-consequences}).
\end{theorem}

\begin{theorem}[The explicit chapter, third cycle V81--V94]
\label{thm:c3v99-explicit}
Every element of the family satisfies the Duhamel formula on finite
windows and in ancient form (Theorems~\ref{thm:c3v81-duhamel},
\ref{thm:c3v81-ancient}).  With \(c_h=2/\sqrt\pi\), \(c_{h,2}\le1.794\), the Oseen constant
\(c_O\le2.836\) and its restriction to symmetric traceless inputs
\(c_O^{s0}\le2.766\) (certified quadrature, Theorem~\ref{thm:c3v89-certified};
hand bound \(4.208\), Theorem~\ref{thm:c3v82-certified}), the Leray
projector's annihilation of trace parts (Theorem~\ref{thm:c3v86-absorptions}),
the strain bound \(\lambda_{\max}(S)\le(2/3)^{1/2}|S|_F\) (Lemma~\ref{lem:c3v91-strain}) and
the transport form of the transport terms (Lemma~\ref{lem:c3v93-bounds}):
every homogeneous type-I point satisfies \((2/3)^{1/2}C_\infty'+\frac{\kappa^2}2\ge1\)
(Theorem~\ref{thm:c3v93-rows}), \(C_\infty'\le\gamma_1^{\rm red}(\kappa)\kappa\),
\(\nu^{1/2}C_\infty''\le\gamma_2^{\rm red}(\kappa)\kappa\), and the certified floors
\(\kappa\ge0.1410\) (velocity Liouville), \(\kappa\ge0.1815\) (enstrophy row), with
\(\sup\int|z|^2|\Omega|^2G\le12\nu C\) at the floor (Theorem~\ref{thm:c3v94-table};
constants sheet Theorem~\ref{thm:c3v94-sheet}).  None is an exclusion.
\end{theorem}

\begin{remark}[Corrections and negative records of the third cycle]
\label{rem:c3v99-corrections}
In-place corrections: Proposition~\ref{prop:c3v98-log}.  Negative
records: Theorems~\ref{thm:c3v97-numbered}, \ref{thm:c3v97-firewalls}.
Manifest: Section~\ref{sec:c3v96-chapters}, 196 labelled results.
\end{remark}

% =============================================================================
\section{Directions recorded for the fourth cycle}
\label{sec:c3v99-directions}
% =============================================================================

\begin{assessment}[Directions]
\label{ass:c3v99-directions}
\begin{enumerate}[label=\textup{(D\arabic*)},leftmargin=3.6em]
\item \emph{The initial term of the gradient bound through the ancient
      form.}  In Theorem~\ref{thm:c3v81-gradient} the initial term
      \(\nabla e^{\nu\theta L\Delta}W(s_0)\) is bounded by \(c_h\|W(s_0)\|_\infty(\nu\theta L)^{-1/2}\).  Writing
      \(W(s_0)\) by the ancient form \eqref{eq:c3v81-ancient} and putting the
      gradient on the Oseen kernel,
      \(\nabla e^{\nu\theta L\Delta}W(s_0)=-\int_{-\infty}^{s_0}\nabla e^{\nu(s-\tau)\Delta}\PP\operatorname{div}(W\otimes W)_0(\tau)\dd\tau\), whose
      sup norm is at most \((2/3)^{1/2}c_{O,2}\kappa^2\log\frac{1+\theta}\theta\) with
      \(c_{O,2}:=\sup_tt\,\|\nabla^2K^O_t\|_{L^1}\) (the second derivative of the Oseen
      kernel; \(\int_{-\infty}^{s_0}(s-\tau)^{-1}(-\tau)^{-1}\dd\tau=L^{-1}\log\frac{1+\theta}\theta\)): quadratic in
      \(\kappa\) and logarithmic in \(\theta\), against linear and \(\theta^{-1/2}\).  The
      gradient bound becomes
      \(C_\infty'\le\bigl((2/3)^{1/2}c_{O,2}\kappa^2\log\frac{1+\theta}\theta\bigr)/\bigl(1-2(8/3)^{1/2}c_O^{s0}\kappa(\theta/(1+\theta))^{1/2}\bigr)\),
      and the floor of \eqref{eq:c3v93-constraint} rises accordingly once
      \(c_{O,2}\) is computed by the method of V82, V87 and V89 (radial form
      of \(\nabla^2K^O_1\), exact pointwise norm, tail, Lipschitz constant,
      midpoint sum).  This is the first acquisition of the fourth
      cycle.
\item \emph{Information the sup-norm argument discards.}  The \(L^6\) bound
      on the far part of the kernel action and the Gaussian energy
      bound on \(\nabla W\) (Assessment~\ref{ass:c3v94-next}), which give
      floors in terms of \(K_6\) and \(C\), i.e.\ of the type-I enstrophy
      constant \(C_I\).
\item \emph{The gate in strain form.}  Gate G1 on the spread anchor
      (Theorem~\ref{thm:c3v68-gate}) with the explicit constants of the
      explicit chapter: the order-one quantities are now numerical
      bounds, and the two orientations are to be examined with them.
\item \emph{The loop method above \(\kappa=1\).}  The palinstrophy row applies
      only for \(\kappa\ge1\) (Proposition~\ref{prop:c3v93-consequences}), where
      the sup-norm absorption fails; a derivative bound there needs
      the class constants directly, and the loops of V74--V78 are to
      be read in that regime.
\item \emph{The two lifts.}  A pointwise use of the equation on a ball
      (to lift N-C3-1) and an explicit Liouville constant for (E9) in
      the class (to lift N-C3-3), Theorem~\ref{thm:c3v97-numbered}.
\end{enumerate}
\end{assessment}

\begin{proposition}[Item (AC2): what the fourth cycle must not redo]
\label{prop:c3v99-not-redo}
The 196 results of the manifest (Section~\ref{sec:c3v96-chapters}); the
kernel constants and their certification (Theorem~\ref{thm:c3v94-sheet});
the Hardy-route bounds of V67, V71, V72 (superseded by V91, V93, not
to be revived); the loop constraints below \(\kappa=1\) (void); the routes
closed by N-C3-1, N-C3-2, N-C3-3 and by the standing firewalls
(Theorem~\ref{thm:c3v97-firewalls}) unless one of the named lifts is
supplied; the bypass of V46.
\end{proposition}

\begin{proof}
The cited records.
\end{proof}

% =============================================================================
\section{Integrated V99 conclusion and next acquisition order}
\label{sec:c3v99-conclusion}
% =============================================================================

\begin{theorem}[What V99 establishes]
\label{thm:c3v99-conclusion}
Relative to the two legacy cycles and V1--V98: the audited summary of
the third cycle (Theorems~\ref{thm:c3v99-early} to
\ref{thm:c3v99-explicit}), the directions
(Assessment~\ref{ass:c3v99-directions}) and the not-to-redo list
(Proposition~\ref{prop:c3v99-not-redo}).  The full Navier--Stokes
stability/global-regularity theorem remains unproved.
\end{theorem}

\begin{proof}
The cited records.
\end{proof}

\begin{programme}[V100 acquisition order]
\label{prog:c3v100}
\begin{enumerate}[label=\textup{(AC\arabic*)},leftmargin=4.8em]
\item The compulsory companion audit: read the current SM and YM
      notes and record their digests and decision.
\item The closure of the third cycle: the position after one
      hundred versions, the freezing of the file as the third legacy
      file, and the opening of the fourth cycle from a new V1 carrying
      the summary of V99; delivery.
\end{enumerate}
\end{programme}

\begin{assessment}[Position after V99]
\label{ass:c3v99-position}
The fourth cycle's opening summary and directions are drafted.  No
full stability or global-regularity claim is made.
\end{assessment}

% =============================================================================
\section*{Version V99 (third cycle) append-only record}
\addcontentsline{toc}{section}{Version V99 (third cycle) append-only record}
% =============================================================================

The complete V98 source of the third cycle was retained before this
continuation.  Its SHA--256 digest was
\begin{center}\scriptsize\ttfamily
edd968b51760e8825d30ff236e48d3d3615d508a827ef57ff4df26f0d9079c2e.
\end{center}
V99 drafted the fourth cycle's opening summary and directions.  No
full regularity claim is made.

% =============================================================================
\section{V100: compulsory companion audit and the closure of the third cycle}
\label{sec:c3v100-entry}
% =============================================================================

This is the hundredth and last version of the third cycle.  The
compulsory review of the two companion programmes is performed
first; the position after one hundred versions is recorded; the
file is then frozen as the third legacy file, and the fourth cycle
opens from a new V1 carrying the summary of V99.  No global
regularity claim is made.

% =============================================================================
\section{Compulsory V100 review of the current companion notes}
\label{sec:c3v100-companion-audit}
% =============================================================================

The current companion files in \texttt{latex\_notes} were inspected
read-only.  Both programmes have moved since the V95 audit: the SM
programme from V50 to V51 of its seventh series, the YM programme
from V27 to V28 of its sixth series.  The frozen checkpoints are
\begin{align}
 &\texttt{Renewal\_SM\_Einstein\_Continuum\_Research\_Notes\_V51.tex},
 \notag\\[-2pt]
 &\qquad
 \texttt{699c607e81d2a1837079d255e56e5d7c0f76e4bee93975197d5bcfc792a01d05},
 \label{eq:c3v100-SM-checkpoint}\\
 &\texttt{Renewal\_Yang\_Mills\_Mass\_Gap\_Research\_Notes\_V28.tex},
 \notag\\[-2pt]
 &\qquad
 \texttt{8ab8b98fe192a04244e048a8022709260e6685b7e34c82892634b9da02d1b60b}.
 \label{eq:c3v100-YM-checkpoint}
\end{align}
(SM V51 has 20468 lines; YM V28 has 12565 lines.  The folder also
held YM V27, the state recorded at V95, superseded.)

\emph{SM V51.}  The SM programme wrote an invariant affine-jet
completion: why its hub must retain the physical coframe, the
affine transverse subspace, the Hermitian completion of the
transverse affine modes, a tangential Fourier-jet completion, row
moments, covariance and the gap, and the remaining heat-flat
derivative problem.

\emph{YM V28.}  The YM programme treated a reflection-positive
suppression: a distance lemma and a local no-go, a Haar-ball lemma
and a lower bound on its partition function, reflection positivity
and a chessboard estimate, same-colour and all-colour theorems, a
link form of its multi-defect suppression card with an explicit
rate, and a cluster corollary.

\begin{theorem}[V100 companion-audit decision]
\label{thm:c3v100-companion-decision}
Nothing is imported from SM V51 or YM V28.  Neither bears on the
third cycle's chapters or on the fourth cycle's first acquisition.
\end{theorem}

\begin{proof}
The SM results concern jet completions and gaps of lattice Dirac
interfaces; the YM results concern reflection positivity and defect
suppression in a lattice gauge theory.  None is a Navier--Stokes
statement, and none concerns the mild formulation, the Oseen
kernels, the vorticity rows or the explicit floors.
\end{proof}

% =============================================================================
\section{Position after one hundred versions, and the closure}
\label{sec:c3v100-position}
% =============================================================================

\begin{assessment}[Position at the close of the third cycle]
\label{ass:c3v100-position-hundred}
The third cycle wrote eighteen chapters and 196 labelled results
(Section~\ref{sec:c3v96-chapters}), three numbered negative records and
the standing firewalls (V97), a constants sheet and a corrections log
(V98), and the audited summary that opens the fourth cycle (V99).
Its explicit result: every homogeneous type-I point of a smooth
periodic solution has rescaled velocity bound \(C_\infty\ge0.1815\,\nu^{1/2}\),
strain constant at least \((3/2)^{1/2}(1-\kappa^2/2)\), and Gaussian second
vorticity moment at most \(12\nu C\) at the floor, with the certification
chain of Theorem~\ref{thm:c3v94-sheet}; its structural results: the
extremal-configuration method and the vorticity ledger's
pressure-free rows, the localization of the gate's cross term on the
anchor, and the exclusion of the nearly constant anchor as its
carrier.  Gate G1 is unproved; the singular point, if it exists, is
described, not excluded.  The fourth cycle begins with direction (D1)
of Assessment~\ref{ass:c3v99-directions}.  No global regularity claim
is made.
\end{assessment}

\begin{theorem}[What V100 establishes]
\label{thm:c3v100-conclusion}
Relative to the two legacy cycles and V1--V99: the companion-audit
decision (Theorem~\ref{thm:c3v100-companion-decision}) and the closure.
The full Navier--Stokes stability/global-regularity theorem remains
unproved.
\end{theorem}

\begin{proof}
The cited records.
\end{proof}

\begin{programme}[Closure and the next acquisition order]
\label{prog:c3v100-closure}
\begin{enumerate}[label=\textup{(AC\arabic*)},leftmargin=4.8em]
\item This file is retained as
      \texttt{Renewal\_NS\_Stability\_Research\_Notes\_V100\_f3.tex}; its
      digest is recorded in the fourth cycle's V1.
\item The fourth cycle's V1 opens with the audited summaries of the
      three cycles, the imports, the gates, and the first acquisition
      (Programme C4-V1-first: the second Oseen constant \(c_{O,2}\) and the
      gradient bound with the initial term through the ancient form).
\end{enumerate}
\end{programme}

% =============================================================================
\section*{Version V100 (third cycle) append-only record}
\addcontentsline{toc}{section}{Version V100 (third cycle) append-only record}
% =============================================================================

The complete V99 source of the third cycle was retained before this
continuation.  Its SHA--256 digest was
\begin{center}\scriptsize\ttfamily
6e99e6fc646fddf539147f1510fb47269961db66ec7a41b46a36dc92639cb2ad.
\end{center}
V100 performed the compulsory companion audit and closed the third
cycle.  No full regularity claim is made.

\end{document}
